% Framework v3.069: verified cubic five-core operator intertwiner; mathematical comparison only.
% Framework source amendment v3.068: scoped geometry/operator additions, physical occurrence unchanged.
\documentclass[11pt]{article}
\usepackage{graphicx}
\usepackage{multirow}
\usepackage{tikz}
\usepackage{pgfplots}
\usepackage{tikz-3dplot}
\usepackage{xcolor}
\usepackage{mathtools}

\DeclareMathOperator{\Hom}{Hom}


\pgfplotsset{compat=1.18}
\usetikzlibrary{3d,calc,arrows.meta,positioning,decorations.markings}


\def\BHMathOnly{}
% Macros-v2_203: v3.077 source-current bindings for unconditionality source integration. No mathematical macro changes.
%% Macros-v2_201: downstream full-carrier geometry update, five current root bindings, no mathematical macro change.
%% Macros-v2_200: Paper III-F v1.280 current-source binding for framework v3.074; mathematical macro definitions unchanged.
%% Macros-v2_198: current project v3.072 III-D and bibliography binding; mathematical macros unchanged.
%% Current project v3.070 source-reference bindings; previous macro bodies unchanged.
% Current II-D reference root for PapersMinimal v3.063; mathematical macros otherwise unchanged.
% Macros-v2_188: re-based current paper identities and A--G family; no change to geometry.
% Macros-v2_184: active Paper III-E v1.200/III-F v1.267 binding for framework v3.031; mathematical macros unchanged.
% Macros-v2_183: updated shared cross-reference bindings for Paper1B/III-E/III-F v3.030; mathematical macros unchanged.
% Macros-v2_182: active Paper II-F/V-A cross-references for framework v3.029; definitions unchanged.
% Paper-series v2.381: new I-B torsion/modular results, bibliography and current source roots.
% Paper-series v2.380: new exceptional bibliography and synchronized paper references.
% Paper-series v2.379 current 31-paper macro/aux registry.
% Paper-series v2.378: current 31-paper AUX/bibliography registry updated for Paper I-B v1.2 literature additions.
% Framework v2.371: current-source reference registry refreshed; old theorem labels retained.
% Framework v2.370: current 30-paper source/AUX root rebinding; native mathematical macros unchanged.
% Framework v2.368: synchronized merged 30-root registry; mathematical notation unchanged.
% Framework v2.358: active 30-paper source registry synchronized; existing mathematical macros unchanged.
% Framework v2.357: selective new source-owned certificates; active reference roots updated for five revised papers.
% Framework v2.356: active 30-root registry synchronized to the 30 papers in this release.
% Framework v2.275: exact positive dense-twin word cubature, finite repeatability and eighth-moment classification, nested sampling frames, joint action and POVM-jet exports; preserves v2.274.
% Framework v2.274: filename and root-source synchronization.
% Framework v2.273: restored missing Problem 3 theorems from v2.268; exact exceptional frames, twin geometry, convolution repeatability and projector/occurrence firewalls; merged on v2.272 without overwriting independent branches.
% Framework v2.272: merged v2.270 covariant native-readout mathematics with v2.271 Problem 4 finite-action Hopf hypercharge/holonomy corrections; no universal bulk-gauge uniqueness claim.
% Independent branches: v2.270 readout/guard exports and v2.270 K16 action-consumed charge closure; all prior scoped assumptions retained.
% Framework v2.269: synchronized 30 active AUX roots; mathematics and bibliography preserved from Macros-v2_161.
% Framework v2.266: Paper II-C converse-derived outcome resources; II-E, III-C, IV-A/B, V-A and VI-B downstream mathematical closures; 30-root registry synchronized.
% Framework v2.265: II-C coordinate-stage constitution and exact first-symbol host inheritance; prior mathematical domains retained.
% Framework v2.264: paper II-B reverse/audit and IV-A observable-minimal exchange; synchronized external references to 30 active roots.
% Framework v2.263: Paper I-C adds bridge-controlled root-carry parity and source-relative old-cap cancellation; other root mathematical bodies retained; active 30-root references synchronized.
% Framework v2.262: integrates Paper II-F radius, Casimir and readout results; active 30-root references synchronized.
% Framework v2.261: merged Paper II-B GP2-level/geometric repairs; preserves v2.260 II-D repairs; active references synchronized.
% Framework v2.260: Paper II-D signed RS, finite-symbol typing, and conditional curvature/Jordan repairs; references synchronized.
% Framework v2.259: conditional native SU(2) atlas and explicit later discharge; v2.258 mathematical bodies retained elsewhere.
% Framework v2.258: preserves parallel II-A/dynamics update and adds complete readout/frequency closures.
% Framework v2.257: exact minimal complete position readout and arbitrary-preparation frequency law.
% Framework v2.256: exact native inertia and sharp bounded ordered-source closure; preserves concurrent v2.255 work.
% Framework v2.255: preserves parallel detector integration and adds exact finite readout closures.
% Framework v2.254: Pasterski bibliography and synchronized 30-paper external-reference registry.
% Framework v2.253: typed raw coordinate defect in I-A and synchronized VI-A import; active references synchronized.
% Framework v2.252: synchronized source registry; optional III-G shell-frame calculation.
% Framework v2.251: forward common-action execution, joint source jets and scoped coupled control.
% Framework v2.251: exact full-heat variations, fermion source jets and explicit physical completion frontier.
% Framework v2.251: certified interacting source jets, pole stability and scoped completion proofs.
% Framework v2.251: synchronized active source registry; mathematical notation unchanged.
% Framework v2.251: synchronized active source registry; mathematical notation unchanged.
% Framework v2.251: synchronized active source registry; mathematical notation unchanged.
% Framework v2.251: synchronized source registry; no mathematical notation changes.
% Framework v2.244: synchronized 30-paper AUX registry and rooted phase-source action-order tests.
% Framework v2.242: signed rooted phase lift, complete transverse return and exact family cubic.
% Framework v2.241: proved rooted conormal geometry, conditional cap realization and scoped action tests.
% Framework v2.239: all 30 active source roots synchronized; no mathematical macro changes.
% Framework v2.238: Paper IV A/B/C split and 30-root reference registry.
% Framework v2.237: Paper VII-D technical split and 28-root reference registry.
% Framework v2.236: Paper VI-C technical-volume reference added; all current root bindings advanced.
% Framework v2.235: 26-root AUX and bibliography bindings synchronized; mathematical macros unchanged.
% BEGIN INLINE SOURCE: Paper00-TeXBuildCompatibility-v1_0.tex
% Rendering-only repairs for inherited long appendices and bitmap text fonts.
% Integer counters through 26 retain the original labels; later values use AA,
% AB, ... . This file changes no action coefficient, state or theorem premise.
\RequirePackage{alphalph}
\makeatletter
\let\BHPrimitiveAlph\@Alph
\newalphalph{\BHExtendedAlph}{\BHPrimitiveAlph}{26}
\let\@Alph\BHExtendedAlph
\AtBeginDocument{%
  \@ifpackageloaded{microtype}{\microtypesetup{expansion=false}}{}%
}
\makeatother

% END INLINE SOURCE: Paper00-TeXBuildCompatibility-v1_0.tex

% Framework v2.164: current reference/bibliography profile; mathematical definitions retained.
% Framework v2.162: finite metric, functional Hilbert and scoped Ward integration.
% Framework v2.161: inherits Macros-v2_57.tex exactly except the current external-reference binding.
% Macros v2.57 / framework v2.156: current reference registry only; mathematical definitions unchanged.
% Framework v2.143 / Macros v2.55: current external-reference bindings; mathematical notation unchanged.
% Framework v2.110 / Macros v2.54: current external-reference registry and bibliography; mathematical notation inherited.
% Framework v2.67 / Macros-v2_53: adds Born-action-cell and joint-EH readout symbols; absolute unit realization is metrological; MATCH_{J/scal}^{EH} remains the sole G24 normalization bridge.
% Framework v2.35 / Macros-v2_52: grouped paper numbering, active AUX registry, and Born-v6_20 bibliography binding synchronized; notation unchanged.
% Framework v2.26: active-source AUX and bibliography bindings synchronized; notation unchanged.
% Framework v2.09: reviewed early-reciprocity integration; scoped physical imports.
% Framework v2.08: Born--Euler mathematics and conditional complete-action potential criterion.
% Framework v2.07: active AUX registry includes Paper I-B v1.2; shared bibliography Born-v6_18.
% Macros-v2_46 / framework v2.04: mathematical atlas and downstream realization correction.
% Macros-v2_46 / framework v2.03: scoped structure integration and current dependency bindings.
% Macros-v2_46 / framework v2.01: shared source bindings; prior revision comments are historical.
% v2.42 / framework v1.99: Synchronized RG/AS owner and bibliography bindings.
% v2.41 / framework v1.99: active manuscript AUX bindings and geometric-closure bibliography synchronized.
% Release-v1.80 centralized source synchronization for scalar-circle/comparison-band integration.
% Release v1.79: Hurwitz hierarchy and kinetic-normalized action response integration; physical mass matching remains conditional.
% Release v1.78: approved M8 closure-batch integration; physical and analytic domains retained.
% Current v1.78 registry: all active TeX roots are synchronized below.
% Earlier release comments are historical scope records, not current status.
% Release-v1.77 synchronization: native same-owner Noether closure, generator--flow/history typing, analytic-certification firewalls, scalaron/K16, LF/QCD, quantum-Ward, and RAR boundary-owner transport are registered; shared mathematical notation is unchanged.
% Release-v1.76 synchronization: registry points to the v1.76 active paper roots.
% Release-v1.76 synchronization: registers the local A+Phi+g Noether/constraint closure and dynamics-frontier propagation; shared mathematical notation is unchanged.
% Release-v1.75 synchronization: registry points to the v1.75 active paper roots.
% Release-v1.75 synchronization: integrates the joint-Dirac Einstein bulk real form and the local odd-scalar heat-coefficient action on explicit analytic first-jet/interior domains; v1.74 Yang--Mills closure is retained. Fermionic reality, pure-metric curvature-squared/exact-nonlocal terms, noncanonical boundaries, physical vacuum selection, and actual solution-history tokens remain open.
% Release-v1.74 synchronization: registry points to the v1.74 active paper roots and retains the v1.73 theorem labels.
% Release-v1.74 synchronization: canonical Macros-v2_34.tex; generated local nonlinear Yang--Mills action, constrained connection-field phase space, source-free Gauss momentum map, and Legendre closure are integrated on their proved domains. The remaining non-gauge ACT_cl sectors and actual solution-history selection remain open.
% Release-v1.71 synchronization: closes the former JOINT_x/JOINT_5
% connection-occurrence gap by canonical structure-group extension of the
% nonflat Born connection, faithful associated representation, and the local
% cotangent--Clifford operator constructor.  Actual fluctuation members,
% actions, vacua, solutions, calibrations, and analytic QFT domains remain
% separately typed.  The active-AUX registry is advanced once.
% Release-v1.71 canonical macro and active-AUX registry.
% ======================================================================
% Macros-v2_39.tex -- canonical publication macro file
% Release v1.71 retains the v1.70 null-ray/QCD boundaries and additionally
% removes the duplicate structural connection actuality premise.  It does not
% claim global bundle topology, select a gauge/Higgs member, or promote the
% finite construction to an interacting continuum theory.
% Release v1.68 constructs the finite light-front total-fraction operators and
% parton number/momentum operator measures, proves their simplex and Hausdorff-limit
% theorems, and separates positive cut measures from scheme-renormalized
% light-ray distributions.  It also separates finite real Schur poles, stable
% norm-resolvent/Riesz pole-residue limits, and continued resonance
% pole/residue limits.  The Paper-VII-C scattering package is claim-indexed:
% local QFT and spectral claims import no S-matrix or infrared premise, while
% asymptotic amplitudes, LSZ, S-matrix poles, completeness, and cross sections
% import only their named massive or massless/long-range clauses.
% Release v1.67 factorizes BND48_H into a same-action Hessian/Schur
% descendant and still-open typed boundary-trace, critical-lift, and
% refinement-limit conditions.  It proves the finite-collar Dirichlet-to-
% Neumann law sqrt(y)tanh(L sqrt(y)), its nonuniform critical limit, and the
% finite-return ultra-infrared turnover.  The Paper-I-A Hardy constant mode is
% an exact geometric near-hit, not yet the physical metric boundary mode.
% The current local f(R) Hessian remains massive and therefore still refutes
% a RAR prediction on the minimal action branch.
% Release v1.66 derives the transverse depth-six normalizer bridge from the
% already generated C6 Reynolds image and Paper-I-A binary-dihedral normalizer.
% The bridge is the unique least 2D3-stable scalar-containing realization of
% the faithful quotient 2D3/C6=C2, has trace defect six and hence K=48, but is
% transverse to the nested active carrier.  Exact K=44 and K=52 contraction,
% semigroup, and energy-midpoint identities are registered as structural
% readouts rather than projector events.  The active E48 carrier gap remains
% unchanged.  A same-action boundary-symbol/Calderon intertwiner is isolated
% as the necessary and still unproved owner bridge to the static DtN sort;
% representation equivalence alone neither derives RAR nor creates a K=48
% action junction, pole, occupation, or cosmological transition.
% Release v1.65 derives the native E_48-centered carrier mass gap from the
% GP1--GP3-generated active spectrum, distinguishes its carrier-free interval
% from a particle located at E_48, and withdraws the unsupported promotion to
% a positive continuum Yang--Mills mass gap.  The latter remains a valid
% conditional transfer theorem on an enlarged continuum branch, but it is not
% an active framework prediction or an unfinished native-gap proof.
% Release v1.64 synchronizes Derivation Batches 14--22.  It registers the
% GP1--GP3 heat-profile orbit, same-Hessian electroweak response and pole
% ownership, K16/K48 action-response corrections, perturbation state and
% non-Gaussianity typing, finite-vacuum/gap/QCD limit repairs, factorized
% position-area ownership, the actual static metric response, and the ordered
% Wilson/string-tension limits.  Strict continuum convergence, positive
% physical gap, surviving higher cumulants, and positive tension remain open.
% Release v1.61 closes TIMEWIDTH16 as an independent direct-production datum.
% JUNC_{2|7} continues to own the finite intrinsic microscopic profile/width,
% while the actual same-action background/production history determines its
% physical-time pullback, support, and duration.  The current Paper-I-A--VII-C AUX
% registry is advanced once; shared mathematical notation is unchanged.
% Release v1.60 derives the final-octet primitive-remnant depletion, canonical
% Hopf-current representative, exact Hopf descent to dOmega/(4pi), and
% closed-FLRW a^{-3} neutral-current shape.  It retypes K=48 as the onset
% of the final one-unit geometric transfer and records cDM/RAR as distinct
% physical projections of that common neutral parent while preserving
% occurrence, Hilbert-owner, stress-matching, and static-kernel premises.
% Release v1.59 transports the v1.58 exceptional-shell/K16 closure through
% Papers III-F and VI-B--VII-C.  It replaces Paper VI-B's obsolete mixed-Hessian
% output typing by the third-variation pair vertex and Bogoliubov occupation
% support, derives the post-seed Higgs-pair feedback vertex and its exact
% alpha--delta kernel factors, propagates the unique K16/K48 action-junction
% split, identifies the QCD delta count with the upstream 2O shell atom only
% as a common-scalar numerical compatibility, and narrows THERM16 to dynamical
% approach/beta selection because the finite Gibbs/KMS target is already
% unique on a supplied simple observable algebra once its Hamiltonian and
% inverse temperature are supplied; direct-sum central weights remain typed.
% Release v1.58 registers the exceptional nearest-shell closure, the derived
% transverse K=16 action event, the unique interior action junction, the
% type-correct ordered pair vertex, the thin-layer squeeze law, and the
% radiation-like/matter-like kinematic split.  REG_0, transfer magnitude,
% physical layer width, and thermalization remain explicitly conditional.
% Release v1.57 registers the active-depth-six obstruction, the native
% K=48 carrier desert, the typed 7|9 owner reduction, and the dynamical
% reheating handoff across Papers III-B--VI-B.
% Release v1.56 registers the least-support n_2 inflation dressing of the
% persistent n_7 Laurent response, the analytic global-convexity and positive
% endpoint-measure proofs, the derived inflation core and K=16 coefficient
% freeze, and the no-particle dark-sector split into static RAR plus the
% conditional homogeneous effective-curvature response at reflected K=48.
% There is no n_6 action, scalaron-DM mass, or second Einstein--Hilbert term.
% Release v1.55 registers the protected-octet resolved/residual reflection
% q_sp+Z_W=8, its j<->8-j and K<->64-K scale covariance, and the derived
% rank-one-exit quadratic grade C_1^2/U=C_2 in Papers III-C, III-E, and VIII.
% The reflection is a typed trace/scale theorem, not a time reversal or a
% physical event-occurrence theorem; reheating transfer remains conditional.
% Release v1.54 registers the exact positive-part Laurent indices, the
% active-domain/terminal-endpoint action split, the no-second-EH firewall, and
% the q_DM=6/K_DM=48 scalaron readout typing across Papers III-D/E/F, VI,
% VIII, and IX.
% Release v1.53 registers the common nonnegative Laurent obstacle projection
% for n_1^+ and n_7^+, discharges the three former exit/freeze assumptions,
% and advances the Paper-III-D/E/F, VIII, and IX consumers.  Shared notation
% and bibliography ownership are unchanged.
% Release v1.51 registers the PQTI derivation batch and its updated
% Paper-III-A/C/D/E/F, VIII, and IX consumers.  Shared notation and
% bibliography ownership are unchanged.
% Release v1.50 advances the Paper-III-B AUX target to v1.40 after the
% rigorous audit repair; shared notation and bibliography ownership are unchanged.
% Release v1.49 advances the Paper-II-E/F AUX target to v3.14 after the
% minimal audit repair; shared notation and bibliography ownership are unchanged.
% Release v1.46 synchronizes Paper II-A's minimal frozen constructor-origin
% theorem and its five dependent AUX targets.  Shared notation is unchanged.
% Release v1.45 updates the Paper-II-B AUX target to v2.23 after the
% GP3-rigidity and appendix-placement repair; shared notation is unchanged.
% Release v1.44 advances the centralized Paper-I-A--VII-C AUX registry after the
% Paper-I-A v5.16 protected-completion scope repair.  Every current
% paper imports this one file; no paper-level opt-in or \externaldocument list
% is permitted.  The bibliography remains centralized here through Born-v6_13.
% Shared bibliography, theorem infrastructure, colors, notation alignment,
% and current Paper-I-A--VII-C external-reference registration.
% This file centralizes common imports, bibliography sorting, colors, theorem
% infrastructure, and the notation shared by the complete series.
% v2.9 (release v1.43) made the Paper-I-A--VII-C AUX registry automatic and
% centralized the bibliography and all external-reference ownership.
% v2.8 release-registry refresh (v1.33) registers the rooted generated-parent
% cascade and the POL/REG/free-law/11+4/Born-group closure revisions.  No
% shared mathematical notation or Paper-II-A fundamental principle is changed.
% v2.8 release-registry refresh (v1.32) registers the universal generated-
% continuation revisions in Papers II-A, III-A, III-C, IV, and VIII.  No
% shared mathematical notation or Paper-II-A fundamental principle is changed.
% v2.8 release-registry refresh (v1.30) registers the CRP4/projective-state
% source revisions and their direct registry-only consumers.  No shared
% mathematical notation or Paper-II-A axiom is changed.
% v2.7 release-registry refresh (v1.27) updates only the current external
% Paper-I-A--VII-C basenames after the placement/assumption synchronization; no
% shared mathematical notation or Paper-II-A axiom is changed.
% v2.7 synchronizes the geometric defect-repair revision: Paper-II-A v3.42,
% Paper-III-C v1.26, Paper-III-D v1.27, Paper-III-E v1.32, and Paper-III-F
% v1.34.  It registers the common trace-defect/spectral-filtration theorem,
% its owner-typed response pullback, and the reduced GRED and TERM interfaces.
% No shared mathematical notation is changed.
% v2.6 synchronizes Paper-III-E v1.31's cone-type repair, minimal-response
% classification, and explicit classical action-role boundary, together with
% Paper-III-F v1.32's downstream response-atlas synchronization.  No shared
% mathematical notation is changed.
% v2.5 synchronizes the Paper-III-D v1.26 faithful-trace-duality proof,
% exact ordered-owner joint locus, and typed alternative-presentation
% bookkeeping repair, together with its Paper-III-E and closure-frontier
% references.  No shared mathematical notation is changed.
% v2.4 synchronizes the Paper-III-C v1.25 orbit/process/owner repair, its
% downstream Papers III-B--F, V, VI, VIII, and IX revisions, and Born-v6_13.
% v2.3 synchronizes the Paper-II-A v3.41 witness/support/principal inheritance
% layer, the Paper-III-A v1.24 joint-realizability/readout firewall, all dependent
% host-rate revisions, and the current Paper-I-A--VII-B source registry.
% v2.2 synchronizes the generated active-realization, determinant-depth,
% dimension-weighted hierarchy, and typed cosmological-event interfaces.
% v2.1 synchronizes the exceptional-support interface revisions: Paper I-A
% v5.9, Paper II-A v3.30, Paper II-E/F v3.7, and Paper IV v4.16.
% v2.0 introduces the stage-safe G0 occurrence notation, synchronizes the
% active Paper-I-A--IV source registry (current rigorous Paper-III closure:
% III-A v1.17, III-B v1.22, III-C--E v1.17, III-F v1.14, IV v4.14), and makes
% external-reference self-skip depend on logical paper role rather than a
% versioned job name.
% ======================================================================

% ======================================================================
% macros.tex — Shared macros for the Born Geometry series
% ======================================================================

% ======================================================================
% 1. PACKAGES, BIBLIOGRAPHY, COLORS, AND UNIT DEFINITIONS
% ======================================================================
% Core packages shared across the complete Paper-I-A--VII-C series.  Loading them here makes
% the macro file a genuine series preamble while remaining compatible with existing
% paper preambles that load the same packages before \input.
\usepackage{amsmath}
\usepackage{amssymb}
\usepackage{amsthm}
\usepackage{mathrsfs}
\usepackage{array}
\usepackage{booktabs}
\usepackage{xcolor}
\makeatletter
\@ifpackageloaded{geometry}{}{\usepackage{geometry}\geometry{margin=1in}}
\makeatother
\usepackage{mdframed}

% Shared support packages required by the macros below and by most papers.
\usepackage{ragged2e}
\usepackage{longtable}
\usepackage{tabularx}
\usepackage{enumitem}
\usepackage[most]{tcolorbox}

% Series-wide numeric bibliography convention using the portable BibTeX path.
% All Papers I-A--VII-C use the ordinary \cite interface.
\usepackage[numbers,sort&compress]{natbib}
\providecommand{\printbibliography}{%
  \bibliographystyle{plainnat}%
  \bibliography{Born-v6_60}%
}

% Shared colors.  \providecolor avoids clobbering a paper-specific override.
\providecolor{postulatecolor}{RGB}{230,242,255}
\providecolor{compare}{RGB}{230,245,230}
\providecolor{definitioncolor}{RGB}{230,242,255}

\providecolor{openriskcolor}{RGB}{255,235,235}
\providecolor{bridgecolor}{RGB}{240,240,255}
\providecolor{gold}{RGB}{255,215,0}

% ======================================================================
% 2. CUSTOM ENVIRONMENTS
% ======================================================================
% Common table column types used in all papers.  Guarded so that a paper may
% define its own versions before loading this file.
\makeatletter
\@ifundefined{NC@find@L}{%
  \newcolumntype{L}[1]{>{\RaggedRight\arraybackslash}p{#1}}%
}{}
\@ifundefined{NC@find@C}{%
  \newcolumntype{C}[1]{>{\Centering\arraybackslash}p{#1}}%
}{}
\@ifundefined{NC@find@Y}{%
  \newcolumntype{Y}{>{\RaggedRight\arraybackslash}X}%
}{}
\makeatother

% ======================================================================
% 3. THEOREM STYLES AND ENVIRONMENTS
% ======================================================================
% Custom underlined theorem style (useful for remarks or notes)
\newtheoremstyle{underlined}
  {3pt}   % Space above
  {3pt}   % Space below
  {}      % Body font (empty means normal/roman text, typical for remarks)
  {}      % Indent amount
  {}      % Theorem head font (can be \itshape, \bfseries, etc.)
  {.}     % Punctuation after theorem head
  { }     % Space after theorem head
  {\underline{\thmname{#1}\thmnumber{ #2}}\thmnote{ (#3)}} % Theorem head formatting

% Standard plain theorem style environments.  Each environment is guarded so
% the shared macro file can be loaded after a paper-specific theorem setup.
\theoremstyle{plain}
\makeatletter
\@ifundefined{theorem}{\newtheorem{theorem}{Theorem}[section]}{}
\@ifundefined{lemma}{\newtheorem{lemma}{Lemma}[section]}{}
\@ifundefined{proposition}{\newtheorem{proposition}{Proposition}[section]}{}
\@ifundefined{corollary}{\newtheorem{corollary}{Corollary}[section]}{}
\makeatother

% Definition style environments (typically non-italicized body)
\theoremstyle{definition}
\makeatletter
\@ifundefined{definition}{\newtheorem{definition}{Definition}[section]}{}
\@ifundefined{postulate}{\newtheorem{postulate}{Postulate}[section]}{}
\@ifundefined{closure}{\newtheorem{closure}{Closure}[section]}{}
\makeatother

% Applying the custom underlined style
\theoremstyle{underlined}
\makeatletter
\@ifundefined{remark}{\newtheorem{remark}{Remark}[section]}{}
\@ifundefined{convention}{\newtheorem{convention}{Convention}[section]}{}
\@ifundefined{principle}{\newtheorem{principle}{Principle}[section]}{}
\makeatother

% Framed proof environment.
% Every \begin{proof} ... \end{proof} is automatically framed/shaded.
% This avoids mdframed/trivlist nesting, which can drop the top border after
% theorem-like environments.
\makeatletter
\renewenvironment{proof}[1][\proofname]{%
  \par\addvspace{6pt}%
  \noindent
  \begin{tcolorbox}[%
    enhanced,
    breakable,
    colback=black!4,
    colframe=black!35,
    boxrule=0.35pt,
    arc=1pt,
    left=6pt,
    right=6pt,
    top=4pt,
    bottom=4pt,
    boxsep=0pt,
    before skip=0pt,
    after skip=6pt,
    parbox=false
  ]%
  \pushQED{\qed}%
  \normalfont\itshape #1\@addpunct{.}\enspace
  \normalfont\ignorespaces
}{%
  \unskip\nobreak\hfill\popQED
  \end{tcolorbox}%
  \par\addvspace{2pt}%
}
\makeatother

\mdfdefinestyle{postulatebox}{
  backgroundcolor=yellow!10,
  linecolor=black,
  linewidth=0.4pt,
  skipabove=\baselineskip,
  skipbelow=\baselineskip,
  nobreak=true
}

% ======================================================================
% PROOF-READABILITY HELPERS
% ======================================================================

\newcommand{\ProofPlan}[1]{%
  \par\smallskip
  \noindent\textbf{Proof plan. }#1%
  \par\smallskip
}

\newlist{proofsteps}{enumerate}{2}
\setlist[proofsteps,1]{
  label=\textbf{(\alph*)},
  ref=\alph*,
  leftmargin=2.2em,
  itemsep=0.45\baselineskip,
  topsep=0.35\baselineskip
}
\setlist[proofsteps,2]{
  label=\textbf{(\roman*)},
  ref=\roman*,
  leftmargin=2.0em,
  itemsep=0.25\baselineskip,
  topsep=0.2\baselineskip
}


% ======================================================================
% 4. NUMBER SYSTEMS AND ALPHABETS
% ======================================================================
\providecommand{\RR}{\mathbb{R}} % Real numbers
\providecommand{\PP}{\mathbb{P}} % Projective space / Probability
\providecommand{\CC}{\mathbb{C}} % Complex numbers
\providecommand{\ZZ}{\mathbb{Z}} % Integers
\providecommand{\NN}{\mathbb{N}} % Natural numbers
\providecommand{\HH}{\mathbb{H}} % Quaternions / Hyperbolic space
\providecommand{\LFRayGeom}{\mathsf{LF\!\mbox{-}RAY\!\mbox{-}GEOM}}
\providecommand{\LFRep}{\mathsf{LF\!\mbox{-}REP}}
\providecommand{\LFState}{\mathsf{LF\!\mbox{-}STATE}}
\providecommand{\LFLim}{\mathsf{LF\!\mbox{-}LIM}}
\providecommand{\LFScat}{\mathsf{LF\!\mbox{-}SCAT}}
% Number systems / projective spaces
\providecommand{\CP}{\mathbb{CP}}
\providecommand{\RP}{\mathbb{RP}}

% ======================================================================
% 5. GROUPS AND LIE THEORY
% ======================================================================
\providecommand{\SO}{\ensuremath{\mathrm{SO}}}     % Special Orthogonal group
\providecommand{\SU}{\ensuremath{\mathrm{SU}}}     % Special Unitary group
\providecommand{\Spin}{\ensuremath{\mathrm{Spin}}} % Spin group
\providecommand{\CS}{\ensuremath{\mathrm{CS}}}     % Chern-Simons (capitalized)
\providecommand{\cs}{\ensuremath{\mathrm{cs}}}     % Chern-Simons (lowercase)

% ======================================================================
% 6. COMMON MANIFOLDS AND TOPOLOGY
% ======================================================================
\providecommand{\Sone}{S^1}                 % Standard circle

% ======================================================================
% 7. MATH OPERATORS (Declared via \DeclareMathOperator)
% ======================================================================
\DeclareMathOperator{\End}{End}   % Endomorphisms
\DeclareMathOperator{\Aut}{Aut}   % Automorphisms
\DeclareMathOperator{\spec}{spec} % Lowercase spectrum
\DeclareMathOperator{\tr}{tr}     % Lowercase trace (if you already have \Tr)
\DeclareMathOperator{\supp}{supp} % Support
\DeclareMathOperator{\TrA}{Tr_{A}} % Partial trace over subsystem A
\DeclareMathOperator{\TrB}{Tr_{B}} % Partial trace over subsystem B
% Lie/algebra operators
\DeclareMathOperator{\Ad}{Ad}
\DeclareMathOperator{\Sym}{Sym}
\DeclareMathOperator{\su}{su}
% ======================================================================
% 8. STANDARD OPERATOR-LIKE SYMBOLS (Safe definitions via \providecommand)
% ======================================================================
\providecommand{\vol}{\operatorname{vol}}   % Volume form
\providecommand{\Vol}{\operatorname{Vol}}   % Total volume
\providecommand{\Ric}{\operatorname{Ric}}   % Ricci tensor
\providecommand{\Rm}{\operatorname{Rm}}     % Riemann tensor
\providecommand{\Tr}{\operatorname{Tr}}     % Uppercase trace
\providecommand{\Spec}{\operatorname{Spec}} % Uppercase spectrum
\providecommand{\diag}{\operatorname{diag}} % Diagonal matrix
\providecommand{\rank}{\operatorname{rank}} % Rank
\providecommand{\scal}{\mathcal R} % Scalar curvature
\providecommand{\Cl}{\operatorname{Cl}}     % Clifford algebra
\providecommand{\diam}{\operatorname{diam}} % Diameter

% ======================================================================
% 9. FUNCTIONS, TRIGONOMETRY, AND CALCULUS
% ======================================================================

% Codifferential naming conventions
% (Avoid \ddag: commonly defined by physics.sty and symbol packages)
\providecommand{\codiff}{\mathrm{d}^{\dagger}} % Codifferential operator

% Quantum mechanics bra-ket notation (needed for basis tables)
\providecommand{\ket}[1]{\lvert #1 \rangle}
\providecommand{\bra}[1]{\langle #1 \rvert}
\providecommand{\braket}[2]{\langle #1 \mid #2 \rangle}

% ======================================================================
% 10. PHYSICAL SCALES AND CONSTANTS (Paper 1c / Born Geometry)
% ======================================================================
\providecommand{\Rp}{\ensuremath{R_{\mathrm P}}}     % matched common EH curvature scale (not defined from observed G)
\providecommand{\hBorn}{\ensuremath{h_{\mathrm B}}} % physical representative 2\pi\hbar of one primitive Born scalar action cell
\providecommand{\UEHjoint}{\ensuremath{U_{J}}} % joint-Dirac EH curvature-coordinate scale 8 kappa_J
\providecommand{\Gread}{\ensuremath{G_{\rm scal}}} % scalar-action coordinate c_op^3/(h_B U); physical matching explicit
\providecommand{\GreadJoint}{\ensuremath{G_{J,\mathrm{raw}}}} % raw ordinary-trace coordinate = G_scal/96
\providecommand{\Rds}{R_{\mathrm{dS}}}    % de Sitter Curvature
\providecommand{\Mpl}{M_{\mathrm{Pl}}}    % Planck mass
\providecommand{\lp}{\ell_{\mathrm{Pl}}}  % Planck length
\providecommand{\Pmax}{P_{\mathrm{max}}}  % Maximum momentum
\providecommand{\Mplbar}{\bar{M}_{\mathrm{Pl}}} % Reduced Planck mass
\providecommand{\eff}{\mathrm{eff}}       % Subscript for effective quantities
\providecommand{\Cosmo}{\ensuremath{\Lambda}}                % Cosmological constant (use \Cosmo in source)
\providecommand{\LCDM}{\ensuremath{\Cosmo\mathrm{CDM}}}  % \LambdaCDM (typesets with \Cosmo)
\providecommand{\Lcut}{\ensuremath{\Lambda_{\mathrm{cut}}}}  % Cutoff scale (brace when used as script)
\providecommand{\Rcs}{\ensuremath{R_{\mathrm{cs}}}}      % Configuration-space radius
\providecommand{\CosmoEff}{\ensuremath{\Cosmo_{\eff}}}   % Effective cosmological constant

% Disambiguation: \Rmom is the momentum-space (S^3) radius. 
% Reserve plain R for the 4D curvature scalar in f(R).
\providecommand{\Rmom}{\ensuremath{R_{\mathrm{mom}}}} 
\providecommand{\Punit}{P_{0}} % Momentum unit set by \Rmom (= \hbar/\Rmom)
\providecommand{\Rvec}{\ensuremath{R_{\mathrm{vec}}}} % invariant vector-space radius in the reduced round sector
% Specific numeric values and expressions
\providecommand{\Vlogcap}{L_{\mathrm{Hopf}}^{(1)}}
\providecommand{\Hopf}{\operatorname{Hopf}}
% Scalar phase (S^1) sector + universal cover notation (Paper I-A)
\providecommand{\tper}{T_\star}      % Time period/scale
\providecommand{\Eper}{E_\star}      % Energy period/scale
\providecommand{\ttilde}{\tilde t}   % Universal cover time
\providecommand{\Etilde}{\tilde E}   % Universal cover energy

% ======================================================================
% 10A. SERIES-WIDE LOGICAL / NOTATIONAL DISAMBIGUATIONS
% ======================================================================
% Use these symbols to avoid collisions among: (i) the Paper--III
% no-insertion/geometric-closure class, (ii) Paper--I-A trimer diagonals,
% (iii) Paper--V-A/V-C family-space diagonals, (iv) scalar-index running,
% and (v) protected shell-clock labels.
\providecommand{\Rprim}{\ensuremath{\mathsf R_{\mathrm{prim}}^{\Sigma}}} % primitive finite structural readouts
\providecommand{\Rloop}{\ensuremath{\mathsf R_{\circ}^{\Sigma}}} % closed-loop readout role
\providecommand{\Rbdry}{\ensuremath{\mathsf R_{\partial}^{\Sigma}}} % boundary-conditioned readout role
\providecommand{\Rgen}{\ensuremath{\mathsf R_{\mathrm{fin}}^{\Sigma}}} % finite generated closure of primitive roles
\providecommand{\PhysOcc}{\operatorname{PhysOcc}}
\providecommand{\Actual}{\operatorname{Actual}}
\providecommand{\GeomLic}{\operatorname{GeomLic}}
\providecommand{\InvDist}{\operatorname{InvDist}}
\providecommand{\Dcl}{\ensuremath{\mathsf{D}_1}}      % unmarked positive Dirac-even closed spectral-readout form class
\providecommand{\Dheat}{\ensuremath{\mathsf D_{\mathrm{heat}}}} % one-parameter heat-semigroup subclass
\providecommand{\DtriI}{\ensuremath{D^{\mathrm{tri}}_1}} % Paper--I-A cyclic-trimer diagonal
\providecommand{\DtriII}{\ensuremath{D^{\mathrm{tri}}_2}}
\providecommand{\DfamI}{\ensuremath{D^{\mathrm{fam}}_1}} % Family-space diagonal selector
\providecommand{\DfamII}{\ensuremath{D^{\mathrm{fam}}_2}}
\providecommand{\alpharun}{\ensuremath{\alpha_{\mathrm{run}}}} % scalar spectral-index running
\providecommand{\Ksh}{\ensuremath{K_{\mathrm{sh}}}} % protected shell-clock coordinate
\providecommand{\ddef}{\ensuremath{\delta_{\mathrm{def}}}} % universal defect scalar after Paper--III realization

% ======================================================================
% 10B. SERIES-WIDE LOGICAL STATUS TAGS
% ======================================================================
% Upright small-caps wrapper for audit/status labels.  It prevents LaTeX from
% trying to use unavailable small-caps italic font shapes inside theorem bodies
% such as TU/lmr/m/scit.
\providecommand{\AuditLabel}[1]{}
\renewcommand{\AuditLabel}[1]{{\normalfont\scshape #1}}

% These tags are bookkeeping labels for consistency audits. They do not
% create new mathematical hypotheses; they record how a datum is being used
% in the proof chain.
\providecommand{\StatusPrimitive}{}
\renewcommand{\StatusPrimitive}{\AuditLabel{Primitive}}
\providecommand{\StatusDerived}{}
\renewcommand{\StatusDerived}{\AuditLabel{Derived}}
\providecommand{\StatusImported}{}
\renewcommand{\StatusImported}{\AuditLabel{Imported}}
\providecommand{\StatusConditional}{}
\renewcommand{\StatusConditional}{\AuditLabel{Conditional}}
\providecommand{\StatusBridge}{}
\renewcommand{\StatusBridge}{\AuditLabel{Bridge-derived}}
\providecommand{\StatusComparison}{}
\renewcommand{\StatusComparison}{\AuditLabel{Comparison-only}}

% ======================================================================
% 10C. SERIES-WIDE CLAIM-WORDING AUDIT LABELS
% ======================================================================
% These are textual audit labels used to keep modal language consistent
% across the series.  They introduce no new mathematical assumptions.
\providecommand{\ClaimForced}{}
\renewcommand{\ClaimForced}{\AuditLabel{forced}}
\providecommand{\ClaimUnique}{}
\renewcommand{\ClaimUnique}{\AuditLabel{unique}}
\providecommand{\ClaimCanonical}{}
\renewcommand{\ClaimCanonical}{\AuditLabel{canonical}}
\providecommand{\ClaimMinimal}{}
\renewcommand{\ClaimMinimal}{\AuditLabel{minimal}}
\providecommand{\ClaimClosed}{}
\renewcommand{\ClaimClosed}{\AuditLabel{closed}}
\providecommand{\ClaimPrediction}{}
\renewcommand{\ClaimPrediction}{\AuditLabel{prediction}}

% ======================================================================
% 10D. SERIES-WIDE THEOREM-CHAIN AUDIT LABELS
% ======================================================================
% These labels mark the direction of load-bearing theorem chains. They are
% audit annotations only; they introduce no new mathematical assumptions.
\providecommand{\ChainInput}{}
\renewcommand{\ChainInput}{\AuditLabel{Chain input}}
\providecommand{\ChainOutput}{}
\renewcommand{\ChainOutput}{\AuditLabel{Chain output}}
\providecommand{\ChainGuard}{}
\renewcommand{\ChainGuard}{\AuditLabel{Chain guard}}
\providecommand{\ChainNoBackflow}{}
\renewcommand{\ChainNoBackflow}{\AuditLabel{No backflow}}

% ======================================================================
% 10E. SERIES-WIDE LOCAL THEOREM-AUDIT LABELS
% ======================================================================
% These labels mark a statement-level dependency audit around individual
% theorem/definition nodes. They are bookkeeping labels only.
\providecommand{\LocalAudit}{}
\renewcommand{\LocalAudit}{\AuditLabel{Local audit}}
\providecommand{\UsesOnly}{}
\renewcommand{\UsesOnly}{\AuditLabel{Uses only}}
\providecommand{\Exports}{}
\renewcommand{\Exports}{\AuditLabel{Exports}}
\providecommand{\DoesNotUse}{}
\renewcommand{\DoesNotUse}{\AuditLabel{Does not use}}

% ======================================================================
% 10F. SERIES-WIDE NUMERICAL-ROW AUDIT LABELS
% ======================================================================
% These labels classify individual rows in numerical, benchmark, status,
% and comparison tables.  They are audit annotations only and do not change
% any theorem, value, formula, or comparison number.
\providecommand{\RowAudit}{}
\renewcommand{\RowAudit}{\AuditLabel{Numerical-row audit}}
\providecommand{\RowTheorem}{}
\renewcommand{\RowTheorem}{\AuditLabel{Theorem-fixed row}}
\providecommand{\RowBridge}{}
\renewcommand{\RowBridge}{\AuditLabel{Bridge-fixed row}}
\providecommand{\RowRepresentative}{}
\renewcommand{\RowRepresentative}{\AuditLabel{Representative row}}
\providecommand{\RowComparison}{}
\renewcommand{\RowComparison}{\AuditLabel{Comparison row}}
\providecommand{\RowCalibrated}{}
\renewcommand{\RowCalibrated}{\AuditLabel{Calibrated row}}
\providecommand{\RowEstimate}{}
\renewcommand{\RowEstimate}{\AuditLabel{Estimate row}}
\providecommand{\RowOpen}{}
\renewcommand{\RowOpen}{\AuditLabel{Open row}}
\providecommand{\RowNoSelector}{}
\renewcommand{\RowNoSelector}{\AuditLabel{No selector}}

% ======================================================================
% 10G. TABLE HELPERS
% ======================================================================
% Table helpers used by Paper V-C.  The unused readout-label preset has been
% removed; only active table helpers remain.
\providecommand{\gpsection}[1]{%
  \addlinespace[0.55em]
  \multicolumn{3}{@{}l}{\textsc{#1}}\\[-0.25em]
  \midrule
}
\providecommand{\mathcell}[1]{%
  \(\begin{gathered}#1\end{gathered}\)%
}

\newcommand{\DerBy}[1]{%
  \begin{tabular}[t]{@{}l@{}}#1\end{tabular}%
}

\newcommand{\TableNotes}[1]{%
  \par\medskip
  \begin{minipage}{0.94\linewidth}
  \footnotesize\emph{Notes.} #1
  \end{minipage}\par
}
% ======================================================================
% 10H. PAPER VI-A / QCD NOTATION
% ======================================================================
% Shared symbols used by Paper VI-A and downstream QCD/Yang--Mills handoffs.
\providecommand{\Pqcd}{\mathfrak P_{\rm VIII}^{\rm QCD}}
\providecommand{\Iqcd}{\mathfrak I_{\rm QCD}^{(7)}}
\providecommand{\muQCD}{\mu_{\rm QCD}^{\rm int}}
\providecommand{\LambdaQCD}{\Lambda_{\rm QCD}}
\providecommand{\LambdaBH}{\Lambda_{\rm BH}^{\rm QCD}}
\providecommand{\HQCD}{\mathcal H_{\rm QCD}^{\rm BH}}
\providecommand{\HQCDfin}{H_{\rm QCD}^{\rm fin}}
\providecommand{\Hraw}{\mathcal H_{\rm raw}}
\providecommand{\HphysQCD}{\mathcal H_{\rm phys}^{\rm QCD}}
\providecommand{\Vcol}{V_C}
\providecommand{\Ps}{P_{\rm sing}}
\providecommand{\PiOct}{\Pi_8}
\providecommand{\MSbar}{\overline{\rm MS}}
\providecommand{\YM}{\mathrm{YM}}
\providecommand{\Ktri}{K_{\triangle}}
\providecommand{\KYM}{K_{\YM}^{\rm BH}}
\providecommand{\aBH}{a_{\rm BH}}
\providecommand{\epsBH}{\varepsilon_{\rm BH}}
\newcommand{\RowConstructive}{\textnormal{\textsc{Constructive}}}
\newcommand{\RowObligation}{\textnormal{\textsc{Bridge obligation}}}
\newcommand{\RowTarget}{\textnormal{\textsc{Target row}}}
\newcommand{\RowValidation}{\textnormal{\textsc{Validation row}}}
\providecommand{\slashed}[1]{\not{#1}}

% ======================================================================
% 10I. PAPER VI-B / CONTINUOUS-QFT NOTATION
% ======================================================================
% Shared symbols used by Paper VI-B and any downstream continuous-QFT or
% Yang--Mills handoff paper.  All definitions are guarded so individual
% papers may override them locally before loading this file if needed.
\providecommand{\A}{\mathfrak A}
\providecommand{\Hphys}{\mathcal H_{\rm phys}}
\providecommand{\QCD}{\mathrm{QCD}}
\providecommand{\BH}{\mathrm{BH}}
\providecommand{\loc}{\mathrm{loc}}

\providecommand{\RBH}{\mathcal R_{\rm BH}}
\providecommand{\DG}{\mathfrak D_G}
\providecommand{\PW}{\mathcal{P\!W}} % Peter--Weyl spectral-symbol subspace
\providecommand{\Toe}{\operatorname{Toe}} % Toeplitz compression/representative
\providecommand{\Hol}{\operatorname{Hol}} % Holonomy

% ======================================================================
% 11. LATE-LOADED PACKAGES (Must be at the end of the preamble)
% ======================================================================
% Load late packages only if a paper has not already loaded them, avoiding
% option clashes in paper-specific preambles.
\makeatletter
\@ifpackageloaded{hyperref}{}{\usepackage{hyperref}}
\makeatother
% --- Series-wide disambiguations (Pass 3) ---
\providecommand{\diracdelta}{\delta_{\rm D}} % Dirac delta distribution (delta reserved for hierarchy)
\providecommand{\Hg}{H_{\mathrm g}} % Hodge harmonic projector (keep H for Hubble)
\providecommand{\Gr}{G_{\mathrm r}} % Green operator (keep G for Newton)
\providecommand{\FD}{\mathrm{D}} % functional/variational derivative symbol (delta reserved)
\providecommand{\Id}{\mathbf{1}} % identity/Kronecker tensor (avoid \delta_{ij})


% BEGIN INLINE SOURCE: Paper00-ActiveExternalReferences-v2_130.tex
% Framework v2.191: current reference binding; existing definitions retained.
% Framework v2.190: current22-root registry with native clock-evolution proofs.
% Framework v2.189: exact current22-root external registry, including retained v2.188 hardening owners.
% Framework v2.187: current registry synchronization only; source release v2.186 preserved.
% Framework v2.186: authoritative current 22-paper registry.
% Framework v2.182: all current 22 logical paper AUX bindings.
% Framework v2.170: current I-B v1.43 external AUX binding.
% Framework v2.167: exact current 22-root registry after CP--BL2 integration.
% Framework v2.166: exact current 22-root compact-release AUX registry.
% Framework v2.165: authoritative current root bindings.
% Framework v2.164: exact current 22 manuscript AUX bindings.
% Framework v2.162: finite metric, functional Hilbert and scoped Ward integration.
% Framework v2.161: current 22 manuscript AUX bindings; source only.
% Framework v2.156 / registry v2.113: synchronized current 22 manuscript roots; source-only registry.
% Framework v2.143 / registry v2.111: active manuscript AUX basenames; no generated AUX or PDF included.
% ======================================================================
% Central external-reference imports, numbers only, no PDF jumping
% ======================================================================

% xr-hyper is loaded centrally so that papers may retain their existing
% \usepackage{xr-hyper} lines without redefining \externaldocument after this
% registry is installed.  If plain xr was loaded earlier, retain it.
\makeatletter
\@ifpackageloaded{xr-hyper}{}{%
  \@ifpackageloaded{xr}{}{\usepackage{xr-hyper}}%
}
\makeatother

% Make external-document registration idempotent.  The registry identifies a
% logical paper by its Paper<number><optional letter> prefix.  Prefix-qualified
% imports remain independent.  Current paper roots contain no local
% \externaldocument calls; this guard also makes accidental duplicates harmless.
\ExplSyntaxOn
\prop_new:N \g_bh_external_refs_prop
\cs_new_eq:NN \BHOriginalExternalDocument \externaldocument

\cs_new_protected:Npn \bh_external_document:nnnn #1#2#3#4
 {
  \tl_set:Nn \l_tmpa_tl {#3}
  \regex_replace_once:nnN
   { \A (?: xref/ )? (Paper[0-9]+[A-G]?) .* \Z }
   { \1 }
   \l_tmpa_tl
  \tl_set:Nx \l_tmpb_tl { #1 | \tl_use:N \l_tmpa_tl }
  \prop_if_in:NVTF \g_bh_external_refs_prop \l_tmpb_tl
   {
    \PackageInfo{BornHopf}
     {Skipping~duplicate~external~refs~for~#3}
   }
   {
    \prop_gput:NVn \g_bh_external_refs_prop \l_tmpb_tl {#3}
    \IfNoValueTF{#4}
     { \BHOriginalExternalDocument[#1][#2]{#3} }
     { \BHOriginalExternalDocument[#1][#2]{#3}[#4] }
   }
 }

\RenewDocumentCommand{\externaldocument}{O{} O{} m o}
 { \bh_external_document:nnnn {#1}{#2}{#3}{#4} }

\NewDocumentCommand{\BHExternalRefs}{m}
 {
  % Do not import the aux file of the logical paper currently being compiled,
  % even when the job and registry use different version suffixes.
  \tl_set:Nx \l_tmpa_tl { \c_sys_jobname_str }
  \regex_replace_once:nnN
   { \A (Paper[0-9]+[A-G]?) .* \Z }
   { \1 }
   \l_tmpa_tl
  \tl_set:Nn \l_tmpb_tl { #1 }
  \regex_replace_once:nnN
   { \A (Paper[0-9]+[A-G]?) .* \Z }
   { \1 }
   \l_tmpb_tl
  \str_if_eq:eeTF { \tl_use:N \l_tmpa_tl } { \tl_use:N \l_tmpb_tl }
   {
    \PackageInfo{BornHopf}
     {Skipping~external~refs~for~#1~because~this~is~the~current~document}
   }
   {
    \IfFileExists{xref/#1.aux}
     { \externaldocument[][nocite]{xref/#1}[#1.pdf] }
     {
      \IfFileExists{#1.aux}
       { \externaldocument[][nocite]{#1} }
       {
        \PackageInfo{BornHopf}
         {External~aux~file~#1.aux~not~found;~no~labels~imported}
       }
     }
   }
 }
\ExplSyntaxOff

% External references are automatic.  This file is the sole owner of both the
% bibliography and the current Paper-I-A--VII-D AUX registry (30 roots including I-C).  Missing auxiliaries
% are informational on the first pass and are populated by the source-only
% two-pass series build.  Basenames deliberately omit temporary download/copy
% suffixes such as ``(1)'' or ``-1''.
  % Paper I-A: round S3 compendium.
  \BHExternalRefs{Paper1A-Geometry-v5_181}

  % Paper I-B: complete exceptional geometry and finite spectral structures.
  \BHExternalRefs{Paper1B-ExceptionalGeometry-v1_27}
  % Paper I-D: technical companion of I-A through I-C.
  \BHExternalRefs{Paper1D-TechnicalMaterial-v1_69}
  % Paper I-C: reciprocal M8 compendium.
  \BHExternalRefs{Paper1C-M8Geometry-v1_146}

\ifdefined\BHMathOnly
  % Papers I-A and I-B import only the mathematical atlas labels.
\else
  % Paper II: five physical-mathematical main papers plus one technical proof volume.
  \BHExternalRefs{Paper2A-GeometricConstitution-v3_227}
  \BHExternalRefs{Paper2B-MomentumCarrier-v2_147}
  \BHExternalRefs{Paper2C-ReciprocalLocalization-v2_177}
  \BHExternalRefs{Paper2D-SymmetrySpectra-v2_145}
  % The former II-F export labels remain in II-E; current II-F is the technical volume.
  \BHExternalRefs{Paper2E-KinematicExports-v3_140}
  % Paper II-F: complete technical proof and calculation owner for II-A--II-E.
  \BHExternalRefs{Paper2F-TechnicalMaterial-v1_71}

  % Paper III: Dynamics and Intrinsic Geometric Constitution (six components).
  \BHExternalRefs{Paper3A-GeneratedDynamics-v1_181}
  \BHExternalRefs{Paper3B-SpectralAction-v1_192}
  \BHExternalRefs{Paper3C-ProtectedHierarchy-v1_196}
  \BHExternalRefs{Paper3D-RunningResponse-v1_168}
  \BHExternalRefs{Paper3E-GravityAction-v1_204}
  \BHExternalRefs{Paper3F-CompleteDynamics-v1_281}
  % Paper III-G: complete preserved technical proof owner.
  \BHExternalRefs{Paper3G-TechnicalMaterial-v1_58}

  % Paper IV-A: generated finite quantum mechanics.
  \BHExternalRefs{Paper4A-FiniteQuantumMechanics-v4_217}
  % Paper IV-B: complete finite quantum processes and physical response.
  \BHExternalRefs{Paper4B-QuantumResponse-v1_55}
  % Paper IV-C: extended proofs and finite technical calculations.
  \BHExternalRefs{Paper4C-TechnicalMaterial-v1_50}

  % Paper V-A: The Standard Model.
  \BHExternalRefs{Paper5A-FiniteMatterGeometry-v3_219}

  % Paper V-B: Finite Spectrum and Flavor Closure.
  \BHExternalRefs{Paper5B-TerminalSpectrum-v5_202}

  % Paper V-C: Finite Interaction Boundaries of the Standard Model.
  \BHExternalRefs{Paper5C-InteractionResponse-v2_185}
  % Paper V-D: preserved technical proof and calculation owner.
  \BHExternalRefs{Paper5D-TechnicalMaterial-v1_58}

  % Paper VI-A: Gravity and Closed Cosmology.
  \BHExternalRefs{Paper6A-GravityCosmology-v2_268}

  % Paper VI-B: Geometric Dark Sector and Cosmological Readouts.
  \BHExternalRefs{Paper6B-CosmologyReadouts-v1_215}

  % Paper VI-C: complete technical proof and comparison owner.
  \BHExternalRefs{Paper6C-TechnicalMaterial-v1_53}

  % Paper VII-A: Finite QCD Kernel and Nonperturbative Closure.
\BHExternalRefs{Paper7A-QCDKernel-v1_185}

  % Paper VII-B: Hadronic, Thermal, and Nuclear QCD Readouts.
  \BHExternalRefs{Paper7B-QCDReadouts-v1_169}

  % Paper VII-C: Continuous Quantum Field Theory.
  \BHExternalRefs{Paper7C-ContinuousQFT-v2_208}

  % Paper VII-D: complete technical material for VII-A through VII-C.
  \BHExternalRefs{Paper7D-TechnicalMaterial-v1_60}

\fi

% END INLINE SOURCE: Paper00-ActiveExternalReferences-v2_130.tex


\hypersetup{colorlinks=true,linkcolor=blue,citecolor=blue,urlcolor=blue}
\emergencystretch=3em
\usepackage{xurl} % permit line breaks in reproducibility filenames


\title{\textbf{Born--Hopf Geometry I-B:}\\[0.4em]\large Exceptional group, polytope, Clifford, finite spectral and cubature geometry}
\author{Peter Kugelmann}
\date{Version 1.27 --- 20 November 2026 \\ Framework v3.077}
\BHExternalRefs{Paper1A-Geometry-v5_180}\BHExternalRefs{Paper1C-M8Geometry-v1_143}\BHExternalRefs{Paper1D-TechnicalMaterial-v1_68}

\begin{document}
\maketitle
\begin{abstract}
This paper develops the binary exceptional subgroups of $SU(2)$, their
explicit quaternionic embeddings, common tetrahedral core, normalizers,
completion lattice, and finite actions.  Over the normalized $S^3$
compendium in Paper~I-A it establishes the exceptional polytope, Hurwitz
and icosian, Clifford, finite spectral and cubature geometry, including
Petersen incidence, shell structures, round space-form spectra, Kleinian
resolutions, and McKay/ADE correspondences.  Mathematical finite-support
and operator constructions are distinguished from any additional physical
selection, coupling, or occurrence statement.  In particular, the smooth
Hopf contact integral $4\pi^2$ belongs to the chosen parent $S^3$ geometry,
not to a finite exceptional point set.
\end{abstract}
\makeatletter
\renewcommand*\l@section[2]{%
  \ifnum \c@tocdepth >\z@
    \addpenalty\@secpenalty
    \addvspace{1.0em \@plus\p@}%
    \setlength\@tempdima{2.5em}%
    \begingroup
      \parindent \z@ \rightskip \@pnumwidth
      \parfillskip -\@pnumwidth
      \leavevmode \bfseries
      \advance\leftskip\@tempdima
      \hskip -\leftskip
      #1\nobreak\hfil
      \nobreak\hb@xt@\@pnumwidth{\hss #2\kern-\p@\kern\p@}\par
    \endgroup
  \fi}
\makeatother
\tableofcontents
\clearpage
\section*{Geometric inputs and scope}
Paper~I-A supplies the normalized smooth $S^3$ geometry and the chosen
Hopf contact integral $4\pi^2$.  This paper supplies the binary
exceptional-group embeddings and develops their finite geometric and
spectral structures.  Paper~I-C treats reciprocal $M_8$ geometry, while
Paper~I-D contains additional technical calculations.  Finite-support
identities alone do not assume a physical spectral measure or a sum over
overlapping exceptional determinants.

\part{Binary exceptional groups and marked subgroup geometry}

\section{Finite Subgroups of SU(2)}
\label{p1:sec:finite}
The 3-sphere $S^3 \cong SU(2)$ admits a complete classification of finite subgroups, 
which provides the natural discrete structures on $S^3$.

\subsection{Classification Theorem}

\begin{theorem}[Finite-subgroup classification]
\label{p1a:theorem:atlas-finite-subgroup-classification-b4f160f}

\mbox{}\\[6pt]
The finite subgroups of $SU(2)$ are precisely the cyclic groups, the binary dicyclic groups, and the three exceptional binary polyhedral groups.  In the table below, $C_n$ in the first row denotes an arbitrary cyclic subgroup of $SU(2)$ of order $n$.  Under the quotient map $SU(2)\to SO(3)$, such a subgroup maps injectively if it does not contain $-1$, while if it contains $-1$ its image has half the order.  By contrast, $2C_n$ denotes the full preimage in $SU(2)$ of an order-$n$ cyclic subgroup of $SO(3)$; it is cyclic of order $2n$, but the notation records its projective origin.  Equivalently, among the subgroups containing the central element $-1$, the cyclic branch is written as the binary cyclic preimages $2C_n$. Here $D_n\subset SO(3)$ denotes the rotational dihedral group of order $2n$ and $2D_n=\pi^{-1}(D_n)$ has order $4n$ (usually $n\ge2$):
\begin{center}
\begin{tabular}{|c|c|c|c|}
\hline
Group type & Notation & Order & Image in $SO(3)$ \\
\hline
Cyclic & $C_n$ & $n$ & cyclic, with central branch $2C_n$ \\
Binary cyclic central branch & $2C_n$ & $2n$ & $C_n \subset SO(3)$ \\
Binary dicyclic & $2D_n$ & $4n$ & $D_n \subset SO(3)$ \\
Binary tetrahedral & $2T$ & $24$ & $A_4 \subset SO(3)$ \\
Binary octahedral & $2O$ & $48$ & $S_4 \subset SO(3)$ \\
Binary icosahedral & $2I$ & $120$ & $A_5 \subset SO(3)$ \\
\hline
\end{tabular}
\end{center}
\cite{conwaysmith2003quaternions,wolf2011spaces}
\end{theorem}

\begin{proof}
The quotient map \(SU(2)\to SO(3)\) sends finite subgroups of \(SU(2)\) to finite rotation groups in \(SO(3)\), whose classification is cyclic, dihedral, tetrahedral, octahedral, and icosahedral. If a finite subgroup contains the central element \(-1\), it is the full preimage of its image in \(SO(3)\), giving the binary cyclic/dicyclic and binary polyhedral groups listed. If it does not contain \(-1\), the covering map restricts injectively to the subgroup; the only such finite subgroups of \(SU(2)\) are cyclic subgroups lying in a maximal torus. See \cite{conwaysmith2003quaternions,wolf2011spaces}.
\end{proof}

\begin{definition}[Unmarked finite-support atlases of the round three-sphere]
\label{p1:def:finite-support-atlases}
For every \(\varepsilon>0\), let
\[
 \mathfrak S_\varepsilon(S^3)
 :=
 \bigl\{[Z]_{\operatorname{Isom}(S^3)}:
 Z\subset S^3\text{ is a finite }\varepsilon\text{-net}\bigr\}.
\]
The group-closed support atlas is
\[
 \mathfrak S_{\rm grp}(S^3)
 :=
 \bigl\{[C_n],[2C_n],[2D_n],[2T],[2O],[2I]\bigr\},
\]
with the integer ranges and overlap conventions of the finite-subgroup
classification above.  Brackets denote unmarked ambient-isometry classes,
not selected labelled subsets.

The first atlas contains all finite point-reference resolutions at a stated
accuracy.  The second exhausts only supports that are closed under a finite
subgroup law; it does not exhaust arbitrary finite spherical codes or finite
subsets of \(S^3\). Both atlases retain their full stated families.
\end{definition}


The groups $2T$, $2O$, and $2I$ are called \textbf{exceptional} as they do not 
belong to infinite families.  The word ``ADE'' refers here only to the corresponding McKay/ADE pattern; the literal subgroup classification includes the ordinary cyclic branch as stated here.

\subsection{The Binary Tetrahedral Group}

The binary tetrahedral group $2T$ has order 24. Its elements, viewed as unit
quaternions in $S^3 \subset \HH$, consist of two sets:

\textbf{8 elements} (quaternion units):
\begin{equation}
    \{\pm 1, \pm i, \pm j, \pm k\}
\end{equation}

\textbf{16 elements} (half-integer quaternions):
\begin{equation}
    \left\{\frac{1}{2}(\varepsilon_0 + \varepsilon_1 i + \varepsilon_2 j + \varepsilon_3 k)
    : \varepsilon_i \in \{+1, -1\}\right\}
\end{equation}

These 24 unit quaternions are exactly the Hurwitz units and, geometrically,
the 24 vertices of the regular 24-cell.
\cite{conwaysmith2003quaternions,coxeter1973regular}
Write \(\mathcal H^{\times}\) for the displayed \(24\)-point set.
It is the unit group of the Hurwitz maximal order considered below.

\subsection{The Binary Octahedral Group}

The binary octahedral group $2O$ has order 48. A concrete realization in
$S^3 \subset \HH$ is obtained by adjoining to $2T$ the 24 unit quaternions
\begin{equation}
\label{p1:eq:2O-explicit}
\frac{1}{\sqrt{2}}\{\pm 1 \pm i,\; \pm 1 \pm j,\; \pm 1 \pm k,\; \pm i \pm j,\; \pm i \pm k,\; \pm j \pm k\},
\end{equation}
where all sign choices are allowed. Thus
\begin{equation}
\label{p1:eq:2O-decomposition}
2O = 2T \sqcup \frac{1}{\sqrt{2}}\{\pm 1 \pm i,\; \pm 1 \pm j,\; \pm 1 \pm k,\; \pm i \pm j,\; \pm i \pm k,\; \pm j \pm k\}.
\end{equation}
The second set contains $6 \times 4 = 24$ distinct unit quaternions, so the
union has cardinality $24+24=48$.
\cite{conwaysmith2003quaternions,coxeter1973regular}

Geometrically, the first 24 points form the Hurwitz 24-cell, while the second
24 points form another 24-cell rotated with respect to the first. Their union
is the spherical $F_4$ root configuration on $S^3$.
\cite{coxeter1973regular,conwaySloane1999sphere}

\begin{proposition}[The octahedral two-coset support is not a direct product]
\label{p1:prop:2O-two-cosets-not-product}
For every \(g\in2O\setminus2T\),
\[
  2O=2T\sqcup g2T,
  \qquad [2O:2T]=2,
  \qquad |2O|=2|2T|=48.
\]
The second coset is the rotated \(24\)-cell displayed in
\eqref{p1:eq:2O-explicit}.  The associated exact sequence
\[
  1\longrightarrow2T\longrightarrow2O\longrightarrow C_2
  \longrightarrow1
\]
does not split.  In particular, \(2O\not\cong C_2\times2T\).
Thus the informal phrase ``two copies of \(2T\)'' may describe the two
cosets, but it does not add a fifth binary group-support class.
\end{proposition}

\begin{proof}
The displayed realization has \(2T\) as an index-two subgroup of \(2O\), so
it is normal and every element outside it generates the unique other coset.
The only element of order two in \(SU(2)\) is \(-1\), and \(-1\in2T\).
A splitting would embed \(C_2\) by an element of order two outside the
kernel \(2T\), which is impossible.  Hence the extension is nonsplit.
\end{proof}

\begin{remark}[Direct products belong to a different ambient role]
The product \(2T\times2T\) is naturally a finite subgroup of
\(SU(2)_L\times SU(2)_R\cong\operatorname{Spin}(4)\), acting by left--right multiplication.
It is not the binary octahedral subgroup of a single \(SU(2)\), and it is not
an additional entry of the one-sided support atlases below.
\end{remark}

\subsection{The Binary Icosahedral Group}

Let
\begin{equation}
\phi = \frac{1+\sqrt{5}}{2}
\end{equation}
be the golden ratio. The binary icosahedral group $2I$ has order 120 and may
be realized as the following 120 unit quaternions in $S^3 \subset \HH$:
\begin{equation}
\label{p1:eq:2I-explicit}
\left\{
\begin{array}{l}
(\pm 1,0,0,0) \text{ and all permutations}, \\
\frac{1}{2}(\pm 1, \pm 1, \pm 1, \pm 1), \\
\frac{1}{2}(0, \pm 1, \pm \phi, \pm \phi^{-1}) \text{ and all even permutations}
\end{array}
\right\},
\end{equation}
where $(a,b,c,d)$ is identified with the quaternion $a+bi+cj+dk$.
The first family contributes 8 points, the second 16 points, and the third
96 points, for a total of 120.
\cite{coxeter1973regular,conwaySloane1999sphere}

Each displayed point has norm one. For the third family,
\begin{equation}
\left|\frac{1}{2}(0,1,\phi,\phi^{-1})\right|^2
= \frac{1}{4}\bigl(1+\phi^2+\phi^{-2}\bigr) = 1,
\end{equation}
using $\phi^2=\phi+1$ and $\phi^{-2}=2-\phi$.
Geometrically these 120 points are the vertices of the regular 600-cell,
that is, the spherical $H_4$ root configuration.
\cite{coxeter1973regular,conwaySloane1999sphere}



\begin{proposition}[Exceptional nesting]
\label{p1:prop:exceptional-nesting}
Under the explicit realizations above one has
\begin{equation}
2T \subset 2O, \qquad 2T \subset 2I.
\end{equation}
Moreover, $2O$ is not a subgroup of $2I$.
\end{proposition}

\begin{proof}
The inclusions follow immediately from \eqref{p1:eq:2O-decomposition} and
\eqref{p1:eq:2I-explicit}, since the first two families in
\eqref{p1:eq:2I-explicit} are exactly the 24 vertices of the 24-cell, i.e.
the elements of $2T$. For the non-inclusion, suppose $2O \subset 2I$. Quotienting
by the common center $\{\pm1\}$ would then give an embedding
$S_4 \cong 2O/\{\pm1\} \hookrightarrow 2I/\{\pm1\} \cong A_5$. But $A_5$ has no
subgroup of order 24, by Lagrange's theorem and the standard classification of
subgroups of $A_5$. Hence such an embedding cannot exist.
\cite{conwaysmith2003quaternions,wolf2011spaces}
\end{proof}

\begin{theorem}[Normalizer fork over the common Hurwitz core]
\label{p1:thm:exceptional-normalizer-fork}
Let
\[
    H:=2T\subset SU(2)
\]
be the explicit Hurwitz subgroup used above, and let
\(I\cong2I\) be the displayed binary-icosahedral subgroup containing the same
copy of \(H\).  Then
\[
    N_{SU(2)}(H)=2O,
    \qquad
    N_I(H)=H,
    \qquad
    2O\cap I=H.
\]
Consequently
\[
    [2O:H]=2,
    \qquad
    [I:H]=5,
    \qquad
    2O/H\cong C_2.
\]
The five-element set \(I/H\) is a homogeneous coset space, not a quotient
group, because \(H\) is not normal in \(I\).
\end{theorem}

\begin{proof}
Let \(\pi:SU(2)\to SO(3)\) be the standard double cover.  Since
\(\pi(H)=A_4\), every element normalizing \(H\) projects to an element
normalizing the tetrahedral rotation group.  The four order-three axes of
\(A_4\) are permuted faithfully by its normalizer, and the octahedral rotation
group \(S_4\subset SO(3)\) realizes all these permutations.  Hence
\[
    N_{SO(3)}(A_4)=S_4
\]
and therefore
\[
    N_{SU(2)}(H)=\pi^{-1}(S_4)=2O.
\]

Inside \(I/\{\pm1\}\cong A_5\), the subgroup
\(H/\{\pm1\}\cong A_4\) is a point stabilizer in the natural five-point
action.  It is self-normalizing: if an element normalizes that point stabilizer,
it preserves its unique fixed point and hence belongs to the stabilizer.  Thus
\(N_I(H)=H\).

Finally, if \(g\in2O\cap I\), then \(g\in2O=N_{SU(2)}(H)\), so \(g\)
normalizes \(H\).  Since also \(g\in I\), one has
\(g\in N_I(H)=H\).  The reverse inclusion is already known, proving
\(2O\cap I=H\).  The index statements follow from the group orders.
\end{proof}






\section{Exceptional configurations, subgroup lattice, and polytopes}
\subsection{Associated Exceptional Spherical Configurations}

The exceptional finite subgroups of $SU(2)$ admit the following concrete
realizations in $S^3 \subset \RR^4$:

\begin{center}
\begin{tabular}{|c|c|c|c|c|c|}
\hline
Group & Order & Spherical configuration & Quotient & Basic generator degrees & ADE type \\
\hline
\(2T\) & \(24\) & 24-cell & \(A_4\) & \(6,8,12\) & \(E_6\) \\
\(2O\) & \(48\) & compound of two 24-cells & \(S_4\) & \(8,12,18\) & \(E_7\) \\
\(2I\) & \(120\) & 600-cell & \(A_5\) & \(12,20,30\) & \(E_8\) \\
\hline
\end{tabular}
\end{center}

Thus \(2T\) is the common minimal exceptional shell, while \(2O\) and \(2I\) are two inequivalent exceptional refinements: the octahedral/\(F_4\) refinement and the golden-ratio/\(H_4\) refinement realized by the 600-cell.

\begin{remark}[Dynkin inclusion is not subgroup inclusion]
\label{p1:rem:dynkin-not-subgroup}
Under the McKay correspondence,
\[
    2T\longleftrightarrow\widetilde E_6,
    \qquad
    2O\longleftrightarrow\widetilde E_7,
    \qquad
    2I\longleftrightarrow\widetilde E_8.
\]
After deleting affine nodes, one may choose graph embeddings
\[
    E_6\subset E_7\subset E_8.
\]
These are inclusions of Dynkin/root data and of corresponding exceptional-curve
configurations.  They do not imply
\(2T\subset2O\subset2I\).  In particular, deformation or partial-resolution
relations among Kleinian singularities change the singularity or its exceptional
divisor data and do not furnish an inclusion of the associated binary groups.
See \cite{McKay1980,Slodowy1980}.
\end{remark}

\subsection{Centers, projective quotients, and element orders}
\label{p1:sec:exceptional-quotients}

For each exceptional binary polyhedral group \(\Gamma\in\{2T,2O,2I\}\), the center is
\[
Z(\Gamma)=\{\pm 1\}.
\]
The projective quotients are the rotational polyhedral groups
\begin{equation}
\label{p1:eq:projective-quotients}
2T/\{\pm1\}\cong A_4,\qquad
2O/\{\pm1\}\cong S_4,\qquad
2I/\{\pm1\}\cong A_5.
\end{equation}

\begin{proposition}[Abelianizations and the perfect exceptional endpoint]
\label{p1:prop:exceptional-abelianizations}
The exceptional binary groups have abelianizations
\[
    (2T)_{\mathrm{ab}}\cong C_3,
    \qquad
    (2O)_{\mathrm{ab}}\cong C_2,
    \qquad
    (2I)_{\mathrm{ab}}=0.
\]
Equivalently,
\[
    [2T,2T]=Q_8,
    \qquad
    [2O,2O]=2T,
    \qquad
    [2I,2I]=2I.
\]
In particular, \(2I\) is the unique perfect group among the three exceptional
binary polyhedral groups.
\end{proposition}

\begin{proof}
The projective commutator subgroups are
\[
    [A_4,A_4]=V_4,
    \qquad
    [S_4,S_4]=A_4,
    \qquad
    [A_5,A_5]=A_5.
\]
Their full preimages in the displayed binary groups are respectively
\(Q_8\), \(2T\), and \(2I\).  Moreover the central element \(-1\) lies in
each binary commutator subgroup because, inside the common quaternion subgroup,
\[
    [i,j]=iji^{-1}j^{-1}=-1.
\]
Hence the binary commutator subgroups are exactly the full preimages listed
above.  Taking the corresponding quotients gives \(C_3\), \(C_2\), and the
trivial group.
\end{proof}

\begin{lemma}[Projective orbit products are one-character semi-invariants]
\label{p1:lem:projective-orbit-product-character}
Let \(\Gamma\subset SU(2)\) be finite and let
\(\mathcal O\subset\mathbb {CP}^1\) be one projective \(\Gamma\)-orbit.
Choose for each \(p\in\mathcal O\) a nonzero linear form \(\ell_p\) whose
zero line is \(p\), and put
\[
 P_{\mathcal O}(u,v):=\prod_{p\in\mathcal O}\ell_p(u,v).
\]
There is a one-dimensional character
\(\chi_{\mathcal O}:\Gamma\to\mathbb C^\times\) such that
\begin{equation}
\label{p1:eq:projective-orbit-semi-invariant}
 \gamma\cdot P_{\mathcal O}
 =\chi_{\mathcal O}(\gamma)P_{\mathcal O}.
\end{equation}
If \(\Gamma_{\rm ab}=0\), then
\(P_{\mathcal O}\in\mathbb C[u,v]^\Gamma\).
\end{lemma}

\begin{proof}
The action permutes the projective zero lines, so
\(\gamma\cdot P_{\mathcal O}\) and \(P_{\mathcal O}\) have the same degree
and divisor on \(\mathbb {CP}^1\).  They differ by a nonzero scalar.
Multiplicativity of the action makes these scalars a character, and every
one-dimensional character factors through the abelianization.
\end{proof}

\begin{theorem}[Nearest-shell derivation of the first \(2I\)-invariant degree]
\label{p1:thm:2I-shell-degree-twelve}
For \(s\in S_{2I}\), define its normalized imaginary direction
\[
 a_s:=\frac{\operatorname{Im}(s)}{\|\operatorname{Im}(s)\|}\in S^2.
\]
Under the standard \(SU(2)\)-equivariant Bloch-sphere identification
\(S^2\simeq\mathbb {CP}^1\), the set
\(\mathcal O_5:=\{a_s:s\in S_{2I}\}\) is the twelve-point \(A_5\)-orbit of
oriented endpoints of the six fivefold axes.  Consequently
\begin{equation}
\label{p1:eq:2I-shell-degree-twelve}
 \boxed{k_0(2I)=12=|S_{2I}|.}
\end{equation}
\end{theorem}

\begin{proof}
The passage \(s\mapsto a_s\) is injective on the fixed positive-real-height
shell, and conjugation of imaginary quaternions is equivariantly the
projective \(SU(2)\)-action on \(\mathbb {CP}^1\).  The projective quotient is
\(2I/\{\pm1\}\cong A_5\).  Point stabilizers in the icosahedral action have
orders \(5,3,2\), or \(1\), giving orbit sizes \(12,20,30\), or \(60\);
\(\mathcal O_5\) is the size-twelve fivefold orbit.

The zero divisor of a nonconstant invariant homogeneous polynomial is a
nonempty invariant effective divisor and therefore a union of projective
orbits, counted with multiplicity.  Its degree is at least twelve.  Conversely,
Lemma~\ref{p1:lem:projective-orbit-product-character} constructs a degree
twelve semi-invariant from \(\mathcal O_5\), and
Proposition~\ref{p1:prop:exceptional-abelianizations} gives
\((2I)_{\rm ab}=0\), so this orbit product is invariant.  The lower bound is
therefore attained.
\end{proof}

\begin{corollary}[One-shell icosahedral onset and chord datum]
\label{p1:cor:2I-one-shell-datum}
For any \(s\in S_{2I}\),
\begin{equation}
\label{p1:eq:2I-one-shell-datum}
 \boxed{\bigl(k_0(2I),\ell_{2I}\bigr)
 =\left(|S_{2I}|,\|1-s\|^2\right)
 =(12,\phi^{-2}).}
\end{equation}
Thus the degree-twelve scalar onset and primitive golden contraction arise
from the same nearest-shell orbit.  The contraction is a dimensionless
squared chord in the normalized finite geometry, not a position
length or a smooth-envelope generator.
\end{corollary}

\begin{corollary}[Character-free exceptional configuration]
\label{p1:cor:character-free-exceptional}
Among \(2T,2O,2I\), only \(2I\) has no nontrivial one-dimensional complex
character.
\end{corollary}

\begin{proof}
One-dimensional complex characters factor through the abelianization.  Apply
Proposition~\ref{p1:prop:exceptional-abelianizations}.
\end{proof}

\begin{theorem}[Axis-free and character-free exceptional support]
\label{p1:thm:axis-character-free-selector}
Let \(\Gamma\subset SU(2)\) be a noncyclic finite subgroup containing
\(-1\).  Suppose
\begin{enumerate}[label=(\roman*)]
    \item its projective image in \(SO(3)\) preserves no unoriented axis; and
    \item \(\Gamma\) has no nontrivial one-dimensional complex character.
\end{enumerate}
Then
\[
    \Gamma\cong2I.
\]
\end{theorem}

\begin{proof}
By the finite-subgroup classification, \(\Gamma\) is binary dihedral or one of
\(2T,2O,2I\).  Every projective dihedral group preserves its principal rotation
axis, so condition (i) excludes the binary-dihedral branch.  The remaining three
polyhedral groups are axis-free, and condition (ii), together with
Proposition~\ref{p1:prop:exceptional-abelianizations}, excludes \(2T\) and
\(2O\).  Hence \(\Gamma\cong2I\).
\end{proof}

\begin{corollary}[Exceptional support-role triple]
\label{p1:cor:conditional-exceptional-role-order}
Let \(H=2T\).  The unique exceptional binary polyhedral group that is both
axis-free and character-free is \(2I\), the full setwise normalizer of \(H\) in
\(SU(2)\) is \(2O\), and the common core is \(H\).  Moreover \([2T]\) is the
unique least exceptional class under structural subgroup comparison: it is
contained in both other exceptional classes, whereas no smaller exceptional
class exists.  Hence the three mathematically distinct roles form the ordered
triple
\[
    (2I,\,2O,\,2T)
    =
    (\text{axis--character-free completion},\,
     \text{normalizer},\,
     \text{core}).
\]
This ordered triple is not a subgroup chain.
\end{corollary}

\begin{proof}
The axis--character-free statement is
Theorem~\ref{p1:thm:axis-character-free-selector}.  The normalizer and
intersection statements are
Theorem~\ref{p1:thm:exceptional-normalizer-fork}.  Finally,
Proposition~\ref{p1:prop:exceptional-nesting} gives
\(2T\subset2O\), \(2T\subset2I\), and \(2O\nsubseteq2I\).  Since the
exceptional part of the finite-subgroup classification has only the three
displayed classes, \([2T]\) is the unique least exceptional class.
\end{proof}

\begin{definition}[Exceptional finite-support role atlas]
\label{p1:def:exceptional-support-role-atlas}
The unmarked exceptional role atlas is
\[
 \mathfrak S_{\rm exc}:=
 \left\{\begin{aligned}
 \mathsf{Supp}_{2T}&:=(\mathrm{core},[2T]),\\
 \mathsf{Supp}_{2O}&:=(\mathrm{normalizer},[2O]),\\
 \mathsf{Supp}_{2I}&:=(\mathrm{axis\text{-}character\text{-}free},[2I])
 \end{aligned}\right\}.
\]
These are distinct outputs of distinct support constructors, not three
competing values of one undifferentiated support variable. The atlas contains
all three outputs. A chosen quaternionic representative, vertex,
diagonal, or trimer is calculation notation inside the corresponding symmetry
orbit and is not an additional entry of the support atlas.
\end{definition}

\begin{theorem}[Exact trimer/axis limitation of binary group supports]
\label{p1:thm:trimer-axis-binary-support-range}
Let \(\pi:SU(2)\to SO(3)\) be the double cover, let
\(\Gamma\subset SU(2)\) be finite with \(-1\in\Gamma\), and write
\(\overline\Gamma:=\pi(\Gamma)\).  Impose the following specialized
admissibility disjunction:
\begin{enumerate}[label=(\roman*)]
  \item \emph{oriented-trimer branch:} there is an oriented
  three-point orbit \(Q\subset S^2\) on which \(\overline\Gamma\) acts
  freely and transitively; or
  \item \emph{axis-free branch:}
  \(\operatorname{Fix}_{\mathbb{R}P^2}(\overline\Gamma)=\varnothing\),
  equivalently no unoriented axis is invariant under all of
  \(\overline\Gamma\).
\end{enumerate}
Then, up to conjugacy,
\[
  \Gamma\in
  \left\{C_6^{\rm grp},\,2T,\,2O,\,2I\right\},
  \qquad
  C_6^{\rm grp}:=\pi^{-1}(C_3)=2C_3.
\]
More precisely, branch~(i) gives \(C_6^{\rm grp}\), whereas branch~(ii)
gives exactly \(2T,2O,2I\).  Conversely each of these four classes satisfies
its indicated branch.  Hence the exact specialized role atlas is
\[
 \mathfrak S_{\rm tri/ax}
 :=
 \left\{
   \mathsf{Supp}_{C_6^{\rm grp}}
      :=(\mathrm{oriented\text{-}trimer\ stabilizer},[C_6^{\rm grp}]),
   \ \mathsf{Supp}_{2T},\ \mathsf{Supp}_{2O},\ \mathsf{Supp}_{2I}
 \right\}.
\]
\end{theorem}

\begin{proof}
Because \(-1\in\Gamma\), the subgroup is the full binary lift
\(\Gamma=\pi^{-1}(\overline\Gamma)\).  In branch~(i), freeness and
transitivity on three points give
\(|\overline\Gamma|=|Q|=3\) by orbit--stabilizer.  Thus
\(\overline\Gamma\cong C_3\) and
\(\Gamma=\pi^{-1}(C_3)=2C_3\cong C_6^{\rm grp}\).

For branch~(ii), the finite rotation-group classification gives
\(C_n,D_n,A_4,S_4,A_5\).  Every cyclic or rotational-dihedral group preserves
its principal unoriented axis, while the three polyhedral rotation groups
have no common invariant axis.  Thus
\(\overline\Gamma\in\{A_4,S_4,A_5\}\), whose full binary lifts are
\(2T,2O,2I\).  Conversely an order-three rotation has a free three-point
orbit off its axis, and the three polyhedral rotation groups are axis-free.
\end{proof}

\begin{remark}[Scope of the trimer/axis classification]
\label{p1:rem:trimer-axis-range-not-occurrence}
Theorem~\ref{p1:thm:trimer-axis-binary-support-range} is a conditional range
theorem on the stated trimer/axis domain, with all four classes retained.
In particular, it does not restrict the generic point-resolution atlas
\(\mathfrak S_\varepsilon(S^3)\), and it does not turn
\(C_6^{\rm grp}\) into a fourth exceptional binary polyhedral group.
\end{remark}

\begin{proposition}[Possible element orders]
\label{p1:prop:element-orders}
The possible element orders in the exceptional binary polyhedral groups are:
\[
\operatorname{ord}(2T)=\{1,2,3,4,6\},
\]
\[
\operatorname{ord}(2O)=\{1,2,3,4,6,8\},
\]
\[
\operatorname{ord}(2I)=\{1,2,3,4,5,6,10\}.
\]
\end{proposition}

\begin{proof}
Let \(\bar x\in \Gamma/\{\pm1\}\) have projective order \(m\), and let \(x\in\Gamma\)
be a lift.  Then \(x^m\in\{\pm1\}\).  If \(x^m=1\), the lift has order \(m\); if
\(x^m=-1\), the lift has order \(2m\).  More concretely, the identity in the
projective group lifts to \(\pm1\), projective elements of even order lift to
elements of doubled order, and projective elements of odd order have the two lifts
of orders \(m\) and \(2m\).  Applying this to the projective order sets
\[
A_4:\{1,2,3\},\qquad S_4:\{1,2,3,4\},\qquad A_5:\{1,2,3,5\}
\]
gives exactly the displayed sets.
\end{proof}

\begin{corollary}[Common tetrahedral core]
\label{p1:cor:common-order3-core}
Both \(2O\) and \(2I\) contain elements of order \(3\), and these lie in tetrahedral
subgroups isomorphic to \(2T\). In this sense the order-\(3\) rotational core of the
exceptional tower is tetrahedral.
\end{corollary}

\begin{proof}
By Proposition~\ref{p1:prop:exceptional-nesting}, one has
\[
2T\subset 2O,\qquad 2T\subset 2I.
\]
Since \(2T/\{\pm1\}\cong A_4\) and \(A_4\) contains \(3\)-cycles, the binary tetrahedral
group contains lifts of \(3\)-cycles; by Proposition~\ref{p1:prop:element-orders},
\(2T\) indeed has elements of order \(3\). Hence both \(2O\) and \(2I\) contain
elements of order \(3\).

To identify the tetrahedral container of such order-\(3\) elements, note first that every
\(3\)-cycle in \(S_4\) lies in the subgroup \(A_4\subset S_4\). Therefore, if
\(g\in 2O\) has order \(3\), its image in
\[
2O/\{\pm1\}\cong S_4
\]
lies in an \(A_4\)-subgroup, and the full preimage of that \(A_4\) under
\[
2O \twoheadrightarrow S_4
\]
is a subgroup isomorphic to \(2T\) containing \(g\).

Similarly, every \(3\)-cycle in \(A_5\) lies in a point-stabilizer subgroup isomorphic
to \(A_4\). Hence if \(h\in 2I\) has order \(3\), its image in
\[
2I/\{\pm1\}\cong A_5
\]
lies in an \(A_4\)-subgroup, and the full preimage of that \(A_4\) under
\[
2I \twoheadrightarrow A_5
\]
is again a subgroup isomorphic to \(2T\) containing \(h\).

Thus the order-\(3\) sectors of both \(2O\) and \(2I\) lie inside binary tetrahedral
subgroups, so the common order-\(3\) rotational core is tetrahedral.
\end{proof}

\subsection{Free actions on \texorpdfstring{\(S^3\)}{S3} and spherical space forms}
\label{p1:sec:space-forms}

Every finite subgroup \(\Gamma\subset SU(2)\cong S^3\) acts freely on \(S^3\) by left
multiplication:
\[
g\cdot q := gq,\qquad g\in \Gamma,\ q\in S^3.
\]
Indeed, if \(gq=q\) for some \(q\in S^3\), then \(g=qq^{-1}=1\).

\begin{proposition}[Exceptional spherical space forms]
\label{p1:prop:exceptional-space-forms}
For \(\Gamma\in\{2T,2O,2I\}\), the quotient
\[
S^3/\Gamma
\]
is a closed spherical \(3\)-manifold with
\[
\pi_1(S^3/\Gamma)\cong \Gamma.
\]
Moreover,
\[
H_1(S^3/2T;\mathbb Z)\cong \mathbb Z/3,\qquad
H_1(S^3/2O;\mathbb Z)\cong \mathbb Z/2,\qquad
H_1(S^3/2I;\mathbb Z)=0.
\]
In particular, \(S^3/2I\) is a Poincar\'e integral homology sphere, up to the orientation convention fixed by the chosen left action.
\end{proposition}

\begin{proof}
Since the action is free, \(S^3/\Gamma\) is a spherical space form and its fundamental
group is \(\Gamma\). The first homology group is the abelianization \(\Gamma_{\mathrm{ab}}\).
For the three exceptional groups one has
\[
(2T)_{\mathrm{ab}}\cong \mathbb Z/3,\qquad
(2O)_{\mathrm{ab}}\cong \mathbb Z/2,\qquad
(2I)_{\mathrm{ab}}=0.
\]
The last statement means that \(S^3/2I\) has the same integral homology as \(S^3\).
\end{proof}

\subsection{Invariant rings and Molien series}
\label{p1:sec:molien-exceptional}

Let \(\Gamma\subset SU(2)\) act on \(\mathbb C^2\) via the defining \(2\)-dimensional
representation. The invariant ring \(\mathbb C[u,v]^\Gamma\) is finitely generated, and
its Hilbert--Poincar\'e series is the Molien series
\[
M_\Gamma(t)=\sum_{n\ge 0}\dim\!\bigl(\mathbb C[u,v]^\Gamma_n\bigr)\,t^n.
\]

For the three exceptional binary polyhedral groups one has the classical formulas \cite{Klein1884,McKay1980,Slodowy1980,conwaysmith2003quaternions}:
\begin{equation}
\label{p1:eq:molien-2T}
M_{2T}(t)=\frac{1-t^{24}}{(1-t^6)(1-t^8)(1-t^{12})},
\end{equation}
\begin{equation}
\label{p1:eq:molien-2O}
M_{2O}(t)=\frac{1-t^{36}}{(1-t^8)(1-t^{12})(1-t^{18})},
\end{equation}
\begin{equation}
\label{p1:eq:molien-2I}
M_{2I}(t)=\frac{1-t^{60}}{(1-t^{12})(1-t^{20})(1-t^{30})}.
\end{equation}

Equivalently,
\[
\mathbb C[u,v]^{2T}\cong \mathbb C[X_6,Y_8,Z_{12}]/(R_{24}),
\]
\[
\mathbb C[u,v]^{2O}\cong \mathbb C[X_8,Y_{12},Z_{18}]/(R_{36}),
\]
\[
\mathbb C[u,v]^{2I}\cong \mathbb C[X_{12},Y_{20},Z_{30}]/(R_{60}),
\]
where the subscript records the weighted degree.
These are the Kleinian surface singularities of types \(E_6,E_7,E_8\), respectively \cite{Klein1884,Slodowy1980}.

\subsection{McKay correspondence and ADE types}
\label{p1:sec:mckay-ade}

The exceptional binary polyhedral groups are exactly the exceptional finite subgroups of
\(SU(2)\), and under the McKay correspondence their representation graphs are the affine
Dynkin diagrams \cite{McKay1980,Slodowy1980}
\[
2T \longleftrightarrow \widetilde E_6,\qquad
2O \longleftrightarrow \widetilde E_7,\qquad
2I \longleftrightarrow \widetilde E_8.
\]
Removing the affine node gives the corresponding finite ADE types
\[
E_6,\qquad E_7,\qquad E_8.
\]

\begin{center}
\begin{tabular}{|c|c|c|c|c|}
\hline
Group & Order & Projective quotient & Basic generator degrees & ADE type \\
\hline
\(2T\) & \(24\) & \(A_4\) & \(6,8,12\) & \(E_6\) \\
\(2O\) & \(48\) & \(S_4\) & \(8,12,18\) & \(E_7\) \\
\(2I\) & \(120\) & \(A_5\) & \(12,20,30\) & \(E_8\) \\
\hline
\end{tabular}
\end{center}

Thus the three exceptional binary polyhedral groups admit parallel descriptions in
terms of spherical geometry, invariant theory, and ADE representation theory.

For later reference, the finite root systems attached by the correspondence
have sizes \(|\Phi|=rh\), with \(r\) the rank and \(h\) the Coxeter number
\cite{Humphreys1990}:
\[
|\Phi(E_6)|=6\cdot12=72,\qquad
|\Phi(E_7)|=7\cdot18=126,\qquad
|\Phi(E_8)|=8\cdot30=240.
\]
These are invariants of the root data alone; no pairing with any finite
carrier product is asserted anywhere in Paper~I-B.

\subsection{Exceptional subgroup lattice}
\label{p1:sec:exceptional-lattice}

The subgroup relations among the three exceptional groups are asymmetric:
\[
2T\subset 2O,\qquad 2T\subset 2I,\qquad 2O\not\subset 2I.
\]

\begin{proposition}[Exceptional subgroup lattice]
\label{p1:prop:exceptional-lattice}
Here ``exceptional subgroup'' means a subgroup of binary tetrahedral, binary octahedral, or binary icosahedral type; cyclic and binary-dihedral subgroups are not being classified in this proposition. Up to conjugacy inside \(SU(2)\), the only proper exceptional subgroup of \(2O\) is \(2T\),
and the only proper exceptional subgroup of \(2I\) is again \(2T\). There is no embedding
\(2O\hookrightarrow 2I\).
\end{proposition}

\begin{proof}
The inclusions \(2T\subset 2O\) and \(2T\subset 2I\) were established above.
Let \(K\) be a proper exceptional subgroup of either \(2O\) or \(2I\).  Since every
binary exceptional group contains the central element \(-1\), quotienting by
\(\{\pm1\}\) identifies \(K/\{\pm1\}\) with an exceptional rotational subgroup of
\(S_4=2O/\{\pm1\}\) or \(A_5=2I/\{\pm1\}\).  The exceptional rotational groups in
\(SO(3)\) are \(A_4,S_4,A_5\).  Inside \(S_4\), the only proper exceptional
subgroup type is \(A_4\) (the rotational tetrahedral subgroup); inside \(A_5\),
the only proper exceptional subgroup type is again \(A_4\).  Lifting these
\(A_4\)-subgroups through the double cover \(SU(2)\to SO(3)\) gives the binary
tetrahedral subgroup \(2T\).  This proves the uniqueness, up to conjugacy, of the
proper exceptional subgroup in both \(2O\) and \(2I\).

If \(2O\subset 2I\), then quotienting by the common center would give an embedding
\[
S_4 \cong 2O/\{\pm1\}\hookrightarrow 2I/\{\pm1\}\cong A_5,
\]
which is impossible: equivalently, \(S_4\) is not a subgroup of \(A_5\) by the
classification just used, and also \(|S_4|=24\) does not divide \(|A_5|=60\).
Thus \(2O\not\subset 2I\).
\end{proof}

\begin{remark}
This subgroup lattice is the purely mathematical reason that the exceptional tower is not a
chain \(2T\subset 2O\subset 2I\), but rather a common tetrahedral core with two inequivalent
refinements.
\end{remark}

\begin{corollary}[Exceptional subgroup structure inside \(2I\)]
\label{p1:cor:exported-exceptional-carrier-skeleton}
Up to conjugacy, the only proper exceptional subgroup type in
\(2I\subset SU(2)\) is \(2T\).  Its Hurwitz \(24\)-cell support has the intrinsic
decompositions
\[
    24=8+16,
    \qquad
    24=8+8+8.
\]
The binary octahedral group \(2O\) is a distinct refinement of the same
\(2T\) core:
\[
2T\subset 2I,\qquad
2T\subset 2O,\qquad
2O\not\subset2I.
\]
Thus the exceptional subgroup lattice is a fork over \(2T\), not a chain.
\end{corollary}

\begin{proof}
The uniqueness of the proper exceptional subgroup type inside \(2I\) is
Proposition~\ref{p1:prop:exceptional-lattice}.  The order and Hurwitz model of
\(2T\) are given in the binary-tetrahedral subsection, where the \(24\) units
split as the quaternionic octet and the sixteen half-integer Hurwitz units.
The triality-octet decomposition is the structural invariant
\(24=8+8+8\) of Theorem~\ref{p1:thm:s3one-structural-invariants}.
The non-nesting statement \(2O\not\subset2I\) is again
Proposition~\ref{p1:prop:exceptional-lattice}.\qedhere
\end{proof}


\subsection{The 24-Cell}
\label{p1:sec:24cell}
The 24-cell is the unique exceptional self-dual regular 4-polytope, distinct from the simplex family and from the hypercube/cross-polytope dual pair.  Its existence is a special feature of four dimensions \cite{coxeter1973regular,conwaySloane1999sphere}.

\begin{center}
\begin{tabular}{|c|c|}
\hline
Property & Value \\
\hline
Schläfli symbol & $\{3, 4, 3\}$ \\
Vertices & 24 \\
Edges & 96 \\
Faces (triangular) & 96 \\
Cells (octahedral) & 24 \\
Vertex figure & Cube \\
Nearest-neighbor chordal edge length (unit $S^3\subset\RR^4$) & $1$ \\
Nearest-neighbor spherical edge angle (unit $S^3$) & $\pi/3$ \\
Circumradius & 1 \\
\hline
\end{tabular}
\end{center}

\begin{proposition}
\label{p1:prop:24cell-adjacency}
In the standard embedding via Hurwitz units on the unit $S^3\subset\RR^4$, the 24-cell has
circumradius $1$ and chordal edge length $1$. Equivalently, for each vertex there
are exactly 8 nearest neighbors at Euclidean distance $1$; for example,
\begin{equation}
\left|1-\frac{1}{2}(1+i+j+k)\right|^2
= \left|\frac{1}{2}(1-i-j-k)\right|^2 = 1.
\end{equation}
The value $\sqrt{2}$ occurs for nonadjacent pairs such as $1$ and $i$, while the
antipodal pair $1$ and $-1$ has distance $2$.
\cite{coxeter1973regular,conwaySloane1999sphere}
\end{proposition}

\begin{proof}
The 24 vertices are the Hurwitz units, hence lie on the unit sphere in
$\RR^4 \cong \HH$, so the circumradius is 1 in the unit Hurwitz normalization. For general radius \(\Rmom\), all chordal lengths scale by \(\Rmom\). For the displayed adjacent pair one
computes
\[
\left|1-\frac{1}{2}(1+i+j+k)\right|^2
= \frac{1}{4}(1^2+1^2+1^2+1^2)=1,
\]
so the chordal edge length is 1. On the unit sphere this corresponds to
spherical geodesic edge length $\pi/3$, since $|p-q|=2\sin(d_{S^3}(p,q)/2)$.
By the standard combinatorics of the 24-cell, each
vertex is incident to 8 edges, hence has 8 nearest neighbors. For a nonadjacent
pair such as $1$ and $i$ one gets $|1-i|^2=2$, while $|1-(-1)|=2$ for the
antipodal pair.
\end{proof}

\begin{proposition}[Explicit Hurwitz-to-\(D_4\) coordinate map]
\label{p1:prop:hurwitz-d4-coordinate-map}
\leavevmode\par
Write a Hurwitz unit as \(q=(a,b,c,d)^{\mathsf T}\) in the quaternionic
coordinate frame \(a+bi+cj+dk\), and set
\[
 B=
 \begin{pmatrix}
 1&1&0&0\\
 1&-1&0&0\\
 0&0&1&1\\
 0&0&1&-1
 \end{pmatrix},
 \qquad
 O:=\frac1{\sqrt2}B.
\]
Then \(O\in O(4)\) and
\[
 B(2T)=\sqrt2\,O(2T)
 =\{\,\pm e_i\pm e_j:1\le i<j\le4\,\}
 =\Phi(D_4).
\]
Equivalently, the unit Hurwitz shell is carried to the unit-sphere
normalization \(\Phi(D_4)/\sqrt2\) by \(O\), whereas the standard roots of
squared length \(2\) are obtained by the additional factor \(\sqrt2\).
Neither the scale alone nor the identity coordinate frame gives the displayed
standard root coordinates.
\end{proposition}

\begin{proof}
The rows of \(B\) are mutually orthogonal and have squared norm \(2\), so
\(O=B/\sqrt2\) is orthogonal.  The quaternion units map to the roots supported
on the coordinate pairs \(\{1,2\}\) and \(\{3,4\}\).  For a half-integer
Hurwitz unit,
\[
 B\frac12
 \begin{pmatrix}\varepsilon_0\\\varepsilon_1\\\varepsilon_2\\\varepsilon_3\end{pmatrix}
 =\frac12
 \begin{pmatrix}
 \varepsilon_0+\varepsilon_1\\
 \varepsilon_0-\varepsilon_1\\
 \varepsilon_2+\varepsilon_3\\
 \varepsilon_2-\varepsilon_3
 \end{pmatrix}.
\]
Exactly one entry in each of the first and last coordinate pairs is nonzero,
and each nonzero entry is \(\pm1\).  These give the remaining sixteen
standard \(D_4\) roots.  The two sets have \(24\) elements, proving equality.
\end{proof}

\begin{proposition}[24-cell vertex figure]
\label{p1:prop:24cell-vertex-figure}
The vertex figure of the 24-cell is a cube with 8 vertices and 12 edges.
\cite{coxeter1973regular}
\end{proposition}

\begin{proof}[Imported standard reference]
The vertex figure of the regular \(24\)-cell is the convex hull of the nearest neighbors of a vertex, and the regular \(24\)-cell Schl\"afli data give this figure as a cube. See \cite{coxeter1973regular}.
\end{proof}


\begin{theorem}[Radial shell stratification of the Hurwitz 24-cell]
\label{p1:thm:hurwitz-radial-shells}
With base vertex \(1\in 2T\subset S^3\subset\HH\), the Hurwitz unit shell decomposes as
\[
2T=
\{1\}
\sqcup \mathcal N_8^+
\sqcup \mathcal M_6
\sqcup \mathcal N_8^-
\sqcup \{-1\},
\]
where
\[
\mathcal N_8^+
=
\left\{\frac12(1+\varepsilon_1 i+\varepsilon_2 j+\varepsilon_3 k):
\varepsilon_a\in\{\pm1\}\right\},
\]
\[
\mathcal M_6=\{\pm i,\pm j,\pm k\},
\]
\[
\mathcal N_8^-
=
\left\{\frac12(-1+\varepsilon_1 i+\varepsilon_2 j+\varepsilon_3 k):
\varepsilon_a\in\{\pm1\}\right\}.
\]
Equivalently, the shells have counts
\[
1,
\quad 8,
\quad 6,
\quad 8,
\quad 1.
\]
After removing the local edge-neighbor cube \(\mathcal N_8^+\), the remaining
sixteen-point global complement is
\[
\mathcal G_{16}
:=\{1\}\sqcup\mathcal M_6\sqcup\mathcal N_8^-\sqcup\{-1\},
\]
and, relative to the chosen base vertex, it carries the near/opposite split
\[
\mathcal G_{16}=\mathcal G_7\sqcup\mathcal G_9^{\rm shell},
\qquad
\mathcal G_7:=\{1\}\sqcup\mathcal M_6,
\qquad
\mathcal G_9^{\rm shell}:=\mathcal N_8^-\sqcup\{-1\}.
\]
Thus the normalized shell-complement weights are
\[
 w_7=\frac{7}{16},
 \qquad
 w_9=\frac{9}{16}.
\]
\end{theorem}

\begin{proof}
The displayed sets are exactly the Hurwitz units grouped by real part
\(1,\frac12,0,-\frac12,-1\).  For unit quaternions the Euclidean distance from
\(1\) is determined by the real part, since
\[
|1-q|^2=2-2\operatorname{Re}(q).
\]
Hence these groups are the radial distance shells about \(1\).  Their
cardinalities are immediate from the displayed formulae: \(1,8,6,8,1\).  Removing
\(\mathcal N_8^+\) leaves the stated sixteen points.  The near block has
\(1+6=7\) points and the opposite block has \(8+1=9\) points, so division by
\(|\mathcal G_{16}|=16\) gives the two normalized weights.
\end{proof}

\begin{corollary}[Tetrahedral shell atom equals the complement counting atom]
\label{p1:cor:2T-shell-complement-counting-atom}
For the canonical shell complement
\(\mathcal G_{16}=2T\setminus S_{2T}
=\mathcal G_7\sqcup\mathcal G_9^{\rm shell}\),
\begin{equation}
\label{p1:eq:2T-shell-complement-counting-atom}
 \boxed{\omega_{2T}=\frac1{16}=\frac1{|\mathcal G_{16}|}.}
\end{equation}
Thus normalized counting measure gives
\begin{equation}
\label{p1:eq:G16-counting-weights}
 \boxed{\mu_{16}(\mathcal G_7)=\frac7{16},\qquad
 \mu_{16}(\mathcal G_9^{\rm shell})=\frac9{16}.}
\end{equation}
These are the weights of normalized counting measure on the stated partition.
\end{corollary}

\begin{proof}
The nearest tetrahedral shell has eight elements and real height \(1/2\), so
\(\omega_{2T}=(1/2)/8=1/16\).  Its complement has
\(|2T|-|S_{2T}|=24-8=16\) elements.  The block weights then follow from
the proved \(7+9\) cardinality split.
\end{proof}

\subsection{\texorpdfstring{Coxeter Orthoschemes and Characteristic Angles on $S^3$}{Coxeter Orthoschemes and Characteristic Angles on S3}}
\label{p1:sec:coxeter}

Regular 4-polytopes provide canonical discrete geometries on $S^3\subset\RR^4$.
Beyond vertex sets (binary polyhedral groups), the \emph{Coxeter orthoscheme} data
encodes the fundamental spherical simplex and its characteristic angles.

\begin{proposition}[Coxeter angle data for the $24$-, $120$-, and $600$-cells]
\label{p1:prop:coxeter-angle-data}
The following two angle datasets are fixed discrete invariants of the regular
four-polytopes most closely tied to $S^3$ discretizations: the dihedral angles between 3-cells and the adjacent essential angles of the characteristic Coxeter orthoschemes.

\smallskip
\noindent\textup{(i) Dihedral angles between cells.}
\begin{center}
\begin{tabular}{|c|c|c|}
\hline
Polytope & Schl\"afli symbol & Dihedral angle between cells \\
\hline
24-cell & $\{3,4,3\}$ & $120^\circ$ \\
120-cell & $\{5,3,3\}$ & $144^\circ$ \\
600-cell & $\{3,3,5\}$ & $164.478\ldots^\circ$ \\
\hline
\end{tabular}
\end{center}

\smallskip
\noindent\textup{(ii) Characteristic orthoscheme angles.}
For the characteristic Coxeter orthoschemes the adjacent essential angles are
$[3,4,3]:60^\circ,45^\circ,60^\circ$ for the 24-cell,
$[3,3,5]:60^\circ,60^\circ,36^\circ$ for the 600-cell, and
$[5,3,3]:36^\circ,60^\circ,60^\circ$ for the 120-cell; non-adjacent facets contribute right angles. These angles are not identical to the cell-cell dihedral angles, e.g.\ the 24-cell cell-cell dihedral angle is $120^\circ$, but they are canonical within the Coxeter/orthoscheme description of the same discrete geometry.
\end{proposition}

\begin{proof}[Imported standard reference]
These are the standard Coxeter--Schl\"afli data for the corresponding regular four-polytopes and their characteristic orthoschemes; see Coxeter's tables for regular polytopes and the Coxeter-system conventions in \cite{coxeter1973regular,Humphreys1990}.  The proposition states the orthoscheme entries separately from the cell--cell dihedral angles.
\end{proof}



The distinction between Coxeter-orthoscheme angles and cell--cell dihedral
angles is part of Proposition~\ref{p1:prop:coxeter-angle-data}.


\subsection{Hurwitz Integers and the Maximal Order}
\label{p1:sec:hurwitz}
\label{p1:subsec:hurwitz-integers-maximal-order}
The quaternion algebra $\HH$ admits discrete subrings called \textbf{orders}.  Maximal orders are unique only up to the standard quaternionic equivalence/conjugacy relation; the displayed Hurwitz order is the distinguished representative fixed by the chosen \(24\)-cell/Hurwitz-unit realization.

\begin{definition}[Hurwitz Integers]
\label{p1:def:hurwitz-integers}
\mbox{}\\[6pt] 
The \textbf{Hurwitz integers} are:
\begin{equation}
    \mathcal{H} = \left\{a + bi + cj + dk : 
    \begin{array}{l}
    a,b,c,d \in \ZZ, \text{ or} \\
    a,b,c,d \in \ZZ + \tfrac{1}{2}
    \end{array}
    \right\}
\end{equation}
Equivalently, $\mathcal{H} = \ZZ\langle 1, i, j, \frac{1}{2}(1+i+j+k)\rangle$.
\end{definition}

\begin{theorem}[Hurwitz, 1896]
\label{p1:thm:hurwitz-maximal-order}
\mbox{}\\[6pt] 
The Hurwitz integers $\mathcal{H}$ form a maximal order in $\HH$; within the chosen \(24\)-cell realization they are the distinguished standard representative, unique up to quaternionic equivalence/conjugacy.
\cite{voight2021quaternion,conwaysmith2003quaternions}
\end{theorem}

\begin{proof}[Imported standard reference]
\leavevmode\par\noindent
The Hurwitz order is the classical maximal order in Hamilton's quaternion
algebra over \(\ZZ\).  The uniqueness assertion used here is uniqueness up to
quaternionic equivalence/\allowbreak conjugacy after choosing the displayed
\(24\)-cell representative.  See
\cite{voight2021quaternion,conwaysmith2003quaternions}.
\end{proof}


This equivalence-class uniqueness, together with the fixed embedded representative, is the number-theoretic input used here; the paper never requires literal uniqueness of a maximal order before choosing the Hurwitz-unit realization.

\begin{theorem}[Units of the Maximal Order]
\label{p1:thm:hurwitz-units-maximal-order}
\mbox{}\\[6pt] 
The group of units (invertible elements) of $\mathcal{H}$ is:
\begin{equation}
    \mathcal{H}^\times = 2T
\end{equation}
the binary tetrahedral group of order 24.
\cite{voight2021quaternion,conwaysmith2003quaternions}
\end{theorem}

\begin{proof}
An element $q \in \mathcal{H}$ is a unit if and only if $|q|^2 = 1$.
For the integer lattice part of \(\mathcal H\), the equation
\(a^2+b^2+c^2+d^2=1\) gives exactly
\(\pm1,\pm i,\pm j,\pm k\), hence 8 units.  For the half-integer coset
\((a,b,c,d)\in(\ZZ+\tfrac12)^4\), the same norm equation forces each
coefficient to have absolute value \(\tfrac12\), giving the 16 elements
\((\pm1\pm i\pm j\pm k)/2\).  Thus exactly \(8+16=24\) elements of
\(\mathcal H\) have norm one, namely the binary tetrahedral group \(2T\).
\end{proof}

\begin{corollary}[Hurwitz units as the 24-cell vertices]
\label{p1:cor:hurwitz-units-24cell-vertices}
The vertices of the 24-cell inscribed in $S^3$ are precisely the units 
of the Hurwitz integers.
\cite{voight2021quaternion,coxeter1973regular}
\end{corollary}

\begin{proof}
The previous theorem identifies the unit group of the Hurwitz order with the \(24\) Hurwitz units. These are exactly the displayed \(24\) unit quaternions, which form the standard vertex set of the regular \(24\)-cell in \(S^3\subset\HH\). See \cite{voight2021quaternion,coxeter1973regular}.
\end{proof}


After fixing the displayed Hurwitz representative, this records the standard relationship between the Hurwitz order, the binary tetrahedral unit group, and the \(24\)-cell vertex set in \(S^3\).


\subsection{Spherical 3-Manifolds as Quotients}
\label{p1:subsec:spherical-quotients}

\begin{theorem}[Elliptization / spherical space form theorem]
\label{p1:thm:elliptization-spherical-space-form}
\mbox{}\\[6pt] 
A closed orientable connected 3-manifold $M$ admits a Riemannian metric of constant positive sectional curvature if and only if $\pi_1(M)$ is finite. In this case $M$ is isometric to a spherical space form $S^3(\Rmom)/\Gamma$, where $\Rmom$ is determined by the sectional curvature $K=1/\Rmom^2$ and $\Gamma \subset \mathrm{SO}(4)$ is a finite subgroup acting freely by isometries on $S^3(\Rmom)$.
\cite{Thurston1997,Perelman2002entropy,Perelman2003surgery,Perelman2003extinction,wolf2011spaces}
\end{theorem}

\begin{proof}[Imported standard reference]
The equivalence ``$\pi_1(M)$ finite $\Leftrightarrow$ $M$ admits a constant positive curvature metric'' in the closed orientable case is the elliptization theorem, proved as part of Thurston's geometrization program via Perelman's Ricci flow work; see \cite{Thurston1997,Perelman2002entropy,Perelman2003surgery,Perelman2003extinction}. If $M$ carries a constant positive curvature metric, its universal cover is $S^3(\Rmom)$ and $M \cong S^3(\Rmom)/\Gamma$ for the deck group $\Gamma \cong \pi_1(M)$ acting freely by orientation-preserving isometries. Hence, after identifying $\mathrm{Isom}^+(S^3(\Rmom))\cong SO(4)$, we have $\Gamma\subset SO(4)$ finite. For the classification of orientable spherical space forms (equivalently, free finite orientation-preserving isometric actions on $S^3$) see \cite{wolf2011spaces}.
\end{proof}

\begin{corollary}[Left-quotient spherical space forms]
\label{p1:cor:left-quotient-spherical-space-forms}
Every finite subgroup \(\Gamma\subset SU(2)\cong S^3\) acts freely on \(S^3\) by left multiplication, and the quotients
\[
S^3/\Gamma
\]
are spherical space forms. Left cyclic subgroups produce the left-cyclic lens-space branch, while the exceptional binary polyhedral groups \(2T,2O,2I\) produce the tetrahedral, octahedral, and Poincar\'e spherical manifolds.
\cite{wolf2011spaces}
\end{corollary}

\begin{proof}
If \(\gamma q=q\) under left multiplication, then \(\gamma=1\), so the action is free. Since left multiplication by unit quaternions is isometric for the round metric, the quotient is a spherical space form. The branch descriptions are the standard cyclic and binary-polyhedral quotient families; see \cite{wolf2011spaces}.
\end{proof}


\begin{proposition}[Quotient branches are equivariant/invariant restrictions of the simply connected model]
\label{p1:prop:quotient-invariant-restriction}
Let \(\Gamma\subset SU(2)\) be a finite subgroup acting freely on \(S^3\), and let
\[
q:S^3\to S^3/\Gamma
\]
be the quotient map. Then, for scalar functions,
\[
U:L^2(S^3/\Gamma)\longrightarrow L^2(S^3)^\Gamma,
\qquad
Uf=|\Gamma|^{-1/2}f\circ q,
\]
is unitary for the quotient and covering Riemannian measures and intertwines the
scalar Laplacians:
\[
U\Delta_{S^3/\Gamma}=\Delta_{S^3}U.
\]
Consequently, for every Borel set $I\subset[0,\infty)$,
\[
U\operatorname{Ran}\mathbf1_I(\Delta_{S^3/\Gamma})
=
\operatorname{Ran}\mathbf1_I(\Delta_{S^3})\cap L^2(S^3)^\Gamma.
\]
Thus Laplace-spectral cutoffs, including finite sums of eigenspaces, are exactly
the invariant restrictions of the corresponding simply connected cutoffs.
For sections of a bundle on the quotient, the analogous statement requires a
specified equivariant lift to the pulled-back bundle; spinor statements also
depend on the chosen quotient spin structure.  No assertion is made here for
an arbitrary finite truncation that is not defined by such an intertwining
functional calculus.
\end{proposition}

\begin{proof}
The covering has degree $|\Gamma|$, so the displayed normalization makes
pullback unitary onto the invariant subspace.  Because $q$ is a local
isometry, the scalar Laplacians intertwine on smooth functions and hence as
self-adjoint operators on their $L^2$ domains.  The spectral theorem then gives
the displayed identity for every spectral projector $\mathbf1_I$.  For vector
bundles and spinor bundles the same reasoning applies after fixing the lifted
fiber action; quotient sections correspond to equivariant lifted sections, and
the spin lift is part of the quotient spin-structure data.
\end{proof}


\subsection{Lens Spaces and Exceptional Spherical Manifolds}

Finite cyclic subgroups acting by left multiplication give the left-cyclic lens-space branch, usually represented by \(L(p,1)\) up to orientation/convention.  The general lens spaces are obtained by allowing the broader diagonal \(SO(4)\)-isometry action
\[
(z_1,z_2)\longmapsto (e^{2\pi i/p}z_1,e^{2\pi i q/p}z_2),
\qquad p\ge2,\quad \gcd(p,q)=1,
\]
on \(S^3\subset\CC^2\), and are denoted
\[
L(p,q)=S^3/\ZZ_p.
\]
Thus \(L(p,q)\) for general \(q\) belongs to the diagonal \(SO(4)\) quotient branch, not only to the pure left-\(SU(2)\) subgroup branch.

The exceptional (non-cyclic) binary polyhedral groups produce three distinguished spherical 3-manifolds:

\begin{itemize}
\item Binary tetrahedral group $2T$ (order 24): tetrahedral spherical manifold,
\item Binary octahedral group $2O$ (order 48): octahedral spherical manifold,
\item Binary icosahedral group $2I$ (order 120): Poincaré homology sphere $\Sigma$.
\end{itemize}

The Poincaré homology sphere has $H_*(\Sigma;\ZZ) \cong H_*(S^3;\ZZ)$ but non-trivial fundamental group of order 120.

\begin{remark}
The Hopf fibration descends to a Seifert fibration only for quotient actions compatible with the chosen fiber foliation. For a left deck action \(\Gamma\subset SU(2)\), the right Hopf fibration \(q\mapsto q e^{i\phi}\) is preserved because left and right multiplications commute. The fixed common-phase convention \(\pi(q)=q^{-1}iq\), whose fibers are left \(U(1)\)-orbits, descends only for deck groups normalizing that chosen left \(U(1)\). This does not mean that a free deck action on \(S^3\) has fixed points upstairs.
\end{remark}

\begin{corollary}[Simply connected positive-curvature model]
\label{p1:cor:simply-connected-positive-curvature-model}
The round 3-sphere $S^3(\Rmom)$ is the unique simply connected compact 3-manifold admitting constant positive sectional curvature.
\cite{wolf2011spaces,do1992riemannian}
\end{corollary}

\begin{proof}
By the spherical space-form theorem, any compact constant-positive-curvature \(3\)-manifold is \(S^3(\Rmom)/\Gamma\) for a finite freely acting group \(\Gamma\). If the quotient is simply connected, \(\Gamma\) is trivial, so the manifold is the round \(S^3(\Rmom)\). See \cite{wolf2011spaces,do1992riemannian}.
\end{proof}



\section{Tetrahedral-core completions, Petersen incidence and finite actions}
\begin{theorem}[Two icosahedral completions of one tetrahedral core]
\label{p1:thm:v273-two-icosahedral-completions}
For a fixed embedded binary tetrahedral subgroup
$H\cong2T\subset G=SU(2)$, put $O=N_G(H)\cong2O$.
Exactly two distinct embedded binary icosahedral subgroups
$I_+,I_-\cong2I$ contain $H$. They satisfy
\[
 I_+\cap I_-=H,\qquad O\cap I_\pm=H,
 \qquad\overline{\langle I_+,I_-\rangle}=G.
\]
Conjugation by $O/H\cong C_2$ exchanges the completions. The space
of pairs $(H,I)$ with $H\subset I$ is $G/H$, while the space of
embedded tetrahedral cores is $G/O$; hence
\[
 G/H\longrightarrow G/O
\]
is a connected two-sheeted cover without a global continuous section.
This statement determines the unordered pair relative to $H$,
not an embedded physical core or a branch choice from the unmarked parent.
\end{theorem}
\begin{proof}
Existence follows from the explicit $2I$ embedding containing $H$.
All $2I$ subgroups are conjugate, and its five tetrahedral subgroups
(the inverse images of the five $A_4$ point stabilizers in $A_5$)
are conjugate inside $2I$. Thus every $I'\supset H$ is the conjugate
of a fixed $I_+$ by an element of $N_G(H)=O$.
The subgroup $I_+$ is self-normalizing in $G$: a normalizer has a
finite image in the automorphism group of the noncentral finite group
$I_+$, its centralizer in $G$ is the center, and the finite-subgroup
classification leaves no larger finite subgroup containing $2I$.
Its $O$-stabilizer is therefore $O\cap I_+=H$ by
Theorem~\ref{p1:thm:exceptional-normalizer-fork}. The orbit has
$[O:H]=2$ elements. Their intersection contains $H$; since $H$
is maximal proper in each $2I$, they intersect exactly in $H$.
The subgroup generated by two distinct $2I$ copies cannot be finite,
again by the finite-subgroup classification. No proper positive-dimensional
closed subgroup of $SU(2)$ can contain a binary icosahedral group,
since such a subgroup lies in a torus normalizer. Therefore its
closure is $G$. The two stabilizers under conjugation are $H$ and $O$,
so the incidence spaces have the stated forms. As $G/H$ is connected,
a section of a two-sheeted cover would disconnect it; none exists.
\end{proof}

\begin{corollary}[Simultaneous exceptional invariance and its no-go]
\label{p1:cor:v273-twins-sharp-common-invariants}
For the left regular $G$-action on $L^2(G)$,
\[
 (L^2(G))^{I_+}\cap(L^2(G))^{I_-}
       =(L^2(G))^{G}=\mathbb C\boldsymbol1.
\]
Thus the two exact finite-group comparison presentations can coexist,
but enforcing both as sharp invariance constraints on one scalar state
would eliminate all nonconstant harmonics. This does not preclude
mathematical co-realization of their projector families.
\end{corollary}
\begin{proof}
The intersection is invariant under the group generated by $I_+$
and $I_-$, whose closure is $G$ by the theorem. Strong continuity of
the left regular representation extends invariance to the closure.
The left-invariant $L^2(G)$ functions are constants.
\end{proof}

\begin{theorem}[Finite mobile tetrahedral completion inside the quaternionic spin group]
\label{p1:thm:v237-mobile-tetrahedral-completion}
Let \(H\cong2T\subset \operatorname{Sp}(1)\cong SU(2)\) be a binary-tetrahedral core and let
\(\mathcal F_H\) be the associated unordered set of its four rank-one
threefold-axis projectors in \(\operatorname{Im}\mathbb H\).  Thus
\[
 \operatorname{Stab}_{\operatorname{Sp}(1)}(\mathcal F_H)=N_{\operatorname{Sp}(1)}(H)=2O.
\]
Let \(\Gamma<\operatorname{Sp}(1)\) be finite, assume \(H\subset\Gamma\), and suppose that
\(\Gamma\) genuinely moves the local tetrahedral projector frame:
\[
 \exists g\in\Gamma:\qquad g\mathcal F_Hg^{-1}\ne\mathcal F_H.
\]
Then
\[
 \boxed{\Gamma\cong2I.}
\]
Equivalently, a finite same-owner symmetry containing one tetrahedral core and
moving its local Clifford projector frame has a unique exceptional completion.
\end{theorem}

\begin{proof}
By the finite-subgroup classification of \(SU(2)\), a finite subgroup
containing \(H\cong2T\) can only have exceptional type \(2T\), \(2O\), or
\(2I\).  The first does not move the frame.  In \(2O=N_{\operatorname{Sp}(1)}(H)\), the
entire projector frame \(\mathcal F_H\) is setwise fixed; equivalently \(H\)
has index two and is normal.  The displayed mobility hypothesis therefore
excludes both \(2T\) and \(2O\), leaving \(2I\) as the unique finite type.
\end{proof}

\begin{corollary}[Generated five-core, ten-axis and local-normalizer packet]
\label{p1:cor:v237-five-core-ten-axis-packet}
Under Theorem~\ref{p1:thm:v237-mobile-tetrahedral-completion}, let
\(I:=\Gamma\cong2I\).  Then the conjugacy orbit of \(H\) in \(I\) consists
of exactly five tetrahedral cores
\[
 \mathcal H_I=\{H_1,\ldots,H_5\},\qquad H_i\cong2T,
\]
and, for \(i\ne j\),
\[
 C_{ij}:=H_i\cap H_j\cong C_6.
\]
The ten unordered pairs \(\{i,j\}\) give ten distinct \(C_6\) subgroups;
triple intersections reduce to the center \(\{\pm1\}\).  For every core,
\[
 O_i:=N_{\operatorname{Sp}(1)}(H_i)\cong2O,
 \qquad O_i\cap I=H_i.
\]
Moreover the union of the noncentral sectors of the ten \(C_{ij}\) has
cardinality \(40\).
\end{corollary}

\begin{proof}
Theorem~\ref{p1:thm:exceptional-normalizer-fork} gives
\(N_I(H)=H\), so orbit--stabilizer gives five conjugate cores.  Projectively,
two distinct \(A_4\) point stabilizers in \(A_5\) intersect in \(C_3\), whose
binary preimage is \(C_6\); three distinct point stabilizers intersect
trivially projectively and hence in \(\{\pm1\}\) upstairs.  The ambient
normalizer statement is again Theorem~\ref{p1:thm:exceptional-normalizer-fork}.
Each \(C_6\) contains four noncentral elements, and distinct axes meet only in
\(\{\pm1\}\), giving \(10\cdot4=40\).
\end{proof}

\begin{corollary}[Stabilizer characterization of the exceptional resolution roles]
\label{p1:cor:v237-exceptional-stabilizer-characterization}
For the generated packet of Corollary~\ref{p1:cor:v237-five-core-ten-axis-packet},
let \(\mathcal F_i\) be the four unoriented threefold-axis projectors of
\(H_i\), let \(\mathcal F_i^+\) denote a coherent oriented tetrahedral lift,
and let \(\mathcal A_{10}\) be the global set of ten axis projectors.  Then,
up to the common central double cover,
\[
 \boxed{
 \begin{aligned}
  \operatorname{Stab}_{\operatorname{Sp}(1)}(\mathcal F_i^+) &\cong 2T,\\
  \operatorname{Stab}_{\operatorname{Sp}(1)}(\mathcal F_i) &\cong 2O,\\
  \operatorname{Stab}_{\operatorname{Sp}(1)}(\mathcal A_{10}) &\cong 2I.
 \end{aligned}}
\]
Thus the local-core, local-normalizer and global-mobile group labels are
outputs of stabilizer computations on one quaternionic projector system.
\end{corollary}

\begin{theorem}[Ten-axis incidence--quadrupole exact sequence]
\label{p1:thm:v285-ten-axis-quadrupole-exact-sequence}
For the generated ten-axis packet of
Corollary~\ref{p1:cor:v237-five-core-ten-axis-packet}, choose unit vectors
\(u_{ij}\) on the unoriented axes and put
\[
 P_{ij}=u_{ij}\otimes u_{ij},
 \qquad
 Q_{ij}=P_{ij}-\frac13I_V,
 \qquad V=\operatorname{Im}\mathbb H.
\]
Then
\[
 \sum_{j\ne i}P_{ij}=\frac43I_V,
 \qquad
 \sum_{j\ne i}Q_{ij}=0,
 \qquad
 \sum_{i<j}P_{ij}=\frac{10}{3}I_V.
\]
Let \(E(K_5)\cong\mathbb R^{10}\) be the real edge module and define
\[
 (Da)_{ij}=a_i+a_j,
 \qquad
 \Phi(w)=\sum_{i<j}w_{ij}Q_{ij}.
\]
Then
\begin{equation}
\label{p1:eq:v285-ten-axis-quadrupole-exact-sequence}
 \boxed{
 0\longrightarrow\mathbb R^5
 \xrightarrow{D}\mathbb R^{10}
 \xrightarrow{\Phi}\operatorname{Sym}^2_0(V)
 \longrightarrow0.}
\end{equation}
In particular
\[
 \boxed{\mathbb R^{10}/D(\mathbb R^5)
 \cong\operatorname{Sym}^2_0(\operatorname{Im}\mathbb H),}
\]
so the ten-axis resolution data split canonically into five incidence/core
directions and a five-dimensional Clifford quadrupole quotient.
\end{theorem}

\begin{proof}
The local tight-frame identity is the standard tetrahedral projector identity;
summing over the five cores counts every pairwise axis twice and gives the
global identity.  Hence \(\Phi D=0\).  The map \(D\) is injective: if
\(a_i+a_j=0\) for every pair, three distinct indices force every \(a_i=0\).
It remains only to prove \(\operatorname{rank}\Phi=5\).

An exact generated coordinate representative is obtained with
\(\varphi=(1+\sqrt5)/2\) and the ten unit axes
\[
\begin{array}{ll}
 u_{12}=(1,1,1)/\sqrt3,&u_{13}=(1,1,-1)/\sqrt3,\\
 u_{14}=(1,-1,1)/\sqrt3,&u_{15}=(1,-1,-1)/\sqrt3,\\
 u_{23}=(\varphi^{-1},-\varphi,0)/\sqrt3,&
 u_{24}=(\varphi,0,-\varphi^{-1})/\sqrt3,\\
 u_{25}=(0,\varphi^{-1},-\varphi)/\sqrt3,&
 u_{34}=(0,\varphi^{-1},\varphi)/\sqrt3,\\
 u_{35}=(\varphi,0,\varphi^{-1})/\sqrt3,&
 u_{45}=(\varphi^{-1},\varphi,0)/\sqrt3.
\end{array}
\]
For a trace-free symmetric matrix use coordinates
\((Q_{11},Q_{22},Q_{12},Q_{13},Q_{23})\).  The five columns associated with
\((12),(13),(14),(23),(24)\) have determinant
\[
 \boxed{-\frac{8}{243}\ne0.}
\]
Thus \(\Phi\) is onto the five-dimensional space
\(\operatorname{Sym}^2_0(V)\).  Since \(\dim\ker\Phi=10-5=5\) and
\(D(\mathbb R^5)\subseteq\ker\Phi\) already has dimension five, the kernel is
exactly \(D(\mathbb R^5)\).
\end{proof}

\begin{remark}[Quadrupole quotient]
The quotient in
\eqref{p1:eq:v285-ten-axis-quadrupole-exact-sequence} is the
\(A_5\)-equivariant five-dimensional trace-free quadratic module.
The quotient map is the displayed weighted sum of projectors.
\end{remark}

\begin{remark}[Cardinality and orbit data]
\label{p1:rem:v237-forty-packet-clock-firewall}
The forty-element noncentral cyclic-axis packet has the counting measure
and conjugation action supplied by its ten distinct \(C_6\) subgroups.
Its cardinality and its representation structure are separate invariants.
\end{remark}

\subsection{Petersen incidence, signed Hopf cycles, and mirror operators}
\label{p1:sec:petersen-finite-resolution}
The results of this subsection are conditional on a specified binary-icosahedral
completion and its exact projective-axis coordinates; no physical branch is selected.

\paragraph{Finite datum and its provenance.}
Continue with the generated five-core and ten-axis packet of
Corollary~\ref{p1:cor:v237-five-core-ten-axis-packet} and the exact unit
representatives $u_{ij}$ in the proof of
Theorem~\ref{p1:thm:v285-ten-axis-quadrupole-exact-sequence}.  The five
binary tetrahedral subgroups are overlapping subgroups of one $2I$,
not five disjoint particles. All signs of $u_{ij}$ are gauge choices;
the fixed $\mathbb Q(\sqrt5)$ coordinate embedding is needed when the
Galois-conjugate representation is explicitly compared. The construction
is not canonical on an unmarked round three-sphere.

\begin{theorem}[Intrinsic Petersen graph and quadrupole spectrum]
\label{p1:thm:petersen-angular-quadrupole}
Let $\mathcal P$ have vertices $ij$ ($1\le i<j\le5$) and join $ij$ to
$kl$ when the underlying pairs are disjoint. For distinct vertices,
\[
 |\langle u_{ij},u_{kl}\rangle|^2=
 \begin{cases}5/9,&\{i,j\}\cap\{k,l\}=\varnothing,\\
 1/9,&|\{i,j\}\cap\{k,l\}|=1.
 \end{cases}
\]
Thus $\mathcal P$ is precisely the Petersen graph, recoverable as the
higher-overlap graph of the projective Hopf axes. If $A$ is its
$10\times10$ adjacency, $\mathsf J_{10}$ is the all-one matrix and
$D:\mathbb R^5\to\mathbb R^{10}$ is the pair-incidence map defined
above, then
\[
 A^2=2I_{10}-A+\mathsf J_{10},\qquad
 \operatorname{Spec}(A)=\{3^{(1)},1^{(5)},(-2)^{(4)}\},
\]
\[
 P_{\rm core}=D(D^{\mathsf T}D)^{-1}D^{\mathsf T},\quad
 P_4=P_{\rm core}-\mathsf J_{10}/10,\quad
 P_5=I_{10}-P_{\rm core}
     ={4I_{10}+2A-\mathsf J_{10}\over6}.
\]
Writing $Q_{ij}=u_{ij}u_{ij}^{\mathsf T}-I_3/3$, their Gram matrix is
$\big(\operatorname{tr}(Q_vQ_w)\big)_{vw}=\tfrac43P_5$.
In particular, the Petersen eigenvalue-$1$ space is exactly the
five-dimensional quadrupole summand of the preceding exact sequence.
\end{theorem}
\begin{proof}
Substitution of the ten displayed vectors gives the two overlap
values, with 15 disjoint and 30 intersecting vertex pairs. Every
vertex has three neighbors, adjacent vertices have zero common
neighbors and distinct nonadjacent vertices have one, proving the
adjacency identity and its eigenvalues. The eigenvalue $-2$ subspace
is the four-dimensional centered core-incidence representation;
the eigenvalue $3$ line is constant. The remaining rank-five projector
is the stated polynomial in $A$. On diagonal axis pairs the trace-free
quadrupole Gram is $2/3$; on distinct pairs it is $2/9$ for disjoint
labels and $-2/9$ for intersecting labels, exactly $4P_5/3$.
\end{proof}

\begin{theorem}[Signed ten-axis module and order-four edge flux]
\label{p1:thm:petersen-signed-and-flux}
Choose representatives $u_v$ and define
\[
 S_{vw}=\begin{cases}\operatorname{sgn}(u_v\cdot u_w),&v\sim w,\\0,&\text{otherwise.}\end{cases}
 \qquad U=[u_v]_{v\in V(\mathcal P)}.
\]
Then
\[
 S^3=5S,\quad US=\sqrt5\,U,\quad UU^{\mathsf T}=\tfrac{10}3I_3,
 \quad\operatorname{Spec}S=\{(\sqrt5)^{(3)},0^{(4)},(-\sqrt5)^{(3)}\}.
\]
Applying the number-field automorphism $\sigma(\sqrt5)=-\sqrt5$ to the
specified exact coordinates yields $U^\sigma$, and
\[
 P_{\mathbf3}=\tfrac3{10}U^{\mathsf T}U
     =\tfrac12(S^2/5+S/\sqrt5),\qquad
 P_{\mathbf3'}=\tfrac3{10}(U^\sigma)^{\mathsf T}U^\sigma
     =\tfrac12(S^2/5-S/\sqrt5),
 \quad U(U^\sigma)^{\mathsf T}=0.
\]
The signed axis module is $\mathbf3\oplus\mathbf4\oplus\mathbf3'$.
Give each Petersen edge an orientation $v\to w$ and define
\[
 k_{vw}=\tfrac32 S_{vw}(u_v\times u_w),\quad
 k_{wv}=-k_{vw},\qquad K=[k_e]_{e\in E(\mathcal P)}.
\]
The $15$ vectors $k_e$ are unit imaginary-quaternion half-turn axes
in $2I$, independent of projective sign gauges. If $\partial$ is the
oriented vertex-edge boundary matrix, then
\[
 KK^{\mathsf T}=K^\sigma(K^\sigma)^{\mathsf T}=5I_3,\qquad
 K(K^\sigma)^{\mathsf T}=0,\qquad
 K\partial^{\mathsf T}=K^\sigma\partial^{\mathsf T}=0.
\]
At each vertex the three outward edge vectors sum to zero and have
mutual inner product $-1/2$, forming one transverse real $A_2$ trimer.
\end{theorem}
\begin{proof}
The signed edges are the fifteen higher-overlap pairs in the previous
theorem; changing a representative axis sign conjugates $S$ by a
diagonal sign matrix. Exact coordinate multiplication gives $S^3=5S$
and $US=\sqrt5U$. The tight-frame sums fix both rank-three spectral
projectors and their Galois-conjugate orthogonality. Since
$|u_v\times u_w|=2/3$ on each edge, $k_e$ has unit length. Its
invariance under independent vertex sign changes follows because
both $S_{vw}$ and $u_v\times u_w$ change sign. Equivariance and the
unit-vector trace yield $KK^{\mathsf T}=5I_3$. The vertex-wise sum is
$\tfrac32u_v\times(\sum_{w\sim v}S_{vw}u_w)=0$; its three unit
transverse vectors therefore form the stated $120^\circ$ trimer.
The Galois-conjugate statements follow in the fixed exact coordinate
frame; inequivalence of the two $A_5$ three-spaces gives the mixed zero.
\end{proof}

\begin{theorem}[Full fifteen-edge Hodge--Clifford resolution]
\label{p1:thm:petersen-fifteen-resolution}
The four operators
\[
 \Pi_4=\tfrac15\partial^{\mathsf T}P_4\partial,\quad
 \Pi_5=\tfrac12\partial^{\mathsf T}P_5\partial,\quad
 \Pi_3=\tfrac15K^{\mathsf T}K,\quad
 \Pi_{3'}=\tfrac15(K^\sigma)^{\mathsf T}K^\sigma
\]
are pairwise orthogonal idempotents of respective ranks $4,5,3,3$;
their sum is $I_{15}$. As real $A_5$-modules,
\[
 \boxed{\mathbb R^{15}=\mathbf4\oplus\mathbf5\oplus\mathbf3\oplus\mathbf3',}
\quad
 0\longrightarrow\mathbb R\longrightarrow\mathbb R^{10}
 \xrightarrow{\partial^{\mathsf T}}\mathbb R^{15}
 \xrightarrow{(K,K^\sigma)}\mathbb R^3\oplus(\mathbb R^3)^\sigma
 \longrightarrow0
\]
is exact. Thus the commutant of this unmarked real $A_5$ edge module
is spanned by the four displayed projectors. This is a completeness
statement for the specified \emph{linear equivariant edge domain}, not
for every composite operator or physical action.
\end{theorem}
\begin{proof}
The identity $\partial\partial^{\mathsf T}=3I_{10}-A$ gives
squared nonzero singular values $5$ and $2$ on the respective
four- and five-dimensional gradients. The previous flux identities
give mutually orthogonal rank-three projectors annihilating all
gradients. Their ranks add to fifteen, proving the resolution and
exactness. The four irreducible real summands are pairwise inequivalent.
\end{proof}

\begin{theorem}[Pentagon holonomy and reciprocal cycle-frame spectrum]
\label{p1:thm:petersen-pentagon-reciprocity}
The signed-axis voltage $s_{vw}=S_{vw}$ gives gauge-invariant loop
values $W(\gamma)=\prod_{vw\in\gamma}s_{vw}\in\{\pm1\}$. The twelve
five-cycles split into six with $W=+1$ and six with $W=-1$.
For the oriented cycle-column matrices $B_\pm\in\mathbb R^{15\times6}$
of these two families put $H_\pm=B_\pm B_\pm^{\mathsf T}$ and
$P_C=\Pi_3+\Pi_{3'}$. With the sign designation fixed by $K$,
\[
 H_++H_-=10P_C,\quad H_+H_-=20P_C,\quad
 (H_+-H_-)^2=20P_C,
\]
\[
 H_+=(5-\sqrt5)\Pi_3+(5+\sqrt5)\Pi_{3'},\quad
 H_-=(5+\sqrt5)\Pi_3+(5-\sqrt5)\Pi_{3'}.
\]
Each of the six perfect matchings complements two disjoint pentagons
$c_+,c_-$ of opposite voltage, and their quaternionic cycle fluxes obey
\[
 Kc_-=\pm\varphi Kc_+,\qquad
 K^\sigma c_-=\pm\varphi^{-1}K^\sigma c_+,
 \qquad\varphi=(1+\sqrt5)/2.
\]
The signs depend on cycle orientations, while the amplitudes are exact.
The signed double cover is the dodecahedral graph; its $+1$ base
pentagons lift to closed pentagons and its $-1$ pentagons to decagons.
\end{theorem}
\begin{proof}
The dodecahedral antipodal quotient identifies precisely the ten
unoriented axes. Its facial cycles yield six positive-voltage base
pentagons; the remaining six do not close in one sheet. Enumerating
the two six-cycle bases in the exact $15$-edge incidence representation
gives $\operatorname{Spec}(B_\pm^{\mathsf T}B_\pm)=
\{(5+\sqrt5)^{(3)},(5-\sqrt5)^{(3)}\}$ and the displayed commuting
identities. Direct application of $K$ to complementary cycles gives
parallel vectors with norm squares $\tfrac52(5\mp\sqrt5)$, hence the
ratio $\varphi$ and its Galois reciprocal.
\end{proof}

\begin{proposition}[Real binary edge holonomy and coloring obstruction]
\label{p1:prop:petersen-binary-and-snark}
The unsigned edge half-turns $q_{vw}=\tfrac32u_v\times u_w$ form a
flat $2I$-valued connection: $q_{vw}=h_vh_w^{-1}$ for suitable
$h_v\in2I$. The signed edge matrices satisfy
$k_{vw}=s_{vw}h_vh_w^{-1}$ and every ordered closed-loop product is
$W(\gamma)\in\{\pm1\}\subset Z(2I)$. Thus this particular signed
connection has a nontrivial central $\mathbb Z_2$ holonomy, while the
unsigned connection is flat. The Petersen graph has chromatic index
$4$: its three locally defined edges cannot be consistently assigned
three globally fixed edge colors.
\end{proposition}
\begin{proof}
Choose a spanning tree and propagate vertex quaternions. Exact
multiplication of the unit half-turns around every fundamental
cycle gives $+1$ for the unsigned connection, proving existence of
the gauges $h_v$. Signed factors are central and commute, yielding
$W(\gamma)$. For the coloring claim, one color class in any proper
three-edge coloring of a cubic graph would be a perfect matching.
Its complement here consists of two odd five-cycles, which cannot
alternate between the two remaining colors.
\end{proof}

\begin{theorem}[Two-parent transverse bisector and core-local mirror simplex]
\label{p1:thm:petersen-bisector-simplex}
At $u_{ij}$ the three signed outward edge vectors form a trimer $N_u$.
The projected root-line trimers from the two tetrahedral parents
$H_i,H_j$ are rotated relative to $N_u$ through angles $\pm\alpha$,
where
\[
 \cos^2\alpha={9+3\sqrt5\over16},\qquad
 \cos(2\alpha)={1+3\sqrt5\over8}.
\]
In particular, $N_u$ is the exact angular bisector of the two local
parent trimers after compatible projective matching.
Let $U_i$ collect the four $u_{ij}$ with $j\ne i$ and set
\[
 J_i={3\over4}U_i^\sigma U_i^{\mathsf T}:\mathbf3\to\mathbf3'.
\]
The five marked-core intertwiners satisfy
\[
 J_i^{\mathsf T}J_i=I_3,\quad\det J_i=1,\quad
 \sum_{i=1}^5J_i=0,\quad
 \operatorname{tr}(J_i^{\mathsf T}J_j)=
 \begin{cases}3,&i=j,\\-3/4,&i\ne j.\end{cases}
\]
They form a regular four-simplex in $\operatorname{Hom}(\mathbf3,\mathbf3')$.
There is no nonzero global $A_5$-intertwiner $\mathbf3\to\mathbf3'$;
a marked core supplies an intertwiner only for its $A_4$ stabilizer.
\end{theorem}
\begin{proof}
Projection of the three other local tetrahedral axes from each
parent into $u_{ij}^\perp$ gives two mutually rotated $A_2$ root
frames. Their exact dot products with the signed Petersen trimer
are $(9+3\sqrt5)/16$ and are opposite rotations, proving the
bisector assertion. Each $U_i$ is a tetrahedral tight frame:
$U_iU_i^{\mathsf T}=4I_3/3$ and its Galois conjugate has the same
Gram matrix; the displayed $J_i$ is therefore orthogonal, with
positive determinant in the given coordinate frame. It sends the
four parent axes to their Galois partners and hence intertwines the
stabilizing $A_4$. The relation $U^\sigma U^{\mathsf T}=0$ implies
$\sum J_i=0$. Equivariance and the common squared norm $3$
force every distinct pair product trace to be $-3/4$.
\end{proof}

\begin{theorem}[All core subsets and their axial polar response]
\label{p1:thm:petersen-subset-polar}
For $T\subseteq\{1,\dots,5\}$ define $J_T=\sum_{i\in T}J_i$.
Then $J_{T^c}=-J_T$. Its singular values depend only on $|T|$:
\[
 \begin{array}{c|ccc}
 |T|&0,5&1,4&2,3\\ \hline
 \operatorname{Spec}(J_T^{\mathsf T}J_T)
   &\{0,0,0\}&\{1,1,1\}&\{4,1/4,1/4\}
 \end{array}.
\]
For a selected pair $T=\{i,j\}$,
\[
 \boxed{J_T^{\mathsf T}J_T
 =4u_{ij}u_{ij}^{\mathsf T}
 +\tfrac14(I_3-u_{ij}u_{ij}^{\mathsf T}),}
 \quad |J_T|=2u_{ij}u_{ij}^{\mathsf T}
 +\tfrac12(I_3-u_{ij}u_{ij}^{\mathsf T}).
\]
Thus complementary two- and three-core markings have opposite linear
maps but identical quadratic metrics. This is a finite-subgroup
selection rule, not a constituent-occupancy or particle-mass theorem.
\end{theorem}
\begin{proof}
The simplex identity gives $J_{T^c}=-J_T$. For two selected
cores the matrix $J_i^{\mathsf T}J_j$ is a rotation fixing their
common Hopf axis. Its trace is $-3/4$, hence its transverse rotation
angle has cosine $-7/8$. Therefore $(J_i+J_j)^{\mathsf T}
(J_i+J_j)$ has eigenvalue $4$ axially and $1/4$ twice transversely.
Complements and the singleton cases exhaust all $32$ subsets.
\end{proof}

\begin{theorem}[The nine-dimensional mirror-operator envelope]
\label{p1:thm:petersen-mirror-envelope}
The ten rank-one operators
$M_{ij}=u_{ij}^\sigma u_{ij}^{\mathsf T}:\mathbf3\to\mathbf3'$
obey $\sum_{i<j}M_{ij}=0$ and span
$\operatorname{Hom}_{\mathbb R}(\mathbf3,\mathbf3')$ of dimension nine.
Their Hilbert--Schmidt Gram matrix is
\[
 G_M={8I_{10}+\mathsf J_{10}-6A\over9}
       ={20\over9}P_4+{2\over9}P_5.
\]
Consequently
$\operatorname{Hom}(\mathbf3,\mathbf3')\simeq\mathbf4\oplus\mathbf5$.
For any $Q\in\operatorname{Sym}^2_0(\mathbb R^3)$ put
\[
 \mathcal B(Q)={3\over4}\sum_{i<j}
  \operatorname{tr}(QQ_{ij})M_{ij}.
\]
Then
\[
 \|\mathcal B(Q)\|_{\rm HS}^2={1\over6}\|Q\|_{\rm HS}^2,
 \quad\operatorname{tr}(J_m^{\mathsf T}\mathcal B(Q))=0,
 \quad(J_i+J_j)^{\mathsf T}\mathcal B(Q_{ij})={1\over3}Q_{ij}.
\]
The first $\mathbf4$ part is generated by marked cores $J_m$ and
the complementary $\mathbf5$ part by quadrupoles $\mathcal B(Q)$.
\end{theorem}
\begin{proof}
The Gram of $M_v$ and $M_w$ is
$(u_v^\sigma\cdot u_w^\sigma)(u_v\cdot u_w)$, which takes the
values $1,1/9,-5/9$ for identical, label-intersecting and disjoint
pairs, proving the displayed polynomial in $A$. It has one zero,
four eigenvalues $20/9$ and five eigenvalues $2/9$, giving rank nine.
The five local $J_m$ form the four-dimensional core summand.
$\mathcal B$ uses only the orthogonal five-dimensional quadrupole
summand, so the norm and trace orthogonality follow from its Gram
scalar; the matched-pair round trip follows by substituting the
coefficients and the axial polar identity.
\end{proof}

\begin{proposition}[Harmonic-cubic portal and the complete edge-space interface]
\label{p1:prop:petersen-cubic-gradient-interface}
\leavevmode\par
For $Q\in\operatorname{Sym}^2_0(\mathbb R^3)$ and $v\in\mathbb R^3$,
\[
 h_{Q,v}(x)=(x^{\mathsf T}Qx)(v\cdot x)
 -{2\over5}|x|^2(Qv\cdot x)
\]
is harmonic. Its evaluation on the ten axis representatives lies in
$\mathbf3'\oplus\mathbf4$, and the mirror projection satisfies
\[
 P_{\mathbf3'}\operatorname{ev}_{10}(h_{Q,v})
 ={2\over5}(U^\sigma)^{\mathsf T}\mathcal B(Q)v.
\]
The oriented edge module is explicitly isomorphic, as an $A_5$
module, to the sum of all nine ordinary-to-mirror operator directions
and six cycle-flux directions:
\[
 \mathbb R^{15}\simeq\operatorname{Hom}(\mathbf3,\mathbf3')
       \oplus\mathbf3\oplus\mathbf3'.
\]
In particular, with $F$ the $9\times10$ column-matrix of the $M_{ij}$
written in the Frobenius basis, the edge-to-operator map
\[
 \mathcal L=F(P_4/5+P_5/2)\partial
\]
annihilates cycles and obeys
$\mathcal L^*\mathcal L=\tfrac49\Pi_4+\tfrac19\Pi_5$.
\end{proposition}
\begin{proof}
Taking the Euclidean Laplacian gives
$\Delta[(x^{\mathsf T}Qx)(v\cdot x)]=4(Qv\cdot x)$ and
$\Delta[|x|^2(Qv\cdot x)]=10(Qv\cdot x)$.
Exact evaluation in the ten inherited axis coordinates yields the
projection formula. The gradient-to-mirror Gram follows by combining
$F^{\mathsf T}F=(20/9)P_4+(2/9)P_5$ with
$\partial\partial^{\mathsf T}=5P_4+2P_5$ on the centered vertex
space. Its kernel is the six-dimensional cycle space, resolved by
$K,K^\sigma$ in the preceding theorem.
\end{proof}


\subsubsection{Shortest-edge realization of the intrinsic Petersen packet}
\label{p1:sec:v371-petersen-600cell-geometry}

\begin{theorem}[Dodecahedral lift and the induced fifty-vertex total graph]
\label{p1:thm:v371-petersen-600cell-total-cover}
In the fixed binary-icosahedral unit-quaternion realization of the
$600$-cell, join $q,r\in2I$ precisely when
$\operatorname{Re}(q^{-1}r)=\varphi/2$, with
$\varphi=(1+\sqrt5)/2$. Put
\[
 D=\{q\in2I:\operatorname{Re}q=1/2\},\qquad
 Z=\{q\in2I:\operatorname{Re}q=0\}.
\]
Then $|D|=20$ and $|Z|=30$. The induced graph on $D$ is the
$dodecahedral$ graph, with $20$ vertices and $30$ edges; its antipodal
quotient under $q\mapsto q^{-1}$ is the intrinsic Petersen graph of
Theorem~\ref{p1:thm:petersen-angular-quadrupole}.
For a dodecahedral edge $e=\{1/2+w,1/2+z\}$, define
$\mu_e=(w+z)/\varphi\in Z$. The map $e\mapsto\mu_e$ is a bijection.
The induced graph on $Z$ is the line graph of the dodecahedron; cross
edges between $D$ and $Z$ are exactly vertex--edge incidences. Hence
\[
 \boxed{600[D\cup Z]\cong T(D_{20}),\qquad
 T(D_{20})/(q\sim q^{-1})\cong T(P_{10}).}
\]
The former is an induced $50$-vertex, $150$-edge, $6$-regular graph; the
latter has $10+15=25$ vertices and $75$ edges. For every $q\in D$, the
three incident midpoint vertices together with $q$ span a regular
$600$-cell tetrahedron. There are twenty such tetrahedra, paired under
quaternion inversion. This is a nearest-edge \emph{cover/quotient}
realization, not a ten-vertex Petersen subgraph of the $600$-cell.
\end{theorem}
\begin{proof}
Write $q=1/2+(\sqrt3/2)u$ and $r=1/2+(\sqrt3/2)v$.
The $600$-cell edge condition is equivalent to
$u\cdot v=\sqrt5/3$, the dodecahedral nearest-neighbor relation.
For an edge, $|w|^2=|z|^2=3/4$ and $w\cdot z=\sqrt5/4$;
therefore $\mu_e$ has norm one. Substitution in the exact $2I$
coordinate list shows that all thirty midpoints are distinct members
of $Z$, and that $\mu_e\sim\mu_f$ precisely when $e,f$ share an
endpoint. The cross-shell edge condition picks exactly the three
midpoints incident at each $q$. These edge counts exhaust the induced
adjacency, and each vertex with its three midpoint neighbors forms
a four-clique. Inversion is a fixed-point-free graph involution on the
fifty vertices, with quotient the asserted total graph. The finite
coordinate and count certificate is recorded in
Paper~I-D, Section~\ref{p1c:sec:v371-petersen-600cell-certificate}.
\end{proof}

\begin{proposition}[Quaternionic edge frames and four angular channel orbits]
\label{p1:prop:v371-petersen-q8-angular}
For $e=\{1/2+w,1/2+z\}$ as above set
\[
 m_e={w+z\over\varphi},\qquad
 d_e=\varphi(w-z),\qquad k_e=2w\times z.
\]
These are orthogonal unit imaginary elements of $2I$ and satisfy
$m_e^2=d_e^2=k_e^2=-1$, $m_ed_e=-k_e$ and $d_em_e=k_e$.
They generate a subgroup $Q_8$; the thirty dodecahedral edges realize
fifteen projective channels, grouped three per each of the five
pre-existing tetrahedral $Q_8$ cores. For every projective $C_6$ axis
$u$ and every projective midpoint channel $e$, exactly four $A_5$
angular orbits occur:
\[
\begin{array}{c|c|c|c}
\text{relation}&\text{number per axis}&|u\cdot m_e|&
\text{distance in the 61-orbit quotient}\\\hline
\text{incident}&3&\varphi/\sqrt3&1\\
\text{near}&6&1/\sqrt3&2\\
\text{cross B}&3&(\varphi-1)/\sqrt3&2\\
\text{cross A}&3&0&3
\end{array}
\]
The crossed pairing classes are exchanged by odd relabelings of the five
abstract cores, not by the realized rotational $A_5$ action. No
physical chirality or decay interpretation follows from this table.
\end{proposition}
\begin{proof}
The quaternionic identities follow from
$w\cdot z=\sqrt5/4$ and the vector-product identity; direct
multiplication gives $m_ed_e=-2w\times z=-k_e$. Exact enumeration
of the fixed $2I$ coordinates checks $m_e,d_e,k_e\in2I$, the five
$Q_8$ cores and the indicated angular counts. Distances are computed
from the full $600$-cell adjacency after inversion, not from an
independently postulated coupling graph. See
Section~\ref{p1c:sec:v371-petersen-600cell-certificate}.
\end{proof}

\subsubsection{Global tangent sequence and rooted Clifford comparisons}
\label{p1:sec:v371-petersen-tangent-sequence}

\begin{theorem}[Projective trimer uniqueness and exact ten-axis tangent sequence]
\label{p1:thm:v371-petersen-tangent-exact-sequence}
For each of the ten unoriented $C_6$ axes $u=ij$, let
$V_u=u^\perp\subset\operatorname{Im}\mathbb H$ and let
$\mathcal V=\bigoplus_{u=1}^{10}V_u$ (real dimension twenty).
The three outward flux directions $k_{u,e}$ attached to the incident
Petersen edges obey $\sum_{e\ni u}k_{u,e}=0$ and
$k_{u,e}\cdot k_{u,f}=-1/2$ for $e\ne f$.
The stabilizer of the \emph{projective} axis is $D_3$; hence
\[
 \dim_{\mathbb R}\operatorname{Hom}_{D_3}(\mathbb R^3,V_u)=1,
 \qquad F_u(x)=\sum_{e\ni u}x_e k_{u,e}
\]
spans its real intertwiner space. This uniqueness would fail if only the
oriented-axis $C_3$ symmetry were imposed.
Define the \emph{scalar/unoriented} channel map and two global readouts
\[
 F:\mathbb R^{15}_{\rm scalar}\to\mathcal V,
 \quad (Fx)_u=\sum_{e\ni u}x_e k_{u,e},
 \qquad
 M_3(t)=\sum_ut_u,
 \quad M_{3'}(t)=\sum_{u=ij}(J_i+J_j)t_u.
\]
Here $J_i$ are \emph{exactly} the five mirror-core intertwiners of
Theorem~\ref{p1:thm:petersen-bisector-simplex}, and $M=(M_3,M_{3'})$.
For the unsigned $10\times15$ Petersen incidence $B$,
\[
 F^\dagger F=3I_{15}-{1\over2}B^{\mathsf T}B,
 \quad\operatorname{Spec}(F^\dagger F)
 =\{0^{(1)},1^{(5)},(5/2)^{(4)},3^{(5)}\},
\]
\[
 MF=0,\qquad MM^\dagger=\operatorname{diag}
 \bigl((20/3)I_3,(5/3)I_{3'}\bigr).
\]
Therefore the following sequence is exact:
\[
 \boxed{0\longrightarrow\mathbb R
 \xrightarrow{\mathbf1_{15}}\mathbb R^{15}_{\rm scalar}
 \xrightarrow{F}\mathcal V\xrightarrow{M}
 \mathbf3\oplus\mathbf3'\longrightarrow0.}
\]
In particular $\ker F=\mathbb R\mathbf1_{15}$,
$\operatorname{im}F=\ker M$ and
$\mathcal V\cong(\mathbf4\oplus2\mathbf5)_{\rm loc}
\oplus(\mathbf3\oplus\mathbf3')_{\rm rec}$ as real $A_5$ modules.
This is \emph{not} the already-established oriented-edge Hodge sequence
of Theorem~\ref{p1:thm:petersen-fifteen-resolution}; the two input
representation types and their rank statements are different.
\end{theorem}
\begin{proof}
The three channel directions have Gram matrix
$(3/2)I_3-(1/2)\mathbf1\mathbf1^\mathsf T$. The projective $D_3$
characters $(3,0,1)$ on channel labels and $(2,-1,0)$ on $V_u$
give a one-dimensional intertwiner space; the displayed trimer map
is nonzero and $D_3$-equivariant. Each pair of distinct edge columns
of $F$ meets in at most one vertex, proving the Gram identity.
The Petersen spectrum gives the displayed four eigenvalue blocks and
$\ker F=\mathbb R\mathbf1_{15}$. The outgoing endpoint signs cancel
in $M_3F$; for the mirror readout the exact core identity
$(J_i+J_j)k_e=(J_k+J_l)k_e$ on every edge $ij\sim kl$ proves
$M_{3'}F=0$. Its verification and the tight-frame sums give the
stated $MM^\dagger$. The latter has rank six, so
$\dim\ker M=14=\operatorname{rank}F$, proving exactness. The
irreducible decomposition follows from the unsigned Petersen
permutation module $\mathbf1\oplus\mathbf4\oplus2\mathbf5$
and the two mirror receiver triplets. All fifteen edge identities
are certified with the inherited core matrices in
Section~\ref{p1c:sec:v371-petersen-600cell-certificate}.
\end{proof}

\begin{theorem}[Canonical recovery and minimal two-axis support]
\label{p1:thm:v371-petersen-source-reconstruction}
Put $X=B^{\mathsf T}B$, $G=F^\dagger F$ and
\[
 G^+=-{(X-6I)(13X^2-2X+25I)\over450}.
\]
Then $G^+$ is the Moore--Penrose inverse of $G$.
Writing $P_3=(3/20)M_3^\dagger M_3$,
$P_{3'}=(3/5)M_{3'}^\dagger M_{3'}$ and
$P_{\rm loc}=I_{\mathcal V}-P_3-P_{3'}$,
\[
 FG^+F^\dagger=P_{\rm loc},\qquad
 \widehat x=G^+F^\dagger t,\qquad
 \|t-F\widehat x\|^2=
 {3\over20}\|M_3t\|^2+{3\over5}\|M_{3'}t\|^2.
\]
Among the $45$ unordered pairs of distinct axes, the $15$ Petersen
edges are \emph{exactly} those pairs that support a nonzero vector in
$\ker M$ with no other tangent components. For a single-axis source,
the squared norm fractions in $(P_{\rm loc},P_3,P_{3'})$ are
$(7/10,3/20,3/20)$, respectively.
\end{theorem}
\begin{proof}
Use the four distinct eigenvalues of $X$ implied by
$BB^\dagger=3I+A_{\mathcal P}$ to evaluate the stated polynomial
on each $G$ eigenspace; it equals $1/\lambda$ on the positive
ones and zero on $\ker G$. Hence $FG^+F^\dagger$ is the orthogonal
projector onto $\operatorname{im}F=\ker M$. The orthogonal receiver
projectors give the residual formula. A tangent supported on $u,v$
with $Mt=0$ must be proportional to $(u\times v,-u\times v)$;
mirror compatibility becomes
$\bigl((J_i+J_j)-(J_k+J_l)\bigr)(u\times v)=0$.
The exact ten-axis matrix test has nonzero solutions for precisely the
fifteen disjoint-pair labels. The single-axis formula follows from
$T_u^\dagger T_u|_{V_u}=I_{V_u}/4$ and the receiver norms.
\end{proof}

\begin{theorem}[Marked $2T$ root leakage and five-root sum rules]
\label{p1:thm:v371-petersen-rooted-leakage}
Let $P_{\rm rec}=P_3+P_{3'}$ and
$P_{\rm loc}=I_{\mathcal V}-P_{\rm rec}$. For each core $H_m\cong2T$,
let $R_m$ project onto the four $C_6$-axis tangent fibres in $H_m$
(rank eight), and define the \emph{geometric} operator
\[
 C_m=P_{\rm rec}R_mP_{\rm loc}:
 \mathcal V_{\rm loc}\longrightarrow\mathcal V_{\rm rec}.
\]
In the normalized receiver basis, the compression of $R_m$ is
\[
 A_m=\begin{pmatrix}
 \tfrac25I_3&\tfrac1{10}J_m^{\mathsf T}\\
 \tfrac1{10}J_m&\tfrac25I_{3'}
 \end{pmatrix},\qquad
 \operatorname{Spec}(A_m)=\{(1/2)^{(3)},(3/10)^{(3)}\}.
\]
Hence
\[
 \boxed{\operatorname{rank}C_m=6,\quad
 \operatorname{Spec}_{>0}(C_mC_m^\dagger)
 =\{(1/4)^{(3)},(21/100)^{(3)}\}.}
\]
The five core operators obey the exact collective identities
\[
 \boxed{\sum_{m=1}^5C_m=0,\quad
 \sum_{m=1}^5C_mC_m^\dagger={23\over20}P_{\rm rec},\quad
 \sum_{m=1}^5C_m^\dagger C_m
 ={3\over10}FF^\dagger-{3\over20}P_{\rm loc}.}
\]
The final sum is strictly positive on $\operatorname{im}F$, with
spectrum $(3/20)^{(5)},(3/5)^{(4)},(3/4)^{(5)}$, so every nonzero
local source is detected by at least one marked-core operator.
These claims concern the specified root projections, \emph{not}
universal $A_4$-equivariant action kernels or physical decay vertices.
\end{theorem}
\begin{proof}
On the four axes of a given $H_m$, the exact core-star sums are
$\sum P_u=8I_3/3$, $\sum T_uP_uT_u^\dagger=2I_{3'}/3$ and
$\sum P_uT_u^\dagger=J_m^\dagger/3$, with $T_{ij}=J_i+J_j$.
Normalize the two receiver inclusions by
$\sqrt{3/20}\,M_3^\dagger$ and
$\sqrt{3/5}\,M_{3'}^\dagger$ to obtain $A_m$.
Since $R_m^2=R_m$, the nonzero leakage eigenvalues are those of
$A_m-A_m^2$, namely $1/4$ and $21/100$, each threefold.
Each axis lies in exactly two core stars, hence
$\sum_mR_m=2I_{\mathcal V}$ and $\sum_mC_m=0$.
Using $\sum_mJ_m=0$ in the receiver compression yields
$\sum_mC_mC_m^\dagger=23P_{\rm rec}/20$.
The source-side identity follows by expanding the same five diagonal
root masks with the exact $F$ Gram; its equality entry by entry over
$\mathbb Q(\sqrt5)$ is supplied in Paper~I-D's certificate.
The displayed spectrum follows from the $F^\dagger F$ spectrum.
\end{proof}

\begin{corollary}[Five-core leakage operators form an exact operator simplex]
\label{p1:cor:v3026-five-core-hs-gram}
For the five \emph{geometric} maps
$C_m=P_{\rm rec}R_mP_{\rm loc}:\mathcal V_{\rm loc}\to
\mathcal V_{\rm rec}$ of
Theorem~\ref{p1:thm:v371-petersen-rooted-leakage}, with the inherited
real Euclidean inner products, the complete Hilbert--Schmidt Gram is
\[
 \boxed{\operatorname{tr}(C_m^{\mathsf T}C_n)
    ={69\over40}\left(\delta_{mn}-{1\over5}\right).}
\]
For real or complex coefficients $a_m$ (using the fixed complexified
maps in the latter case),
\[
 \boxed{\left\|\sum_{m=1}^5a_mC_m\right\|_{\rm HS}^2
       ={69\over200}\sum_{m<n}|a_m-a_n|^2.}
\]
Thus the sole relation among the five maps is their uniform sum;
the relative four-dimensional amplitude sector has Gram eigenvalue
$69/40$. This is an \emph{operator-overlap $K_5$ Dirichlet form},
not a physical hadron transition matrix or decay-rate law.
\end{corollary}
\begin{proof}
The six nonzero squared singular values in the preceding theorem give
$\|C_m\|_{\rm HS}^2=3(1/4+21/100)=69/50$.
The inherited $A_5$ action permutes the five cores two-transitively and
commutes with $P_{\rm loc},P_{\rm rec}$; hence all off-diagonal trace
products have one common value $c$. Taking the inner product of
$\sum_m C_m=0$ with $C_n$ yields $69/50+4c=0$, so
$c=-69/200$. This is the Gram of a regular centered simplex.
The identity $\sum_{m<n}|a_m-a_n|^2
=5\sum_m|a_m|^2-|\sum_m a_m|^2$ proves the formula and rank.
\end{proof}

% BEGIN v3.071 exact five-core resistance closure (mathematical comparison only).
\begin{theorem}[Canonical five-core resistance and reversible comparison kernel]
\label{p1b:thm:v371-five-core-resistance-rigidity}
Let $C:\mathbb R^5\to\operatorname{HS}(\mathcal V_{\rm loc},\mathcal V_{\rm rec})$
be the \emph{existing} core-leakage analysis map $Ce_m=C_m$ from
Corollary~\ref{p1:cor:v3026-five-core-hs-gram}.  Its positive
coefficient-space Gram operator has the exact, parameter-free form
\[
 \boxed{L:=C^{\mathsf T}C=\frac{69}{40}P_\perp,
 \qquad x^{\mathsf T}Lx=\frac{69}{200}
                   \sum_{1\le a<b\le5}(x_a-x_b)^2.}
\]
It is the graph Laplacian of the complete five-core graph with every
edge conductance $c_{ab}=69/200$. Its Moore--Penrose inverse and every
pairwise effective resistance are
\[
 \boxed{L^+=\frac{40}{69}P_\perp,\qquad
 \mathcal R_{ab}=(e_a-e_b)^{\mathsf T}L^+(e_a-e_b)
   =\frac{80}{69}\quad(a\ne b).}
\]
The canonical mathematical heat semigroup and its entries are
\[
 e^{-tL}=P_0+e^{-69t/40}P_\perp,\qquad P_0=J_5/5,
\]
\[
 (e^{-tL})_{aa}=\frac{1+4e^{-69t/40}}5,\qquad
 (e^{-tL})_{ab}=\frac{1-e^{-69t/40}}5\quad(a\ne b).
\]
These entries are positive for $t>0$, symmetric and stochastic. Every
real self-adjoint $A_5$-equivariant five-core operator annihilating the
constant vector is proportional to $P_\perp$; the preceding value of
its scale is selected by the original geometric maps $C_m$.
The comparison semigroup is \emph{reversible and not a physical decay
law}; no channel energy, source current, open continuum, or lifetime has
been inferred.
\end{theorem}
\begin{proof}
The source Gram identity gives $C^{\mathsf T}C=(69/40)P_\perp$.
Since $P_\perp=I_5-J_5/5$, its diagonal entries are $4/5$ and its
off-diagonal entries $-1/5$; hence every off-diagonal entry of $L$
is $-69/200$, proving the Laplacian and Dirichlet identities.
The spectrum is $\{0^{(1)},(69/40)^{(4)}\}$, which gives the
pseudoinverse, resistance, exponential and stochastic formulas.
The core permutation representation of $A_5$ decomposes over
$\mathbb R$ as the trivial line plus its real irreducible
four-dimensional augmentation representation.  A real self-adjoint
intertwiner is scalar on that augmentation and scalar on the trivial
line; annihilating constants sets the latter to zero.  Symmetry
of the kernel gives detailed balance with the uniform stationary
measure, and therefore does not provide a preferred decay direction.
\end{proof}
% END v3.071 exact five-core resistance closure.

% BEGIN v3.072 aligned five-core resistance/orientation identity (mathematical).
\begin{corollary}[Exact resistance-to-directed-overlap bridge on the same five-core labels]
\label{p1b:cor:v372-resistance-orientation-bridge}
In the five-core coefficient space of
Theorem~\ref{p1b:thm:v371-five-core-resistance-rigidity},
let $E=\{(a,b):1\le a,b\le5,\ a\ne b\}$ and define the
normalized outgoing and incoming incidence isometries by
\[
 U_{\rm out}|a\rangle=\frac12\sum_{b\ne a}|a,b\rangle,
 \qquad
 U_{\rm in}|b\rangle=\frac12\sum_{a\ne b}|a,b\rangle.
\]
With $P_0=J_5/5$, $P_\perp=I_5-P_0$ and the geometric leakage
Laplacian $L=(69/40)P_\perp$, the complete cross-port Gram is
\[
 \boxed{U_{\rm out}^{\dagger}U_{\rm in}
 =\frac{J_5-I_5}{4}
 =P_0-\frac14P_\perp
 =I_5-\frac{50}{69}L.}
\]
It has eigenvalues $1$ on the constant core line and $-1/4$
on the four-dimensional augmentation. Thus the resistance geometry
and the directed incidence geometry have the \emph{same} exact
core-label augmentation projector. This statement does not identify
the reversible heat generator $L$ with the oriented commutator
$[P_{\rm out},P_{\rm in}]$ acting on the different $20$-pair space,
or with a physical irreversible decay generator.
\end{corollary}
\begin{proof}
The two isometries have one nonzero coefficient $1/2$ per
outgoing/incoming incidence. Their cross Gram has diagonal entries
zero and off-diagonal entries $1/4$. Hence it is
$(J_5-I_5)/4=P_0-P_\perp/4=I_5-(5/4)P_\perp$.
The preceding theorem gives $P_\perp=(40/69)L$, proving the
last equality and the stated eigenvalues.
\end{proof}
% END v3.072 aligned five-core resistance/orientation identity.

% BEGIN v3.074 -- source-cyclic oriented five-core Schur return (finite mathematics only).
\begin{theorem}[Least outgoing-source carrier and exact Laplacian Schur self-energy]
\label{p1b:thm:v374-outgoing-source-cyclic-schur}
Use the five-core incidence isometries $U_{\rm out},U_{\rm in}$ of
Corollary~\ref{p1b:cor:v372-resistance-orientation-bridge} and define, on
$\mathbb C^{20}$, the Hermitian oriented comparison operator
\[
 P_{\rm out}=U_{\rm out}U_{\rm out}^{\dagger},\quad
 P_{\rm in}=U_{\rm in}U_{\rm in}^{\dagger},\quad
 H_{\rm ori}=\frac{16i}{\sqrt{15}}[P_{\rm out},P_{\rm in}].
\]
Let $P_0=J_5/5$, $P_\perp=I_5-P_0$ and let $L=(69/40)P_\perp$
be the \emph{independently owned} five-core leakage Gram. Then
\[
 \boxed{H_{\rm ori}^3=H_{\rm ori},\quad
 U_{\rm out}^{\dagger}H_{\rm ori}U_{\rm out}=0,\quad
 U_{\rm out}^{\dagger}H_{\rm ori}^2U_{\rm out}=P_\perp.}
\]
The unique smallest $H_{\rm ori}$-reducing subspace containing the
outgoing source image is
\[
 \mathcal K_{\rm out}=\operatorname{Ran}U_{\rm out}
 \oplus\operatorname{Ran}(H_{\rm ori}U_{\rm out}P_\perp),
 \qquad\boxed{\dim\mathcal K_{\rm out}=5+4=9.}
\]
In this \emph{source-owned} $P\oplus Q$ splitting the diagonal blocks
of $H_{\rm ori}$ vanish, while the off-diagonal block $B=QH_{\rm ori}P$
satisfies $B^{\dagger}B=P_\perp$ in outgoing-source coordinates.
For $z\notin\{0,-1,1\}$ the exact Schur self-energy and complete
outgoing-source resolvent are therefore
\[
 \boxed{\Sigma_{\rm out}(z)=\frac{P_\perp}{z}
       =\frac{40}{69z}L,\qquad
 G_{\rm out}(z)=\frac{P_0}{z}+\frac{z}{z^2-1}P_\perp.}
\]
The source-visible spectral masses are $0,+1,-1$, with matrix
weights $P_0,P_\perp/2,P_\perp/2$, respectively.  The mixed
outgoing--incoming return retains an exactly determined phase:
\[
 \boxed{U_{\rm out}^{\dagger}(z-H_{\rm ori})^{-1}U_{\rm in}
 =\frac{P_0}{z}-\frac{z+i\sqrt{15}}{4(z^2-1)}P_\perp.}
\]
The reverse mixed return is its adjoint for real $z$ off the spectrum;
hence this finite model does not create unequal forward/reverse
\emph{probabilities}, an absorptive width, or a physical hadron pole.
\end{theorem}
\begin{proof}
The cross Gram $U_{\rm out}^{\dagger}U_{\rm in}
=P_0-P_\perp/4$ gives principal-angle cosines $1$ and $1/4$.
On the four augmentation principal-angle planes the commutator
$[P_{\rm out},P_{\rm in}]$ has squared singular value $15/256$;
it vanishes on their orthogonal complement. Thus the indicated scale
makes $H_{\rm ori}^3=H_{\rm ori}$. Direct substitution of the
incidence relations yields the two compressed moments
$U_{\rm out}^{\dagger}HU_{\rm out}=0$ and
$U_{\rm out}^{\dagger}H^2U_{\rm out}=P_\perp$.
Consequently $H U_{\rm out}P_\perp$ has rank four, is orthogonal to
$\operatorname{Ran}U_{\rm out}$, and its image is taken back to
$U_{\rm out}P_\perp$ by $H$. The displayed nine-dimensional sum
is invariant under $H$ and, since $H$ is Hermitian, reducing; no
smaller reducing space can contain the source and its first iterate.
Both diagonal blocks are zero and the off-diagonal return is
$P_\perp$. The Schur formula then gives
$\Sigma=P_\perp/z$ and $G=(zI_5-P_\perp/z)^{-1}$.
Using the spectral projectors $(H^2\pm H)/2$ and $I-H^2$ gives the
three source weights. Their mixed contractions with
$U_{\rm out},U_{\rm in}$ yield the last formula. The equality to
$(40/69z)L$ uses only the earlier leakage Gram, not a newly
selected coupling coefficient.
\end{proof}
% END v3.074 source-cyclic oriented five-core Schur return.

% BEGIN FRAMEWORK v3.069 -- finite cubic five-core operator connection.
\begin{theorem}[First symmetric-cubic five-core augmentation intertwiner]
\label{p1b:thm:v369-five-core-cubic-intertwiner}
Let $V=\operatorname{Im}\mathbb H\cong\mathbb R^3$ carry the real
icosahedral rotation representation of $A_5=2I/\{\pm1\}$, and let
$F_5=\mathbb R^5$ carry its permutation action $P_5(g)$ on the five
conjugate Hurwitz $2T$ cores.  Write
$P_\perp=I_5-J_5/5$, with $J_5$ the all-ones matrix.  For this
\emph{specified} symmetric source representation,
\[
 \operatorname{Sym}^1(V)=\mathbf3,\qquad
 \operatorname{Sym}^2(V)=\mathbf1\oplus\mathbf5,\qquad
 \operatorname{Sym}^3(V)=\mathbf3\oplus\mathbf3'\oplus\mathbf4.
\]
Thus the multiplicity of the five-core augmentation $\mathbf4$ in the
first two symmetric degrees is zero and its multiplicity in the cubic
is precisely one.  Put $D_3(g)=\operatorname{Sym}^3R(g)$ on the
orthonormal ten-dimensional symmetric tensor cube.  Mark a core $a_0$;
let $h_{a_0}$ be the unit $A_4$-invariant harmonic cubic $xyz$ in
orthonormal coordinates adapted to its tetrahedral stabilizer.  Define
\[
 W_{\rm raw}=P_\perp\left({1\over120}\sum_{g\in2I}
 P_5(g)|a_0\rangle\langle h_{a_0}|D_3(g)^\dagger\right)
 P_{\rm harm},\qquad W_3=\sqrt{20}\,W_{\rm raw},
\]
where $P_{\rm harm}$ is the trace-free degree-three projector.  Then
\[
 \boxed{W_3W_3^\dagger=P_\perp,\quad
 W_3^\dagger W_3=P_{\mathbf4}^{(\operatorname{Sym}^3V)},\quad
 P_5(g)W_3=W_3D_3(g).}
\]
The five rows of $W_3$ have Gram $\delta_{ab}-1/5$; all four
nonzero singular values equal one.  Here ``first cubic'' means first
within $\operatorname{Sym}^k(V)$ for the chosen $V$, not first in every
possible Clifford tensor construction and not a cubic physical baryon
interaction theorem.
\end{theorem}
\begin{proof}
The character formula for symmetric powers of the ordinary
icosahedral triplet gives the three displayed decompositions and hence
multiplicity-one of $\mathbf4$.  Group averaging makes $W_{\rm raw}$
equivariant and the harmonic/augmentation projections restrict its
image and source support to the two matching copies of $\mathbf4$.
The marked tetrahedral stabilizer $A_4$ fixes the normalized harmonic
$xyz$, and $\|P_\perp|a_0\rangle\|^2=4/5$; averaging the corresponding
rank-one seed gives $W_{\rm raw}W_{\rm raw}^\dagger=P_\perp/20$
by finite-group orthogonality.  Normalizing yields the first identity;
Schur orthogonality and equality of the two rank-four supports yield
the second.  The stated Gram and covariance follow.  Direct replay of
the original $120$ quaternion coordinates gives Gram and covariance
Frobenius residuals below $2\cdot10^{-15}$.
\end{proof}

\begin{corollary}[Cubic map into the existing five-core operator simplex]
\label{p1b:cor:v369-five-core-operator-cubic}
For the \emph{actual} maps $C_m:\mathcal V_{\rm loc}\to
\mathcal V_{\rm rec}$ of Corollary~\ref{p1:cor:v3026-five-core-hs-gram},
let $C:\mathbb C^5\to\operatorname{HS}(\mathcal V_{\rm loc}^{\mathbb C},
\mathcal V_{\rm rec}^{\mathbb C})$ send $|m\rangle$ to $C_m$.
Their printed Gram gives $C^\dagger C=(69/40)P_\perp$.
Consequently
\[
 \boxed{J_{\rm geom}=\sqrt{40/69}\,CW_3,\qquad
 J_{\rm geom}^\dagger J_{\rm geom}
 =P_{\mathbf4}^{(\operatorname{Sym}^3V)}.}
\]
Thus the cubic augmentation embeds isometrically into an \emph{existing
Hilbert--Schmidt geometric operator space}; this does not assign a
physical axial-current coefficient or a decay amplitude.
\end{corollary}
\begin{proof}
Compute $(40/69)W_3^\dagger C^\dagger CW_3
=W_3^\dagger P_\perp W_3=W_3^\dagger W_3$ and apply the theorem.
\end{proof}
% END FRAMEWORK v3.069 -- finite cubic five-core operator connection.

% BEGIN DOWNSTREAM FLAVOR GEOMETRIC COMPLETION v1.1.
\subsection{Protected dodecahedral flavor algebra and finite charge completions}
\label{p1b:sec:fullcarrier-flavor-completion}

\begin{definition}[Directed-core-pair flavor frame and local matrix lift]
\label{p1b:def:fullcarrier-corepair-flavor-lift}
Take the twenty oriented pairs of distinct conjugate $2T$ cores of
$2I$ and their already constructed dodecahedral coordinate vectors
$u_j\in\mathbb R^3$, $1\le j\le20$, with $|u_j|^2=3/4$.
Let $A$ be their unit-edge dodecahedral adjacency and set
\[
 Q_{ja}=u_j^a/\sqrt5,\qquad P_3=QQ^\dagger,\qquad R_{17}=I_{20}-P_3.
\]
The exact native edge-frame identities are
\[
 \sum_{j=1}^{20}u_ju_j^{\mathsf T}=5I_3,\qquad
 \sum_{k\sim j}u_ku_k^{\mathsf T}={1\over2}I_3+u_ju_j^{\mathsf T}.
\]
For \emph{every} $H\in M_3(\mathbb C)$ define the nearest-neighbor,
site-and-edge, complex-linear operator
\[
 \boxed{\mathscr K(H)_{jk}=2A_{jk}u_j^{\mathsf T}Hu_k
 -2\delta_{jk}u_j^{\mathsf T}Hu_j.}
\]
The definition uses a fixed geometric flavor-frame chart and is not
by itself a physical QCD current or a statement of charge occurrence.
\end{definition}

\begin{theorem}[Protected local flavor matrix calculus]
\label{p1b:thm:fullcarrier-protected-local-flavor}
The local lift is $*$-preserving and satisfies
\[
 \boxed{Q^\dagger Q=I_3,\quad
 \mathscr K(H)Q=QH,\quad
 \mathscr K(H)^\dagger=\mathscr K(H^\dagger),\quad
 \mathscr K(H)=QHQ^\dagger+R_{17}\mathscr K(H)R_{17}.}
\]
Consequently $\operatorname{Ran}Q$ is reducing for every Hermitian
source and carries the fundamental matrix algebra exactly.  On this
sector the three imaginary-skew flavor rotations and five real-symmetric
traceless quadrupoles arise from the \emph{same} one-hop formula.  The
rotational current in the established convention obeys
$\mathscr K(\lambda_a/2)=-C_a$.  No full $20$-state Lie representation
of this specific local lift is asserted.
\end{theorem}
\begin{proof}
The first edge-frame identity gives $Q^\dagger Q=I_3$.
For each $j$ the other edge-frame identity yields
\[
 (\mathscr K(H)Q)_{j\bullet}={2\over\sqrt5}
 u_j^{\mathsf T}H\left({1\over2}I_3+u_ju_j^{\mathsf T}\right)
 -{2\over\sqrt5}(u_j^{\mathsf T}Hu_j)u_j^{\mathsf T}
 ={u_j^{\mathsf T}H\over\sqrt5}=(QH)_{j\bullet}.
\]
Reality of $A,u_j$ gives the adjoint identity.  Apply it also to
$H^\dagger$ to obtain $Q^\dagger\mathscr K(H)=HQ^\dagger$;
therefore both mixed $P_3/R_{17}$ blocks vanish.  For
$(\lambda_a)_{bc}=-i\varepsilon_{abc}$, the edge entries of
$\mathscr K(\lambda_a/2)$ equal $-iA_{jk}(u_j\times u_k)_a$
and its on-site entries vanish, proving the last identification.
\end{proof}

\begin{theorem}[Intrinsic exact zero-defect flavor-three carrier]
\label{p1b:thm:fullcarrier-lie-defect-kernel}
For elementary matrices $E_{ab}$ the local Lie defect
\[
 D_{12}=[\mathscr K(E_{12}),\mathscr K(E_{21})]
 -\mathscr K(E_{11}-E_{22})
\]
has exact rank $17$, squared Hilbert--Schmidt norm $124$, and
\[
 \boxed{\ker D_{12}=\operatorname{Ran}Q.}
\]
For any orthonormal real basis $F_A$ of traceless Hermitian flavor
matrices, set
\[
 D_{AB}=[\mathscr K(F_A),\mathscr K(F_B)]
 -\mathscr K([F_A,F_B]),\qquad
 \mathcal E_{\rm Lie}=\sum_{A<B}D_{AB}^\dagger D_{AB}.
\]
Then $\mathcal E_{\rm Lie}\succeq0$ has exactly the same kernel, and
\[
 \boxed{P_3=\lim_{t\to\infty}\exp(-t\mathcal E_{\rm Lie}).}
\]
This is a geometric algebra-consistency selector, not an assertion
that the physically occurring Hamiltonian minimizes $\mathcal E_{\rm Lie}$.
\end{theorem}
\begin{proof}
The protected lift identity annihilates $\operatorname{Ran}Q$
for every Lie defect.  On the specified $20$ dodecahedral points the
exact $\mathbb Z[\sqrt5]$ computation has
$\operatorname{tr}(D_{12}^\dagger D_{12})=124$.
After scaling $64D_{12}$ to an integral quadratic-number matrix and
using $\sqrt5\mapsto4$ modulo $11$, its leading $17$-minor
(indices $0,\ldots,16$) is $7\pmod{11}$, hence rank at least $17$.
The protected rank-three kernel bounds the rank above by $17$.
A full matrix-unit replay, scaling convention, and minor witness are
included in the named exact finite-geometry certificate.
Each $D_{AB}$ annihilates $Q$; their squared sum is positive and the
single-defect rank witness forces its kernel to be no larger than
$\operatorname{Ran}Q$.  Finite spectral calculus proves the limit.
\end{proof}

\begin{theorem}[Minimum symmetry-compatible full-charge representation]
\label{p1b:thm:fullcarrier-minimal-su3-completion}
The twenty directed pairs carry the exact $A_5$ module
\[
 \mathbb C^{20}\big|_{A_5}\cong
 \mathbf1\oplus\mathbf3\oplus\mathbf3'\oplus2\mathbf4\oplus\mathbf5.
\]
Among finite unitary $SU(3)$ representation spaces containing this
$A_5$ module \emph{and retaining the protected fundamental three as
an invariant $SU(3)$ subrepresentation}, the minimal possible total
dimension is $28$.  An explicit witness is
\[
 \boxed{\mathcal V_{28}=\operatorname{Sym}^1V\oplus
 \operatorname{Sym}^3V\oplus\operatorname{Sym}^4V,
 \quad 3+10+15=28,\quad V=\mathbb C^3.}
\]
Under the real icosahedral subgroup, its harmonic degrees split as
\[
 \bigl(\mathcal H_1\bigr)\oplus
 \bigl(\mathcal H_3\oplus\mathcal H_1\bigr)\oplus
 \bigl(\mathcal H_4\oplus\mathcal H_2\oplus\mathcal H_0\bigr).
\]
The original module is $\mathcal H_0\oplus\mathcal H_1\oplus
\mathcal H_3\oplus\mathcal H_4$; the eight extra \emph{representation
modes} are $\mathcal H_1\oplus\mathcal H_2\cong\mathbf3\oplus\mathbf5$.
The induced $SU(3)$ generators obey the complete charge algebra.
The original harmonic evaluation on twenty vertices is invertible
(in the fixed exact certificate chart, determinant $8\pmod{11}$).
\end{theorem}
\begin{proof}
Decomposition of symmetric tensors into traceless harmonics gives the
stated $A_5$ restrictions and an explicit $A_5$-equivariant
$20$-submodule.  Invertibility of its evaluation and polar
normalization yield the isometry from the original graph carrier.
The finite character census of $SU(3)$ irreducibles through dimension
$28$, intersected with the required $A_5$ character and the invariant
fundamental-three condition, excludes dimensions below $28$; the
same exact character certificate lists all candidates.  On the direct
sum of symmetric tensor powers, $SU(3)$ acts by its canonical tensor
representation, hence its generators satisfy the charge Lie algebra.
\end{proof}

\begin{proposition}[Unchanged one-hop current is impossible in the minimum 28]
\label{p1b:prop:fullcarrier-minimal-compression-no-go}
The above minimal full-charge representation is \emph{not} a
$20$-state exact compression of the unchanged $\mathscr K$.
For each normalized rotation flavor generator $F_a$,
\[
 \boxed{\|\mathscr K(F_a)\|_{\rm HS}^2=5,
 \qquad\|P_{20}\widehat Q_{28}(F_a)P_{20}\|_{\rm HS}^2={45\over2}.}
\]
In fact the forced harmonic-three block contributes $7$ and the
preserved fundamental contributes $1/2$, already exceeding $5$.
The objection applies to the specified minimal-representation class,
not to all finite dilations of $\mathscr K$.
\end{proposition}
\begin{proof}
The one-hop norm is the exact edge sum over the $30$ dodecahedral
edges.  The normalized rotation restricted to $\mathcal H_3$ has
$\tfrac14\sum_{m=-3}^{3}m^2=7$ and the invariant $\mathcal H_1$
fundamental has norm squared $1/2$.  The complete direct-sum
Casimir/harmonic projector calculation gives $45/2$; the two norms
cannot agree under any unitary intertwiner of the required class.
\end{proof}

\begin{theorem}[Finite exact-current dilation without minimality claim]
\label{p1b:thm:fullcarrier-finite-exact-current-dilation}
The original one-hop $\mathscr K$ is the compression of genuine
unitary $U(3)$ generators on a finite $99\,470$-dimensional auxiliary
representation.  More precisely, there are an isometry
$\mathcal J:\mathbb C^{20}\to\mathcal V_{24}$ and a genuine unitary
$U(3)$-representation $\rho$ such that, for every Hermitian $H$,
\[
 \boxed{\mathcal J^\dagger\,d\rho(H)\,\mathcal J=\mathscr K(H).}
\]
The protected flavor-three maps into the invariant fundamental block.
The construction does not claim dimension minimality or a fully
$A_5$-equivariant $\mathcal J$.
\end{theorem}
\begin{proof}
Let $w_s$ be the nine rays consisting of $e_a$ and
$(e_a+e_b)/\sqrt2,(e_a+ie_b)/\sqrt2$ for $a<b$.
Their $P_s=w_sw_s^\dagger$ form a Hermitian-matrix basis.
With Gram-duals $D_s=\sum_t(G^{-1})_{st}P_t$,
$G_{st}=\operatorname{tr}(P_sP_t)$, define
$B_s=R_{17}\mathscr K(D_s)R_{17}$ on its $17$-space,
$B_s^{\pm}=(|B_s|\pm B_s)/2$, and $T=\sum_s|B_s|$.
The exact native data give
\[
 \bigl(\operatorname{tr}B_s^2\bigr)_{s=1}^{9}
 =(57,57,57,42,18,42,18,42,18).
\]
The integer/quadratic-number row-sum certificate establishes
\[
 \left\|\sum_s B_s^2\right\|<64,
 \qquad |B_s|\preceq\tfrac3{16}B_s^2+\tfrac43 I_{17},
 \qquad T\prec24I_{17}.
\]
The middle inequality follows by $(|x|-8/3)^2\ge0$.
Take $n=24$ and the representation
\[
 \mathcal V_n=V\oplus\bigoplus_{s=1}^{9}
 \left[(\operatorname{Sym}^nV\otimes\mathbb C^{17})
 \oplus(\operatorname{Sym}^n\overline V\otimes\mathbb C^{17})\right]
 \oplus\mathbb C^{17}_{\rm triv}.
\]
On the orthogonal $17$-space put the $s$th isometry components
$n^{-1/2}w_s^{\otimes n}\otimes(B_s^+)^{1/2}$,
$n^{-1/2}\overline w_s^{\otimes n}\otimes(B_s^-)^{1/2}$,
and the singlet component $(I_{17}-T/n)^{1/2}$.
Together with the fundamental $Q^\dagger$ component their squared
norms sum to $I_{20}$.  The coherent-state generator expectation is
$\pm n w_s^\dagger H w_s$ on the plus/minus symmetric-tensor
blocks; their compression gives
$\sum_s(w_s^\dagger Hw_s)B_s=R_{17}\mathscr K(H)R_{17}$.
Add the protected $QHQ^\dagger$ block to obtain the assertion.
The dimension is $3+18\cdot17\binom{26}{2}+17=99\,470$.
The exact bound, tomography, and symmetric-tensor matrix certificates
are independently replayable in the geometry network.
\end{proof}

\begin{remark}[Representation completions are not occurring particles]
\label{p1b:rem:fullcarrier-geometry-physical-firewall}
The compact $28$-state construction preserves the original
$A_5$-equivariant representation but changes its off-flavor current;
the larger dilation preserves the entire original current as a
compression but has no established $A_5$-intertwining isometry.
Neither comparison imports a new fundamental axiom, selects the
physical Born/QCD action, adds eight particles or generations, or
produces a hadron pole, retarded decay channel, partial width or
lifetime.  Any use of the auxiliary carriers as physical carriers
requires the separately owned same-action and local-Gauss certificates.
\end{remark}
% END DOWNSTREAM FLAVOR GEOMETRIC COMPLETION v1.1.


\begin{proposition}[Quadrupole--flux tensor contraction]
\label{p1:prop:v371-petersen-quadrupole-flux}
With $Q_v=u_vu_v^{\mathsf T}-I_3/3$ and the oriented edge fluxes
$k_e$ from Theorem~\ref{p1:thm:petersen-signed-and-flux}, define
\[
 \Phi(e=(v,w))=(Q_v+Q_w)\otimes k_e
 \in\operatorname{Sym}^2_0(\mathbb R^3)\otimes\mathbb R^3.
\]
The fixed-coordinate $A_5$-equivariant map $\Phi$ is invertible and
\[
 \Phi^\dagger\Phi={2\over9}\Pi_4
 +{20\over9}\Pi_5+{4\over3}\Pi_3+{32\over9}\Pi_{3'}.
\]
For the natural contraction $c(Q\otimes k)=Qk$,
\[
 \boxed{c\circ\Phi=-{2\over3}K.}
\]
The contraction accesses the ordinary $\mathbf3$ flux sector, not
$\mathbf3'$, and its existence does not insert a physical source.
\end{proposition}
\begin{proof}
The tensor target has dimension fifteen. Equivariance makes the
Gram scalar on the four inequivalent oriented Petersen summands;
its exact finite edge inner products yield the four displayed
strictly positive eigenvalues. Orthogonality
$u_v\cdot k_e=u_w\cdot k_e=0$ gives
$(Q_v+Q_w)k_e=-2k_e/3$, proving the contraction identity directly.
\end{proof}

\begin{proposition}[The marked-core 120-element strata]
\label{p1:prop:petersen-120-strata}
For a selected core $H_m\cong2T$ among the five subgroups of
$I\cong2I$, the union of all five cores has $72$ distinct elements:
its center has two elements, the noncentral portions of the ten
threefold cyclic intersections have $40$ elements, and the fifteen
binary lifts of half-turn axes contribute $30$ order-four elements.
The disjoint marked-core partition is
\[
 I=H_m\;\sqcup\;
 \left(\bigcup_{i=1}^5H_i\setminus H_m\right)
 \;\sqcup\;\left(I\setminus\bigcup_{i=1}^5H_i\right),
 \qquad 120=24+48+48.
\]
The middle $48$ has $24$ order-four and $24$ noncentral threefold
sector elements. The outer $48$ consists of binary fivefold elements;
in the displayed icosian coordinate realization it splits into the
real-part shells
\[
 48=12_{\Re q=\varphi/2}
    +24_{\Re q=\pm1/(2\varphi)}
    +12_{\Re q=-\varphi/2}.
\]
The partition is a partition of group elements, not a sum of
mutually independent subgroup supports or physical mass terms.
\end{proposition}
\begin{proof}
The ten cyclic $C_6$ intersections contribute four noncentral
elements each, meeting only at the common center. The fifteen
half-turn axes in $A_5$ lift to two order-four elements apiece,
giving $2+40+30=72$. Every $2T$ contains the center, sixteen
noncentral threefold lifts and six order-four elements. Removing one
selected $2T$ leaves $24+24=48$ in the union. The remaining $48$
are the fivefold-sector lifts and have the displayed real parts in
the explicit $120$-quaternion list.
\end{proof}

\begin{remark}[Exact extent and physical firewall]
The Petersen $5\to10\to15$ incidence and the icosahedral symmetry
are known in graph theory and icosahedral Jahn--Teller models; they
are not claimed as newly discovered structures. The contribution here
is the explicitly typed realization in the already owned Hopf-axis
geometry and its exact compatibility with the earlier quadrupole
packet. The twelve pentagonal cycles and their reciprocal fluxes
are graph-space data, not hadron masses. The real transverse $A_2$
plane is not the full complex color fundamental. The $15$-edge
module here must not be conflated with the independent level-one
finite carrier $15=1+4+10$ merely because the dimensions agree. Bare Petersen
permutation modules factor through $A_5$; only a separately specified
binary quaternionic coefficient action can see the central $-1$.
The later protected midpoint source acts on a two-dimensional
$u_{ij}^{\perp}$ plane, so its existence cannot be inferred from
these geometric results; any dynamic or mass inference needs a
same-owner source and the actual finite action.
\end{remark}


\subsection{Two-core mirror projector and graded transport algebra}
\label{p1:sec:v327-two-core-mirror-clifford}
For any two distinct marked cores $m,n$ let
$W_{mn}=V_m\oplus V_n$, $D_{mn}=P_{\{m,n\}}DP_{\{m,n\}}$
with $D=\Pi_3-\Pi_{3'}$ as above.  The labels $m,n$ are
\emph{geometric core labels}, not valence-quark or color labels.

\begin{theorem}[Exact two-core mirror-active projection and transport]
\label{p1:thm:v327-pair-mirror-projector}
On the real six-space $W_{mn}$ define
\[
 c=\sqrt{\frac3{20}},\qquad
 Q_{mn}=\frac{20}{3}D_{mn}^{\,2},\qquad
 U_{mn}=\sqrt{\frac{20}{3}}D_{mn},\qquad
 C_{mn}=P_m-P_n.
\]
Then $Q_{mn}$ is an orthogonal projector of rank four, with
$\operatorname{rank}(P_m Q_{mn})=
 \operatorname{rank}(P_n Q_{mn})=2$. More precisely,
\[
\begin{aligned}
 Q_{mn}^2&=Q_{mn}=Q_{mn}^{\mathsf T},&
 U_{mn}^2&=Q_{mn},\\
 U_{mn}Q_{mn}&=Q_{mn}U_{mn}=U_{mn},&
 \{C_{mn},U_{mn}\}&=0,\quad [C_{mn},Q_{mn}]=0.
\end{aligned}
\]
The spectrum of $U_{mn}$ is
$(-1)^{(2)}\oplus0^{(2)}\oplus(+1)^{(2)}$.
On $\operatorname{ran}Q_{mn}\cong\mathbb R^4$,
$U_{mn}^{\,2}=C_{mn}^{\,2}=I$ and
$\{U_{mn},C_{mn}\}=0$, giving two copies of the real
rank-two Clifford involution pair.
Changing the convention $D\mapsto-D$ reverses $U_{mn}$ but preserves
$Q_{mn}$.
\end{theorem}
\begin{proof}
Construct the fifteen oriented Petersen-edge columns $K$ from
the ten exact axis representatives displayed in
Theorem~\ref{p1:thm:v285-ten-axis-quadrupole-exact-sequence},
as in Theorem~\ref{p1:thm:petersen-signed-and-flux}, and form
$D=(K^{\mathsf T}K-(K^\sigma)^{\mathsf T}K^\sigma)/5$.
For one selected pair of omitted-core triples, with harmless
edge sign choices, its off-diagonal block $X=P_mDP_n$ obeys
\[
 XX^{\mathsf T}=\frac1{20}
 \begin{pmatrix}2&-1&-1\\-1&2&-1\\-1&-1&2\end{pmatrix},
 \qquad \operatorname{Spec}(XX^{\mathsf T})
 =\{(3/20)^{(2)},0\}.
\]
The other nine pairs have the same singular values by the native
signed $A_5$ core action; alternatively all ten are checked by
exact substitution in the attached axis-coordinate certificate.
The two-core compression has block form
$D_{mn}=\left(\begin{smallmatrix}0&X\\X^{\mathsf T}&0\end{smallmatrix}\right)$,
so its eigenvalues are $\pm c$ twice each and zero twice.
Its square divided by $c^2$ is the orthogonal projector onto the
four-dimensional nonzero singular space, rank two in each core.
$U_{mn}$ is its normalized off-diagonal partial involution.
Conjugating this off-diagonal block by
$C_{mn}=\operatorname{diag}(I_3,-I_3)$ reverses its sign, while
$D_{mn}^2$ is block diagonal. These identities prove the assertions.
\end{proof}

\begin{corollary}[Eligibility is not coherence]
\label{p1:cor:v327-pair-eligibility-not-coherence}
$Q_{mn}$ cannot witness inter-core phase coherence: it is block
 diagonal in $V_m\oplus V_n$ and a block-incoherent state may have
$\operatorname{Tr}(\rho Q_{mn})>0$.
In contrast $U_{mn}$ is purely off diagonal, so
$\operatorname{Tr}(\rho U_{mn})=0$ for every state $\rho$ block
 diagonal in the pair. This statement is about mirror transport on the
finite geometric edge module, not an observed meson state or energy.
\end{corollary}
\begin{proof}
The block forms in the theorem prove both trace statements. A pure
state in $\operatorname{ran}(P_mQ_{mn})$ is block incoherent yet has
unit $Q_{mn}$ expectation.
\end{proof}


\subsection{Finite Kneser--Hodge completion of the five-core packet}
\label{p1:sec:kneser-hodge-closure}

Continue with the occurring, marked binary-icosahedral group of
Corollary~\ref{p1:cor:v237-five-core-ten-axis-packet} and its five conjugate
binary-tetrahedral cores. All finite incidence statements below refer to that
same packet. They do not assert a new universally occurring constituent or an
additional factor of the Born space. The ordinary Kneser graphs $KG(n,k)$
should not be confused with the distinct generalized Petersen graph family
$GP(n,k)$.

\begin{theorem}[Dual six-axis triangle systems and the second Petersen graph]
\label{p1:thm:kneser-six-axis-triangle-systems}
\leavevmode\par
Let $G=2I/\{\pm1\}\cong A_5$ act on its six Sylow-$5$ subgroups,
equivalently the six projective fivefold axes. On the twenty unordered
three-element subsets of these axes there are exactly two $G$-orbits
$\mathcal F_+,\mathcal F_-$, each of size ten. Complementation exchanges
the two orbits. In either orbit every unordered pair of the six axes lies
in exactly two triples and the vertex link of each axis is a five-cycle.
Thus both face systems are triangulations of $\mathbb{RP}^2$ with
\[
 (f_0,f_1,f_2)=(6,15,10),\qquad \chi=6-15+10=1.
\]
The dual graph of either triangulation is the Petersen graph.
The six-axis pair graph defined by disjointness is the ordinary Kneser
graph $KG(6,2)$; its fifteen vertices identify with the fifteen edges of
each projective-plane triangulation.
\end{theorem}
\begin{proof}
Conjugation on the six Sylow-$5$ subgroups gives the faithful degree-six
action of $A_5$. An exact orbit computation in this finite action gives
subset-orbit sizes $1,6,15,(10,10),15,6,1$ for subset cardinalities
$0,1,2,3,4,5,6$ respectively. One explicit representative ten-face
system, in a consistent labeling of the six axes, is
\[
\mathcal F_+=\{123,124,135,146,156,236,245,256,345,346\},
\qquad \mathcal F_-=\{\{1,\dots,6\}\setminus t:t\in\mathcal F_+\}.
\]
Direct inspection shows that each pair occurs in exactly two faces.
By transitivity, every vertex has five incident triangles; its link
is a simple two-regular graph on five vertices, hence a five-cycle.
The resulting connected, closed simplicial surface has Euler
characteristic one, so it is $\mathbb{RP}^2$. Its ten-face dual graph
is cubic; direct common-neighbor enumeration gives
$A^2=2I-A+J$, the Petersen relation. The fifteen edges are all
pairs of the six labels. Disjointness among those pairs defines
$KG(6,2)$. All stated finite checks are independently reproducible
from the six Sylow-$5$ subgroups, without assuming the diagram.
\end{proof}

\begin{theorem}[Equivariant Petersen duality and signed-Hodge orientation]
\label{p1:thm:kneser-five-six-hodge-duality}
For each sign, there is an $A_5$-equivariant bijection
\[
 f_\pm:\binom{\{1,\dots,5\}}2\longrightarrow\mathcal F_\pm,
\]
unique once the core and six-axis labelings and the sign orbit are
fixed, such that
\[
 \{i,j\}\cap\{k,l\}=\varnothing
 \quad\Longleftrightarrow\quad
 |f_\pm(ij)\cap f_\pm(kl)|=2.
\]
Hence the fifteen Petersen edges correspond bijectively to the
fifteen unordered pairs of fivefold axes. The two maps are related
by triple complementation.

Set $V=\mathbb R^5$, $c=e_1+\cdots+e_5$ and
\[
 \mathsf S_H=*_5(c\wedge\cdot):\Lambda^2V\to\Lambda^2V,
 \qquad B=\iota_c:\Lambda^2V\to V.
\]
The matrix of $\mathsf S_H$ in the oriented pair basis is a signed
Petersen adjacency. It satisfies
\[
 B\mathsf S_H=0,\qquad B^{\mathsf T}B+\mathsf S_H^2=5I_{10},
 \qquad\mathsf S_H^3=5\mathsf S_H.
\]
Orienting the triangular faces coherently in local charts of each
projective plane induces signed dual-graph matrices $\mathsf S_\pm$.
For diagonal matrices $G_\pm$ with entries in $\{\pm1\}$ and a
suitable assignment of orbit names,
\[
 \boxed{\mathsf S_+=G_+\mathsf S_HG_+,\qquad
 \mathsf S_-=-G_-\mathsf S_HG_-.}
\]
Thus the complementary projective-plane systems exchange the
rank-three $\pm\sqrt5$ Hodge eigenspaces. In this marked packet these
are precisely the two conjugate icosahedral triplets $\mathbf3$ and
$\mathbf3'$ of Theorem~\ref{p1:thm:petersen-signed-and-flux}.
\end{theorem}
\begin{proof}
A pair of five core labels has stabilizer of order six in $A_5$.
In each ten-triangle orbit it fixes exactly one triangle. Transporting
that triangle equivariantly defines $f_\pm$; direct enumeration of the
ten images verifies the incidence equivalence and that the maps differ
by complement. The corresponding fifteen shared pairs exhaust all
fifteen edges of $K_6$. For the exterior-algebra identities, put
$W=c^\perp$. The orthogonal splitting
$\Lambda^2V=(c\wedge W)\oplus\Lambda^2W$ has dimensions $4+6$.
The contraction $B$ has $B^{\mathsf T}B=5I$ on $c\wedge W$ and
zero on $\Lambda^2W$, while $\mathsf S_H$ vanishes on $c\wedge W$
and is $\sqrt5$ times the four-dimensional Hodge star on
$\Lambda^2W$. These facts give all three identities and the
$4+3+3$ spectrum. The local triangle boundary signs yield the signed
dual-graph matrices. An explicit ten-by-ten sign-gauge check proves
the two displayed conjugacies; the same check is encoded by the
integer certificate accompanying this theorem.
\end{proof}

\begin{corollary}[Twisted-cycle roots and the complete quintet bridge]
\label{p1:cor:kneser-twisted-cycle-root-bridge}
For either projective-plane triangulation, orient its fifteen
$K_6$ edges and let $F\in\mathbb R^{15\times6}$ be their signed
vertex-edge incidence matrix. Let $\partial_s\in\mathbb R^{10\times15}$
be the twisted boundary operator on the dual Petersen graph induced by
local triangle orientations. Then
\[
 \partial_sF=0,\qquad F^{\mathsf T}F=6I_6-J_6,\qquad
 \partial_s\partial_s^{\mathsf T}=3I_{10}-\mathsf S_\pm,
\]
\[
 \dim\ker\partial_s=5,\qquad
 \boxed{P_{\rm tw}=\frac16FF^{\mathsf T}}.
\]
The normalized fifteen root lines are
$r_{ab}=(e_a-e_b)/\sqrt2$ in the sum-zero five-space of
$\mathbb R^6$. Their orthogonality graph is $KG(6,2)$.
This five-dimensional twisted cycle module is isometric, as an
$A_5$ module, to the trace-free quadrupole module of
Theorem~\ref{p1:thm:v285-ten-axis-quadrupole-exact-sequence}.
\end{corollary}
\begin{proof}
At a dual vertex (a triangle $abc$), the three edges
$ab,ac,bc$ satisfy the oriented identity
$(e_b-e_a)-(e_c-e_a)+(e_c-e_b)=0$.
This is exactly $\partial_sF=0$ with the chosen local signs.
The incidence matrix of $K_6$ has
$F^{\mathsf T}F=6I_6-J_6$, so $F$ has rank five.
The signed Petersen vertex Laplacian $3I-\mathsf S_\pm$ is positive
definite since its eigenvalues are $3,3\pm\sqrt5$;
thus $\partial_s$ has rank ten and its kernel is precisely
$\operatorname{im}F$. The displayed normalized projector follows.
Distinct root vectors are orthogonal exactly when their index pairs
are disjoint. The induced $A_5$ quintet agrees with the unique
five-dimensional real irreducible quadrupole representation, and the
isometry is fixed by the normalized tight frames.
\end{proof}

\begin{theorem}[Universal signed odd-Kneser Hodge operator]
\label{p1:thm:kneser-universal-odd-hodge}
For any integer $k\ge1$ set $n=2k+1$, $V=\mathbb R^n$,
$c=\sum_{i=1}^ne_i$, and
\[
 B_k=\iota_c:\Lambda^k V\to\Lambda^{k-1}V,
 \quad \mathsf S_k=*_V(c\wedge\cdot):\Lambda^kV\to\Lambda^kV.
\]
Its nonzero matrix entries, all $\pm1$, join precisely disjoint
$k$-subsets, so its unsigned support is the odd graph $KG(2k+1,k)$.
It satisfies the exact identities
\[
 \mathsf S_k^{\mathsf T}=(-1)^k\mathsf S_k,\quad
 B_k\mathsf S_k=0,\quad
 \boxed{B_k^{\mathsf T}B_k+(-1)^k\mathsf S_k^2=nI},
 \quad \mathsf S_k^3=(-1)^k n\mathsf S_k.
\]
The zero eigenvalue has multiplicity
$\binom{2k}{k-1}$. For even $k$, the other eigenvalues are
$\pm\sqrt n$, each of multiplicity $\tfrac12\binom{2k}{k}$;
for odd $k$ they are $\pm i\sqrt n$ with the same multiplicities.
The operator transforms by the orientation sign under permutations
of the $n$ labels, and is consequently equivariant under $A_n$.
\end{theorem}
\begin{proof}
Write $c=\sqrt n\,e_0$ and $W=e_0^\perp$, so $\dim W=2k$.
On $\Lambda^kV=(e_0\wedge\Lambda^{k-1}W)\oplus\Lambda^kW$,
$B_k^{\mathsf T}B_k$ is $nI$ on the first summand and zero on
the second. Conversely, $\mathsf S_k$ is zero on the first summand
and restricts, up to the orientation convention, to
$\sqrt n\,*_W$ on the second. As $*_W^2=(-1)^kI$ on
middle-degree $k$-forms of the $2k$-space $W$, all stated operator
identities follow. The nonzero block has dimension
$\binom{2k}{k}$ and vanishing trace, giving equal multiplicities.
A basis coefficient of $\mathsf S_k$ is nonzero exactly when its two
$k$-subsets have disjoint labels, leaving one unused label.
\end{proof}

\begin{corollary}[Universal incidence--Hodge Dirac square]
\label{p1:cor:kneser-universal-dirac}
Put $H_k=\mathsf S_k$ for even $k$ and $H_k=i\mathsf S_k$ for odd
$k$. The Hermitian operator
\[
 \mathscr D_k=\begin{pmatrix}0&B_k\\B_k^{\mathsf T}&H_k\end{pmatrix}
 \quad\text{on }\Lambda^{k-1}V\oplus\Lambda^kV
\]
obeys
\[
 \boxed{\mathscr D_k^2=\begin{pmatrix}
 B_kB_k^{\mathsf T}&0\\0&nI_{\binom nk}
 \end{pmatrix}.}
\]
Its spectrum consists of $+\sqrt n$ and $-\sqrt n$, each of
multiplicity $\binom{2k}{k-1}+\tfrac12\binom{2k}{k}$, and a zero
of multiplicity $\binom{2k}{k-2}$ (zero by convention when $k=1$).
For $k=2$ this gives the $(7,1,7)$ multiplicity pattern in the
fifteen-dimensional five-core incidence--Hodge operator.
\end{corollary}
\begin{proof}
The previous theorem gives $B_kH_k=0$ and
$B_k^{\mathsf T}B_k+H_k^2=nI$. On the lower-degree space,
$B_kB_k^{\mathsf T}$ acts as $nI$ on
$\Lambda^{k-1}W$ and as zero on
$e_0\wedge\Lambda^{k-2}W$. The positive and negative
multiplicities are equal by the displayed block structure and the
zero trace of $H_k$. Counting dimensions proves the claim.
\end{proof}

\begin{proposition}[No forced higher odd-graph tower from one $2I$]
\label{p1:prop:kneser-five-six-native-limit}
Within one marked $2I$, the homogeneous five-core and six-fivefold-axis
sets supply the two finite subset constructions above. The first
higher odd graph beyond Petersen, $KG(7,3)$, is a valid abstract
Hodge construction but is not forced by a new homogeneous seven-label
icosahedral orbit: a transitive $A_5$ action on seven objects is
impossible because $7\nmid60$. The same obstruction applies to
nine labels. Subset sets assembled from composite or differently
marked orbits do not discharge an occurring-source requirement.
In particular neither the universal odd-graph series nor the finite
$KG(6,2)$ root system supplies a new color operator, hadron mass pole,
or coefficient in the physical action.
\end{proposition}
\begin{proof}
An orbit of a finite group has size equal to the group order divided
by a stabilizer order. Neither seven nor nine divides $|A_5|=60$.
The generalized exterior-algebra theorem requires only abstract
labels and an orientation; these are not a same-owner physical
constitution proof. The identified geometric modules have no
established isometric reducing embedding into the finite color-singlet
Hamiltonian of Paper~VII-A. This is exactly the typed action firewall
already imposed by the Petersen source analysis.
\end{proof}



\begin{theorem}[Exceptional sibling generation is dense in \(SU(2)\)]
\label{p1:thm:2O-2I-generate-SU2}
For the representatives containing the common Hurwitz core \(2T\),
\begin{equation}
\label{p1:eq:2O-2I-generate-SU2}
 \boxed{\overline{\langle2O,2I\rangle}=SU(2).}
\end{equation}
Consequently, if one finite-dimensional continuous \(SU(2)\)-carrier carries
one and the same typed operator \(X\) satisfying
\[
 \rho(g)X\rho(g)^{-1}=X\quad(g\in2O),
 \qquad
 \rho(h)X\rho(h)^{-1}=X\quad(h\in2I),
\]
then \(X\) is \(SU(2)\)-invariant.  On an irreducible complex
\(SU(2)\)-carrier, \(X\) is therefore scalar.
\end{theorem}

\begin{proof}
Let \(K=\overline{\langle2O,2I\rangle}\).  The classification of finite
subgroups of \(SU(2)\), together with \(2O\not\subset2I\), shows that \(K\)
is not finite.  If a proper closed subgroup of \(SU(2)\) has positive
dimension, its identity component is a maximal torus and the subgroup lies
in that torus's normalizer.  Such a normalizer is solvable and cannot contain
the perfect nonabelian group \(2I\).  Hence \(K=SU(2)\).  Continuity extends
invariance from the generated dense subgroup to its closure, and Schur's
lemma gives the irreducible-carrier statement.
\end{proof}

\begin{remark}[Same carrier and same operator are essential]
The theorem does not combine quantities merely because one has an
octahedral label and another an icosahedral label.  Both finite actions must
first be proved to occur on the same typed carrier and on the same operator
or tensor.  In particular, this result neither creates an owner map nor
equates amplitudes living on different response blocks.
\end{remark}



\begin{theorem}[Shared-axis octahedral--icosahedral completion is dense in \(\operatorname{Pin}(2)\)]
\label{p1:thm:v285-shared-axis-pin2}
Let
\[
 h=\frac{1+\mathbf i+\mathbf j+\mathbf k}{2}
 =\cos\frac{\pi}{3}+n\sin\frac{\pi}{3},
 \qquad
 n=\frac{\mathbf i+\mathbf j+\mathbf k}{\sqrt3},
\]
so \(C_6=\langle h\rangle\) is a common generated axis subgroup of the
standard aligned \(2O\) and \(2I\) representatives.  Then
\[
 \boxed{
 \overline{\left\langle
 N_{2O}(C_6),N_{2I}(C_6)
 \right\rangle}
 =N_{SU(2)}(T_n)\simeq\operatorname{Pin}(2),}
\]
where \(T_n=\{\cos t+n\sin t\}\) is the axial maximal circle.
\end{theorem}

\begin{proof}
Take the transverse pure unit quaternions
\[
 q_O=\frac{\mathbf i-\mathbf j}{\sqrt2}\in2O,
 \qquad
 q_I=\frac{\mathbf i-\varphi\mathbf j+\varphi^{-1}\mathbf k}{2}\in2I.
\]
Both imaginary axes are perpendicular to \(n\), so conjugation by either
sends \(n\mapsto-n\) and therefore \(h\mapsto h^{-1}\).  Thus
\(q_O\in N_{2O}(C_6)\) and \(q_I\in N_{2I}(C_6)\), and both lie in the
non-toral component of \(N_{SU(2)}(T_n)\).

The product \(r=q_Oq_I\) lies in \(T_n\).  If
\(r=\cos\alpha+n\sin\alpha\), direct multiplication gives
\[
 2\cos\alpha=-\frac{\varphi^2}{\sqrt2}
 =-\frac{3+\sqrt5}{2\sqrt2}.
\]
This algebraic number has irreducible polynomial
\[
 4x^4-14x^2+1,
\]
so it is not an algebraic integer.  If \(\alpha/\pi\) were rational, then
\(2\cos\alpha=\zeta+\zeta^{-1}\) for a root of unity \(\zeta\), hence would
be an algebraic integer, a contradiction.  Therefore \(r\) has infinite
order, and its powers are dense in the circle \(T_n\).  Since the generated
closure also contains the transverse element \(q_O\), it contains both
components of the normalizer.  Conversely both finite normalizers already lie
in \(N_{SU(2)}(T_n)\), proving equality.
\end{proof}

\begin{remark}[Local/global continuous completions]
Theorem~\ref{p1:thm:v285-shared-axis-pin2} is the axis-local companion of
Theorem~\ref{p1:thm:2O-2I-generate-SU2}.  The former completes one shared
exceptional axis to its full axial normalizer; the latter completes the two
exceptional siblings globally to \(SU(2)\).  Both statements concern the
topological closures of the displayed generated subgroups.
\end{remark}

\begin{definition}[Nearest-shell walk]
\label{p1:def:nearest-shell-walk}
Let \(\Gamma\subset SU(2)\simeq S^3\) be a finite subgroup written in quaternion
coordinates.  Its nearest nontrivial real-height shell is
\[
    S_\Gamma:=\{g\in \Gamma\setminus\{1\}:\operatorname{Re}(g)
    \text{ is maximal}\}.
\]
The normalized nearest-shell walk is the convolution operator
\[
    T_\Gamma f(x):=\frac{1}{|S_\Gamma|}
    \sum_{s\in S_\Gamma} f(s^{-1}x),
    \qquad f\in \ell^2(\Gamma).
\]
This is a finite group-algebra datum only.  No physical interpretation is imposed on \(T_\Gamma\).
\end{definition}

\begin{theorem}[Nearest-shell spectra of the exceptional binary groups]
\label{p1:thm:exceptional-nearest-shell-spectra}
For the explicit quaternionic realizations of \(2T,2O,2I\) above, the normalized
nearest-shell walks have the following spectra on \(\ell^2(\Gamma)\).

For the binary tetrahedral group \(2T\),
\[
    |S_{2T}|=8,\qquad \operatorname{Re}(S_{2T})=\frac12,
\]
and
\[
\operatorname{Spec}(T_{2T})
=
1^{(1)}
\oplus
\left(\frac12\right)^{(4)}
\oplus
0^{(9)}
\oplus
\left(-\frac14\right)^{(8)}
\oplus
\left(-\frac12\right)^{(2)} .
\]

For the binary octahedral group \(2O\),
\[
    |S_{2O}|=6,\qquad \operatorname{Re}(S_{2O})=\frac{1}{\sqrt2},
\]
and
\[
\operatorname{Spec}(T_{2O})
=
1^{(1)}
\oplus
\left(\frac1{\sqrt2}\right)^{(4)}
\oplus
\left(\frac13\right)^{(9)}
\oplus
0^{(20)}
\oplus
\left(-\frac13\right)^{(9)}
\oplus
\left(-\frac1{\sqrt2}\right)^{(4)}
\oplus
(-1)^{(1)} .
\]

For the binary icosahedral group \(2I\),
\[
    |S_{2I}|=12,\qquad \operatorname{Re}(S_{2I})=\frac{\phi}{2},
\]
and
\[
\begin{aligned}
\operatorname{Spec}(T_{2I})
={}&
1^{(1)}
\oplus
\left(\frac{\phi}{2}\right)^{(4)}
\oplus
\left(\frac{\phi}{3}\right)^{(9)}
\oplus
\left(\frac14\right)^{(16)}
\oplus
0^{(25)}
\\
&\oplus
\left(-\frac16\right)^{(36)}
\oplus
\left(-\frac{1}{3\phi}\right)^{(9)}
\oplus
\left(-\frac14\right)^{(16)}
\oplus
\left(-\frac{1}{2\phi}\right)^{(4)} .
\end{aligned}
\]
\end{theorem}

\begin{proof}
The shells \(S_\Gamma\) are inverse-closed real-height shells in the displayed
quaternionic groups, hence \(T_\Gamma\) is a self-adjoint convolution operator.
Since each \(S_\Gamma\) is a union of conjugacy classes in \(\Gamma\), the
corresponding group-algebra element is central.  By the regular representation
decomposition, a central convolution element acts on the \(\rho\)-isotypic block by
the scalar
\[
    \lambda_\rho
    =
    \frac{1}{|S_\Gamma|\dim\rho}
    \sum_{s\in S_\Gamma}\chi_\rho(s),
\]
with multiplicity \((\dim\rho)^2\).  For completeness, direct evaluation on
the displayed quaternion shells gives the following full character-sum
certificate.  In the last column
\(\Sigma_\rho:=\sum_{s\in S_\Gamma}\chi_\rho(s)
=|S_\Gamma|(\dim\rho)\lambda_\rho\); repeated dimensions denote distinct
irreducibles.
\[
\begin{array}{c|c|c|c}
\Gamma&\lambda_\rho&\dim\rho&\Sigma_\rho\\
\hline
2T&1&1&8\\
&\frac12&2&8\\
&0&3&0\\
&-\frac14&2,2&-4,-4\\
&-\frac12&1,1&-4,-4\\
\hline
2O&1&1&6\\
&\frac1{\sqrt2}&2&6\sqrt2\\
&\frac13&3&6\\
&0&2,4&0,0\\
&-\frac13&3&-6\\
&-\frac1{\sqrt2}&2&-6\sqrt2\\
&-1&1&-6\\
\hline
2I&1&1&12\\
&\frac\phi2&2&12\phi\\
&\frac\phi3&3&12\phi\\
&\frac14&4&12\\
&0&5&0\\
&-\frac16&6&-12\\
&-\frac1{3\phi}&3&-12/\phi\\
&-\frac14&4&-12\\
&-\frac1{2\phi}&2&-12/\phi
\end{array}
\]
The squared dimensions in the three blocks sum to \(24,48,120\),
respectively.  Multiplying each row by \((\dim\rho)^2\) therefore gives
exactly the displayed regular-representation multiplicities and accounts for
every dimension.  The same table is obtained from the standard character
tables of the binary polyhedral groups
\cite{conwaysmith2003quaternions,wolf2011spaces}.
\end{proof}

\begin{definition}[Uniform exceptional nearest-shell scalar atom]
\label{p1:def:exceptional-shell-scalar-atom}
For \(\Gamma\in\{2T,2O,2I\}\), write
\[
 n_\Gamma:=|S_\Gamma|,
 \qquad
 r_\Gamma:=\operatorname{Re}(s)\quad(s\in S_\Gamma).
\]
The uniform nearest-shell scalar atom is
\begin{equation}
\label{p1:eq:exceptional-shell-scalar-atom}
 \boxed{\omega_\Gamma:=\frac{r_\Gamma}{n_\Gamma}.}
\end{equation}
It is the scalar part of one uniformly weighted nearest-shell summand.
Because the real height is constant on \(S_\Gamma\), no shell element is
marked.  The symbol \(\omega_\Gamma\) is a finite group-algebra scalar; it is
not, by this definition, a smooth determinant correction, coupling,
transition probability, or thermal weight.
\end{definition}

\begin{theorem}[Exceptional \(3,4,5\) shell-angle, atom, and chord table]
\label{p1:thm:exceptional-345-shell-table}
Put \(m_T=3\), \(m_O=4\), and \(m_I=5\).  Then
\begin{equation}
\label{p1:eq:exceptional-345-real-height}
 r_{2T}=\cos\frac{\pi}{3}=\frac12,\qquad
 r_{2O}=\cos\frac{\pi}{4}=\frac1{\sqrt2},\qquad
 r_{2I}=\cos\frac{\pi}{5}=\frac{\phi}{2},
\end{equation}
so the defining lifted traces are
\begin{equation}
\label{p1:eq:exceptional-345-lifted-trace}
 \boxed{2r_{2T}=1,\qquad2r_{2O}=\sqrt2,\qquad2r_{2I}=\phi.}
\end{equation}
The corresponding uniform atoms and identity-based squared chords are
\begin{align}
\label{p1:eq:exceptional-shell-atom-table}
 \boxed{\omega_{2T}=\frac1{16},\qquad
 \omega_{2O}=\frac1{6\sqrt2},\qquad
 \omega_{2I}=\frac{\phi}{24},}\\
\label{p1:eq:exceptional-nearest-chords}
 \boxed{\ell_{2T}=1,\qquad
 \ell_{2O}=2-\sqrt2,\qquad
 \ell_{2I}=2-\phi=\phi^{-2},}
\end{align}
where \(\ell_\Gamma:=\|1-s\|^2\) for \(s\in S_\Gamma\).  Finally,
\begin{equation}
\label{p1:eq:exceptional-shell-atom-order}
 \boxed{\omega_{2T}<\omega_{2I}<\omega_{2O}.}
\end{equation}
\end{theorem}

\begin{proof}
A rotation of angle \(2\pi/m\) in \(SO(3)\) lifts to a unit quaternion of
real part \(\cos(\pi/m)\).  The shell cardinalities and real heights are those
of Theorem~\ref{p1:thm:exceptional-nearest-shell-spectra}; division gives
\eqref{p1:eq:exceptional-shell-atom-table}.  For a unit quaternion,
\(\|1-s\|^2=2-2\operatorname{Re}(s)\), and
\(2-\phi=\phi^{-2}\).  The strict ordering follows from
\(2\phi>3\) and \(2\sqrt2>\phi\).
\end{proof}

\begin{corollary}[Defining trace-field fork]
\label{p1:cor:exceptional-defining-trace-fields}
For the defining two-dimensional representations,
\begin{equation}
\label{p1:eq:exceptional-defining-trace-fields}
 \boxed{K_{2T}=\mathbb Q,\qquad
 K_{2O}=\mathbb Q(\sqrt2),\qquad
 K_{2I}=\mathbb Q(\sqrt5)=\mathbb Q(\phi),}
\end{equation}
where \(K_\Gamma:=\mathbb Q(\operatorname{tr}g:g\in\Gamma)\).  Hence
\[
 K_{2O}\cap K_{2I}=\mathbb Q=K_{2T}.
\]
This arithmetic intersection is compatible with the geometric fork
\(2O\cap2I=2T\).  Field membership alone is not a reverse support-owner
selector.
\end{corollary}

\begin{proof}
In the defining \(SU(2)\) representation,
\(\operatorname{tr}g=2\operatorname{Re}(g)\).  The explicit quaternion
lists show that the three trace sets lie respectively in the displayed
fields, while the nearest shells contain the primitive values
\(1,\sqrt2,\phi\).  Distinct quadratic extensions
\(\mathbb Q(\sqrt2)\) and \(\mathbb Q(\sqrt5)\) intersect in \(\mathbb Q\).
\end{proof}

\begin{proposition}[Nearest-shell traces as McKay adjacency eigenvalues]
\label{p1:prop:exceptional-shell-mckay-eigenvalues}
Let \(A_\Gamma\) be the McKay adjacency matrix for tensoring irreducible
\(\Gamma\)-representations by the defining representation \(V\).  For each
conjugacy class \(C\subset\Gamma\), the character column
\((\chi_\rho(C))_\rho\) is an eigenvector with eigenvalue \(\chi_V(C)\).
Every class contained in the nearest shell therefore has eigenvalue
\begin{equation}
\label{p1:eq:exceptional-shell-mckay-eigenvalues}
 \boxed{2T\leftrightarrow\widetilde E_6:\ 1,\qquad
 2O\leftrightarrow\widetilde E_7:\ \sqrt2,\qquad
 2I\leftrightarrow\widetilde E_8:\ \phi.}
\end{equation}
\end{proposition}

\begin{proof}
Evaluating
\(V\otimes\rho_i=\bigoplus_j(A_\Gamma)_{ij}\rho_j\) on \(C\) gives
\[
 \chi_V(C)\chi_{\rho_i}(C)
 =\sum_j(A_\Gamma)_{ij}\chi_{\rho_j}(C).
\]
Thus the character column is an eigenvector.  On the nearest shell,
\(\chi_V(C)=2\operatorname{Re}(C)=2r_\Gamma\), and
Theorem~\ref{p1:thm:exceptional-345-shell-table} gives the values.
\end{proof}

\begin{theorem}[Octahedral nearest-shell tangent frame]
\label{p1:thm:2O-nearest-shell-tangent-frame}
For \(s\in S_{2O}\), define
\[
 a_s:=\frac{\operatorname{Im}(s)}{\|\operatorname{Im}(s)\|}
 \in\mathfrak{su}(2)\cong\mathbb R^3.
\]
In the round orthonormal basis \((e_1,e_2,e_3)\),
\[
 S_{2O}=\left\{
 \frac{1\pm e_1}{\sqrt2},
 \frac{1\pm e_2}{\sqrt2},
 \frac{1\pm e_3}{\sqrt2}\right\},
\]
and hence
\begin{equation}
\label{p1:eq:2O-shell-tight-frame}
 \boxed{\sum_{s\in S_{2O}}a_s\otimes a_s=2I_3,\qquad
 \frac1{6}\sum_{s\in S_{2O}}a_s\otimes a_s=\frac13I_3.}
\end{equation}
The shell-frame form
\[
 F_O(X,Y):=\sum_{s\in S_{2O}}
 g_E(a_s,X)g_E(a_s,Y)
\]
satisfies
\begin{equation}
\label{p1:eq:2O-shell-frame-metric}
 \boxed{F_O=2g_E,\qquad F_O|_H=\gamma_0.}
\end{equation}
Moreover the defining and adjoint nearest-shell eigenvalues are
\begin{equation}
\label{p1:eq:2O-def-ad-eigenvalues}
 \boxed{\lambda_{\rm def}(T_{2O})=\frac1{\sqrt2},\qquad
 \lambda_{\operatorname{Ad}}(T_{2O})=\frac13.}
\end{equation}
\end{theorem}

\begin{proof}
The six unit directions are \(\pm e_1,\pm e_2,\pm e_3\), so every coordinate
projector occurs twice.  This proves the tight-frame and metric identities.
Restriction to \(H\), together with the existing identity
\(\gamma_0=2g_E|_H\), gives \(F_O|_H=\gamma_0\).  The defining eigenvalue is
the positive two-dimensional entry in
Theorem~\ref{p1:thm:exceptional-nearest-shell-spectra}.  The positive
three-dimensional entry is adjoint: a shell element projects to a
\(\pi/2\) rotation in \(SO(3)\), whose adjoint character is
\(1+2\cos(\pi/2)=1\), and normalization by dimension three gives \(1/3\).
\end{proof}

\begin{corollary}[Exact \(E_6\) rank--Coxeter numerical compatibility]
\label{p1:cor:2O-shell-E6-reconstruction}
Under the standard McKay identification
\(2T\leftrightarrow\widetilde E_6\) and the proved normalizer identity
\(N_{SU(2)}(2T)=2O\),
\begin{align}
\label{p1:eq:2O-shell-E6-rank-coxeter}
 |S_{2O}|&=6=\operatorname{rank}(E_6),&
 \frac{|S_{2O}|}{r_{2O}^{\,2}}&=12=h(E_6),\\
\label{p1:eq:2O-shell-E6-root-count}
 \boxed{\omega_{2O}^{-2}
 =\left(\frac{|S_{2O}|}{r_{2O}}\right)^2
 =\operatorname{rank}(E_6)h(E_6)
 =|\Phi(E_6)|=72.}
\end{align}
This is an exact compatibility of independently defined finite invariants,
not a reverse derivation of the \(E_6\) owner and not an identification with
a smooth determinant germ.
\end{corollary}

\begin{corollary}[Pure inclusion consequences of the nearest-shell walks]
\label{p1:cor:nearest-shell-inclusion-consequences}
The explicit shells and spectra in
Theorem~\ref{p1:thm:exceptional-nearest-shell-spectra} imply the following
finite inclusion facts.

First, \(S_{2O}\subset 2O\setminus2T\).  Therefore the \(2O\) nearest-shell walk
alternates between the two cosets
\[
    2T,\qquad 2O\setminus2T,
\]
and the quotient two-state walk has eigenvalues \(1\) and \(-1\).

Second, under the inclusion \(2T\subset2I\), the set \(2I\) decomposes into five
right \(2T\)-cosets.  The \(2I\) nearest-shell walk induces on these five cosets the
complete-graph walk
\[
    \frac14(J_5-I_5),
\]
with quotient spectrum \(1^{(1)}\oplus(-\frac14)^{(4)}\).

Third, if \(T_{2T}\) is lifted to \(\ell^2(2I)\) by the same nearest-shell
convolution over the subgroup \(2T\subset2I\), then
\[
    [T_{2T},T_{2I}]=0,
\]
and the two kernels have the finite dimensions
\[
    \dim\ker T_{2I}=25,\qquad
    \dim(\ker T_{2T}\cap\ker T_{2I})=15.
\]
\end{corollary}

\begin{proof}
The first assertion follows because the nearest nontrivial elements of \(2O\) have
real part \(1/\sqrt2\), and these are precisely the six positive-real elements in
the octahedral coset \(2O\setminus2T\).  Multiplication by such a coset element
exchanges the two \(2T\)-cosets.  Thus the unnormalized two-coset incidence
matrix is
\[
  6\begin{pmatrix}0&1\\1&0\end{pmatrix},
\]
and normalization gives the two-point complete-graph walk.

For \(2T\subset2I\), the index is \(120/24=5\).  Directly from the displayed
\(2I\) shell, multiplication of the twelve shell elements against five right
\(2T\)-coset representatives gives the exact incidence matrix
\[
  3(J_5-I_5).
\]
Thus every coset is sent three times to each of the other four cosets and
never to itself, and normalization gives
\(\frac14(J_5-I_5)\).

The \(2I\) shell is conjugacy-invariant, so its group-algebra element is
central in \(\CC[2I]\).  It therefore commutes with the lifted
\(2T\)-shell element, proving \([T_{2T},T_{2I}]=0\) without a numerical
diagonalization.  The character certificate in the preceding theorem shows
that the zero eigenspace of \(T_{2I}\) is the regular isotypic block
\[
    V_5\otimes\CC^5,
\]
where \(V_5\) is the unique five-dimensional irreducible of \(2I\); hence
its dimension is \(5^2=25\).  On restriction to the common \(2T\) subgroup,
\[
    V_5\!\downarrow_{2T}
    \cong \chi_\omega\oplus\chi_{\bar\omega}\oplus V_3.
\]
Indeed, this is the lift of the elementary restriction of the
five-dimensional \(A_5\)-representation to a point-stabilizer \(A_4\):
the two nontrivial one-dimensional characters plus its three-dimensional
rotation representation.  The \(2T\)-shell eigenvalues on these summands are,
from the same certificate,
\[
    -\frac12,\qquad-\frac12,\qquad0.
\]
Consequently \(T_{2T}\) has a three-dimensional kernel on each copy of
\(V_5\), and the regular multiplicity of \(V_5\) is five.  Therefore
\[
    \dim(\ker T_{2T}\cap\ker T_{2I})=3\cdot5=15,
\]
which closes the joint-kernel calculation.
\end{proof}

\subsection{Exceptional visibility and recentered readouts over the Hurwitz core}
\label{p1:subsec:exceptional-core-visibility}

Let
\[
    H:=2T,
    \qquad
    I:=2I,
    \qquad
    N:=N_{SU(2)}(H)=2O
\]
be the fixed subgroups of
Theorem~\ref{p1:thm:exceptional-normalizer-fork}.  Thus
\[
    I\cap N=H,
    \qquad
    [I:H]=5,
    \qquad
    [N:H]=2,
    \qquad
    N/H\cong C_2,
\]
whereas \(H\) is not normal in \(I\).  The purpose of this subsection is to
make precise the sense in which the same fixed representative of the unmarked Hurwitz core has an
icosahedral global readout and an octahedral normalizer readout.  These are
two homogeneous representations over the common \(H\)-datum; they are not a
subgroup chain and they are not geometric zoom levels.

\begin{definition}[Homogeneous core readout and visible stabilizer]
\label{p1:def:homogeneous-core-readout}
For \(R\in\{H,N,I\}\), define the scalar homogeneous readout carrier
\[
    X_R:=R/H,
    \qquad
    \mathcal X_R:=\ell^2(X_R),
\]
with its standard orthonormal delta basis and the left permutation action
\[
    \lambda_R(g)\delta_{rH}:=\delta_{grH}.
\]
Its Hurwitz-relative readout rank is
\[
    \rho_H(R):=\dim_{\CC}\mathcal X_R=[R:H].
\]
For a readout vector \(v\in\mathcal X_R\), its visible stabilizer is the
pointwise stabilizer
\[
    \operatorname{Vis}_R(v)
    :=\{g\in R:\lambda_R(g)v=v\}.
\]
This definition concerns a vector, not merely its complex line.  In
particular, a quotient-odd line can be invariant as a line while its nonzero
vectors have the kernel of the quotient character as their visible
stabilizer.
\end{definition}

\begin{theorem}[Icosahedral global--local visibility and octahedral recentering]
\label{p1:thm:exceptional-core-visibility}
Put
\[
    X:=X_I=I/H,
    \qquad
    x_0:=H\in X,
    \qquad
    \mathcal X:=\CC[X].
\]
Then the following statements hold.

\begin{enumerate}[label=\textup{(\roman*)}]
\item The scalar \(I\)-module is the five-point permutation module
\begin{equation}
\label{p1:eq:five-cell-permutation-branching}
    \mathcal X\cong\mathbf1_I\oplus V_4,
\end{equation}
where \(V_4\) is the irreducible four-dimensional standard module of the
natural \(A_5\)-action.  The central element \(-1\in I\) acts trivially, so
this scalar action factors through \(I/\{\pm1\}\cong A_5\).

Let
\[
    u_0:=\delta_{x_0},
    \qquad
    u_{\mathrm{out}}
    :=\frac12\sum_{x\ne x_0}\delta_x.
\]
Both vectors have norm one.  The normalized global and marked-contrast
vectors are
\begin{equation}
\label{p1:eq:global-marked-icosahedral-vectors}
    u_I:=\frac{u_0+2u_{\mathrm{out}}}{\sqrt5},
    \qquad
    v_I:=\frac{2u_0-u_{\mathrm{out}}}{\sqrt5}.
\end{equation}
They satisfy
\[
    \operatorname{Vis}_I(u_I)=I,
    \qquad
    \operatorname{Vis}_I(v_I)=H.
\]
Thus the unmarked scalar readout has full \(2I\)-symmetry, whereas marking
the chosen Hurwitz cell reduces the pointwise stabilizer inside \(I\) exactly
to \(2T\).

\item If \(T_I\) denotes the five-coset walk of
Corollary~\ref{p1:cor:nearest-shell-inclusion-consequences}, then
\[
    T_I=\frac14(J_5-I_5).
\]
The global and contrast projectors are therefore the explicit spectral
projectors
\begin{equation}
\label{p1:eq:icosahedral-global-contrast-projectors}
    P_{I,\mathrm{glob}}
    =\frac15J_5
    =\frac15(I_5+4T_I),
    \qquad
    P_{I,\mathrm{con}}
    =I_5-P_{I,\mathrm{glob}}
    =\frac45(I_5-T_I),
\end{equation}
with ranks \(1\) and \(4\), respectively.  Moreover
\[
    \operatorname{Ran}P_{I,\mathrm{glob}}=\CC u_I,
    \qquad
    v_I\in\operatorname{Ran}P_{I,\mathrm{con}}.
\]

\item Under restriction to \(H\), the five-point set has exactly two
\(H\)-orbits, namely
\[
    \{x_0\},
    \qquad
    X\setminus\{x_0\}.
\]
Consequently
\begin{equation}
\label{p1:eq:H-fixed-recentered-plane}
    \mathcal X^H
    =\operatorname{span}_{\CC}\{u_0,u_{\mathrm{out}}\},
    \qquad
    \mathcal X\!\downarrow_H
    \cong\mathbf1_H\oplus\mathbf1_H\oplus V_3,
\end{equation}
where
\[
    V_3
    :=\left\{f\in\CC[X]:
      f(x_0)=0,\ 
      \sum_{x\ne x_0}f(x)=0
      \right\}
\]
is the three-dimensional tetrahedral module.

Let \(Y:=N/H=\{H,nH\}\), where \(nH\) is the unique nontrivial coset, and
define the unitary identification
\begin{equation}
\label{p1:eq:normalizer-transport-map}
    \Theta:\CC[Y]\longrightarrow\mathcal X^H,
    \qquad
    \Theta\delta_H=u_0,
    \qquad
    \Theta\delta_{nH}=u_{\mathrm{out}}.
\end{equation}
It is the unique unitary that sends the base coset to the selected-cell
vector and the nontrivial coset to the positive normalized sum over the
outer \(H\)-orbit.
The recentered normalizer action is the transported representation
\[
    \widetilde\lambda_N(g)
    :=\Theta\lambda_N(g)\Theta^{-1},
    \qquad g\in N.
\]
The image of the nontrivial element of \(N/H\) is the involution
\[
    S_Ou_0=u_{\mathrm{out}},
    \qquad
    S_Ou_{\mathrm{out}}=u_0.
\]
By Corollary~\ref{p1:cor:nearest-shell-inclusion-consequences}, this is also
the transported two-coset nearest-shell walk.
Its even and odd projectors and normalized eigenvectors are
\begin{equation}
\label{p1:eq:recentered-octahedral-projectors}
\begin{split}
    P_{O,\pm}
    &:=\frac12\bigl(I_{\mathcal X^H}\pm S_O\bigr),
    \\
    u_+&:=\frac{u_0+u_{\mathrm{out}}}{\sqrt2},
    \qquad
    u_-:=\frac{u_0-u_{\mathrm{out}}}{\sqrt2}.
\end{split}
\end{equation}
Both projectors have rank one.  For the transported action, define
\[
    \operatorname{Vis}^{\mathrm{rec}}_N(w)
    :=\{g\in N:\widetilde\lambda_N(g)w=w\}.
\]
Then
\[
    \operatorname{Vis}^{\mathrm{rec}}_N(u_+)=N,
    \qquad
    \operatorname{Vis}^{\mathrm{rec}}_N(u_-)=H.
\]
Equivalently, \(u_-\) carries the unique nontrivial character of
\(N/H\cong C_2\).

The icosahedral and recentered-octahedral bases of the same \(H\)-fixed plane
are related by
\begin{equation}
\label{p1:eq:icosahedral-octahedral-basis-change}
    u_+=\frac{3u_I+v_I}{\sqrt{10}},
    \qquad
    u_-=\frac{-u_I+3v_I}{\sqrt{10}}.
\end{equation}
In particular, neither octahedral parity line is the icosahedral global line.
The \(2O\)-readout is a transported representation obtained after recentering
on the selected \(H\)-cell; it is not the restriction of the \(2I\)-action to
a nonexistent subgroup \(2O\subset2I\).
\end{enumerate}
\end{theorem}

\begin{proof}
The action of \(I\) on \(I/H\) factors through
\(I/\{\pm1\}\cong A_5\) because the center lies in \(H\).  Under the natural
five-point action of \(A_5\), the stabilizer of \(x_0\) is
\(H/\{\pm1\}\cong A_4\).  The action is doubly transitive, so its permutation
module is the direct sum of the constant line and an irreducible
four-dimensional complement.  This proves
\eqref{p1:eq:five-cell-permutation-branching}.

The vector \(u_I\) is the normalized constant vector on \(X\), and hence is
fixed by all of \(I\).  Relative to the delta basis, \(v_I\) has one
coefficient \(2/\sqrt5\), at \(x_0\), and four equal coefficients
\(-1/(2\sqrt5)\).  Left translation moves the unique distinguished
coefficient from \(x_0\) to \(gH\).  Therefore
\(\lambda_I(g)v_I=v_I\) holds exactly when \(gH=H\), that is, exactly when
\(g\in H\).  This proves the two stabilizer statements.

The quotient-walk formula gives \(J_5=I_5+4T_I\).  The normalized constant
projector is \(J_5/5\); subtracting it from the identity gives the second
formula in \eqref{p1:eq:icosahedral-global-contrast-projectors}.  Its ranks
follow from the eigenvalue multiplicities
\(1^{(1)}\oplus(-1/4)^{(4)}\) of \(T_I\).

The point stabilizer \(A_4\) fixes \(x_0\) and acts transitively on the other
four points.  Hence an \(H\)-fixed function is determined by its value at
\(x_0\) and its common value on the outer orbit, proving
\eqref{p1:eq:H-fixed-recentered-plane}.  The zero-sum complement on that
four-point orbit is the standard irreducible three-dimensional tetrahedral
module.

Since \(H\triangleleft N\) and \(N/H\cong C_2\), the two-coset permutation
module is the regular module of \(C_2\).  The map \(\Theta\) sends its
orthonormal delta basis to the orthonormal basis
\(\{u_0,u_{\mathrm{out}}\}\), so it is unitary.  Transport of the nontrivial
quotient element gives the displayed swap.  Its \(+1\) and \(-1\)
eigenspaces are precisely \(\CC u_+\) and \(\CC u_-\).  The even vector is
fixed by \(N\), whereas the pointwise stabilizer of the odd vector is the
kernel of the sign character, namely \(H\).  Finally, substitution of
\eqref{p1:eq:global-marked-icosahedral-vectors} into
\eqref{p1:eq:recentered-octahedral-projectors} gives
\eqref{p1:eq:icosahedral-octahedral-basis-change}.
\end{proof}

The positive recentering also fixes the two rank-one contrast projectors'
overlap to \(9/10\).  Appendix~\ref{p1:app:v179-exceptional-two-plane}
derives the resulting quadratic pencil and dimensionless walk spectrum.


\begin{corollary}[Explicit exceptional readout ranks]
\label{p1:cor:exceptional-aggregation-order}
The common Hurwitz core has the three explicitly constructed homogeneous
readout carriers
\[
    \mathcal X_H=\CC[H/H],
    \qquad
    \mathcal X_N=\CC[N/H],
    \qquad
    \mathcal X_I=\CC[I/H],
\]
and therefore
\begin{equation}
\label{p1:eq:explicit-exceptional-readout-ranks}
    \rho_H(H)=1,
    \qquad
    \rho_H(N)=2,
    \qquad
    \rho_H(I)=5.
\end{equation}
The comparison \(5>2>1\) is a comparison of the dimensions of these displayed
carriers.  It is not a filtration, a subgroup chain, a temporal order, or a
partition of \(I\) containing \(N\setminus H\).
Only the displayed carriers and their explicit maps enter the conclusion.
\end{corollary}

\begin{proof}
The dimension of a finite homogeneous permutation carrier is the cardinality
of its coset space.  Apply
Theorem~\ref{p1:thm:exceptional-normalizer-fork} to obtain
\[
    [H:H]=1,
    \qquad
    [N:H]=2,
    \qquad
    [I:H]=5.
\]
The final exclusions follow from \(I\cap N=H\) and \(N\nsubseteq I\).
\end{proof}

\begin{theorem}[Projective-trimer resolution of the five-cell scalar carrier]
\label{p1:thm:five-cell-trimer-resolution}
Retain the base cell \(x_0=H\), and choose a projective order-three subgroup
\[
    \overline C\cong C_3
    \subset H/\{\pm1\}\cong A_4;
\]
equivalently, choose its binary preimage \(2C_3\subset H\).  On the outer
four-point orbit \(X\setminus\{x_0\}\), the group \(\overline C\) fixes one
point \(x_C\) and cyclically permutes the remaining three-point set
\(\mathcal O_C\).  Define
\begin{equation}
\label{p1:eq:trimer-axis-vector-and-plane}
    u_C
    :=\frac{1}{2\sqrt3}
      \left(3\delta_{x_C}
      -\sum_{x\in\mathcal O_C}\delta_x\right),
    \qquad
    E_C
    :=\left\{f\in\RR[X]:
    \begin{array}{l}
      f(x_0)=f(x_C)=0,\\[-0.2em]
      \displaystyle\sum_{x\in\mathcal O_C}f(x)=0
    \end{array}
    \right\}.
\end{equation}
Then \(u_C\) has norm one, \(E_C\) is a real two-plane, and the scalar carrier
has the orthogonal resolution
\begin{equation}
\label{p1:eq:one-one-one-two-resolution}
    \RR[X]
    =\RR u_+
     \oplus\RR u_-
     \oplus\RR u_C
     \oplus E_C,
    \qquad
    5=1+1+1+2.
\end{equation}
Equivalently, extend \(P_{O,+}\) and \(P_{O,-}\) by zero on
\(V_3\cap\RR[X]\), and set
\[
    P_C:=|u_C\rangle\langle u_C|,
    \qquad
    P_{C,\perp}
    :=I_{\RR[X]}-P_{O,+}-P_{O,-}-P_C.
\]
Then
\begin{equation}
\label{p1:eq:one-one-one-two-projectors}
    P_{O,+}+P_{O,-}+P_C+P_{C,\perp}=I_{\RR[X]},
    \qquad
    \bigl(
      \rank P_{O,+},
      \rank P_{O,-},
      \rank P_C,
      \rank P_{C,\perp}
    \bigr)
    =(1,1,1,2).
\end{equation}
After complexification,
\begin{equation}
\label{p1:eq:C3-five-cell-character-branching}
    \CC[X]\!\downarrow_{\overline C}
    \cong
    \mathbf1^{\oplus3}
    \oplus\chi_{\omega}
    \oplus\chi_{\bar\omega},
    \qquad
    E_C\otimes_{\RR}\CC
    \cong\chi_{\omega}\oplus\chi_{\bar\omega},
\end{equation}
where \(\omega=e^{2\pi i/3}\).  Thus the final rank-two term is a real
rotation plane whose complex eigenlines form a conjugate pair.
\end{theorem}

\begin{proof}
The action of \(H\) on the outer four points is the natural tetrahedral
\(A_4\)-action.  A three-cycle in this action fixes one point and permutes the
other three cyclically.  Hence \(x_C\) and \(\mathcal O_C\) have the stated
properties.

The coefficients of \(u_C\) sum to zero on the outer orbit, so \(u_C\) is
orthogonal to \(u_{\mathrm{out}}\), and it is also orthogonal to \(u_0\).
Moreover
\[
    \|u_C\|^2=\frac{9+3}{12}=1.
\]
The space \(E_C\) is the zero-sum subspace of the three-cycle permutation
module and therefore has real dimension two.  It is orthogonal to
\(u_0\), \(u_{\mathrm{out}}\), and \(u_C\).  Because
\[
    \RR u_+\oplus\RR u_-
    =\RR u_0\oplus\RR u_{\mathrm{out}},
\]
the four mutually orthogonal summands in
\eqref{p1:eq:one-one-one-two-resolution} have total dimension five and hence
exhaust \(\RR[X]\).  This proves the projector identities and ranks.

On \(E_C\), a generator of \(\overline C\) is a rotation through \(2\pi/3\).
Its complex eigenvalues are \(\omega\) and \(\bar\omega\), while the three
displayed real lines are fixed by \(\overline C\).  This gives
\eqref{p1:eq:C3-five-cell-character-branching}.
\end{proof}

\begin{proposition}[Center-sensitive induced completion fork]
\label{p1:prop:center-sensitive-induced-exceptional-fork}
Let \(W\) be a finite-dimensional complex \(H\)-module with representation
\(\rho_W\).  For \(R\in\{H,N,I\}\), define
\[
    \mathcal V_R(W):=\operatorname{Ind}_H^R W
    =\left\{f:R\to W:
      f(rh)=\rho_W(h^{-1})f(r)
      \text{ for }r\in R,\ h\in H
      \right\}.
\]
Then
\begin{equation}
\label{p1:eq:induced-completion-ranks}
    \dim\mathcal V_H(W)=\dim W,
    \qquad
    \dim\mathcal V_N(W)=2\dim W,
    \qquad
    \dim\mathcal V_I(W)=5\dim W.
\end{equation}
The functions supported on the base coset give canonical \(H\)-equivariant
embeddings
\begin{equation}
\label{p1:eq:induced-completion-span}
    \mathcal V_I(W)
    \xleftarrow{\ j_I\ }
    W
    \xrightarrow{\ j_N\ }
    \mathcal V_N(W).
\end{equation}
This span is the representation-theoretic exceptional fork over the same
\(H\)-owned datum.

If the central element acts on \(W\) as
\[
    \rho_W(-1)=-I_W,
\]
then it acts as \(-I\) on both induced completions.  Hence the induced fork is
genuinely center-sensitive and does not factor through \(A_5\) or \(S_4\).
In particular, for the defining two-dimensional \(SU(2)\)-module \(U\), the
tensor identity gives
\begin{equation}
\label{p1:eq:defining-module-induced-fork}
    \operatorname{Ind}_H^R(U|_H)
    \cong U|_R\otimes\CC[R/H],
    \qquad R\in\{N,I\},
\end{equation}
so the relative readout ranks remain \(2\) and \(5\) while the binary center
is retained by \(U\).
\end{proposition}

\begin{proof}
Choose left-coset representatives for \(R/H\).  An induced function is
determined freely by one vector of \(W\) on each representative, proving
\[
    \dim\operatorname{Ind}_H^R W=[R:H]\dim W
\]
and hence \eqref{p1:eq:induced-completion-ranks}.

For \(w\in W\), define \(j_Rw\) to vanish off \(H\subset R\) and, on the base
coset, set
\[
    (j_Rw)(h):=\rho_W(h^{-1})w,
    \qquad h\in H.
\]
For \(k\in H\), left translation gives
\[
    k\cdot j_R(w)=j_R\bigl(\rho_W(k)w\bigr),
\]
so \(j_R\) is \(H\)-equivariant.  This proves the span
\eqref{p1:eq:induced-completion-span}.

Let \(z=-1\).  Since \(z\) is central in \(R\) and belongs to \(H\), the
induced action satisfies
\[
    (z\cdot f)(r)
    =f(z^{-1}r)
    =f(rz^{-1})
    =\rho_W(z)f(r).
\]
Thus the central action on the induced module is the same as its action on
\(W\), proving center sensitivity.  The last isomorphism is the standard
tensor identity for induction applied to the \(R\)-module \(U|_R\).
\end{proof}

\subsection{Hurwitz-link carrier module}
\label{p1:sec:hurwitz-link-carrier}

The preceding 24-cell geometry supplies a local finite carrier module.  This
subsection records the finite supports and rational projector defects fixed by
the stated Hurwitz choices.

\begin{definition}[Hurwitz cube link]
\label{p1:def:hurwitz-cube-link}
At the Hurwitz vertex \(1\in\mathcal H^\times\), the eight adjacent vertices are
\[
    h_\varepsilon
    =\frac12(1+\varepsilon_1 i+\varepsilon_2 j+\varepsilon_3 k),
    \qquad
    \varepsilon=(\varepsilon_1,\varepsilon_2,\varepsilon_3)\in\{\pm1\}^3.
\]
We denote this neighbor set by \(N(1)\).  After left multiplication by an
arbitrary Hurwitz unit, the same construction gives the neighbor set \(N(v)\) at
any vertex \(v\).
\end{definition}

\begin{proposition}[The link is a sign cube]
\label{p1:prop:hurwitz-link-cube}
For each Hurwitz vertex \(v\), the edge-neighbor set \(N(v)\) is, relative to that chosen vertex, a
cube canonical up to the finite cube symmetry.  At \(v=1\) the coordinates above identify
\[
    N(1)\cong\{\pm1\}^3\cong(\ZZ_2)^3
\]
as a set with the standard sign cube. This is a sign-coordinate identification: geometrically, $N(1)$ lies in the affine hyperplane $\operatorname{Re}(q)=1/2$, with imaginary coordinates $(\varepsilon_1,\varepsilon_2,\varepsilon_3)/2$, and its cube side length is $1$ in the ambient Euclidean metric.
\end{proposition}

\begin{proof}
The adjacency statement is Proposition~\ref{p1:prop:24cell-adjacency}.  The eight
neighbors of \(1\) are exactly the displayed half-integer Hurwitz units.  Their
three independent signs give the standard cube coordinates.  Left multiplication
by a Hurwitz unit is an isometry of the unit-quaternion sphere and carries the
link at \(1\) to the link at any vertex \(v\).
\end{proof}

\begin{theorem}[Quaternionic Walsh--Clifford carrier]
\label{p1:thm:hurwitz-walsh-clifford}
Let \(G=(\ZZ_2)^3\), write \(a=(a_1,a_2,a_3)\in\{0,1\}^3\), and let
\(\chi_a\) be the Walsh basis of real functions on \(N(1)\cong G\).  Define the
quaternionic Clifford cocycle
\[
    \sigma_-(a,b)
    =(-1)^{\sum_i a_i b_i+\sum_{i>j}a_i b_j},
\]
and the twisted product
\[
    \chi_a\star_{\sigma_-}\chi_b=\sigma_-(a,b)\chi_{a+b},
    \qquad a+b\text{ taken mod }2.
\]
Then
\[
    \bigl(\RR^{N(1)},\star_{\sigma_-}\bigr)\cong \Cl(0,3).
\]
The product \(\star_{\sigma_-}\) is not the pointwise product of functions on the cube and is not induced by the cube as a bare finite set. It is the choice-relative twisted group-algebra product determined by the ordered sign coordinates and the quaternionic Clifford cocycle.
In particular the Hurwitz cube link realizes the protected eight-dimensional
Clifford carrier and its scalar-free seven-dimensional vector-space quotient:
\[
    \dim_\RR\RR^{N(1)}=8,
    \qquad
    \dim_\RR\bigl(\RR^{N(1)}/\RR\mathbf1\bigr)=7.
\]
\end{theorem}

\begin{proof}
\ProofPlan{We verify the carrier in three steps: associativity of the twisted Walsh product, the Clifford generator relations, and the quaternionic realization.}

\begin{proofsteps}
\item \textbf{Associativity from the cocycle identity.}
The exponent defining \(\sigma_-(a,b)\) is bilinear over \(\mathbb F_2\).  Hence
\[
  \sigma_-(a,b)\sigma_-(a+b,c)
  =
  \sigma_-(b,c)\sigma_-(a,b+c),
\]
so the twisted product is associative and has unit \(\chi_0\).

\item \textbf{Clifford generator relations.}
Let
\(e_i:=\chi_{\mathbf e_i}\), where \(\mathbf e_i\) is the \(i\)-th coordinate
vector in \((\ZZ_2)^3\).  The diagonal factor in \(\sigma_-\) gives
\[
    e_i\star_{\sigma_-} e_i=-1.
\]
For \(i\ne j\), the triangular factor \(\sum_{i>j}a_i b_j\) gives
\[
    e_i\star_{\sigma_-} e_j=-e_j\star_{\sigma_-}e_i.
\]
Thus the generators satisfy the defining relations of \(\Cl(0,3)\).  The eight
basis elements \(\chi_a\), indexed by subsets of \(\{1,2,3\}\), give the usual
vector-space grading
\[
    1+3+3+1
    =\dim\Lambda^0\RR^3+\dim\Lambda^1\RR^3+\dim\Lambda^2\RR^3+
      \dim\Lambda^3\RR^3.
\]

\item \textbf{Quaternionic realization.}
The cocycle is not an arbitrary sign choice: it is the Walsh form of the local
quaternionic multiplication rules
\[
    i^2=j^2=k^2=-1,
    \qquad
    ij=-ji,
    \quad jk=-kj,
    \quad ki=-ik.
\]
Equivalently, the assignment
\[
    e_1\mapsto(i,i),\qquad e_2\mapsto(j,j),\qquad e_3\mapsto(k,-k)
\]
realizes the algebra in \(\HH\oplus\HH\), with volume element sent to the central
splitter \((-1,1)\).  Hence the twisted Walsh algebra is the framework's
\(\Cl(0,3)\) carrier.
\end{proofsteps}
\end{proof}

\begin{corollary}[Full nonconstant Walsh evaluation support]
\label{p1:cor:walsh-evaluation-support-56}
Let
\[
    G=(\ZZ_2)^3,
    \qquad
    \widehat G:=\Hom(G,\{\pm1\}),
    \qquad
    \widehat G_*:=\widehat G\setminus\{\mathbf1\}.
\]
Then
\[
    |G|=8,
    \qquad
    |\widehat G_*|=7,
\]
and the full nonconstant Walsh evaluation support
\[
    \mathcal E_*:=G\times\widehat G_*
\]
has cardinality
\[
    |\mathcal E_*|=8\cdot7=56.
\]
Equivalently, the table of all nonconstant Walsh characters evaluated at all
cube-link points has \(56\) entries, and
\[
    \RR^{\mathcal E_*}
    \cong
    \Hom_\RR\!\left(\RR^{\widehat G_*},\RR^G\right)
\]
has real dimension \(56\).  For every subset
\(B\subset\widehat G_*\) with \(|B|=q\), the partial support
\(G\times B\) has cardinality \(8q\).
\end{corollary}

\begin{proof}
The character group of a finite elementary abelian group of rank three is again
an elementary abelian group of rank three, hence has eight characters.  Removing
the trivial character leaves seven.  The product and function-space counts are
then immediate.
\end{proof}

\begin{remark}[Symmetry status of the intermediate octet rungs]
\label{p1:rem:walsh-flag-noncanonical}
The full automorphism group
\[
    \Aut(G)\cong GL(3,2)
\]
acts transitively on the seven nontrivial characters.  Consequently no
nonempty proper subset of \(\widehat G_*\) is invariant under the full group
structure.  The endpoint \(56\) is therefore intrinsic to the full evaluation
support, whereas a one-character-at-a-time ladder
\(8,16,\ldots,56\) requires a chosen ordering or flag.  Every proper partial
support depends on the selected subset of nontrivial characters.
\end{remark}

In the trimer and antiprism statements below we use the rescaled sign-cube coordinates
\[
(\varepsilon_1,\varepsilon_2,\varepsilon_3)\in\{\pm1\}^3,
\]
rather than the actual imaginary Hurwitz-link coordinates
\[
(\varepsilon_1/2,\varepsilon_2/2,\varepsilon_3/2).
\]
Consequently, displayed cube lengths in this subsection are sign-cube lengths; the corresponding ambient Hurwitz-link lengths are smaller by a factor of \(2\).

\begin{theorem}[Diagonal trimer and antiprism decomposition]
\label{p1:thm:hurwitz-diagonal-trimer}
Choose an oriented body diagonal of the cube link.  For the representative
\[
    u_+=(+,+,+),\qquad u_-=(-,-,-),
\]
the link decomposes as
\[
    N(1)=\{u_+,u_-\}\sqcup T_+\sqcup T_-,
\]
where
\[
    T_+=\{(-,+,+),(+,-,+),(+,+,-)\},
\]
\[
    T_-=\{(+,-,-),(-,+,-),(-,-,+)\}.
\]
The two triples lie in the parallel planes
\[
    x+y+z=+1,
    \qquad
    x+y+z=-1,
\]
and each is an equilateral triangle.  The six non-axial vertices form a
\(D_{3d}\)-symmetric staggered triangular antiprism.  It is not a regular
octahedron in the raw cube metric; its top/bottom triangle edges have length
\(\sqrt8\), while the short cross-edges have length \(2\).
\end{theorem}

\begin{proof}
The displayed sets are the vertices of the sign cube after deleting the chosen
opposite pair.  The sum \(x+y+z\) equals \(+1\) on \(T_+\) and \(-1\) on
\(T_-\), proving the two-plane statement.  Pairwise distances inside either
triple equal \(\sqrt8\), so each triple is equilateral.  The projection along the
body diagonal alternates the two triples around a regular hexagon, producing the
staggered triangular-antiprism combinatorics.  The edge lengths stated above
show that this antiprism is not regular in the ambient cube metric.
\end{proof}

\begin{theorem}[Oriented axial--opposite-trimer quaternionic frame]
\label{p1:thm:axial-opposite-trimer-quaternionic-frame}
For the Hurwitz-link quaternions
\[
 h_\varepsilon
 =\frac12(1+\varepsilon_1i+\varepsilon_2j+\varepsilon_3k),
 \qquad \varepsilon\in\{\pm1\}^3,
\]
equip \(\HH\) with
\(\langle a,b\rangle_{\HH}:=\operatorname{Re}(a\bar b)\).
Choose the oriented axial vertex \(u_+=(+,+,+)\), and let \(T_-\) be the
opposite trimer of Theorem~\ref{p1:thm:hurwitz-diagonal-trimer}.  Then
\[
 \mathcal E_{u_+}:=
 \{h_{u_+}\}\sqcup\{h_\eta:\eta\in T_-\}
\]
is an orthonormal real basis of \(\HH\).  Left multiplication by
\(\bar h_{u_+}\) is an orientation-preserving isometry and gives
\[
 L_{\bar h_{u_+}}h_{u_+}=1,
 \qquad
 L_{\bar h_{u_+}}
 \operatorname{span}_{\RR}\{h_\eta:\eta\in T_-\}
 =\operatorname{Im}\HH.
\]
Consequently the oriented link itself generates the complex carrier split
\begin{equation}
\label{p1:eq:generated-oriented-one-plus-three}
 Q_{u_+}
 :=\operatorname{span}_{\CC}\mathcal E_{u_+}
 =\ell_{u_+}\oplus C_{u_+}
 \cong \CC1\oplus\operatorname{Im}(\HH)_{\CC},
 \qquad
 \dim_{\CC}(\ell_{u_+},C_{u_+})=(1,3),
\end{equation}
where
\[
 \ell_{u_+}:=\CC h_{u_+},
 \qquad
 C_{u_+}:=\operatorname{span}_{\CC}\{h_\eta:\eta\in T_-\}.
\]
The remaining trimer supplies the independent rank-three support
\begin{equation}
\label{p1:eq:generated-oriented-complementary-trimer}
 G_{u_+}:=\CC^{T_+}\cong\CC^3.
\end{equation}
Thus, relative to the fixed oriented axial branch, the incidence
\(Q=\ell\oplus C\) and the complementary multiplicity trimer are generated
by orthogonality and complement; they are not an additional assignment.
Changing the axial representative by an orientation-preserving Hurwitz/cube
symmetry gives a unitarily equivalent pair.  Reversing the orientation gives
the conjugate representative and exchanges the two displayed trimer
presentations, not a second unoriented constitution.
\end{theorem}

\begin{proof}
For sign vectors \(\varepsilon,\delta\in\{\pm1\}^3\), direct quaternionic
multiplication gives
\begin{equation}
\label{p1:eq:hurwitz-link-inner-product}
 \langle h_\varepsilon,h_\delta\rangle_{\HH}
 =\frac{1+\varepsilon\cdot\delta}{4}.
\end{equation}
Every \(\eta\in T_-\) differs from \(u_+\) in two signs, so
\(u_+\cdot\eta=-1\), and distinct members of \(T_-\) also have dot product
\(-1\).  Equation~\eqref{p1:eq:hurwitz-link-inner-product} therefore makes
the four displayed unit quaternions pairwise orthogonal.  They form a real
basis of the four-dimensional space \(\HH\).

Left multiplication by a unit quaternion is an orientation-preserving real
isometry.  It sends \(h_{u_+}\) to \(1\), hence sends its three-dimensional
orthogonal complement onto \(1^\perp=\operatorname{Im}\HH\).  Complexifying
proves \eqref{p1:eq:generated-oriented-one-plus-three}.  The diagonal-trimer
decomposition leaves exactly \(T_+\) after the axial line and opposite trimer
have been used, proving \eqref{p1:eq:generated-oriented-complementary-trimer}.
The final equivalence statements follow because the oriented cube/Hurwitz
symmetry acts isometrically and orientation reversal conjugates the
quaternionic frame.
\end{proof}

\begin{theorem}[Rooted trimer embeddings and inverse incidence]
\label{p1:thm:rooted-trimer-vector-overlap}
Use the unit Hurwitz coordinates in
Theorem~\ref{p1:thm:axial-opposite-trimer-quaternionic-frame},
with root \(v=1=(1,0,0,0)\), oriented axial pair
\[
 a_+=\tfrac12(1,1,1,1),\qquad
 a_-=\tfrac12(1,-1,-1,-1),
\]
and the ordered complementary trimers
\[
\begin{aligned}
 p_0&=\tfrac12(1,-1,1,1),&
 p_1&=\tfrac12(1,1,-1,1),&
 p_2&=\tfrac12(1,1,1,-1),\\
 m_0&=\tfrac12(1,1,-1,-1),&
 m_1&=\tfrac12(1,-1,1,-1),&
 m_2&=\tfrac12(1,-1,-1,1).
\end{aligned}
\]
Let \(G_{\mathbb R}=\mathbb R^3\) have its counting inner product,
and let \(V=\mathbb H\) have its real Euclidean inner product.
Define
\[
 L_+e_j=p_j,\qquad L_-e_j=m_j,
 \qquad L_\pm:G_{\mathbb R}\longrightarrow V.
\]
The root and metric determine the orthogonal involution
\[
 \kappa_v(q)=2\langle v,q\rangle v-q.
\]
On the root link it is \(q\mapsto v-q\), and therefore
\(\kappa_v(p_j)=m_j\), \(\kappa_v(a_+)=a_-\), and
\(\kappa_vL_+=L_-\). This pairing uses no choice of a short-edge
half-step.

Put \(J=\mathbf1\mathbf1^{\mathsf T}\),
\(Q_0=J/3\), and \(Q_\perp=I-Q_0\).
Then \(L_+\) and \(L_-\) are isometries, with
\[
 L_\pm^{\mathsf T}L_\pm=I,\qquad
 L_+L_+^{\mathsf T}=I-a_-a_-^{\mathsf T},\qquad
 L_-L_-^{\mathsf T}=I-a_+a_+^{\mathsf T}.
\]
In the aligned labels, the cross-trimer short-edge incidence is
\[
 (C_\triangle)_{ij}
 =\mathbf1_{\{\langle p_i,m_j\rangle=1/2\}},
 \qquad C_\triangle=J-I=X+X^{-1},
 \qquad Xe_j=e_{j+1\pmod3}.
\]
Its inverse is the signed metric overlap:
\begin{equation}
\label{p1:eq:trimer-vector-inverse-incidence}
 K_V:=L_+^{\mathsf T}L_-
 =\tfrac12J-I=\tfrac12Q_0-Q_\perp=C_\triangle^{-1}.
\end{equation}
If \(u=(1,1,1)/\sqrt3\) and
\(P=\{0\}\oplus u^\perp\subset V\), then
\[
 \operatorname{im}L_+\cap\operatorname{im}L_-=P,
 \qquad L_+(a)=(0,-a),\quad L_-(a)=(0,a)
 \quad(a\in Q_\perp G_{\mathbb R}).
\]
Thus the augmentation plane is isometrically realized as the common
two-plane of the two trimer spans.
\end{theorem}

\begin{proof}
Every link point has \(\langle v,q\rangle=1/2\), which gives the
formula for \(\kappa_v\). The displayed coordinates give
\[
 L_+=\begin{pmatrix}\tfrac12\mathbf1^{\mathsf T}\\J/2-I\end{pmatrix},
 \qquad
 L_-=\begin{pmatrix}\tfrac12\mathbf1^{\mathsf T}\\I-J/2\end{pmatrix}.
\]
Using \(J^2=3J\), their self-Gram matrices are \(I\) and their
cross-Gram matrix is \(J/2-I\). Orthogonality to the respective
axial vector gives the range projectors. The cross-Gram entries
are \(-1/2\) on the diagonal and \(1/2\) off the diagonal,
so the short-edge incidence is \(J-I\). Furthermore,
\[
 (J-I)(J/2-I)=I.
\]
The two axial vectors span \(\operatorname{span}\{1,(0,u)\}\);
their common orthogonal complement is \(P\). Setting \(Ja=0\)
in the displayed matrices gives the final formulas.
\end{proof}

The construction is covariant under a common orthogonal transport of
the rooted configuration: \(\kappa_{Ov}=O\kappa_vO^{-1}\).
A simultaneous cyclic relabelling conjugates the overlap matrices.
Consequently the formulas depend on the existing rooted geometric
data, not on the chosen coordinate representative.

\begin{corollary}[Exterior-square overlap and reciprocal returns]
\label{p1:cor:rooted-trimer-bivector-overlap}
Under the hypotheses of
Theorem~\ref{p1:thm:rooted-trimer-vector-overlap},
orient \(G_{\mathbb R}\) by the rooted cyclic ordering and let
\(\star_G:G_{\mathbb R}\to\Lambda^2G_{\mathbb R}\) be its real
Hodge isometry. Explicitly,
\[
 \star_Ge_0=e_1\wedge e_2,\qquad
 \star_Ge_1=e_2\wedge e_0,\qquad
 \star_Ge_2=e_0\wedge e_1.
\]
Equip \(\Lambda^2V\) with the induced Euclidean inner product and set
\[
 B_\pm=(\Lambda^2L_\pm)\star_G:
 G_{\mathbb R}\longrightarrow\Lambda^2V.
\]
Both maps are isometries. Their overlap satisfies
\begin{equation}
\label{p1:eq:trimer-bivector-half-incidence}
 K_{\Lambda^2}:=B_+^{\mathsf T}B_-
 =\tfrac12C_\triangle=Q_0-\tfrac12Q_\perp,
 \qquad K_VK_{\Lambda^2}=\tfrac12I.
\end{equation}
Consequently the positive overlap returns are
\begin{equation}
\label{p1:eq:trimer-reciprocal-positive-returns}
 H_V:=K_V^{\mathsf T}K_V=\tfrac14Q_0+Q_\perp,
 \qquad
 H_{\Lambda^2}:=K_{\Lambda^2}^{\mathsf T}K_{\Lambda^2}
 =Q_0+\tfrac14Q_\perp=\tfrac14H_V^{-1}.
\end{equation}
These statements also hold after complexification, with transpose
replaced by Hermitian adjoint.
\end{corollary}

\begin{proof}
Exterior powers of a real isometry are isometries. Functoriality
of the exterior square and its compatibility with transpose give
\[
 B_+^{\mathsf T}B_-
 =\star_G^{-1}\Lambda^2(L_+^{\mathsf T}L_-)\star_G.
\]
For an invertible real three-dimensional map \(K\), the cofactor
identity is
\[
 \star_G^{-1}(\Lambda^2K)\star_G=(\det K)K^{-\mathsf T}.
\]
Here \(K_V\) is symmetric with eigenvalues \(1/2,-1,-1\), so
\(\det K_V=1/2\) and \(K_V^{-1}=C_\triangle\).
This proves the overlap formula. Squaring on the orthogonal
projectors \(Q_0,Q_\perp\) proves the positive-return identities.
\end{proof}

The vector embeddings extend to the orthogonal frames
\((a_-,p_0,p_1,p_2)\) and \((a_+,m_0,m_1,m_2)\).
Their exterior squares therefore give explicit maps into the
vector and bivector spaces already realized by the rooted tetrad
and coordinate-pair constructions. The factors \(1/2\) and \(1/4\)
are fixed by the unit Hurwitz and counting metrics.

The vector and bivector image pairs have principal-angle lists
\((\pi/3,0,0)\) and \((0,\pi/3,\pi/3)\), respectively;
their incidence and overlap maps share a polar factor while
retaining different positive factors. Appendix~
\ref{p1:app:trimer-overlap-spectra} records these comparisons
and the cyclic character decompositions.

\begin{proposition}[Projected metric invariants of the diagonal trimer]
\label{p1:prop:projected-trimer-metric}
Normalize the sign cube by writing
\[
\widehat\varepsilon
:=\frac1{\sqrt3}(\varepsilon_1,\varepsilon_2,\varepsilon_3),
\qquad
\nu:=\frac1{\sqrt3}(1,1,1).
\]
For every non-axial vertex \(\widehat\varepsilon\in T_+\sqcup T_-\), the absolute
height over the plane \(\nu^\perp\) is
\[
 h=|\widehat\varepsilon\cdot\nu|=\frac13.
\]
Therefore the projected six-point orbit in \(\nu^\perp\) has circumradius
\[
R_{\rm hex}=\sqrt{1-h^2}=\frac{2\sqrt2}{3},
\]
forms a regular hexagon, and each alternating trimer has inradius
\[
r_\triangle=\frac12 R_{\rm hex}=\frac{\sqrt2}{3}.
\]
\end{proposition}

\begin{proof}
For a sign vertex,
\[
\widehat\varepsilon\cdot\nu
=\frac{\varepsilon_1+\varepsilon_2+\varepsilon_3}{3}.
\]
On the six non-axial vertices the numerator is \(\pm1\), hence the height is
\(1/3\).  Since \(|\widehat\varepsilon|=1\), orthogonal projection to \(\nu^\perp\)
has radius \(\sqrt{1-1/9}=2\sqrt2/3\).  Cyclically ordering the six projected
vertices gives the regular hexagon already used in the proof of
Theorem~\ref{p1:thm:hurwitz-diagonal-trimer}.  The alternating vertices of a
regular hexagon form an equilateral triangle whose inradius is one half of the
hexagon circumradius, giving \(r_\triangle=\sqrt2/3\).
\end{proof}

\begin{proposition}[Choice-relative antiprism and transport-prism matchings]
\label{p1:prop:antiprism-transport-matchings}
Relative to the chosen Hurwitz vertex and oriented cube-link diagonal,
the native six-point metric support is the staggered antiprism
\(T_+\sqcup T_-\).
The root involution of
Theorem~\ref{p1:thm:rooted-trimer-vector-overlap}
pairs opposite vertices of its projected hexagon.
This pairing does not choose a nearest-neighbor transport.
After quotienting by common cyclic relabelling of both trimers,
there are exactly two nearest-neighbor cyclic matchings
\[
 T_+\longrightarrow T_-,\qquad +60^\circ\quad\text{or}\quad-60^\circ.
\]
The axial inversion interchanges these two nearest-neighbor
matching classes.
\end{proposition}

\begin{proof}
Projection along the chosen body diagonal gives a regular hexagon with vertices
alternating between \(T_+\) and \(T_-\).  A nearest-neighbor matching that respects
the cyclic trimer order must send each vertex of \(T_+\) to one of its two
adjacent \(T_-\)-vertices in this hexagon.  These are the two possible half-step
choices, clockwise and counterclockwise.  Simultaneously rotating both trimers by
the common \(C_3\) action only changes labels, so the two nearest-neighbor
half-step choices are the only two classes under this short-edge condition.
Reversing the oriented body diagonal swaps the two trimers and reverses the
hexagon orientation, hence exchanges the two classes.
\end{proof}

\begin{proposition}[Diagonal equivalence and stabilizer]
\label{p1:prop:diagonal-equivalence}
The rotational symmetry group of the cube acts transitively on the four
unoriented body diagonals.  The stabilizer of one unoriented diagonal is
\(S_3\cong D_3\).  After orienting the diagonal, the oriented-diagonal-preserving
stabilizer is the cyclic subgroup \(C_3\subset S_3\); the remaining three
elements of the \(S_3\)-stabilizer reverse the oriented diagonal and exchange
\(T_+\) with \(T_-\).
\end{proposition}

\begin{proof}
The rotation group of the cube is isomorphic to \(S_4\), acting on the four
unoriented body diagonals.  The action is transitive, and the stabilizer of one
diagonal has order \(24/4=6\), hence is \(S_3\).  In the coordinates of
Theorem~\ref{p1:thm:hurwitz-diagonal-trimer}, the order-three subgroup cyclically
permutes the three off-axis coordinates and preserves the chosen orientation,
while the remaining three elements reverse the oriented diagonal and swap the two
planes \(x+y+z=\pm1\).
\end{proof}

\begin{theorem}[Binary lifts of cube-diagonal stabilizers]
\label{p1:thm:binary-dihedral-diagonal-lift}
Let \(\pi:2O\to2O/\{\pm1\}\cong S_4\) be the projective quotient acting as
the rotational symmetry group of the cube link.  For the chosen body diagonal,
let \(C_3\) be its oriented stabilizer and \(D_3\cong S_3\) its unoriented
stabilizer.  Then
\[
    \pi^{-1}(C_3)=2C_3\cong C_6^{\rm grp},
    \qquad
    \pi^{-1}(D_3)=2D_3,
\]
where \(|2D_3|=12\).  Moreover,
\[
    N_{2O}(2C_3)=2D_3,
    \qquad
    2C_3\triangleleft2D_3,
    \qquad
    2D_3/2C_3\cong C_2,
    \qquad
    2D_3/\{\pm1\}\cong D_3\cong S_3.
\]
A presentation is
\[
    2D_3
    =\left\langle r,s\ \middle|\
      r^6=1,\ s^2=r^3,\ srs^{-1}=r^{-1}
      \right\rangle,
\]
with \(2C_3=\langle r\rangle\).
Here \(C_6^{\rm grp}\) denotes the cyclic group of order six; it is not the
six-cycle graph \(C_6^{\rm link}\) used below.
\end{theorem}

\begin{proof}
Proposition~\ref{p1:prop:diagonal-equivalence} identifies the oriented and
unoriented projective stabilizers as \(C_3\) and \(D_3\).  Taking their full
preimages through the double cover doubles their orders and gives the binary
cyclic and binary dicyclic groups.  The normalizer of a three-cycle subgroup in \(S_4\) is its
\(D_3\cong S_3\) stabilizer, so taking full preimages also gives
\(N_{2O}(2C_3)=2D_3\).  Since \(C_3\triangleleft D_3\) with quotient
\(C_2\), the binary preimages satisfy
\(2C_3\triangleleft2D_3\) with the same quotient.  The displayed presentation
is the standard binary-dihedral presentation for \(n=3\).
\end{proof}

\begin{corollary}[Generated diagonal-stabilizer and six-cycle bridge]
\label{p1:cor:generated-c6-stabilizer-link-bridge}
Fix the unmarked support class \(\mathsf{Supp}_{2T}\). Let
\(\sigma=(H,v,\delta)\) range over the complete orbit of calculation
representatives consisting of \(H\in[2T]\), a vertex \(v\in H\), and an
oriented body diagonal \(\delta=(u_+,u_-)\) in the cube link \(N(v)\).
For every such representative define
\[
 G_\delta:=\pi^{-1}(\operatorname{Stab}^{+}(\delta))
      \cong C_6^{\rm grp},
 \qquad
 C_6^{\rm link}(\sigma)
      :=N(v)[N(v)\setminus\{u_+,u_-\}],
\]
where the second expression is the induced link graph.  Then
\[
 V\bigl(C_6^{\rm link}(\sigma)\bigr)=T_+\sqcup T_-,
\]
and \(C_6^{\rm link}(\sigma)\) is a six-cycle whose alternating vertex
classes are the two oriented trimers.  The complete unmarked families
\([G_\delta]\) and \([C_6^{\rm link}(\sigma)]\) are equivariant structural
descendants determined by \(\mathsf{Supp}_{2T}\) (using its canonical
normalizer class \(\mathsf{Supp}_{2O}\)). A labelled vertex and oriented
diagonal specify a representative of this equivariant family.
\end{corollary}

\begin{proof}
The Hurwitz support makes each vertex link a cube, and ambient symmetry is
transitive on the complete family of vertices and body diagonals.
Theorem~\ref{p1:thm:binary-dihedral-diagonal-lift} identifies the oriented
stabilizer lift with \(C_6^{\rm grp}\).  Removing the two axial vertices from
the cube leaves the induced six-cycle of
Theorem~\ref{p1:thm:hurwitz-diagonal-trimer}, with alternating classes
\(T_+\) and \(T_-\).  All constructions commute with the support
automorphisms, so retaining their complete orbit adds no mark and proves the
generated-descendant statement.
\end{proof}

\begin{theorem}[Binary-dihedral Weyl skeleton of the trimer]
\label{p1:thm:binary-dihedral-a2-weyl}
The projected six non-axial cube-link vertices form, after an overall rescaling,
the root hexagon \(\Phi(A_2)\).  Its Weyl group is
\[
    W(A_2)\cong D_3\cong S_3
    \cong2D_3/\{\pm1\}.
\]
For one oriented trimer \(T=\{a,b,c\}\), let
\(C:=\CC^T\).  Then
\[
    \mathfrak{sl}(C)
    =\mathfrak h_{A_2}
     \oplus
     \bigoplus_{x\ne y}\CC E_{xy},
\]
where \(\dim_\CC\mathfrak h_{A_2}=2\), the six off-diagonal lines have weights
\(e_x-e_y\in\Phi(A_2)\), and therefore
\[
    \dim_\CC\mathfrak{sl}(C)=2+6=8.
\]
The compact real form \(\mathfrak{su}(C)\) has real dimension \(8\).  Thus
\(2D_3\) is the binary lift of the finite Weyl symmetry organizing the trimer's
rank-two root data; it is not itself the continuous group \(SU(3)\).
\end{theorem}

\begin{proof}
By Proposition~\ref{p1:prop:projected-trimer-metric}, the six projected points
are the vertices of a regular centered hexagon, which is the \(A_2\) root system
up to scale.  The subgroup generated by reflections in the root hyperplanes is
the Weyl group \(W(A_2)\cong D_3\cong S_3\), and
Theorem~\ref{p1:thm:binary-dihedral-diagonal-lift} identifies its binary lift.
(The full automorphism group of the unoriented hexagon is larger by the central
inversion; no identification with that larger group is used here.)  The standard matrix-unit decomposition of
\(\mathfrak{sl}_3(\CC)\) has a two-dimensional diagonal Cartan subalgebra and
six one-dimensional off-diagonal root spaces, indexed by the ordered pairs
\(x\ne y\).  The dimension statements follow.
\end{proof}

\begin{proposition}[Two orientations of the binary-tetrahedral extension]
\label{p1:prop:two-oriented-2T-presentations}
Let \(\varphi\in\Aut(Q_8)\) cyclically permute the three quaternionic axes.
Then
\[
    Q_8\rtimes_{\varphi}C_3\cong2T,
    \qquad
    Q_8\rtimes_{\varphi^{-1}}C_3\cong2T.
\]
The two actions correspond to the two choices of orientation of the cyclic
generator and are exchanged by an automorphism of \(Q_8\) conjugating
\(\varphi\) to \(\varphi^{-1}\).  Consequently the two oriented presentations
are isomorphic copies of the same abstract group.
\end{proposition}

\begin{proof}
The standard Hurwitz multiplication gives
\(2T\cong Q_8\rtimes C_3\), with the generator of \(C_3\) cyclically permuting
\(i,j,k\).  Replacing that generator by its inverse replaces \(\varphi\) by
\(\varphi^{-1}\).  The two order-three automorphisms are conjugate in
\(\Aut(Q_8)\cong S_4\), so the corresponding semidirect products are isomorphic.
\end{proof}

\begin{theorem}[Cyclic trimer representation and rank-three operator space]
\label{p1:thm:cyclic-trimer-operator-grammar}
Fix the Hurwitz vertex, oriented cube-link diagonal, and induced cyclic order on the trimers.  Let \(T\) be either oriented trimer \(T_+\) or \(T_-\), and let \(X\) denote the
orientation-preserving cyclic shift of its three vertices.  Then
\[
\CC^{T}\cong\CC[Z_3]
\]
as a \(Z_3\)-module.  In a cyclic vertex basis the distinguished diagonal defect
plane is spanned by
\[
\DtriI=\operatorname{diag}\!\left(\frac23,-\frac13,-\frac13\right),
\qquad
\DtriII=\operatorname{diag}(0,1,-1),
\]
and the cyclic transport plane is spanned by
\[
C_1=X+X^\dagger,
\qquad
C_2=-i(X-X^\dagger).
\]
Consequently the primitive first-order Hermitian operator space associated with the trimer
projectors and one cyclic transport orbit is
\[
\mathfrak T_{\rm tri}
=\operatorname{span}_{\RR}\{I,\DtriI,\DtriII,C_1,C_2\},
\]
so every element of this five-dimensional real operator space has the form
\[
H=\alpha I+\beta \DtriI+\gamma \DtriII+\rho C(\phi),
\qquad
C(\phi)=\cos\phi\,C_1+\sin\phi\,C_2
       =e^{-i\phi}X+e^{i\phi}X^2,
\]
with \(\alpha,\beta,\gamma\in\RR\), \(\rho\ge0\), and \(\phi\in\RR/2\pi\ZZ\)
(for instance \(0\le\phi<2\pi\)).  Here \((\rho,\phi)\) are polar coordinates on
the two-dimensional real transport plane; the representation is unique for
\(\rho>0\), while the zero transport vector \(\rho=0\) has arbitrary \(\phi\).
\end{theorem}

\begin{proof}
The oriented trimer is a free transitive \(Z_3\)-set under the cyclic subgroup of
the diagonal stabilizer, so its complex span is the regular representation
\(\CC[Z_3]\).  The three vertex projectors have real Hermitian span equal to the
real diagonal algebra.  Its scalar-free part is two-dimensional, and the displayed
\(\DtriI,\DtriII\) form a basis of that plane.  The oriented one-step transport is the
unitary cyclic shift \(X\); Hermiticity of a real two-dimensional transport plane
forces the real span of \(X+X^\dagger\) and \(-i(X-X^\dagger)\).  Thus the stated
five real directions are exactly the scalar, two diagonal defect directions, and
two cyclic transport directions retained in the primitive first-order slice
generated by a primitive trimer plus one cyclic transport orbit.  Combining the two transport basis vectors in polar form with \(\rho\ge0\) and
\(\phi\) taken modulo \(2\pi\) gives
\(C(\phi)=e^{-i\phi}X+e^{i\phi}X^2\), since
\(X^\dagger=X^2\) and
\(\cos\phi(X+X^\dagger)+\sin\phi[-i(X-X^\dagger)]
=e^{-i\phi}X+e^{i\phi}X^2\).
\end{proof}



\begin{remark}[Scope of the trimer operator space]
The full Hermitian \(3\times3\) matrix space is nine-dimensional over
\(\RR\), and the \(*\)-algebra generated by the diagonal projectors and a cyclic
shift is \(M_3(\CC)\).  The subspace \(\mathfrak T_{\rm tri}\) is only the
five-dimensional real slice generated by the scalar direction, two centered
diagonal defects, and one cyclic transport plane.
\end{remark}

\begin{corollary}[Finite trimer-pair supports]
\label{p1:cor:finite-trimer-pair-supports}
An oriented Hurwitz cube-link diagonal supplies two rank-three trimer supports,
\[
    E_+:=\CC^{T_+},
    \qquad
    E_-:=\CC^{T_-}.
\]
The first trimer gives the rank-three pair carrier
\[
    \mathcal O_9^{\rm tri}
    :=\End(E_+)
    \cong E_+\otimes E_+^*,
    \qquad
    \dim_\CC\mathcal O_9^{\rm tri}=9.
\]
The displayed trimer-pair construction gives the geometric source of the
nine-dimensional carrier.  Its dimension has an arithmetic compatibility with
the Clifford octet count, but it is not a Clifford scalar completion.  This is a carrier-support statement only:
after the chosen Hurwitz vertex and oriented body diagonal are fixed, the finite
cube symmetry supplies the support and its permutation skeleton, not a continuous
Lie-group action.
\end{corollary}

\begin{proof}
The vector spaces \(\CC^{T_+}\) and \(\CC^{T_-}\) have complex dimension three; either oriented trimer gives the same rank-three support.  Choosing \(T_+\) fixes a representative.
Hence the pair carrier on \(T_+\) has dimension
\(\dim_\CC\mathcal O_9^{\rm tri}=\dim_\CC(E_+\otimes E_+^*)=3^2=9\).  After the chosen Hurwitz vertex and oriented body diagonal are fixed, the cube
supplies a finite support and its permutation skeleton, canonical only up to the finite Hurwitz/cube symmetry; it does not by itself define a
continuous Lie-group action.  Changing those choices only acts by the finite
Hurwitz/cube symmetry.
\end{proof}



\begin{corollary}[Trimer projector defects and diagonal operator space]
\label{p1:cor:opposite-trimer-defects}
Let \(P_a,P_b,P_c\in\End(\CC^T)\) be the coordinate rank-one projectors onto the three vertices of one oriented trimer
\(E_-=\CC^{T_-}\), listed in the induced cyclic order, and let
\(I=P_a+P_b+P_c\).  Use the finite primitive-idempotent trace
\(\operatorname{tr}_F\), normalized by
\(\operatorname{tr}_F(P_a)=\operatorname{tr}_F(P_b)=\operatorname{tr}_F(P_c)=1\).
The scalar-centered apex defect and the oriented opposite-edge defect are
\[
    \DtriI=P_a-\frac13I,
    \qquad
    \DtriII=P_b-P_c.
\]
Then
\[
    P_a=\frac13I+\DtriI,
    \qquad
    \operatorname{tr}_F\!\left(\left(\frac13I\right)^2\right)=\frac13,
    \qquad
    \tau_D:=\operatorname{tr}_F({\DtriI}^2)=\frac23,
\]
while the oriented edge defect satisfies
\[
    \operatorname{tr}_F({\DtriII}^2)=2,
    \qquad
    \operatorname{tr}_F(\DtriI\DtriII)=0,
    \qquad
    \operatorname{tr}_F\!\left((\DtriI+\DtriII)^2\right)=\frac83.
\]
\end{corollary}

\begin{proof}
In the ordered basis \((a,b,c)\),
\[
    \DtriI=\operatorname{diag}\left(\frac23,-\frac13,-\frac13\right),
    \qquad
    \DtriII=\operatorname{diag}(0,1,-1).
\]
The scalar and centered apex pieces are trace-orthogonal and satisfy
\[
    \operatorname{tr}_F\!\left(\left(\frac13I\right)^2\right)=\frac19+\frac19+\frac19=\frac13,
    \qquad
    \operatorname{tr}_F({\DtriI}^2)=\frac49+\frac19+\frac19=\frac23.
\]
Similarly,
\[
    \operatorname{tr}_F({\DtriII}^2)=1+1=2,
    \qquad
    \operatorname{tr}_F(\DtriI\DtriII)=-\frac13+\frac13=0,
\]
and therefore
\[
    \operatorname{tr}_F\!\left((\DtriI+\DtriII)^2\right)
    =\operatorname{tr}_F({\DtriI}^2)+\operatorname{tr}_F({\DtriII}^2)
    =\frac23+2=\frac83.
\]
\end{proof}

\begin{corollary}[Specific finite carrier weights]
\label{p1:cor:specific-finite-carrier-weights}
For the full Hurwitz cube-link carrier, its scalar-free reduced carrier, and
the rank-three pair carrier, Definition~\ref{p1:def:finite-carrier-weight} gives
\[
    \wp_8:=\wp(\RR^{N(1)})=\frac18,
    \qquad
    \wp_7:=\wp(\RR^{N(1)}/\RR\mathbf 1)=\frac17,
    \qquad
    \wp_9:=\wp(E_+\otimes E_+^{*})=\frac19.
\]
The products \(\wp_8\wp_7\) and \(\wp_8\wp_9\) are finite trace products
between the displayed finite carriers.
\end{corollary}

\begin{proof}
Apply the vector-space clause of
Definition~\ref{p1:def:finite-carrier-weight} to the displayed
finite-dimensional carriers.  Their dimensions are \(8\), \(7\), and \(9\),
respectively.
\end{proof}

\begin{definition}[Normalized octet trace and rank-one multiplicative factor]
\label{p1:def:normalized-octet-boundary-model}
Let \(O_8\) be an eight-dimensional real Hilbert space and set
\[
   \tau_8(A):=\frac1{8}\operatorname{Tr}_{O_8}(A).
\]
For an orthogonal rank-one projection \(P\in\End(O_8)\) and \(X>0\), define
the rank-one multiplicative factor
\[
   D_X(P):=I_{O_8}+(X-1)P.
\]
\end{definition}

\begin{proposition}[Octet-resolved multiplicative substep]
\label{p1:prop:octet-resolved-multiplicative-substep}
Let \(O_8\) be the full Hurwitz cube-link Walsh--Clifford octet of
Corollary~\ref{p1:cor:specific-finite-carrier-weights}, let \(P\) be an
orthogonal rank-one projection on \(O_8\), and let \(X>0\).  The normalized
logarithmic determinant readout of the rank-one factor \(D_X(P)\) is
\[
    \mathfrak s_8(X;P)
    :=
    \exp\!\left(\tau_8(\log D_X(P))\right)
    =
    X^{1/8}.
\]
This is independent of the selected rank-one line and is the multiplicative
form of the carrier weight \(\wp_8=1/8\); accordingly we may abbreviate it by
\(\mathfrak s_8(X)\).  By contrast, a scalar action
\(XI_{O_8}\) on the entire octet has normalized logarithmic readout \(X\), not
\(X^{1/8}\).

\end{proposition}

\begin{proof}
The eigenvalues of \(D_X(P)\) are \(X\) and seven copies of \(1\).
Functional calculus therefore gives \(\log D_X(P)=(\log X)P\).  Hence
\[
   \tau_8(\log D_X(P))
   =\frac{\operatorname{Tr}(P)}8\log X
   =\frac18\log X,
\]
and exponentiation gives \(X^{1/8}\).  For comparison,
\(\tau_8(\log(XI_{O_8}))=\log X\).
\end{proof}



\begin{corollary}[Two octet carrier products]
\label{p1:cor:two-octet-descendants}
For the chosen Hurwitz cube-link representative, pairing the octet weight with
the independently defined nine- and seven-dimensional carrier weights gives
\[
(\wp_8\wp_9)^{-1}=8\cdot9=72,
\qquad
(\wp_8\wp_7)^{-1}=8\cdot7=56.
\]
The first uses the rank-three pair carrier.  The second uses the scalar-free
reduced Walsh--Clifford carrier and gives the finite septet index
\[
D_\star=\wp_\star^{-1}=56.
\]
Thus \(72\) and \(56\) are reciprocal products of declared carrier weights.
The word ``product'' refers to these numerical tensor-dimension identities; no
functorial descendant object or canonical map between the carriers is asserted.
\end{corollary}

\begin{proof}
This is immediate from Definition~\ref{p1:def:finite-carrier-weight}:
\(\wp_8=1/8\), \(\wp_9=1/9\), and \(\wp_7=1/7\).  The identity
\(\wp_\star=\wp_8\wp_7\) is Corollary~\ref{p1:cor:carrier-q0}.  The distinction is
structural: \(9\) is the rank-three pair support of
Corollary~\ref{p1:cor:finite-trimer-pair-supports}, whereas \(7\) is the
scalar-free Walsh--Clifford quotient of Theorem~\ref{p1:thm:hurwitz-walsh-clifford}.
\end{proof}





\part{Exceptional quotients, McKay--ADE geometry and topology}
\section{The exceptional groups across four mathematical categories}
\label{p1b:sec:v12-literature-math}

A finite binary-polyhedral subgroup $\Gamma\subset SU(2)$, the round
space form $M_\Gamma=\Gamma\backslash S^3$, the invariant affine surface
$X_\Gamma=\mathbb C^2/\Gamma$, and the McKay root datum are distinct
objects.  The literature supplies additional exact bridges between them
\cite{Klein1884,McKay1980,Slodowy1980,wolf2011spaces,
BertschGyengeSzendroi2024Kleinian,Suter2007QuantumAffine,
MartinMerinoPoddarRoby2026SpaceForms}.  This part retains those bridges
as mathematical results even where no subsequent Born--Hopf physical
paper consumes them.  In particular the intrinsic Hopf contact integral
$4\pi^2$ remains a theorem about the original, specified $S^3$ host.

\subsection{Explicit Kleinian hypersurfaces, Milnor numbers and resolution matrices}
\label{p1b:subsec:v12-kleinian}

\begin{theorem}[Klein's exceptional hypersurfaces and the Milnor-number test]
\label{p1b:thm:v12-kleinian-normalforms}
For $\Gamma=2T,2O,2I$, the corresponding affine invariant surface
$X_\Gamma=\operatorname{Spec}\mathbb C[u,v]^\Gamma$ is, after a
weighted complex coordinate change, respectively the hypersurface
$f_\Gamma=0$ in $\mathbb C^3$ with
\[
\begin{array}{c|c|c|c|c}
\Gamma&f_\Gamma(x,y,z)&(w_x,w_y,w_z;d)&
\text{ADE type}&\mu(f_\Gamma)\\ \hline
2T&x^2+y^3+z^4&(12,8,6;24)&E_6&6\\
2O&x^2+y^3+yz^3&(18,12,8;36)&E_7&7\\
2I&x^2+y^3+z^5&(30,20,12;60)&E_8&8
\end{array}
\]
Here $d$ is the common weighted degree and $\mu$ the local Milnor
number (complex dimension of the Jacobian algebra).  Consequently
\[
\operatorname{Hilb}(\mathbb C[u,v]^\Gamma;t)
=\frac{1-t^d}{(1-t^{w_x})(1-t^{w_y})(1-t^{w_z})},
\qquad
\mu(f_\Gamma)=\prod_{w\in\{w_x,w_y,w_z\}}
\left(\frac d w-1\right).
\]
The polynomial normal forms are classical \emph{imported} invariant
geometry, not conclusions of the Hilbert series alone
\cite{Klein1884,Slodowy1980,BertschGyengeSzendroi2024Kleinian}.
\end{theorem}
\begin{proof}
For the classical hypersurface identifications use Klein's
classification as cited.  Once the weighted forms are fixed, the
principal homogeneous relation of degree $d$ gives the Hilbert series
by the graded non-zero-divisor exact sequence.  Each displayed $f$
has an isolated critical point at $0$.  Its three weighted-homogeneous
partials therefore form a regular sequence in the polynomial ring,
and the quotient by the Jacobian ideal has Hilbert series
$\prod_w(1-t^{d-w})/(1-t^w)$.  Taking $t\to1$ gives the stated
formula and the values $6,7,8$.  Thus the Milnor-number computation
is independently reproducible from the displayed weights.
\end{proof}

\begin{theorem}[Exceptional curves, negative Cartan form and links]
\label{p1b:thm:v12-resolution-link}
Let $\pi_\Gamma:Y_\Gamma\to X_\Gamma$ be the minimal resolution
of the corresponding Kleinian singularity.  Its exceptional divisor
consists of $r=6,7,8$ smooth rational $(-2)$-curves for
$\Gamma=2T,2O,2I$, respectively, with intersection matrix
\[
   \bigl(C_i\cdot C_j\bigr)_{1\leq i,j\leq r}=-C_{E_r},
\]
where $C_{E_r}$ is the finite ADE Cartan matrix.  The link of the
singularity is \emph{diffeomorphic} to the corresponding spherical
space form $S^3/\Gamma$ (orientation and left/right conventions fixed),
with fundamental group $\Gamma$.  The asserted diffeomorphism does
\emph{not} identify the hypersurface's induced metric with the round
quotient metric.
\end{theorem}
\begin{proof}[Imported standard theorem and topological identification]
The resolution-curve and intersection statements are the classical
Kleinian resolution theorem \cite{Slodowy1980,
BertschGyengeSzendroi2024Kleinian}.  For the link, the natural
radial projection $\mathbb C^2\setminus\{0\}\to S^3$ is
$\Gamma$-equivariant.  The quotient of a sufficiently small
spherical link has underlying smooth three-manifold $S^3/\Gamma$
because $\Gamma$ acts freely on the unit sphere.  The quotient
singularity is analytically isomorphic to $X_\Gamma$ by the
preceding theorem; the link is preserved up to diffeomorphism.
\end{proof}

\begin{remark}[Finite, affine and geometric ADE data are not interchangeable]
\label{p1b:rem:v12-ade-type-separation}
The \emph{affine} Dynkin graph $\widetilde E_r$ records tensoring
irreducible $\Gamma$-modules with the defining two-dimensional
module.  The \emph{finite} graph $E_r$ records the exceptional curves
on $Y_\Gamma$.  The $H_4$ vertices of the $600$-cell are neither the
$240$ roots of the $E_8$ lattice nor the $E_8$ resolution curves.
Any map between these data must be stated explicitly.
\end{remark}

\begin{theorem}[Derived McKay bridge, imported]
\label{p1b:thm:v12-derived-mckay}
For $\Gamma\subset SL(2,\mathbb C)$ finite, let
$Y_\Gamma=\Gamma\text{-Hilb}(\mathbb C^2)$, identified with the
minimal crepant resolution of $\mathbb C^2/\Gamma$.  There is an
exact equivalence of bounded derived categories
\[
  D^b\!\operatorname{Coh}(Y_\Gamma)
   \simeq D^b_\Gamma\!\operatorname{Coh}(\mathbb C^2).
\]
This is the established geometric/derived McKay correspondence,
quoted as an imported theorem \cite{BridgelandKingReid2001,
BertschGyengeSzendroi2024Kleinian}; it is not an identification of
finite-group state spaces with the particular Born--Hopf physical
Hilbert space.
\end{theorem}

\subsection{Seifert data, central fibres and the exact spin-structure count}
\label{p1b:subsec:v12-seifert}

\begin{theorem}[The exceptional Hopf--Seifert quotient atlas]
\label{p1b:thm:v12-seifert-atlas}
Use the left-invariant Hopf connection $\eta_0$ of Paper~I-A and its
\emph{right} circle fibres $q\mapsto qe^{i\theta}$.
For $\Gamma\in\{2T,2O,2I\}$ the left $\Gamma$-action commutes with
the right circle action; hence $M_\Gamma=\Gamma\backslash S^3$
has a Seifert fibration with base orbifold $S^2/\bar\Gamma$.
In the convention $e=b+\sum_i\beta_i/\alpha_i$, one may choose
Seifert symbols
\[
\begin{array}{c|c|c|c|c}
\Gamma& \text{orbifold orders}&
(b;(\alpha_i,\beta_i)_{i=1}^{3})& e&|H_1(M_\Gamma;\mathbb Z)|\\ \hline
2T&(2,3,3)&(-1;(2,1),(3,1),(3,1))&1/6&3\\
2O&(2,3,4)&(-1;(2,1),(3,1),(4,1))&1/12&2\\
2I&(2,3,5)&(-1;(2,1),(3,1),(5,1))&1/30&1
\end{array}
\]
The rational Euler quantity in this convention equals the orbifold
Euler characteristic $\chi_{\mathrm{orb}}=4/|\Gamma|$ for these three
specified left quotients.  The finite first-homology orders satisfy
\[
 |H_1(M_\Gamma;\mathbb Z)|
   =(\alpha_1\alpha_2\alpha_3)|e|.
\]
The symbols, which depend on orientations and Seifert conventions,
are not a new choice of Hopf contact normalization
\cite{wolf2011spaces,BorotEynardWeisse2017Seifert}.
\end{theorem}
\begin{proof}
Left and right quaternionic multiplication commute.  The induced
left action on the right-Hopf base is the projective rotation action
of $\bar\Gamma=\Gamma/\{\pm1\}$, hence the three triangle-orbifold
signatures.  The displayed symbols are the standard Seifert
presentations of these specified spherical space forms; substituting
$b=-1$ and $\beta_i=1$ gives the three rational Euler values.
The abelian presentation matrix of the Seifert fundamental group has
absolute determinant $\alpha_1\alpha_2\alpha_3|e|$;
its values $3,2,1$ agree independently with the known
abelianizations of $2T,2O,2I$.
\end{proof}

\begin{corollary}[The central half-fibre and contact Fubini test]
\label{p1b:cor:v12-half-fibre}
The central element $-1\in\Gamma$ translates a right Hopf fibre by
$\pi$ in the normalized period-$2\pi$ circle coordinate.  Thus the
generic descended fibre has period $\pi$ for the descended form
$\bar\eta_\Gamma$.  The curvature flux through the orbifold base is
$2\pi/|\bar\Gamma|$ under the unit-Chern normalization upstairs.
Therefore
\[
  \int_{M_\Gamma}\bar\eta_\Gamma\wedge d\bar\eta_\Gamma
   =\pi\,\frac{2\pi}{|\bar\Gamma|}
   =\frac{4\pi^2}{|\Gamma|}.
\]
This recovers Theorem~\ref{p1b:thm:v11-quotient-contact} without
modifying $\int_{S^3}\eta_0\wedge d\eta_0=4\pi^2$.
\end{corollary}
\begin{proof}
The right circle path from $q$ to $-q$ has parameter difference
$\pi$.  The noncentral quotient group acts on the base as a covering
in the orbifold sense, scaling the curvature integral by
$1/|\bar\Gamma|$.  Multiply by the generic fibre period; singular
fibres form a measure-zero locus for this integral.
\end{proof}

\begin{corollary}[Number of spin structures on exceptional spherical quotients]
\label{p1b:cor:v12-spin-count}
The oriented three-manifolds $M_{2T},M_{2O},M_{2I}$ have, respectively,
\[
 \boxed{\#\operatorname{Spin}(M_{2T})=1,\qquad
 \#\operatorname{Spin}(M_{2O})=2,\qquad
 \#\operatorname{Spin}(M_{2I})=1.}
\]
\end{corollary}
\begin{proof}
Every oriented three-manifold is spin, and its spin structures form
an affine space over $H^1(M;\mathbb Z/2)\cong
\operatorname{Hom}(H_1(M;\mathbb Z),\mathbb Z/2)$.
The respective abelianizations are $C_3,C_2,0$, producing the counts
$1,2,1$.  Therefore a quotient \emph{spinor} spectral claim involving
$2O$ must declare which of its two spin structures is used.
\end{proof}

\subsection{Complete scalar round spectra of the exceptional space forms}
\label{p1b:subsec:v12-peter-weyl}

\begin{theorem}[Peter--Weyl--Molien multiplicities]
\label{p1b:thm:v12-full-scalar-spectrum}
Let $S^3(R)$ carry its round metric, $\Gamma\in\{2T,2O,2I\}$
act by left translations, and $V=\mathbb C^2$ be the defining
$SU(2)$ representation.  Put
$a_n^\Gamma=\dim(\operatorname{Sym}^n V)^\Gamma$.
The nonnegative scalar Laplacian on $M_\Gamma=\Gamma\backslash S^3(R)$
has eigenvalues and multiplicities
\[
 \boxed{\lambda_n=\frac{n(n+2)}{R^2},\qquad
 d_n^\Gamma=(n+1)a_n^\Gamma,\quad n\ge0,}
\]
with $d_0^\Gamma=1$, and generating function
\[
 \sum_{n\ge0}d_n^\Gamma t^n
 =\left(1+t\frac{d}{dt}\right)M_\Gamma(t),\qquad
 M_\Gamma(t)=\sum_{n\ge0}a_n^\Gamma t^n.
\]
Equivalently, for $g\in\Gamma$ with defining eigenangles
$e^{\pm i\theta_g}$,
\[
 a_n^\Gamma=\frac1{|\Gamma|}\sum_{g\in\Gamma}
 \frac{\sin((n+1)\theta_g)}{\sin\theta_g},
\]
with the continuous limits $n+1$ and $(-1)^n(n+1)$ at $g=1$
and $g=-1$, respectively.  All odd $n$ have $a_n^\Gamma=0$.
\end{theorem}
\begin{proof}
By Peter--Weyl, the $n$th scalar harmonic eigenspace on $S^3(R)$
is, as an $SU(2)_L\times SU(2)_R$ module,
$\operatorname{Sym}^nV\otimes(\operatorname{Sym}^nV)^*$.
The left $\Gamma$-invariant part has dimension
$\dim(\operatorname{Sym}^nV)^\Gamma(n+1)$.
The scalar Laplacian acts by $n(n+2)/R^2$ on this band.
Taking the derivative of the Molien generating series gives the
claimed multiplicity series.  The character-average projector
$|\Gamma|^{-1}\sum_g \rho_n(g)$ proves the finite sum, and the
$SU(2)$ Weyl character formula gives the sine quotient.  Pairing
$g$ and $-g$ shows the odd-degree vanishing.
This is also the invariant-sector method used in the recent
space-form spectral analysis \cite{MartinMerinoPoddarRoby2026SpaceForms}.
\end{proof}

\begin{corollary}[Exact first nonconstant quotient scalar modes]
\label{p1b:cor:v12-quotient-scalar-gaps}
The first nonzero scalar eigenvalue and its multiplicity on the
\emph{unit} round quotients, followed by the next nonzero mode, are
\[
\begin{array}{c|c|c|c|c}
\Gamma& n_1 & (\lambda_{n_1},d_{n_1})&n_2&(\lambda_{n_2},d_{n_2})\\\hline
2T&6&(48,7)&8&(80,9)\\
2O&8&(80,9)&12&(168,13)\\
2I&12&(168,13)&20&(440,21)
\end{array}.
\]
For radius $R$, divide these eigenvalues by $R^2$.
\end{corollary}
\begin{proof}
The existing invariant-ring series reduce to
\[
 M_{2T}(t)=\frac{1+t^{12}}{(1-t^6)(1-t^8)},\quad
 M_{2O}(t)=\frac{1+t^{18}}{(1-t^8)(1-t^{12})},\quad
 M_{2I}(t)=\frac{1+t^{30}}{(1-t^{12})(1-t^{20})}.
\]
Expand the denominators as geometric series to find the first two
positive degrees and their coefficients, which in every listed case
are one.  Multiply by $n+1$ according to the preceding theorem.
\end{proof}

\begin{remark}[This is not the previous spinorial fixed-charge ladder]
\label{p1b:rem:v12-spectrum-owner-distinction}
The scalar quotient eigenvalues $48,80,168$ are eigenvalues of the
\emph{full round scalar Laplacian on three distinct quotient
manifolds}.  The values $34,62,142$ in
Section~\ref{p1b:sec:new-exceptional-spinorial-comparison}
refer instead to the first excited modes of a \emph{fixed-charge,
spinorially compensated horizontal operator} on the original $S^3$.
Neither table is a correction to $4\pi^2$, and their eigenvalue
lists must not be merged or interpreted as one ordered spectrum.
\end{remark}

\begin{theorem}[Exact scalar heat trace and local-to-global separation]
\label{p1b:thm:v12-quotient-heat-trace}
Let $K_R(t;x,y)$ be the scalar heat kernel of the positive round
Laplacian on $S^3(R)$, and $H_\Gamma(t)$ the heat trace on
$M_\Gamma$.  Then, for every $t>0$,
\[
 H_\Gamma(t)
 =\sum_{n=0}^\infty d_n^\Gamma
    e^{-t n(n+2)/R^2}
 =\frac1{|\Gamma|}\sum_{g\in\Gamma}
   \int_{S^3(R)}K_R(t;x,gx)\,d\operatorname{vol}(x).
\]
As $t\downarrow0$, the identity term gives
\[
 H_\Gamma(t)=\frac{2\pi^2R^3}{|\Gamma|(4\pi t)^{3/2}}
 \left(1+\frac{t}{R^2}+O(t^2)\right)
 +O(t^{-N}e^{-c_\Gamma/t})
\]
for some $c_\Gamma>0$ and finite $N$.  The exponential remainder
here refers to the \emph{nonidentity group-image part}; it does not
replace the further local heat coefficients in $O(t^2)$.
\end{theorem}
\begin{proof}
The covering heat kernel is summed over deck transformations, and
integration over the quotient yields the group-average image formula.
The identity term is $1/|\Gamma|$ of the round $S^3(R)$ heat trace.
Its first local coefficients are the Riemannian volume
$2\pi^2R^3$ and scalar curvature $6/R^2$, hence $a_1=1/R^2$.
For every $g\ne1$ the free left action has positive minimum
geodesic displacement on the compact sphere.  Standard off-diagonal
heat-kernel estimates bound the corresponding integral by an
exponentially small quantity.  No contact determinant is inferred
from this \emph{Riemannian} heat-trace identity.
\end{proof}

\begin{definition}[Independent scalar massive determinant family]
\label{p1b:def:v12-massive-scalar-zeta}
For $\mathfrak m^2>0$, the complete scalar quotient spectral zeta is
\[
 \zeta^{\rm round}_{\Gamma}(s;\mathfrak m^2)
 =\sum_{n=0}^\infty (n+1)a_n^\Gamma
 \left(\frac{n(n+2)}{R^2}+\mathfrak m^2\right)^{-s},
\qquad
 \det_\zeta(\Delta_\Gamma+\mathfrak m^2)
 =\exp\bigl[-\zeta^{\rm round\,\prime}_\Gamma(0;\mathfrak m^2)\bigr].
\]
At $\mathfrak m=0$ the constant mode must be omitted to form a
positive nonzero-mode determinant.  The massive determinant's
closed polylogarithmic forms and complementary theta/Bessel
representations are studied in
\cite{MartinMerinoPoddarRoby2026SpaceForms} (November 2026
preprint); they are separate from the horizontal and finite-shell
operators of Paper~III.  The displayed definition is self-contained
and makes no physical owner-selection claim.
\end{definition}

\subsection{McKay recursion and flat-bundle spectral characters}
\label{p1b:subsec:v12-mckay-spectral}

\begin{theorem}[All-isotypic McKay generating-series identity]
\label{p1b:thm:v12-mckay-resolvent}
Let $\{\rho_i\}_{i=0}^{r}$ be the irreducible complex modules
of $\Gamma$, with $\rho_0=\mathbf1$, and let
\[
 A_{ij}=\dim\operatorname{Hom}_\Gamma
  (\rho_i,V\otimes\rho_j)
\]
be its affine McKay adjacency matrix.  If
$b_n^i=\dim\operatorname{Hom}_\Gamma
(\rho_i,\operatorname{Sym}^nV)$ and
$P_i(t)=\sum_{n\ge0}b_n^i t^n$, then
\[
 \boxed{\mathbf P(t)=(I-tA+t^2 I)^{-1}\mathbf e_0.}
\]
For every irreducible $\rho_i$, the associated flat rank-$\dim\rho_i$
local-system Laplacian on $M_\Gamma$ has scalar eigenvalues
$n(n+2)/R^2$ with multiplicities $(n+1)b_n^i$.
\end{theorem}
\begin{proof}
The restricted $SU(2)$ Clebsch--Gordan decomposition
$V\otimes\operatorname{Sym}^nV\cong
\operatorname{Sym}^{n+1}V\oplus\operatorname{Sym}^{n-1}V$
gives $\mathbf b_{n+1}=A\mathbf b_n-\mathbf b_{n-1}$,
with $\mathbf b_0=\mathbf e_0$ and $\mathbf b_{-1}=0$.
Multiplication by powers of $t$ and summing gives
$(I-tA+t^2I)\mathbf P=\mathbf e_0$.
The associated flat bundle is
$(S^3(R)\times\rho_i)/\Gamma$ with its chosen equivariant
lift.  Its sections correspond to $\rho_i$-equivariant functions
upstairs.  Projecting the Peter--Weyl left factor onto the
appropriate isotypic component supplies the $(n+1)b_n^i$
multiplicity.  This is an exact representation-theoretic identity,
not a claimed derivation of a physical bundle.
\end{proof}

\begin{corollary}[Central parity and the defining flat-bundle ground level]
\label{p1b:cor:v12-flat-parity}
For each irreducible $\rho_i$, the central $-1$ acts by a scalar
$\epsilon_i\in\{+1,-1\}$.  Thus $b_n^i=0$ unless
$(-1)^n=\epsilon_i$.  The defining two-dimensional
$\Gamma$-module $V$ has $\epsilon=-1$ and
$b_0^V=0$, $b_1^V=1$, so its associated flat rank-two local-system
Laplacian has first eigenvalue $3/R^2$ with multiplicity two,
for \emph{each} of $2T,2O,2I$.
\end{corollary}
\begin{proof}
Central compatibility follows from Schur's lemma and the action
$-1\mapsto(-1)^n$ on $\operatorname{Sym}^nV$.
The defining module is irreducible for each exceptional group and
occurs once in $\operatorname{Sym}^1V=V$.  Apply the previous theorem.
\end{proof}

\begin{remark}[Three different uses of a spinorial representation]
\label{p1b:rem:v12-spin-owners}
A flat bundle whose fibre is $V$, a spinor bundle associated with a
chosen spin structure on $M_\Gamma$, and the separately defined
fixed-Hopf-charge compensated space
$(\mathcal V_{+1}\otimes V)^\Gamma$ are \emph{not} the same spectral
owner.  In particular, choosing a flat $V$ bundle does not
implicitly choose between the two spin structures on $M_{2O}$.
The all-isotypic McKay resolvent is a classification mechanism, not a
physical action functional.  See also
\cite{Suter2007QuantumAffine,BridgelandKingReid2001}.
\end{remark}

\subsection{Icosahedral spinor roots, the icosian rank-eight completion and an $E_8$ firewall}
\label{p1b:subsec:v12-icosian}

\begin{proposition}[Dense quaternionic icosians versus discrete eight-dimensional data]
\label{p1b:prop:v12-icosian-minkowski}
Let $K=\mathbb Q(\phi)$ with $\phi=(1+\sqrt5)/2$ and let
$\mathcal I$ be the icosian $\mathbb Z$-order generated by the
$120$ unit quaternions of $2I$.  Then $\mathcal I$ is a finite
rank-eight $\mathbb Z$ module, but its image in the original
Euclidean quaternion space $\mathbb H\simeq\mathbb R^4$ is
\emph{not} discrete.  Under the two real embeddings of $K$,
its coordinatewise Minkowski image is a discrete rank-eight
lattice in $\mathbb R^8$.
\end{proposition}
\begin{proof}
An icosahedral unit quaternion has real part $\phi/2$; its inverse
is also a group element.  Their sum is $\phi\,1$, so
$\mathbb Z[\phi]1\subset\mathcal I$.  Because $\phi$ is
irrational, the subgroup $\mathbb Z+\mathbb Z\phi\subset\mathbb R$
contains arbitrarily small nonzero elements (Dirichlet
approximation), so $\mathcal I$ is not discrete in $\mathbb R^4$.
All icosian coordinates belong to $\tfrac12\mathbb Z[\phi]$;
$1,i,j,k\in\mathcal I$ supply
$\mathbb Z[\phi]^4\subset\mathcal I$.
Consequently
$\mathbb Z[\phi]^4\subset\mathcal I\subset
(\tfrac12\mathbb Z[\phi])^4$,
so $\mathcal I$ is a full rank-eight $\mathbb Z$-module.
The diagonal map using the two conjugate real embeddings of $K$
takes $\mathbb Z[\phi]$ to a rank-two lattice in $\mathbb R^2$,
and hence takes $\mathcal I$ to a rank-eight lattice in
$\mathbb R^8$.
\end{proof}

\begin{remark}[Spinor-induced root sets do not identify the groups with Lie algebras]
\label{p1b:rem:v12-h4-e8-boundary}
The spinor-root constructions in the mathematical literature
associate the three-dimensional reflection systems
$A_3,B_3,H_3$ with the four-dimensional root sets
$D_4,F_4,H_4$, respectively.  Here $2I$ itself supplies the
$120$ unit points of the \emph{noncrystallographic} $H_4$
root system, whereas the crystallographic $E_8$ root system
has $240$ roots in eight dimensions.  Standard icosian and
Clifford constructions do connect these objects, but an
$E_8$ \emph{lattice} requires an explicitly chosen eight-dimensional
integral form/twist or a specified Clifford identification:
merely calling the original dense quaternionic icosian set a
four-dimensional lattice, or equating its $120$ vertices with
$240$ $E_8$ roots, is false.  The distinct constructions are
recorded in \cite{Stembridge1999H4,Dechant2016SpinorialE8,
KocaKocAlBarwani2001H4,Baez2017IcosahedronE8}.
\end{remark}


\subsection{Four-periodic group cohomology from the free three-sphere action}
\label{p1b:subsec:v12-four-periodic}

\begin{theorem}[The exact periodic integral cohomology of the binary exceptions]
\label{p1b:thm:v12-periodic-cohomology}
For $\Gamma\in\{2T,2O,2I\}$, cup-product with the orientation
Euler class of the free $\Gamma$-action on $S^3$ gives a period-four
isomorphism in Tate cohomology.  With trivial integral coefficients,
\[
\boxed{
\widehat H^k(\Gamma;\mathbb Z)\cong
\begin{cases}
\mathbb Z/|\Gamma|,&k\equiv0\pmod4,\\
0,&k\equiv1\pmod4,\\
\operatorname{Hom}(\Gamma_{\rm ab},\mathbb Q/\mathbb Z),
  &k\equiv2\pmod4,\\
0,&k\equiv3\pmod4.
\end{cases}}
\]
In particular,
\[
H^4(2T;\mathbb Z)\cong\mathbb Z/24,\quad
H^4(2O;\mathbb Z)\cong\mathbb Z/48,\quad
H^4(2I;\mathbb Z)\cong\mathbb Z/120,
\]
whereas $H^2$ is $\mathbb Z/3,\mathbb Z/2,0$, respectively.
The period-four assertion is the standard periodic-group-cohomology
theorem for free $S^3$ actions
\cite{Nicholson2024SwanCancellation}.
\end{theorem}
\begin{proof}
The free orientation-preserving $\Gamma$ action gives the
oriented spherical fibration
$S^3\to E\Gamma\times_\Gamma S^3\to B\Gamma$,
whose total space is homotopy equivalent to $M_\Gamma$.
Its degree-four Euler class induces the $4$-periodic cup-product
isomorphisms in Tate cohomology (the standard sphere-action
periodicity theorem).  Since
$\widehat H^0(\Gamma;\mathbb Z)=\mathbb Z/|\Gamma|$,
$\widehat H^{-1}(\Gamma;\mathbb Z)=0$,
$\widehat H^{-2}(\Gamma;\mathbb Z)=H_1(\Gamma;\mathbb Z)=\Gamma_{\rm ab}$,
and $\widehat H^1(\Gamma;\mathbb Z)=0$ for finite $\Gamma$,
periodicity yields the displayed table.
The $H^2$ expression is the dual of the finite abelianization;
for the cyclic orders $3,2,1$ already obtained above it gives the
quoted groups.  This is an independent, exact topological invariant
of the embedded finite groups, not a spectral determinant weight.
\end{proof}

\subsection{Operator domains and quotient-geometry scope}
\label{p1b:subsec:v12-frontiers}
The quotient, invariant-ring, ADE and cohomological statements above apply to
precisely the objects specified in each theorem.  The smooth Hopf contact
integral $4\pi^2$ remains an invariant of the marked parent $S^3$; it does
not supply a finite-group spectral measure or select a physical operator.
Spinorial and flat-bundle results require the stated lift, representation and
operator domain.  No result in this part identifies a physical source without
an independent observable map.


\clearpage

\part{Exceptional Clifford, shell and finite-carrier geometry}
\section{Hopf--Clifford exceptional carrier and finite supports}
\subsection{Hopf--Clifford carrier structure}
\label{p1:subsec:hopf-clifford-carrier}

The local/stable distinction in the exceptional Hopf geometry determines a finite Clifford carrier.
Its dimensions and the relation between the group and Hopf-circle actions follow
from the Hopf-fibered round sphere after the stated framing and sign conventions
are fixed.
In this subsection \(\Cl_{\mathrm{car}}\) denotes the common eight-dimensional vector-space carrier supplied by either sign convention, \(\Cl_{3,0}\) or \(\Cl_{0,3}\).  The algebra isomorphism is stated explicitly whenever it matters: \(\Cl_{3,0}\cong M_2(\CC)\), while the Riemannian/quaternionic product convention is \(\Cl_{0,3}\cong\HH\oplus\HH\).  The carrier dimensions \(8\) and \(7\) are signature-independent.  In odd dimension an irreducible spinor representation need not be faithful on the full real Clifford algebra: for \(\Cl_{0,3}\cong\HH\oplus\HH\), an irreducible real spinor module factors through one \(\HH\)-summand, whereas the finite-carrier construction here retains the full eight-dimensional Clifford vector-space carrier.  Thus the octet used below is a Clifford-carrier count, not the real dimension of an irreducible spinor module.  The quaternionically twisted Walsh product on the Hurwitz cube link realizes the \(\Cl_{0,3}\) carrier in Sec.~\ref{p1:sec:hurwitz-link-carrier}.

\begin{theorem}[Hopf--Clifford finite carrier after fixed conventions]
\label{p1:thm:hopf-clifford-finite-carrier}
\leavevmode\par\noindent
Let \(S^3(\Rmom)\) be equipped with its round metric, chosen Hopf fibration,
and unique spin structure.  Set
\[
    r_{\rm HC}:=\dim_{\RR}\Cl_{\mathrm{car}}=8,
    \qquad
    \Cl_{\mathrm{car}}\in\{\Cl_{3,0},\Cl_{0,3}\}.
\]
Then the following Hopf--\allowbreak Clifford carrier data are available after
the stated Hopf, framing, and Clifford-sign conventions have been fixed:
\begin{enumerate}
    \item the Lie-group identity \(S^3\cong\SU(2)\) supplies the non-abelian
    rank-one compact Lie-group carrier and its commuting left and right
    translation actions;
    \item the chosen Hopf fiber supplies an embedded principal circle
    \(U(1)_H\subset SU(2)_L\), generated infinitesimally by the dimensionless
    Reeb/principal generator \(X_H\) for \(\eta_0\); the metric-normalized Reeb
    field for \(\eta\) is \(\xi_\eta=\Rmom^{-1}X_H\);
    \item the chosen left Hopf-circle action commutes with the opposite right
    action \(SU(2)_R\).  Consequently
    \[
       \rho:U(1)_H\times SU(2)_R\longrightarrow
       \operatorname{Isom}^+(S^3),
       \qquad
       \rho(u,b)(q)=u q b^{-1},
    \]
    is a homomorphism with kernel
    \(\{(1,1),(-1,-1)\}=\Delta\ZZ_2\).  Its effective image is therefore
    \[
       \frac{U(1)_H\times SU(2)_R}{\Delta\ZZ_2}\cong U(2).
    \]
    \item the globally trivial Clifford bundle
    \[
        \Cl_{\mathrm{car}}(TS^3)\cong S^3\times\Cl_{\mathrm{car}}
    \]
    supplies an intrinsic rank \(r_{\rm HC}=8\) real carrier;
    \item the scalar/non-scalar Clifford split supplies
    \[
        \Cl_{\mathrm{car}}=\RR\mathbf 1\oplus \Cl_{\mathrm{car},\perp},
        \qquad
        \dim_\RR\Cl_{\mathrm{car},\perp}=r_{\rm HC}-1=7;
    \]
    here \(\Cl_{\mathrm{car},\perp}\) denotes the exterior-degree-positive summand under the Chevalley vector-space identification \(\Cl_{\mathrm{car}}\cong\Lambda^*\RR^3\); independently of this chosen complement, the quotient \(\Cl_{\mathrm{car}}/\RR\mathbf1\) is canonical as a seven-dimensional carrier.
\end{enumerate}
Consequently, at the level of finite carrier geometry, the Hopf--Clifford
package supplies the carrier dimensions \(8\) and \(7\).  Its intrinsic
same-factor group relation is the nested pair
\[
    U(1)_H\subset SU(2)_L,
\]
whereas its commuting left--right symmetry is represented by the covering
product \(U(1)_H\times SU(2)_R\) with the displayed central \(\ZZ_2\) kernel.
Thus a literal faithful direct product is not inferred from a group and one of
its subgroups.
\end{theorem}

\begin{proof}
The first assertion is the standard identification of the unit quaternions with
\(\SU(2)\).  The Hopf circle generator and its normalization are those of
Proposition~\ref{p1:prop:contact-normalization}.  By the fixed common-phase
convention its orbits are left \(U(1)_H\)-orbits.
Proposition~\ref{p1:prop:unit-quaternion-left-right-actions}
shows that these commute with right translations.  If \(u q b^{-1}=q\) for
every unit quaternion \(q\), then \(u=q bq^{-1}\) for every \(q\); hence
\(b\) is central and \((u,b)=(1,1)\) or \((-1,-1)\).  This proves the kernel
statement and the standard quotient identification with \(U(2)\).
Parallelizability of \(S^3\)
trivializes the tangent bundle, hence the Clifford bundle, giving
\(\Cl_{\mathrm{car}}(TS^3)\cong S^3\times\Cl_{\mathrm{car}}\).  Since both \(\Cl_{3,0}\) and \(\Cl_{0,3}\) have real dimension \(8\), and since \(\Cl_{3,0}\cong M_2(\CC)\) while \(\Cl_{0,3}\cong\HH\oplus\HH\), the carrier dimension is \(8\) in either sign convention.  Removing the scalar line leaves a \(7\)-dimensional non-scalar complement.
\end{proof}

\begin{remark}[Carrier dimensions]
The dimensions \(8\) and \(7\) are supplied by the Clifford carrier and its
scalar-free quotient.  The genuine Clifford scalar completion is
\(7+1=8\).  Separately, adjoining a new external line to the full octet gives
the dimension identity \(8+1=9\); no such extra line is canonical, and no
Clifford completion is asserted by that second equality.  The independently
constructed oriented Hurwitz-trimer endomorphism carrier below also has
dimension \(9\).
\end{remark}

\begin{definition}[Finite carrier weight]
\label{p1:def:finite-carrier-weight}
For a nonempty finite support \(S\), and separately for a nonzero
finite-dimensional vector space \(V\) over
\(\mathbb F\in\{\RR,\CC\}\), define
\[
    \wp(S):=\frac{1}{|S|},
    \qquad
    \wp(V):=\frac{1}{\dim_{\mathbb F}V}.
\]
If \(V\) is equipped with an inner product and
\(P\in\End_{\mathbb F}(V)\) is a rank-one orthogonal projection, then the
normalized trace state
\[
    \tau_V(A):=\frac{\Tr_{\mathbb F}(A)}{\dim_{\mathbb F}V}
\]
satisfies
\[
    \tau_V(P)=\frac{\Tr_{\mathbb F}(P)}{\dim_{\mathbb F}V}
    =\frac{1}{\dim_{\mathbb F}V}
    =\wp(V).
\]
No identification between a support cardinality and a vector-space dimension
is implied unless a separate construction supplies it.
\end{definition}


\begin{theorem}[Septet-to-octet resolution-data module]
\label{p1:thm:septet-octet-resolution-module}
Let
\[
    \overline{\Cl}_{\mathrm{car}}
    :=\Cl_{\mathrm{car}}/\RR\mathbf1,
    \qquad
    \dim_\RR\overline{\Cl}_{\mathrm{car}}=7.
\]
The complete linear resolution-data module from the scalar-free carrier into the
full Clifford carrier is
\[
    \mathcal D_\star
    :=\Hom_\RR\!\left(\overline{\Cl}_{\mathrm{car}},
                       \Cl_{\mathrm{car}}\right)
    \cong
    \overline{\Cl}_{\mathrm{car}}^{\!*}\otimes_\RR
    \Cl_{\mathrm{car}}.
\]
Consequently
\[
    D_\star:=\dim_\RR\mathcal D_\star=7\cdot8=56.
\]
More generally, for every subspace
\(Q\subset\overline{\Cl}_{\mathrm{car}}\) of dimension \(q\),
\[
    \mathcal D(Q):=\Hom_\RR(Q,\Cl_{\mathrm{car}})
\]
has dimension
\[
    \dim_\RR\mathcal D(Q)=8q.
\]
Thus any complete flag
\[
    0=U_0\subset U_1\subset\cdots\subset U_7
      =\overline{\Cl}_{\mathrm{car}},
    \qquad \dim U_q=q,
\]
induces the exact dimension ladder
\[
    0,8,16,24,32,40,48,56.
\]
The endpoint \(D_\star=56\) is independent of a flag.  The individual
intermediate subspaces are not selected by the Clifford carrier alone.  The
Hurwitz cube-link Walsh basis gives a concrete \(8\times7\) evaluation-table
realization in Corollary~\ref{p1:cor:walsh-evaluation-support-56}.
\end{theorem}

\begin{proof}
For finite-dimensional real vector spaces,
\(\Hom_\RR(U,V)\cong U^*\otimes_\RR V\), and therefore
\[
    \dim_\RR\Hom_\RR(U,V)
    =(\dim_\RR U)(\dim_\RR V).
\]
Insert \(U=\overline{\Cl}_{\mathrm{car}}\), \(V=\Cl_{\mathrm{car}}\),
and then \(U=Q\).  This gives respectively \(7\cdot8=56\) and \(8q\).
A complete flag exists in every seven-dimensional vector space, but no preferred
flag follows from the vector-space dimensions alone.
\end{proof}

\begin{corollary}[Complete resolution-data weight]
\label{p1:cor:carrier-q0}
With the inverse-dimension convention of
Definition~\ref{p1:def:finite-carrier-weight}, the complete resolution-data
module has carrier weight
\[
    \wp_\star:=\wp(\mathcal D_\star)=D_\star^{-1}=\frac1{56}.
\]
\end{corollary}

\begin{proof}
Theorem~\ref{p1:thm:septet-octet-resolution-module} gives
\(D_\star=\dim_\RR\mathcal D_\star=56\).  Applying
Definition~\ref{p1:def:finite-carrier-weight} to \(\mathcal D_\star\) gives
\(\wp_\star=1/\dim_\RR\mathcal D_\star=1/56\).
\end{proof}

\section{Finite shell operators and their boundary structures}
\subsection{Level-one shear and tensor operator exports}
\label{p1:sec:level-one-tensor-operators}
\begin{corollary}[Level-one object and operator types]
\label{p1:cor:level-one-shear-tt-operators}
On the level-one package of Corollary~\ref{p1:cor:H1-structure}, the
sector-preserving algebra $\mathcal A_1$ has
\[
 \mathcal A_1\cong\mathbb C\oplus M_2(\mathbb C)\oplus M_2(\mathbb C)
       \oplus M_4(\mathbb C)\oplus M_3(\mathbb C)\oplus M_3(\mathbb C),
 \qquad \dim_{\mathbb C}\mathcal A_1=43.
\]
Under the conjugation action of $\operatorname{Spin}(4)$ it decomposes as
\[
 \mathcal A_1\cong6(0,0)\oplus3(1,0)\oplus3(0,1)
                 \oplus(1,1)\oplus(2,0)\oplus(0,2).
\]
Thus it contains one shear operator type $(1,1)$ and one of each lowest
TT operator type $(2,0),(0,2)$. Neither $\mathcal H_1$ itself nor
$\operatorname{End}(\mathcal K_p)$, where
$\mathcal K_p=(\tfrac12,0)\oplus(0,\tfrac12)$ is the protected momentum
spinor, contains any of these three types. In contrast,
$\operatorname{End}(\mathcal H_1)$ has dimension $225$ and contains three
copies of $(1,1)$ and one of each TT type.
\end{corollary}
\begin{proof}
For each irreducible block, Clebsch--Gordan gives
\[
 \operatorname{End}(V_{j_L,j_R})
  \cong\bigoplus_{r=0}^{2j_L}\bigoplus_{s=0}^{2j_R}(r,s).
\]
Summing the six level-one blocks gives the displayed decomposition and
$6+9+9+9+5+5=43$. The list of object sectors contains none of the three
types. For the protected spinor, the diagonal endomorphisms give
$2(0,0)\oplus(1,0)\oplus(0,1)$ and the two cross blocks give
$2(\tfrac12,\tfrac12)$, of total dimension $16$.
In the full endomorphism algebra the vector block supplies one $(1,1)$;
the two ordered cross blocks between $(1,0)$ and $(0,1)$ supply two more.
Only the diagonal spin-one blocks can supply $(2,0)$ and $(0,2)$.
\end{proof}

For spin-one generators $J_i$, the five-dimensional quadrupole subspace is
\[
 Q_{ij}=\tfrac12\{J_i,J_j\}-\tfrac23\delta_{ij}I,
 \qquad Q_{ii}=0.
\]
It spans the spin-two summand of each $M_3$ block. Equivariance and any
nonzero quadrupole establish its rank five. The TT spectrum of
Theorem~\ref{p1:thm:tt-lichnerowicz-spectrum} has these same two types at
$L=2$. The matrix quadrupole and TT spaces have matching
\(\operatorname{Spin}(4)\) representation types.  An identification of
their operators still requires an explicit intertwiner; the level of the
finite carrier and the label of a tensor eigenspace are different indices.




\begin{proposition}[Symmetric and exterior forms of the level-one module]
\label{p1:prop:h1-symmetric-exterior-form}
With \(S=S_+\oplus S_-\),
\[
 \operatorname{Sym}^nS
 \cong
 \bigoplus_{p+q=n}
 \operatorname{Sym}^pS_+\otimes\operatorname{Sym}^qS_-,
 \qquad
 \dim_\CC\operatorname{Sym}^nS=\binom{n+3}{3}.
\]
In particular,
\[
 \operatorname{Sym}^2S
 \cong
 (1,0)\oplus(\tfrac12,\tfrac12)\oplus(0,1)
 \cong \Lambda^2_\CC\oplus\Lambda^1_\CC,
\]
and therefore
\[
 \boxed{
 \mathcal H_1
 \cong \mathbf1\oplus S\oplus\operatorname{Sym}^2S
 \cong \Lambda^0_\CC\oplus S\oplus\Lambda^1_\CC\oplus\Lambda^2_\CC,
 \qquad 1+4+10=15.}
\]
Define the nine-dimensional submodule
\[
 \mathcal E_9
 :=\mathbf1\oplus S\oplus(S_+\otimes S_-)
 \cong \Lambda^0_\CC\oplus S\oplus\Lambda^1_\CC.
\]
Then
\[
 \mathcal H_1/\mathcal E_9
 \cong\Lambda^2_\CC
 \cong\operatorname{Sym}^2S_+\oplus\operatorname{Sym}^2S_-.
\]
\end{proposition}

\begin{proof}
The symmetric-power decomposition is the standard identity for a direct sum.
Since \(\dim S=4\), its degree-\(n\) symmetric power has dimension
\(\binom{n+3}{3}\).  For \(n=2\), Clebsch--Gordan gives
\(\operatorname{Sym}^2S_+=(1,0)\),
\(S_+\otimes S_-=(\tfrac12,\tfrac12)\), and
\(\operatorname{Sym}^2S_-=(0,1)\).  Corollary~\ref{p1:cor:H1-structure}
then identifies the resulting six irreducible sectors with \(\mathcal H_1\).
\end{proof}

\begin{proposition}[Box versus simplex and the terminal quadratic antichain]
\label{p1:prop:box-simplex-antichain}
In highest-weight coordinates \((a,b)=(2j_L,2j_R)\in\NN^2\), the supports are
\[
 \operatorname{supp}\mathcal E_9
 =\{(0,0),(1,0),(0,1),(1,1)\}
 =\{0\le a\le1,\ 0\le b\le1\},
\]
while
\[
 \operatorname{supp}\mathcal H_1
 =\{(a,b)\in\NN^2:a+b\le2\}.
\]
Thus \(\mathcal E_9\) is the unit box downset and \(\mathcal H_1\) is the
degree-two simplex downset.  The terminal degree-two antichain is
\[
 A_2=\{(2,0),(1,1),(0,2)\},
\]
and its module is precisely
\[
 \Lambda^2_\CC\oplus\Lambda^1_\CC
 \cong\operatorname{Sym}^2S.
\]
If all three points of \(A_2\) are required, then the degree-two prefix, the
componentwise-divisibility downclosure, and the least Casimir sublevel
containing them all have support \(\mathcal H_1\).  If the common lower data
\((0,0),(1,0),(0,1)\) are retained and only the middle terminal point
\((1,1)\) is required, then the degree-two prefix gives \(\mathcal H_1\),
whereas componentwise divisibility and the least Casimir sublevel give
\(\mathcal E_9\).
\end{proposition}

\begin{proof}
The first two displays are direct enumeration.  The divisibility downclosure
of all three points of \(A_2\) is the six-point simplex, while the
divisibility downclosure of \((1,1)\) is the four-point box.  For the Casimir
\(c=j_L(j_L+1)+j_R(j_R+1)\), the least sublevel containing all of \(A_2\) is
\(c\le2\), whose irreducible labels are exactly the six labels of
\(\mathcal H_1\).  The least sublevel containing only \((1,1)\) is
\(c\le\tfrac32\), whose labels are exactly those of \(\mathcal E_9\).
Finally, a total-degree prefix containing any degree-two point contains all
labels with \(a+b\le2\).
\end{proof}

\begin{definition}[Total-weight shell filtration]
\label{p1:def:total-weight-shell-filtration}
For \(s\in\frac12\ZZ_{\ge0}\), define the \(\operatorname{Spin}(4)\)-invariant total-weight
shell inside \(\mathcal H_m\) by
\[
\mathcal S_s:=
\bigoplus_{j_L+j_R=s}V_{j_L}\otimes V_{j_R}.
\]
Then
\[
\mathcal H_m=
\bigoplus_{0\le s\le m}\mathcal S_s.
\]
\end{definition}

\begin{proposition}[Shell dimensions and coordinate-parity grading]
\label{p1:prop:shell-dimensions-parity}
For every integer \(r\ge0\),
\[
    \dim_\CC\mathcal S_{r/2}
    =\sum_{a+b=r}(a+1)(b+1)
    =\binom{r+3}{3}.
\]
Consequently
\[
    \dim_\CC\mathcal H_m
    =\sum_{r=0}^{2m}\binom{r+3}{3}
    =\binom{2m+4}{4}.
\]
The parity of \(r=2(j_L+j_R)\) defines a coordinate-invariant grading
\[
    \mathcal H_m
    =\mathcal H_m^{\mathrm{even}}
     \oplus\mathcal H_m^{\mathrm{odd}},
\]
with
\[
\dim_\CC\mathcal H_m^{\mathrm{even}}
=\frac{(m+1)(m+2)(2m^2+6m+3)}6,
\]
\[
\dim_\CC\mathcal H_m^{\mathrm{odd}}
=\frac{m(m+1)^2(m+2)}3.
\]
At \(m=1\),
\[
    \dim\mathcal H_1^{\mathrm{even}}=11,
    \qquad
    \mathcal H_1^{\mathrm{odd}}
    =(\tfrac12,0)\oplus(0,\tfrac12),
    \qquad
    \dim\mathcal H_1^{\mathrm{odd}}=4.
\]
Thus the protected block of
Definition~\ref{p1:def:s3one-protected-block} is exactly the full odd-parity
component of the projected level-one coordinate-coupling graph.
\end{proposition}

\begin{proof}
For \(a=2j_L\), \(b=2j_R\), and \(a+b=r\), the sector dimension is
\((a+1)(b+1)\).  The coefficient identity
\[
    \sum_{r\ge0}\left(\sum_{a+b=r}(a+1)(b+1)\right)t^r
    =\frac1{(1-t)^4}
\]
gives \(\dim\mathcal S_{r/2}=\binom{r+3}{3}\), and the hockey-stick identity
gives the formula for \(\dim\mathcal H_m\).  Separating the finite sum into
even and odd \(r\) and evaluating the two polynomial sums gives the displayed
dimensions.

Tensoring by \((\tfrac12,\tfrac12)\) replaces \((a,b)\) by
\((a\pm1,b\pm1)\), omitting negative spins.  Hence the total \(a+b\) changes
by \(-2,0\), or \(2\), so its parity is preserved before and after projection
to \(\mathcal H_m\).  At \(m=1\), the only odd shell is \(r=1\), which is
\((\tfrac12,0)\oplus(0,\tfrac12)\).
\end{proof}

\begin{corollary}[Minimal-level shell counts]
\label{p1:cor:m1-shell-counts}
At \(m=1\), the total-weight filtration is
\[
\mathcal H_1=
\mathcal S_0\oplus\mathcal S_{1/2}\oplus\mathcal S_1,
\]
with shell dimensions
\[
\dim\mathcal S_0=1,
\qquad
\dim\mathcal S_{1/2}=4,
\qquad
\dim\mathcal S_1=10.
\]
Thus the minimal \(\operatorname{Spin}(4)\)-equivariant truncation has the canonical shell
count
\[
15=1+4+10.
\]
This total-weight filtration is distinct from the real Clifford structural
count \(15=7+8\).  The former is canonical within \(\mathcal H_1\); the latter
compares separately declared Clifford carriers and supplies neither a
decomposition nor a map on \(\mathcal H_1\).
\end{corollary}

\begin{proof}
At total weight \(s=0\) only \((0,0)\) occurs, with dimension \(1\).  At
\(s=\frac12\), the sectors \((\frac12,0)\oplus(0,\frac12)\) occur, with
dimensions \(2+2=4\).  At \(s=1\), the sectors
\((\frac12,\frac12)\oplus(1,0)\oplus(0,1)\) occur, with dimensions \(4+3+3=10\).
Their sum is \(15\), as in Corollary~\ref{p1:cor:H1-structure}.  The
\(7+8\) split instead comes from the scalar-free/full real Clifford carrier, not
from a decomposition of \(\mathcal H_1\) into Hilbert subspaces.
\end{proof}

\begin{theorem}[Intrinsic Casimir scale-space of the level-one carrier]
\label{p1:thm:s3one-casimir-scale-space}
Let \(C_L\) and \(C_R\) denote the quadratic Casimirs on the two
\(SU(2)\) factors and define the positive \(\operatorname{Spin}(4)\)-equivariant operator
\[
    K_1:=C_L+C_R
\]
on \(\mathcal H_1\).  For \(t\ge0\), set
\[
    \mathcal V_t^{(1)}:=e^{-tK_1}.
\]
Then \(\mathcal V_t^{(1)}\) is a positive contraction semigroup satisfying
\[
    \mathcal V_0^{(1)}=I_{\mathcal H_1},
    \qquad
    \mathcal V_{t+s}^{(1)}
    =\mathcal V_t^{(1)}\mathcal V_s^{(1)},
\]
and
\[
    \lim_{t\to\infty}\mathcal V_t^{(1)}=P_{(0,0)}.
\]
Its trace is
\[
    \Tr_{\mathcal H_1}\mathcal V_t^{(1)}
    =1+4e^{-3t/4}+4e^{-3t/2}+6e^{-2t}.
\]

For \(\Lambda\ge0\), define the sharp Casimir projector
\[
    \Pi_\Lambda:=\mathbf1_{[0,\Lambda]}(K_1).
\]
Then
\[
    \Pi_{\Lambda_1}\Pi_{\Lambda_2}
    =\Pi_{\min\{\Lambda_1,\Lambda_2\}},
\]
so the visible spaces are nested.  At the distinct level-one Casimir thresholds
the cumulative ranks are
\[
    1,\qquad5,\qquad9,\qquad15.
\]
\end{theorem}

\begin{proof}
On an irreducible sector \(V_{j_L}\otimes V_{j_R}\), the operator \(K_1\)
acts by
\[
    \kappa(j_L,j_R)
    =j_L(j_L+1)+j_R(j_R+1).
\]
For the six sectors of \(\mathcal H_1\), the distinct eigenvalues and total
multiplicities are
\[
\begin{array}{c|c}
\kappa&\text{multiplicity}\\
\hline
0&1\\
\frac34&4\\
\frac32&4\\
2&6.
\end{array}
\]
Functional calculus gives positivity, the semigroup law, the large-\(t\) limit,
and the displayed trace.  The spectral-projector identity is immediate from
indicator-function multiplication, and summing multiplicities yields the ranks
\(1,5,9,15\).
\end{proof}

\subsection{Distinguished exceptional shell at the minimal level}
\label{p1:subsec:24cell-canonical}

\begin{proposition}[Minimal exceptional shell parallel to \(\mathcal H_1\)]
\label{p1:prop:minimal-exceptional-shell-H1}
At the minimal nontrivial truncation \(m=1\), Corollary~\ref{p1:cor:H1-structure} gives
\[
\mathcal H_1
=
V_{0,0}\oplus V_{1/2,0}\oplus V_{0,1/2}\oplus V_{1/2,1/2}\oplus V_{1,0}\oplus V_{0,1},
\]
of total dimension \(15\).  Within the ambient \(S^3\) geometry, the 24-cell vertices --- equivalently the 24 units
\(\mathcal H^\times=2T\) of the Hurwitz maximal order --- provide a distinguished exceptional shell
parallel to, but not derived from, the \(15\)-dimensional minimal representation truncation \(\mathcal H_1\). This \(24\)-point shell is not a subspace of \(\mathcal H_1\); it is a separate minimal exceptional finite subgroup configuration in the ambient \(S^3\) geometry.  It is minimal among the exceptional binary-polyhedral configurations and furnishes the standard tetrahedral core shared by the larger exceptional enrichments \(2O\) and \(2I\).
\end{proposition}

\begin{proof}
The decomposition and dimension of \(\mathcal H_1\) are Corollary~\ref{p1:cor:H1-structure}.  The identification \(\mathcal H^\times=2T\) is Theorem~\ref{p1:thm:hurwitz-units-maximal-order}, and the realization of these units as the vertices of the regular \(24\)-cell is Corollary~\ref{p1:cor:hurwitz-units-24cell-vertices}; the shell stratification is Theorem~\ref{p1:thm:hurwitz-radial-shells}.  The binary tetrahedral group \(2T\) is the smallest exceptional binary polyhedral group in the classification used in Section~\ref{p1:sec:finite}; the representation truncation \(\mathcal H_1\) and the vertex shell live in different mathematical categories, so no subspace relation is asserted.
\end{proof}

\begin{corollary}[Exceptional refinements]
\label{p1:cor:exceptional-refinements}
After choosing the common embedded tetrahedral core \(2T\), the larger exceptional configurations \(2O\) and \(2I\) refine rather than replace the 24-cell shell:
\[
2O = 2T \sqcup X_{24}, \qquad 2I = 2T \sqcup Y_{96}.
\]
Thus \(2T\) is the minimal exceptional shell, while \(2O\) and \(2I\) are two forked exceptional
refinements in the same ambient \(S^3\) geometry.  No inclusion between the two refinement branches is asserted.
\end{corollary}

\begin{proof}
Both \(2O\) and \(2I\) contain the common tetrahedral core \(2T\), with the displayed complementary counts.  Proposition~\ref{p1:prop:exceptional-nesting} proves that the two branches do not continue as a chain.
\end{proof}


\begin{definition}[Coordinate-sector tensoring and compressed action]
\label{p1:def:coordinate-sector-action}
Let \(V_{\mathrm{coord}}:=V_{1/2}\otimes V_{1/2}\).  In the finite
\(\operatorname{Spin}(4)\)-equivariant package \(\mathcal H_m\), a symbol \(M_x\) with
\(x\in V_{\mathrm{coord}}\) denotes the abstract Clebsch--Gordan tensoring map by
the \((\tfrac12,\tfrac12)\)-sector, followed by projection onto the retained
sectors of the chosen truncation.  Thus \(M_x\) is ordinary multiplication by a
coordinate function only on the diagonal scalar Peter--Weyl subspace; on the
non-diagonal homogeneous-bundle sectors it is representation-theoretic tensoring,
not pointwise scalar multiplication.  For a retained submodule \(\mathcal K\) with
orthogonal projection \(\Pi_{\mathcal K}\), the compressed Toeplitz action is
\[
    T_x:=\Pi_{\mathcal K}M_x\Pi_{\mathcal K}.
\]
Here \(M_x\) is already a retained operator.  Raw tensor insertion into
\(V_{\mathrm{coord}}\otimes\mathcal K\) will instead be denoted by
\(\iota_x\).  The comparison defect in
Definition~\ref{p1:def:compact-coordinate-closure-defect} is formed from this
raw map before retained projection.  The reduced matrix elements of a
chosen retained coordinate operator remain separately specified.
\end{definition}

\begin{proposition}[Level-\(1\) coordinate closure]
\label{p1:prop:closure-m1}
Let
\[
\mathcal K\cong (\tfrac12,0)\oplus(0,\tfrac12)
\]
denote the classical Killing-spinor module inside the finite truncation picture,
and let \(\Pi_{\mathcal K}\) be the projection onto \(\mathcal K\).  At \(m=1\),
the level-\(1\) retained coordinate action sends the protected block only through the retained chiral partner channels.  Equivalently, the compressed Toeplitz operator
\[
    T_x:=\Pi_{\mathcal K}M_x\Pi_{\mathcal K}
\]
associated with any \((\tfrac12,\tfrac12)\) coordinate-sector carrier operation \(M_x\) is the retained coordinate action on the protected block at level \(1\).
\end{proposition}

\begin{proof}
The uncompressed tensor product can send the protected spinor block into the
first exterior spinor channels \((\tfrac12,1)\) and \((1,\tfrac12)\), both of
which satisfy \(j_L+j_R=\tfrac32>1\) and are therefore absent from the retained
minimal truncation \(\mathcal H_1\).  Thus the level-\(1\) carrier retains the protected chiral channels and omits the first exterior spinor channels.
\end{proof}

\begin{proposition}[First exterior channels for \(m\ge 2\)]
\label{p1:prop:closure-failure}
For \(m\ge 2\), the first exterior spinor channels from \(\mathcal K\) under the coordinate sector are retained, namely
\[
(\tfrac12,1),\qquad (1,\tfrac12)
\]
are then present in the truncation.
\end{proposition}

\begin{proof}
This follows from the standard tensor-product decomposition
\[
(\tfrac12,\tfrac12)\otimes(0,\tfrac12)\cong(\tfrac12,0)\oplus(\tfrac12,1),
\]
and its left-right symmetric counterpart.
\end{proof}

\begin{theorem}[Smallest truncation excluding the first exterior spinor channels]
\label{p1:thm:minimal-coordinate-closed-truncation}
\leavevmode\par
Among the integer-level finite \(\operatorname{Spin}(4)\)-equivariant truncations \(\mathcal H_m\), \(m\in\ZZ_{\ge0}\), the level \(m=1\) is the unique truncation satisfying all of the following --- not merely the smallest one, since every \(m\ge2\) fails condition~(3) outright:
\begin{enumerate}
    \item it contains the protected half-integer sectors
    \[
    \left(\tfrac12,0\right),\qquad \left(0,\tfrac12\right);
    \]
    \item it contains the vector sector
    \[
    \left(\tfrac12,\tfrac12\right);
    \]
    \item the \(\mathcal H_m\)-retained \((\tfrac12,\tfrac12)\) coordinate action retains no first exterior spinor channels out of the protected spinor module
    \[
    \mathcal K\cong \left(\tfrac12,0\right)\oplus\left(0,\tfrac12\right).
    \]
\end{enumerate}
For \(m=0\), items (1) and (2) fail. For \(m\ge 2\), the first exterior spinor channels in Proposition~\ref{p1:prop:closure-failure} are retained.
\end{theorem}

\begin{proof}
At \(m=0\), only \((0,0)\) is present, so neither the protected half-integer sectors nor the vector sector occur. At \(m=1\), Corollary~\ref{p1:cor:H1-structure} shows that all sectors listed in (1) and (2) are present, and Proposition~\ref{p1:prop:closure-m1} gives the stated exclusion of the first exterior spinor channels. Proposition~\ref{p1:prop:closure-failure} lists the additional exterior spinor channels retained for \(m\ge 2\).
\end{proof}

\begin{remark}
This identifies the first integer level satisfying the displayed
retained-channel criteria.
\end{remark}

\subsection{Symmetric-coordinate boundary geometry of the selected cell}
\label{p1:subsec:symmetric-coordinate-boundary}

Let \(X:=V_{1/2}\otimes V_{1/2}\) be the four-dimensional coordinate sector.
The following construction is purely finite representation geometry; no
occupation rule is included in this mathematical definition.

\begin{definition}[Raw symmetric coordinate degree]
\label{p1:def:raw-symmetric-coordinate-degree}
For \(r\in\ZZ_{\ge0}\), define
\[
  \mathcal B_r:=\Sym^r(X).
\]
Then
\[
  \dim\mathcal B_r=\binom{r+3}{3},
\]
so the first values are \(1,4,10,20,35,56,84\).
\end{definition}

\begin{proposition}[Equal shell and symmetric-coordinate counts do not define a bridge]
\label{p1:prop:shell-boundary-count-guard}
For every \(r\ge0\),
\[
    \dim_\CC\mathcal S_{r/2}
    =\dim_\CC\mathcal B_r
    =\binom{r+3}{3}.
\]
This equality of dimensions does not give a \(\operatorname{Spin}(4)\)-equivariant
identification.  Already
\[
    \mathcal S_{1/2}
    =(\tfrac12,0)\oplus(0,\tfrac12),
    \qquad
    \mathcal B_1=(\tfrac12,\tfrac12),
\]
and
\[
    \mathcal S_1=(1,0)\oplus(\tfrac12,\tfrac12)\oplus(0,1),
\]
whereas
\[
    \mathcal B_2=\Sym^2(\tfrac12,\tfrac12)
    =(1,1)\oplus(0,0).
\]
Thus the common initial counts \(1,4,10,\ldots\) do not supply a map between
the carrier shells and the symmetric-coordinate boundary construction.
\end{proposition}

\begin{proof}
The shell dimension is
Proposition~\ref{p1:prop:shell-dimensions-parity}, while
Definition~\ref{p1:def:raw-symmetric-coordinate-degree} gives the same
binomial coefficient for \(\mathcal B_r\).  The displayed decompositions
follow directly from the definition of \(\mathcal S_s\) and from
\[
\Sym^2(A\otimes B)
\cong
\bigl(\Sym^2A\otimes\Sym^2B\bigr)
\oplus
\bigl(\Lambda^2A\otimes\Lambda^2B\bigr)
\]
for the two defining \(\SU(2)\)-modules.  The irreducible contents differ, so
no equivariant isomorphism follows from the common dimensions.
\end{proof}

\begin{proposition}[First excluded chiral coordinate boundary]
\label{p1:prop:first-excluded-chiral-boundary}
For the protected block
\[
  \mathcal K=(V_{1/2}\otimes V_0)\oplus(V_0\otimes V_{1/2}),
\]
coordinate tensoring gives
\[
\begin{aligned}
X\otimes(V_{1/2}\otimes V_0)
  &\cong(V_0\otimes V_{1/2})\oplus(V_1\otimes V_{1/2}),\\
X\otimes(V_0\otimes V_{1/2})
  &\cong(V_{1/2}\otimes V_0)\oplus(V_{1/2}\otimes V_1).
\end{aligned}
\]
The returned summands lie in \(\mathcal H_1\).  The two exterior summands have
total weight \(3/2\), lie outside \(\mathcal H_1\), and each has dimension
\(6\).
\end{proposition}

\begin{proof}
Apply the Clebsch--Gordan rule in each \(\SU(2)\) factor.  The exterior
summands have dimension \(3\cdot2=6\) and total left-right weight \(3/2\).
\end{proof}

\begin{definition}[Raw compact coordinate-closure defect]
\label{p1:def:compact-coordinate-closure-defect}
Write
\[
 S_+:=V_{1/2}\otimes V_0,\qquad
 S_-:=V_0\otimes V_{1/2},\qquad
 \mathcal K=S_+\oplus S_-,\qquad
 \mathcal B_{\mathcal K}:=X\otimes\mathcal K.
\]
Proposition~\ref{p1:prop:first-excluded-chiral-boundary} gives the orthogonal
Clebsch--Gordan decomposition
\[
 \mathcal B_{\mathcal K}
 \cong
 S_-\oplus(V_1\otimes V_{1/2})
 \oplus S_+\oplus(V_{1/2}\otimes V_1).
\]
Let \(\mathcal B_{\mathrm{ret}}\subset\mathcal B_{\mathcal K}\) be the
sum of its \(S_-\) and \(S_+\) summands, and let
\(p_{\mathrm{ret}}\) be the orthogonal projector within
\(\mathcal B_{\mathcal K}\) onto \(\mathcal B_{\mathrm{ret}}\).
Thus \(\mathcal B_{\mathrm{ret}}\cong\mathcal K\); it is not a copy of
the whole fifteen-dimensional \(\mathcal H_1\).
For \(x\in X\), define raw tensor insertion and its comparison defect by
\[
 \iota_x:\mathcal K\longrightarrow\mathcal B_{\mathcal K},
 \qquad \iota_x(k)=x\otimes k,
\]
\[
 \Delta_x^{\mathrm{raw}}
 :=(I_{\mathcal B_{\mathcal K}}-p_{\mathrm{ret}})\iota_x:
 \mathcal K\longrightarrow\mathcal E_{\mathrm{ext}},
 \qquad
 \mathcal E_{\mathrm{ext}}
 :=(V_1\otimes V_{1/2})\oplus(V_{1/2}\otimes V_1).
\]
Subsequent references to \(\Delta_x\) mean this raw comparison defect.
If a domain of \(\mathcal H_1\) is required, its extension is
\(\Delta_x^{\mathrm{raw}}\Pi_{\mathcal K}\).
The defect records the exterior component of raw coordinate tensoring.
The retained operator \(M_x\) of
Definition~\ref{p1:def:coordinate-sector-action}, and its compression
\(T_x\), are not used as the raw map in this definition.
\end{definition}

\begin{proposition}[Representability status of the raw defect]
\label{p1:prop:compact-coordinate-defect-status}
For every \(x\in X\), the image of \(\Delta_x^{\mathrm{raw}}\) lies in
the two exterior chiral summands of
Proposition~\ref{p1:prop:first-excluded-chiral-boundary}.
The defect is nonzero for suitable \(x\) and \(k\in\mathcal K\).
It measures failure of raw tensoring to remain in the returned summands;
the twelve-dimensional exterior comparison space is not an additional
retained sector of \(\mathcal H_1\).
The retained level-\(1\) action of
Proposition~\ref{p1:prop:closure-m1} still maps the protected block through
its retained chiral partner channels.
\end{proposition}

\begin{proof}
The displayed orthogonal decomposition identifies the complement of
\(\mathcal B_{\mathrm{ret}}\) with \(\mathcal E_{\mathrm{ext}}\), proving
the image statement.  Choose nonzero coordinate and \(S_+\) inputs both
of highest left weight.  Their left tensor has weight \(1\), so it lies
in the spin-\(1\) summand rather than the spin-\(0\) returned summand.
Their raw tensor therefore lies in \(V_1\otimes V_{1/2}\) and is unchanged
by the complementary projector; this proves the nonzero statement.
Finally,
\[
 (I_{\mathcal B_{\mathcal K}}-p_{\mathrm{ret}})p_{\mathrm{ret}}=0.
\]
Thus a map already projected onto the returned summands has zero exterior
defect.  Raw comparison and retained action have different types and stages,
so the former does not alter the latter's level-\(1\) closure statement.
\end{proof}

\begin{definition}[Oriented \texorpdfstring{\(1\oplus3\)}{1+3} boundary datum]
\label{p1:def:oriented-boundary-face}
Choose a complex normal-symbol line \(N_\sigma=\CC\eta_\sigma\), an oriented
three-dimensional complex space \(F_\sigma\), and an orientation-compatible
linear identification
\[
  X\cong N_\sigma\oplus F_\sigma.
\]
The multiplication map in the symmetric algebra then embeds
\[
  N_\sigma\odot\Sym^{r-1}(F_\sigma)
  \hookrightarrow\Sym^r(X).
\]
This multiplication map exists for every \(r\ge1\).  The following choice to
set the image to zero in degrees \(1\) and \(2\) is an additional
\emph{degree-three onset convention}; it is not forced by the normal line,
orientation, or dimension count.
For \(r\ge3\), define its image to be
\[
  \mathcal L_r^\sigma
  :=N_\sigma\odot\Sym^{r-1}(F_\sigma)
  \subset\mathcal B_r,
\]
and set \(\mathcal L_r^\sigma=0\) for \(r<3\).  Its rank is independent of
the chosen bases.  A later Hurwitz trimer space may be attached as a model for
\(F_\sigma\) only after an explicit isomorphism is chosen; no Hurwitz geometry
enters this definition or the rank proof below.
\end{definition}

\begin{theorem}[Boundary rank and quotient-reduced finite support]
\label{p1:thm:boundary-reduced-ranks}
For \(r\ge3\),
\[
  \dim_\CC\mathcal L_r^\sigma=\binom{r+1}{2}.
\]
Define
\[
  \mathcal N_r:=
  \begin{cases}
    \mathcal B_r,&0\le r\le2,\\[2pt]
    \mathcal B_r/\mathcal L_r^\sigma,&r\ge3.
  \end{cases}
\]
Then
\[
  \dim_\CC\mathcal N_r=
  \begin{cases}
    \binom{r+3}{3},&0\le r\le2,\\[4pt]
    \binom{r+3}{3}-\binom{r+1}{2},&r\ge3,
  \end{cases}
\]
and hence
\[
  \dim_\CC\mathcal N_0,\ldots,\dim_\CC\mathcal N_6
  =1,4,10,14,25,41,63.
\]
These are ranks of a finite quotient after the stated orientation and
normal-line choices.
\end{theorem}

\begin{proof}
Since \(\dim_\CC N_\sigma=1\) and \(\dim_\CC F_\sigma=3\), the direct-sum identification gives an injective symmetric multiplication map and
\[
  \dim_\CC\mathcal L_r^\sigma
  =\dim_\CC\Sym^{r-1}(F_\sigma)
  =\binom{r+1}{2}.
\]
Subtracting this boundary rank from
\(\dim\mathcal B_r=\binom{r+3}{3}\) gives the formula and the displayed
sequence.
\end{proof}

\begin{proposition}[Degreewise-only status of the boundary quotients]
\label{p1:prop:boundary-quotient-not-algebra}
The subspaces
\[
    \mathcal L^\sigma:=\bigoplus_{r\ge0}\mathcal L_r^\sigma
    \subset\Sym(X)
\]
do not form a homogeneous ideal.  Consequently the spaces
\(\mathcal N_r=\mathcal B_r/\mathcal L_r^\sigma\) are degreewise complex
vector-space quotients only.  They do not inherit a multiplication from
\(\Sym(X)\), and Paper~I-B defines no quotient coordinate algebra
\(\bigoplus_r\mathcal N_r\).
\end{proposition}

\noindent
The failure of ideal closure, and hence of inherited quotient
multiplication, is proved in
Appendix~\ref{p1:app:boundary-quotient-obstruction}.

\clearpage
\part{Finite level-one completion, coordinates and exceptional observables}
\label{p1:part:m1-finite-geometry}

\section{The finite datum \texorpdfstring{\(S^3_1\)}{S³}}
\label{p1:sec:s3one-finite-datum}

\subsection{The \texorpdfstring{\(m=1\)}{m=1} level}
\label{p1:subsec:s3one-level}

This final part records the finite geometry of \(S^3_1\), the level \(m=1\)
specialization of the integer-level family \(S^3_m\).  It first describes the
derived module skeleton and then separates the family of compatible finite
coordinate completions from one convenient Paper-I-A representative.  The
Hurwitz, normalizer, and icosahedral supports belong instead to the unmarked
support atlases of Definitions~\ref{p1:def:finite-support-atlases}
and~\ref{p1:def:exceptional-support-role-atlas}.  A local representative may
be fixed for a calculation, but it is not an attached selected datum and it
does not follow from the integer label \(m=1\) alone.

\begin{definition}[\texorpdfstring{\(S^3_1\)}{S3_1} skeleton and coordinate-completion family]
\label{p1:def:s3one-finite-geometry-datum}
The derived level-one skeleton is
\[
    \mathsf B_{\rm skel}
    =
    (S^3(\Rmom),\mathcal H_\bullet,\mathcal H_1),
\]
where \(\mathcal H_\bullet=\{\mathcal H_m\}_{m\ge0}\) is the integer-level
finite \(\operatorname{Spin}(4)\)-equivariant module family.  Let
\(\mathfrak M_{\rm coord}(\mathcal H_1)\) be the family of Hermitian
\(\operatorname{Spin}(4)\)-equivariant vector quadruples on \(\mathcal H_1\) that match the
normalized scalar multiplication edge.  Corollary~\ref{p1:cor:level-one-coordinate-completion-moduli}
proves that, after quotienting by sector rephasings, the four-edge forest
leaves exactly three nonnegative reduced magnitudes:
\[
 \mathfrak M_{\rm coord}(\mathcal H_1)/{\sim}
 \;\cong\;
 \mathbb R_{\ge0}^{3},
 \qquad
 [\mathsf C]\longleftrightarrow
 (r_{1,0},r_{0,1},\varrho_{\rm sp}).
\]
Here the first two entries are the two geometric off-diagonal magnitudes and
\(\varrho_{\rm sp}\) is the phase-invariant protected magnitude.  Frames,
bases, and removable edge phases are presentation conventions.

For every completion class define its protected radius
\[
 \rho_{\mathsf C}:=2\varrho_{\rm sp}(\mathsf C)\ge0.
\]
This notation is defined on the full completion cone and may vanish.  On the
Paper-I-A family below it agrees with the supplied positive scale,
\(\rho_{\mathsf C}=\rho_{\rm prot}\), but \(\rho_{\rm prot}\) remains
reserved for that special family.

For a completion \(\mathsf C\in\mathfrak M_{\rm coord}(\mathcal H_1)\), write
\[
 \mathsf B(\mathsf C):=(\mathsf B_{\rm skel},\mathsf C),
 \qquad \mathsf B:=\mathsf B(\mathsf C).
\]
The notation \(\mathsf C_{\rm PI}(\rho_{\rm prot})\) denotes the Paper-I-A
Toeplitz/Clifford family after its geometric reduced matrix elements have been
fixed and its protected scale \(\rho_{\rm prot}>0\) supplied; on that member
\(\varrho_{\rm sp}=\rho_{\rm prot}/2\).  Once a value of
\(\rho_{\rm prot}\) is fixed we may abbreviate the member by
\(\mathsf C_{\rm PI}\).  The scale-matched member
\[
 \mathsf C_{\rm cone}:=
 \mathsf C_{\rm PI}(\rho_{\rm prot}=\Rmom)
\]
is the metric-cone Clifford-symbol point defined below.  Thus
\(\mathsf C_{\rm PI}(\rho_{\rm prot})\) is a one-parameter family inside the
raw completion atlas, while \(\mathsf C_{\rm cone}\) is one distinguished
member of it; neither is a consequence of the carrier alone.

The second primitive round-symbol export is the intrinsic fixed-radius class
\(\mathsf C_D\) of
Theorem~\ref{p1:thm:intrinsic-tangential-coordinate-class}.  It is generated
from the retained Hodge/spin model and the intrinsic tangential metric
Clifford symbol, rather than from the ambient cone symbol.  In particular,
\(\mathsf C_D\) is not a member of the protected-scale family merely by
renaming \(\rho_{\rm prot}\); its distinct quadratic fingerprint fixes its
position independently in the raw atlas.

The specialization \(\mathcal H_\bullet\mapsto\mathcal H_1\) determines the
skeleton and its module-theoretic consequences, but it selects neither a point
of \(\mathfrak M_{\rm coord}\) nor a support from
\(\mathfrak S_\varepsilon\), \(\mathfrak S_{\rm grp}\), or
\(\mathfrak S_{\rm exc}\).

For this datum we set
\[
    m_{\mathsf B}:=1,
    \qquad
    \mathcal H_{\mathsf B}:=\mathcal H_1.
\]
This specialization fixes the carrier inside the defined family; the
coordinate-completion family remains as stated.
\end{definition}

\subsection{The \texorpdfstring{\(m=1\)}{m=1} carrier}

\begin{proposition}[Carrier of \texorpdfstring{\(S^3_1\)}{S3_1}]
\label{p1:prop:s3one-carrier}
Let \(\mathsf B\) be an \(S^3_1\) finite-geometry datum.  Then
\[
    m_{\mathsf B}=1,
    \qquad
    \mathcal H_{\mathsf B}=\mathcal H_1,
\]
and
\[
    \dim\mathcal H_{\mathsf B}=15.
\]
\end{proposition}

\begin{proof}
The level and carrier are fixed in Definition~\ref{p1:def:s3one-finite-geometry-datum}.  The dimension statement is Corollary~\ref{p1:cor:H1-structure} applied to \(\mathcal H_1\).
\end{proof}

\begin{corollary}[Sector content of \texorpdfstring{\(S^3_1\)}{S3_1}]
\label{p1:cor:s3one-sector-content}
The level-\(1\) carrier has the sector decomposition
\[
\mathcal H_{\mathsf B}
=
(0,0)
\oplus
(\tfrac12,0)
\oplus
(0,\tfrac12)
\oplus
(\tfrac12,\tfrac12)
\oplus
(1,0)
\oplus
(0,1),
\]
so that
\[
    \dim\mathcal H_{\mathsf B}
    =1+2+2+4+3+3=15.
\]
\end{corollary}

\begin{proof}
This is Corollary~\ref{p1:cor:H1-structure} with \(\mathcal H_{\mathsf B}=\mathcal H_1\).
\end{proof}

\begin{corollary}[Shell filtration of \texorpdfstring{\(S^3_1\)}{S3_1}]
\label{p1:cor:s3one-shell-filtration}
The level-\(1\) carrier inherits the total-weight shell decomposition
\[
    \mathcal H_{\mathsf B}
    =
    \mathcal S_0\oplus\mathcal S_{1/2}\oplus\mathcal S_1
\]
with
\[
    \dim \mathcal S_0=1,
    \qquad
    \dim \mathcal S_{1/2}=4,
    \qquad
    \dim \mathcal S_1=10.
\]
Thus
\[
    15=1+4+10.
\]
\end{corollary}

\begin{proof}
This is Corollary~\ref{p1:cor:m1-shell-counts} restricted to \(\mathcal H_{\mathsf B}=\mathcal H_1\).
\end{proof}

\section{Protected block and coordinate closure}
\label{p1:sec:protected-block-coordinate-closure}

\subsection{Protected block and retained coordinate channels}

\begin{definition}[Protected block of \texorpdfstring{\(S^3_1\)}{S3_1}]
\label{p1:def:s3one-protected-block}
On \(S^3_1\) define
\[
    \mathcal K_{\mathsf B}
    :=
    (\tfrac12,0)\oplus(0,\tfrac12)
    \subset \mathcal H_{\mathsf B}.
\]
Let \(\Pi_{\mathcal K}\) denote the orthogonal projection onto this four-dimensional block.
\end{definition}

\begin{corollary}[Protected block of \texorpdfstring{\(S^3_1\)}{S3_1}]
\label{p1:cor:s3one-protected-block}
The level \(m=1\) contains the canonical four-dimensional half-integer block
\[
    \mathcal K_{\mathsf B}
    =
    (\tfrac12,0)\oplus(0,\tfrac12).
\]
It is the retained block at level \(m_{\mathsf B}=1\) consisting of the two chiral half-integer sectors of total weight \(1/2\).
\end{corollary}

\begin{proof}
The sector list of Corollary~\ref{p1:cor:s3one-sector-content} contains the two half-integer sectors of total weight \(1/2\), namely \((\tfrac12,0)\) and \((0,\tfrac12)\). Their dimensions are \(2\) and \(2\), hence their direct sum has dimension \(4\).
\end{proof}

\begin{theorem}[Retained coordinate channels at \texorpdfstring{\(m=1\)}{m=1}]
\label{p1:thm:s3one-coordinate-closure}
Let \(x\) transform in the coordinate sector \((\tfrac12,\tfrac12)\). Tensoring by \(x\) can produce the first exterior spinor sectors
\[
    (\tfrac12,1),
    \qquad
    (1,\tfrac12),
\]
but both have total weight \(3/2\) and therefore are not retained in \(\mathcal H_{\mathsf B}=\mathcal H_1\).  Hence \(S^3_1\) has no retained coordinate channel from \(\mathcal K_{\mathsf B}\) into those exterior sectors.
\end{theorem}

\begin{proof}
This is Proposition~\ref{p1:prop:closure-m1} with \(\mathcal H_{\mathsf B}=\mathcal H_1\).  The exterior sectors listed have \(j_L+j_R=3/2>1\), so they lie outside the level-\(1\) carrier.
\end{proof}

\subsection{Finite coordinate quadruple}
Before using the symbols below, fix a completion
\(\mathsf C\in\mathfrak M_{\rm coord}(\mathcal H_1)\), together with
normalized Clebsch--Gordan intertwiners representing its reduced matrix
elements, for
tensoring by the coordinate representation
\[
V_{\rm coord}=\left(\tfrac12,\tfrac12\right).
\]
For a coordinate symbol \(x_a\in V_{\rm coord}\), the notation \(M_{x_a}\)
denotes this fully specified representation-theoretic coordinate-sector
operator.  On scalar Peter--Weyl sectors it agrees with ordinary multiplication
by the corresponding embedding coordinate.  On non-scalar sectors it is the
declared Clebsch--Gordan carrier operator, not pointwise scalar multiplication.

\begin{definition}[Finite coordinate operators on \(S^3_1\)]
\label{p1:def:s3one-toeplitz-coordinates}
Let \(x_a\), \(a=1,\ldots,4\), be a normalized round embedding-coordinate
quadruple on \(S^3(\Rmom)\), transforming as \((\tfrac12,\tfrac12)\).  On the
\(S^3_1\) cell define
\[
    \widehat X_a^{\mathsf B}:=\Pi_1 M_{x_a}\Pi_1
    \in \End(\mathcal H_{\mathsf B}),
\]
where \(\Pi_1\) is the projection onto \(\mathcal H_1\).  On the protected block define
\[
    \widehat X_{a,\mathcal K}^{\mathsf B}
    :=
    \Pi_{\mathcal K}M_{x_a}\Pi_{\mathcal K}
    \in \End(\mathcal K_{\mathsf B}).
\]
\end{definition}

\begin{remark}[What the sector data do and do not determine]
\label{p1:rem:s3one-coordinate-implementation-status}
The Clebsch--Gordan rule fixes four undirected retained channels,
\[
\begin{gathered}
 (0,0)\leftrightarrow(\tfrac12,\tfrac12),\qquad
 (\tfrac12,\tfrac12)\leftrightarrow(1,0),\\
 (\tfrac12,\tfrac12)\leftrightarrow(0,1),\qquad
 (\tfrac12,0)\leftrightarrow(0,\tfrac12),
\end{gathered}
\]
but the sector list does not fix their four reduced magnitudes.  Scalar
multiplication fixes the first; the other three are coordinates on
\(\mathfrak M_{\rm coord}(\mathcal H_1)/{\sim}\).  The particular
representative \(\mathsf C_{\rm PI}\) fixes them only after its reduced
matrix elements and protected normalization have been supplied.
On the protected block the multiplicity-one maps
have, in spinor indices and up to the declared normalization and phase, the
form
\[
 T_+(v_{\alpha\dot\beta}\otimes s_\gamma)
 =c_+\,\varepsilon_{\alpha\gamma}\,
   \widetilde s_{\dot\beta},
 \qquad
 T_-(v_{\alpha\dot\beta}\otimes\widetilde s_{\dot\gamma})
 =c_-\,\varepsilon_{\dot\beta\dot\gamma}\,s_\alpha.
\]
Each is the projection onto a multiplicity-one retained summand of an
eight-dimensional tensor-product domain, not a unitary map from that entire
domain to its two-dimensional codomain.  The self-adjoint real structure and
the completion radius \(\rho_{\mathsf C}=2\varrho_{\rm sp}\) fix the protected
common magnitude.  On \(\mathsf C_{\rm PI}\) this is the positive algebraic
scale \(\rho_{\rm prot}\); neither quantity determines the two middle channel
magnitudes.
\end{remark}

\begin{proposition}[Multiplicity-one protected coordinate coupling]
\label{p1:prop:protected-coordinate-multiplicity-one}
Let
\[
 V:=V_{\frac12,\frac12},\qquad
 S_+:=V_{\frac12,0},\qquad S_-:=V_{0,\frac12}.
\]
Then
\[
 \dim_{\CC}\operatorname{Hom}_{\operatorname{Spin}(4)}(V\otimes S_+,S_-)=1,
 \qquad
 \dim_{\CC}\operatorname{Hom}_{\operatorname{Spin}(4)}(V\otimes S_-,S_+)=1.
\]
Thus before imposing Hermiticity there are two one-dimensional complex
coordinate directions, one in each chirality direction.  For a self-adjoint
coordinate quadruple the two directions are adjoint to one another.  After a
conventional relative rephasing of the two chiral summands, the protected
Hermitian coupling is therefore parametrized by one nonnegative real
magnitude.  Any two nonzero equivariant realizations on this protected channel
are proportional; a compressed-coordinate representative and a Clifford
representative coincide on the channel only after the corresponding
normalization is fixed.
\end{proposition}

\begin{proof}
Clebsch--Gordan gives
\[
 (\tfrac12,\tfrac12)\otimes(\tfrac12,0)
 =(1,\tfrac12)\oplus(0,\tfrac12),
\]
and the retained target \((0,\tfrac12)=S_-\) occurs once.  The
left--right exchanged decomposition gives the second Hom-space.  Schur's
lemma therefore makes each equivariant channel one-dimensional.  Self-adjoint
coordinates identify the reverse reduced coefficient with the Hermitian
adjoint of the forward one, and a relative phase convention makes the common
magnitude real and nonnegative.
\end{proof}

\begin{corollary}[Level-one coordinate-completion moduli]
\label{p1:cor:level-one-coordinate-completion-moduli}
After matching the normalized scalar edge and quotienting by sector
rephasings, the compatible Hermitian equivariant coordinate actions on
\(\mathcal H_1\) satisfy
\[
 \mathfrak M_{\rm coord}(\mathcal H_1)/{\sim}
 \cong \mathbb R_{\ge0}^{3},
 \qquad
 [\mathsf C]\longleftrightarrow
 (r_{1,0},r_{0,1},\varrho_{\rm sp}).
\]
Thus the carrier and channel graph are derived, and the three reduced
magnitudes are the free parameters of the classified completion cone.
A particular point of this cone requires additional data.
\end{corollary}

\begin{proof}
The Clebsch--Gordan rule gives the four unoriented edges listed in
Remark~\ref{p1:rem:s3one-coordinate-implementation-status}, each with
multiplicity one.  Hermiticity identifies the two orientations of every
edge.  Their graph is a forest, so independent rephasings of its vertices
remove all edge phases and leave one nonnegative magnitude per edge.  Scalar
multiplication fixes the \((0,0)\leftrightarrow(\tfrac12,\tfrac12)\) edge,
leaving the displayed three magnitudes.
\end{proof}

\begin{proposition}[Protected anticommutator on the full completion cone]
\label{p1:prop:general-protected-clifford-relation}
Let \(\mathsf C\in\mathfrak M_{\rm coord}(\mathcal H_1)\) be arbitrary and
write
\[
 X_a:=\widehat X_{a,\mathcal K}^{\mathsf B},
 \qquad
 \rho_{\mathsf C}:=2\varrho_{\rm sp}\ge0.
\]
With the reduced-magnitude convention
\[
 \sum_{a=1}^4X_a^2=4\varrho_{\rm sp}^{\,2}\mathbf1_{\mathcal K}
 =\rho_{\mathsf C}^{\,2}\mathbf1_{\mathcal K},
\]
the protected coordinates satisfy on the full cone
\[
 \{X_a,X_b\}
 =2\varrho_{\rm sp}^{\,2}\delta_{ab}\mathbf1_{\mathcal K}
 =\frac{\rho_{\mathsf C}^{\,2}}2\delta_{ab}\mathbf1_{\mathcal K}.
\]
If \(\rho_{\mathsf C}=0\), Hermiticity forces \(X_a=0\) for all \(a\).
If \(\rho_{\mathsf C}>0\), the normalized operators
\[
 \Gamma_a^{\mathsf C}:=\frac{2}{\rho_{\mathsf C}}X_a
\]
obey
\(\{\Gamma_a^{\mathsf C},\Gamma_b^{\mathsf C}\}
=2\delta_{ab}\mathbf1_{\mathcal K}\).  Thus inverse-normalized gamma,
nonzero two-cluster spectrum, Hilbert--Schmidt frame, and Clifford--Fierz
statements apply precisely on the positive protected-magnitude stratum.
\end{proposition}

\begin{proof}
The protected coordinate is odd for the chiral splitting
\(\mathcal K=S_+\oplus S_-\).  Hence every anticommutator is block diagonal.
As a \(\operatorname{Spin}(4)\)-module,
\[
 \operatorname{Sym}^2(\tfrac12,\tfrac12)
 =(0,0)\oplus(1,1),
\]
whereas
\[
 \End(S_+)\oplus\End(S_-)
 =2(0,0)\oplus(1,0)\oplus(0,1).
\]
There is no \((1,1)\) target, so equivariance makes the symmetric tensor
\(\{X_a,X_b\}\) a scalar multiple of \(\delta_{ab}\) on each chirality.
Writing the off-diagonal block of \(X_a\) as \(A_a:S_+\to S_-\), cyclicity
gives
\(\Tr_{S_+}(A_a^*A_a)=\Tr_{S_-}(A_aA_a^*)\).  Since both chiralities have
dimension two, the two scalar coefficients agree.  Thus
\(\{X_a,X_b\}=c\delta_{ab}\mathbf1_{\mathcal K}\).

Summing the diagonal cases yields
\[
 4c\mathbf1_{\mathcal K}
 =2\sum_aX_a^2
 =8\varrho_{\rm sp}^{\,2}\mathbf1_{\mathcal K},
\]
so \(c=2\varrho_{\rm sp}^{\,2}=\rho_{\mathsf C}^{\,2}/2\).
At \(\rho_{\mathsf C}=0\), the positive operators \(X_a^2\) vanish, hence
each Hermitian \(X_a\) vanishes.  For positive radius division by
\(\rho_{\mathsf C}/2\) gives the normalized Clifford relation.
\end{proof}

\begin{definition}[Exterior/Clifford level-one coordinate representative]
\label{p1:def:clifford-coordinate-representative}
Let
\[
 \mathcal G_1:=\Lambda^{\le2}(\RR^4)\otimes_\RR\CC
 =\Lambda^0_\CC\oplus\Lambda^1_\CC\oplus\Lambda^2_\CC,
\]
so that Proposition~\ref{p1:prop:h1-symmetric-exterior-form} identifies
\(\mathcal G_1\oplus S\) with \(\mathcal H_1\) as a \(\operatorname{Spin}(4)\)-module.
For an orthonormal basis \(e_a\) of \(\RR^4\), put
\[
 c_a:=e_a\wedge+\iota_{e_a}
\]
on the full exterior algebra, let \(P_{\le2}\) be the degree-\(\le2\)
orthogonal projection, and let \(c_a^S\) be the ordinary Euclidean Clifford
action on \(S\).  Define
\[
 \widehat X_a^{\rm Cl}
 :=\frac{\Rmom}{2}
 \bigl(P_{\le2}c_aP_{\le2}\oplus c_a^S\bigr)
 \quad\text{on }\mathcal G_1\oplus S.
\]
Write \(\mathsf C_{\rm cone}\) for the class of this quadruple.  With the
raw exterior-edge convention fixed by
\[
 Q|_{(1,0)}=2r_{1,0}^2I,\qquad
 Q|_{(0,1)}=2r_{0,1}^2I,\qquad
 Q|_{S_+\oplus S_-}=4\varrho_{\rm sp}^2I,
\]
it has
\[
 (r_{1,0},r_{0,1},\varrho_{\rm sp})_{\rm cone}
 =\left(\frac{\Rmom}{2},\frac{\Rmom}{2},
               \frac{\Rmom}{2}\right).
\]
This is a particular \(\operatorname{Spin}(4)\)-equivariant representative of the four
retained channel types of
Remark~\ref{p1:rem:s3one-coordinate-implementation-status}.  It is not
identified with \(\widehat X_a^{\mathsf B}\) unless the completion
\(\mathsf C\) is chosen to have these reduced matrix elements.
\end{definition}

\paragraph{General terminal-degree obstruction.}
The terminal-degree formula for truncated Clifford multiplication and its
proof are recorded in
Proposition~\ref{p1:prop:terminal-degree-leakage} of
Appendix~\ref{p1:app:coordinate-clifford-obstructions}.  Its level-one
consequence remains in the main chain below.

\begin{corollary}[Level-one Clifford values and recompression defect]
\label{p1:cor:clifford-level-one-values}
For the representative of
Definition~\ref{p1:def:clifford-coordinate-representative},
\[
 \sum_{a=1}^4(\widehat X_a^{\rm Cl})^2
 =\Rmom^2 I
 \quad\text{on }\Lambda^0_\CC\oplus\Lambda^1_\CC\oplus S,
\]
while
\[
 \sum_{a=1}^4(\widehat X_a^{\rm Cl})^2
 =\frac12\Rmom^2 I
 \quad\text{on }\Lambda^2_\CC.
\]
The anticommutator
\[
 \{\widehat X_a^{\rm Cl},\widehat X_b^{\rm Cl}\}
 =\frac12\Rmom^2\delta_{ab}I
\]
holds on \(\Lambda^0_\CC\oplus\Lambda^1_\CC\oplus S\) and fails on the
terminal bivector degree.  If
\(P_{\mathcal E_9}\) denotes projection to
\(\mathcal E_9=\Lambda^0_\CC\oplus\Lambda^1_\CC\oplus S\) and
\[
 \widehat X_a^{\mathcal E}:=
 P_{\mathcal E_9}\widehat X_a^{\rm Cl}P_{\mathcal E_9},
\]
then
\[
 \sum_a(\widehat X_a^{\mathcal E})^2
 =\frac14\Rmom^2 I
 \quad\text{on }\Lambda^1_\CC.
\]
Thus removing the bivector sector causes a quadratic-radius loss of
\(3\Rmom^2/4\) on the vector sector for this recompressed representative.
\end{corollary}

\paragraph{Second-order scalarity is not a selector.}
For an equivariant vector quadruple the contraction
\(\sum_a\widehat X_a^2\) is invariant and therefore scalar on each
irreducible summand; that scalarity cannot select the level-one ceiling.
Proposition~\ref{p1:prop:quadratic-coordinate-scalarity-no-selector} and its
proof are in Appendix~\ref{p1:app:coordinate-clifford-obstructions}.

\paragraph{Global Clifford scope.}
Because \(\Cl_4(\CC)\cong M_4(\CC)\), an exact unital complex
\(\Cl_4\)-module has dimension divisible by four.  Hence the
15-dimensional carrier is not itself such a module, although its protected
four-dimensional block is.  The precise obstruction is
Proposition~\ref{p1:prop:global-clifford-divisibility-obstruction} in
Appendix~\ref{p1:app:coordinate-clifford-obstructions}.

\begin{remark}[What the Clifford representative does and does not determine]
\label{p1:rem:clifford-representative-scope}
\leavevmode\par
The retained coordinate-channel graph has two connected components: the
geometric component
\[
 \Lambda^0_\CC-\Lambda^1_\CC-\Lambda^2_\CC
\]
and the protected spinor component \(S_+-S_-\).  Consequently constraints imposed only on the
geometric component cannot determine the protected magnitude.  In the
notation of Remark~\ref{p1:rem:s3one-coordinate-implementation-status}, the
finite coordinate datum therefore retains three independent magnitudes: the
two middle geometric reduced magnitudes and the protected magnitude fixed by
\(\rho_{\rm prot}\).  The exterior/Clifford representative above is precisely
the scale-matched point \(\mathsf C_{\rm cone}\), so within that branch its
spin block has \(\rho_{\rm prot}=\Rmom\).  The broader declared family
\(\mathsf C_{\rm PI}(\rho_{\rm prot})\) retains the independent protected
scale; the round carrier alone selects neither that family member nor the cone
ancestry.
Likewise, the two middle geometric magnitudes are equal for a fixed
completion that is itself \(O(4)\)-equivariant (equivalently, carries the
appropriate factor-exchange invariance), but family-level exchange symmetry
alone does not select such a fixed point.
\end{remark}

\begin{definition}[Protected-to-parent dimensionless scale ratio]
\label{p1:def:protected-parent-scale-ratio}
Define
\[
 \eta_{\rm prot}:=\frac{\rho_{\rm prot}}{\Rmom}.
\]
Here \(\Rmom\) is the unfixed overall round-sphere homothety/scale parameter,
whereas \(\rho_{\rm prot}\) is the independent protected algebraic scale of
the family member \(\mathsf C_{\rm PI}(\rho_{\rm prot})\).  Their ratio is dimensionless and remains a
completion modulus unless an independent matching principle is supplied.  The
Clifford representative of
Definition~\ref{p1:def:clifford-coordinate-representative} realizes the
optional point \(\eta_{\rm prot}=1\); it does not select that value.
\end{definition}


\begin{remark}[Symmetric versus Clifford quadratic content]
\label{p1:rem:symmetric-vs-clifford-quadratic}
For \(V=\Lambda^1_\CC=(\tfrac12,\tfrac12)\),
\[
 V\otimes V
 =\underbrace{\mathbf1\oplus(1,1)}_{\operatorname{Sym}^2V}
 \oplus\underbrace{\Lambda^2_\CC}_{\Lambda^2V}.
\]
Ordinary commuting coordinate multiplication probes the symmetric part; its
traceless component \((1,1)\) has total highest-weight degree four and is not
retained as a state sector of \(\mathcal H_1\).  The uncompressed Clifford
product has scalar-plus-bivector quadratic content, as does the exact
protected Clifford block.  Products of the full compressed representative
\(\widehat X_a^{\rm Cl}\) in
Definition~\ref{p1:def:clifford-coordinate-representative}, however, also
contain a nonzero \((1,1)\) operator component on the terminal bivector
sector.  Their quadratic operator image is
\[
 \operatorname{span}_{\CC}
 \{\widehat X_a^{\rm Cl}\widehat X_b^{\rm Cl}:1\le a,b\le4\}
 \cong\mathbf1\oplus(1,1)\oplus\Lambda^2_\CC,
 \qquad \dim_\CC=16.
\]
This is content in \(\End(\mathcal H_1)\), not an additional retained
state sector.  Appendix~\ref{p1:app:truncated-quadratic-content} gives the
terminal witness and the image calculation.  No uniqueness claim for the
coordinate completion is implied.
\end{remark}


\begin{corollary}[Frame and basis covariance of a fixed coordinate datum]
\label{p1:cor:s3one-coordinate-uniqueness}
Once a completion \(\mathsf C\in\mathfrak M_{\rm coord}(\mathcal H_1)\) is fixed, the quadruple
\[
    (\widehat X_1^{\mathsf B},\ldots,\widehat X_4^{\mathsf B})
\]
transforms by \(SO(4)\)-rotation of the oriented embedding frame and unitary
conjugation under changes of finite Clebsch--Gordan basis.  Allowing handedness
reversal enlarges the frame action to \(O(4)\).  The sector list alone does not
determine the quadruple.
\end{corollary}

\begin{proof}
The coordinate symbols form the irreducible vector sector \((\tfrac12,\tfrac12)\).
Changing the oriented orthonormal embedding frame acts by \(SO(4)\), while changing finite
Clebsch--Gordan bases conjugates the resulting matrices unitarily; allowing
handedness reversal enlarges the frame ambiguity to \(O(4)\).  The reduced
matrix elements themselves are the coordinates of the chosen completion, as
declared in Definition~\ref{p1:def:s3one-finite-geometry-datum}.
\end{proof}

\begin{proposition}[Finite-rank coordinate compression has a radius defect]
\label{p1:prop:finite-coordinate-compression-defect}
Let \(P\) be a nonzero finite-rank orthogonal projection on
\(L^2(S^3(\Rmom))\), let \(M_{x_a}\) be multiplication by the four real
embedding coordinates, and put \(X_a=PM_{x_a}P\) on \(\operatorname{Ran}P\).
Then
\[
  \sum_{a=1}^4X_a^2
  =\Rmom^2P-
   \sum_{a=1}^4PM_{x_a}(1-P)M_{x_a}P,
\]
and the second term cannot vanish identically.  In particular, no nonzero
finite-dimensional compression of radius-\(\Rmom\) coordinate multiplication
satisfies \(\sum_aX_a^2=\Rmom^2P\).
\end{proposition}

\noindent
The positivity and nonvanishing proof is deferred to
Appendix~\ref{p1:app:coordinate-clifford-obstructions}; the exact defect
identity above is the load-bearing main-text outcome.

\begin{theorem}[Protected Cliffordization from normalized Clebsch--Gordan intertwiners]
\label{p1:thm:s3one-cliffordization}
Let
\[
    \mathcal K_{\mathsf B}
    =
    S_+\oplus S_-
    =
    (\tfrac12,0)\oplus(0,\tfrac12),
\]
and let
\[
    V_{\rm coord}=(\tfrac12,\tfrac12)
\]
be the round coordinate sector.  Fix
\(\mathsf C=\mathsf C_{\rm PI}(\rho_{\rm prot})\) with
\(\rho_{\rm prot}>0\), and set
\(\mathsf B=\mathsf B(\mathsf C)\).  Choose the multiplicity-one Clebsch--Gordan
intertwiners
\[
    V_{\rm coord}\otimes S_+\longrightarrow S_-,
    \qquad
    V_{\rm coord}\otimes S_-\longrightarrow S_+
\]
with the self-adjoint real structure and protected algebraic normalization
\[
    \sum_{a=1}^{4}
    \left(\widehat X_{a,\mathcal K}^{\mathsf B}\right)^2
    =
    \rho_{\rm prot}^2\mathbf1_{\mathcal K}.
\]
Here \(\rho_{\rm prot}>0\) is the independent algebraic scale of the
Paper-I-A representative \(\mathsf C_{\rm PI}\).  It is not inherited from compression of
radius-\(\Rmom\) multiplication; Proposition~\ref{p1:prop:finite-coordinate-compression-defect}
rules out that interpretation.  The optional equality
\(\rho_{\rm prot}=\Rmom\) is only a scale-matching convention.
This is the positive-family specialization of
Proposition~\ref{p1:prop:general-protected-clifford-relation}, with
\(\rho_{\mathsf C}=\rho_{\rm prot}\).
Then the protected-block coordinate quadruple satisfies
\[
    \left\{
    \widehat X_{a,\mathcal K}^{\mathsf B},
    \widehat X_{b,\mathcal K}^{\mathsf B}
    \right\}
    =
    \frac{\rho_{\rm prot}^2}{2}\,\delta_{ab}\,\mathbf1_{\mathcal K}.
\]
Equivalently, with
\[
    \Gamma_a:=\frac{2}{\rho_{\rm prot}}\widehat X_{a,\mathcal K}^{\mathsf B},
\]
one has the normalized Euclidean Clifford relation
\[
    \{\Gamma_a,\Gamma_b\}=2\delta_{ab}\mathbf1_{\mathcal K}.
\]
\end{theorem}

\begin{proof}
Proposition~\ref{p1:prop:general-protected-clifford-relation} applies to the
specified equivariant Hermitian completion with
\(\rho_{\mathsf C}=\rho_{\rm prot}>0\).  Dividing its anticommutator by
\(\rho_{\rm prot}^2/4\) gives the normalized Clifford relation.  The detailed
Clebsch--Gordan derivation is retained in
Appendix~\ref{p1:app:protected-cliffordization-proof}.
\end{proof}

\begin{proposition}[Choice-relative transport to the negative Clifford carrier]
\label{p1:prop:s3one-even-clifford-sign-transport}
Fix the ordered oriented protected frame used in
Theorem~\ref{p1:thm:s3one-cliffordization} and the ordered sign coordinates
used in Theorem~\ref{p1:thm:hurwitz-walsh-clifford}.  In the real protected
Clifford algebra set
\[
    E_i:=\Gamma_i\Gamma_4,
    \qquad i=1,2,3.
\]
Then
\[
    E_i^2=-\mathbf1,
    \qquad
    E_iE_j=-E_jE_i\quad(i\ne j).
\]
Consequently the assignment
\[
    E_i\longmapsto\chi_{\mathbf e_i}
\]
extends to a real-algebra isomorphism
\[
    \Cl^0_{4,0}
    \xrightarrow{\ \cong\ }
    \Cl_{0,3}
    \xrightarrow{\ \cong\ }
    \bigl(\RR^{N(1)},\star_{\sigma_-}\bigr).
\]
This identifies the protected even octet with the Hurwitz--Walsh octet only
after the frame, real structure, ordered Walsh generators, cocycle, and
orientation have been fixed.  It is not a canonical identification of either
octet with \(\mathcal H_1\), \(\End_0(\mathcal K_{\mathsf B})\), or a
Bose--Mesner sector.
\end{proposition}

\begin{proof}
For \(i\ne4\), the positive Clifford relations give
\[
 (\Gamma_i\Gamma_4)^2
 =-\Gamma_i^2\Gamma_4^2=-\mathbf1.
\]
For distinct \(i,j\le3\), moving \(\Gamma_4\) through \(\Gamma_j\) gives
\(E_iE_j=-E_jE_i\).  These are the defining relations of
\(\Cl_{0,3}\); both real algebras have dimension eight, so the induced
surjective homomorphism is an isomorphism.  The second isomorphism is
Theorem~\ref{p1:thm:hurwitz-walsh-clifford}.
\end{proof}

\begin{remark}[Exceptional action classes]
\label{p1:rem:s3one-exceptional-action-contract}
The protected Clifford system is \(\operatorname{Spin}(4)\)-covariant.  If
\(\rho_{\mathcal K}(g)\) is the protected action and \(R(g)\) the orthogonal
coordinate action, then
\[
 \rho_{\mathcal K}(g)\Gamma(v)\rho_{\mathcal K}(g)^{-1}
 =\Gamma(R(g)v).
\]
Thus every declared subgroup \(2T\), \(2O\), or \(2I\) of a \(\operatorname{Spin}(4)\)
factor preserves the Clifford relation, although it need not fix the ordered
quadruple \(\Gamma_a\) term by term.  This covariance must not be conflated
with an action preserving the chosen \(24\)-point support.  Left multiplication
by an element outside \(2T\) sends \(2T\) to a different coset.  Conjugation
by \(2O=N_{SU(2)}(2T)\) does preserve \(2T\), whereas for the chosen
\(I\cong2I\) one has \(2O\cap I=2T\), so elements of \(I\setminus2T\)
do not preserve that core by the same normalizer action.  Any exceptional
bridge must therefore specify protected \(\operatorname{Spin}(4)\) covariance, left
translation, normalizer conjugation, or coset transport as its action class.
\end{remark}

\subsection{The protected commutant line}

\begin{theorem}[Protected commutant line]
\label{p1:thm:s3one-commutant-line}
On
\[
    \mathcal K_{\mathsf B}
    =
    (\tfrac12,0)\oplus(0,\tfrac12),
\]
one has
\[
    \End_{\operatorname{Spin}(4)}^{\CC}(\mathcal K_{\mathsf B})
    =
    \CC P_L\oplus\CC P_R
    =
    \CC\mathbf 1_{\mathcal K}\oplus\CC\Gamma_{\mathcal K},
\]
where \(P_L\) and \(P_R\) project onto the two chiral summands and
\[
    \Gamma_{\mathcal K}:=P_L-P_R.
\]
Consequently the non-scalar equivariant commutant direction is the unique complex
line
\[
    \CC\Gamma_{\mathcal K}.
\]
\end{theorem}

\begin{proof}
The summands \((\tfrac12,0)\) and \((0,\tfrac12)\) are inequivalent irreducible
\(\operatorname{Spin}(4)\)-modules.  Schur's lemma gives scalar endomorphisms on each summand
and no equivariant off-diagonal maps between them.  Therefore the complex-linear equivariant
endomorphism algebra is \(\CC P_L\oplus\CC P_R\).  Rewriting the two idempotents as
\(\mathbf 1_{\mathcal K}=P_L+P_R\) and \(\Gamma_{\mathcal K}=P_L-P_R\) gives the
second decomposition.
\end{proof}

\begin{corollary}[Level-one operator-algebra dimensions]
\label{p1:cor:level-one-operator-algebra-dimensions}
At \(m=1\), the relevant complex operator spaces have dimensions
\[
\begin{array}{c|ccccc}
\text{space}
&\End(\mathcal H_1)
&\mathcal A_1
&\End_{\operatorname{Spin}(4)}(\mathcal H_1)
&\End(\mathcal K_{\mathsf B})
&\End_{\operatorname{Spin}(4)}(\mathcal K_{\mathsf B})\\
\hline
\dim_\CC&225&43&6&16&2.
\end{array}
\]
These are distinct algebras.  If \(\rho_{\mathsf C}>0\), the protected
coordinate quadruple has nonzero off-diagonal action between the two chiral
sectors of \(\mathcal K_{\mathsf B}\), and its zero extension to
\(\mathcal H_1\) is not contained in the sector-preserving algebra
\(\mathcal A_1\).  If \(\rho_{\mathsf C}=0\), all four protected
coordinates vanish and their zero extensions belong to \(\mathcal A_1\).
The displayed algebra dimensions hold on both strata.
\end{corollary}

\begin{proof}
Proposition~\ref{p1:prop:block-algebra-dimensions} at \(n=2\) gives
\[
\dim\mathcal H_1=15,\qquad
\dim\End(\mathcal H_1)=15^2=225,\qquad
\dim\mathcal A_1=43,\qquad
\dim\End_{\operatorname{Spin}(4)}(\mathcal H_1)=6.
\]
The protected block has complex dimension \(4\), so its full endomorphism
algebra has dimension \(16\), while
Theorem~\ref{p1:thm:s3one-commutant-line} gives equivariant-commutant
dimension \(2\).  Proposition~\ref{p1:prop:general-protected-clifford-relation}
gives the off-diagonal chiral action and its zero/positive split.  On the
positive stratum \((\widehat X_{a,\mathcal K}^{\mathsf B})^2=\rho_{\mathsf C}^2I/4\), so each
protected coordinate is nonzero and cannot preserve the individual chiral
sectors.  On the zero stratum each protected coordinate is zero.
\end{proof}




\section{Finite Hurwitz faces and rank-three support}
\label{p1:sec:finite-hurwitz-faces-rank-three-support}

\subsection{Finite faces}

\begin{remark}[One decomposition and one separate count identity]
\label{p1:rem:s3one-two-fifteens}
The equality
\[
    15=1+4+10
\]
is the shell decomposition of the level-\(1\) \(\operatorname{Spin}(4)\)-module package
\(\mathcal H_{\mathsf B}=\mathcal H_1\).  The equality
\[
    15=7+8
\]
is the Clifford--Hurwitz finite-count identity formed by the scalar-free
septet and the full octet.  The latter is not a vector-space decomposition of
\(\mathcal H_1\): no inclusion of either Clifford carrier in \(\mathcal H_1\)
is asserted.  It is only an equality between \(\dim\mathcal H_1\) and the sum
of two independently defined finite carrier dimensions, all included in the
enriched \(S^3_1\) datum.
Two further complex fifteen-dimensional spaces occur below:
\(\End_0(\mathcal K_{\mathsf B})\) and the Bose--Mesner sector
\(\mathcal E_{15}^{\rm BM}\subset\CC^{2T}\).  They are not identified with
each other or with \(\mathcal H_1\).


Three nine-dimensional objects must likewise be kept separate.  The carrier
\(\mathcal E_9\) of Proposition~\ref{p1:prop:h1-symmetric-exterior-form} is a
faithful reducible \(\operatorname{Spin}(4)\)-module with four simple summands.  The
irreducible \((1,1)\cong S^2_0(\RR^4)_\CC\) is not faithful and first occurs
at total highest-weight degree four.  The Bose--Mesner kernel
\(\mathcal E_9^{\rm BM}=\ker A_+\subset\CC^{2T}\) is a third object.  Under
the natural left-factor embedding \(2T\subset SU(2)_L\),
\(\dim(\mathcal E_9)^{2T}=3\), whereas the later Bose--Mesner computation
gives \(\dim(\mathcal E_9^{\rm BM})^{2T}=0\).

There is also a fourth natural complex fifteen-dimensional module that must
not be confused with \(\mathcal H_1\): for the Dirac-spinor block embedding
\(\operatorname{Spin}(4)\to SU(4)\) on \(S_+\oplus S_-\), the complexified adjoint obeys
\[
 \mathfrak{sl}_4(\CC)\!\downarrow_{\operatorname{Spin}(4)}
 \cong
 (1,0)\oplus(0,1)\oplus2(\tfrac12,\tfrac12)\oplus(0,0).
\]
It has two vector copies and no chiral spinor sector, so it is not
\(\mathcal H_1\).  The complexified conformal algebra
\(\mathfrak{so}(4,1)_\CC\cong\mathfrak{so}(5,\CC)\) appearing in
Theorem~\ref{p1:thm:quadratic-shell-conformal-killing} is ten-dimensional,
not fifteen; no \(\mathfrak{so}(6)\cong\mathfrak{sl}_4\) identification is
established by the quadratic-shell theorem.
\end{remark}

\begin{proposition}[Count comparison between the level-one carrier and the Hurwitz shell]
\label{p1:prop:scalar-shell-visibility-lift}
Let \(\mathcal H_1\) be the level-one carrier with
\(\dim\mathcal H_1=15\), and let \(2T\subset S^3\) be the
twenty-four-point Hurwitz unit set of
Theorem~\ref{p1:thm:hurwitz-radial-shells}.  These are separate objects:
\(\mathcal H_1\) is a complex \(\operatorname{Spin}(4)\)-representation carrier, whereas
\(2T\) is a finite subset and group in the ambient quaternionic sphere.
Paper~I-B defines no canonical, equivariant, or isometric map from
\(\mathcal H_1\) to \(2T\) or to \(\CC^{2T}\).
Proposition~\ref{p1:prop:h1-scalar-hurwitz-evaluation-bridge} later supplies
a natural isometric evaluation map only on the five-dimensional scalar
subcarrier \(\mathcal H_1^{\rm sc}\); it is not a map with domain all of
\(\mathcal H_1\).

The numerical comparison is
\[
\frac{|2T|}{\dim\mathcal H_1}=\frac{24}{15}=\frac{8}{5}.
\]
It is only a comparison of two counts, not an
amplitude-support factor, a trace Jacobian, or a visibility map.  Moreover,
the real linear span of the twenty-four Hurwitz units is \(\HH\cong\RR^4\);
its dimension is \(4\), not \(24\).  The dimension \(24\) belongs instead to
the function space \(\CC^{2T}\).

One exact trace-normalization use of the ratio \(24/15\) in Paper~I-B is the
trace-normalization conversion \(\tau_{15}\leftrightarrow\tau_{24}\) along
the inclusion \(\iota_{15}:\mathcal E_{15}^{\rm BM}\hookrightarrow\CC^{2T}\) established
in Theorem~\ref{p1:thm:s3one-numerical-data}; that conversion lives entirely
inside \(\CC^{2T}\) and involves neither \(\mathcal H_1\) nor any
amplitude.
\end{proposition}

\begin{proof}
The construction and dimension of \(\mathcal H_1\) are given by
Corollary~\ref{p1:cor:H1-structure}.  The set \(2T\) and its cardinality are
given by Theorem~\ref{p1:thm:hurwitz-units-maximal-order}.  Their ratio is
the displayed arithmetic identity.  The units contain
\(1,i,j,k\), so their real span is all of \(\HH\), which has real dimension
four.  No construction in either cited result supplies a map between the
representation carrier and the vertex-function space.
\end{proof}

\begin{corollary}[\(m=1\) finite carrier atlas]
\label{p1:cor:s3one-finite-carrier-atlas}
The level-one skeleton and the core-support output
\(\mathsf{Supp}_{2T}\) provide the following finite structures on their
respective typed domains.  The first five belong to the module skeleton; the
last three are orbit-invariant consequences of the unmarked Hurwitz core:
\[
\begin{array}{c|c|c}
\text{structure} & \text{finite count} & \text{source within Paper I-B}\\
\hline
\operatorname{Spin}(4)\text{-equivariant module package} & \dim\mathcal H_1=15 & \text{declared finite family}\\
\text{shell filtration} & 15=1+4+10 & \text{total weight filtration}\\
\text{protected block} & 4=2+2 & \text{half-integer sectors}\\
\text{coordinate sector} & 4 & (\tfrac12,\tfrac12)\\
\text{self-dual and anti-self-dual sectors} & 3+3 & (1,0)\oplus(0,1)\\
\text{Clifford carrier} & 8 & \text{Hopf--Clifford carrier skeleton}\\
\text{scalar-free Clifford carrier} & 7 & \text{reduced septet}\\
\text{trimer operator carrier }\mathcal O_9^{\rm tri} & 3^2=9 & \text{oriented Hurwitz trimer support}.
\end{array}
\]
\end{corollary}

\begin{proof}
Each count was proved in the finite truncation or Hurwitz--Clifford sections.
Definition~\ref{p1:def:s3one-finite-geometry-datum} and
Definition~\ref{p1:def:exceptional-support-role-atlas} preserve their typed
dependency split.  The Hurwitz/Clifford entries are not inferred from
\(m=1\); they apply when the core-support role output is present.
\end{proof}

\subsection{Finite rank-three Hurwitz support}

\begin{definition}[Level-\(1\) rank-three support]
\label{p1:def:s3one-rank-three-support}
On \(\mathsf{Supp}_{2T}\), choose for calculation any Hurwitz representative,
vertex, and oriented cube-link diagonal from
Theorem~\ref{p1:thm:hurwitz-diagonal-trimer}.  Let
\[
    T_+\sqcup T_-
\]
be the associated oriented trimer pair, and set
\[
    E_+:=\CC^{T_+},
    \qquad
    E_-:=\CC^{T_-}.
\]
On the \(S^3_1\) cell this gives a finite rank-three support template
\[
    E_+\oplus E_-
\]
and, after choosing one oriented trimer,
\[
    E_+\cong \CC[Z_3],
    \qquad
    \mathcal O_9^{\rm tri}:=\End(E_+)\cong E_+\otimes E_+^*,
    \qquad
    \dim_\CC\mathcal O_9^{\rm tri}=3^2=9.
\]
The ambient (24)-cell symmetry acts transitively on the allowed marked
templates of this declared vertex--diagonal--trimer incidence type.  Hence the
unmarked object is the full symmetry orbit of this construction; the chosen
vertex, diagonal, and ordered trimer are calculation marks and are not
additional entries of the unmarked object.
\end{definition}

On the Hurwitz vertex-function space \(\CC^{2T}\), write
\(I,A_+,A_0,A_-,A_\pi\) for the relation matrices of the inner products
\(1,1/2,0,-1/2,-1\), respectively.

\begin{definition}[Trimer zero-extension maps]
\label{p1:def:s3one-trimer-zero-extension}
Use the ordered representatives
\[
\begin{aligned}
(t_1^+,t_2^+,t_3^+)
&=(h_{(-,+,+)},h_{(+,-,+)},h_{(+,+,-)}),\\
(t_1^-,t_2^-,t_3^-)
&=(h_{(+,-,-)},h_{(-,+,-)},h_{(-,-,+)}).
\end{aligned}
\]
For the coordinate bases \(e_i^\sigma\) of \(E_\sigma=\CC^{T_\sigma}\),
define
\[
 \iota_\sigma:E_\sigma\longrightarrow\CC^{2T},
 \qquad
 \iota_\sigma e_i^\sigma=\delta_{t_i^\sigma},
 \qquad \sigma\in\{+,-\}.
\]
These are isometries for the counting inner products.  Set
\[
 \iota_W:=\iota_+\oplus\iota_-,
 \qquad
 \mathcal W_6:=\operatorname{ran}\iota_W\subset\CC^{2T}.
\]
\end{definition}

\begin{proposition}[Nearest-neighbor compression on the trimer pair]
\label{p1:prop:s3one-trimer-adjacency-compression}
In the ordered basis of
Definition~\ref{p1:def:s3one-trimer-zero-extension},
\[
 \iota_W^*A_+\iota_W
 =\begin{pmatrix}0&C\\ C&0\end{pmatrix},
 \qquad
 C=J_3-I_3
 =\begin{pmatrix}0&1&1\\1&0&1\\1&1&0\end{pmatrix}.
\]
This is the adjacency matrix of the induced six-cycle.  It is a compression,
not a restriction: every trimer vertex has two \(A_+\)-neighbors in
\(T_+\sqcup T_-\) and six outside it, so
\[
 A_+\mathcal W_6\not\subset\mathcal W_6.
\]
The two off-diagonal perfect matchings whose sum is \(C\) are the two
half-step transports of
Proposition~\ref{p1:prop:antiprism-transport-matchings}; choosing either one
separately remains extra transport data.
\end{proposition}

\begin{proof}
Two vertices in the link are \(24\)-cell neighbors exactly when their sign
vectors have Hamming distance one.  Each \(t_i^+\) has this relation with
the two \(t_j^-\) for \(j\ne i\), and with no vertex of its own trimer.  This
gives the displayed block matrix.  Since the \(24\)-cell graph has valency
eight, the remaining six neighbors lie outside the six-point support, proving
non-invariance.
\end{proof}

\begin{theorem}[Hurwitz--trimer Bose--Mesner interface]
\label{p1:thm:s3one-trimer-bose-mesner-interface}
On \(\CC^{2T}\), set
\[
 \mathcal E_9^{\rm BM}:=\ker A_+,
 \qquad
 \mathcal E_{15}^{\rm BM}:=(\mathcal E_9^{\rm BM})^\perp,
\]
and let \(P_9^{\rm BM}\), \(P_{15}^{\rm BM}\) be their orthogonal
projectors.  Explicitly,
\[
 P_9^{\rm BM}=\frac18(3I-A_0+3A_\pi),
 \qquad
 P_{15}^{\rm BM}=\frac18(5I+A_0-3A_\pi).
\]
Then
\[
 \iota_W^*P_9^{\rm BM}\iota_W
 =\begin{pmatrix}B_9&0\\0&B_9\end{pmatrix},
 \qquad
 B_9=\frac18(4I_3-J_3),
\]
and
\[
 \iota_W^*P_{15}^{\rm BM}\iota_W
 =\begin{pmatrix}B_{15}&0\\0&B_{15}\end{pmatrix},
 \qquad
 B_{15}=\frac18(4I_3+J_3).
\]
With \(Q_0:=J_3/3\) and \(Q_\perp:=I_3-Q_0\), equivalently
\[
 B_9=\frac18Q_0+\frac12Q_\perp,
 \qquad
 B_{15}=\frac78Q_0+\frac12Q_\perp.
\]
Both are positive definite.  Consequently
\[
 \rank(P_r^{\rm BM}\iota_W)=6,
 \qquad
 \mathcal W_6\cap\mathcal E_r^{\rm BM}=\{0\},
 \qquad r\in\{9,15\}.
\]
The polar-normalized maps
\[
 U_{9,\sigma}
 :=P_9^{\rm BM}\iota_\sigma
   (\sqrt8\,Q_0+\sqrt2\,Q_\perp),
\]
\[
 U_{15,\sigma}
 :=P_{15}^{\rm BM}\iota_\sigma
   (\sqrt{8/7}\,Q_0+\sqrt2\,Q_\perp)
\]
are \(C_3\)-equivariant isometries satisfying
\[
 U_{r,\sigma}^*U_{r,\sigma}=I_3,
 \qquad
 U_{r,+}^*U_{r,-}=0,
 \qquad r\in\{9,15\}.
\]
Thus each Bose--Mesner sector contains two orthogonal, choice-relative
trimer copies.  For the oriented-diagonal action
\((a,b,c,d)\mapsto(a,d,b,c)\), the exact restrictions are
\[
 \mathcal E_9^{\rm BM}\!\downarrow_{C_3}
 \cong3\,\CC[C_3],
 \qquad
 \mathcal E_{15}^{\rm BM}\!\downarrow_{C_3}
 \cong6\,\mathbf1\oplus3\,\CC[C_3].
\]
These statements are internal to \(\CC^{2T}\); they do not construct a map
from the full level-one carrier.
\end{theorem}

\begin{proof}
The Bose--Mesner multiplication table and joint spectrum in
Theorem~\ref{p1:thm:s3one-numerical-data}\textup{(iv)} show that the two
displayed formulas are complementary self-adjoint idempotents with ranges
\(\ker A_+\) and \((\ker A_+)^\perp\), respectively.
Within either trimer, distinct vertices have inner product zero; between the
two trimers, no pair is equal, orthogonal, or antipodal.  Substitution in
\[
 P_9^{\rm BM}=\frac18(3I-A_0+3A_\pi),
 \qquad
 P_{15}^{\rm BM}=\frac18(5I+A_0-3A_\pi)
\]
therefore gives the two block Gram matrices.  Their eigenvalues on the
constant and centered trimer subspaces are respectively
\[
 \operatorname{spec}(B_9)=\left\{\frac18,\frac12,\frac12\right\},
 \qquad
 \operatorname{spec}(B_{15})
 =\left\{\frac78,\frac12,\frac12\right\}.
\]
Positive definiteness gives the ranks.  If a vector in \(\mathcal W_6\) lay
in one Bose--Mesner sector, its complementary-projector quadratic form would
vanish; positivity of the complementary Gram matrix forces the vector to be
zero, proving both intersection statements.

The displayed right factors are exactly \(B_r^{-1/2}\), so the \(U_{r,\sigma}\)
are isometries.  The off-diagonal Gram blocks vanish, proving mutual
orthogonality.  All Gram operators are circulant, giving \(C_3\)-equivariance.
Finally, the coordinate permutation has six fixed Hurwitz vertices and six
three-cycles.  Its characters on \(\CC^{2T}\) are therefore \((24,6,6)\);
tracing the two projectors against the permutation gives \((9,0,0)\) and
\((15,6,6)\), which yield the stated restrictions.
\end{proof}

\begin{corollary}[Trimer operator corner and normalized trace]
\label{p1:cor:s3one-trimer-operator-corner}
For either sign, \(\End(E_\sigma)\) is a copy of the nine-dimensional
trimer operator algebra (with \(\End(E_+)=\mathcal O_9^{\rm tri}\) in the
chosen notation), and zero extension induces the injective nonunital
\(*\)-homomorphism
\[
 \jmath_\sigma:\End(E_\sigma)
 \longrightarrow\End(\CC^{2T}),
 \qquad
 \jmath_\sigma(A)=\iota_\sigma A\iota_\sigma^*.
\]
With
\[
 \tau_3(A):=\frac13\Tr_{E_\sigma}A,
 \qquad
 \tau_{24}(B):=\frac1{24}\Tr_{\CC^{2T}}B,
\]
one has the exact conversion
\[
 \tau_{24}(\jmath_\sigma(A))=\frac18\tau_3(A).
\]
Cross-trimer maps are typed separately as
\(A\mapsto\iota_+A\iota_-^*\); they are not elements of either diagonal
corner.
\end{corollary}

\begin{proof}
The identities \(\iota_\sigma^*\iota_\sigma=I_{E_\sigma}\) give
multiplicativity, adjoint preservation, and injectivity.  Moreover
\(\Tr_{\CC^{2T}}(\iota_\sigma A\iota_\sigma^*)=\Tr_{E_\sigma}A\), so division
by \(24=8\cdot3\) gives the normalized trace formula.
\end{proof}


\section{The \texorpdfstring{\(m=1\)}{m=1} finite cell}
\label{p1:sec:s3one-finite-cell}

\begin{definition}[\(m=1\) finite cell]
\label{p1:def:s3one-finite-cell}
Let \(\mathsf B=\mathsf B(\mathsf C)\) be a level-one skeleton with a fixed
coordinate completion.  For calculation choose any representative
\(\sigma=(H,v,\delta)\) of the unmarked core support
\(\mathsf{Supp}_{2T}\), where \(H\in[2T]\), \(v\) is a vertex, and \(\delta\)
is an oriented cube-link diagonal.  Define the representative finite cell
\[
\mathfrak C_{\mathsf B;\sigma}
:=
\left(
S^3(\Rmom),
\mathcal H_{\mathsf B},
\mathcal K_{\mathsf B},
(\widehat X^{\mathsf B}_1,\ldots,\widehat X^{\mathsf B}_4),
(\widehat X^{\mathsf B}_{1,\mathcal K},\ldots,\widehat X^{\mathsf B}_{4,\mathcal K}),
\CC\Gamma_{\mathcal K},
\mathcal C_8,
\mathcal C_7,
E_+\oplus E_-,
\mathsf C,
\sigma
\right),
\]
where \(\mathcal C_8\) denotes the Hopf--Clifford octet carrier constructed in
the intrinsic part, \(\mathcal C_7\) its scalar-free reduced vector-space
carrier, and \(E_+\oplus E_-\) the oriented Hurwitz trimer-pair support.  The
unmarked finite cell is the ambient-symmetry equivalence class
\([\mathfrak C_{\mathsf B;\sigma}]\).  Thus \(\sigma\) is proof notation for a
representative of the same unmarked datum.
\end{definition}

\begin{theorem}[Structure and dependency split of the \(m=1\) finite cell]
\label{p1:thm:s3one-finite-cell-structure}
For any \(S^3_1\) skeleton \(\mathsf B_{\rm skel}\) and completion
\(\mathsf C\in\mathfrak M_{\rm coord}(\mathcal H_1)\), the module-derived part
of \(\mathfrak C_{\mathsf B;\sigma}\) is determined by
\((\mathcal H_1,\mathsf C)\) up to
\begin{enumerate}[label=(\roman*)]
    \item an \(SO(4)\)-rotation of the oriented round embedding-coordinate frame;
    \item unitary changes of finite Clebsch--Gordan bases;
\end{enumerate}
On the core-support role, the Hurwitz octet, reduced septet, and trimer pair
are determined up to equivariant isomorphism by
\(\mathsf{Supp}_{2T}\).  A concrete calculation additionally fixes
\begin{enumerate}[label=(\roman*),resume]
    \item a Hurwitz vertex and oriented cube-link diagonal.
\end{enumerate}
Changing that marking within its symmetry orbit changes only the
representative \(\sigma\), not the unmarked cell. The carrier label and
finite-support class remain separate inputs to the cell construction.
\end{theorem}

\begin{proof}
Definition~\ref{p1:def:s3one-finite-geometry-datum} gives
\(\mathcal H_{\mathsf B}=\mathcal H_1\) and explicitly separates the derived
skeleton from its coordinate-completion family and support atlas.  The protected block,
retained coordinate channels, finite coordinate quadruple, protected
anticommutator, and commutant line are respectively
Corollary~\ref{p1:cor:s3one-protected-block},
Theorem~\ref{p1:thm:s3one-coordinate-closure},
Definition~\ref{p1:def:s3one-toeplitz-coordinates},
Proposition~\ref{p1:prop:general-protected-clifford-relation}, and
Theorem~\ref{p1:thm:s3one-commutant-line}; coordinate-dependent statements
use the chosen \(\mathsf C\).  The inverse-normalized Clifford system is an
additional conclusion precisely when \(\rho_{\mathsf C}>0\).
The Clifford--Hurwitz carriers and trimer support are supplied by
Corollary~\ref{p1:cor:s3one-finite-carrier-atlas} and
Definition~\ref{p1:def:s3one-rank-three-support} on any representative of
\(\mathsf{Supp}_{2T}\).  Ambient symmetry carries any such representative to
an equivalent marked cell, proving the displayed dependency split.  The
zero-extension, Bose--Mesner spectral,
and operator-corner maps of
Theorem~\ref{p1:thm:s3one-trimer-bose-mesner-interface} and
Corollary~\ref{p1:cor:s3one-trimer-operator-corner} are then derived from that
fixed support; they add no new finite choice and no map from all of
\(\mathcal H_1\).
\end{proof}

\begin{theorem}[\(m=1\) finite numerical data]
\label{p1:thm:s3one-numerical-data}
After the \(S^3_1\) cell has fixed \(\mathcal H_{\mathsf B}=\mathcal H_1\), the
following numerical data are determined either by the completion-independent
finite geometry or, where displayed, by the chosen coordinate completion and
the representative-independent orbit data of
\([\mathfrak C_{\mathsf B;\sigma}]\).  Clause~\textup{(ii)} applies exactly
when the chosen completion lies on the nondegenerate protected stratum
\(\rho_{\mathsf C}=2\varrho_{\rm sp}>0\).  At
\(\rho_{\mathsf C}=0\), Proposition~\ref{p1:prop:general-protected-clifford-relation}
gives \(X_a=0\), so no inverse-normalized gamma or nonzero frame is asserted.

\smallskip
\noindent\textup{(i) \(\operatorname{Spin}(4)\)-sector Casimirs.}
For a sector \((j_L,j_R)\), set
\[
    C_L=j_L(j_L+1),\qquad
    C_R=j_R(j_R+1),\qquad
    C_{\rm tot}=C_L+C_R.
\]
On the level-\(1\) carrier one obtains
\[
\begin{array}{c|c|c|c|c}
\text{sector} & \dim & C_L & C_R & C_{\rm tot}\\
\hline
(0,0) & 1 & 0 & 0 & 0\\
(\tfrac12,0) & 2 & \tfrac34 & 0 & \tfrac34\\
(0,\tfrac12) & 2 & 0 & \tfrac34 & \tfrac34\\
(\tfrac12,\tfrac12) & 4 & \tfrac34 & \tfrac34 & \tfrac32\\
(1,0) & 3 & 2 & 0 & 2\\
(0,1) & 3 & 0 & 2 & 2.
\end{array}
\]

\smallskip
\noindent\textup{(ii) Protected Clifford spectra and Gram matrix.}
Assume \(\rho_{\mathsf C}>0\).  Writing
\(X_a:=\widehat X^{\mathsf B}_{a,\mathcal K}\), the protected Clifford relation
implies
\[
    X_a^2=\frac{\rho_{\mathsf C}^2}{4}\mathbf1_{\mathcal K},
    \qquad
    \operatorname{spec}(X_a)=
    \left\{-\frac{\rho_{\mathsf C}}{2},+\frac{\rho_{\mathsf C}}{2}\right\}
\]
with multiplicities \(2+2\).  Moreover
\[
    \operatorname{Tr}_{\mathcal K}(X_a)=0,
    \qquad
    \operatorname{Tr}_{\mathcal K}(X_aX_b)=\rho_{\mathsf C}^2\delta_{ab}.
\]
Thus the four protected coordinate matrices form an orthogonal Hilbert--Schmidt
frame; for \(a\ne b\), their Hilbert--Schmidt angle is \(\pi/2\).  Equivalently,
with
\[
    \Gamma_a:=\frac{2}{\rho_{\mathsf C}}X_a,
\]
one has the normalized Clifford system
\[
    \{\Gamma_a,\Gamma_b\}=2\delta_{ab}\mathbf1_{\mathcal K},
    \qquad
    \sum_{a=1}^4\Gamma_a^2=4\mathbf1_{\mathcal K}.
\]
The chirality element
\[
    \Gamma_5:=\Gamma_1\Gamma_2\Gamma_3\Gamma_4
\]
satisfies
\[
    \Gamma_5^2=\mathbf1_{\mathcal K},
    \qquad
    \operatorname{Tr}_{\mathcal K}(\Gamma_5)=0,
    \qquad
    \operatorname{rank}\frac{1\pm\Gamma_5}{2}=2.
\]
Thus the protected Clifford representation encodes the grade-count package
\[
    \dim\Lambda^\bullet\RR^4=(1,4,6,4,1),
    \qquad
    \dim\Cl^0_4=8=1+7,
    \qquad
    \dim\Lambda^2\RR^4=6=3+3.
\]
This octet is the even \(4\)-dimensional Clifford shadow of the protected gamma system.  It has the same \(8=1+7\) carrier count as the Hopf--Clifford octet used earlier, but the two are not identified as the same labeled object without the chosen Clifford convention.
The six Clifford bivector products \(X_aX_b\), \(a<b\), form an orthogonal Hilbert--Schmidt basis of the bivector space, each with norm \(\rho_{\mathsf C}^2/2\). After choosing the orientation, the Hodge combinations of these bivectors split \(\Lambda^2\RR^4\) into self-dual and anti-self-dual triples \(3+3\).  In normalized
form, the skew-Hermitian products \(\Gamma_a\Gamma_b\) have eigenvalues
\(\pm i\) with multiplicities \(2+2\), or equivalently
\(i\Gamma_a\Gamma_b\) has spectrum \(\{ -1,+1\}\) with multiplicities \(2+2\).

\smallskip
\noindent\textup{(iii) Hurwitz radial shell angle spectrum.}
For the Hurwitz support, normalized first on the unit sphere and then rescaled by
\(\Rmom\) if embedded in \(S^3(\Rmom)\), the base vertex \(1\) has shell data
\[
\begin{array}{c|c|c|c|c}
\text{shell} & \# & \operatorname{Re}(q) & |1-q| & d_{S^3}(1,q)\\
\hline
\{1\} & 1 & 1 & 0 & 0\\
\mathcal N_8^+ & 8 & \tfrac12 & 1 & \tfrac\pi3\\
\mathcal M_6 & 6 & 0 & \sqrt2 & \tfrac\pi2\\
\mathcal N_8^- & 8 & -\tfrac12 & \sqrt3 & \tfrac{2\pi}{3}\\
\{-1\} & 1 & -1 & 2 & \pi.
\end{array}
\]
On the radius-\(\Rmom\) sphere, chordal distances are multiplied by \(\Rmom\) and
geodesic distances by \(\Rmom\), while the displayed angular values are
unchanged.

\smallskip
\noindent\textup{(iv) Full \(24\)-cell incidence, distance, and metric package.}
The Hurwitz support is the regular \(24\)-cell with
\[
    (f_0,f_1,f_2,f_3)=(24,96,96,24),
    \qquad
    \operatorname{Lk}(v)\cong\text{cube},
    \quad
    (8,12,6).
\]
Equivalently, each vertex is incident to \(8\) edges, \(12\) triangular faces,
and \(6\) octahedral cells.  Each edge is incident to \(3\) triangular faces
and \(3\) octahedral cells, and each triangular face is incident to \(2\)
octahedral cells.  The complete unordered two-point distance distribution is
\[
\begin{array}{c|c|c|c}
|p-q| & \langle p,q\rangle & d_{S^3}(p,q) & \#\{\{p,q\}\}\\
\hline
1 & \tfrac12 & \tfrac\pi3 & 96\\
\sqrt2 & 0 & \tfrac\pi2 & 72\\
\sqrt3 & -\tfrac12 & \tfrac{2\pi}{3} & 96\\
2 & -1 & \pi & 12.
\end{array}
\]
The relation matrices of this five-relation antipodal scheme are also fixed.
Let \(I,A_+,A_0,A_-,J\) denote the matrices for inner products
\(1,1/2,0,-1/2,-1\), respectively.  Their valencies are
\[
    1,\quad 8,\quad 6,\quad 8,\quad 1,
\]
and their spectra, with multiplicities shown as superscripts, are
\[
\begin{array}{c|c}
\text{matrix} & \text{spectrum}\\
\hline
A_+ & 8^{(1)},\;4^{(4)},\;0^{(9)},\;(-2)^{(8)},\;(-4)^{(2)}\\
A_0 & 6^{(3)},\;0^{(12)},\;(-2)^{(9)}\\
A_- & 8^{(1)},\;2^{(8)},\;0^{(9)},\;(-4)^{(6)}\\
J & 1^{(12)},\;(-1)^{(12)}.
\end{array}
\]
The corresponding multiplication rules are
\[
\begin{aligned}
A_+^2&=8I+3A_++4A_0+A_-,\\
A_+A_0&=3A_++3A_-,\\
A_+A_-&=A_++4A_0+3A_-+8J,\\
A_+J&=A_-,\\
A_0^2&=6I+4A_0+6J,\\
A_0A_-&=3A_++3A_-,\\
A_0J&=A_0,\\
A_-^2&=8I+3A_++4A_0+A_-,\\
A_-J&=A_+,\\
J^2&=I.
\end{aligned}
\]
In particular, two vertices with inner product \(1/2,0,-1/2,-1\) have,
respectively, \(3,4,1,0\) common edge-neighbors in the \(24\)-cell graph.
The graph therefore has diameter \(3\), with graph-distance layers from any
vertex
\[
    1,
    \quad 8,
    \quad 14,
    \quad 1.
\]
The Hurwitz vertices also form a tight spherical frame in \(\RR^4\).  Writing
vertices as unit vectors \(q\in 2T\subset\RR^4\),
\[
    \sum_{q\in2T}q=0,
    \qquad
    \sum_{q\in2T} qq^{\mathsf T}=6I_4,
\]
and the full \(24\times24\) Gram matrix \(G_{pq}=\langle p,q\rangle\) has
spectrum
\[
    6^{(4)},\qquad 0^{(20)}.
\]
Moreover the fourth moment is isotropic:
\[
    \sum_{q\in2T} q_iq_jq_kq_\ell
    =\delta_{ij}\delta_{k\ell}
     +\delta_{ik}\delta_{j\ell}
     +\delta_{i\ell}\delta_{jk},
\]
equivalently
\[
    \sum_{q\in2T}\langle u,q\rangle^2=6|u|^2,
    \qquad
    \sum_{q\in2T}\langle u,q\rangle^4=3|u|^4.
\]
In the same edge-one normalization, each cell is a regular octahedron of edge
\(1\).  The \(24\)-cell inradius is
\[
    r_{24}=\frac1{\sqrt2},
\]
the octahedral cell has circumradius \(1/\sqrt2\), inradius \(1/\sqrt6\), and
three-volume \(\sqrt2/3\), the total boundary three-volume is \(8\sqrt2\),
the convex four-volume is \(2\), and the cell-cell dihedral angle is
\(2\pi/3\).

\smallskip
\noindent\textup{(v) Cube-link and trimer metric data.}
For two Hurwitz-link vertices
\(h_\varepsilon,h_\eta\in\mathcal N_8^+\), let \(k\) be the Hamming distance
between \(\varepsilon\) and \(\eta\).  Then
\[
    \langle h_\varepsilon,h_\eta\rangle=1-\frac{k}{2},
\]
and hence the distinct link angles are
\[
\begin{array}{c|c|c|c}
 k & \#\text{ unordered link pairs} & \langle h_\varepsilon,h_\eta\rangle &
 d_{S^3}(h_\varepsilon,h_\eta)\\
\hline
1 & 12 & \tfrac12 & \tfrac\pi3\\
2 & 12 & 0 & \tfrac\pi2\\
3 & 4 & -\tfrac12 & \tfrac{2\pi}{3}.
\end{array}
\]
The Hamming association scheme on the cube link is thereby fixed.  If \(B_r\)
denotes the relation matrix for Hamming distance \(r\) on \(\{\pm1\}^3\), then
\[
\begin{array}{c|c}
\text{matrix} & \text{spectrum}\\
\hline
B_1 & 3^{(1)},\;1^{(3)},\;(-1)^{(3)},\;(-3)^{(1)}\\
B_2 & 3^{(2)},\;(-1)^{(6)}\\
B_3 & 1^{(4)},\;(-1)^{(4)}.
\end{array}
\]
In particular, the ordinary cube-link graph has Laplacian spectrum
\[
    0^{(1)},\qquad 2^{(3)},\qquad 4^{(3)},\qquad 6^{(1)}.
\]
For the oriented diagonal trimer decomposition, the raw sign-cube lengths
\(\sqrt8\) and \(2\) become, in the Hurwitz coordinates,
\[
    \ell_{\rm tri}=\sqrt2,
    \qquad
    \ell_{\rm cross}=1.
\]
After normalizing sign vertices by \(1/\sqrt3\) and projecting perpendicular to
\(\nu=(1,1,1)/\sqrt3\), the six non-axial vertices have
\[
    h=\frac13,
    \qquad
    R_{\rm hex}=\frac{2\sqrt2}{3},
    \qquad
    r_\triangle=\frac{\sqrt2}{3}.
\]
Equivalently, the angle from a non-axial vertex to the chosen diagonal axis is
\(\arccos(1/3)\), and its elevation angle above the projected plane is
\(\arcsin(1/3)\).  The two trimer planes are separated by \(2/3\) in the
normalized sign cube, the projected hexagon has side length \(2\sqrt2/3\), the
alternating triangular faces have side length \(2\sqrt6/3\), and the nearest
cross-trimer side length is \(2/\sqrt3\).  Hence
\[
    \frac{s_{\rm tri}}{s_{\rm cross}}=\sqrt2.
\]
The two nearest-neighbor cyclic trimer-transport classes are precisely the two
half-step rotations \(\pm\pi/3\).

\smallskip
\noindent\textup{(vi) Rank-three trimer projector and transport trace invariants.}
For one oriented trimer with coordinate rank-one projectors \(P_a,P_b,P_c\in\End(\CC^T)\) and
\(I=P_a+P_b+P_c\), the scalar-centered apex defect and the oriented opposite-edge
defect
\[
    \DtriI=P_a-\frac13I,
    \qquad
    \DtriII=P_b-P_c
\]
are fixed by the trimer geometry.  With the primitive-idempotent trace
\(\operatorname{tr}_F(P_i)=1\), they satisfy
\[
    P_a=\frac13I+\DtriI,
    \qquad
    \operatorname{tr}_F\!\left(\left(\frac13I\right)^2\right)=\frac13,
    \qquad
    \operatorname{tr}_F({\DtriI}^2)=\frac23,
\]
\[
    \operatorname{tr}_F({\DtriII}^2)=2,
    \qquad
    \operatorname{tr}_F(\DtriI\DtriII)=0,
    \qquad
    \operatorname{tr}_F\!\left((\DtriI+\DtriII)^2\right)=\frac83.
\]
Thus
\[
    \|\DtriI\|_F=\sqrt{\frac23},
    \qquad
    \|\DtriII\|_F=\sqrt2,
    \qquad
    \|\DtriI\pm \DtriII\|_F=\sqrt{\frac83}.
\]
The diagonal spectra are fixed as
\[
    \operatorname{spec}(\DtriI)=\left\{\frac23,-\frac13,-\frac13\right\},
    \qquad
    \operatorname{spec}(\DtriII)=\{0,1,-1\},
\]
\[
    \operatorname{spec}(\DtriI\pm \DtriII)=
    \left\{\frac23,\frac23,-\frac43\right\}.
\]
Together with the scalar line and the cyclic transport plane below, the primitive
trimer operator basis
\[
    \left(\frac13I,\DtriI,\DtriII,C_1,C_2\right)
\]
has diagonal trace Gram matrix
\[
    \operatorname{diag}\left(\frac13,\frac23,2,6,6\right),
    \qquad
    \det=16.
\]
If \(X\) is the cyclic shift on \(\CC^T\cong\CC[Z_3]\), then
\[
    \operatorname{spec}(X)=\{1,\omega,\omega^2\},
    \qquad
    \omega=e^{2\pi i/3}.
\]
For
\[
    C_1=X+X^\dagger,
    \qquad
    C_2=-i(X-X^\dagger),
    \qquad
    C(\phi)=\cos\phi\,C_1+\sin\phi\,C_2,
\]
the transport plane is trace-orthogonal to the diagonal defect plane and satisfies
\[
    \operatorname{tr}_F(C_1)=\operatorname{tr}_F(C_2)=0,
    \qquad
    \operatorname{tr}_F(C_1^2)=\operatorname{tr}_F(C_2^2)=6,
    \qquad
    \operatorname{tr}_F(C_1C_2)=0,
\]
\[
    \operatorname{tr}_F(\DtriI C_j)=\operatorname{tr}_F(\DtriII C_j)=0
    \quad (j=1,2),
    \qquad
    \operatorname{tr}_F(C(\phi)^2)=6.
\]
The exact transport spectrum is
\[
    \operatorname{spec}(C(\phi))=
    \left\{2\cos\phi,
    2\cos\left(\frac{2\pi}{3}-\phi\right),
    2\cos\left(\frac{4\pi}{3}-\phi\right)\right\}.
\]
In particular
\[
    \operatorname{spec}(C_1)=\{2,-1,-1\},
    \qquad
    \operatorname{spec}(C_2)=\{0,\sqrt3,-\sqrt3\},
\]
and
\[
    \det C(\phi)=2\cos(3\phi),
    \qquad
    C(\phi)\text{ satisfies }\lambda^3-3\lambda-2\cos(3\phi)=0.
\]
These are projector and cyclic-transport identities internal to the selected
trimer support.

\smallskip
\noindent\textup{(vii) Uniform shell weights and the Bose--Mesner \(15+9\) rank split.}
\par\noindent
With the finite shell-complement counting convention, distinct from the
inverse-dimension carrier weights \(\wp(V)=1/\dim V\), the Hurwitz complement
split gives
\[
    w_7=\frac{7}{16},
    \qquad
    w_9=\frac{9}{16},
    \qquad
    w_9-w_7=\frac18.
\]
On the twenty-four-dimensional vertex-function space \(\CC^{2T}\), let
\(\mathcal E_9^{\rm BM}=\ker A_+\) and
\(\mathcal E_{15}^{\rm BM}=(\mathcal E_9^{\rm BM})^\perp\) be the
Bose--Mesner sectors introduced in
Theorem~\ref{p1:thm:s3one-trimer-bose-mesner-interface}.
The spectrum in part~\textup{(iv)} gives the genuine orthogonal
decomposition
\[
    \CC^{2T}=\mathcal E_{15}^{\rm BM}\oplus\mathcal E_9^{\rm BM},
    \qquad
    \dim\mathcal E_{15}^{\rm BM}=15,
    \qquad
    \dim\mathcal E_9^{\rm BM}=9.
\]
Equivalently,
\[
    P_{15}^{\rm BM}=\mathbf1_{\RR\setminus\{0\}}(A_+)
\]
is the orthogonal projector onto \(\mathcal E_{15}^{\rm BM}\).
The subspace \(\mathcal E_{15}^{\rm BM}\) is not itself a primitive idempotent of
the scheme: by Proposition~\ref{p1:prop:math-twentyfourcell-bose-mesner} it
is the sum of the primitive sectors of ranks \(1,4,8,2\), while
\(\mathcal E_9^{\rm BM}\) is the rank-\(9\) primitive sector.
Define the Bose--Mesner rank fraction
\[
    \chi_{24}:=\frac{\dim\mathcal E_9^{\rm BM}}{24}
    =\frac{24-15}{24}=\frac38,
    \qquad
    1-\chi_{24}
    =\frac{\dim\mathcal E_{15}^{\rm BM}}{24}
    =\frac58.
\]
Although \(\dim\mathcal H_1=\dim\mathcal E_{15}^{\rm BM}=15\), Paper~I-B specifies no
canonical or equivariant identification between those spaces.  Thus
\(24=15+9\) is a genuine rank decomposition of \(\CC^{2T}\), not a statement
that \(\mathcal H_1\) sits inside the Hurwitz vertex set or its function
space.

There is, however, an exact trace-normalization conversion wholly inside the
single space \(\CC^{2T}\).  For the inclusion
\(\iota_{15}:\mathcal E_{15}^{\rm BM}\hookrightarrow\CC^{2T}\), define
\[
    \tau_{24}(A):=\frac1{24}\Tr_{\CC^{2T}}A,
    \qquad
    \tau_{15}(B):=\frac1{15}\Tr_{\mathcal E_{15}^{\rm BM}}B.
\]
Then every \(B\in\End(\mathcal E_{15}^{\rm BM})\) satisfies
\[
    \tau_{24}(\iota_{15} B\iota_{15}^*)
    =\frac{15}{24}\tau_{15}(B),
    \qquad
    \tau_{15}(B)
    =\frac{24}{15}\tau_{24}(\iota_{15} B\iota_{15}^*).
\]
This is the precise common-space setting in which \(24/15\) is a trace
conversion.  It neither identifies \(\mathcal H_1\) with
\(\mathcal E_{15}^{\rm BM}\) nor determines an amplitude normalization.

\end{theorem}

\begin{proof}
\ProofPlan{We verify the seven displayed finite-cell data in order: Casimir values, protected Clifford data, Hurwitz shell distances, \(24\)-cell combinatorics, link-cube data, trimer defects, and the finite shell and Bose--Mesner rank counts.}

\begin{proofsteps}[label=\textbf{(\roman*)}]
\item \textbf{Casimir values.}
This is the standard \(SU(2)_L\times SU(2)_R\) quadratic-Casimir value
\(j(j+1)\) applied to the six sectors of
Corollary~\ref{p1:cor:s3one-sector-content}.

\item \textbf{Protected Clifford data.}
The case \(a=b\) in
Proposition~\ref{p1:prop:general-protected-clifford-relation} gives
\(X_a^2=\rho_{\mathsf C}^2\mathbf1_{\mathcal K}/4\).  On the stipulated
positive stratum the \(X_a\) are self-adjoint and anticommute with the other
invertible Clifford generators; hence their traces vanish and the two eigenvalues have
equal multiplicity on \(\dim\mathcal K_{\mathsf B}=4\). Taking traces in the
Clifford relation gives the displayed Hilbert--Schmidt Gram matrix. The
normalized generators \(\Gamma_a\) obey the standard rank-four Clifford
relations; the pseudoscalar \(\Gamma_5\), chiral projectors, grade counts,
bivector norms, and the \(3+3\) split are the corresponding Euclidean Clifford
\(\Lambda^2\RR^4=\Lambda^2_+\oplus\Lambda^2_-\) decomposition on the protected
spinor block.

\item \textbf{Hurwitz shell distances.}
This is Theorem~\ref{p1:thm:hurwitz-radial-shells} with
\(|1-q|^2=2-2\operatorname{Re}(q)\) and
\(d_{S^3}(1,q)=\arccos(\operatorname{Re}q)\) on the unit sphere.

\item \textbf{\(24\)-cell combinatorics.}
This combines the regular \(24\)-cell combinatorics of
Section~\ref{p1:sec:24cell} with the distance shell counts of the previous step,
dividing vertex-based counts by two to obtain unordered pair counts. The
relation-matrix spectra and multiplication rules are obtained by multiplying the
five inner-product relation matrices of the Hurwitz unit set; equivalently, they
are the Bose--Mesner relations of this antipodal \(D_4\)-root configuration.
To make the diagonalization explicit, simultaneous evaluation of those
relations gives the five joint characters
\[
\begin{array}{c|rrrr|c}
\text{rank}&A_+&A_0&A_-&J&\text{antipodal parity}\\
\hline
1&8&6&8&1&+\\
4&4&0&-4&-1&-\\
9&0&-2&0&1&+\\
8&-2&0&2&-1&-\\
2&-4&6&-4&1&+
\end{array}
\]
Indeed, the multiplication table is commutative, \(A_- = A_+J\), and its
one-dimensional algebra characters are exactly the five rows above.
Applying the ordinary matrix trace to the corresponding Lagrange
idempotents---using
\(\Tr I=24\) and zero trace for each nonidentity relation
matrix---gives the ranks \(1,4,9,8,2\).  These ranks sum to \(24\), and each
column now yields the complete spectrum and multiplicity stated in
part~\textup{(iv)}.  This is also an exact reconstruction certificate for
all displayed relation spectra, rather than an appeal to an unspecified
diagonalization.
The tight-frame and moment identities follow by summing the explicit Hurwitz
coordinates, with antipodal symmetry eliminating odd moments. The supporting
hyperplanes of the Hurwitz realization are of the form \(\pm x_a\pm x_b=1\);
their distance from the origin is \(1/\sqrt2\), and each supports a regular
octahedral cell of edge \(1\). The octahedral metric values and volumes are the
standard edge-one octahedron formulae, and the \(24\)-cell volume follows from
pyramid decomposition using inradius \(1/\sqrt2\) and total boundary
three-volume \(24(\sqrt2/3)=8\sqrt2\).

\item \textbf{Hurwitz-link cube data.}
This follows from the displayed Hurwitz-link coordinates:
\[
\langle h_\varepsilon,h_\eta\rangle
 =\frac{1+\varepsilon\cdot\eta}{4}
 =1-\frac{k}{2},
\]
together with Theorem~\ref{p1:thm:hurwitz-diagonal-trimer},
Proposition~\ref{p1:prop:projected-trimer-metric}, and
Proposition~\ref{p1:prop:antiprism-transport-matchings}. The pair counts
\(12,12,4\), the Hamming-scheme spectra, and the cube Laplacian spectrum are the
standard Fourier/Walsh eigenvalues of \((\ZZ_2)^3\). The additional antiprism
lengths are obtained by resolving the normalized sign cube into the chosen axis
and its orthogonal plane.

\item \textbf{Trimer defects and diagonal operator space.}
This is Corollary~\ref{p1:cor:opposite-trimer-defects} together with the
cyclic-shift relations \(X^3=1\) and \(X^\dagger=X^2\) in
Theorem~\ref{p1:thm:cyclic-trimer-operator-grammar}. The diagonal matrices have
zero trace against the off-diagonal transport matrices, while the Fourier basis
of \(Z_3\) gives the listed spectra, characteristic equation, determinant, and
transport-plane traces.

\item \textbf{Finite shell and Bose--Mesner rank counts.}
This is the normalized form of the \(7+9\) Hurwitz-shell complement in
Theorem~\ref{p1:thm:hurwitz-radial-shells}, relative to its displayed Hurwitz
representative, base vertex, and counting measure.  From part
\textup{(iv)}, \(A_+\) is real symmetric and its zero eigenvalue has
multiplicity \(9\).  Therefore
\(\CC^{2T}=\ker(A_+)^\perp\oplus\ker(A_+)\) has dimensions \(15+9\), proving
the rank statements and \(\chi_{24}=9/24=3/8\).  Moreover,
\(\Tr_{\CC^{2T}}(\iota_{15}B\iota_{15}^*)
=\Tr_{\mathcal E_{15}^{\rm BM}}B\), which proves the
trace-conversion formula after applying the two normalized traces.  No identification with \(\mathcal H_1\) is used.
\end{proofsteps}
\end{proof}


\begin{theorem}[Geometric structural invariants of the \(S^3_1\) cell]
\label{p1:thm:s3one-structural-invariants}
In addition to the numerical data of
Theorem~\ref{p1:thm:s3one-numerical-data}, the \(m=1\) finite
cell carries the following structural invariants.  They are fixed by the
ordinary geometry and finite representation theory of
\[
    \mathcal H_1,
    \qquad
    \mathcal K_{\mathsf B},
    \qquad
    2T\subset S^3,
    \qquad
    \text{the chosen vertex and oriented cube-link trimer},
\]
viewed as a finite Hurwitz--Clifford--trimer cell.  Entries
\textup{(i)}--\textup{(xiv)} use only the displayed module/support data.
Entry~\textup{(xv)} additionally assumes
\(\rho_{\mathsf C}=2\varrho_{\rm sp}>0\); at zero radius only the
unnormalized relation and the chirality commutant remain.

\smallskip
\noindent\textup{(i) Stabilizer chain of the \(m=1\) finite frame.}
The regular \(24\)-cell has full symmetry group
\[
    |\Aut(24\text{-cell})|=1152,
    \qquad
    |\Aut^+(24\text{-cell})|=576.
\]
Fixing a Hurwitz vertex leaves the cube symmetry of the vertex figure:
\[
    1152\longrightarrow 48,
    \qquad
    576\longrightarrow 24.
\]
Fixing an oriented body diagonal in that cube leaves
\[
    48\longrightarrow 6,
    \qquad
    24\longrightarrow 3.
\]
Thus the choice-relative stabilizer chains are
\[
    1152\supset 48\supset 6,
    \qquad
    576\supset 24\supset 3,
\]
and the number of oriented vertex-diagonal local frames is
\[
    24\cdot 8=192.
\]

\smallskip
\noindent\textup{(ii) Binary tetrahedral group structure.}
The Hurwitz vertex set is not only a spherical code; under quaternion
multiplication it is the binary tetrahedral group
\[
    2T\cong Q_8\rtimes C_3,
\]
where
\[
    Q_8=\{\pm1,\pm i,\pm j,\pm k\}
\]
is normal.  Hence
\[
    2T/Q_8\cong C_3,
    \qquad
    2T/\{\pm1\}\cong A_4.
\]
The metric shells about the identity also refine by group order:
\[
\begin{array}{c|c|c}
\text{subset} & \# & \text{group order}\\
\hline
\{1\} & 1 & 1\\
\{-1\} & 1 & 2\\
\{\pm i,\pm j,\pm k\} & 6 & 4\\
\Re(q)=\tfrac12 & 8 & 6\\
\Re(q)=-\tfrac12 & 8 & 3.
\end{array}
\]

\smallskip
\noindent\textup{(iii) Triality, octets, and the antipodal quotient graph.}
The \(24\)-cell vertices form \(12\) antipodal pairs.  The root-system triality
admits a decomposition of these pairs into three orthonormal \(4\)-frames,
\[
    12=4+4+4,
\]
or, equivalently, a decomposition of the \(24\) vertices into three disjoint
\(16\)-cell octets,
\[
    24=8+8+8.
\]
After choosing a \(D_4\) triality frame, the vertices split into three \(16\)-cell octets.  Different triality frames are related by the \(S_3\) triality action, so only the existence of such a three-octet decomposition is invariant, not a preferred labeling of the three octets.  In the
antipodal quotient, the edge graph is
\[
    K_{4,4,4}.
\]
Thus the level-\(1\) Hurwitz support contains a purely geometric source of three
maximum independent octets.

\smallskip
\noindent\textup{(iv) \(D_4\), \(F_4\), and Coxeter invariant data.}
Under the explicit map of
Proposition~\ref{p1:prop:hurwitz-d4-coordinate-map}, the Hurwitz unit model is
the unit-sphere normalization of the \(D_4\) root system:
\[
    O(2T)=\frac1{\sqrt2}\Phi(D_4),
    \qquad
    \sqrt2\,O(2T)=\Phi(D_4),
    \qquad
    |\Phi(D_4)|=24.
\]
The Coxeter number and exponents are
\[
    h(D_4)=6,
    \qquad
    1,3,3,5,
\]
and
\[
    |W(D_4)|=192.
\]
The full \(24\)-cell symmetry is the triality extension
\[
    W(F_4)\cong W(D_4)\rtimes S_3,
    \qquad
    192\cdot 6=1152.
\]
The polynomial invariant degrees of the full \(F_4\)-symmetry are
\[
    2,\;6,\;8,\;12.
\]
In particular, after the quadratic invariant, the first full-symmetry invariant
degree is \(6\).

\smallskip
\noindent\textup{(v) Spherical design strength and first failure.}
The Hurwitz \(24\)-cell is a spherical \(5\)-design on \(S^3\) and is not a
spherical \(6\)-design.  Equivalently, for every nonconstant spherical harmonic
\(Y_\ell\) of degree \(1\le \ell\le5\),
\[
    \sum_{q\in2T}Y_\ell(q)=0.
\]
The first failure occurs in degree \(6\).  In the coordinate system of the
Hurwitz units, for instance,
\[
    \sum_{q\in2T}\langle e_1,q\rangle^6=\frac94,
\]
whereas the corresponding uniform \(S^3\)-average multiplied by \(24\) would be
\[
    24\,\mathbb E_{S^3}(x_1^6)=\frac{15}{8}.
\]

\smallskip
\noindent\textup{(vi) Spherical packing, covering, and self-duality.}
As a spherical code in \(S^3\), the Hurwitz support has minimum angle
\[
    \theta_{\min}=\frac\pi3,
\]
so its spherical packing radius is
\[
    \rho_{\rm pack}=\frac\pi6.
\]
Its covering radius is
\[
    \rho_{\rm cover}=\frac\pi4,
\]
equivalently the farthest points have nearest-vertex inner product
\[
    \frac1{\sqrt2}.
\]
The \(24\)-cell is self-dual.  The normalized facet centers form another
\(24\)-cell, and for a primal vertex \(p\) and a dual vertex \(n\),
\[
    \langle p,n\rangle\in
    \left\{\frac1{\sqrt2},0,-\frac1{\sqrt2}\right\}
\]
with counts
\[
    6,
    \qquad
    12,
    \qquad
    6.
\]
Thus the primal--dual incidence matrix is a \(24\times24\) matrix with row and
column sum \(6\).  The radial spherical tessellation of \(S^3\) has \(24\)
congruent spherical octahedra, each of spherical volume
\[
    \frac{\operatorname{vol}(S^3)}{24}=\frac{2\pi^2}{24}=\frac{\pi^2}{12}.
\]

\smallskip
\noindent\textup{(vii) Great polygons and rank-two root subsystems.}
The \(D_4\) root system contains the rank-two subsystems
\[
    \#A_2=16,
    \qquad
    \#(A_1\times A_1)=18.
\]
Geometrically, the \(24\)-cell contains
\[
    16\quad\text{great regular hexagons},
    \qquad
    18\quad\text{great squares}.
\]
Each edge lies in exactly one great hexagon,
\[
    16\cdot6=96,
\]
and each vertex lies in four great hexagons.  The three triality
\(16\)-cells are the three \(A_1^4\) decompositions.

\smallskip
\noindent\textup{(viii) Spherical face angles and curvature closure.}
The spherical edge length of every triangular face is \(\pi/3\).  Hence the
spherical triangle face angle is
\[
    \alpha_\triangle=\arccos\!\left(\frac13\right),
\]
and the spherical area of each triangular face is
\[
    3\arccos\!\left(\frac13\right)-\pi.
\]
At a vertex the link is a cube; the six incident octahedral cells therefore
partition the unit tangent \(S^2\) into six congruent square solid-angle regions,
each of area
\[
    \frac{4\pi}{6}=\frac{2\pi}{3}.
\]
The dihedral angle of the Euclidean \(24\)-cell between adjacent octahedral
cells, measured across their common triangular face, is \(2\pi/3\).  In the
radial spherical tessellation of \(S^3\) by the \(24\) octahedral cells, each
triangular face is shared by exactly two cells, while each edge is incident to
exactly three cells; each cell has edge dihedral angle \(2\pi/3\), and the
three-cell angular closure holds around every edge:
\[
    3\cdot\frac{2\pi}{3}=2\pi.
\]

\smallskip
\noindent\textup{(ix) Cube-link diagonal removal and the fixed
\(C_6^{\rm link}\).}
After a vertex and oriented body diagonal are fixed, the cube-link layers are
\[
    1,3,3,1.
\]
Removing the two axial vertices leaves a fixed six-cycle
\[
    C_6^{\rm link}.
\]
The two trimers are the alternating vertices of this cycle:
\[
    V(C_6^{\rm link})=T_+\sqcup T_-.
\]
The two nearest-neighbor trimer transports are precisely the two perfect
matchings of this \(C_6^{\rm link}\) compatible with the cyclic orientation.

\smallskip
\noindent\textup{(x) \(A_2\) root and weight geometry of one trimer.}
\par\noindent
For coordinate rank-one trimer projectors
\[
    P_a,P_b,P_c\in\End(\CC^T),
\]
set
\[
    \alpha_1=P_a-P_b,
    \qquad
    \alpha_2=P_b-P_c.
\]
Then
\[
    \operatorname{tr}_F(\alpha_i^2)=2,
    \qquad
    \operatorname{tr}_F(\alpha_1\alpha_2)=-1,
\]
so the trimer carries the \(A_2\) Cartan matrix
\[
    A_{A_2}=\begin{pmatrix}2&-1\\-1&2\end{pmatrix},
    \qquad
    A_{A_2}^{-1}=\frac13\begin{pmatrix}2&1\\1&2\end{pmatrix}.
\]
The scalar-centered defect
\[
    \DtriI=P_a-\frac13 I
\]
is the first fundamental weight of this \(A_2\) system:
\[
    \DtriI=\omega_1,
    \qquad
    \operatorname{tr}_F({\DtriI}^2)=\frac23.
\]
Thus the trace constants \(1/3\), \(2/3\), \(2\), and \(3\) in the trimer block
are root/weight-lattice constants of \(A_2\).

\smallskip
\noindent\textup{(xi) Weyl--Heisenberg algebra of the cyclic trimer.}
Once the cyclic orientation of a trimer is fixed, let \(X\) be the cyclic shift
and let
\[
    Z=\operatorname{diag}(1,\omega,\omega^2),
    \qquad
    \omega=e^{2\pi i/3}.
\]
Then
\[
    ZX=\omega XZ.
\]
The nine operators
\[
    X^aZ^b,
    \qquad
    a,b\in\ZZ_3,
\]
form a trace-orthogonal basis of \(M_3(\CC)\):
\[
    \operatorname{tr}\bigl((X^aZ^b)^\dagger X^{a'}Z^{b'}\bigr)
    =3\delta_{a,a'}\delta_{b,b'}.
\]
Thus the oriented trimer carries the finite Weyl algebra of a cyclic
three-point system.

\smallskip
\noindent\textup{(xii) Binary-tetrahedral McKay invariant.}
For the defining two-dimensional representation of \(2T\subset SU(2)\), the
McKay graph is the affine \(E_6\) diagram.  The irreducible complex
representation dimensions are
\[
    1,1,1,2,2,2,3,
\]
whose squares sum to \(24\).  The corresponding McKay dimension vector is the
positive null root of \(\widetilde E_6\).  Equivalently, the quotient
\(\CC^2/2T\) is the Kleinian surface singularity of type \(E_6\).  This is a
finite-group invariant of the Hurwitz support itself.

\smallskip
\noindent\textup{(xiii) Coordinate-coupling selection graph on \(\mathcal H_1\).}
The coordinate representation is
\[
    V=(\tfrac12,\tfrac12).
\]
Tensoring each retained \(\operatorname{Spin}(4)\)-sector by \(V\), then projecting back to
\(\mathcal H_1\), gives the retained coupling channels
\[
    (0,0)\leftrightarrow(\tfrac12,\tfrac12),
\]
\[
    (\tfrac12,0)\leftrightarrow(0,\tfrac12),
\]
\[
    (\tfrac12,\tfrac12)\leftrightarrow(1,0),
    \qquad
    (\tfrac12,\tfrac12)\leftrightarrow(0,1).
\]
The full tensor decompositions are
\[
\begin{aligned}
V\otimes(0,0)&=(\tfrac12,\tfrac12),\\
V\otimes(\tfrac12,0)&=(0,\tfrac12)\oplus(1,\tfrac12),\\
V\otimes(0,\tfrac12)&=(\tfrac12,0)\oplus(\tfrac12,1),\\
V\otimes(\tfrac12,\tfrac12)&=(0,0)\oplus(1,0)\oplus(0,1)\oplus(1,1),\\
V\otimes(1,0)&=(\tfrac12,\tfrac12)\oplus(\tfrac32,\tfrac12),\\
V\otimes(0,1)&=(\tfrac12,\tfrac12)\oplus(\tfrac12,\tfrac32).
\end{aligned}
\]
Thus the projected coordinate action has exactly four retained channel types;
the first exterior sectors are also fixed by representation theory.

\smallskip
\noindent\textup{(xiv) Full \(\mathcal H_1\) Casimir trace package.}
The total-Casimir spectrum on the level-\(1\) carrier is
\[
    \operatorname{spec}_{\mathcal H_1}(C_{\rm tot})
    =0^{(1)}
    \oplus\left(\frac34\right)^{(4)}
    \oplus\left(\frac32\right)^{(4)}
    \oplus2^{(6)}.
\]
Consequently
\[
    \operatorname{Tr}_{\mathcal H_1}(C_{\rm tot})=21,
    \qquad
    \frac1{15}\operatorname{Tr}_{\mathcal H_1}(C_{\rm tot})=\frac75,
\]
\[
    \operatorname{Tr}_{\mathcal H_1}(C_{\rm tot}^2)=\frac{141}{4},
    \qquad
    \operatorname{Tr}_{\mathcal H_1}(C_L)
    =\operatorname{Tr}_{\mathcal H_1}(C_R)=\frac{21}{2}.
\]

\smallskip
\noindent\textup{(xv) Protected Clifford--Fierz completeness.}
If \(\rho_{\mathsf C}>0\), the protected Clifford block carries the
trace-orthogonal Fierz basis of
\(\End(\mathcal K_{\mathsf B})\)
\[
    \mathbf1,
    \qquad
    \Gamma_a,
    \qquad
    \Gamma_{ab},
    \qquad
    \Gamma_5\Gamma_a,
    \qquad
    \Gamma_5,
\]
with grade dimensions
\[
    1+4+6+4+1=16.
\]
Here \(\Gamma_{ab}:=\Gamma_a\Gamma_b\) for \(a<b\). With the standard Clifford trace normalization,
\[
    \operatorname{Tr}_{\mathcal K}(\Gamma_I^\dagger\Gamma_J)=4\delta_{IJ}.
\]
The even Clifford algebra splits chirally as
\[
    \Cl^0_4(\CC)\cong M_2(\CC)\oplus M_2(\CC),
\]
with central idempotents
\[
    P_\pm=\frac{1\pm\Gamma_5}{2},
    \qquad
    \operatorname{rank}P_\pm=2.
\]
The traceless part of \(\End(\mathcal K_{\mathsf B})\) has dimension \(15\),
matching \(\dim\mathcal H_1=15\).  This last equality is only a dimension
identity unless further structure is supplied.
\end{theorem}

\begin{proof}
\ProofPlan{We verify the structural entries by reducing each item to either an explicit Hurwitz-coordinate calculation, a standard \(D_4/F_4\) root-system fact, or a finite matrix identity on the protected block.}

\begin{proofsteps}[label=\textbf{(\roman*)}]
\item \textbf{Stabilizers.}
The full symmetry group of the regular \(24\)-cell is the Coxeter group \(W(F_4)\), of order \(1152\), and its orientation-preserving subgroup has order \(576\).  The vertex figure at any Hurwitz vertex is a cube.  Hence the stabilizer of one vertex is the full cube symmetry group of order \(48\), with orientation-preserving subgroup of order \(24\).  A cube has four unoriented body diagonals and eight oriented body diagonals; fixing one oriented diagonal leaves the permutations of the three transverse coordinate directions, of order \(6\), and in the orientation-preserving subgroup leaves the cyclic order-three rotations about the diagonal.  This gives
\[
1152\supset48\supset6,
\qquad
576\supset24\supset3,
\]
and the number of oriented vertex--diagonal local frames is \(24\cdot8=192\).

\item \textbf{Binary tetrahedral structure.}
The Hurwitz units are
\[
\{\pm1,\pm i,\pm j,\pm k\}
\quad\text{and}\quad
\frac12(\pm1\pm i\pm j\pm k).
\]
They are closed under quaternion multiplication and inversion, giving the binary tetrahedral group \(2T\).  The subgroup \(Q_8=\{\pm1,\pm i,\pm j,\pm k\}\) is normal, the quotient has order three, and quotienting by the center \(\{\pm1\}\) gives the ordinary tetrahedral rotation group \(A_4\).  The shell table follows by reading the real parts \(1,0,\pm1/2,-1\) from the displayed units; the corresponding cyclic orders are the orders of \(\cos\theta+u\sin\theta\) in \(SU(2)\).

\item \textbf{Triality octets and quotient graph.}
In the standard \(D_4\) model
\[
\Phi(D_4)=\{\pm e_i\pm e_j:1\le i<j\le4\},
\]
the twelve antipodal axes split into the three tetrads associated with the three pairings of \(\{1,2,3,4\}\).  Each tetrad lifts to a \(16\)-cell octet, so \(24=8+8+8\).  Axes in the same tetrad are orthogonal, while axes from two different tetrads have representatives with inner product \(\pm1/2\).  Therefore the antipodal quotient edge graph is \(K_{4,4,4}\), and the outer \(D_4\) triality action permutes the three tetrads.

\item \textbf{\(D_4/F_4\) data.}
The \(D_4\) root system has \(24\) roots, Coxeter number \(6\), exponents \(1,3,3,5\), and Weyl-group order
\[
|W(D_4)|=2^{3}\,4!=192.
\]
Adjoining the triality symmetry gives the full \(24\)-cell symmetry
\[
W(F_4)\cong W(D_4)\rtimes S_3,
\qquad
192\cdot6=1152.
\]
The fundamental invariant degrees of \(W(F_4)\) are \(2,6,8,12\), so the first non-radial full-symmetry invariant on the unit sphere occurs in degree \(6\).

\item \textbf{Design strength and first failure.}
For the Hurwitz units, antipodality eliminates all odd moments.  The second and fourth coordinate moments agree with the corresponding moments of the uniform measure on \(S^3\), which is the moment form of the spherical \(5\)-design property.  The degree-six failure is explicit: along a coordinate axis, the values are two entries \(\pm1\), six entries \(0\), and sixteen entries \(\pm1/2\), hence
\[
\sum_{q\in2T} q_1^6=2+16\left(\frac1{2^6}\right)=\frac94.
\]
For the uniform measure on \(S^3\subset\RR^4\),
\[
\mathbb E(x_1^6)=\frac{15}{4\cdot6\cdot8}=\frac5{64},
\]
so \(24\mathbb E(x_1^6)=15/8\), as displayed.

\item \textbf{Packing, covering, duality, and face geometry.}
The maximum inner product among distinct Hurwitz units is \(1/2\), giving minimum spherical distance \(\pi/3\) and packing radius \(\pi/6\).  The dual \(24\)-cell facet centers are at nearest-vertex inner product \(1/\sqrt2\), giving covering radius \(\pi/4\).  Self-duality gives another \(24\)-cell of facet normals, and the signs of the defining octahedral hyperplanes yield the \(6,12,6\) primal--dual incidence counts.  For a spherical equilateral triangle with side \(\pi/3\), the spherical law of cosines gives
\[
\cos\alpha_\triangle=
\frac{\cos(\pi/3)-\cos^2(\pi/3)}{\sin^2(\pi/3)}=\frac13,
\]
so its area is \(3\arccos(1/3)-\pi\).  The cube link at a vertex partitions the unit tangent two-sphere into six square solid-angle regions of area \(2\pi/3\).  Each edge of the radial spherical tessellation is incident to exactly three octahedral cells by the incidence data of Theorem~\ref{p1:thm:s3one-numerical-data}(iv); by flag-transitivity of the symmetry group, the stabilizer of an edge permutes these three cells transitively, so each contributes an equal dihedral angle along the edge, and since the total transverse angle around an interior point of a geodesic edge in \(S^3\) is \(2\pi\), each cell has edge dihedral angle \(2\pi/3\) and the three-cell closure \(3\cdot(2\pi/3)=2\pi\) holds around every edge.

\item \textbf{Great polygons.}
For every edge \(\{p,q\}\), normalized so
\(\langle p,q\rangle=1/2\), the six roots
\[
    \{\pm p,\ \pm q,\ \pm(p-q)\}
\]
form the unique \(A_2\) subsystem in their two-plane.  Thus every one of the
\(96\) edges lies in one great root hexagon, and every such hexagon contains
six edges.  Hence
\[
    \#A_2=\frac{96}{6}=16.
\]
In the antipodal quotient, each of the \(12\) root axes is orthogonal to the
other three axes in its triality tetrad and to no axis outside that tetrad.
Therefore the number of unordered orthogonal axis pairs is
\[
    \#(A_1\times A_1)=\frac{12\cdot3}{2}=18.
\]
These subsystems are respectively the great regular hexagons and great
squares.  Finally, \(16\cdot6=96\) and
\((16\cdot6)/24=4\) give the asserted edge and vertex incidences.

\item \textbf{Cube-link trimer and \(A_2\)/Weyl algebra.}
Deleting the chosen oriented antipodal pair from the cube vertex figure leaves
the induced six-cycle \(C_6^{\rm link}\).  Its alternating vertices are the
two oriented trimers, and the two cyclic nearest-neighbor matchings are the
two half-step rotations.  This graph is distinct from the cyclic stabilizer
group \(C_6^{\rm grp}\).  On one trimer, the projectors \(P_a,P_b,P_c\)
satisfy \(P_iP_j=\delta_{ij}P_i\), so the two differences \(P_a-P_b\) and
\(P_b-P_c\) have Gram matrix the \(A_2\) Cartan matrix.  The cyclic shift
\(X\) and clock matrix \(Z\) obey \(ZX=\omega XZ\), and the nine operators
\(X^aZ^b\) are trace-orthogonal because their traces are the finite Fourier
sums on \(\ZZ_3\).

\item \textbf{McKay, coupling, Casimirs, and Clifford completeness.}
The binary tetrahedral subgroup of \(SU(2)\) has McKay graph \(\widetilde E_6\), with irreducible dimensions \(1,1,1,2,2,2,3\) and squared-dimension sum \(24\).  The coordinate-coupling graph is obtained by applying the Clebsch--Gordan rule for \(SU(2)_L\times SU(2)_R\) to multiplication by \((1/2,1/2)\) and retaining only summands in \(\mathcal H_1\).  The Casimir trace package is the dimension-weighted sum of the sector values in Theorem~\ref{p1:thm:s3one-numerical-data}.  On the stipulated positive protected stratum, four gamma matrices generate the complex rank-four Clifford algebra on \(\mathcal K_{\mathsf B}\), whose grade dimensions are \(1,4,6,4,1\); Schur's lemma on \((1/2,0)\oplus(0,1/2)\) leaves exactly the commutant line \(\CC1\oplus\CC\Gamma_5\).
\end{proofsteps}
\end{proof}

\begin{definition}[Primitive Hurwitz angular cell and scalarizer]
\label{p1:def:primitive-hurwitz-cell-scalarizer}
For the unit Hurwitz \(24\)-cell, the nearest-neighbour angular edge is
\[
 \boxed{\theta_H:=\frac{\pi}{3}.}
\]
The preceding finite-data theorem also proves that every such edge belongs to
a unique great \(A_2\) root hexagon with six congruent edges.  Hence
\[
 \boxed{6\theta_H=2\pi.}
\]
Define the dimensionless primitive finite-cell scalarizer
\[
 \boxed{\rho_H:=\theta_H^{-1}=\frac3\pi=\frac6{2\pi}.}
\]
Equivalently, if the primitive edge-cell count is normalized to unit mass,
\(d\mu_H=\rho_H\,d\varphi\) on the associated smooth great-circle angular
coordinate.  This is a finite-cell normalization certificate, not a
three-dimensional Riemannian volume density.
\end{definition}

\begin{remark}[Typing of the Hurwitz scalarizer]
\label{p1:rem:primitive-hurwitz-cell-scalarizer-scope}
The scalarizer \(\rho_H\) is the inverse primitive edge angle of the
unmarked Hurwitz geometry.  It is distinct from a Hilbert-space
multiplicity, a heat-kernel moment, the three-dimensional volume density,
and the Hopf/contact integral \(4\pi^2\).
\end{remark}

\begin{theorem}[Classical completions carried by the level-\(1\) Hurwitz support]
\label{p1:thm:s3one-classical-completions}
The same level-\(1\) Hurwitz support also carries the following classical geometric
completions.  These are consequences of recognizing the vertex set as the
Hurwitz unit group, the $D_4$ root system, and the regular $24$-cell.

\smallskip
\noindent\textup{(i) Optimal kissing support, but not universal optimality.}
In its $D_4$ root-system normalization, $2T$ is the unique optimal kissing
configuration in four dimensions, up to orthogonal equivalence, and is an
optimal spherical code of $24$ points on $S^3$, as proved in the cited 2024
preprint \cite{deLaatLeijenhorstDeMuinckKeizer2024}:
\[
    |2T|=24,
    \qquad
    \min_{q\ne q'} d_{S^3}(q,q')=\frac\pi3.
\]
This statement is deliberately the kissing/spherical-code optimality statement;
it is not the stronger statement of universal optimality.  In fact the $D_4$
root system is not universally optimal among $24$-point codes in $S^3$.

\smallskip
\noindent\textup{(ii) $D_4$ lattice Voronoi cell and the $24$-cell honeycomb.}
Let
\[
    D_4=\{x\in\ZZ^4:\;x_1+x_2+x_3+x_4\equiv0\pmod2\}.
\]
After the standard Euclidean scaling, the Voronoi cell of the $D_4$ lattice is
the regular $24$-cell.  Translating this Voronoi cell by $D_4$ gives the regular
$24$-cell honeycomb
\[
    \{3,4,3,3\}.
\]
Each Voronoi $24$-cell has $24$ facet-neighbours, one across each octahedral
facet.  After choosing the standard $D_4$ lattice embedding and Euclidean scaling, the finite level-\(1\) support extends to the local
Voronoi/tessellation rule of the $D_4$ lattice.

\smallskip
\noindent\textup{(iii) Hurwitz order and unit closure.}
Let $\mathcal O_H$ denote the Hurwitz order in the quaternions.  Its norm-one
unit group is precisely the Hurwitz vertex set:
\[
    \mathcal O_H^\times=2T.
\]
Consequently the level-\(1\) $24$-point support is closed under quaternion
multiplication and inversion:
\[
    q,q'\in2T\quad\Longrightarrow\quad qq'\in2T,
    \qquad
    q^{-1}=\bar q\in2T.
\]
The $24$ vertices therefore carry, simultaneously, the metric structure of the
regular $24$-cell and the multiplication table of the binary tetrahedral group.

\smallskip
\noindent\textup{(iv) Primal--dual union and the $F_4$ root system.}
The regular $24$-cell is self-dual, but the primal and dual copies are rotated
relative to one another.  With the standard $F_4$ root-system scaling, the level-\(1\) unit $D_4$ shell is rescaled to the long-root shell, while the rotated dual $24$-cell is placed at the short-root scale; their union is the $F_4$ root system:
\[
    \Phi(F_4)=\Phi_{\mathrm{long}}(D_4)\sqcup\Phi_{\mathrm{short}},
    \qquad
    |\Phi(F_4)|=48=24+24,
\]
where \(\Phi_{\mathrm{short}}\) is the rotated dual $24$-cell, consisting of the
short roots \(\{\pm e_i\}\) together with the half-sum vertices
\(\frac12(\pm e_1\pm e_2\pm e_3\pm e_4)\).  This gives a second classical way to see why the full symmetry group of the
cell is $W(F_4)$ while the level-\(1\) vertex shell itself is the $D_4$ root shell.

\smallskip
\noindent\textup{(v) Reye configuration of axes and great hexagons.}
Passing to antipodal pairs, the $24$ vertices determine $12$ axes.  Together
with the $16$ great-hexagon planes of the $24$-cell, these form the Reye
configuration
\[
    12_4\,16_3.
\]
Equivalently, each antipodal axis lies in exactly four great hexagon planes,
and each great hexagon plane contains exactly three antipodal axes.  The
incidence count closes as
\[
    12\cdot4=16\cdot3=48.
\]
Thus the \(S^3_1\) cell contains a projective incidence geometry on its
antipodal quotient, not only a metric graph.

\smallskip
\noindent\textup{(vi) Design hierarchy and the degree-six obstruction.}
The spherical $5$-design statement of
Theorem~\ref{p1:thm:s3one-structural-invariants} is the cubature side
of the $D_4/F_4$ invariant theory.  The invariant degrees of $W(F_4)$ are
\[
    2,6,8,12,
\]
so after the quadratic form the first nonconstant full-symmetry invariant occurs
in degree $6$.  This matches the first failure of the $24$-point support to be a
spherical design: it is exact through degree $5$ and obstructed at degree $6$.

\smallskip
\noindent\textup{(vii) Classical reference package.}
The regular-polytope, root-system, lattice, quaternionic, spherical-code, and
McKay identifications used above are the standard ones recorded in the classical
literature on regular polytopes, lattices, reflection groups, Hurwitz
quaternions, spherical designs, and finite subgroups of $SU(2)$
\cite{coxeter1973regular,conwaySloane1999sphere,Humphreys1990,conwaysmith2003quaternions,DelsarteGoethalsSeidel1977,CohnConwayElkiesKumar2007,deLaatLeijenhorstDeMuinckKeizer2024,McKay1980}.
\end{theorem}

\begin{proof}
\ProofPlan{We separate the classical completions from the \(S^3_1\) finite-geometry datum: each completion is a standard structure carried by the Hurwitz unit set after it is identified with the \(D_4\) root shell and the regular \(24\)-cell.}

\begin{proofsteps}[label=\textbf{(\roman*)}]
\item \textbf{Kissing and spherical-code optimality.}
The Hurwitz unit set is the \(D_4\) root shell on \(S^3\), whose nearest-neighbor inner product is \(1/2\), hence whose minimum spherical distance is \(\pi/3\).  The optimality and uniqueness statement used here is the kissing/spherical-code theorem for the \(D_4\) root system; the text separately records the theorem of Cohn--Conway--Elkies--Kumar that this configuration is not universally optimal \cite{CohnConwayElkiesKumar2007,deLaatLeijenhorstDeMuinckKeizer2024}.

\item \textbf{Voronoi cell and honeycomb.}
For the checkerboard lattice
\[
D_4=\{x\in\ZZ^4:x_1+x_2+x_3+x_4\equiv0\pmod2\},
\]
the minimal vectors are the \(D_4\) roots.  The perpendicular bisectors to these minimal vectors bound the Voronoi region, and in the standard scaling that region is the regular \(24\)-cell.  Translating it by the lattice gives the \(24\)-cell honeycomb \(\{3,4,3,3\}\), with one adjacent Voronoi cell across each of the \(24\) octahedral facets.

\item \textbf{Hurwitz order.}
The Hurwitz order consists of the Lipschitz integer quaternions together with the half-integer coset.  Its norm-one elements are exactly the \(24\) units
\[
\{\pm1,\pm i,\pm j,\pm k\}
\cup
\left\{\frac12(\pm1\pm i\pm j\pm k)\right\}.
\]
Since the norm is multiplicative and quaternionic conjugation is inversion on the unit sphere, these vertices are closed under multiplication and inverse; this is precisely the binary tetrahedral group.

\item \textbf{Dual cell and \(F_4\).}
In standard root coordinates the long-root shell of \(F_4\) is \(D_4\), while the short-root shell consists of the coordinate roots \(\{\pm e_i\}\) together with the half-sum vertices.  Geometrically this short-root shell is the rotated dual \(24\)-cell at the appropriate scale.  Thus the primal--dual union gives
\[
\Phi(F_4)=\Phi_{\rm long}(D_4)\sqcup\Phi_{\rm short},
\qquad
48=24+24.
\]

\item \textbf{Reye incidence.}
The antipodal quotient has \(12\) axes.  The \(16\) great regular hexagons are the rank-two \(A_2\) subsystems of \(D_4\).  Each such plane contains three antipodal axes and each axis lies in four such planes, giving
\[
12\cdot4=16\cdot3=48.
\]
This proves the displayed \(12_4\,16_3\) incidence structure.

\item \textbf{Design hierarchy.}
The degree-through-five cubature statement is the spherical-design statement already verified in Theorem~\ref{p1:thm:s3one-structural-invariants}.  The invariant-theoretic side is the fundamental-degree list \(2,6,8,12\) for \(W(F_4)\): the degree-two invariant is radial and constant on \(S^3\), so degree six is the first possible nonconstant full-cell invariant.  The explicit sixth-moment discrepancy in Theorem~\ref{p1:thm:s3one-structural-invariants} shows that this possible obstruction actually occurs.
\end{proofsteps}
\end{proof}


\section{Finite-cell corollaries and summary}
\label{p1:sec:finite-cell-corollaries-summary}

\subsection{Mathematical finite-cell corollaries}
\label{p1:sec:math-finite-cell-corollaries}

This section gives compact labels for finite geometry, finite representation theory, and finite trace identities that follow from the level-\(1\) construction above.

\paragraph{Supplementary degree-six calculation.}
The Hurwitz support is exact through degree five and has first scalar moment
defect \(3/8\) in degree six.  The explicit invariant harmonic sextic and its
orbit-sum expansion are recorded in
Proposition~\ref{p1:prop:math-hurwitz-cubature-six-defect} of
Appendix~\ref{p1:app:hurwitz-supplementary}; they are not needed to follow the
main finite-cell proof chain.

\begin{proposition}[Bose--Mesner shell algebra of the \(24\)-cell]
\label{p1:prop:math-twentyfourcell-bose-mesner}
Let \(X=2T\) be the Hurwitz \(24\)-cell support and define relation matrices \(I,A_+,A_0,A_-,A_\pi\) by the inner products
\[
  1,\quad \frac12,\quad 0,\quad -\frac12,\quad -1.
\]
Then
\[
  \mathcal A_{24}:=\operatorname{span}\{I,A_+,A_0,A_-,A_\pi\}
\]
is the Bose--Mesner algebra of the antipodal shell relations.  Its primitive idempotent ranks are
\[
  1,\;4,\;9,\;8,\;2,
\]
and the nearest-neighbor matrix \(A_+\) has eigenvalues
\[
  8,\;4,\;0,\;-2,\;-4
\]
with those multiplicities.  Hence the normalized nearest-neighbor operator \(T=A_+/8\) has spectrum
\[
  1,\quad \frac12,\quad 0,\quad -\frac14,\quad -\frac12.
\]
In particular, the nine-dimensional primitive sector is the zero eigenspace of \(T\).
\end{proposition}

\begin{proof}
The shell valencies at any Hurwitz vertex are \(1,8,6,8,1\), corresponding to the five inner products listed above.  The association scheme of the regular \(24\)-cell is commutative with five primitive idempotents.  Direct multiplication of the relation matrices, equivalently diagonalization of the \(24\)-cell distance scheme, gives the displayed eigenvalue table for \(A_+\).  Division by the valency \(8\) gives the spectrum of \(T\).
The full joint-eigenvalue table and the trace computation of the five ranks
are displayed in the proof of
Theorem~\ref{p1:thm:s3one-numerical-data}\textup{(iv)}.
\end{proof}

\paragraph{Supplementary projector formulas.}
The rank-\(9\) kernel projector and its complementary rank-\(15\) projector
are explicit polynomials in the Bose--Mesner relations.  Their formulas and
direct idempotence check are recorded in
Corollary~\ref{p1:cor:math-twentyfourcell-explicit-projectors} of
Appendix~\ref{p1:app:hurwitz-supplementary}; the main chain uses only the
rank split and spectral subspaces already proved above.

\begin{proposition}[Natural left-\(2T\) no-go for a full Hurwitz carrier]
\label{p1:prop:h1-e15-left-2t-obstruction}
Restrict the \(SU(2)_L\)-factor of the declared \(\operatorname{Spin}(4)\)-action on
\(\mathcal H_1\) to \(2T\), and let \(2T\) act on
\(\CC^{2T}\) by the left regular action.  Then
\[
    \dim_\CC\mathcal H_1^{2T}=6,
    \qquad
    \dim_\CC(\CC^{2T})^{2T}=1,
    \qquad
    \dim_\CC(\mathcal E_{15}^{\rm BM})^{2T}=1.
\]
For every left-\(2T\)-equivariant linear map
\(F:\mathcal H_1\to\CC^{2T}\),
\[
    \dim\ker F\ge5.
\]
Consequently no such map is injective.  In particular, there is no equivariant
isomorphism between \(\mathcal H_1\) and
\(\mathcal E_{15}^{\rm BM}\) for these natural left actions.
This does not exclude arbitrary non-equivariant maps or a comparison made
after declaring different actions.
\end{proposition}

\noindent
The fixed-space calculation and rank--nullity proof are in
Appendix~\ref{p1:app:left-2t-carrier-obstruction}.  The conclusion limits the
scalar vertex-function target with the stated left action; it does not
exclude the vector-valued sampling construction below.

\begin{proposition}[Scalar evaluation bridge and exact Casimir intertwiner]
\label{p1:prop:h1-scalar-hurwitz-evaluation-bridge}
Use the scalar Peter--Weyl realization of
\[
 \mathcal H_1^{\rm sc}
 :=(0,0)\oplus(\tfrac12,\tfrac12)
 \cong\operatorname{span}_{\CC}\{1,u_1,u_2,u_3,u_4\},
 \qquad u_a:=x_a/\Rmom,
\]
and evaluate the normalized coordinate functions on the unit Hurwitz shell:
\[
 \operatorname{ev}_{2T}:\mathcal H_1^{\rm sc}\longrightarrow\CC^{2T},
 \qquad
 f\longmapsto\bigl(f(q)\bigr)_{q\in2T}.
\]
For the normalized round \(L^2\) inner product and normalized counting
inner product
\[
 24^{-1}\sum_{q\in2T}\overline{F(q)}G(q),
\]
this is a left-\(2T\)-equivariant isometric embedding.  Its image is exactly
\[
 \ker(A_+-8I)\oplus\ker(A_+-4I)
 \subset\mathcal E_{15}^{\rm BM}.
\]
On this five-dimensional domain it satisfies the operator identities
\[
 \operatorname{ev}_{2T}K_1
 =\frac38(8I-A_+)\operatorname{ev}_{2T},
\]
and, for the nonnegative round scalar Laplacian
\(\Delta=2K_1/\Rmom^2\),
\[
 \operatorname{ev}_{2T}\Delta
 =\frac{3}{4\Rmom^2}(8I-A_+)\operatorname{ev}_{2T}.
\]
Thus evaluation preserves the constant/coordinate eigenspace split and the
first two cumulative ranks \(1,5\) after the displayed normalization.  It does
not define the later rank-\(9\) and rank-\(15\) stages and does not extend to a
full-carrier injection under the natural left action.
\end{proposition}

\begin{proof}
The Hurwitz shell is a spherical \(5\)-design.  Hence normalized counting on
\(2T\) agrees with normalized round integration on products of functions of
degree at most one, proving the isometry.  Equivariance follows from
restriction of scalar functions under left translation.

The constant function satisfies \(A_+1=8\,1\).  At the base vertex \(1\), the
sum of the eight neighboring Hurwitz units is \(4\cdot1\), because their real
parts sum to \(4\) and their sign coordinates cancel.  Left multiplication
transports this identity to every vertex, so each coordinate column satisfies
\(A_+u_a=4u_a\).  The corresponding eigenspaces have dimensions \(1\) and
\(4\), proving the image statement and injectivity independently.

Finally, \(K_1\) has eigenvalue \(0\) on constants and \(3/2\) on
\((1/2,1/2)\), while \(8I-A_+\) has eigenvalues \(0\) and \(4\) on their
images.  This proves the first intertwining identity; multiplication by
\(2/\Rmom^2\) gives the Laplacian identity.  The full-carrier limitation is
Proposition~\ref{p1:prop:h1-e15-left-2t-obstruction}.
\end{proof}

\begin{proposition}[Protected Clifford--Fierz basis]
\label{p1:prop:math-protected-clifford-fierz}
On the \(S^3_1\) cell, let
\[
  \mathcal K_{\mathsf B}=\left(\frac12,0\right)\oplus\left(0,\frac12\right)
\]
be the protected four-dimensional complex block with its fixed
\(\operatorname{Spin}(4)\) action.  Choose an equivariant Hermitian coordinate completion
\(\mathsf C\) with \(\rho_{\mathsf C}>0\), and set
\[
 \Gamma_a:=\frac{2}{\rho_{\mathsf C}}
 \widehat X_{a,\mathcal K}^{\mathsf B},\qquad a=1,\ldots,4.
\]
Thus the \(\Gamma_a\) are Hermitian, transform as the vector representation
for this same \(\operatorname{Spin}(4)\) action, and satisfy
\[
  \{\Gamma_a,\Gamma_b\}=2\delta_{ab}1,
\]
by Proposition~\ref{p1:prop:general-protected-clifford-relation}.
With \(\Gamma_5=\Gamma_1\Gamma_2\Gamma_3\Gamma_4\) and
\(\Gamma_{ab}:=\Gamma_a\Gamma_b\) for \(a<b\),
\[
  \operatorname{End}(\mathcal K_{\mathsf B})
  =\langle 1,\Gamma_a,\Gamma_{ab},\Gamma_5\Gamma_a,\Gamma_5\rangle,
\]
with grade dimensions
\[
  1+4+6+4+1=16.
\]
Moreover, for complex-linear equivariant endomorphisms,
\[
  \operatorname{End}_{\operatorname{Spin}(4)}^{\CC}(\mathcal K_{\mathsf B})=\CC 1\oplus\CC\Gamma_5.
\]
\end{proposition}

\begin{proof}
The complex Clifford algebra generated by the four gamma matrices is
\(M_4(\CC)\), with exterior-grade dimensions \(1,4,6,4,1\).  Its unital
action on the four-dimensional block is faithful, so the displayed sixteen
elements form a basis of \(\End(\mathcal K_{\mathsf B})\).  Hermiticity and
the Clifford relations give their Hilbert--Schmidt trace orthogonality.  The
grade-three elements are Hodge-dual to \(\Gamma_5\Gamma_a\).

Covariance for the fixed oriented vector action makes \(\Gamma_5\)
\(\operatorname{Spin}(4)\)-invariant.  It satisfies \(\Gamma_5^2=1\) and
\(\operatorname{Tr}\Gamma_5=0\), so its two eigenspaces have dimension two.
Schur's lemma on the inequivalent summands \((1/2,0)\) and \((0,1/2)\)
identifies them with the chiral blocks, in the order fixed by the volume
convention.  Hence the commutant is \(\CC1\oplus\CC\Gamma_5\).
The fixed-action conclusion uses covariance as well as the Clifford
relations.
\end{proof}

\begin{proposition}[Cyclic trimer Weyl algebra and \(A_2\) defect plane]
\label{p1:prop:math-trimer-weyl-a2}
Let \(E=\{e_0,e_1,e_2\}\) be an oriented trimer and let \(P_j=|e_j\rangle\langle e_j|\), \(j\in\ZZ_3\).  For \(\omega=e^{2\pi i/3}\), define
\[
  Ze_j=\omega^j e_j,\qquad Xe_j=e_{j+1}.
\]
Then \(X^3=Z^3=1\), \(ZX=\omega XZ\), and the nine operators \(X^aZ^b\), \(a,b\in\ZZ_3\), form a trace-orthogonal basis of \(M_3(\CC)\).  The centered defects \(D_j=P_j-I/3\) satisfy
\[
  \operatorname{tr}(D_j^2)=\frac23,
  \qquad
  \operatorname{tr}(D_iD_j)=-\frac13\quad(i\ne j).
\]
For \(\alpha_1=P_0-P_1\), \(\alpha_2=P_1-P_2\), one obtains the \(A_2\) Cartan matrix
\[
  \bigl(\operatorname{tr}(\alpha_i\alpha_j)\bigr)_{i,j=1}^2
  =
  \begin{pmatrix}2&-1\\-1&2\end{pmatrix}.
\]
\end{proposition}

\begin{proof}
The relation \(ZX=\omega XZ\) follows by applying both sides to \(e_j\).  The
standard trace calculation gives
\[
  \operatorname{tr}\bigl((X^aZ^b)^*X^cZ^d\bigr)
  =
  3\delta_{ac}\delta_{bd},
\]
so the nine Weyl operators form a basis of \(M_3(\CC)\).

For the defect identities, use
\[
  D_0
  =
  \operatorname{diag}
  \left(
    \frac{2}{3},
    -\frac{1}{3},
    -\frac{1}{3}
  \right)
\]
and cyclic permutation.  Finally, the root identities follow by expanding
\(P_iP_j=\delta_{ij}P_i\).
\end{proof}

\begin{proposition}[Projected coordinate-coupling graph on \(\mathcal H_1\)]
\label{p1:prop:math-h1-coordinate-coupling-graph}
Let
\[
\mathcal H_1=(0,0)\oplus\left(\frac12,\frac12\right)\oplus(1,0)\oplus(0,1)\oplus\left(\frac12,0\right)\oplus\left(0,\frac12\right),
\]
and let the coordinate representation be \(X=(1/2,1/2)\).  Tensoring by \(X\) using the fixed Clebsch--Gordan carrier convention, and then projecting back to \(\mathcal H_1\), has exactly the retained channels
\[
  (0,0)\leftrightarrow\left(\frac12,\frac12\right),
\]
\[
  \left(\frac12,\frac12\right)\leftrightarrow(1,0),
  \qquad
  \left(\frac12,\frac12\right)\leftrightarrow(0,1),
\]
and
\[
  \left(\frac12,0\right)\leftrightarrow\left(0,\frac12\right).
\]
No other coordinate-coupling channel remains inside \(\mathcal H_1\).
\end{proposition}

\begin{proof}
For \(\operatorname{Spin}(4)\cong SU(2)_L\times SU(2)_R\), tensoring \((j_L,j_R)\) with \((1/2,1/2)\) changes each spin by \(\pm1/2\), omitting negative spins.  Applying this Clebsch--Gordan rule to each summand of \(\mathcal H_1\) using the fixed coordinate-sector tensoring operation gives exactly the displayed retained terms.  The additional terms \((1,1)\), \((3/2,1/2)\), \((1/2,3/2)\), \((1,1/2)\), and \((1/2,1)\) lie outside \(\mathcal H_1\) and are therefore projected out.
\end{proof}

\paragraph{Supplementary invariant and triality calculations.}
For the full unmarked \(24\)-cell, degree six is the first nonconstant
spherical invariant degree, and an explicit \(D_4\) triality frame realizes
the quotient graph \(K_{4,4,4}\) and \(24=8+8+8\).  The detailed
invariant-ring and coordinate proofs are recorded in
Propositions~\ref{p1:prop:math-f4-sextic-obstruction}
and~\ref{p1:prop:math-triality-octets-k444} of
Appendix~\ref{p1:app:hurwitz-supplementary}.


\section{Collected \texorpdfstring{\(m=1\)}{m=1} finite package}
\label{p1:sec:s3one-summary}

\begin{corollary}[Typed \(S^3_1\) finite package]
\label{p1:cor:s3one-finite-package}
The conclusions split according to their actual hypotheses.

\smallskip
\noindent\emph{Module-derived clause.}
For every level-one completion
\(\mathsf B(\mathsf C)=(\mathsf B_{\rm skel},\mathsf C)\),
\[
\mathcal H_{\mathsf B}=\mathcal H_1,
\qquad
\dim\mathcal H_{\mathsf B}=15,
\qquad
\mathcal K_{\mathsf B}=(\tfrac12,0)\oplus(0,\tfrac12),
\]
and the finite coordinate operators obey
\[
\left\{
\widehat X_{a,\mathcal K}^{\mathsf B},
\widehat X_{b,\mathcal K}^{\mathsf B}
\right\}
=
\frac{\rho_{\mathsf C}^2}{2}\delta_{ab}\mathbf1_{\mathcal K},
\qquad
\rho_{\mathsf C}:=2\varrho_{\rm sp}\ge0.
\]
It also has
\[
15=1+4+10,
\qquad
\operatorname{End}_{\operatorname{Spin}(4)}^{\CC}(\mathcal K_{\mathsf B})
=\CC1\oplus\CC\Gamma_{\mathcal K},
\]
with the displayed unnormalized protected relation.  These statements require
no finite-orbit support choice.  If \(\rho_{\mathsf C}=0\), then every
\(\widehat X_{a,\mathcal K}^{\mathsf B}\) vanishes.  If
\(\rho_{\mathsf C}>0\), the normalized matrices
\(\Gamma_a^{\mathsf C}=2\widehat X_{a,\mathcal K}^{\mathsf B}/\rho_{\mathsf C}\)
and the associated nonzero spectra, Hilbert--Schmidt frame, and
Clifford--Fierz package apply.

\smallskip
\noindent\emph{Support-derived clause.}
Only after a finite-cell representative over \(\mathsf{Supp}_{2T}\) is also
specified do the independent identities
\[
 15=7+8,
 \qquad
 \dim_\CC\mathcal O_9^{\rm tri}=9
\]
and the trimer/Hurwitz interfaces apply.
For a rooted oriented Hurwitz-link diagonal, the point maps
\(L_\pm\) into \(V\) and \(B_\pm\) into \(\Lambda^2V\)
also have the exact overlaps
\(K_V=C_\triangle^{-1}\) and
\(K_{\Lambda^2}=C_\triangle/2\), with
\(H_{\Lambda^2}=H_V^{-1}/4\)
(Theorem~\ref{p1:thm:rooted-trimer-vector-overlap} and
Corollary~\ref{p1:cor:rooted-trimer-bivector-overlap}).
On any such calculation
representative, the zero
extensions \(\iota_\pm\), normalized spectral embeddings
\(U_{9,\pm},U_{15,\pm}\), and operator corners \(\jmath_\pm\) give typed
interfaces internal to \(\CC^{2T}\).  Independently, the scalar subcarrier
\(\mathcal H_1^{\rm sc}\) evaluates isometrically into
\(\mathcal E_{15}^{\rm BM}\) and intertwines \(K_1\) with
\(\frac38(8I-A_+)\).  Neither statement supplies an injective natural
left-\(2T\)-equivariant map from all of \(\mathcal H_1\).
Its remaining numerical, structural, and classical-completion invariants are
those of
Theorems~\ref{p1:thm:s3one-numerical-data},
\ref{p1:thm:s3one-structural-invariants}, and
\ref{p1:thm:s3one-classical-completions}.
\end{corollary}

\begin{proof}
The module-derived clause follows from the level-one module and
Proposition~\ref{p1:prop:general-protected-clifford-relation} and is uniform
over the completion family; the proposition also proves the zero/positive
stratum split.  The
support-derived clause follows separately from
Theorem~\ref{p1:thm:s3one-finite-cell-structure},
Theorem~\ref{p1:thm:s3one-trimer-bose-mesner-interface}, and
Proposition~\ref{p1:prop:h1-scalar-hurwitz-evaluation-bridge}.
The rooted overlap exports follow from
Theorem~\ref{p1:thm:rooted-trimer-vector-overlap} and
Corollary~\ref{p1:cor:rooted-trimer-bivector-overlap}. The typed
hypothesis split prevents the latter conclusions from being inferred from
\(m=1\) alone.
\end{proof}


\section{Round-symbol ancestry and the level-one coordinate exports}
\label{p1:sec:round-symbol-ancestry}
\label{p1:sec:metric-cone-clifford-ancestry}

The raw coordinate atlas of
Corollary~\ref{p1:cor:level-one-coordinate-completion-moduli} classifies all
compatible equivariant finite quadruples, but does not explain their geometric
ancestry.  This section exports the two primitive round-symbol descendants:
the intrinsic fixed-radius Hodge/spin class \(\mathsf C_D\) and the ambient
metric-cone class \(\mathsf C_{\rm cone}\).  It then proves exhaustion relative
to the explicitly specified primitive ancestry grammar.

The ambient metric-cone parent used for comparison is
\[
 \mathsf C(S^3_{\Rmom})=((0,\infty)\times S^3_{\Rmom},
 dr^2+r^2\Rmom^{-2}g_{\Rmom}).
\]
Since $\Rmom^{-2}g_{\Rmom}$ has sectional curvature one, this cone is flat.  Its
completion is Euclidean four-space, uniquely up to the already acting
orientation-preserving orthogonal group.

On the exterior-plus-spinor Clifford bundle set
\[
 c(\xi)=\bigl(\varepsilon(\xi)+\iota_{\xi^\sharp}\bigr)\oplus c_S(\xi).
\]
For any first-order operator $\mathscr D$ with principal symbol
$\sigma_1(\mathscr D)(\xi)=-i c(\xi)$ and every scalar $f$,
\begin{equation}
 \label{p1:eq:principal-symbol-commutator}
 i[\mathscr D,M_f]=c(df).
\end{equation}
Thus every zero-order perturbation of $\mathscr D$ drops out of the coordinate
commutator.

\begin{proposition}[Ambient radial component and the intrinsic tangential quotient]
\label{p1:prop:cone-tangential-symbol-split}
Let \(i_{\Rmom}:S^3_{\Rmom}\hookrightarrow\mathbb R^4\) be the fixed-radius
inclusion and let \(e^a\) be an ambient orthonormal constant covector.  Along
the hypersurface,
\[
 e^a=d_{S^3}x_a+\frac{x_a}{\Rmom}\,dr,
 \qquad i_{\Rmom}^*e^a=d_{S^3}x_a.
\]
Consequently the ambient cone symbol \(c(e^a)\) and the intrinsic tangential
symbol \(c(d_{S^3}x_a)\) are different symbol descendants: the former retains
the radial Clifford component and the latter quotients it out.  Neither
ancestry map canonically identifies its finite coordinate class with the
other.
\end{proposition}

\begin{proof}
Write an ambient point as \(y=(r/\Rmom)x\), with
\(x\in S^3_{\Rmom}\).  Differentiating
\(y_a=(r/\Rmom)x_a\) at \(r=\Rmom\) gives the displayed decomposition.
Pullback kills \(dr\), proving the tangential formula.  Since
\(\sum_a x_a^2=\Rmom^2\), the radial coefficients cannot all vanish; Clifford
multiplication by the radial term is therefore nonzero on the cone fibre.
The two functorial symbol constructions have different targets, so an
identification would require an additional comparison theorem rather than a
sector rephasing.
\end{proof}

\subsection{Canonical intrinsic fixed-radius export}

\begin{theorem}[Canonical retained smooth model of the level-one carrier]
\label{p1:thm:canonical-retained-smooth-model}
\leavevmode\par
Let \(V_R:=\operatorname{Vol}(S^3(\Rmom))\), let
\(x_a\), \(a=1,\ldots,4\), be the embedding-coordinate module, and put
\[
 e_0:=V_R^{-1/2},\qquad
 u_a:=\frac{2}{\sqrt{3V_R}}\,dx_a,\qquad
 w_{ab}:=\sqrt{\frac{2}{V_R}}\,dx_a\wedge dx_b
 \quad(a<b).
\]
With \(S_+\) and \(S_-\) the lowest positive and negative Dirac eigenspaces,
define
\[
 \mathcal L_R
 :=\operatorname{span}\{e_0\}
 \oplus\operatorname{span}\{u_a\}_{a=1}^4
 \oplus\operatorname{span}\{w_{ab}\}_{a<b}
 \oplus S_+\oplus S_-.
\]
The displayed modes are orthonormal, \(\dim_{\CC}\mathcal L_R=15\), and
there is a natural unitary \(\operatorname{Spin}(4)\)-equivariant equivalence class
\[
 J_R:\mathcal H_1\xrightarrow{\ \simeq\ }\mathcal L_R.
\]
Any representative of this class sends the scalar, vector, bivector, and
half-spin sectors to the displayed summands.  It is unique up to the already
unmarked phases on the six multiplicity-one irreducibles.  Thus
``canonical'' here means this equivariant unitary class, not one literally
distinguished phase convention or embedded basis.
\end{theorem}

\begin{proof}
The round identities
\[
 \int_{S^3}\!\langle dx_a,dx_b\rangle\,d\vol
 =\frac{3V_R}{4}\delta_{ab}
\]
and
\[
 \int_{S^3}\!\langle dx_a\wedge dx_b,dx_c\wedge dx_d\rangle\,d\vol
 =\frac{V_R}{2}
  (\delta_{ac}\delta_{bd}-\delta_{ad}\delta_{bc})
\]
give the normalizations.  The three form summands transform as
\((0,0)\), \((\tfrac12,\tfrac12)\), and
\((1,0)\oplus(0,1)\); adjoining the lowest spinors gives exactly the six
summands of Corollary~\ref{p1:cor:H1-structure}.  Schur's lemma gives the
phase-qualified uniqueness.  Reversing orientation exchanges the two
bivector labels and changes no unitary class.
\end{proof}

\begin{corollary}[Normalized constant density is Haar probability]
\label{p1:cor:normalized-constant-haar-density}
The normalized constant in
Theorem~\ref{p1:thm:canonical-retained-smooth-model} satisfies
\[
 |e_0|^2d\vol_{S^3(\Rmom)}
 =\frac{d\vol_{S^3(\Rmom)}}{\operatorname{Vol}(S^3(\Rmom))}
 =\mu_H,
\]
where \(\mu_H\) is normalized Haar measure on the round group
\(S^3\cong\SU(2)\).
\end{corollary}

\begin{proof}
The first equality is the definition \(e_0=V_R^{-1/2}\).  The round metric
is bi-invariant, so its normalized Riemannian volume is the unique normalized
Haar measure.
\end{proof}

\begin{theorem}[Intrinsic tangential round-symbol coordinate class]
\label{p1:thm:intrinsic-tangential-coordinate-class}
Choose any representative \(J_R\) of the unitary class in
Theorem~\ref{p1:thm:canonical-retained-smooth-model}.  On forms plus spinors
over the fixed-radius sphere put
\[
 \mathscr D_{\rm LC}:=i(d-d^\dagger)\oplus D_{S^3}
\]
and define
\[
 X_a^D:=\frac{\Rmom}{\sqrt3}\,
 J_R^*i[\mathscr D_{\rm LC},M_{x_a}]J_R,
 \qquad a=1,\ldots,4.
\]
The phase class \(\mathsf C_D:=[X^D]\) is independent of the representative
\(J_R\), of the overall Clifford-sign convention, and of every zero-order
term added to \(\mathscr D_{\rm LC}\).  It is the unique class on the
intrinsic fixed-radius metric-symbol branch after the normalized
scalar-to-vector edge is matched.  In the sector order
\[
 (0,0),\quad(\tfrac12,\tfrac12),\quad(1,0),\quad(0,1),
 \quad(\tfrac12,0),\quad(0,\tfrac12),
\]
its invariant quadratic fingerprint is
\[
 \boxed{
 Q_D:=\sum_{a=1}^4(X_a^D)^2
 =\left(
 \Rmom^2,\ \frac{11\Rmom^2}{12},\
 \frac{4\Rmom^2}{9},\ \frac{4\Rmom^2}{9},\
 \frac{3\Rmom^2}{4},\ \frac{3\Rmom^2}{4}
 \right).}
\]
With the raw exterior-edge atlas convention
\[
 Q|_{(1,0)}=2r_{1,0}^2I,\qquad
 Q|_{(0,1)}=2r_{0,1}^2I,\qquad
 Q|_{S_+\oplus S_-}=4\varrho_{\rm sp}^2I,
\]
the reduced tuple is
\[
 \boxed{
 (r_{1,0},r_{0,1},\varrho_{\rm sp})_D
 =\left(\frac{\sqrt2\Rmom}{3},
        \frac{\sqrt2\Rmom}{3},
        \frac{\sqrt3\Rmom}{4}\right).}
\]
In particular every allowed forest edge is nonzero.
\end{theorem}

\begin{proof}
On forms, the formal-adjoint identities
\[
 [d,M_f]=\varepsilon(df),\qquad
 [d^\dagger,M_f]=-\iota_{\nabla f}
\]
give, for the displayed differential operator,
\[
 i[i(d-d^\dagger),M_f]
 =-\bigl(\varepsilon(df)+\iota_{\nabla f}\bigr).
\]
Thus the exterior summand has the opposite symbol sign from the hypothesis
of Equation~\eqref{p1:eq:principal-symbol-commutator}.  Let
\(\mathcal U\) act by \((-1)^p\) on \(p\)-forms and by the identity on
spinors.  It preserves the retained smooth model, commutes with
\(\operatorname{Spin}(4)\), and conjugates the odd exterior Clifford action to its
negative.  Replacing \(J_R\) by \(\mathcal UJ_R\) is an allowed sector-phase
change.  It gives the positive exterior edge representative used below,
without changing \(\mathsf C_D\), its quadratic fingerprint, or the retained
square.  Zero-order terms commute with scalar multipliers on both summands.
The form coordinate matrices in this phase choice are therefore computed
from the positive Clifford action in the original normalized modes:
\[
 c(dx_a)e_0=\frac{\sqrt3}{2}u_a,
 \qquad
 P_{\mathcal L_R}c(dx_a)u_b
 =\frac{\sqrt3}{2}\delta_{ab}e_0
  +\sqrt{\frac23}\,w_{ab},
\]
where \(w_{ba}=-w_{ab}\).  Multiplication by \(\Rmom/\sqrt3\) gives scalar
edge \(\Rmom/2\) and raw vector--bivector edge
\(b_D=\sqrt2\Rmom/3\).  Summing outgoing squared edge norms gives
\[
 Q_D|_{(0,0)}=\Rmom^2,
 \quad
 Q_D|_{(\frac12,\frac12)}
 =\frac{\Rmom^2}{4}+3\frac{2\Rmom^2}{9}
 =\frac{11\Rmom^2}{12},
\]
and
\[
 Q_D|_{(1,0)}=Q_D|_{(0,1)}=2b_D^2
 =\frac{4\Rmom^2}{9}.
\]

For a normalized \(\psi_+\in S_+\), write
\[
 A:=\sum_a\|P_{S_-}M_{x_a}\psi_+\|^2,
 \qquad
 B:=\sum_a\|P_{(1,1/2)}M_{x_a}\psi_+\|^2.
\]
The sphere relation and the two Dirac gaps \(3/\Rmom\) and \(1/\Rmom\)
give
\[
 A+B=\Rmom^2,
 \qquad
 \frac9{\Rmom^2}A+\frac1{\Rmom^2}B=3,
\]
because \(\sum_a|dx_a|^2=3\).  Hence
\(A=\Rmom^2/4\), and the retained unscaled commutator radius is
\(9/4\).  The common factor \(\Rmom^2/3\) gives
\(Q_D|_{S_\pm}=3\Rmom^2/4\).  The atlas definitions yield the boxed tuple.
All edge spaces have multiplicity one, so the remaining choices are exactly
the already quotiented sector phases.
\end{proof}

\begin{remark}[Raw and coisometric edge normalization]
\label{p1:rem:intrinsic-edge-normalization-conversion}
For the exterior map \(T(v\otimes u)=v\wedge u\), one has
\(TT^*=2I_{\Lambda^2}\).  Relative to this raw map, the intrinsic edge is
\(\sqrt2\Rmom/3\).  Relative to the coisometric intertwiner
\(T_0=2^{-1/2}T\), the same operator edge is written as \(2\Rmom/3\).
The invariant fingerprint and the defining equations for
\(r_{1,0},r_{0,1}\) above are unchanged. Thus no geometric
branch depends on which convention is called ``normalized.''
\end{remark}

\begin{theorem}[Additional exact eight-dimensional Clifford strata]
\label{p1:thm:v214-exact-eight-clifford-strata}
The raw level-one completion atlas contains the two geometric points
\[
 (r_{1,0},r_{0,1})=(\Rmom/\sqrt2,0),
 \qquad (r_{1,0},r_{0,1})=(0,\Rmom/\sqrt2).
\]
At the first point the connected geometric block
\[
 K_8^+=\Lambda^0_{\mathbb C}\oplus V_{\mathbb C}
                         \oplus\Lambda^2_+V_{\mathbb C}
 \cong(S_+\oplus S_-)\otimes S_+
\]
has exact coordinates \(X(v)=(\Rmom/2)\gamma_+(v)\), where
\[
 \gamma_+(v)(\lambda,w,\omega)=
 \bigl(\langle v,w\rangle,\,
       \lambda v+\sqrt2\,\iota_v\omega,\,
       \sqrt2\,P_+(v\wedge w)\bigr).
\]
Here \(v\) is real, \(P_+\) is the orthogonal self-dual projection, and
the formulas extend complex linearly. They satisfy
\[
 \gamma_+(v)^*=\gamma_+(v),\qquad
 \{\gamma_+(v),\gamma_+(w)\}=2\langle v,w\rangle I_{K_8^+}.
\]
The other geometric three-space has zero coordinate action. Replacing
\(P_+\) by \(P_-\) gives the exchanged point and \(K_8^-\).
The protected spinor magnitude is independent of these choices.
This constructs two strata; it does not classify all exact Clifford
subblocks of every completion.
\end{theorem}
\begin{proof}
Exterior multiplication and contraction are adjoints, so the displayed
operator is Hermitian. For a unit \(v\), the Hodge identities are
\[
 2\iota_vP_\pm(v\wedge w)=w-\langle v,w\rangle v,\qquad
 2P_\pm(v\wedge\iota_v\omega)=\omega
 \quad(\omega\in\Lambda^2_\pm V).
\]
The first follows from
\(\iota_v*(v\wedge w)=0\); the second follows by splitting a
self-dual or anti-self-dual form into its \(v\)-containing part and
its Hodge partner. Substitution gives \(\gamma_\pm(v)^2=I\).
Homogeneity and polarization prove the Clifford relation for all real
\(v,w\). The scalar edge of \(X\) is \(\Rmom/2\), and its raw
vector--bivector edge is \(\sqrt2\,\Rmom/2\), proving the atlas values.
Independently, outgoing-path traces give the general geometric fingerprint
\[
 Q_0=\Rmom^2,\qquad
 Q_V=\Rmom^2/4+\tfrac32(r_{1,0}^2+r_{0,1}^2),\qquad
 Q_{\Lambda^2_\pm}=2r_\pm^2.
\]
Indeed each bivector sector has dimension three, so its total return
trace \(6r_\pm^2\) contributes \(3r_\pm^2/2\) to the
four-dimensional vector sector. At either displayed point \(Q=\Rmom^2I\)
on the connected eight-space and vanishes on the inactive triplet.
The representation identity follows from
\(S_+\otimes S_+=(0,0)\oplus(1,0)\) and
\(S_-\otimes S_+=(1/2,1/2)\).
\end{proof}

\begin{corollary}[Canonical retained Levi--Civita square and scalar projector]
\label{p1:cor:canonical-retained-dirac-square}
For
\[
 K_{\rm LC}:=J_R^*\mathscr D_{\rm LC}^{\,2}J_R,
\]
the sectorwise eigenvalues in the order used above are
\[
 \boxed{
 K_{\rm LC}\sim
 \left(0,\frac3{\Rmom^2},\frac4{\Rmom^2},
 \frac4{\Rmom^2},\frac9{4\Rmom^2},\frac9{4\Rmom^2}\right).}
\]
Thus the zero eigenspace is exactly the normalized scalar line and
\[
 \Pi_{\rm sc}:=\mathbf1_{\{0\}}(K_{\rm LC})
\]
is its canonical geometric spectral projector.  This statement is about the
explicit Levi--Civita representative; the principal symbol alone does not
determine the square.
\end{corollary}

\begin{proof}
On forms, \(\mathscr D_{\rm LC}^{\,2}\) is the Hodge Laplacian.  Constants
have eigenvalue zero, exact first modes have the scalar \(k=1\) value
\(3/\Rmom^2\), and the lowest coexact/Killing modes have value
\(4/\Rmom^2\) by Theorem~\ref{p1:thm:hodge-one-form-spectrum}; Hodge duality
gives the same value on the retained two-forms.  The lowest spinor values are
\(\pm3/(2\Rmom)\) by
Theorem~\ref{p1:thm:representation-harmonic-package}.  Squaring gives the
displayed list, whose only zero is the scalar entry.
\end{proof}




\begin{theorem}[First-order closure and the exact compression leakage]
\label{p1:thm:v214-first-order-retained-closure}
Write \(R=\Rmom\). Let \(D=\mathscr D_{\rm LC}\), let \(P\) project onto
\(\mathcal L_R\), and put \(D_F=J_R^*DJ_R\).
The retained model is invariant under \(D^2\), but not under \(D\).
In the sector order of
Corollary~\ref{p1:cor:canonical-retained-dirac-square},
\[
 D_F^2\sim R^{-2}(0,0,0,0,9/4,9/4),\qquad
 K_{\rm LC}-D_F^2
 =J_R^*D(1-P)DJ_R
 \sim R^{-2}(0,3,4,4,0,0).
\]
Thus \(D_F^2\) has an eleven-dimensional kernel, whereas
\(K_{\rm LC}=J_R^*D^2J_R\) has only the constant kernel.
The minimal \(D\)-invariant smooth subspace containing \(\mathcal L_R\) is
\[
 \widetilde{\mathcal L}_R=\mathcal L_R+D\mathcal L_R,\qquad
 \dim_{\mathbb C}\widetilde{\mathcal L}_R=25.
\]
It adds the four scalar first harmonics and the six Killing one-forms.
As a \(\operatorname{Spin}(4)\)-module it is
\[
 \widetilde{\mathcal L}_R\cong
 \mathbf1\oplus2(\tfrac12,\tfrac12)
 \oplus2(1,0)\oplus2(0,1)
 \oplus(\tfrac12,0)\oplus(0,\tfrac12).
\]
For its orthogonal projection \(\widetilde P\), the finite self-adjoint
operator \(\widetilde D=D|_{\widetilde{\mathcal L}_R}\) has spectrum
\[
 \operatorname{spec}\widetilde D
 =0^{(1)}\oplus(\sqrt3/R)^{(4)}
 \oplus(-\sqrt3/R)^{(4)}
 \oplus(2/R)^{(6)}\oplus(-2/R)^{(6)}
 \oplus(3/(2R))^{(2)}\oplus(-3/(2R))^{(2)}.
\]
For every smooth scalar \(f\), with
\(\widetilde M_f=\widetilde P M_f\widetilde P\) on this space,
\[
 [\widetilde D,\widetilde M_f]
 =\widetilde P[D,M_f]\widetilde P.
\]
Compressing this commutator back to \(\mathcal L_R\) recovers exactly
the intrinsic coordinate construction. No such equality with \(D_F\)
alone is asserted on the original fifteen-space.
\end{theorem}
\begin{proof}
The constant is killed by \(D\). The differential of a first scalar
harmonic lies in the retained exact one-form space, and its codifferential
returns to that scalar space. Consequently \(D\) exchanges the retained
four exact one-forms with four omitted scalar first harmonics.
The retained two-forms are closed and lie in the eigenvalue \(4/R^2\)
Hodge space. Their codifferentials are the six coexact Killing
one-forms; these are orthogonal to all exact one-forms. Their differentials
return to the retained two-forms. The spinor summands are already
\(D\)-invariant. Therefore \(PDP\) vanishes on all eleven form modes
and has the two stated lowest spinor eigenvalues on the spin block.
Inserting \(P+(1-P)\) between the two factors of \(D^2\) proves the
leakage identity and the displayed fingerprints.

The ten additional modes are linearly independent and orthogonal to
\(\mathcal L_R\). Since \(D^2\mathcal L_R\subseteq\mathcal L_R\),
the space \(\mathcal L_R+D\mathcal L_R\) is \(D\)-invariant.
Every invariant space containing \(\mathcal L_R\) must contain its
\(D\)-image, proving minimality. On each positive Hodge eigenspace pair,
\(D\) is self-adjoint, exchanges the two form degrees, and has square
\(\lambda I\). It therefore has eigenvalues \(\pm\sqrt\lambda\)
with equal multiplicities. This gives the spectrum and its unique zero.
The finite invariant space is reducing for the self-adjoint \(D\);
hence \([\widetilde P,D]=0\) on its smooth domain. Expanding the
finite commutator proves the last identity. Multiplication by
\(R/\sqrt3\) and compression by \(J_R\) recovers the definition of
\(X_a^D\). The operation is square-then-compress on the original
fifteen-space and a genuine square on its stated first-order closure.
\end{proof}
\begin{remark}[The first-order closure has its own carrier]
The twenty-five-dimensional closure is a reducing comparison carrier for
the specified smooth differential operator. It does not change the defined
level-one module \(\mathcal H_1\), identify it with a different
\(\mathcal H_m\), or turn its compressed square into the square of
\(D_F\). The additional multiplicities are precisely the first-order
partners displayed in the theorem.
\end{remark}

\subsection{Exact fibre sampling of the full retained carrier}
\label{p1:sec:v203-full-carrier-sampling}
The following construction uses the existing forms-plus-spinors bundle and
the retained Hurwitz cell. Its image has dimension fifteen; sample-fibre
coordinates represent the same fifteen-dimensional subspace. The scalar comparison band of
Corollary~\ref{p1:cor:v185-2t-degree-two-sampling-isometry} remains a distinct
fourteen-dimensional representation.

\begin{theorem}[Exact Hurwitz sampling of the full smooth carrier]
\label{p1:thm:v203-full-carrier-sampling}
Write $R=\Rmom$, $Q=2T\subset S^3(1)$ and
$\mathcal E=(\Lambda^*T^*S^3(R))_{\mathbb C}\oplus\mathbb S$, with its
geometric Hermitian metric. On the canonical smooth carrier of
Theorem~\ref{p1:thm:canonical-retained-smooth-model}, define
\[
 \mathsf S_R:\mathcal L_R\longrightarrow\bigoplus_{q\in Q}\mathcal E_{Rq},
 \qquad (\mathsf S_Rs)_q=\sqrt{V_R/24}\,s(Rq).
\]
Then $\mathsf S_R^*\mathsf S_R=I$ and
$\dim\mathcal V_R=15$, where $\mathcal V_R=\operatorname{Ran}\mathsf S_R$.
The full ten-dimensional fibre permits the intrinsic Clifford symbol to
act; the retained carrier has zero three-form component. Changing fibre
frames changes only the unitary coordinates of the target.
\end{theorem}
\begin{proof}
Different form degrees and the spinor summand are pointwise orthogonal.
In dimensionless coordinates $q_a=x_a/R$,
$\langle dx_a,dx_b\rangle=\delta_{ab}-q_aq_b$. The corresponding
two-by-two determinant for two-form pairings has degree at most two,
because its quartic terms cancel. In a left spin trivialization, the
lowest spinor spaces are represented by constant spinors and by $q^{-1}$
times constant spinors. Their Hermitian pairings have degree at most one.
Thus every retained-mode pairing has polynomial degree at most two.
The Hurwitz 5-design of
Theorem~\ref{p1:thm:v180-exceptional-scalar-moment-ranges} integrates each
pairing exactly. The factor $V_R/24$ converts this quadrature to the
continuum inner product, proving the identity and rank.
\end{proof}

\begin{corollary}[Finite amplitude action and intrinsic coordinate compiler]
\label{p1:cor:v203-full-symbol-transport}
For any retained geometric unitary $U(g)$, including the Spin(4) action and
its Hopf circle,
\[
 \widetilde U(g)=\mathsf S_RU(g)\mathsf S_R^*|_{\mathcal V_R},\qquad
 \mathsf S_RU(g)=\widetilde U(g)\mathsf S_R.
\]
For the pointwise symbol
$C_a(q)=R\,i[\mathscr D_{\rm LC},M_{x_a}](Rq)/\sqrt3$ on
$\mathcal E_{Rq}$, the existing intrinsic coordinate class satisfies
\[
 \boxed{X_a^D=J_R^*\mathsf S_R^*\operatorname{diag}_{q\in Q}C_a(q)
                  \mathsf S_RJ_R.}
\]
In particular it has the same sector fingerprint
$R^2(1,11/12,4/9,4/9,3/4,3/4)$ as
Theorem~\ref{p1:thm:intrinsic-tangential-coordinate-class}.
For scalar multiplication, its compressed sampled matrix is also exact
for every scalar symbol of polynomial degree at most three.
\end{corollary}
\begin{proof}
An isometry onto its image transports a unitary representation exactly.
In a left orthonormal frame, retained one-form, two-form and spinor
components have respective degrees at most one, two and one in $q$.
The first symbol has degree one. Each symbol matrix element therefore has
degree at most five and is integrated exactly by $Q$. This equals the
defining continuum compression of $X_a^D$ and hence its already proved
fingerprint. Scalar multiplication instead uses the degree-two bound on
mode pairings, so symbols through degree three are exact as well.
\end{proof}

The construction consists of finite bundle evaluation, the induced Hermitian
pairing, adjoint, image and operator transport. The phase freedom of $J_R$
is retained. Proposition~\ref{p1:prop:h1-e15-left-2t-obstruction} concerns
the distinct scalar left-regular target. The present fibre-valued isometry
and its continuous amplitude action have the explicit domains above.


\begin{corollary}[Eight-fibre realization of the full retained carrier]
\label{p1:cor:v233-eight-fibre-sampling}
For \(Q_8=\{\pm1,\pm i,\pm j,\pm k\}\subset S^3(1)\), define
\[
 (T_8s)_q=\sqrt{V_R/8}\,s(Rq),\qquad
 T_8:\mathcal L_R\longrightarrow\bigoplus_{q\in Q_8}\mathcal E_{Rq}.
\]
Then \(T_8^*T_8=I\), and its image has complex dimension fifteen in
eight ten-dimensional fibres. On that image every already defined
retained operator \(A\) has the exact representative \(T_8AT_8^*\).
Products, adjoints, spectra, states, actions and their tensor or block
transports are unchanged.
\end{corollary}
\begin{proof}
The eight coordinate-axis points have second moments \(\delta_{ab}/4\)
and vanishing odd moments, hence are a spherical 3-design.
The proof of Theorem~\ref{p1:thm:v203-full-carrier-sampling} shows that
every retained-mode pairing has degree at most two. The quadrature is
therefore exact, proving the isometry and image dimension.
The identities \(T_8^*T_8=I\) and
\((T_8AT_8^*)(T_8BT_8^*)=T_8ABT_8^*\) prove transport.
Its image identity is \(T_8T_8^*\).
\end{proof}

This is fibre-valued evaluation, not a scalar fifteen-dimensional
subrepresentation of \(\mathbb C^{Q_8}\) or \(\mathbb C^{2T}\).
It also does not assert that the pointwise diagonal-symbol compiler in
Corollary~\ref{p1:cor:v203-full-symbol-transport} remains exact with
eight nodes: that compiler uses degree-five matrix elements.
Operator transport requires only the isometry once the operator is
already constituted. It cannot supply omitted support-dependent
incidence data.

\begin{corollary}[Geometric ground state and a quantified heat limit]
\label{p1:cor:v233-geometric-ground-state}
For the precise retained Levi--Civita square \(K_{\rm LC}\) of
Corollary~\ref{p1:cor:canonical-retained-dirac-square}, a positive
trace-one density with \(\operatorname{Tr}(\rho K_{\rm LC})=0\)
is uniquely \(\Pi_{\rm sc}\). For \(\beta\ge0\), write
\[
 S_\beta=4e^{-3\beta/R^2}+6e^{-4\beta/R^2}
                   +4e^{-9\beta/(4R^2)},\qquad
 \rho_\beta=\frac{e^{-\beta K_{\rm LC}}}{1+S_\beta}.
\]
Then
\[
 \|\rho_\beta-\Pi_{\rm sc}\|_1
 =\frac{2S_\beta}{1+S_\beta}
 \le28e^{-9\beta/(4R^2)}.
\]
The principal symbol and coordinate commutators alone do not force
this ground state; Appendix~\ref{p1:app:v233-square-owner}
gives an explicit zero-order counterexample.
\end{corollary}
\begin{proof}
The positive eigenspaces have dimensions \(4,6,4\), with the three
positive eigenvalues displayed in the cited corollary. Zero expected
energy forces their diagonal weights to vanish; positivity also
annihilates their off-diagonal density entries. Only the one-dimensional
kernel remains. The heat state's scalar discrepancy and its total
excited weight both equal \(S_\beta/(1+S_\beta)\).
All fourteen excited terms are bounded by
\(e^{-9\beta/(4R^2)}\), giving the inequality.
\end{proof}



\begin{theorem}[Exact algebra and spectral transport onto an isometric image]
\label{p1:thm:v204-isometric-algebra-transport}
Let $T:W\to V$ be a unitary map between finite-dimensional Hermitian
spaces, including an isometry regarded as a map onto its image. Then
$\mathcal T_T(A)=TAT^*$ is a unital star-algebra isomorphism
$\operatorname{End}(W)\to\operatorname{End}(V)$, and
\[
 \mathcal T_T(AB)=\mathcal T_T(A)\mathcal T_T(B),\qquad
 \mathcal T_T(A^*)=\mathcal T_T(A)^*,\qquad
 \mathcal T_T([A,B])=[\mathcal T_T(A),\mathcal T_T(B)].
\]
It preserves operator norm, rank, trace, determinant, positivity, spectrum
and spectral projections. For a normal $A$ and every function $f$ on its
spectrum, $\mathcal T_T(f(A))=f(\mathcal T_T(A))$.
Subalgebras, their commutants and kernels of operator families are carried
to the corresponding objects on $V$. In particular these statements hold
for $T=\mathsf S_RJ_R$ and $V=\mathcal V_R$.
\end{theorem}
\begin{proof}
Insert $T^*T=I_W$ between factors and use $TT^*=I_V$; conjugation by $T^*$
is the inverse. Transport an orthonormal basis to prove trace, determinant,
norm and rank invariance. The spectral resolution of a normal operator
transports termwise, proving functional calculus and positivity. The
commutator identity and $T(\ker A)=\ker(TAT^*)$ prove the final statements.
In a larger coordinate host the unit is the image projector, not the host
identity; all commutants and spectra here are taken on $V$.
\end{proof}


\subsection{Intrinsic joint Hopf--Hurwitz comparison}

The intrinsic fixed-radius class supplies an even comparison between the
protected spin block and an oriented quaternionic support frame.  All the
objects in this subsection are geometric exports of the same round parent.
Their domains, codomains and equivariance are specified below.

\begin{theorem}[Intrinsic even Hopf--Hurwitz comparison]
\label{p1:thm:intrinsic-even-hopf-hurwitz-comparison}
Use the intrinsic class of
Theorem~\ref{p1:thm:intrinsic-tangential-coordinate-class}, so that
\(\varrho_D^2=3\Rmom^2/16>0\).  In a compatible protected spin frame write
the normalized chiral coordinate map as
\(c(q):W_R\to W_L\), for \(q\in\HH\).  For a retained real unit
\(u\in\HH\), the even word
\[
 A_D(u,q):=\varrho_D^{-2}P_RX_D(u)X_D(q)P_R
          =c(u)^\dagger c(q)\in\operatorname{End}(W_R)
\]
is independent of the removable relative chiral phase and the coordinate
scale.  Under a common spin-frame transport
\(c(q')=g_Lc(q)g_R^{-1}\), with the same transport applied to \(u\),
\[
 A_D(u',q')=g_RA_D(u,q)g_R^{-1}.
\]
The restriction \(c:\HH_{\CC}\otimes W_R\to W_L\) is a nonzero
geometric Clifford evaluation; \(c(h)\) is unitary for every unit
quaternion \(h\).  The relative product obeys
\(c(h)^\dagger c(q)=\rho(\bar hq)\) in a compatible quaternionic model,
and its scalar contraction is
\(\tfrac12\operatorname{Tr}(c(h)^\dagger c(q))=\langle h,q\rangle_{\RR}\)
on real inputs.
\end{theorem}

\begin{proof}
On the protected chiral sum the intrinsic coordinate has off-diagonal blocks
\(\varrho_Dc(q)\) and its Hermitian adjoint.  Multiplication of two such
blocks followed by \(P_R\) gives the displayed formula.  The common factor
\(\varrho_D^2\) cancels.  Rephasing the two spin summands multiplies the
two chiral factors by inverse phases, which also cancel.  Substituting the
transport law cancels \(g_L^\dagger g_L\), leaving conjugation by \(g_R\).
The real quaternion representation satisfies
\(\rho(\bar h)=\rho(h)^\dagger\),
\(\rho(\bar h)\rho(q)=\rho(\bar hq)\), and
\(\tfrac12\operatorname{Tr}\rho(a)=\operatorname{Re}a\).
These identities prove the remaining assertions.
\end{proof}

\begin{remark}[The parent unit is transported with the axial frame]
\label{p1:rem:joint-comparison-unit-transport}
The axial isometry \(\varphi=L_{\bar h}\) sends \(h\) to the coordinate
\(1\), but sends the old algebraic unit to \(\bar h\).  The transported
multiplication is
\[
 \mu'(a,b)=ahb,\qquad 1_{\mu'}=\bar h,
 \qquad A_D(\bar h,1)=\rho(h).
\]
Indeed \(\varphi^{-1}(a)=ha\), hence
\(\varphi(\varphi^{-1}(a)\varphi^{-1}(b))=ahb\).
Replacing the transported unit by \(1\) and the transported multiplication
by the standard product changes the input operation.  The nonzero intrinsic
scale used here is already fixed by the intrinsic theorem; no claim is made
for every point of the raw coordinate-completion cone.
In the axial/opposite-trimer orthonormal frame of the retained Hurwitz
vertex, the actual parent unit has coordinates
\((1/2,1/2,1/2,1/2)\); it is not identical with the first basis vector.
\end{remark}



\begin{theorem}[Quadratic Hopf--Hurwitz comparison]
\label{p1:thm:hopf-hurwitz-quadratic-comparison}
Let \(\varepsilon(s,t)=s^TEt\) be the invariant alternating form on the
defining spin representation, where
\(E=\left(\begin{smallmatrix}0&1\\-1&0\end{smallmatrix}\right)\).
For fixed real unit \(u\), complexify the variable \(q\) and put
\[
 \mathcal T_u(q;s)
 :=\operatorname{Tr}(A_D(u,q))\,\varepsilon(s,A_D(u,q)s).
\]
This is homogeneous of degree two in \(q\) and two in \(s\), and
\[
 \mathcal T_{u'}(q';g_Rs)=\mathcal T_u(q;s)
\]
under the full two-input transport of the preceding theorem.  In the frame
\(u=1\), use
\[
 \rho(r+x\mathbf i+y\mathbf j+z\mathbf k)
 =\begin{pmatrix}r+ix&y+iz\\-y+iz&r-ix\end{pmatrix}.
\]
Then \(\mathcal T_1(q;s)=s^TS(q)s\), with
\begin{equation}
\label{p1:eq:joint-quadratic-matrix}
 S(q)=2r\begin{pmatrix}-y+iz&-ix\\-ix&-y-iz\end{pmatrix}.
\end{equation}
Among homogeneous quadratic maps
\(\HH_{\CC}\to\operatorname{Sym}^2(W_R^*)\) using no additional
Cartan, trimer or spinorial input, full \(SU(2)\) conjugation covariance
gives a one-dimensional space, spanned by this map.
\end{theorem}

\begin{proof}
Both occurrences of \(A_D(u,q)\) are complex linear in \(q\); its adjoint
is applied only to the fixed real input \(u\).  This proves the bidegree.
Trace invariance and \(g_R^TEg_R=E\) prove covariance.  Taking the
symmetric part of \(EA_D\) gives
\(S=\operatorname{Tr}(A_D)(EA_D+(EA_D)^T)/2\), and multiplication gives
\eqref{p1:eq:joint-quadratic-matrix}.  The alternating scalar term has
zero quadratic contraction.

Identify symmetric spinor forms with the complex three-vector representation.
A quadratic vector covariant on \(q=r+v\) has the form
\(r^2w+rAv+B(v,v)\).  Rotation invariance forces \(w=0\) and makes
\(A\) scalar.  The stabilizer of a nonzero real vector \(v\) forces
\(B(v,v)\) to be parallel to \(v\).  A half-turn sending \(v\) to
\(-v\) negates that output, while quadratic dependence leaves it unchanged;
therefore \(B(v,v)=0\).  Polynomial continuation gives the complex result.
Only \(rv\) remains, proving the stated classification.
\end{proof}

\begin{corollary}[Joint C4 incidence and nonvanishing]
\label{p1:cor:joint-c4-quadratic-nonvanishing}
For the explicit Hurwitz units \(H_{\rm Hur}=2T\) and
\(U(1)_n=\{\cos t+n\sin t\}\),
\[
 H_{\rm Hur}\cap U(1)_n\cong
 \begin{cases}
 C_4,&n\text{ is a coordinate-axis direction},\\
 C_6,&n\text{ is a body-diagonal direction},\\
 C_2,&\text{otherwise}.
 \end{cases}
\]
Let \(L_+(n)\subset W_R\) be the positive oriented Hopf line and
\(\ell_h=\CC h\) the complexification of an actual real unit axial line.
Restriction and polarization give
\[
 m_+\in (L_+(n)\otimes\ell_h)^{*\otimes2},\qquad
 \|m_+\|^2=4r^2\bigl(|v|^2-(n\cdot v)^2\bigr),
 \quad A_D(u,h)=r+v.
\]
For
\(h_\epsilon=(1+\epsilon_1\mathbf i+\epsilon_2\mathbf j+
\epsilon_3\mathbf k)/2\), \(\epsilon_j=\pm1\), the joint C4 branch has
\[
 \|m_\epsilon\|^2=\frac12,
 \qquad
 m_\epsilon(e_+\otimes h_\epsilon,e_+\otimes h_\epsilon)
       =\frac{-\epsilon_2+i\epsilon_3}{2}
 \quad(n=\mathbf i).
\]
\end{corollary}

\begin{proof}
Membership in \(U(1)_n\) requires imaginary part parallel to \(n\).
The nonreal Hurwitz units have imaginary parts only along coordinate axes
or body diagonals.  Counting with \(\pm1\) gives respectively four, six,
or two elements, necessarily a cyclic group.  The restriction to a tensor
line follows from the two homogeneous degrees.  In the Cartan frame its
coefficient is \(2r(-y+iz)\), so its squared norm is
\(4r^2(y^2+z^2)\).  Covariance gives the invariant formula.  Substitution
of the eight half-integer units yields
\(f_\epsilon(n)=[3-(\epsilon\cdot n)^2]/4\); at a coordinate axis
\((\epsilon\cdot n)^2=1\).
\end{proof}

\begin{theorem}[Rooted Hurwitz phase, shell--trimer returns, and full family algebra]
\label{p1:thm:v556-rooted-hurwitz-family-return-algebra}
Fix the oriented axial vertex
\[
 u_+=\frac12(1+\mathbf i+\mathbf j+\mathbf k)
\]
and the cyclic family order on the opposite trimer
\[
 T_+=\left\{
 \frac12(1-\mathbf i+\mathbf j+\mathbf k),\,
 \frac12(1+\mathbf i-\mathbf j+\mathbf k),\,
 \frac12(1+\mathbf i+\mathbf j-\mathbf k)
 \right\}.
\]
For the joint quadratic of
Corollary~\ref{p1:cor:joint-c4-quadratic-nonvanishing} with
\(n=\mathbf i\), divide the three nonzero coefficients on this ordered
trimer by the first.  The resulting relative phase operator on
\(G=\CC^{T_+}\) is
\begin{equation}
\label{p1:eq:v556-rooted-family-phase}
 R_G=\operatorname{diag}(1,-i,i),
\end{equation}
up to a common phase and simultaneous cyclic relabelling.  Reversing the
orientation conjugates this representative.

Let \(J_G\) be the all-ones operator on \(G\), and put
\[
 Q_0=\frac13J_G,\qquad Q_\perp=I_G-Q_0.
\]
For a retained Hurwitz subset \(S\subset 2T\), define the nearest-neighbour
incidence map
\[
 \mathcal I_S:\CC^S\longrightarrow G,
 \qquad
 (\mathcal I_S)_{t,s}
 =\mathbf 1_{\{\langle t,s\rangle_{\RR^4}=1/2\}}.
\]
Then the canonical shell returns are
\begin{align}
 \mathcal I_{\mathcal M_6}\mathcal I_{\mathcal M_6}^\dagger
   &=5Q_0+2Q_\perp=2I_G+J_G,\nonumber\\
 K_7:=\mathcal I_{\mathcal G_7}\mathcal I_{\mathcal G_7}^\dagger
   &=8Q_0+2Q_\perp=2I_G+2J_G,\label{p1:eq:v556-shell-trimer-grams}\\
 K_9:=\mathcal I_{\mathcal G_9^{\rm shell}}
       \mathcal I_{\mathcal G_9^{\rm shell}}^\dagger
   &=Q_0+Q_\perp=I_G,\nonumber\\
 \mathcal I_{\mathcal N_8^+}\mathcal I_{\mathcal N_8^+}^\dagger
   &=7Q_0+Q_\perp=I_G+2J_G,\nonumber\\
 \mathcal I_{\mathcal G_{16}}\mathcal I_{\mathcal G_{16}}^\dagger
   &=9Q_0+3Q_\perp=3I_G+2J_G,\nonumber\\
 \mathcal I_{2T}\mathcal I_{2T}^\dagger
   &=16Q_0+4Q_\perp=4I_G+4J_G.\nonumber
\end{align}
For the already constituted local split
\(Q=\ell\oplus C\), with \(\ell\) represented by the axial column \(u_+\)
and \(C\) by the opposite trimer \(T_-\), the corresponding family
incidence Grams are
\begin{equation}
\label{p1:eq:v556-local-q-family-grams}
 K_\ell=J_G,
 \qquad
 K_C=I_G+J_G.
\end{equation}
Finally,
\begin{equation}
\label{p1:eq:v556-full-family-algebra}
 C^*(K_7,K_9,R_G)=M_3(\CC)=\End(G).
\end{equation}
Thus the same rooted Hurwitz geometry contains a two-generator unphased
return algebra together with a phase that generates the complete
noncommutative family matrix algebra.  The displayed incidence returns and rooted phase generate
the full matrix algebra, with no preferred matrix member.
\end{theorem}

\begin{proof}
For the displayed ordering of \(T_+\), the three coefficients from
Corollary~\ref{p1:cor:joint-c4-quadratic-nonvanishing} are
\[
 \frac{-1+i}{2},\qquad \frac{1+i}{2},\qquad \frac{-1-i}{2}.
\]
Their ratios to the first are \(1,-i,i\), proving
\eqref{p1:eq:v556-rooted-family-phase}.  Orientation reversal conjugates the
joint coefficient and exchanges the oriented trimer presentation, hence
conjugates the relative phase representative.

For every listed shell the three rows of \(\mathcal I_S\) have a common row
degree \(d_S\) and a common pairwise overlap \(c_S\).  Exact Hurwitz
incidence gives
\[
\begin{array}{c|rrrrrr}
 S&\mathcal M_6&\mathcal G_7&\mathcal G_9^{\rm shell}&
 \mathcal N_8^+&\mathcal G_{16}&2T\\ \hline
 d_S&3&4&1&3&5&8\\
 c_S&1&2&0&2&2&4
\end{array}
\]
and therefore
\[
 \mathcal I_S\mathcal I_S^\dagger=(d_S-c_S)I_G+c_SJ_G.
\]
This is exactly \eqref{p1:eq:v556-shell-trimer-grams}.  The axial vertex
\(u_+\) is adjacent to all three vertices of \(T_+\), giving the all-ones
axial column and hence \(K_\ell=J_G\).  Each vertex of \(T_+\) is adjacent
to the two complementary vertices of \(T_-\); after the cyclic basis is
fixed the opposite-trimer incidence block is \(J_G-I_G\), so
\((J_G-I_G)^2=I_G+J_G\), proving
\eqref{p1:eq:v556-local-q-family-grams}.

Since \(K_9=I_G\) and
\(J_G=(K_7-2K_9)/2\), the unphased Bose--Mesner return algebra is already
contained in \(C^*(K_7,K_9)\).  The phase \(R_G\) has simple spectrum, so
its three rank-one spectral projectors \(P_a,P_b,P_c\) are polynomials in
\(R_G\).  In the rooted basis
\(P_rJ_GP_s=|e_r\rangle\langle e_s|\); hence all matrix units belong to
\(C^*(K_7,K_9,R_G)\), proving \eqref{p1:eq:v556-full-family-algebra}.
\end{proof}

On real quaternionic inputs,
\(S^\dagger S=\det(S)I=4r^2|v|^2I\) and \(E^TSE=\bar S\).
The displayed coefficient is a frame expression for a tensor, not a selected
spinor phase.  Finite Hurwitz conjugation alone permits a second quadratic
covariant \((yz,zx,xy)\); the exact finite-only classification and its
axial restriction are proved in
Appendix~\ref{p1:app:joint-axis-comparison-readouts}.  In particular,
continuous uniqueness must not be asserted under finite symmetry alone.


\begin{corollary}[Rooted primal--dual incidence compresses to the family singlet]
\label{p1:cor:v570-rooted-primal-dual-singlet-return}
Let $D$ be the normalized dual $24$-cell in the $F_4$ realization, written in
an orthonormal real frame as
\[
 D=\left\{\frac{\varepsilon_i e_i+\varepsilon_j e_j}{\sqrt2}:
 1\le i<j\le4,\ \varepsilon_i,\varepsilon_j\in\{\pm1\}\right\}.
\]
For a primal unit vertex $p$ put
\[
 \lambda_p(n):=\sqrt2\,\langle p,n\rangle,
 \qquad s_p:=(\lambda_p(n))_{n\in D}\in\mathbb R^D.
\]
For the rooted axial vertex
\[
 u_+=\tfrac12(1,1,1,1)
\]
and the ordered family trimer $T_+=\{p_0,p_1,p_2\}$ of
Theorem~\ref{p1:thm:v556-rooted-hurwitz-family-return-algebra}, define
$S:G\to\mathbb C^D$ by $Se_j=s_{p_j}$.  Then
\[
 \boxed{S^\dagger s_{u_+}=6(1,1,1)^T=6\sqrt3\,c,}
 \qquad c:=\frac1{\sqrt3}(1,1,1)^T,
\]
and hence
\[
 \boxed{Q_\perp S^\dagger s_{u_+}=0,\qquad
 Q_0S^\dagger s_{u_+}=S^\dagger s_{u_+}.}
\]
Moreover the three level sets
\[
 X_\sigma:=\{n\in D:\lambda_{u_+}(n)=\sigma\},
 \qquad \sigma\in\{+1,0,-1\},
\]
have cardinalities $(6,12,6)$ and are preserved by the rooted $C_3$ stabilizer
of $u_+$.  Consequently every $C_3$-equivariant family readout formed from the
unmarked layer data lands in the fixed line $\mathbb Cc=Q_0G$.

The common-parent return therefore lies in the fixed line of the rooted
cyclic action, with the coefficient determined by the primal--dual Gram.
\end{corollary}
\begin{proof}
Sign symmetry gives
\[
 \sum_{n\in D}|n\rangle\langle n|=6I_4.
\]
Every $p_j\in T_+$ satisfies $\langle p_j,u_+\rangle=1/2$, so
\[
 \langle s_{p_j},s_{u_+}\rangle
 =2\sum_{n\in D}\langle p_j,n\rangle\langle u_+,n\rangle
 =12\langle p_j,u_+\rangle=6.
\]
This proves the displayed cross return.  For each of the six coordinate pairs,
the four sign choices give one $+1$, two zeroes and one $-1$ value of
$\lambda_{u_+}$, proving the $(6,12,6)$ count.  The rooted $C_3$ fixes $u_+$,
preserves the dual cell and cycles the three $p_j$, so it preserves every
$X_\sigma$ and its equivariant image lies in $G^{C_3}=\mathbb Cc$.
\end{proof}




\begin{proposition}[Natural axial relative determinant line]
\label{p1:prop:axial-relative-determinant-line}
For a complex line \(L\), the ordered split \(Q=\ell\oplus C\) with
\(\dim_\CC C=3\), and \(E_{\det}=L^{*\otimes2}\oplus\Lambda^2C\),
there is a natural isomorphism
\begin{equation}
\label{p1:eq:axial-relative-determinant-line}
 (L\otimes\ell)^{*\otimes2}
 \cong\det E_{\det}\otimes(\det Q)^{-2}
 \cong\operatorname{Hom}((\det Q)^{\otimes2},\det E_{\det}).
\end{equation}
It holds fibrewise for bundles of these ranks on a common base.  No volume
form or nonzero determinant element is chosen.
\end{proposition}

\begin{proof}
The ordered split gives \(\det Q=\ell\otimes\det C\).  In rank three,
\[
 \Lambda^2C\longrightarrow C^*\otimes\det C,
 \quad v\wedge w\longmapsto[u\mapsto u\wedge v\wedge w]
\]
is a natural isomorphism.  Taking its determinant yields
\(\det(\Lambda^2C)=(\det C)^{-1}\otimes(\det C)^3=(\det C)^2\).
Substitution cancels the two rank-three determinant factors by evaluation and
leaves \(L^{*\otimes2}\otimes\ell^{*\otimes2}\).  Under independent
frame changes \(z,\lambda,g\) on \(L,\ell,C\), the numerator and
denominator characters are \(z^{-2}(\det g)^2\) and
\(\lambda^2(\det g)^2\); their ratio is \(z^{-2}\lambda^{-2}\).
Thus the entire phase action is retained.
\end{proof}

The rank-four \(E_{\det}\) is a comparison carrier expressing this line.
Its fibres have the displayed representation type; it is not a tangent-bundle
identification.
The identity identifies the target of a nonzero relative map; it does not by
itself choose such a map.



\begin{theorem}[Complete joint Hopf--axial orbit]
\label{p1:thm:joint-hopf-axial-orbit}
Retain the tuple
\[
 D_0=(u=1,H_{\rm Hur}=2T,n_0=\mathbf i,
 h_0=(1+\mathbf i+\mathbf j+\mathbf k)/2)
\]
with common \(SU(2)\) conjugation on its quaternionic inputs and defining
spin transport on \(W_R\).  Its orbit is
\(\mathcal O=SU(2)/\{\pm1\}\cong SO(3)\).
The fibres
\[
 L_g=gL_+(n_0),\quad \ell_g=\CC gh_0g^{-1},\quad N_g=L_g\otimes\ell_g
\]
form the Hermitian sign line
\(N_{\mathcal O}=SU(2)\times_{\{\pm1\}}\CC_{\rm sign}\).
\end{theorem}

\begin{proof}
The imaginary part \(v_0\) of \(h_0\) is not parallel to \(n_0\).
A rotation fixing both fixes their cross product and hence a full frame,
so it is the identity in \(SO(3)\).  Its two spin lifts fix every
conjugation input, including the support and unit.  This proves the orbit
claim.  The axial unit \(h_g=gh_0g^{-1}\) is globally defined on the orbit.
For a temporary unit \(s_0\in L_+(n_0)\), the lifted tensor frame
\(e_g=(gs_0)\otimes h_g\) satisfies \(e_{-g}=-e_g\), giving precisely
the sign-associated line.  The temporary frame is not an extra chosen
geometric input.
\end{proof}

\begin{theorem}[Minimal quadratic descent on the joint orbit]
\label{p1:thm:joint-orbit-quadratic-descent}
For the complete conjugation-only tuple above,
\[
 \dim_\CC\Gamma_{\rm eq}(\mathcal O,N_{\mathcal O}^{*\otimes p})
 =\begin{cases}0,&p\text{ odd},\\1,&p\text{ even}.
 \end{cases}
\]
The intrinsic word \(\mathcal T\) supplies a nowhere-zero quadratic
\(m_{\mathcal O}\) with \(\|m_{\mathcal O}\|^2=1/2\).  Via
\eqref{p1:eq:axial-relative-determinant-line} it is also a nonzero relative
determinant isomorphism
\[
 (\det Q_{\mathcal O})^{\otimes2}
 \xrightarrow{\ \sim\ }
 \det(L_{\mathcal O}^{*\otimes2}\oplus\Lambda^2C_{\mathcal O}).
\]
Two is the first permitted positive tensor degree on this signature.
\end{theorem}

\begin{proof}
An equivariant section over \(G/H\) is determined by its value over the
identity coset, and that value must be \(H\)-fixed.  Conversely every such
value extends by transport.  The isotropy character is \((-1)^p\),
proving the dimension formula.  This argument alone does not supply a
nonzero value.  The intrinsic polynomial does: at the reference point,
\[
 S(h_0)=\begin{pmatrix}(-1+i)/2&-i/2\\-i/2&(-1-i)/2\end{pmatrix},
 \quad m_0(e_0,e_0)=\mu_0=(-1+i)/2.
\]
Conjugation covariance gives \(m_g(e_g,e_g)=\mu_0\).  The two deck signs
cancel in the square, so this tensor descends and has constant squared norm
\(1/2\).  The determinant claim follows by the natural isomorphism.
\end{proof}

\begin{corollary}[Sign-line torsion and actual pullbacks]
\label{p1:cor:joint-orbit-sign-line-torsion}
Let \(\alpha\ne0\) classify \(SU(2)\to SO(3)\).  Then
\[
 c_1(N_{\mathcal O})=\beta(\alpha)\ne0
 \quad\hbox{in }H^2(\mathcal O;\mathbb Z)=\mathbb Z_2,
 \qquad 2c_1(N_{\mathcal O})=0.
\]
Thus \(N_{\mathcal O}\) has no global nonzero linear section, even without
equivariance, and \(m_{\mathcal O}\) has no global linear square root.
An actual retained map \(f:Z\to\mathcal O\) admitting a lift to \(SU(2)\)
trivializes \(f^*N_{\mathcal O}\).  A spinorial lift changes the input
signature of the minimal-degree theorem.
\end{corollary}

\begin{proof}
The usual one-cell-in-each-dimension CW complex of
\(SO(3)\cong\mathbb{RP}^3\) has cellular boundary multiplication by two
in even positive degree and zero in odd degree.  Consequently
\(H^1(\mathcal O;\mathbb Z)=0\),
\(H^1(\mathcal O;\mathbb Z_2)=\mathbb Z_2\), and
\(H^2(\mathcal O;\mathbb Z)=\mathbb Z_2\).
The coefficient exact sequence for multiplication by two makes
\(\beta:H^1(\mathcal O;\mathbb Z_2)\to H^2(\mathcal O;\mathbb Z)\)
an isomorphism.  For sign transitions \((-1)^{a_{ij}}\), choose integer
lifts of \(a_{ij}\).  The triple-overlap cocycle \(\delta a/2\) represents
both the Chern class of the complex sign line and \(\beta(\alpha)\).
A nowhere-zero linear section would trivialize that class.  A retained lift
provides the frame \(e_g\) on the pullback, proving the final assertion.
\end{proof}

The orbit is a parameter base of a hosted geometric construction.  Its
torsion is not a statement about a line whose base is the completed smooth
\(M_8\), and geometric hosting supplies neither a map to the orbit nor its
lift.  With the axial input fixed while only the Hopf line varies, the
quadratic instead defines an \(\mathcal O(2)\) section on \(\mathbb{CP}^1\)
with two zeros; Appendix~\ref{p1:app:joint-axis-comparison-readouts} records
that different restriction problem.



\begin{theorem}[Unitary connection of the joint comparison]
\label{p1:thm:joint-comparison-unitary-connection}
The constant-norm \(m_{\mathcal O}\) determines a unique Hermitian
connection \(\nabla^m\) on \(N_{\mathcal O}\) satisfying
\(\nabla^m m_{\mathcal O}=0\).  It is flat and its holonomy around the
generator of \(\pi_1(\mathcal O)\) is \(-1\).
\end{theorem}

\begin{proof}
Locally choose a unit frame in which the coefficient of the normalized
quadratic is a fixed positive real number.  Such frames differ by constant
signs.  Declaring them parallel is consistent and gives a flat unitary
connection.  The nontrivial sign cover returns a lifted frame with its
negative around the generator.  If another unitary connection differs by
\(ia\), its induced action on the nonzero dual-square tensor differs by
\(-2ia\,m_{\mathcal O}\), forcing \(a=0\).
\end{proof}

\begin{proposition}[Axial connection from the parent vectors]
\label{p1:prop:joint-parent-axial-connection}
On the same orbit write \(h=r+v\), \(|n|=1\),
\(c=n\cdot v\), and \(s^2=|n\times v|^2>0\).  With the Hopf spinor
projector connection fixed, the existing axial line with its actual unit
\(h\) has the geometrically constructed connection
\begin{equation}
\label{p1:eq:joint-parent-axial-connection}
 \nabla^{\rm axial}h=iA_{\rm axial}h,
 \qquad
 A_{\rm axial}
 =-\frac{(n\times v)\cdot(dv-c\,dn)}{2|n\times v|^2}.
\end{equation}
Its tensor product with the spinor projector connection is \(\nabla^m\).
It differs from the ambient real-projector connection of the axial line,
whose coefficient in the unit \(h\) is zero.
\end{proposition}

\begin{proof}
In a lifted spin frame the projector coefficient is
\(A_{\mathrm{proj}}=-is_0^\dagger g^\dagger dg\,s_0\).
Set \(e_2=(v-cn)/s\) and \(e_3=n\times e_2\).  Differentiation gives
\[
 \frac12e_3\cdot de_2
 =\frac{(n\times v)\cdot(dv-c\,dn)}{2s^2}=A_{\mathrm{proj}}.
\]
For a joint angular velocity \(a\),
\(dn=2a\times n\), \(dv=2a\times v\); the numerator is
\(2s^2(a\cdot n)\), equal to twice the half-spin projector coefficient
times \(s^2\).  This fixes the factor \(1/2\) by the spin cover.
Both coefficients are unchanged under \(g\mapsto-g\) and descend to
the orbit.  Their sum on the tensor line is zero in the lifted frame,
where the quadratic coefficient is constant, hence equals \(\nabla^m\).
The real axial projector coefficient is zero because
\(\langle h,dh\rangle=0\).  Finally, in the left-invariant quaternionic
frame, \(A_{\mathrm{proj}}(X_{\mathbf i})=1\),
\(A_{\mathrm{proj}}(X_{\mathbf j})=A_{\mathrm{proj}}(X_{\mathbf k})=0\), and
\([X_{\mathbf j},X_{\mathbf k}]=2X_{\mathbf i}\) gives
\(dA_{\mathrm{proj}}(X_{\mathbf j},X_{\mathbf k})=-2\).
Thus the projector connection alone does not preserve the quadratic.
\end{proof}

The uniqueness concerns the tensor-line connection.  If the spinor
connection is also varied, parallelism does not fix its allocation between
the two factors.  These connections live on the stated geometric parameter
orbit.  Identification with a connection used by a later action requires
its actual base map and representation intertwiner.

\subsection{Ambient metric-cone export}

\begin{theorem}[Metric-cone Clifford-symbol coordinate class]
\label{p1:thm:metric-cone-clifford-coordinate-class}
\leavevmode\par
Use the already derived multiplicity-free identification
\[
 \mathcal H_1\cong
 \Lambda^{\le2}_{\CC}V^*\oplus S_+\oplus S_-.
\]
Compress the common ambient-cone Clifford symbol to this carrier and match the
normalized scalar-to-vector edge.  Then, modulo the independent phase changes
of the six irreducible $\operatorname{Spin}(4)$ sectors already quotiented in the coordinate
classification, there is exactly one coordinate class on this declared
ambient common-cone-symbol branch.  It is
\[
 \widehat X_a^*
 =\frac{\Rmom}{2}\left(P_{\le2}c_aP_{\le2}\oplus c_a^S\right),
 \qquad [\widehat X^*]=\mathsf C_{\rm cone}.
\]
Consequently the three raw reduced completion magnitudes are not independent
data on the ambient common-symbol branch.  No uniqueness is asserted across
the raw atlas or across the intrinsic tangential branch of
Proposition~\ref{p1:prop:cone-tangential-symbol-split}.
\end{theorem}

\begin{proof}
The cone construction and its Clifford symbol are functorial in the oriented
round metric.  The scalar edge gives the coefficient $\Rmom/2$: if $\phi_0$ is
the normalized constant and $\phi_a$ the normalized first harmonic, then
$M_{x_a}\phi_0=(\Rmom/2)\phi_a$.  Because all six irreducibles occur with
multiplicity one, any two equivariant identifications of $\mathcal H_1$ with
the compressed Clifford fibre differ by a phase on each irreducible; these are
precisely the previously quotiented rephasings.  A relative continuous
rescaling of exterior and spinor Clifford pieces would change their common
metric Clifford relation, so no such modulus remains.  Equation
\eqref{p1:eq:principal-symbol-commutator} shows that uniqueness of a complete
first-order operator is not needed.
\end{proof}

\begin{corollary}[Invariant quadratic fingerprint of the cone point]
\label{p1:cor:metric-cone-quadratic-fingerprint}
Let \(Q_{\rm cone}:=\sum_{a=1}^4(\widehat X_a^*)^2\).  In the sector order
\[
 (0,0),\quad(\tfrac12,\tfrac12),\quad(1,0),\quad(0,1),
 \quad(\tfrac12,0),\quad(0,\tfrac12),
\]
its scalar values are
\[
 \boxed{
 Q_{\rm cone}=
 \left(\Rmom^2,\ \Rmom^2,\ \frac{\Rmom^2}{2},\
       \frac{\Rmom^2}{2},\ \Rmom^2,\ \Rmom^2\right).}
\]
This fingerprint is invariant under all sector rephasings and therefore
distinguishes \(\mathsf C_{\rm cone}\) from any finite class with different
sectorwise values.
\end{corollary}

\begin{proof}
On \(\Lambda^0\oplus\Lambda^1\) the four ambient Clifford directions retain
all outgoing degrees, whereas on terminal \(\Lambda^2\) only the two
contraction channels remain after projection.  Thus the exterior contribution
is respectively \(4\), \(4\), and \(2\) before multiplication by
\(\Rmom^2/4\).  The full spin Clifford module contributes \(4\), giving the
displayed values.  This is also the content of
Corollary~\ref{p1:cor:clifford-level-one-values} written sectorwise.
\end{proof}

\begin{corollary}[Cone--tangential inequivalence]
\label{p1:cor:cone-tangential-coordinate-inequivalence}
The two primitive round-symbol exports are distinct points of the raw atlas:
\[
 \boxed{\mathsf C_D\ne\mathsf C_{\rm cone}.}
\]
No sector rephasing, overall Clifford sign, or raw/coisometric
Clebsch--Gordan convention identifies them.
\end{corollary}

\begin{proof}
The fingerprints of
Theorem~\ref{p1:thm:intrinsic-tangential-coordinate-class} and
Corollary~\ref{p1:cor:metric-cone-quadratic-fingerprint} differ on the vector,
bivector, and spin sectors.  Sector rephasings act by unitary conjugation and
cannot change the scalar eigenvalues of \(Q=\sum_aX_a^2\).  A change between
raw and coisometric reduced intertwiners changes coefficient notation, not
the operator \(Q\).
\end{proof}

\subsection{Primitive ancestry grammar and exhaustion}

\begin{definition}[Primitive round-symbol ancestry grammar]
\label{p1:def:primitive-round-symbol-ancestry-grammar}
Let \(\mathfrak G_{\rm round}^{\rm prim}\) be the following one-symbol
constructor grammar.  Its frozen inputs are the oriented round sphere
\((S^3(\Rmom),g_{\Rmom})\), its standard fixed-radius inclusion into the flat
metric cone, the metric Clifford symbols on the intrinsic and ambient
cotangent targets, the unitary class \(J_R\), compression to \(\mathcal H_1\),
and the normalized scalar edge. Each word in this defined class uses one of
the two anchor
words
\[
 V^*\xrightarrow{i_{\Rmom}^*}T^*S^3(\Rmom)
 \xrightarrow{c_T}\operatorname{End}(\mathcal L_R)
\]
or
\[
 V^*\xrightarrow{\mathrm{id}}T^*\mathbb R^4|_{S^3(\Rmom)}
 \xrightarrow{c_4}\operatorname{End}(\Lambda^*V^*\oplus S),
\]
followed by the stated compression, normalization, and quotient by sector
phases. The class excludes sums of the two target words, an independent
tangential/normal weight, a nonparallel ambient anchor, a twist, quotient,
chosen Hopf direction, deformation, extra bundle copy, or a zero-order term
as a coordinate constructor.
\end{definition}

\begin{theorem}[Primitive round-symbol ancestry exhaustion]
\label{p1:thm:primitive-round-symbol-ancestry-exhaustion}
Relative to the constructor grammar
\(\mathfrak G_{\rm round}^{\rm prim}\), the complete essential image in
\(\mathfrak M_{\rm coord}(\mathcal H_1)/\!\sim\) is
\[
 \boxed{
 \operatorname{EssIm}(\mathfrak G_{\rm round}^{\rm prim})
 =\{\mathsf C_D,\mathsf C_{\rm cone}\}.}
\]
More precisely:
\begin{enumerate}[label=(\roman*)]
\item the intrinsic fixed-radius target has exactly the class
\(\mathsf C_D\);
\item the flat parallel ambient-cone target has exactly the class
\(\mathsf C_{\rm cone}\); and
\item the two targets are inequivalent, so the grammar has neither omission
nor duplication.
\end{enumerate}
Every other point of the raw three-modulus atlas requires a constructor or
datum outside this frozen grammar.  This is ancestry exhaustion relative to
the displayed search basis, not an absolute classification of every enlarged
notion of a natural symbol.
\end{theorem}

\begin{proof}
Fix a point \(p\in S^3(\Rmom)\).  Its stabilizer is
\(H\cong\SO(3)\), and the ambient vector representation restricts as
\[
 V^*|_H\cong W^*\oplus\mathbf1,
 \qquad T_p^*S^3\cong W^*.
\]
Therefore
\[
 \operatorname{Hom}_H(V^*|_H,T_p^*S^3)
 \cong\operatorname{Hom}_H(W^*,W^*)\cong\RR.
\]
Every equivariant intrinsic anchor is a scalar multiple of tangential
projection.  The nonzero scalar is fixed by the normalized scalar edge, and
metric Clifford uniqueness plus multiplicity-one compression leaves only the
sector phases already quotiented.  This proves (i) and reproduces
\(\mathsf C_D\).

For the ambient word, parallelness makes the anchor a constant endomorphism
\(A\in\operatorname{End}(V^*)\).  Equivariance requires \(A\) to commute with
the irreducible \(\SO(4)\) vector action, so \(A=aI\).  Compatibility with the
ambient metric Clifford relation gives \(|a|=1\); the two signs differ only by
the overall Clifford sign already quotiented, and scalar-edge normalization
fixes the same finite class.  This proves (ii) and gives
\(\mathsf C_{\rm cone}\).

Corollary~\ref{p1:cor:cone-tangential-coordinate-inequivalence} proves (iii).
The grammar contains no further anchor word and permits no linear combination
of its two differently typed targets, so any additional raw atlas point has
no typed constructor word in \(\mathfrak G_{\rm round}^{\rm prim}\).
\end{proof}

\begin{remark}[Sharpness of the exhaustion scope]
\label{p1:rem:ancestry-exhaustion-sharpness}
If ambient parallelness and the primitive-word rule are dropped, equivariance
alone does \emph{not} give exhaustion.  At each point the two isotypically
distinct summands permit
\[
 A_{\alpha,\beta}=\alpha P_T+\beta P_N,
\]
where \(P_T\) and \(P_N\) are tangential and normal projection.  The ratio
\(\beta/\alpha\) is an extra dimensionless radial-weight datum; allowing it
enlarges the ancestry grammar and can generate additional finite classes.
Likewise chosen Hopf/CR structure, twisting, quotienting, Berger deformation,
or an extra bundle owner defines a different search basis.  The theorem
therefore closes exactly the unextended round grammar and no more.
\end{remark}

\begin{remark}[The two mathematical ancestry branches]
The intrinsic and ambient target spaces each have one normalized primitive
output in the defined class. Both belong to the atlas. The larger raw
completion cone and the extended operation classes retain their own moduli.
\end{remark}

\begin{remark}[Compression rather than a finite full Clifford module]
The theorem does not contradict the finite $\Cl_4$ divisibility obstruction.
The full Clifford action lives on the cone parent.  The $15$-dimensional
carrier is its $\Lambda^{\le2}\oplus S$ compression, so the previously proved
terminal bivector leakage remains exactly the expected compression defect.
\end{remark}



\subsection{Exceptional finite moment resolution of the scalar comparison band}
\label{p1:sec:v180-exceptional-scalar-moments}

Let $d\sigma$ be probability-normalized round measure on the unit
three-sphere, and use the explicit $2T$, $2O$, and $2I$ vertex sets already
constructed in this paper.

\begin{theorem}[Exact exceptional polynomial moment ranges]
\label{p1:thm:v180-exceptional-scalar-moment-ranges}
Uniform evaluation on $2T$, $2O$, and $2I$ agrees with the round integral
for every polynomial of total degree at most $5$, $7$, and $11$,
respectively.  Equivalently, for every even monomial one has
\[
 \int_{S^3}\prod_{i=0}^{3}x_i^{2a_i}\,d\sigma
 =\frac{\prod_i(2a_i)!}
 {4^{\sum_i a_i}(\sum_i a_i+1)!\prod_i a_i!},
\]
and the three explicit vertex sets reproduce this value for
$\sum_i a_i\le2,3,5$, respectively; odd moments vanish by their independent
coordinate-sign symmetries.
\end{theorem}
\begin{proof}
The displayed round moment is the standard beta-integral on $S^3$.
Substitution of the explicit coordinate orbits reduces the comparison to
$15$, $35$, and $126$ even monomials.  The orbit sign symmetries remove every
odd-exponent case.  The resulting finite equalities establish the stated
degrees without introducing an additional quadrature or action premise.
\end{proof}

\begin{corollary}[Exceptional groups resolve successive scalar action orders]
\label{p1:cor:v180-exceptional-scalar-action-orders}
For a real scalar profile in the complete degree-two position band
$\mathcal P_x(2)=E_0^x\oplus E_1^x\oplus E_2^x$, uniform $2T$ evaluation is
exact for the quadratic and Dirichlet pairings, uniform $2O$ evaluation is
also exact for every cubic scalar vertex, and uniform $2I$ evaluation is also
exact for every quartic scalar vertex.  These inclusions are sharp on this
band.  For $f=4x_0^2-1$,
\[
\begin{array}{c|ccc}
 &\langle f^2\rangle&\langle f^3\rangle&\langle f^4\rangle\\ \hline
 2T&1&2&7\\
 2O&1&1&4\\
 2I&1&1&3\\
 S^3&1&1&3 .
\end{array}
\]
Thus the exceptional fork has a precise finite-action readout role:
$2T$ resolves the quadratic level, $2O$ the cubic level, and $2I$ the
quartic level, on the stated scalar comparison family.  This is not a subgroup
chain $2T\subset2O\subset2I$.
\end{corollary}

\begin{theorem}[Sharp complete-band interaction frontier]
\label{p1:thm:v181-exceptional-complete-band-frontier}
Let $t_{2T}=5$, $t_{2O}=7$, and $t_{2I}=11$ be the exact uniform-design
strengths of Theorem~\ref{p1:thm:v180-exceptional-scalar-moment-ranges}.
For an $r$-field local scalar vertex on the complete band
$\mathcal P_x(m)=\bigoplus_{k=0}^mE_k^x$, uniform evaluation on
$\Gamma\in\{2T,2O,2I\}$ is exact for every such vertex whenever
\[
 rm\le t_\Gamma.
\]
For $r=2,3,4$ this bound is sharp, and hence
\[
 \boxed{m_{\max}(\Gamma,r)=\left\lfloor\frac{t_\Gamma}{r}\right\rfloor}
\]
with
\[
\begin{array}{c|ccc}
 &2T&2O&2I\\ \hline
 r=2&2&3&5\\
 r=3&1&2&3\\
 r=4&1&1&2
\end{array}.
\]
\end{theorem}
\begin{proof}
A product of $r$ restricted scalar polynomials from $\mathcal P_x(m)$ has
total coordinate degree at most $rm$, so the positive statements follow
from the exact design strengths.  Sharp witnesses immediately above all nine
frontiers are recorded in Appendix~\ref{p1:app:v181-exceptional-frontier-witnesses};
they are ordinary scalar profiles in the next complete band and require no
additional quadrature assumption.
\end{proof}

\begin{corollary}[Common degree-two saturation of the exceptional action fork]
\label{p1:cor:v181-exceptional-common-degree-two-saturation}
The three action roles meet on one and the same maximal complete band:
\[
 \boxed{
 \left\lfloor\frac52\right\rfloor
 =\left\lfloor\frac73\right\rfloor
 =\left\lfloor\frac{11}{4}\right\rfloor=2.}
\]
Thus $2T$ is maximally complete at degree two for quadratic/kinetic pairings,
$2O$ is maximally complete at degree two for cubic vertices, and $2I$ is
maximally complete at degree two for quartic vertices. These are polynomial
evaluation bounds on the scalar band; that band is distinct from the
fifteen-dimensional carrier $\mathcal H_1$.
\end{corollary}

\begin{corollary}[Unique exceptional support for the complete quartic comparison band]
\label{p1:cor:v569-unique-quartic-comparison-support}
Let $H\cong2T$ be the occurring Hurwitz core and let
$H\subseteq\Gamma\subset\operatorname{Sp}(1)$ be a finite same-owner
comparison support.  If uniform evaluation on $\Gamma$ is exact for every
fourfold product
\[
 f_1f_2f_3f_4,\qquad f_j\in\mathcal P_x(2),
\]
then
\[
 \boxed{\Gamma\cong2I.}
\]
Thus the $2I$ entry of the exceptional fork is the unique finite exceptional
support containing the rooted Hurwitz core that resolves the complete quartic
$m_{\rm cmp}=2$ comparison band.  The statement concerns the finite
comparison support; it does not promote $2I$ to an unrestricted microscopic
symmetry of every field owner.
\end{corollary}
\begin{proof}
A finite subgroup of $\operatorname{Sp}(1)$ containing $2T$ has exceptional
type $2T$, $2O$, or $2I$.  By
Theorem~\ref{p1:thm:v181-exceptional-complete-band-frontier}, exactness of
all fourfold products on $\mathcal P_x(2)$ requires design strength at least
$4\cdot2=8$.  The exact strengths are $5,7,11$, respectively, so only $2I$
qualifies.
\end{proof}

\begin{corollary}[Comparison-support mobility is dependent]
\label{p1:cor:v569-comparison-support-mobility}
Under Corollary~\ref{p1:cor:v569-unique-quartic-comparison-support}, the
finite comparison support contains an element moving the rooted tetrahedral
projector frame.  Hence the mobility hypothesis of
Theorem~\ref{p1:thm:v237-mobile-tetrahedral-completion} is a dependent
consequence on this comparison branch.
\end{corollary}
\begin{proof}
Let $I\cong2I$ be the forced support.  By
Theorem~\ref{p1:thm:exceptional-normalizer-fork},
$N_{\operatorname{Sp}(1)}(H)=2O$ and $2O\cap I=H$.  Any $g\in I\setminus H$
therefore lies outside the frame stabilizer $2O$ and moves the frame.
\end{proof}


\begin{corollary}[Exact $2T$ sampling isometry on the common degree-two band]
\label{p1:cor:v185-2t-degree-two-sampling-isometry}
Let
\[
 \mathsf S_{2T,2}:\mathcal P_x(2)\longrightarrow\mathbb C^{24},
 \qquad
 (\mathsf S_{2T,2}f)(q):=\frac{f(q)}{\sqrt{24}}.
\]
With the round $L^2(S^3,d\sigma)$ inner product on $\mathcal P_x(2)$ and the
standard Hermitian inner product on $\mathbb C^{24}$,
\[
 \boxed{
 \langle \mathsf S_{2T,2}f,\mathsf S_{2T,2}g\rangle
 =\int_{S^3}\overline f g\,d\sigma .}
\]
Hence $\mathsf S_{2T,2}$ is an isometric embedding of the $14$-dimensional
complete degree-two scalar band.  If $U_H(s)$ is the retained Hopf-circle
unitary action on $\mathcal P_x(2)$, then
\[
 \widetilde U_H(s)
 :=\mathsf S_{2T,2}U_H(s)\mathsf S_{2T,2}^{-1}
 \quad\hbox{on }\operatorname{Ran}\mathsf S_{2T,2}
\]
defines an exact unitary amplitude action on the sampled image and obeys
\[
 \boxed{\mathsf S_{2T,2}U_H(s)=\widetilde U_H(s)\mathsf S_{2T,2}.}
\]
It may be extended by the identity on the orthogonal complement in
$\mathbb C^{24}$.  This is an amplitude-space intertwiner; it does not claim
that the connected Hopf circle permutes the finite $2T$ vertex set.
\end{corollary}
\begin{proof}
For $f,g\in\mathcal P_x(2)$ the polynomial $\overline f g$ has total
coordinate degree at most four.  Theorem~\ref{p1:thm:v180-exceptional-scalar-moment-ranges}
therefore makes uniform $2T$ evaluation exact, which gives the displayed
inner-product identity and injectivity.  Conjugating the already unitary Hopf
action by this isometry gives the stated unitary action on the image and the
intertwining identity.  A continuous nontrivial circle action need not act by
permutations on a finite set, so no vertex-permutation claim is used.
\end{proof}


\begin{theorem}[Exact design evaluation of compressed scalar observables]
\label{p1:thm:v203-scalar-observable-compiler}
Let $\Pi_m$ project onto the complete scalar band $\mathcal P_x(m)$,
and let $S_{Q,m}$ be normalized evaluation on a spherical $t$-design.
For a polynomial scalar symbol $f$, put
\[
 Q_m(f)=\Pi_mM_f\Pi_m|_{\mathcal P_x(m)},\qquad
 Q_{Q,m}(f)=S_{Q,m}^*\operatorname{diag}(f(q))S_{Q,m}.
\]
If $2m+\deg f\le t$, then $Q_{Q,m}(f)=Q_m(f)$. At $m=2$, the
guaranteed symbol degrees for $2T,2O,2I$ are respectively $1,3,7$.
These are degree bounds for evaluating the displayed compressed operators.
For every geometric isometry preserving the complete band,
\[
 U(g)Q_m(f)U(g)^*=Q_m(f\circ g^{-1}).
\]
Writing $U_H(s)=\exp(-isD_H)$ gives
$[D_H,Q_m(f)]=-i Q_m(\mathcal L_{N_H}f)$, with $N_H$ the
geometric Hopf vector field. The same identities hold for exact sampled
matrices on the stated symbol domain.
\end{theorem}
\begin{proof}
A matrix element integrates $\overline u f v$, of degree at most
$2m+\deg f$. Design exactness proves the compiler formula. The strengths
$5,7,11$ are those of
Theorem~\ref{p1:thm:v180-exceptional-scalar-moment-ranges}.
Geometric isometries preserve the band and its orthogonal complement, so
they commute with $\Pi_m$. Conjugate multiplication and then compress;
differentiation at the identity proves the commutator formula.
\end{proof}

\begin{corollary}[Quadratic correction confined to the top comparison block]
\label{p1:cor:v203-quadratic-top-block-correction}
On the fourteen-dimensional band $\mathcal P_x(2)$ and $Q=2T$, let
$E_f=Q_{Q,2}(f)-Q_2(f)$ for a quadratic symbol. If $\Pi_2^{\rm hom}$
projects onto the degree-two harmonic subspace, then
\[
 E_f=\Pi_2^{\rm hom}E_f\Pi_2^{\rm hom},\qquad
 \operatorname{rank}E_f\le9.
\]
The exact correction $Q_{Q,2}(f)-E_f$ restores the preceding covariance
without enlarging the carrier. It is fixed by the degree-six moment
discrepancy. Appendix~\ref{p1:app:v203-sampling-reference} exhibits a
nonzero discrepancy; direct quadratic quadrature must not be assumed exact.
\end{corollary}
\begin{proof}
If either test mode has degree at most one, the total degree of
$\overline u f v$ is at most five. Every such matrix element is exact.
Only the degree-two by degree-two block can differ, proving the formula.
Subtracting that actual discrepancy gives $Q_2(f)$ by definition.
\end{proof}




\clearpage
\section{Exceptional axes, angular cubature, and the Hopf interface}
\label{p1b:sec:new-exceptional-axis-cubature}
\noindent
The following comparisons use the exceptional groups embedded above
and do not posit new dynamical or physical integration measures.  Write
$\phi=(1+\sqrt5)/2$ and let $\mathcal V$ be the twelve normalized directions
proportional to $(0,\pm1,\pm\phi)$ and their cyclic coordinate
permutations.  Let $\mathcal E$ be the thirty unit edge-midpoint directions
of the regular icosahedron with vertex set $\mathcal V$, and let
$\mathcal F$ be the twenty unit face-centre directions.
The three sets consist of oriented endpoints of the six fivefold, fifteen
twofold and ten threefold axes in the $A_5$ rotation system, respectively.
These are \emph{axis orbits of $A_5$}, not copies of the different binary
groups $2T$, $2O$ and $2I$.

\begin{theorem}[Exact positive icosahedral 62-point angular cubature]
\label{p1b:thm:new-62-point-angular-cubature}
For every polynomial $f$ on $\mathbb R^3$ of total degree at most eleven,
restricted to $S^2$, one has
\[
 \frac1{4\pi}\int_{S^2} f\,d\Omega
 =\frac{81}{462\cdot20}\sum_{u\in\mathcal F}f(u)
 +\frac{256}{462\cdot30}\sum_{u\in\mathcal E}f(u)
 +\frac{125}{462\cdot12}\sum_{u\in\mathcal V}f(u).
\]
The orbit coefficients $(81,256,125)/462$ are the unique real weights of
sum one that annihilate the invariant harmonic tests in degrees six and ten,
and are all positive.  The formula is not exact on every degree-twelve
polynomial.
\end{theorem}
\begin{proof}
Each of the three endpoint sets is a full antipodal $A_5$ orbit.  By the
icosahedral harmonic-invariant calculation, there are no nonconstant
$A_5$-invariant harmonics of degrees $1,\ldots,11$ except for one line
in each of degrees six and ten.  Therefore invariant averaging reduces the
cubature identity to its constant and two invariant harmonic coefficients.
Choose an icosahedral vertex $n$ and test $P_6(n\cdot u)$ and
$P_{10}(n\cdot u)$, where $P_k$ is the degree-$k$ Legendre polynomial.
Direct evaluation from the displayed vertex, midpoint and face-centre
coordinates yields, for $\mathcal F,\mathcal E,\mathcal V$ in that order,
\[
\begin{array}{c|ccc}
 &\mathcal F&\mathcal E&\mathcal V\\\hline
 P_6&-11/45&-11/80&11/25\\
 P_{10}&247/1215&-247/1920&247/1875
\end{array}.
\]
Solving the two zero-moment equations together with the normalization
$w_F+w_E+w_V=1$ gives $(w_F,w_E,w_V)=(81,256,125)/462$.
Each individual orbit is centrally symmetric, so the odd harmonic moments
vanish.  The invariant tests and the $A_5$ symmetry then prove the full
polynomial formula through degree eleven.  For degree twelve the exact
weighted average of $P_{12}(n\cdot u)$ is $7429/33000\ne0$, whereas its
spherical mean vanishes.  This supplies a counterexample at the first
higher degree.
\end{proof}

\begin{corollary}[Hopf-weighted pullback observables]
\label{p1b:cor:new-hopf-pullback-11-cubature}
Let $\pi_H:S^3\to S^2$ be the chosen normalized Hopf fibration of
Paper~I-A, with $\int_{S^3}\eta_0\wedge d\eta_0=4\pi^2$.
For every such polynomial $f$, one has
\[
\int_{S^3}(f\circ\pi_H)\,\eta_0\wedge d\eta_0
 =4\pi^2\left[
 {81\over462}\langle f\rangle_{\mathcal F}
 +{256\over462}\langle f\rangle_{\mathcal E}
 +{125\over462}\langle f\rangle_{\mathcal V}\right].
\]
The embedded groups do not change the underlying $4\pi^2$: this is an
exact \emph{finite evaluation of selected pullback observables}, not a new
contact volume or a physical determinant correction.
\end{corollary}
\begin{proof}
The fibre pushforward of the normalized Hopf contact measure is
$4\pi^2$ times normalized spherical area on $S^2$; apply the preceding
theorem to the base polynomial.
\end{proof}

\section{Exceptional spinorial fixed-charge spectra: a comparison}
\label{p1b:sec:new-exceptional-spinorial-comparison}
For $\Gamma\in\{2T,2O,2I\}$ let $V$ be the defining complex
$SU(2)$ two-dimensional representation and let
$\mathcal V_{+1}$ denote the positive unit-charge Hopf sector with
$\lambda_j=4j^2-2$ and original multiplicities $2j$, $j\ge1$.
The central element $-1\in\Gamma$ acts by $-1$ on $\mathcal V_{+1}$,
so $\mathcal V_{+1}^\Gamma=0$ in the scalar sector.  Its tensor action
on $V$ has a compensating central sign; thus define the \emph{new,
explicitly separate} twisted spectral owner
\[
 \mathscr H_\Gamma=(\mathcal V_{+1}\otimes V)^\Gamma,
 \qquad L_\Gamma=L_{+1}\otimes\mathrm{id}\big|_{\mathscr H_\Gamma}.
\]
These are mathematical comparison sectors, not a replacement for the
primitive scalar determinant used in the downstream common action.

\begin{proposition}[Exact twisted multiplicity and zeta invariant]
\label{p1b:prop:new-twisted-exceptional-zeta}
Put $a_n^\Gamma=\dim(\operatorname{Sym}^nV)^\Gamma$.  For the above
fixed-charge comparison one has
\[
 m_j^\Gamma=a_{2j}^\Gamma+a_{2j-2}^\Gamma,
\qquad
 G_\Gamma(z):=\sum_{j\ge1}m_j^\Gamma z^j
 ={(1+z)(1+z^d)\over(1-z^a)(1-z^b)}-1,
\]
where $(a,b,d)$ is $(3,4,6)$ for $2T$, $(4,6,9)$ for $2O$, and
$(6,10,15)$ for $2I$.  The meromorphic continuation of
$Z_\Gamma(s)=\sum_{j\ge1}m_j^\Gamma(4j^2-2)^{-s}$ satisfies
\[
 Z_\Gamma(0)={a^2+b^2-3a-3b+6\over3ab},
 \qquad
 (Z_{2T}(0),Z_{2O}(0),Z_{2I}(0))=
 \left({5\over18},{7\over18},{47\over90}\right).
\]
\end{proposition}
\begin{proof}
The Clebsch--Gordan decomposition
$\operatorname{Sym}^kV\otimes V\cong
\operatorname{Sym}^{k+1}V\oplus\operatorname{Sym}^{k-1}V$
identifies the multiplicities.  Substitute the classical even-degree
binary-polyhedral Molien series, with $z=t^2$, to obtain the displayed
three generating functions.  For $D_\Gamma(s)=\sum_jm_j^\Gamma j^{-s}$
the Euler--Mellin continuation gives
$D_\Gamma(0)=(a^2+b^2-3a-3b+3)/(3ab)$ and
$\operatorname{Res}_{s=2}D_\Gamma(s)=4/(ab)$.
Expanding $(4j^2-2)^{-s}$ at $s=0$ yields
$Z_\Gamma(0)=D_\Gamma(0)+(1/4)\operatorname{Res}_{s=2}D_\Gamma(s)$,
which is the stated rational expression.
\end{proof}

\begin{remark}[Why neither averaging nor an amplitude follows]
The primitive scalar fixed-charge invariant sector is zero, while
$\mathscr H_\Gamma$ has a nonzero compensating spinorial fibre.  The
normalizer $2O$ has an independently derived scalar action interface in
Paper~III; it must not silently be replaced by this spinorial owner.
In particular, an arithmetic mean, axis-weighted mean, or weighted sum of
these $Z_\Gamma(0)$ values is not a new physical correction to
$4\pi^2$ without an action-derived common operator and a treatment of
overlapping invariant subspaces.
\end{remark}

\clearpage
\section{Complete comparative mathematics of the exceptional Hopf branches}
\label{p1b:sec:v11-complete-exceptional-comparisons}

The purpose of this section is mathematical completeness rather than physical
selection.  In particular, exact identities, failed equalities, different
normalization conventions, and explicitly conditional numerical comparisons
are all retained.  No conclusion below is removed because it is not currently
consumed by a later Born--Hopf paper.  The primitive smooth contact integral
$4\pi^2$ and its chosen normalization remain results of Paper~I-A.

\subsection{Embeddings, spherical quotients, and the unchanged contact host}
\label{p1b:subsec:v11-quotient-contact}

\begin{theorem}[Contact measure of the binary-polyhedral spherical quotients]
\label{p1b:thm:v11-quotient-contact}
Let $G=SU(2)\simeq S^3$ have the normalized left-invariant Hopf contact
form $\eta_0$ of Paper~I-A, oriented so that
$\int_G\eta_0\wedge d\eta_0=4\pi^2$.  Let
$\Gamma\in\{2T,2O,2I\}$ act freely by left multiplication, and let
$p_\Gamma:G\to\Gamma\backslash G$ be the quotient projection.  There is
a unique descended contact form $\bar\eta_\Gamma$ satisfying
$p_\Gamma^*\bar\eta_\Gamma=\eta_0$, and
\[
\boxed{\int_{\Gamma\backslash G}\bar\eta_\Gamma\wedge
 d\bar\eta_\Gamma=\frac{4\pi^2}{|\Gamma|}.}
\]
For the corresponding spherical rotation orbifold
$S^2/\bar\Gamma$, where $\bar\Gamma=\Gamma/\{\pm1\}$, one also has
\[
\chi_{\rm orb}(S^2/\bar\Gamma)
 =\frac12+\frac13+\frac1m-1=\frac4{|\Gamma|},
\qquad
\boxed{\int_{\Gamma\backslash G}\bar\eta_\Gamma\wedge
 d\bar\eta_\Gamma=\pi^2\chi_{\rm orb}(S^2/\bar\Gamma),}
\]
with $m=3,4,5$ for $2T,2O,2I$, respectively.  Explicitly,
\[
\begin{array}{c|c|c|c}
\Gamma&|\Gamma|&\chi_{\rm orb}&\displaystyle
 \int_{\Gamma\backslash G}\bar\eta_\Gamma\wedge d\bar\eta_\Gamma\\\hline
2T&24&1/6&\pi^2/6\\
2O&48&1/12&\pi^2/12\\
2I&120&1/30&\pi^2/30
\end{array}
\]
\end{theorem}
\begin{proof}
Left invariance of $\eta_0$ makes it invariant under the finite left
$\Gamma$ action.  A differential form invariant under a free, properly
discontinuous group action descends uniquely to the quotient.  As
$p_\Gamma$ is an orientation-preserving covering of degree $|\Gamma|$,
its change-of-variables formula divides the integral by $|\Gamma|$.
The induced finite rotation groups are $A_4,S_4,A_5$, with spherical
triangle signatures $(2,3,3)$, $(2,3,4)$, $(2,3,5)$.  Orbifold
Gauss--Bonnet gives
$\chi_{\rm orb}=1/2+1/3+1/m-1$; the indicated
orders give $\chi_{\rm orb}=4/|\Gamma|$.
\end{proof}

\begin{remark}[Three distinct meanings of ``putting a group in $S^3$'']
\label{p1b:rem:v11-embedding-quotient-sampling}
An embedded finite subset $\Gamma\subset S^3$ has zero continuous
three-dimensional contact measure and does not modify $4\pi^2$.
The contact integral of the different manifold $\Gamma\backslash S^3$
is $4\pi^2/|\Gamma|$.  Finally, a weighted finite evaluation
$\frac{4\pi^2}{|\Gamma|}\sum_{g\in\Gamma}f(g)$ is a quadrature
comparison on the \emph{original} $S^3$; for a general non-invariant
$f$ it is not an integral of a function on $\Gamma\backslash S^3$.
Nor is the orbifold Hopf fibration automatically a unit-Chern ordinary
circle bundle after quotienting.  These distinctions are part of the
mathematical data, not physical caveats to be dropped.
\end{remark}

\subsection{All three nearest-shell contractions and their exact failure tests}
\label{p1b:subsec:v11-shell-contractions}

\begin{proposition}[Defining, adjoint, and scalar-atom comparisons]
\label{p1b:prop:v11-three-contractions}
For $\Gamma\in\{2T,2O,2I\}$, write $n_\Gamma=|S_\Gamma|$,
$r_\Gamma=\operatorname{Re}(s)$ for $s\in S_\Gamma$, and
$\omega_\Gamma=r_\Gamma/n_\Gamma$ as in
Theorem~\ref{p1:thm:exceptional-345-shell-table}.  The normalized
nearest-shell eigenvalues on the defining and adjoint representations are
\[
\lambda_{\rm def}(T_\Gamma)=r_\Gamma,\qquad
\lambda_{\rm Ad}(T_\Gamma)=\frac{4r_\Gamma^2-1}{3}.
\]
In particular, the following complete comparison table holds:
\[
\begin{array}{c|c|c|c|c|c}
\Gamma&n_\Gamma&r_\Gamma&\lambda_{\rm Ad}&
 \omega_\Gamma&\frac12\lambda_{\rm Ad}\lambda_{\rm def}\\\hline
2T&8&1/2&0&1/16&0\\
2O&6&1/\sqrt2&1/3&1/(6\sqrt2)&1/(6\sqrt2)\\
2I&12&\varphi/2&\varphi/3&\varphi/24&\varphi^2/12
\end{array}
\quad \left(\varphi=\frac{1+\sqrt5}{2}\right).
\]
Thus the explicitly defined trial contraction
$\delta^{\rm trial}_\Gamma:=\frac12\lambda_{\rm Ad}(T_\Gamma)
\lambda_{\rm def}(T_\Gamma)$ satisfies
\[
\boxed{\frac{\delta_\Gamma^{\rm trial}}{\omega_\Gamma}
 =\frac{n_\Gamma(4r_\Gamma^2-1)}6
 \in\{0,1,2\varphi\},}
\]
with the displayed values in the order $2T,2O,2I$.
Consequently the equality $\delta_\Gamma^{\rm trial}=\omega_\Gamma$
holds for $2O$ and fails for $2T$ and $2I$.
\end{proposition}
\begin{proof}
The defining trace of $s\in SU(2)$ is $2r_\Gamma$, so the
central-shell scalar in dimension two is $r_\Gamma$.  The real
three-dimensional adjoint character is
$1+2\cos(2\arccos r_\Gamma)=4r_\Gamma^2-1$, and division by three
gives the adjoint eigenvalue.  Insert the shell data
$(n_\Gamma,r_\Gamma)=(8,1/2),(6,1/\sqrt2),(12,\varphi/2)$.
Since $4(\varphi/2)^2-1=\varphi$, division by
$\omega_\Gamma=r_\Gamma/n_\Gamma$ gives the stated ratios.
\end{proof}

\begin{remark}[Comparison with the downstream primitive smooth coefficient]
\leavevmode\par
The primitive smooth Hopf line used downstream has the separately
computed scale coefficient $\zeta_{\rm red}(0)=1/3$.  The finite
adjoint comparison above gives
\[
\lambda_{\rm Ad}(T_\Gamma)=\zeta_{\rm red}(0)
\quad\Longleftrightarrow\quad\Gamma=2O
\qquad(\Gamma\in\{2T,2O,2I\}).
\]
On the already anchored downstream octahedral branch the numerical
hierarchy expression is $4\pi^2+1/(6\sqrt2)$.  This comparison does
not redefine the $S^3$ contact residue, nor derive the downstream
action or give that smooth coefficient to the other two groups.
\end{remark}

\begin{remark}[Which equality is actually owned by a smooth operator?]
The octahedral equality is compatible with the separately established
normalizer/smooth-germ interface in Paper~III.  Applying the \emph{same
algebraic expression} to $2T$ and $2I$ is an exact finite test, not a
proof that their smooth determinant germs are the same operator or
receive their shell atoms as physical corrections.  Both failures are
retained as mathematical information.
\end{remark}

\subsection{Second-moment frames and the unique positive raw-frame match}
\label{p1b:subsec:v11-shell-frames}

\begin{theorem}[Complete exceptional nearest-shell frame comparison]
\label{p1b:thm:v11-three-shell-frames}
For $s\in S_\Gamma$ let
$a_s=\operatorname{Im}(s)/\|\operatorname{Im}(s)\|\in\mathbb R^3$
and define the \emph{unnormalized} frame metric
\[
 F_\Gamma(X,Y)=\sum_{s\in S_\Gamma}
  g_E(a_s,X)g_E(a_s,Y).
\]
Each set of imaginary shell directions is a unit tight frame, with
\[
 \sum_{s\in S_\Gamma}a_s\otimes a_s=\frac{n_\Gamma}{3}I_3,
 \qquad
\boxed{F_{2T}=\frac83g_E,\quad F_{2O}=2g_E,\quad
F_{2I}=4g_E.}
\]
For the original Hopf horizontal plane $H$ with Levi metric
$\gamma_0=2g_E|_H$,
\[
 F_{2T}|_H=\frac43\gamma_0,\qquad
 F_{2O}|_H=\gamma_0,\qquad
 F_{2I}|_H=2\gamma_0.
\]
If $w_T,w_O,w_I\ge0$, $w_T+w_O+w_I=1$, and
$\sum_\Gamma w_\Gamma F_\Gamma|_H=\gamma_0$, then necessarily
\[
\boxed{(w_T,w_O,w_I)=(0,1,0).}
\]
By contrast, each \emph{individually trace-normalized} metric
$\widehat F_\Gamma=(6/n_\Gamma)F_\Gamma$ equals $2g_E$,
so second moments alone do not distinguish the groups after that
normalization.
\end{theorem}
\begin{proof}
The $2T$ shell directions are the eight cube diagonals
$(\pm1,\pm1,\pm1)/\sqrt3$; their outer products sum to
$(8/3)I_3$.  The $2O$ directions are $\{\pm e_1,\pm e_2,\pm e_3\}$,
whose outer products sum to $2I_3$.  The twelve $2I$ directions form
the centrally symmetric icosahedral vertex orbit; its second-moment
tensor commutes with the irreducible three-dimensional icosahedral
rotation representation, hence is scalar, and its trace is twelve;
thus it equals $4I_3$.  Restrict to $H$ to obtain the displayed
ratios.  The weighted Levi condition reads
\[
 \frac83 w_T+2w_O+4w_I=2,
 \qquad w_T+w_O+w_I=1,
\]
so $\frac23w_T+2w_I=0$.  Nonnegativity forces
$w_T=w_I=0$ and $w_O=1$.  Multiplying each tight frame by
$6/n_\Gamma$ proves the last assertion.
\end{proof}

\begin{remark}[Why the positivity qualification matters]
If signed weights are allowed, the raw-frame equation instead admits
the one-parameter family $w_T=-3w_I$ and $w_O=1+2w_I$.
The uniqueness statement is therefore a theorem about \emph{convex}
mixtures of the unrescaled shell frames, not a general uniqueness
principle for every spectral or signed response operator.
\end{remark}

\subsection{A scalar test observable: exact moment defects of all branches}
\label{p1b:subsec:v11-moment-comparison}

\begin{proposition}[Complete $f=4x_0^2-1$ moment comparison]
\label{p1b:prop:v11-scalar-moment-comparison}
On the unit $S^3\subset\mathbb R^4$, put $f(x)=4x_0^2-1$.
Let $\langle\cdot\rangle_{S^3}$ be normalized Haar/contact
averaging and $\langle\cdot\rangle_\Gamma$ uniform averaging over
the embedded group vertices.  Then
\[
\begin{array}{c|c|c|c|c}
 &S^3&2T&2O&2I\\\hline
\langle f^2\rangle&1&1&1&1\\
\langle f^3\rangle&1&2&1&1\\
\langle f^4\rangle&3&7&4&3
\end{array}.
\]
In particular, the continuous Hopf observable and the three finite
sampling evaluations are
\[
\int_{S^3}f^4\,\eta_0\wedge d\eta_0=12\pi^2,
\qquad
\boxed{(Q_{2T}(f^4),Q_{2O}(f^4),Q_{2I}(f^4))
 =(28\pi^2,16\pi^2,12\pi^2),}
\]
where $Q_\Gamma(h):=(4\pi^2/|\Gamma|)\sum_{g\in\Gamma}h(g)$.
\end{proposition}
\begin{proof}
The normalized $S^3$ marginal moments are
$\langle x_0^{2k}\rangle=(\frac12)_k/(2)_k$; expanding
$(4x_0^2-1)^j$ for $j=2,3,4$ gives $1,1,3$.
For $2T$ the scalar coordinates $x_0$ have multiplicities
$2$ at $\pm1$, $6$ at $0$, and $16$ at $\pm\frac12$.
Thus the corresponding $f$ values are $3,-1,0$ with those
multiplicities and give $1,2,7$.  For $2O$ the multiplicities
are $2,18,16,12$ at the magnitudes
$1,0,\frac12,1/\sqrt2$, respectively, yielding $f$ values
$3,-1,0,1$ and averages $1,1,4$.  The $2I$ vertex set is an
11-design on $S^3$; the polynomial degrees of $f^2,f^3,f^4$ are
$4,6,8$, so all three finite averages equal their $S^3$ averages.
Finally multiply the normalized averages by the unchanged $4\pi^2$.
\end{proof}

\begin{remark}[Observable selection rather than a universal correction]
The three groups agree on this second moment; the tetrahedral defect
first appears at the cubic and the octahedral defect at the quartic
power of $f$, while the icosahedral samples are exact for all three.
These are comparisons for the \emph{declared polynomial observables};
they neither rescale $4\pi^2$ nor establish a universal group correction.
\end{remark}

\subsection{Thirty-one axes, the global quadratic frame, and distinct measures}
\label{p1b:subsec:v11-axis-measures}

\begin{proposition}[Axis-frame identity and exact per-point cubature coefficients]
\label{p1b:prop:v11-axis-frame-and-perpoint}
Let $\mathcal A_{31}$ consist of the $10$ threefold, $15$ twofold,
and $6$ fivefold \emph{unoriented} axes of the $A_5$ icosahedral
rotation group, and choose any unit representative $u_\ell$ of each.
Then
\[
\boxed{\sum_{\ell\in\mathcal A_{31}}
 u_\ell\otimes u_\ell=\frac{31}{3}I_3,\qquad
 \frac6{31}\sum_{\ell\in\mathcal A_{31}}
 u_\ell\otimes u_\ell=2I_3.}
\]
These axes have $20,30,12$ oriented endpoints in the three respective
orbits $\mathcal F,\mathcal E,\mathcal V$.  The 62-point degree-eleven
rule of Theorem~\ref{p1b:thm:new-62-point-angular-cubature}
can equivalently be written using the \emph{integer per-endpoint} weights
\[
\boxed{\frac1{27720}\left(
 243\sum_{u\in\mathcal F}f(u)
 +512\sum_{u\in\mathcal E}f(u)
 +625\sum_{u\in\mathcal V}f(u)\right).}
\]
The uniform 31-axis average assigns orbit weights
$(10,15,6)/31$, and is \emph{not} the degree-eleven rule: for an
icosahedral vertex $n$ its degree-six Legendre moment equals
\[
\boxed{\frac{10}{31}\langle P_6(n\cdot u)\rangle_{\mathcal F}
+\frac{15}{31}\langle P_6(n\cdot u)\rangle_{\mathcal E}
+\frac6{31}\langle P_6(n\cdot u)\rangle_{\mathcal V}
=-\frac{6721}{111600}\ne0.}
\]
\end{proposition}
\begin{proof}
Each axis orbit is an irreducible icosahedral second-moment frame:
the sum of its rank-one projectors is the orbit size divided by three
times $I_3$.  Summing the three orbits proves the first identity.
The per-endpoint coefficient identities are
$81/(462\cdot20)=243/27720$,
$256/(462\cdot30)=512/27720$, and
$125/(462\cdot12)=625/27720$.
The degree-six value follows by inserting the three entries
$(-11/45,-11/80,11/25)$ of the Legendre table in
Theorem~\ref{p1b:thm:new-62-point-angular-cubature}.
\end{proof}

\begin{remark}[What the angular weights do and do not count]
The weights $81:256:125$ follow from exactness through degree eleven
and select a positive \emph{angular cubature measure}.  Symmetry alone
only forces equality \emph{inside} each orbit; it does not force these
three orbit totals to coincide.  The axis orbit classes of $A_5$ are
not themselves three separate binary groups.  The twofold axes
inside $A_5$ are not an embedded binary octahedral subgroup
$2O\subset2I$, an inclusion which is false.  Applying the orbit
weights as if they were weights on three different group determinants
would change the mathematical object and requires an additional map.
\end{remark}

\subsection{Fixed-charge obstruction, full low-mode ladder, and spectral overlap}
\label{p1b:subsec:v11-spinorial-completion}

\begin{lemma}[No scalar or one-dimensional-character rescue of unit charge]
\label{p1b:lem:v11-unit-charge-obstruction}
For the primitive $q=+1$ scalar Hopf line, and for every
$\Gamma\in\{2T,2O,2I\}$,
\[
 \mathcal V_{+1}^{\Gamma}=0.
\]
Furthermore, twisting by a one-dimensional complex character of
$\Gamma$ cannot make this scalar line invariant: every such character
sends the common central element $-1$ to $+1$.
\end{lemma}
\begin{proof}
The central element $-1\in SU(2)$ acts on fixed Hopf charge $q$
by $(-1)^q$, hence by $-1$ at $q=+1$.  Invariance under any group
containing $-1$ would require the same vector to equal its negative.
For the character claim, $-1$ belongs to the commutator subgroup of
each binary group (their abelianizations are respectively $C_3,C_2$
and the trivial group).  Every one-dimensional complex character is
trivial on commutators.  The defining two-dimensional spinor $V$,
which also carries central sign $-1$, is a genuinely different twist
and gives the already defined nonzero $\mathscr H_\Gamma$.
\end{proof}

\begin{proposition}[Explicit initial spinorial multiplicity ladder and gaps]
\label{p1b:prop:v11-spinorial-low-modes}
For the sectors $\mathscr H_\Gamma$ defined before
Proposition~\ref{p1b:prop:new-twisted-exceptional-zeta}, the first
multiplicities of $\lambda_j=4j^2-2$ are
\[
\begin{array}{r|r|r|r|r}
 j&\lambda_j&m_j^{2T}&m_j^{2O}&m_j^{2I}\\\hline
 1&2&1&1&1\\
 2&14&0&0&0\\
 3&34&1&0&0\\
 4&62&2&1&0\\
 5&98&1&1&0\\
 6&142&2&1&1\\
 7&194&3&1&1\\
 8&254&2&1&0\\
 9&322&3&2&0\\
10&398&4&2&1\\
11&482&3&1&1\\
12&574&4&2&1
\end{array}.
\]
Each comparison sector has one singlet ground mode at $\lambda_1=2$,
while the first strictly higher occupied eigenvalues are
\[
\boxed{\lambda_{\rm first\ excited}^{2T}=34,\qquad
\lambda_{\rm first\ excited}^{2O}=62,\qquad
\lambda_{\rm first\ excited}^{2I}=142.}
\]
The generating functions in the preceding proposition determine the
entire ladder, not just the displayed rows.
\end{proposition}
\begin{proof}
Expand
\[
 G_\Gamma(z)=\frac{(1+z)(1+z^d)}{(1-z^a)(1-z^b)}-1
\]
for $(a,b,d)=(3,4,6),(4,6,9),(6,10,15)$ through $z^{12}$.
The first nonconstant even-degree invariants occur at degrees
$6,8,12$, respectively, which proves the three excited-level
identities without reliance on a numerical fit.
\end{proof}

\begin{theorem}[Exact intersection, union, and residual spectral zeta]
\label{p1b:thm:v11-spinorial-overlap-zeta}
Use the aligned embeddings $H=2T\subset2O,2I$ with
$2O\cap2I=H$.  In each fixed-charge eigenspace the spinorial
invariant subspaces obey
\[
\mathscr H_{2O}\subset\mathscr H_{2T},\quad
\mathscr H_{2I}\subset\mathscr H_{2T},\quad
\boxed{\mathscr H_{2O}\cap\mathscr H_{2I}=\mathbb C\psi_1,}
\]
where $\psi_1$ is the unique $SU(2)$-singlet at $\lambda_1=2$.
Denote by $\mathscr U$ the closed eigenspace-wise linear span of the
two subspaces inside $\mathscr H_{2T}$, and by
$\mathscr R=\mathscr H_{2T}\ominus\mathscr U$ its orthogonal
$L$-invariant complement.  Then, on their meromorphic domains,
\[
\boxed{Z_{\mathscr U}(s)=Z_{2O}(s)+Z_{2I}(s)-2^{-s},\qquad
Z_{\mathscr R}(s)=Z_{2T}(s)-Z_{2O}(s)-Z_{2I}(s)+2^{-s}.}
\]
In particular,
\[
\boxed{Z_{\mathscr U}(0)=-\frac4{45},\qquad
Z_{\mathscr R}(0)=\frac{11}{30}.}
\]
The individual heat traces have leading coefficients
\[
\operatorname{Tr}_{\mathscr H_\Gamma}e^{-tL_\Gamma}
=\frac1{|\Gamma|t}+Z_\Gamma(0)+O(t),\qquad t\downarrow0,
\]
while those of the union and residual sector have respective leading
coefficients $7/240$ and $1/80$.  The subtraction of the common
singlet prevents double counting of the intersection.
\end{theorem}
\begin{proof}
A vector fixed by both $2O$ and $2I$ is fixed by their generated
dense subgroup of $SU(2)$, by
Theorem~\ref{p1:thm:2O-2I-generate-SU2} and continuity.
At level $j$ the left $SU(2)$ representation is
$\operatorname{Sym}^{2j-1}V\otimes V\cong
\operatorname{Sym}^{2j}V\oplus\operatorname{Sym}^{2j-2}V$.
It has an $SU(2)$ singlet precisely for $j=1$, occurring once.
The dimension formula
$\dim(U+W)=\dim U+\dim W-\dim(U\cap W)$ at each eigenvalue
proves the zeta identities for sufficiently large $\operatorname{Re}s$;
meromorphic continuation proves them everywhere defined.  Insert the
three established rational values $5/18,7/18,47/90$ at $s=0$.
For the heat coefficient, the leading singular term of
$G_\Gamma(e^{-t})$ is $4/(ab)t^{-2}$ and
$2ab=|\Gamma|$ for the displayed triples; the corresponding
$4j^2-2$ heat ladder yields $1/(|\Gamma|t)$.
The union and residual leading coefficients follow by subtraction:
$1/48+1/120=7/240$ and
$1/24-1/48-1/120=1/80$.
\end{proof}

\begin{remark}[Zeta value versus zeta determinant]
\label{p1b:rem:v11-zeta-scaling}
$Z_\Gamma(0)$ is a scale coefficient and not by itself the full
regularized logarithmic determinant.  Whenever the determinant is
defined, its scaling identity is
\[
\log\det_{\zeta}(cL_\Gamma)
 =\log\det_{\zeta}L_\Gamma+Z_\Gamma(0)\log c,
\qquad\log\det_{\zeta}L_\Gamma=-Z_\Gamma'(0).
\]
A physical amplitude would additionally require a specified owner,
normalization, and relative determinant prescription.
\end{remark}

\subsection{All numerical combinations retained as conditional comparisons}
\label{p1b:subsec:v11-weighted-comparisons}

\begin{definition}[Two spectral and two shell-scalar comparison triples]
\label{p1b:def:v11-comparison-triples}
As a purely formal comparison, define the \emph{same-convention}
spinorial scale coefficients
\[
\boxed{c_{\Gamma}^{\rm spin}:=\frac{Z_\Gamma(0)}{2\sqrt2},\qquad
(c_{2T}^{\rm spin},c_{2O}^{\rm spin},c_{2I}^{\rm spin})
=\left(\frac5{36\sqrt2},\frac7{36\sqrt2},
\frac{47}{180\sqrt2}\right).}
\]
The factor $1/(2\sqrt2)$ is an \emph{explicitly chosen comparison
conversion}, not a derivation that the three twisted sectors enter a
common physical amplitude with that coefficient.  To compare with the
separately established original octahedral scalar/germ value, define
\[
\boxed{c^{\rm mixed}
 =\left(\frac5{36\sqrt2},\frac1{6\sqrt2},
\frac{47}{180\sqrt2}\right).}
\]
The mixed triple intentionally uses a \emph{different operator owner}
in its middle entry.  Finally, retain separately the shell atom triple
$\omega=(1/16,1/(6\sqrt2),\varphi/24)$ and the shell
contraction triple
$\delta^{\rm trial}=(0,1/(6\sqrt2),\varphi^2/12)$.
Their comparison is arithmetic, not an asserted equality of observables.
\end{definition}

\begin{proposition}[Exact arithmetic means and non-equivalent weight choices]
\label{p1b:prop:v11-complete-average-table}
For any row vector of weights $w=(w_T,w_O,w_I)$ of sum one let
$C_w(c)=\sum_\Gamma w_\Gamma c_\Gamma$.  The equal-weight
arithmetic means are
\[
\boxed{C_{(1,1,1)/3}(c^{\rm mixed})=\frac{17}{90\sqrt2},\qquad
 C_{(1,1,1)/3}(c^{\rm spin})=\frac{107}{540\sqrt2},}
\]
\[
\boxed{C_{(1,1,1)/3}(\omega)
 =\frac13\left(\frac1{16}+\frac1{6\sqrt2}+
 \frac\varphi{24}\right),\qquad
 C_{(1,1,1)/3}(\delta^{\rm trial})
 =\frac1{18\sqrt2}+\frac{\varphi^2}{36}.}
\]
The following table retains other exact \emph{chosen} weights and the
resulting values for both spectral coefficient triples:
\[
\begin{array}{c|c|c|c}
\text{weighting convention}&(w_T,w_O,w_I)&
 C_w(c^{\rm mixed})&C_w(c^{\rm spin})\\\hline
\text{equal names}&(1,1,1)/3&17/(90\sqrt2)&107/(540\sqrt2)\\
\text{group orders}&(1,2,5)/8&2/(9\sqrt2)&11/(48\sqrt2)\\
\text{inverse orders}&(10,5,2)/17&247/(1530\sqrt2)&173/(1020\sqrt2)\\
\text{uniform axis counts}&(10,15,6)/31&491/(2790\sqrt2)&1057/(5580\sqrt2)\\
\text{degree-11 axis cubature}&(81,256,125)/462&779/(4158\sqrt2)&281/(1386\sqrt2)
\end{array}.
\]
For the originally discussed mixed arithmetic mean, the purely
numerical scalar combination with the full Hopf residue is
\[
\boxed{4\pi^2+\frac{17}{90\sqrt2}\approx39.6119822186.}
\]
\end{proposition}
\begin{proof}
Each identity is the indicated finite weighted sum of the exact triples
in the preceding definition.  The group-order weights come from
$(24,48,120)$ and the inverse-order weights from
$(1/24,1/48,1/120)$, both normalized to total one.
The axis weights are the raw counts of the \emph{three $A_5$ axis
orbits}; the degree-eleven weights solve the two invariant harmonic
moment constraints.  For instance,
\[
\frac1{3}\left(\frac5{36\sqrt2}+
 \frac1{6\sqrt2}+\frac{47}{180\sqrt2}\right)
 =\frac{17}{90\sqrt2},
\]
whereas
\[
\frac{81}{462}\frac5{36\sqrt2}+
\frac{256}{462}\frac7{36\sqrt2}+
\frac{125}{462}\frac{47}{180\sqrt2}
 =\frac{281}{1386\sqrt2}.
\]
The other rows follow by the same exact rational arithmetic.
\end{proof}

\begin{remark}[Classification of these comparisons, including the failure]
\label{p1b:rem:v11-no-canonical-three-determinant-weight}
Equal group-name weighting, order weighting, and inverse-order
weighting are mathematically defined \emph{different choices}.
The $31$-axis and degree-eleven rows are valid as scalar arithmetic
only after an \emph{extra arbitrary pairing} of $A_5$ rotational
orbit labels (threefold, twofold, fivefold) with the three different
binary groups.  Such a pairing does not turn the octahedral group
into a subgroup of the icosahedral group, and it does not supply an
operator-valued map between their different spectral owners.
Furthermore the fixed-charge spinorial invariant spaces overlap, as
Theorem~\ref{p1b:thm:v11-spinorial-overlap-zeta} computes exactly.
Thus none of the tabulated arithmetic combinations is, on the current
mathematical premises, a uniquely required summed physical correction.
This no-selection conclusion is recorded alongside the positive
mathematical identities rather than suppressing them.
\end{remark}

\begin{remark}[Completeness boundary]
The present section records the entire comparison investigated here:
(1) unchanged $S^3$ contact residue and exact exceptional quotients;
(2) all three shell atoms, trial contractions, and failure ratios;
(3) all three unnormalized and normalized tight frames and their
convex-mixture theorem; (4) an explicit scalar-moment discrimination;
(5) 31-axis/projector and degree-eleven angular integration with its
first degree-twelve defect; (6) scalar fixed-charge obstruction,
spinorial multiplicities, gaps, zeta values and exact overlap;
(7) all arithmetical averaging conventions tested and their limits.
None is discarded for lack of a known physical application.
\end{remark}


\section{Finite-carrier and exceptional-support synthesis}
\label{p1:sec:conclusion}

Paper~I-A fixes a single normalization for the round Hopf-fibered sphere
\(S^3(\Rmom)\), including its metric, orientation, bundle connection, contact
and CR structures, the positive primitive lattice of the chosen Reeb flow,
harmonic analysis, spin and Clifford conventions, and the intrinsic representation bank.  Paper~I-B owns the
exceptional binary subgroup classification, embeddings and finite constructions.

The finite part defines the integer-level \(\operatorname{Spin}(4)\)-carrier family
\(S^3_m\), specializes the declared family at \(m=1\), and proves the
resulting consequences within that skeleton.  The unrestricted compatible
equivariant coordinate actions form the explicit three-modulus atlas
\(\mathfrak M_{\rm coord}/\!\sim\).  Every point of that cone carries the
unnormalized protected relation with
\(\rho_{\mathsf C}=2\varrho_{\rm sp}\ge0\); inverse-normalized gamma matrices,
nonzero protected spectra, and Clifford--Fierz data apply only when
\(\rho_{\mathsf C}>0\).  The canonical retained smooth model
exports the intrinsic fixed-radius class \(\mathsf C_D\), its exact quadratic
fingerprint, the canonical Levi--Civita square and scalar spectral projector,
and the normalized-constant/Haar identity.  Independently, the ambient
metric-cone common-symbol branch exports
\(\mathsf C_{\rm cone}=\mathsf C_{\rm PI}(\rho_{\rm prot}=\Rmom)\).
Their different fingerprints prove
\(\mathsf C_D\ne\mathsf C_{\rm cone}\).
Theorem~\ref{p1:thm:primitive-round-symbol-ancestry-exhaustion} proves that
these are exactly the two outputs of the frozen primitive round-symbol
grammar.  Its scope is sharp: admitting an independent normal weight or any
other enlarged ancestry data creates a different grammar and restores further
possibilities.  These are compression theorems, not a claim that the finite
carrier is a full \(\Cl_4\)-module.  The exceptional
subgroup classification supplies the distinct unmarked support-role outputs
\(\mathsf{Supp}_{2T}\), \(\mathsf{Supp}_{2O}\), and
\(\mathsf{Supp}_{2I}\), rather than one selected support datum.
On any Hurwitz-core representative, the trimer zero extensions,
Bose--Mesner spectral embeddings, and operator corners give an explicit
internal realization in \(\CC^{2T}\).  The constant-plus-linear scalar sector
also has an exact five-dimensional evaluation/Casimir bridge. The natural
left-\(2T\) obstruction still excludes a scalar-valued injective extension
from all of \(\mathcal H_1\) into the left regular target. In contrast,
Theorem~\ref{p1:thm:v203-full-carrier-sampling} retains the existing form and
spinor fibres and gives an exact rank-fifteen sampled realization; its
intrinsic coordinate compiler is Corollary~\ref{p1:cor:v203-full-symbol-transport}.

These finite constructions retain their stated hypotheses and
conditional comparison domains.  None selects a physical exceptional support
without a separate source/observable identification.

\section{Complete quadratic four-generator content}
\label{p1:sec:v189-complete-quadratic-content}
Let $V\cong\mathbb R^4$ carry the already fixed invariant inner product and
orientation of the four-generator parent geometry.  Before any Clifford
quotient is taken, the complete quadratic representation is
\[
 \boxed{V\otimes V=\mathbb R g\oplus S^2_0V\oplus\Lambda^2V,}
 \qquad 16=1+9+6.
\]
The three summands are respectively the trace, symmetric-traceless and
antisymmetric channels, with the displayed representation types.

\begin{proposition}[Trace-free quadratic harmonics]
\label{p1:prop:v189-tracefree-quadratic-harmonics}
For $Q\in S^2_0(\mathbb R^4)$ define
\[
 f_Q(x)=Q_{ab}x^ax^b.
\]
Then $\Delta_{\mathbb R^4}f_Q=2\operatorname{tr}Q=0$, and restriction to the
unit round three-sphere gives a $\operatorname{Spin}(4)$-equivariant isomorphism
\[
 \boxed{S^2_0(\mathbb R^4)_\mathbb C\cong E_2(S^3),}
\]
with both sides carrying the $(1,1)$ representation and having dimension
nine.
\end{proposition}


\begin{corollary}[Quadratic first-harmonic product]
\label{p1:cor:v191-sym2-e1-e0e2}
Let $E_1(S^3)$ be the degree-one scalar harmonic carrier on the unit round
three-sphere.  Pointwise multiplication followed by harmonic decomposition
defines a $\operatorname{Spin}(4)$-equivariant map
\[
 \mu_2:\operatorname{Sym}^2E_1(S^3)
 \longrightarrow E_0(S^3)\oplus E_2(S^3).
\]
It is an isomorphism.  Equivalently,
\[
 \boxed{\operatorname{Sym}^2E_1(S^3)\cong E_0(S^3)\oplus E_2(S^3),}
 \qquad 10=1+9.
\]
The scalar summand is the trace $\sum_a x_a^2=1$, while the trace-free
quadratic summand is the harmonic space of
Proposition~\ref{p1:prop:v189-tracefree-quadratic-harmonics}.
\end{corollary}

\begin{proof}
Every symmetric quadratic polynomial on $\mathbb R^4$ decomposes uniquely as
its trace multiple of $|x|^2$ plus a trace-free quadratic polynomial.  On the
unit sphere the first term restricts to $E_0$, while the second restricts,
by Proposition~\ref{p1:prop:v189-tracefree-quadratic-harmonics}, to $E_2$.
The dimensions are $\dim\operatorname{Sym}^2\mathbb R^4=10$ and
$\dim(E_0\oplus E_2)=1+9=10$, so the equivariant multiplication/decomposition
map is bijective.
\end{proof}


\begin{corollary}[Level-one carrier boundary]
\label{p1:cor:v189-level-one-shear-exclusion}
The $(1,1)$ symmetric-traceless branch is not a sector of the fixed
level-one carrier $\mathcal H_1$. It belongs to the displayed quadratic
representation and scalar harmonic space.
\end{corollary}

\section{Rooted tetrahedral complement and antipodal representation}
\label{p1:sec:v224-icosian-comparison}
\begin{proposition}[Four free complement orbits]
\label{p1:prop:v224-icosian-complement}
Let the rooted binary tetrahedral subgroup be $H=2T\subset I=2I$
and let $Y=I\setminus H$. Then $|Y|=96$, left multiplication by
$H$ preserves $Y$ and acts freely, and
\[
 \mathbb C[Y]\!\downarrow_H\cong4\operatorname{Reg}(2T).
\]
The antipodal quotient $\bar Y=Y/\{\pm1\}$ has 48 elements and,
for $\bar H=H/\{\pm1\}\cong A_4$,
\[
 \mathbb C[\bar Y]\cong4\operatorname{Reg}(A_4).
\]
These are rooted permutation representations, with fixed-space dimension
four for their respective left actions.
\end{proposition}
\begin{proof}
If $hy\in H$ for $h\in H$, then $y\in H$, a contradiction.
If $hy=y$, cancellation gives $h=1$. Thus all orbits have size
24, and $96/24=4$. The central element $-1$ pairs opposite vertices.
If $[hy]=[y]$, then $h=\pm1$, so the quotient action is free;
its four orbits have size 12. A free transitive permutation module is
the regular module. Each orbit contributes one constant invariant.
\end{proof}
For the specified comparison extension $Q=\mathbf1\oplus\mathbf3$,
$G=\mathbf1\oplus\mathbf1'\oplus\mathbf1''$ and an inert
four-dimensional $W$, character multiplication gives
\[
 Q\otimes G=\mathbf1\oplus\mathbf1'\oplus\mathbf1''
                  \oplus3\mathbf3=\operatorname{Reg}(A_4),
 \qquad W\otimes Q\otimes G\cong\mathbb C[\bar Y].
\]
This proves a representation comparison under those extended actions.
It supplies no preferred intertwiner of a complete gauge, Dirac or source
word: the multiplicities $4,4,4,12$ give an equivariant endomorphism
algebra of dimension $3\cdot4^2+12^2=192$. The subgroup $H$ is not
normal in $I$, so this left complement representation is not an
$I$-invariant subrepresentation of its full regular representation.
The comparison does not turn the antipodal permutation action into a
spinorial central sign. Its matter application is stated in
Paper~V-A's geometric particle/source comparison.
An isometric intertwiner preserves the ordinary trace and cannot by
itself rescale an action by $1/96$.



\part{Exceptional quotient spectra, spin geometry and heat invariants}
\section{One invariant-series input for several different spectral operators}
The spectral operators in this part act on the stated round spherical quotient
or explicitly defined unitary flat bundle, with the specified spin lift when
needed.  They are not the primitive full-$S^3$ fixed-charge determinant of
Paper~III-C.
\label{p1b:sec:v13-owners}

Fix the unit round group manifold $S^3=SU(2)$ and its left quotient
$M_\Gamma=\Gamma\backslash S^3$ for
$\Gamma\in\{2T,2O,2I\}$.  Radius $R>0$ is reinstated by the metric
$R^2g_{\rm round}$.  This part constructs \emph{new mathematical
observables} from the quotient and does not modify the intrinsic,
unit-Chern full-host identity
$\int_{S^3}\eta_0\wedge d\eta_0=4\pi^2$ of Paper~I-A.
It uses the Dirac character formulae of B\"ar
\cite{Baer1996DiracSpaceForms,Baer2000SpinStructure}, the spherical
form decomposition of Ikeda--Taniguchi
\cite{ikedaTaniguchi1978eigenforms}, and the twisted spectral-action
results of \cite{CacicMarcolliTeh2014CosmicTopology}.

Let $V$ be the complex defining two-dimensional $SU(2)$ module and
$W$ a finite-dimensional unitary $\Gamma$-module. Define
\begin{equation}
\label{p1b:eq:v13-isotypic-series}
 b_n^W:=\dim(\operatorname{Sym}^n V\otimes W)^\Gamma,
 \qquad B_W(z):=\sum_{n\geq0}b_n^Wz^n
 =\frac1{|\Gamma|}\sum_{g\in\Gamma}
 \frac{\chi_W(g)}{1-\chi_V(g)z+z^2}.
\end{equation}
The flat bundle convention is the diagonal quotient of the trivial
bundle on $S^3$ carrying $W$; using the dual convention replaces $W$
by $W^*$.  This matters for complex non-real characters but not for the
trivial and sign characters used explicitly below.  For $W=\mathbf1$,
write $A_\Gamma=B_{\mathbf1}$ and $a_n^\Gamma=b_n^{\mathbf1}$.
The previously established Molien series are
\begin{equation}
\label{p1b:eq:v13-molien-three}
 A_{2T}=\frac{1+z^{12}}{(1-z^6)(1-z^8)},\qquad
 A_{2O}=\frac{1+z^{18}}{(1-z^8)(1-z^{12})},\qquad
 A_{2I}=\frac{1+z^{30}}{(1-z^{12})(1-z^{20})}.
\end{equation}
The same $B_W$ can feed scalar, Dirac, Hodge, and torsion calculations,
but their operators, multiplicities and regularizations are different.

\begin{proposition}[The second octahedral spin lift has a difference-of-invariants series]
\label{p1b:prop:v13-octa-sign-molien}
Let $\sigma:2O\to\{\pm1\}$ be the character with kernel $2T$.
Then, exactly,
\begin{equation}
\label{p1b:eq:v13-octa-sign-series}
 B_\sigma(z)=A_{2T}(z)-A_{2O}(z),
 \qquad B_\sigma(z)=z^6+O(z^{12}).
\end{equation}
The two spin structures on $M_{2O}$ are represented by the natural
spin lift and its product with $\sigma$.
\end{proposition}
\begin{proof}
The index-two normal inclusion $2T\triangleleft2O$ gives
$\operatorname{Ind}_{2T}^{2O}\mathbf1\cong\mathbf1\oplus\sigma$.
By Frobenius reciprocity,
$\dim(\operatorname{Sym}^nV)^{2T}
 =\dim(\operatorname{Sym}^nV)^{2O}
 +\dim(\operatorname{Sym}^nV\otimes\sigma)^{2O}$.
Taking generating functions proves the identity.  The first nonzero
coefficients follow by expanding the three rational series in
\eqref{p1b:eq:v13-molien-three}.  A spin lift can be changed by
exactly a $\{\pm1\}$ character; the existing homology/spin calculation
shows there are precisely two such choices for $2O$.
\end{proof}

\section{Full signed Dirac spectrum and exact eta invariants}
\label{p1b:sec:v13-dirac}

\begin{theorem}[Signed Dirac multiplicities from one Molien series]
\label{p1b:thm:v13-dirac-generating}
Choose the canonical lift $\epsilon(g)=(g,1)$ in
$\operatorname{Spin}(4)\cong SU(2)_L\times SU(2)_R$ for the left
$\Gamma$ action, and choose the spin orientation so that its invariant
Killing spinors have \emph{positive} Dirac eigenvalue.  For the
round Dirac operator $D_{\Gamma,W}$ on the spin bundle twisted by the
flat unitary module $W$, its possible eigenvalues are
$\pm(k+\tfrac32)/R$, $k\geq0$, with generating functions
\begin{equation}
\label{p1b:eq:v13-dirac-poincare}
 \boxed{F_+^W(z)=2B_W(z)+zB_W'(z),\qquad
 F_-^W(z)=B_W'(z).}
\end{equation}
Equivalently the actual multiplicities are
\begin{equation}
\label{p1b:eq:v13-dirac-multiplicities}
 \boxed{m_{+,k}^W=(k+2)b_k^W,\qquad
 m_{-,k}^W=(k+1)b_{k+1}^W.}
\end{equation}
A possible eigenvalue is absent if its multiplicity vanishes.  In
particular, no $D_{\Gamma,W}$ on these positive-curvature round
quotients has a zero mode.
\end{theorem}
\begin{proof}
The two $\operatorname{Spin}(4)$ eigenspace modules on the round
three-sphere at $\pm(k+3/2)/R$ are respectively
$\operatorname{Sym}^{k}V_L\otimes\operatorname{Sym}^{k+1}V_R$
and
$\operatorname{Sym}^{k+1}V_L\otimes\operatorname{Sym}^{k}V_R$,
with the displayed choice of sign.  Their invariant dimensions after
left restriction and tensoring by $W$ are $(k+2)b_k^W$ and
$(k+1)b_{k+1}^W$.  This is the $SU(2)_L\times SU(2)_R$ specialization
of B\"ar's spherical-space-form character theorem
\cite{Baer1996DiracSpaceForms,Baer2000SpinStructure}.
Summation proves \eqref{p1b:eq:v13-dirac-poincare}.
The Lichnerowicz identity
$D^2=\nabla^*\nabla+\mathrm{Scal}/4$ with
$\mathrm{Scal}=6/R^2>0$ rules out zero modes independently.
Changing orientation exchanges the two signs and sends the eta
invariant below to its negative.
\end{proof}

\begin{corollary}[Exact earliest signed spinor levels]
\label{p1b:cor:v13-first-dirac-levels}
For $W=\mathbf1$ the common natural ground level is
$+3/(2R)$ with multiplicity two.  The first \emph{additional}
negative and positive levels are, respectively,
\[
\begin{array}{c|c|c|c}
\Gamma & \text{least positive invariant degree }d_\Gamma
  & \text{first negative }(R\lambda,\mathrm{mult})
  & \text{next positive }(R\lambda,\mathrm{mult})\\\hline
2T&6&(-13/2,6)&(15/2,8)\\
2O&8&(-17/2,8)&(19/2,10)\\
2I&12&(-25/2,12)&(27/2,14)
\end{array}
\]
For the \emph{other} $2O$ spin structure the $+3/(2R)$ ground
level is absent; the first negative and positive levels are
$-13/(2R)$ with multiplicity six and $+15/(2R)$ with multiplicity
eight. These levels are distinct from both the primitive Paper III-C
fixed-charge ladder and the $2T/2O/2I$ spinorial-sector trial
spectra constructed earlier in this paper.
\end{corollary}
\begin{proof}
In \eqref{p1b:eq:v13-molien-three}, the smallest positive exponent
with a nonzero coefficient is $d_\Gamma=6,8,12$.
Apply \eqref{p1b:eq:v13-dirac-multiplicities} to $a_0=1$ and
$a_{d_\Gamma}=1$.  For the second $2O$ spin lift use
\eqref{p1b:eq:v13-octa-sign-series} instead.
\end{proof}

\begin{theorem}[Finite character formula for all flat-twisted Dirac eta invariants]
\label{p1b:thm:v13-eta-character}
With the orientation and canonical lift of Theorem~
\ref{p1b:thm:v13-dirac-generating}, and for any unitary
$\Gamma$-module $W$, the Atiyah--Patodi--Singer eta invariant is
\begin{equation}
\label{p1b:eq:v13-eta-character}
 \boxed{\eta_{D_{\Gamma,W}}(0)
   =\frac1{|\Gamma|}\sum_{g\in\Gamma\setminus\{1\}}
      \frac{\chi_W(g)}{1-\operatorname{Re}g}.}
\end{equation}
The right side is real, and is invariant under scaling the round
radius. The relative eta quantity for a rank-$N$ module is exactly
\begin{equation}
\label{p1b:eq:v13-eta-rho}
 \rho_W(D_\Gamma)
 :=\eta_{D_{\Gamma,W}}(0)-N\eta_{D_{\Gamma,\mathbf1}}(0)
 =\frac1{|\Gamma|}\sum_{g\ne1}
   \frac{\chi_W(g)-N}{1-\operatorname{Re}g}.
\end{equation}
This provides an explicit finite expression for \emph{every}
irreducible exceptional flat representation, not merely for one-dimensional
characters.  It is an eta/rho invariant of this declared operator and
round metric, not a universal spin-structure-independent scalar.
\end{theorem}
\begin{proof}
In B\"ar's general eta formula for spherical spin space forms,
\[
 \eta(0)=2|\Gamma|^{-1}\sum_{g\ne1}
 \frac{\chi_W(g)(\chi^-(\epsilon(g))-\chi^+(\epsilon(g)))}{\det(I_4-g)}
\]
\cite{Baer2000SpinStructure,
BaerSchopka2003DiracDeterminant}.  For the selected left lift,
$\chi^-=2$, $\chi^+=\chi_V(g)=2\operatorname{Re}g$.
Left quaternion multiplication on $\mathbb R^4$ has
$\det(I_4-g)=(2-\chi_V(g))^2$.  Cancelling gives
\eqref{p1b:eq:v13-eta-character}.  The $g,g^{-1}$ terms are
complex conjugates for unitary $W$, so the sum is real.
No eigenvalue is zero by the preceding theorem, and a positive
rescaling of the spectrum leaves eta at $s=0$ unchanged.
Subtracting $N$ times the trivial-character formula proves
\eqref{p1b:eq:v13-eta-rho}.
\end{proof}

\begin{corollary}[Exact exceptional eta fractions and the $2O$ spin difference]
\label{p1b:cor:v13-exact-eta-values}
For the natural spin structure, the three eta invariants are
\begin{equation}
\label{p1b:eq:v13-eta-fractions}
 \boxed{\eta_{2T}=\frac{167}{144},\qquad
 \eta_{2O}=\frac{383}{288},\qquad
 \eta_{2I}=\frac{1079}{720}.}
\end{equation}
For the second spin structure on $M_{2O}$,
\begin{equation}
\label{p1b:eq:v13-eta-octa-second}
 \boxed{\eta_{2O,\sigma}=-\frac{49}{288},\qquad
   \eta_{2O}-\eta_{2O,\sigma}=\frac32.}
\end{equation}
All values reverse sign if the chosen orientation convention is
reversed.
\end{corollary}
\begin{proof}
The exact element-height census, with the identity $\operatorname{Re}1=1$
excluded from the sum, is
\[
\begin{array}{c|c|c|c|c|c}
 \Gamma&\operatorname{Re}(-1)=-1&0&\pm\frac12
 &\pm\frac1{\sqrt2}
 &\pm\frac{\varphi}{2},\ \pm\frac{\varphi-1}{2}\\\hline
 2T&1&6&8\text{ each}&0&0\\
 2O&1&18&8\text{ each}&6\text{ each}&0\\
 2I&1&30&20\text{ each}&0&12\text{ each for both pairs}
\end{array}
\]
where $\varphi=(1+\sqrt5)/2$. These multiplicities sum to
$23,47,119$, respectively, after excluding the identity.  Insert
them in \eqref{p1b:eq:v13-eta-character}.
For example,
\[
 \eta_{2T}=\frac1{24}
 \left(\frac12+6+\frac8{1-1/2}+\frac8{1+1/2}\right)
 =\frac{167}{144}.
\]
For $2O$, the $24$ elements outside $2T$ contribute
$12+6/(1-1/\sqrt2)+6/(1+1/\sqrt2)=36$ to the numerator.
Consequently $\eta_{2O}=(24\eta_{2T}+36)/48=383/288$,
while the sign twist reverses precisely those contributions:
$\eta_{2O,\sigma}=(24\eta_{2T}-36)/48=-49/288$.
For $2I$, pairing each $\pm$ height and using
$\varphi/2=\cos(\pi/5)$ and $(\varphi-1)/2=\cos(2\pi/5)$
produces
\[
 \eta_{2I}=\frac1{120}
 \left(\frac12+30+40+\frac{40}{3}+96\right)
 =\frac{1079}{720}.
\]
The last computation can be independently verified without
regularizing an infinite eigenvalue sum: it uses a finite class
census and rationalization of the two golden height pairs.
\end{proof}

\begin{theorem}[The complete Dirac eta vector is an affine-McKay Green column]
\label{p1b:thm:v13-eta-cartangreen}
Let $\{\rho_i\}_{i=0}^{r-1}$ be the irreducible complex
$\Gamma$-representations with $\rho_0=\mathbf1$, and let
$d_i=\dim\rho_i$.  Write the McKay tensor-product matrix as
$V\otimes\rho_i\cong\bigoplus_j A_{ij}\rho_j$, and set
$C_\Gamma=2I-A_\Gamma$, the symmetric affine $\widetilde E_{6,7,8}$
Cartan matrix in the corresponding case.  The vector of \emph{all}
round, naturally lifted, flat-twisted Dirac eta invariants
$\boldsymbol\eta=(\eta_{D_{\Gamma,\rho_i}}(0))_i$
is the unique solution to
\begin{equation}
\label{p1b:eq:v13-eta-mckay-system}
 \boxed{C_\Gamma\boldsymbol\eta
 =2e_0-\frac{2}{|\Gamma|}\boldsymbol d,
 \qquad \boldsymbol d^{\mathsf T}\boldsymbol\eta=0.}
\end{equation}
Equivalently,
\begin{equation}
\label{p1b:eq:v13-eta-mckay-pseudoinverse}
 \boxed{\boldsymbol\eta=2C_\Gamma^+e_0,\qquad
 \eta_{D_{\Gamma,\mathbf1}}(0)=2(C_\Gamma^+)_{00},}
\end{equation}
where $C_\Gamma^+$ denotes the Moore--Penrose pseudoinverse.
The equivalence applies to these finite $SU(2)$ subgroups and the
explicit round Dirac/flat-bundle spectral owner, not to an
arbitrary positive operator or the full-host Born--Hopf determinant.
\end{theorem}
\begin{proof}
Define the inversion-invariant real central function on $\Gamma$
\[
 K(g)=\begin{cases}
 2/(2-\chi_V(g)),&g\ne1,\\
 0,&g=1.
 \end{cases}
\]
Since $K(g)=K(g^{-1})$, its irreducible-character Fourier
coefficient is
$|\Gamma|^{-1}\sum_g K(g)\chi_i(g)
 =\eta_{D_{\Gamma,\rho_i}}(0)$ by
\eqref{p1b:eq:v13-eta-character}. Thus
$K=\sum_i\eta_i\chi_i$ in the space of central class functions.
The identity $(2-\chi_V)K=2(1-\delta_1)$ holds pointwise,
including at the identity.  Using
$\chi_V\chi_i=\sum_jA_{ij}\chi_j$ and the regular-character identity
$|\Gamma|\delta_1=\sum_j d_j\chi_j$ yields the first equation
in \eqref{p1b:eq:v13-eta-mckay-system}.  Evaluating the Fourier
expansion at $g=1$ gives
$0=K(1)=\sum_i d_i\eta_i$.
The connected affine ADE matrix $C_\Gamma$ is symmetric positive
semidefinite with one-dimensional kernel spanned by
$\boldsymbol d$, and Burnside's identity gives
$\boldsymbol d^{\mathsf T}\boldsymbol d=|\Gamma|$.
The first equation is consistent and determines $\boldsymbol\eta$
modulo $\boldsymbol d$; the second fixes its component in that
kernel to zero. By the defining property of the Moore--Penrose
inverse, $C_\Gamma^+\boldsymbol d=0$ and
$C_\Gamma C_\Gamma^+e_0=e_0-\boldsymbol d/|\Gamma|$.
This gives \eqref{p1b:eq:v13-eta-mckay-pseudoinverse} exactly.
\end{proof}

\begin{corollary}[Complete exact eta tables on $\widetilde E_6$, $\widetilde E_7$, $\widetilde E_8$]
\label{p1b:cor:v13-all-irrep-etas}
The following vertex orderings fix the labels and make the
pseudoinverse formula independently reproducible. For $2T$, use
\[
 d=(1,1,1,2,2,2,3),\quad\text{edges }\{03,14,25,36,46,56\};
\]
for $2O$ use
\[
 d=(1,2,3,4,3,2,1,2),\quad\text{edges }\{01,12,23,34,45,56,37\};
\]
and for $2I$ use
\[
 d=(1,2,3,4,5,6,4,2,3),\quad\text{edges }\{01,12,23,34,45,56,67,58\}.
\]
Vertex $0$ always denotes the trivial representation, and the
unoriented edge $ij$ contributes $-1$ to entries $(i,j),(j,i)$ of
$C_\Gamma$. With this ordering, the exact eta vectors are
\[
\begin{aligned}
\boldsymbol\eta_{2T}&=
 (\tfrac{167}{144},-\tfrac{25}{144},-\tfrac{25}{144},
   \tfrac{29}{72},-\tfrac{19}{72},-\tfrac{19}{72},-\tfrac{3}{16}),\\[2pt]
\boldsymbol\eta_{2O}&=
 (\tfrac{383}{288},\tfrac{101}{144},\tfrac{5}{32},
   -\tfrac{19}{72},-\tfrac{11}{32},-\tfrac{43}{144},
   -\tfrac{49}{288},-\tfrac{25}{144}),\\[2pt]
\boldsymbol\eta_{2I}&=
 (\tfrac{1079}{720},\tfrac{73}{72},\tfrac{9}{16},
   \tfrac{29}{180},-\tfrac{25}{144},-\tfrac{17}{40},
   -\tfrac{61}{180},-\tfrac{67}{360},-\tfrac{19}{80}).
\end{aligned}
\]
In particular the trivial-representation entries agree with the
independent quaternionic height sums in Corollary~
\ref{p1b:cor:v13-exact-eta-values}; the second one-dimensional
character of $2O$ occupies vertex $6$ and gives $-49/288$.
The entire vector is determined by affine Cartan/mark data, without
needing to re-enumerate the individual finite-group elements.
\end{corollary}
\begin{proof}
Build the indicated matrices $C=2I-A$ from their edge lists.
The displayed dimension vectors satisfy $C\boldsymbol d=0$ and
$\boldsymbol d^{\mathsf T}\boldsymbol d=24,48,120$.
Each displayed rational vector satisfies both equations in
\eqref{p1b:eq:v13-eta-mckay-system}; uniqueness from the
preceding theorem proves the result.  These are exact rational
linear-algebra certificates, not fitted numerical spectra.
\end{proof}

\begin{corollary}[All flat Dirac rho invariants from the same affine Green column]
\label{p1b:cor:v13-all-rho-green}
For every irreducible $\rho_i$ of dimension $d_i$, the normalized
relative eta value is
\begin{equation}
\label{p1b:eq:v13-rho-green}
 \boxed{\rho_{\rho_i}(D_\Gamma)=
 2\bigl[(C_\Gamma^+)_{i0}
       -d_i(C_\Gamma^+)_{00}\bigr].}
\end{equation}
In particular, all $r-1$ nontrivial relative eta numbers follow
from a single affine-Cartan Green column.  No claim is made that
these round-metric spectral values select a Born--Hopf coupling.
\end{corollary}
\begin{proof}
Subtract $d_i$ times the trivial eta entry in
\eqref{p1b:eq:v13-eta-mckay-pseudoinverse} and compare
with \eqref{p1b:eq:v13-eta-rho}.
\end{proof}

\begin{remark}[What these eta data do and do not determine]
\label{p1b:rem:v13-eta-owner}
The values above measure the spectral \emph{asymmetry} of the round,
quotient Dirac operator; none is an additive correction to the
full-host $4\pi^2$ Hopf residue. The eta invariant depends on the
orientation, spin lift, flat twist and generally on the Riemannian
metric.  The difference $3/2$ is exact for the specified round
$2O$ quotient; a generic metric-variation invariance is not claimed.
For a conventionally chosen spectral cut, the phase of a zeta
regularized Dirac determinant may involve $\exp(\pm i\pi\eta/2)$;
it must not be identified with the positive determinant of $D^2$.
\end{remark}

\section{The complete round Hodge spectrum of zero- through three-forms}
\label{p1b:sec:v13-hodge}

\begin{theorem}[Exact flat-twisted scalar and coexact one-form spectra]
\label{p1b:thm:v13-hodge-spectra}
For the flat unitary $W$-bundle defined above, the Hodge Laplacian
on $M_\Gamma$ has the following positive scalar spectrum:
\begin{equation}
\label{p1b:eq:v13-hodge-scalar}
 \lambda_{n,0}=\frac{n(n+2)}{R^2},\qquad
 d_{n,0}=(n+1)b_n^W,\qquad n\ge1.
\end{equation}
Its kernel in degree zero has dimension $b_0^W$.
On \emph{coexact} one-forms, the positive eigenvalues and
multiplicities are
\begin{equation}
\label{p1b:eq:v13-hodge-coexact}
 \boxed{\lambda_{k,1}^{\rm coex}=\frac{(k+1)^2}{R^2},\qquad
 d_{k,1}^{\rm coex}=k\,b_{k+1}^W+(k+2)b_{k-1}^W,
 \quad k\ge1.}
\end{equation}
Exact one-forms duplicate the positive scalar eigenvalues and
multiplicities in \eqref{p1b:eq:v13-hodge-scalar}; Hodge star then
transfers the complete one-form spectrum to two-forms and the scalar
spectrum to three-forms. Thus these two displayed series determine
the \emph{entire} Hodge spectrum in degrees $p=0,1,2,3$.
\end{theorem}
\begin{proof}
The scalar Peter--Weyl calculation is Theorem~
\ref{p1b:thm:v12-full-scalar-spectrum} with the $W$ isotypic
projection.  Ikeda--Taniguchi's round-sphere eigenform decomposition
\cite{ikedaTaniguchi1978eigenforms} gives the two coexact
$SU(2)_L\times SU(2)_R$ summands at eigenvalue
$(k+1)^2/R^2$:
\[
 \operatorname{Sym}^{k+1}V_L\otimes\operatorname{Sym}^{k-1}V_R,
 \qquad
 \operatorname{Sym}^{k-1}V_L\otimes\operatorname{Sym}^{k+1}V_R.
\]
Restricting the left tensor factor to $\Gamma$, applying the
$W$ projector and taking the right dimensions $k$ and $k+2$
gives \eqref{p1b:eq:v13-hodge-coexact}.  A nonzero scalar eigenform
$f$ maps injectively and isometrically up to its eigenvalue to the
exact form $df$, with the same eigenvalue; Hodge decomposition gives
the rest. On an oriented three-manifold, Hodge star interchanges
degrees $p$ and $3-p$ and commutes with the Laplacian.
\end{proof}

\begin{corollary}[The shared coexact ground and the first exceptional splitting]
\label{p1b:cor:v13-hodge-gap-table}
On each untwisted exceptional quotient, the \emph{first} coexact
one-form eigenvalue is $4/R^2$ with multiplicity three, despite the
very different scalar first positive eigenvalues. The next coexact
one-form levels are
\begin{equation}
\label{p1b:eq:v13-hodge-gaps}
 \begin{array}{c|c|c|c}
\Gamma & \text{first scalar positive }R^2\lambda_{0,+}
& \text{coexact ground }(R^2\lambda,d)
& \text{next coexact }(R^2\lambda,d)\\\hline
2T&48&(4,3)&(36,5)\\
2O&80&(4,3)&(64,7)\\
2I&168&(4,3)&(144,11)
\end{array}
\end{equation}
Here the subscript $0,+$ means the first \emph{positive} scalar
Laplacian eigenvalue; no scalar zero mode is being ignored.
\end{corollary}
\begin{proof}
All three groups contain $-1$, and their least positive invariant
symmetric-power degree is $d_\Gamma=6,8,12$.  At $k=1$,
\eqref{p1b:eq:v13-hodge-coexact} gives
$1\cdot a_2^\Gamma+3a_0^\Gamma=3$.
For $1<k<d_\Gamma-1$, both $a_{k-1}$ and $a_{k+1}$
vanish.  At $k=d_\Gamma-1$, the only contribution is
$k a_{d_\Gamma}=d_\Gamma-1$.  The eigenvalue is
$(k+1)^2=d_\Gamma^2$.  The scalar eigenvalue comes from
$n=d_\Gamma$ and equals $d_\Gamma(d_\Gamma+2)$.
\end{proof}

\begin{proposition}[Twisted cohomology and the exact zero-zeta distinctions]
\label{p1b:prop:v13-cohomology-zeta}
For any unitary $W$,
\begin{equation}
\label{p1b:eq:v13-twisted-cohomology}
 H^p(M_\Gamma;E_W)\cong
 \begin{cases}
 W^\Gamma,&p=0,3,\\
 0,&p=1,2.
 \end{cases}
\end{equation}
In particular every nontrivial irreducible $W$ gives an acyclic
flat bundle.  For the \emph{untwisted} quotient, the spectral zeta
functions with zero modes omitted obey
\begin{equation}
\label{p1b:eq:v13-zeta-zero}
 \boxed{\zeta_{\Delta_0,+}(0)=-1,\qquad
 \zeta_{\Delta_1,\mathrm{coex}}(0)=1.}
\end{equation}
For nontrivial irreducible $W$, both corresponding values are zero.
These statements concern three-dimensional Hodge Laplacians, not the
separate Paper III-C fixed-charge spectral zeta value $1/3$.
\end{proposition}
\begin{proof}
The finite covering $S^3\to M_\Gamma$ identifies twisted de Rham
cohomology with $\Gamma$-invariants of
$H^p(S^3;\mathbb C)\otimes W$.  The two nonzero sphere cohomology
groups are trivial $\Gamma$ modules in degrees zero and three;
thus the claimed formula follows by averaging over the finite group.
For any positive Laplace-type operator on a closed odd-dimensional
manifold, the heat expansion has no $t^0$ bulk term, so
$\zeta_{\Delta_p,+}(0)=-\dim\ker\Delta_p$. The ordinary quotient
has $b_0=1$, $b_1=0$. The positive exact-one-form spectrum equals
the positive scalar spectrum. Consequently
$\zeta_{\Delta_1,\mathrm{coex}}(0)
 =\zeta_{\Delta_1,+}(0)-\zeta_{\Delta_0,+}(0)=1$.
The same argument for an acyclic irreducible $W$ gives zero in
both sectors.
\end{proof}

\section{Ray--Singer torsion as a complete exceptional spectral observable}
\label{p1b:sec:v13-torsion}

\begin{theorem}[Exact spectral representation of acyclic analytic torsion]
\label{p1b:thm:v13-torsion-algorithm}
Let $W$ be any nontrivial irreducible unitary $\Gamma$-module.  Define
meromorphically continued zeta functions from the two exact ladders
\begin{align}
 \zeta_{0,W}(s)
 &=\sum_{n\ge1}(n+1)b_n^W
   \left(\frac{n(n+2)}{R^2}\right)^{-s},
 \label{p1b:eq:v13-zeta-scalar}\\
 \zeta_{c,W}(s)
 &=\sum_{k\ge1}
   \bigl(kb_{k+1}^W+(k+2)b_{k-1}^W\bigr)
   \left(\frac{(k+1)^2}{R^2}\right)^{-s}.
 \label{p1b:eq:v13-zeta-coexact}
\end{align}
The positive Ray--Singer analytic torsion has the exact expression
\begin{equation}
\label{p1b:eq:v13-torsion-formula}
 \boxed{\log T(M_\Gamma,E_W)
   =\frac12\Bigl[\zeta_{c,W}'(0)-2\zeta_{0,W}'(0)\Bigr].}
\end{equation}
It is a well-defined positive, metric-independent invariant for
this acyclic unitary local system, and equals its Reidemeister torsion
(with the standard compatible torsion conventions) by the
Cheeger--M\"uller theorem \cite{Mueller1978AnalyticTorsion}.
This is an explicit \emph{algorithmic spectral determination} by
finite-group characters and zeta continuation; no closed-form
numeric torsion values for every irrep are being asserted here.
\end{theorem}
\begin{proof}
Acyclicity follows from Proposition~
\ref{p1b:prop:v13-cohomology-zeta}. In dimension three,
Hodge star identifies the degree-three and degree-zero spectra,
and degree-two with degree-one.  Moreover
$\zeta_{1,W}=\zeta_{0,W}+\zeta_{c,W}$ by exact/coexact splitting.
Substitute these identities into the Ray--Singer definition
$\log T=\frac12\sum_{p=0}^{3}(-1)^p p\zeta'_{p,W}(0)$:
\[
 2\log T=-\zeta_1'+2\zeta_2'-3\zeta_3'
 =\zeta_1'-3\zeta_0'=\zeta_c'-2\zeta_0'.
\]
The standard analytic-torsion theorem for odd-dimensional closed
manifolds with acyclic unitary flat coefficients gives metric
independence and the Cheeger--M\"uller identification.
All the coefficients of the two defining zeta functions are exactly
provided by $B_W(z)$ in \eqref{p1b:eq:v13-isotypic-series}.
\end{proof}

\begin{remark}[Torsion is neither an eta correction nor a contact-volume correction]
\label{p1b:rem:v13-torsion-owner}
The torsion is the alternating \emph{full de Rham} determinant
combination, while eta is a signed Dirac asymmetry and $4\pi^2$
is a full-$S^3$ contact integral. Their different operator domains
preclude substituting any one for another without an explicit
mathematical map.  The nontrivial-irrep acyclicity condition is
essential when presenting torsion as a scalar without a separate
choice of cohomology determinant-line basis.
\end{remark}

\section{Shortest exceptional displacements and the finite-cover heat trace}
\label{p1b:sec:v13-heat}

\begin{theorem}[Exact minimum displacement and exponentially small orbital heat terms]
\label{p1b:thm:v13-heat-orbit}
For the left action of $\Gamma$ on $S^3(R)$, each $g\in\Gamma$
has constant displacement
$\operatorname{dist}(q,gq)=R\arccos(\operatorname{Re}g)$.
The least nonzero displacement is
\begin{equation}
\label{p1b:eq:v13-displacements}
 \boxed{\ell_{2T}=\frac{\pi R}{3},\qquad
 \ell_{2O}=\frac{\pi R}{4},\qquad
 \ell_{2I}=\frac{\pi R}{5}.}
\end{equation}
For any unitary flat $W$ and its round spin Dirac operator,
\begin{equation}
\label{p1b:eq:v13-heat-identity}
 \operatorname{Tr}e^{-tD_{\Gamma,W}^2}
 =\frac{\dim W}{|\Gamma|}\operatorname{Tr}e^{-tD_{S^3}^2}
 +\mathcal R_{\Gamma,W}(t),
 \qquad
 |\mathcal R_{\Gamma,W}(t)|
 \le C\,t^{-N}e^{-c_\Gamma/t}
\end{equation}
for all sufficiently small $t>0$ and suitable positive constants
$C,N,c_\Gamma$, with $c_\Gamma$ controlled by
$\ell_\Gamma^2$.  In particular every ordinary local heat
coefficient of the quotient equals $\dim W/|\Gamma|$ times the
corresponding full-sphere heat coefficient; the nonidentity
group-dependent orbital information survives beyond the power-series
heat expansion.
\end{theorem}
\begin{proof}
The bi-invariant round metric gives
$\operatorname{dist}(q,gq)=\operatorname{dist}(1,q^{-1}gq)$;
conjugation preserves $\operatorname{Re}g$, and geodesic distance
on unit quaternions is $R\arccos(\operatorname{Re}g)$.
The maximal real height of a nonidentity element is the established
nearest-shell height $\cos(\pi/m_\Gamma)$ for
$m_\Gamma=3,4,5$, respectively; this proves the first identity.
For the second, the finite projection onto $\Gamma$-equivariant
sections is an average over $g\in\Gamma$. The identity term in the
heat trace is $\dim W/|\Gamma|$ times the full-sphere trace.
Every $g\ne1$ displaces every point at least $\ell_\Gamma>0$.
The off-diagonal Gaussian heat-kernel bound for a compact manifold
then bounds the remaining integrals by a finite sum of
$t^{-N}e^{-c/t}$ terms. This also follows from the finite-cover
heat-kernel method for twisted Dirac-type operators
\cite{CacicMarcolliTeh2014CosmicTopology}.
\end{proof}

\begin{corollary}[A precise smooth-cutoff regime of group-order spectral-action scaling]
\label{p1b:cor:v13-spectral-action-scale}
Let $f(x)=\int_0^\infty\psi(s)e^{-s x^2}\,ds$ with
$\psi\in C_c^\infty((0,\infty))$. Then as $\Lambda\to\infty$,
\begin{equation}
\label{p1b:eq:v13-spectral-action-scale}
 \boxed{\operatorname{Tr}f(D_{\Gamma,W}/\Lambda)
 =\frac{\dim W}{|\Gamma|}\operatorname{Tr}f(D_{S^3}/\Lambda)
 +O(e^{-c\Lambda^2})}
\end{equation}
for some $c>0$ depending on $\psi,\Gamma,R$.  For broader classes
of smooth test functions, the published spectral-action theorem
provides corresponding asymptotics
\cite{CacicMarcolliTeh2014CosmicTopology}.
This is a \emph{finite-cover quotient} statement: no replacement of
the original Hopf value $4\pi^2$ or Born--Hopf coupling correction
is licensed by it.
\end{corollary}
\begin{proof}
The compact support of $\psi$ lies in $[s_-,s_+]$ with
$0<s_-<s_+<\infty$.  Insert the finite-cover heat identity with
$t=s/\Lambda^2$ under the integral. Uniformly in this interval the
error is bounded by a polynomial in $\Lambda$ times
$\exp[-c_\Gamma\Lambda^2/s_+]$, which can be absorbed into
$O(e^{-c\Lambda^2})$ for a smaller $c>0$.
\end{proof}

\begin{remark}[A hierarchy of mathematically different exceptional observables]
\label{p1b:rem:v13-observable-atlas}
The following independent readouts are now explicit:
\[
\begin{array}{c|l}
\text{Owner} & \text{Mathematical quantity}\\\hline
S^3\text{ with fixed Hopf normalization} & 4\pi^2\\
\Gamma\backslash S^3\text{ contact quotient} & 4\pi^2/|\Gamma|\\
\text{round scalar quotient} & 48,80,168\text{ as first positive }R^2\lambda\\
\text{round coexact one-forms} & 4\text{ common ground; }36,64,144\text{ next}\\
\text{round Dirac quotient} & \eta_D(0)\text{ and spin-lift dependence}\\
\text{nontrivial flat unitary module} & \text{acyclic Ray--Singer torsion}\\
\text{round finite-cover heat trace} & (\dim W)/|\Gamma|\text{ local scaling}\\
\end{array}
\]
Here the array's final row is a group-order \emph{factor}, not a
universally selected physical determinant weight.  These distinctions
are retained even if no downstream model uses them; a later physical
application must select its actual operator, representation, metric,
spin lift, and action-owned normalization.
\end{remark}


\section{Conjugacy-resolved exceptional geodesic and representation observables}
\label{p1b:sec:v14-conjugacy-double-gauge}
The following constructions use the quaternionic embeddings and finite-group
structures fixed above.  They are included as mathematical objects in their
own right.  None changes the full-cover Hopf number $4\pi^2$, and none asserts
that the physical spectral action contains a quantum double or a new flat gauge
sector.  The finite group doubles and twisted variants belong to the standard
orbifold/quantum-double theory \cite{Kirillov2002Modular,DijkgraafPasquierRoche1991};
the torsion comparison below is informed by the explicit tetrahedral resolution
work of \cite{ChiriviSpreafico2018}.

\begin{definition}[Height, displacement angle and class-resolved finite measure]
\label{p1b:def:v14-class-geodesic}
For $g\in\Gamma\subset SU(2)\cong S^3$ write
$h(g)=\operatorname{Re}(g)$, $\theta(g)=\arccos h(g)\in[0,\pi]$,
$\ell(g)=R\theta(g)$, and
\[
 \nu_\Gamma:=\sum_{[g]\in\operatorname{Conj}(\Gamma)}
                  |[g]|\,\delta_{\theta(g)}.
\]
This is an \emph{element/displacement-angle measure}, not by definition the
primitive closed-geodesic length spectrum.  We denote the centralizer by
$C_\Gamma(g)$ and the number of irreducible complex representations by $k(\Gamma)$.
\end{definition}

\begin{theorem}[Complete exceptional conjugacy-height and displacement-angle distributions]
\label{p1b:thm:v14-class-geodesic-atlas}
Put $\varphi=(1+\sqrt5)/2$. For the standard embeddings the full weighted
height measures, including class multiplicity, are
\begin{align}
\label{p1b:eq:v14-class-height}
 \nu_{2T}&=\delta_0+\delta_\pi+6\delta_{\pi/2}
                      +8\delta_{\pi/3}+8\delta_{2\pi/3},\\
 \nu_{2O}&=\delta_0+\delta_\pi+18\delta_{\pi/2}
                      +8\delta_{\pi/3}+8\delta_{2\pi/3}
                      +6\delta_{\pi/4}+6\delta_{3\pi/4},\\
 \nu_{2I}&=\delta_0+\delta_\pi+30\delta_{\pi/2}
                      +20\delta_{\pi/3}+20\delta_{2\pi/3}
                      +12\sum_{j=1}^{4}\delta_{j\pi/5}.
\end{align}
The individual conjugacy-class sizes, ordered by height, are as follows:
\[
\begin{array}{c|c|c|c}
 \Gamma & h(g)&\text{conjugacy-class size(s)}&|C_\Gamma(g)|\\\hline
 2T &1,-1&1,1&24,24\\
    &0&6&4\\
    &+\tfrac12,-\tfrac12&4+4,\ 4+4&6,6\\\hline
 2O &1,-1&1,1&48,48\\
    &0&6+12&8,4\\
    &+\tfrac12,-\tfrac12&8,8&6,6\\
    &+1/\sqrt2,-1/\sqrt2&6,6&8,8\\\hline
 2I &1,-1&1,1&120,120\\
    &0&30&4\\
    &+\tfrac12,-\tfrac12&20,20&6,6\\
    &\pm\varphi/2,\ \pm1/(2\varphi)&12\ \text{each}&10\ \text{each}
\end{array}
\]
The respective total masses of $\nu_\Gamma$ are exactly $24,48,120$, and
the shortest nonidentity displacement lengths are
$\pi R/3,\pi R/4,\pi R/5$. For every $g$, left multiplication by $g$
is a Clifford translation of the round $S^3(R)$ with displacement $\ell(g)$.
\end{theorem}
\begin{proof}
The quaternionic coordinates from the previous exceptional atlas have
real parts $\{\pm1,0,\pm\tfrac12\}$ for $2T$,
those values together with $\pm1/\sqrt2$ for $2O$, and
$\{\pm1,0,\pm\tfrac12,\pm\varphi/2,\pm(2\varphi)^{-1}\}$
for $2I$.  Counting the actual unit quaternions at each real height gives
respectively $(1,1,6,8,8)$, $(1,1,18,8,8,6,6)$ and
$(1,1,30,20,20,12,12,12,12)$.
The noncentral centralizer lies in the circle of unit quaternions commuting
with $g$ and is cyclic. Its order follows from orbit--stabilizer
$|C_\Gamma(g)|=|\Gamma|/|[g]|$; the two central elements have whole-group
centralizer.  The known $2T$ class splitting of the $\pm\tfrac12$ heights,
and $2O$ splitting of the height-zero elements, yield the displayed class
multiplicities.  Since $\cos\theta=\operatorname{Re}(g)$,
$\operatorname{dist}_{S^3(R)}(x,gx)=R\arccos\operatorname{Re}(g)$ for every
$x\in S^3(R)$; taking the smallest positive angle gives the gaps.
The angles enumerate all displacement types, but primitive-orbit and
holonomy weights are a separate quotient-geodesic counting problem.
\end{proof}

\begin{proposition}[Exact character-weighted finite-cover orbital trace]
\label{p1b:prop:v14-orbital-image-trace}
Let $W$ be a finite-dimensional unitary $\Gamma$-representation and $P$
a Laplace-type operator on the covering round $S^3$ commuting with the
left $\Gamma$ action. For its induced operator $P_{\Gamma,W}$ on the associated
flat bundle, the heat trace is exactly
\[
 \boxed{\operatorname{Tr}e^{-tP_{\Gamma,W}}
  =\frac1{|\Gamma|}\sum_{g\in\Gamma}
   \chi_W(g^{-1})\operatorname{Tr}(L_g e^{-tP})}. 
\]
Each term has a specified displacement $R\theta(g)$ from the preceding
atlas.  For a class-invariant heat kernel, this is a finite sum of class
contributions with the class sizes above.  The identity term is
$(\dim W)/|\Gamma|$ times the full-cover heat trace; the others record
global exceptional information. A fixed-charge or spinorial spectrum is
obtained only after choosing its corresponding covering operator and lift.
\end{proposition}
\begin{proof}
The equivariant projection on covering $W$-valued fields is
$|\Gamma|^{-1}\sum_g L_g\otimes W(g^{-1})$.  It commutes with the heat
operator and its image is the twisted quotient field space. Taking the
operator trace and using tensor-factorization proves the formula. The
Clifford displacement identity supplies the geometric classification of
its nonidentity summands.  No asymptotic expansion is needed.
\end{proof}

\begin{definition}[Untwisted finite-group quantum double]
\label{p1b:def:v14-drinfeld-double}
For a finite group $\Gamma$ write $D(\Gamma)$ for its ordinary, untwisted
Drinfeld double. Its simple complex modules are indexed by $(C,\rho)$,
where $C$ is a conjugacy class, $g\in C$ is fixed and
$\rho\in\operatorname{Irr}(C_\Gamma(g))$.  Their quantum dimensions and
untwisted ribbon twists are, respectively,
\[
 d_{C,\rho}=|C|\dim\rho,\qquad
 \vartheta_{C,\rho}={\chi_\rho(g)\over\dim\rho}.
\]
The second expression is well-defined because $g$ is central in its own
centralizer and thus acts scalarly in every irreducible $\rho$.
\end{definition}

\begin{theorem}[Complete exceptional untwisted double-sector counts and dimensions]
\label{p1b:thm:v14-quantum-double-counts}
For $\Gamma=2T,2O,2I$ the numbers of simple $D(\Gamma)$ modules are
\[
 \boxed{N_D(2T)=42,\quad N_D(2O)=56,\quad N_D(2I)=74.}
\]
Their total squared quantum dimensions and the numbers of invertible
simples are
\[
\begin{array}{c|ccc}
 &2T&2O&2I\\\hline
 \sum_{(C,\rho)} d_{C,\rho}^2&576&2304&14400\\
 \mathcal D=\sqrt{\sum d^2}&24&48&120\\
 \#\{(C,\rho):d_{C,\rho}=1\}&6&4&2
\end{array}
\]
These are \emph{category-theoretic} observables, not additional physical
particle species assigned by the present framework.
\end{theorem}
\begin{proof}
Both central classes have $C_\Gamma(g)=\Gamma$ and each contributes
$k(\Gamma)=7,8,9$ simples, respectively.  The centralizers of all
noncentral elements are cyclic, contributing one simple per centralizer
element. From the class table the noncentral contributions are
\[
\begin{array}{c|c}
2T&4+4(6)=28\\
2O&8+4+8+8+6+6=40\\
2I&4+6+6+4(10)=56.
\end{array}
\]
Adding $2k(\Gamma)$ gives $42,56,74$.  For any class $C$,
$\sum_\rho (|C|\dim\rho)^2=|C|^2|C_\Gamma(g)|=|C|\,|\Gamma|$.
Summing over classes gives $|\Gamma|^2$. Invertibility requires
$|C|=\dim\rho=1$, so the class must be either central element and
$\rho$ one-dimensional. Since the abelianizations are
$2T_{\mathrm{ab}}\cong C_3$, $2O_{\mathrm{ab}}\cong C_2$ and
$2I_{\mathrm{ab}}=1$, the numbers are $2(3),2(2),2(1)$.
\end{proof}

\begin{corollary}[Finite cohomological family of twisted exceptional doubles]
\label{p1b:cor:v14-twisted-double-classes}
For the standard trivial action on $U(1)$, group cohomology gives
\[
 \boxed{H^3(2T;U(1))\cong\mathbb Z/24,\quad
 H^3(2O;U(1))\cong\mathbb Z/48,\quad
 H^3(2I;U(1))\cong\mathbb Z/120.}
\]
These classify $3$-cocycle \emph{cohomology classes} available for the
quasi-Hopf construction $D^\omega(\Gamma)$, not necessarily distinct
braided-equivalence classes after possible group automorphisms.
The untwisted module counts above are not asserted for arbitrary
$\omega$; nontrivial twists can change projective centralizer
representation theory \cite{DijkgraafPasquierRoche1991}.
\end{corollary}
\begin{proof}
The exponential coefficient exact sequence
$0\to\mathbb Z\to\mathbb R\to U(1)\to0$ and averaging over a
finite group give $H^j(\Gamma;\mathbb R)=0$ for $j>0$, hence
$H^3(\Gamma;U(1))\cong H^4(\Gamma;\mathbb Z)$.
The four-periodic group-cohomology computation already established in
this paper has $H^4(\Gamma;\mathbb Z)\cong\mathbb Z/|\Gamma|$.
The final distinction concerns equivalence of quasi-Hopf categories,
which is not the same as equality of cocycle classes.
\end{proof}

\begin{proposition}[A finite and observable-typed flat-connection moduli problem]
\label{p1b:prop:v14-flat-su2-moduli}
For the spherical space form $M_\Gamma=\Gamma\backslash S^3$,
flat $SU(2)$ connections modulo gauge equivalence are parametrized by
\[
 \boxed{\mathscr M_{\rm flat}(M_\Gamma,SU(2))
       \cong\operatorname{Hom}(\Gamma,SU(2))/SU(2).}
\]
In particular, the moduli set is finite. One may compute it from the
complex two-dimensional representation table by retaining exactly
those representations with determinant one, including reducible pairs
of one-dimensional characters.  This is a \emph{holonomy-space} result;
its Chern--Simons invariants, eta weights and torsions require separate
normalization and evaluation and do not follow from $|\Gamma|$ alone.
\end{proposition}
\begin{proof}
A flat principal connection has a holonomy homomorphism from
$\pi_1(M_\Gamma)=\Gamma$, well-defined up to $SU(2)$ conjugation;
conversely the universal-cover quotient $(S^3\times SU(2))/\Gamma$
with holonomy $\rho$ carries a flat connection. Every principal
$SU(2)$ bundle over a three-dimensional CW complex is topologically
trivial since $BSU(2)$ is 3-connected.  Every homomorphism is a
unitary two-dimensional representation with determinant one, and
unitary equivalence can be implemented by $SU(2)$ because a unitary
intertwiner can be multiplied by a central phase to have determinant
one.  A finite group has only finitely many two-dimensional
representations up to isomorphism, proving finiteness.
\end{proof}


\begin{corollary}[Exact finite flat-$SU(2)$ holonomy counts]
\label{p1b:cor:v14-flat-su2-counts}
For the binary tetrahedral, octahedral and icosahedral space forms,
respectively, the numbers of flat $SU(2)$ holonomy classes are
\[
\boxed{\#\mathscr M_{\rm flat}(M_{2T},SU(2))=3,\quad
       \#\mathscr M_{\rm flat}(M_{2O},SU(2))=4,\quad
       \#\mathscr M_{\rm flat}(M_{2I},SU(2))=3.}
\]
The counts include the trivial connection. They enumerate gauge-equivalence
classes, not the Chern--Simons values or Floer generators.
\end{corollary}
\begin{proof}
Use the irreducible complex representation degrees and determinant
characters in the preceding McKay/character atlas.  For $2T$, the
abelianization $C_3$ gives one nontrivial determinant-one reducible
pair $\chi\oplus\chi^{-1}$, in addition to $\mathbf1\oplus\mathbf1$;
there is one faithful determinant-one irreducible of dimension two.
The remaining two-dimensional irreducibles are its nontrivial $C_3$
character twists and have nontrivial determinant.  For $2O$, the
abelianization $C_2$ supplies the two scalar pairs
$\mathbf1\oplus\mathbf1$ and $\varepsilon\oplus\varepsilon$;
there are two inequivalent faithful determinant-one spinorial
representations.  The remaining two-dimensional irreducible factors
through the rotational quotient and has nontrivial determinant.
For the perfect group $2I$, there are no nontrivial one-dimensional
characters and there are exactly two inequivalent faithful
quaternionic two-dimensional irreducibles, each of determinant one.
The preceding proposition identifies these representation classes with
flat $SU(2)$ holonomy classes.
\end{proof}

\begin{proposition}[Acyclicity and torsion as a separate quotient invariant]
\label{p1b:prop:v14-acyclic-flat-torsion}
If $\rho\ne\mathbf1$ is an irreducible finite-dimensional unitary
complex representation of $\Gamma$, then the associated flat bundle
$E_\rho$ on $M_\Gamma$ is acyclic:
\[
 H^q(M_\Gamma;E_\rho)=0\qquad(q=0,1,2,3).
\]
Thus its Reidemeister and Ray--Singer torsions are defined as positive,
metric-independent numbers with an agreed cochain and volume
normalization, and the classical Cheeger--M\"uller equality relates
their consistent conventions. The group resolution of
\cite{ChiriviSpreafico2018} gives an explicit calculation route in the
tetrahedral family, but does not determine all octahedral and
icosahedral torsion values without further computation.
\end{proposition}
\begin{proof}
The universal cover is $S^3$ with cohomology only in degrees zero
and three. The transfer/Cartan--Leray spectral sequence for the finite
cover has $E_2^{p,q}=H^p(\Gamma;H^q(S^3;\mathbb C)\otimes\rho)$.
Averaging kills positive-degree group cohomology over $\mathbb C$;
for $p=0$ the invariants $\rho^\Gamma$ vanish by nontrivial
irreducibility. Thus every entry vanishes and so does the total
cohomology.  The torsion statements use the standard acyclic flat
unitary setting; no equality of torsions for different $\Gamma$ is
claimed.
\end{proof}

\begin{remark}[Precise separation of geometric owners]
\label{p1b:rem:v14-observable-and-research-scope}
The six distinct objects used here are: the embedded finite subgroup
$\Gamma\subset S^3$; the quotient manifold $\Gamma\backslash S^3$;
the Clifford-translation length multiset; a twisted quotient
Laplace/Dirac/Hodge operator; the modular category
$\operatorname{Rep}(D^\omega(\Gamma))$; and a flat-$SU(2)$ holonomy
class. Their exact finite calculations do not identify them with a
single determinant, correction coefficient, or physical decay map.
In particular $4\pi^2$ remains the inherited $S^3$ Hopf number;
$\frac{4\pi^2}{|\Gamma|}$ is a separately specified quotient contact
integral.  The open mathematical computations include the complete
irreducible flat-bundle torsion values, the twisted-double modular
matrices for each cohomology class, and normalized Chern--Simons and
Floer invariants for all flat holonomies. They are not silently
promoted to completed results.
\end{remark}

\part{Exceptional flat-bundle torsions and modular data}
\label{p1b:part:v15-torsion-modular}
\section{Algebraic evaluation of exceptional flat-bundle torsions}
\label{p1b:sec:v15-closed-torsion}
The exact zeta-function construction in Theorem~\ref{p1b:thm:v13-torsion-algorithm}
provides an analytic definition for every nontrivial irreducible flat unitary bundle.
Closed evaluations of precisely these spherical-space-form torsions were previously
obtained by Dowker and Chang \cite{DowkerChang2009AnalyticTorsion}; the following
records their results in a notation that distinguishes representations having
the same dimension and keeps their real/complex convention explicit.

\begin{theorem}[Complete algebraic exceptional torsion atlas]
\label{p1b:thm:v15-exceptional-torsion-table}
Let $M_\Gamma=\Gamma\backslash S^3(R)$, with the standard round metric,
and let $W$ be a nontrivial irreducible flat unitary representation.
For the acyclic complex torsion $T_{\mathbb C}(M_\Gamma,W)$ defined
by \eqref{p1b:eq:v13-torsion-formula}, the complete values are given by
\[
\begin{array}{c|l|c}
 \Gamma&\text{representation }W& T_{\mathbb C}(M_\Gamma,W)\\\hline
 2T&1',\,1''\ (\text{nontrivial }C_3\text{ characters})&3/2\\
   &3&2\\
   &2_s\ (\text{defining spinor})&1/2\\
   &2_s',\,2_s''\ (\text{character twists of }2_s)&2\\\hline
 2O&1'\ (\text{sign})&4/3\\
   &2&3/2\\
   &3,\,3'&2,\,1\ \text{respectively}\\
   &4_s&2\\
   &2_s,\,2_s'&(2-\sqrt2)/2,\,(2+\sqrt2)/2\\\hline
 2I&4&5/3\\
   &2_s,\,2_s'&(3-\sqrt5)/4,\,(3+\sqrt5)/4\\
   &6_s&2\\
   &3,\,3'&1+1/\sqrt5,\,1-1/\sqrt5\\
   &5&3/2\\
   &4_s&1
\end{array}
\tag{T15.1}\label{p1b:eq:v15-torsion-table}
\]
There are respectively six, seven and eight nontrivial irreducible
complex representations, exhausting $k(2T)=7$, $k(2O)=8$, $k(2I)=9$.
The primes in this table follow the conventions of
\cite{DowkerChang2009AnalyticTorsion}; a later identification with another
character table must be made by characters, not by dimensions alone.
For the two nonreal-conjugate tetrahedral pairs, that source reports
$T(1')=T(1'')=(3/2)^2$ and $T(2_s')=T(2_s'')=2^2$
when the conjugate pair is counted together as a real bundle.
The displayed values $3/2$ and $2$ are the corresponding
\emph{individual complex-flat-bundle} torsions; direct-sum torsion is
multiplicative. All other rows are transcribed in the stated unitary
representation normalization.
\end{theorem}
\begin{proof}
Dowker--Chang, Eq.~(45), evaluates the exceptional torsions by
character averaging and an effective zeta cancellation, then checks
selected products by induction from cyclic lens-space torsions.
For nontrivial complex line holonomy $\zeta_m^r$, the one-sided lens-space
torsion in the compatible normalization is
$T(L(m,1),\zeta_m^r)=|1-\zeta_m^r|^2$.
The Ray--Singer covering identity gives
$T(M_\Gamma,\operatorname{Ind}_H^\Gamma\psi)
 =T(M_H,\psi)$ for $H\subset\Gamma$.
Multiplicativity on direct sums and the character tables yield the table.
For the two nonreal tetrahedral pairs, the published doubled real-bundle
numbers are square roots away from the individual complex values,
which accounts for the displayed convention conversion.
\end{proof}

\begin{proposition}[Independent icosahedral induction certificate]
\label{p1b:prop:v15-icosahedral-induction}
Let $x=T(2_s)$, $y=T(2_s')$, $u=T(4_s)$ and $v=T(6_s)$
in the $2I$ quotient. The cyclic $C_{10}$ and $C_6$ induction
identities imply
\[
 xuv=\frac{3-\sqrt5}{2},\qquad
 yuv=\frac{3+\sqrt5}{2},\qquad
 v^2=4,\qquad xyuv^2=1.
\tag{T15.2}\label{p1b:eq:v15-induction-system}
\]
Their unique positive solution is
\[
\boxed{v=2,\quad u=1,\quad
x=\frac{3-\sqrt5}{4},\quad
 y=\frac{3+\sqrt5}{4}.}
\tag{T15.3}\label{p1b:eq:v15-induction-result}
\]
\end{proposition}
\begin{proof}
Under $C_{10}\subset2I$, the inductions of its characters of
exponents $1,3,5$ are respectively
$2_s\oplus4_s\oplus6_s$,
$2_s'\oplus4_s\oplus6_s$ and $6_s\oplus6_s$.
The corresponding lens torsions are
$|1-e^{\pi i/5}|^2=(3-\sqrt5)/2$,
$|1-e^{3\pi i/5}|^2=(3+\sqrt5)/2$ and $4$.
A $C_6$ character of exponent one induces
$2_s\oplus 2_s'\oplus4_s\oplus2(6_s)$ and its torsion is
$|1-e^{\pi i/3}|^2=1$.
These are the four equations. Positivity gives $v=2$;
multiplying the first two shows $xyu^2v^2=1$, whereas the last
gives $xyuv^2=1$, hence $u=1$ and the displayed $x,y$ follow.
This independently reconstructs the nontrivial golden-ratio torsions
without fitting a spectral correction.
\end{proof}

\begin{remark}[Nonacyclic trivial holonomy requires a volume convention]
\label{p1b:rem:v15-trivial-torsion}
The trivial representation has nonzero cohomology and does not belong to
the acyclic table. In the cohomology-basis/volume convention of
\cite{DowkerChang2009AnalyticTorsion} its torsion is, respectively,
$\mathrm{Vol}(S^3)/(72)$, $\mathrm{Vol}(S^3)/(96)$ and
$\mathrm{Vol}(S^3)/(120)$ for $2T,2O,2I$ at unit radius
(with the appropriate radius-dependent volume in general).
These volume-dependent quantities are not the unit-Chern Hopf contact
residue $4\pi^2$, and no numerical comparison with that residue is
an action-derived coupling formula.
\end{remark}

\section{Untwisted exceptional quantum-double modular data}
\label{p1b:sec:v15-modular-calculus}
Write $z=-1\in\Gamma$. For each conjugacy class select $a\in\Gamma$,
write $C_a=C_\Gamma(a)$, and let $\chi$ be the character of an
irreducible $C_a$-module. Put $X=(a,\chi)$ and similarly
$Y=(b,\psi)$. Conventions for complex conjugation in the modular
$S$ matrix are fixed below, consistently with
\cite{CoquereauxZuber2013Doubles}.

\begin{theorem}[Explicit finite character formula for all untwisted modular entries]
\label{p1b:thm:v15-untwisted-ST}
The entire modular matrices of $\operatorname{Rep}(D(\Gamma))$ are
specified, without a numerical fit, by
\begin{equation}
 \boxed{\mathsf S_{X,Y}
 =\frac{1}{|C_a|\,|C_b|}
 \sum_{\substack{x\in\Gamma\\ [a,xbx^{-1}]=1}}
 \overline{\chi(xbx^{-1})}\,
 \overline{\psi(x^{-1}ax)}}
 \tag{T15.4}\label{p1b:eq:v15-S}
\end{equation}
\begin{equation}
 \boxed{\mathsf T_{X,Y}
 =\delta_{X,Y}\frac{\chi(a)}{\chi(1)}}
 \tag{T15.5}\label{p1b:eq:v15-T}
\end{equation}
The matrices have sizes $42$, $56$, and $74$ for $2T$, $2O$, and $2I$.
For the vacuum $0=(1,\mathbf1)$ and any simple $X$,
\[
\mathsf S_{0X}=\frac{d_X}{|\Gamma|},\qquad
\mathsf T_{00}=1,\qquad
\sum_Xd_X^2=|\Gamma|^2.
\tag{T15.6}\label{p1b:eq:v15-vacuum-row}
\]
With the customary ribbon convention, $\mathsf S$ is symmetric unitary,
$\mathsf S^2=\mathsf C$ is charge conjugation, and
$(\mathsf S\mathsf T)^3=\mathsf C$.
\end{theorem}
\begin{proof}
The character formulas are the standard finite-double modular
character transform, as established in
\cite{CoquereauxZuber2013Doubles}; the preceding section already
constructs its simple labels $(a,\chi)$ and their dimensions.
For $a=1$ and the trivial character, all $x$ contribute and the
character sum reduces to
$|\Gamma|\chi(1)/(|\Gamma||C_b|)
=\chi(1)/|C_b|=d_X/|\Gamma|$.
The modular relations follow from the ribbon structure of
$\operatorname{Rep}(D(\Gamma))$; they are conditions for a
full finite evaluation of the displayed character sums, not
consequences of the bare counts alone.
\end{proof}

\begin{corollary}[The complete central-flux modular block]
\label{p1b:cor:v15-central-block}
For ordinary irreducible $\Gamma$-characters $\chi_\rho,\chi_\sigma$,
write $d_\rho=\chi_\rho(1)$ and
$\chi_\rho(z)=p_\rho d_\rho$ with $p_\rho\in\{\pm1\}$.
For $a,b\in\{0,1\}$, the centrally supported simple
labels $X=(z^a,\rho)$, $Y=(z^b,\sigma)$ satisfy
\[
\boxed{\mathsf S_{X,Y}=
 \frac{d_\rho d_\sigma}{|\Gamma|}
 p_\rho^{\,b}p_\sigma^{\,a},\qquad
 \mathsf T_{X,X}=p_\rho^{\,a}.}
\tag{T15.7}\label{p1b:eq:v15-central-block}
\]
In particular the entire central block is determined by the
center-parity characters of the affine McKay representation bank.
\end{corollary}
\begin{proof}
As $z$ is central, the commuting condition in
\eqref{p1b:eq:v15-S} is automatic and both character values are
constant over the group. Since $C_{z^a}=\Gamma$, the sum of
$|\Gamma|$ identical terms gives the stated entry. The twist
formula is immediate from \eqref{p1b:eq:v15-T}.
\end{proof}

\begin{theorem}[Exact modular-twist exponents and trace sequences]
\label{p1b:thm:v15-T-exponents}
Write $\mathbf1_{m\mid k}=1$ if $m$ divides $k$ and zero otherwise,
and let $r_\Gamma=|\operatorname{Irr}(\Gamma)|=7,8,9$.
The noncentral conjugacy-class centralizer multisets are
\[
\begin{array}{c|c|c|c}
 \Gamma&\{|C_\Gamma(g)|:[g]\text{ noncentral}\}
  &\sum_\rho p_\rho&\operatorname{ord}(\mathsf T)\\\hline
 2T&\{4,6,6,6,6\}&1&12\\
 2O&\{4,6,6,8,8,8\}&2&24\\
 2I&\{4,6,6,10,10,10,10\}&1&60
\end{array}
\tag{T15.8}\label{p1b:eq:v15-centralizers}
\]
For every positive integer $k$ the traces of the untwisted modular
$\mathsf T$ are consequently
\begin{align}
\operatorname{Tr}\mathsf T_{2T}^{\,k}
 &=7+(7\mathbf1_{2\mid k}+\mathbf1_{2\nmid k})
   +4\mathbf1_{4\mid k}+12\mathbf1_{3\mid k}
   +12\mathbf1_{6\mid k},\\
\operatorname{Tr}\mathsf T_{2O}^{\,k}
 &=8+(8\mathbf1_{2\mid k}+2\mathbf1_{2\nmid k})
   +12\mathbf1_{4\mid k}+6\mathbf1_{3\mid k}
   +6\mathbf1_{6\mid k}+16\mathbf1_{8\mid k},\\
\operatorname{Tr}\mathsf T_{2I}^{\,k}
 &=9+(9\mathbf1_{2\mid k}+\mathbf1_{2\nmid k})
   +4\mathbf1_{4\mid k}+6\mathbf1_{3\mid k}
   +6\mathbf1_{6\mid k}+20\mathbf1_{5\mid k}
   +20\mathbf1_{10\mid k}.
\tag{T15.9}\label{p1b:eq:v15-T-traces}
\end{align}
The noncentral terms at $k=3$ already distinguish the corrected formula
from a centralizer-order-only rule.  At $k=12,24,60$ the respective
traces remain exactly $42,56,74$.
\end{theorem}
\begin{proof}
For a noncentral $g$, its centralizer $C_\Gamma(g)$ is cyclic of
order $m$, but $g$ need not be a generator: write $g=h^a$ with $h$
a generator.  The $m$ ordinary characters $\chi_j(h)=e^{2\pi i j/m}$
give ribbon twists $\chi_j(g)=e^{2\pi i ja/m}$, hence
\[
  \sum_{j=0}^{m-1}\chi_j(g)^k
   =m\,\mathbf1_{\operatorname{ord}(g)\mid k}.
\]
One must use the \emph{element order}, not the centralizer order.
The noncentral $(\operatorname{ord}(g),|C_\Gamma(g)|)$ pairs,
counted by conjugacy class, are
\[
\begin{array}{c|l}
2T&(4,4)^{(1)},(3,6)^{(2)},(6,6)^{(2)}\\
2O&(4,4)^{(1)},(4,8)^{(1)},(3,6)^{(1)},
    (6,6)^{(1)},(8,8)^{(2)}\\
2I&(4,4)^{(1)},(3,6)^{(1)},(6,6)^{(1)},
    (5,10)^{(2)},(10,10)^{(2)}.
\end{array}
\]
These follow from the angles and class sizes in
Theorem~\ref{p1b:thm:v14-class-geodesic-atlas}; in particular an
element at angle $2\pi/3$ has order three despite a cyclic
centralizer of order six. The identity-flux sector contributes
$r_\Gamma$ unit twists, while the central $z$-flux contributes
$\sum_\rho p_\rho^k$, equal to $r_\Gamma$ for even $k$ and
$1,2,1$ for odd $k$.  The resulting least common multiples of the
\emph{actual} ribbon orders are still $12,24,60$: the classes of
orders $4,6$; $8,3$; and $4,3,5$, respectively, supply the necessary
prime powers. No assertion that every element generates its
centralizer is used.
\end{proof}

\begin{corollary}[Untwisted exceptional doubles have trivial topological central charge]
\label{p1b:cor:v15-gauss-sum}
For $\theta_X=\mathsf T_{XX}$ and $d_X=\mathsf S_{0X}/\mathsf S_{00}$,
\[
\boxed{\tau_+:=\sum_Xd_X^2\theta_X=|\Gamma|,
\qquad \tau_-:=\sum_Xd_X^2\theta_X^{-1}=|\Gamma|.}
\tag{T15.10}\label{p1b:eq:v15-gauss}
\]
In particular $\tau_+/\sqrt{\sum_Xd_X^2}=1$.
\end{corollary}
\begin{proof}
At fixed nonidentity flux $a$, summing over all irreducible
centralizer characters gives
$|[a]|^2\sum_\chi\chi(1)\chi(a)=0$, since this is
$|[a]|^2$ times the regular character of $C_a$ evaluated at
$a\ne1$.  At identity flux every twist is one and
$\sum_{\rho\in\operatorname{Irr}(\Gamma)}(\dim\rho)^2
=|\Gamma|$. The inverse-twist sum is its complex conjugate.
\end{proof}

\section{A twist-stability theorem and a genuine odd-twist effect}
\label{p1b:sec:v15-all-twists}

\begin{theorem}[Vanishing Schur multipliers and all-twist stable simple dimensions]
\label{p1b:thm:v15-Schur-twist-stability}
For every finite $\Gamma\subset SU(2)$ acting freely by left
multiplication on $S^3$, not only the exceptional groups,
\[
\boxed{H^2(\Gamma;U(1))=0.}
\tag{T15.11}\label{p1b:eq:v15-Schur-zero}
\]
For every normalized $3$-cocycle $\omega\in Z^3(\Gamma;U(1))$,
all the transgressed centralizer $2$-cocycles $\alpha_g$
are cohomologically trivial.  Consequently the simple objects
of $\operatorname{Rep}(D^\omega(\Gamma))$ have the same numbers
and dimensions, class by class, as those of the untwisted double.
In particular, for \emph{every} cohomology class of twists on the
three exceptional groups,
\[
\boxed{
\begin{array}{c|ccc}
\Gamma&2T&2O&2I\\\hline
\#\operatorname{Irr}D^\omega(\Gamma)&42&56&74\\
\sum_Xd_X^2&576&2304&14400
\end{array}}
\tag{T15.12}\label{p1b:eq:v15-all-twist-counts}
\]
The modular matrices and ribbon twists need not agree between
noncohomologous $\omega$; only the stated sector-count and dimension
invariants are fixed.
\end{theorem}
\begin{proof}
The Cartan--Leray homology spectral sequence of the covering
$S^3\to M_\Gamma$ has
$E^2_{p,q}=H_p(\Gamma,H_q(S^3;\mathbb Z))$.
Because $H_1(S^3)=H_2(S^3)=0$, its total-degree-two edge gives
$H_2(\Gamma;\mathbb Z)\cong H_2(M_\Gamma;\mathbb Z)$.
The manifold $M_\Gamma$ is closed oriented and has finite
$H_1(M_\Gamma)$; by Poincar\'e duality,
$H_2(M_\Gamma;\mathbb Z)\cong H^1(M_\Gamma;\mathbb Z)=0$.
For the divisible trivial coefficient group $U(1)$ the universal
coefficient theorem gives
$H^2(\Gamma;U(1))\cong\operatorname{Hom}(H_2(\Gamma;\mathbb Z),U(1))=0$.
The centralizers of noncentral elements of a finite subgroup of
$SU(2)$ are cyclic (intersection with a maximal torus), while a
central element has full centralizer $\Gamma$; cyclic Schur
multipliers also vanish.
The Dijkgraaf--Pasquier--Roche slant factor $\alpha_g$
is therefore a coboundary on every centralizer.  Projective
centralizer modules can be rephased to ordinary modules without
changing their irreducible counts or dimensions.  The class sizes
are unchanged by twisting, completing the proof.
\end{proof}

\begin{proposition}[Odd cohomological twists have imaginary central-flux spins]
\label{p1b:prop:v15-odd-central-flux}
Fix one of the exceptional groups and identify
$H^3(\Gamma;U(1))\cong\mathbb Z/|\Gamma|$ using the periodic
Euler/Chern generator of the defining $SU(2)$ module.
Let $k\in\mathbb Z/|\Gamma|$ label a cocycle class and let
$z=-1\in\Gamma$. For the central flux $z$, every irreducible
projective centralizer representation has a scalar ribbon twist
$\vartheta_{z,\rho}$. Then
\[
\boxed{\vartheta_{z,\rho}^{\,2}=(-1)^k.}
\tag{T15.13}\label{p1b:eq:v15-odd-spin}
\]
Thus for even $k$ these central-flux twists belong to
$\{+1,-1\}$, while for odd $k$ they belong to $\{+i,-i\}$.
This is a \emph{twist-dependent} invariant despite the all-twist
sector-count stability of the preceding theorem.
\end{proposition}
\begin{proof}
The degree-four periodicity generator can be represented by the
second Chern class of the complex defining doublet. Under
restriction to $\langle z\rangle=C_2$ the doublet is
$L\oplus L$, with $L$ the sign line; hence its $c_2$ becomes
$c_1(L)^2\ne0\in H^4(C_2;\mathbb Z)=\mathbb Z/2$.
The restriction $H^3(\Gamma;U(1))\to H^3(C_2;U(1))$
is therefore reduction of $k$ modulo two.
For a normalized $3$-cocycle representative $\omega$,
its centralizer slant cocycle is
\[
 \alpha_z(x,y)=
 \frac{\omega(z,x,y)\,\omega(x,y,z)}{\omega(x,z,y)}
 \quad(x,y\in\Gamma),
\qquad \alpha_z(z,z)=\omega(z,z,z)=(-1)^k.
\]
Because $z$ is central, $\alpha_z(x,z)=\alpha_z(z,x)$,
so its projective representative commutes with all matrices and is
scalar by irreducibility.  The projective multiplication law gives
$\rho(z)^2=\alpha_z(z,z)I$, and the ribbon twist is
$\rho(z)$ as a scalar. This proves the claimed alternatives.
\end{proof}

\begin{remark}[Scope: complete untwisted modular formula versus explicit twisted matrices]
\label{p1b:rem:v15-completeness}
Equations \eqref{p1b:eq:v15-S}--\eqref{p1b:eq:v15-T}
are an entrywise finite algorithm specifying the full \emph{untwisted}
matrices; their central block, twist-order traces and Gauss sums are
evaluated above. For a \emph{nontrivial} cocycle class,
\eqref{p1b:eq:v15-all-twist-counts} proves the unchanged rank and
dimensions, and \eqref{p1b:eq:v15-odd-spin} gives an explicit new
spin effect. An entrywise table of all nontrivial-twist modular
matrices still requires choosing normalized cocycle representatives,
projective character identifications and their full finite sums;
none is being silently claimed.  All these structures are mathematical
operators on finite-group/quotient data and do not replace the
original Hopf contact integral $4\pi^2$.
\end{remark}


\clearpage
\part{Finite exceptional reconstruction, cubature and graph readouts}
\label{p1b:part:v16-combined}
The following are \emph{mathematical} consequences of the
finite-cell, McKay, quotient-spectral and topological constructions.
They add no G0/GP principle, physical channel assignment, gauge owner,
or particle mass or lifetime.  The matrix identities can be checked by exact finite arithmetic on
the explicitly defined group and incidence matrices.

\section{Exact register transport, memory and effective resistance}
\label{p1b:sec:v16-transport}
\begin{theorem}[Full register-source resolvent and its non-lumpability defect]
\label{p1b:thm:v16-register-resolvent}
Take the $600$-chamber face adjacency $A$, the constructed
$600\times25$ chamber partition $U$, and
$R=(J_5-I_5)\otimes(J_5-I_5)$ of Theorems
\ref{p1:thm:v366-chamber-transporter-complement}--
\ref{p1:thm:v366-petersen-face-walk}. Define
\[
 S=U/\sqrt{24},\quad P=J_5/5,\quad Q=I_5-P,\quad
 P_0=P\otimes P,\quad
 P_a=P\otimes Q+Q\otimes P,\quad P_c=Q\otimes Q.
\]
Then $S^\mathsf TS=I_{25}$, $S^\mathsf TAS=R/4$, and the
$25\times25$ \emph{compressed full-chamber resolvent} is exactly
\begin{equation}
\boxed{
S^\mathsf T(zI_{600}-A)^{-1}S
 =\frac{P_0}{z-4}+\frac{P_a}{z+1}
 +\frac{z+5/4}{z^2+z-5}\,P_c.}
\tag{T16.1}\label{p1b:eq:v16-register-G}
\end{equation}
This holds outside the spectrum of $A$ (and extends as a rational
identity at all regular points of the compressed resolvent).
On each of the sixteen centred source directions the $A$-cyclic
space is two-dimensional, and in a source-first orthonormal basis
$A$ is represented by
\[
\boxed{M_c=
\begin{pmatrix}
 1/4&5\sqrt3/4\\5\sqrt3/4&-5/4
\end{pmatrix}},\quad
\lambda_\pm=\frac{-1\pm\sqrt{21}}2,
\qquad
w_\pm=\frac12\pm\frac{3}{4\sqrt{21}}.
\]
Consequently, with the source unit vector $s$ in any centred
register direction, its $n$-step face-walk return moment is
$\langle s,A^ns\rangle=w_+\lambda_+^n+w_-\lambda_-^n$.
For $\mathsf P=A/4$ the exact failure of homogeneous register
lumpability occurs already at two steps:
\begin{equation}
\boxed{S^\mathsf T\mathsf P^2S-
             (S^\mathsf T\mathsf P S)^2=\frac{75}{256}P_c.}
\tag{T16.2}\label{p1b:eq:v16-nonlump}
\end{equation}
The induced rook-complement transition $R/16$ matches a single
block-uniform one-step observation, but is \emph{not} the exact
multi-step marginal of the $600$-chamber walk.
\end{theorem}
\begin{proof}
The existing partition identity gives $U^\mathsf TU=24I$
and $U^\mathsf TAU=6R$.  The three projectors displayed above
are orthogonal and exhaust the register space; $R$ acts on them
as $16,-4,1$.  On $P_0$, $AS=4S$ since every chamber has four
face-neighbors.  On an additive centred register array
$x_{ij}=a_i+b_j$, $\sum_i a_i=\sum_j b_j=0$, the four
destination rows and columns of each tetrahedron form the
bijection of Proposition~\ref{p1:prop:v366-local-face-permutation};
therefore $(AUx)_f=-a_i-b_j=-(Ux)_f$ for $f\in\mathscr F_{ij}$.
This proves the uniform and eight additive resolvent terms.

For any unit vector $v$ in $P_c\mathbb R^{25}$ let $s=Sv$.
The pre-existing exact identity
$(A^2+A-5I)SP_c=0$ and $\langle s,As\rangle=1/4$
give $\langle s,A^2s\rangle=19/4$ and
$\|(A-\tfrac14)s\|^2=75/16$.  The normalized partner
$t=(4/(5\sqrt3))(A-\tfrac14)s$ is perpendicular to the
register source.  The monic polynomial $x^2+x-5$ gives the displayed
matrix $M_c$; different $P_c$ directions stay orthogonal because
the first two moments are scalar on $P_c$.  Inverting $zI-M_c$
proves \eqref{p1b:eq:v16-register-G}; its residues at
$\lambda_\pm$ give $w_\pm$.  The two-step identity follows
by subtracting $(R/16)^2$ from $S^\mathsf TA^2S/16$.
All source equalities are exact, not small-step approximations.
\end{proof}

\begin{theorem}[Canonical resistor graph and exact register hitting times]
\label{p1b:thm:v16-register-resistance}
Assign unit conductance to every face-contact pair between
\emph{different} chamber blocks and then sum conductances between
registers.  The resulting connected weighted register graph has
conductance matrix $W=6R$, weighted degree $96$ and Laplacian
\[L=96I-6R=6(16I-R).
\]
Its Moore--Penrose inverse is
\[\boxed{L^+=\frac1{120}P_a+\frac1{90}P_c.}\]
For distinct register vertices $u=(i,j)$ and $v=(k,l)$,
\begin{equation}
\boxed{R_{\rm eff}(u,v)=
\begin{cases}
19/900,&\text{exactly one index agrees},\\
1/50,&i\ne k\text{ and }j\ne l.
\end{cases}}
\tag{T16.3}\label{p1b:eq:v16-resistance}
\end{equation}
For the \emph{independently defined} simple random walk
$R/16$ on this $25$-register graph, the mean hitting times
from $u$ to $v$ are respectively $76/3$ and $24$ steps, and
the corresponding commute times are $152/3$ and $48$.
Neither electrical resistance nor these times is an observed decay
width or lifetime.
\end{theorem}
\begin{proof}
$R$ has spectrum $16,-4,1$ on $P_0,P_a,P_c$.
The positive Laplacian eigenvalues are thus $120$ on $P_a$ and
$90$ on $P_c$, establishing $L^+$.  The effective resistance is
$(e_u-e_v)^\mathsf TL^+(e_u-e_v)$.  For one common index
the squared norms of the $P_a$ and $P_c$ components of $e_u-e_v$
are $2/5$ and $8/5$; for no common index they are $4/5$ and
$6/5$.  These give the two rational numbers in
\eqref{p1b:eq:v16-resistance}.  The graph has total weighted
edge conductance $25\cdot96/2=1200$.  The standard electrical
commute-time identity gives $2\cdot1200 R_{\rm eff}$; register
automorphisms exchange any ordered pair and its reverse, so the
one-direction hitting time is half the commute time.
This walk is defined from the aggregate network and should not be
identified with the generally non-Markovian block observation of
Theorem~\ref{p1b:thm:v16-register-resolvent}.
\end{proof}

\section{McKay spectrum to exceptional geodesic heat reconstruction}
\label{p1b:sec:v16-spectral-reconstruction}
\begin{theorem}[Adjacency-spectrum and conjugacy-angle reconstruction]
\label{p1b:thm:v16-mckay-class-angles}
Let $A_\Gamma$ be the symmetric affine McKay adjacency from
Theorem~\ref{p1b:thm:v12-mckay-resolvent}, with defining
$V=\mathbb C^2$, and let $g_C$ range over conjugacy classes.
Its class-character vectors
$v_C=(\chi_\rho(g_C))_{\rho\in\widehat\Gamma}$ satisfy
\[A_\Gamma v_C=2\cos\theta_C\,v_C.
\]
In particular, the \emph{full} characteristic polynomials are
\begin{equation}
\begin{aligned}
\det(xI-A_{2T})&=x(x^2-4)(x^2-1)^2,\\
\det(xI-A_{2O})&=x^2(x^2-4)(x^2-1)(x^2-2),\\
\det(xI-A_{2I})&=x(x^2-4)(x^2-1)(x^4-3x^2+1).
\end{aligned}
\tag{T16.4}\label{p1b:eq:v16-mckay-class-poly}
\end{equation}
Thus the affine-McKay eigenvalues determine the class-angle
\emph{multiset} through $\theta_C=\arccos(\lambda_C/2)$.
Repeated eigenvalues retain distinct class labels only if the
character-table eigenvectors are retained; the polynomial alone does
not recover class sizes.
\end{theorem}
\begin{proof}
The representation-ring identity $V\otimes\rho_i
=\bigoplus_j (A_\Gamma)_{ij}\rho_j$ implies, on evaluation
at $g_C$, $\sum_j A_{ij}\chi_j(g_C)
=\chi_V(g_C)\chi_i(g_C)$.  The defining $SU(2)$ character
is $2\operatorname{Re}g_C=2\cos\theta_C$, proving the eigenvector
identity.  The irreducible character table is square and nonsingular
by character orthogonality, so these eigenvectors exhaust the
spectrum with multiplicity. Insert the independently established
class heights of Theorem~\ref{p1b:thm:v14-class-geodesic-atlas}
and multiply the resulting linear factors; in the $2I$ case
$\varphi^2+\varphi^{-2}=3$ gives the quartic.
\end{proof}

\begin{theorem}[Exact inverse character transform of all flat scalar spectra]
\label{p1b:thm:v16-flat-heat-inversion}
For each irreducible $\rho$, let $b_n^\rho$ be the McKay
multiplicity and let the \emph{flat-bundle scalar} heat trace
be
\[H_\rho(t)=\sum_{n\ge0}(n+1)b_n^\rho
                         e^{-tn(n+2)/R^2},\quad t>0.
\]
For every conjugacy class $C$, the complete family of these
flat-bundle traces satisfies the \emph{inverse character transform}
\begin{equation}
\boxed{
K_C(t):=\sum_\rho\chi_\rho(g_C)H_\rho(t)
 =\sum_{n\ge0}(n+1)U_n(\cos\theta_C)
           e^{-tn(n+2)/R^2}
 =\operatorname{Tr}(L_{g_C}e^{-t\Delta_{S^3(R)}}).}
\tag{T16.5}\label{p1b:eq:v16-flat-inversion}
\end{equation}
$U_n$ is the Chebyshev polynomial of the second kind, and $L_g$
uses the same left-action convention as
Proposition~\ref{p1b:prop:v14-orbital-image-trace}.
The scalar heat kernel is strictly positive; since left
multiplication is a Clifford translation, $K_C(t)$ is positive
for every $t>0$ and its short-time exponential determines the
round displacement exactly:
\begin{equation}
\boxed{R^2\theta_C^2=
       \lim_{t\downarrow0}\bigl(-4t\log K_C(t)\bigr).}
\tag{T16.6}\label{p1b:eq:v16-distance-heat}
\end{equation}
The inverse reconstruction requires the \emph{known}
representation labels and character table.  Distinct conjugacy
classes with the same quaternionic height share their recovered
round displacement; neither a unique abstract group nor unlabelled
class sizes are reconstructed from a scalar heat spectrum alone.
\end{theorem}
\begin{proof}
The group-restricted $SU(2)$ character of
$\operatorname{Sym}^n V$ decomposes as
$\chi_{\operatorname{Sym}^n V}
 =\sum_\rho b_n^\rho\chi_\rho$.
At $g_C$, the left side is $U_n(\operatorname{Re}g_C)$;
substitution proves the first equality.  By Peter--Weyl,
the scalar eigenspace of eigenvalue $n(n+2)/R^2$ on the full
$S^3$ is $\operatorname{Sym}^nV_L\otimes
(\operatorname{Sym}^nV_R)^*$, so the left $g_C$ trace is
$(n+1)\chi_{\operatorname{Sym}^nV}(g_C)$.
This proves the orbital identity and absolute convergence for $t>0$.
In kernel form the final trace is
$\int_{S^3}k_t(x,g_Cx)\,dx$; bi-invariance and the Clifford
translation property make its integrand the same positive radial
value for each $x$.  The standard local heat-kernel exponential
(or Varadhan distance limit on the compact round sphere) yields
$\lim_{t\downarrow0}-4t\log k_t(x,g_Cx)
 =d(x,g_Cx)^2=R^2\theta_C^2$; the constant sphere volume drops out.
\end{proof}

\section{ADE discriminant topology and exceptional spin counts}
\label{p1b:sec:v16-ade-discriminant}
\begin{theorem}[Resolution plumbing determines the spherical-link pairing]
\label{p1b:thm:v16-ade-link-discriminant}
Let $W_\Gamma$ be the compact simply connected plumbing
neighborhood of the exceptional curves of the minimal Kleinian
resolution of Theorem~\ref{p1b:thm:v12-resolution-link}, with
intersection matrix $Q_W=-C_{E_r}$ and boundary $M_\Gamma$
with its boundary orientation. Then
\begin{equation}
\boxed{H_1(M_\Gamma;\mathbb Z)
 \cong\operatorname{coker}(-C_{E_r})
 \cong\Gamma_{\mathrm{ab}}.}
\tag{T16.7}\label{p1b:eq:v16-ade-h1}
\end{equation}
Under the boundary linking convention
$\lambda_{\partial W}(x,y)=-\widetilde x^\mathsf T
 Q_W^{-1}\widetilde y\pmod{\mathbb Z}$, the discriminant pairing
is represented by $C_{E_r}^{-1}$ modulo integers. Its Smith
invariants and resulting spin-structure counts are
\[
\begin{array}{c|c|c|c|c}
\Gamma&E_r&\operatorname{SNF}(C_{E_r})&H_1(M_\Gamma)&
\#\operatorname{Spin}(M_\Gamma)\\\hline
2T&E_6&\operatorname{diag}(1^{(5)},3)&\mathbb Z/3&1\\
2O&E_7&\operatorname{diag}(1^{(6)},2)&\mathbb Z/2&2\\
2I&E_8&\operatorname{diag}(1^{(8)})&0&1.
\end{array}
\]
Orientation reversal negates the linking pairing but leaves its
underlying group and spin count unchanged.
\end{theorem}
\begin{proof}
The $E_r$ plumbing is a tree of disk bundles over
$S^2$, so $H_1(W)=0$ and $H_2(W)=\mathbb Z^r$; Poincar\'e--Lefschetz
duality identifies $H_2(W,\partial W)$ with $\mathbb Z^r$.
The segment of the long exact sequence of $(W,\partial W)$ is
\[H_2(W)\xrightarrow{\,-C_{E_r}\,}H_2(W,\partial W)
 \longrightarrow H_1(\partial W)\longrightarrow 0.
\]
The discriminant linking formula is the usual boundary pairing
of a nonsingular integral plumbing intersection lattice;
its displayed sign follows the declared boundary convention.
Removing the affine trivial-representation vertex from the
McKay graphs in Corollary~\ref{p1b:cor:v13-all-irrep-etas} gives
finite $E_6,E_7,E_8$ matrices whose exact integral row/column
reductions have the stated Smith forms.  Since
$\pi_1(M_\Gamma)=\Gamma$, the group is also the abelianization.
Every oriented three-manifold is spin, and its spin structures
form a torsor for $H^1(M_\Gamma;\mathbb Z/2)
=\operatorname{Hom}(H_1(M_\Gamma),\mathbb Z/2)$.
This proves the final column and independently reproduces
Corollary~\ref{p1b:cor:v12-spin-count}.
\end{proof}

\section{All-cocycle Gauss sums and protected central charge}
\label{p1b:sec:v16-all-twist-gauss}
\begin{theorem}[Complete first Gauss-sum stability under three-cocycle twists]
\label{p1b:thm:v16-all-twist-gauss}
For every finite group $\Gamma$ and any unitary normalized
$\omega\in Z^3(\Gamma;U(1))$, give the twisted double
its canonical unitary spherical/ribbon structure. Then
\begin{equation}
\boxed{\tau_+(D^\omega(\Gamma))
 =\tau_-(D^\omega(\Gamma))=|\Gamma|,
 \qquad \frac{\tau_+}{\mathcal D}=1,
 \quad\mathcal D=|\Gamma|.}
\tag{T16.8}\label{p1b:eq:v16-all-twist-gauss}
\end{equation}
For the exceptional groups the two Gauss sums remain
$24,48,120$, even when the central-flux twists are imaginary
as in Proposition~\ref{p1b:prop:v15-odd-central-flux}.
This stability does \emph{not} imply that individual twists,
$\operatorname{Tr}\mathsf T^k$, higher Gauss sums, or full
braiding matrices are independent of $\omega$.
\end{theorem}
\begin{proof}[Imported Drinfeld-center theorem, with normalization]
The standard categorical identification is
$\operatorname{Rep}(D^\omega(\Gamma))
 \simeq Z(\operatorname{Vec}_\Gamma^\omega)$.
The pointed unitary spherical fusion category
$\operatorname{Vec}_\Gamma^\omega$ has global dimension
$|\Gamma|$; its Drinfeld center has global dimension
$|\Gamma|^2$ and trivial multiplicative central charge.
The canonical-center Gauss-sum theorem of M\"uger
(\url{https://doi.org/10.1016/S0022-4049(02)00248-7}) thus gives
$\tau_+=\sqrt{|\Gamma|^2}=|\Gamma|$.
Since the ribbon category is unitary, $\tau_-=\overline{\tau_+}$.
The result does not use triviality of the centralizer Schur
multipliers and is valid beyond the binary-polyhedral cases.
\end{proof}

\section{A joint defining-flat-connection spectral/topological readout}
\label{p1b:sec:v16-flat-joint}
\begin{theorem}[Normalized Chern--Simons class and defining-doublet rho identity]
\label{p1b:thm:v16-flat-CS-eta}
Fix the standard left inclusion $\iota:\Gamma\hookrightarrow SU(2)$,
orient $S^3=SU(2)$ so the identity map has degree $+1$, and
use the natural left-quotient spin lift/orientation of
Theorem~\ref{p1b:thm:v13-eta-character}.
For the flat connection with holonomy $\iota$ define the
Chern--Simons normalization on a descended gauge potential by
\[
\operatorname{CS}(A)=\frac1{8\pi^2}\int_{M_\Gamma}
  \operatorname{Tr}\!\left(A\wedge dA+\frac23A\wedge A\wedge A\right)
 \pmod{\mathbb Z},
\]
with fundamental $2\times2$ matrix trace and the compatible
pure-gauge winding convention
$\deg(f)=(24\pi^2)^{-1}\int_{S^3}
\operatorname{Tr}((f^{-1}df)^3)$.
Then the connection has the exact class
\begin{equation}
\boxed{\operatorname{CS}(\iota)=-\frac1{|\Gamma|}
       \pmod{\mathbb Z}.}
\tag{T16.9}\label{p1b:eq:v16-CS-inc}
\end{equation}
Its defining flat-doublet Dirac eta and rho satisfy
\begin{equation}
\boxed{\eta_{D_{\Gamma,V}}(0)
 =2\eta_{D_{\Gamma,\mathbf1}}(0)-2+\frac2{|\Gamma|},\qquad
\rho_V(D_\Gamma)=-2+\frac2{|\Gamma|}.}
\tag{T16.10}\label{p1b:eq:v16-rho-inc}
\end{equation}
For these matched conventions this produces a joint congruence
\[\boxed{\rho_V(D_\Gamma)+2\operatorname{CS}(\iota)
 \equiv0\pmod{2\mathbb Z},}
\]
where $2\operatorname{CS}$ is read modulo $2\mathbb Z$.
Combining the already established \emph{acyclic} complex torsion
atlas with the defining-doublet identification yields
\[
\begin{array}{c|c|c|c|c}
\Gamma&\operatorname{CS}(\iota)&\eta_{D_{\Gamma,V}}(0)&
\rho_V(D_\Gamma)&T_{\mathbb C}(M_\Gamma,V)\\\hline
2T&-1/24&29/72&-23/12&1/2\\
2O&-1/48&101/144&-47/24&(2-\sqrt2)/2\\
2I&-1/120&73/72&-119/60&(3-\sqrt5)/4.
\end{array}
\]
The torsion column uses the defining-spinor character convention
of Theorem~\ref{p1b:thm:v15-exceptional-torsion-table};
when comparing independently named spinorial irreps the
characters, not just their dimensions, must first be matched.
\end{theorem}
\begin{proof}
Represent the flat bundle as the quotient of
$S^3\times SU(2)$ by the diagonal left $\Gamma$ action.
The map $f(q)=q$ is equivariant,
$f(\gamma q)=\iota(\gamma)f(q)$, and supplies a global
trivialization; its induced gauge potential is
$A=f^{-1}df$.  The Maurer--Cartan equation is $dA=-A\wedge A$.
Hence the normalized Chern--Simons 3-form on the cover is
$-(24\pi^2)^{-1}\operatorname{Tr}(A^3)$.
The degree of $f$ is one and the cover has degree $|\Gamma|$,
proving \eqref{p1b:eq:v16-CS-inc}.
A change of global trivialization shifts this value by an integer,
so only its residue class is intrinsic.

For the eta identity use the previously proved formula
\eqref{p1b:eq:v13-eta-character} and
$\chi_V(g)=2\operatorname{Re}g$:
\[
\eta_V-2\eta_{\mathbf1}
 =\frac1{|\Gamma|}\sum_{g\ne1}
   \frac{2\operatorname{Re}g-2}{1-\operatorname{Re}g}
 =-\frac{2(|\Gamma|-1)}{|\Gamma|}.
\]
This proves \eqref{p1b:eq:v16-rho-inc}; combining it with
\eqref{p1b:eq:v16-CS-inc} gives the congruence (the displayed
representatives sum to $-2$).
Insert the existing exact eta fractions of
Corollary~\ref{p1b:cor:v13-exact-eta-values} and the
character-matched torsion values in \eqref{p1b:eq:v15-torsion-table}
for the three rows.  No index-theoretic comparison of an arbitrary
connection with an unrelated spectral operator is assumed.
\end{proof}
\begin{remark}[Scope of the joint readout]
\label{p1b:rem:v16-joint-owner-boundary}
The Chern--Simons entry is an element of $\mathbb R/\mathbb Z$,
not a canonically real-valued energy; the eta entries are for the
specified round spin lift; and torsion is an acyclic-flat-bundle
positive number.  The other reducible and irreducible $SU(2)$
holonomies enumerated in Corollary~\ref{p1b:cor:v14-flat-su2-counts}
are not assigned uncomputed Chern--Simons invariants or Floer groups.
In particular no numerical relation between this joint mathematical
readout and a physical decay amplitude is asserted.
\end{remark}

\section{A first exact sampled spinorial Dirac band}
\label{p1b:sec:v16-finite-spin}
\begin{theorem}[Quaternionic Hurwitz sampling of both lowest Dirac signs]
\label{p1b:thm:v16-hurwitz-spin-band}
Use the $24$ unit Hurwitz vertices $H=2T$, the adjacency $A_H$
with $q\sim q'$ iff $\langle q,q'\rangle=1/2$, and the
standard left spinor trivialization of the round $S^3(R)$.
For $u,v\in\mathbb C^2$ define sampled spinors
\[
\Psi_+(u)(q)=u,\qquad
\Psi_-(v)(q)=q^{-1}v,\qquad q\in H,
\]
where $q$ is its fundamental unitary $2\times2$ matrix.
The normalized sampling map $\Psi\mapsto24^{-1/2}\Psi|_H$
is an isometry from the orthogonal direct sum of these two
$2$-dimensional lowest round Dirac eigenspaces into
$\mathbb C^{H}\otimes\mathbb C^2$.  On their four-dimensional
sampled image the \emph{finite} operator
\begin{equation}
\boxed{D_{H,0}=\frac{3}{4R}
  (A_H\otimes I_2-6I_{48})\big|_{\mathscr K_+\oplus\mathscr K_-}}
\tag{T16.11}\label{p1b:eq:v16-first-spin-band}
\end{equation}
acts as $+3/(2R)$ with multiplicity two on
$\mathscr K_+$ and $-3/(2R)$ with multiplicity two on
$\mathscr K_-$.  The two sampled projectors have entrywise kernels
\[
  (P_+)_{q,q'}=I_2/24,\qquad
  (P_-)_{q,q'}=q^{-1}q'/24.
\]
This is a \emph{four-dimensional band intertwining}, not a
formula for the full $48$-dimensional finite spin Dirac operator.
It does not establish higher spinorial bands, a covariant discrete
spin connection, or an autonomous physical fermion action.
\end{theorem}
\begin{proof}
The five-design result of
Theorem~\ref{p1:thm:v350-hurwitz-five-design} includes
$\sum_{q\in H}q=0$.  Every $q$ is unitary, so the two
sampling maps have individual identity Gram matrices and mutual
Gram matrix zero.  Their orthogonal projectors are the two
stated rank-two kernels.
The neighbor-moment identity from
Theorem~\ref{p1:thm:v350-hurwitz-scalar-laplacian}
gives $\sum_{q'\sim q}1=8$ and
$\sum_{q'\sim q}(q')^{-1}=4q^{-1}$, since quaternionic
conjugation is real-linear.  Hence $A_H\otimes I_2$ has
eigenvalues $8$ and $4$ on the two sampled spaces, respectively.
The affine polynomial in \eqref{p1b:eq:v16-first-spin-band}
therefore has the stated signed eigenvalues.
Under $SU(2)_L\times SU(2)_R$ the continuous functions
$u$ and $q^{-1}v$ transform as the distinct
$\operatorname{Sym}^0 V_L\otimes\operatorname{Sym}^1 V_R$
and $\operatorname{Sym}^1 V_L\otimes\operatorname{Sym}^0 V_R$
lowest spinorial types; the sign assignment agrees with the
chosen convention of
Theorem~\ref{p1b:thm:v13-dirac-generating}.
No statement about the remaining sampled modes follows.
\end{proof}
\begin{remark}[Scalar polynomial coincidence versus genuine spinorial extension]
\label{p1b:rem:v16-spin-scope}
The operator in \eqref{p1b:eq:v16-first-spin-band} is obtained
by restricting a scalar adjacency polynomial to two already
identified spinor eigenspaces.  Its successful lowest-band signs
must not be extrapolated to the full spin spectrum: from the next
band onward the left/right spinor types overlap scalar polynomial
degrees and require an explicit connection or further finite
operator structure.  The previous scope warning
\ref{p1:rem:v350-hurwitz-spectral-domain} is preserved rather
than silently replaced by a global Dirac equivalence.
\end{remark}



\section{Left-invariant continuation of the first Hurwitz spinor band}
\label{p1b:sec:v17-anisotropic-hurwitz-spin-bridge}
\begin{theorem}[Exact two-band Hurwitz sampling for every Milnor triple]
\label{p1b:thm:v17-general-metric-hurwitz-spin-band}
Adopt the positively oriented \emph{left-invariant} metric-adapted
spinor trivialization of
Theorem~\ref{p1:thm:general-left-invariant-hodge-dirac-compiler},
with $g_{a,b,c}$ for any $a,b,c>0$ and
\[
  \kappa=\frac{a^2+b^2+c^2}{4abc},\qquad
  s=\frac1a+\frac1b+\frac1c.
\]
Regard the unit quaternion $q$ as its fundamental $2\times2$ unitary
matrix, choose $T_i=-i\tau_i/2$ with $[T_1,T_2]=T_3$ cyclically,
and set $E_i=T_i/a_i$.  For constant $u,v\in\mathbb C^2$, define
\[
  \psi_+(u)(q)=u,\qquad \psi_-(v)(q)=q^{-1}v.
\]
The two function spaces are invariant under the \emph{full continuous}
signed Dirac operator of the specified metric, with eigenvalues
\begin{equation}\label{p1b:eq:v17-aniso-spinor-eigenvalues}
\boxed{\lambda_+(a,b,c)=\kappa,\qquad
       \lambda_-(a,b,c)=\kappa-\frac{s}{2}
       =\frac{a^2+b^2+c^2-2(ab+bc+ca)}{4abc}.}
\tag{T16.12}
\end{equation}
Each eigenvalue has multiplicity two \emph{within the displayed
representation-graded bands}; there may be coincident eigenvalues
elsewhere in the full spectrum.  They correspond, respectively, to the
$j=0$ spinor block and to the singlet in the $j=\frac12$ spinor block
of Paper~I-A's complete compiler.

Let $H=2T$ be the same fixed Hurwitz $24$-point set as in
Theorem~\ref{p1b:thm:v16-hurwitz-spin-band}, with $A_H$ its
\emph{round-reference} degree-eight quaternionic-neighbor adjacency.
Normalized sampling of $\psi_\pm$ into
$\mathscr K_+\oplus\mathscr K_-\subset\mathbb C^H\otimes\mathbb C^2$
is an isometry for the normalized left-invariant metric volume measure.
For every positive Milnor triple, the four-dimensional restriction of
\begin{equation}\label{p1b:eq:v17-sampled-aniso-dirac}
\boxed{D_{H,0}(a,b,c)=
\left[\kappa I_{48}+\frac{s}{8}(A_H\otimes I_2-8I_{48})\right]
\big|_{\mathscr K_+\oplus\mathscr K_-}}
\end{equation}
intertwines \emph{exactly} with the continuous Dirac operator on these
two bands.  No assertion of isometry of the non-round metric by the
fixed neighbor graph, or of equality with the Dirac operator on the
other $44$ sampled dimensions, is made.

On the round slice $a=b=c=R/2$ this reduces precisely to
\eqref{p1b:eq:v16-first-spin-band}, with signed eigenvalues
$\pm3/(2R)$.  On the axisymmetric left-invariant slice
$a=b=r$, $c=t$, the sampled $j=\frac12$ singlet eigenvalue is
\[
   \lambda_-(r,r,t)=\frac{t-4r}{4r^2};
\]
thus this \emph{specified geometric Dirac band} has a zero crossing
exactly at $t=4r$.  Such an eigenvalue crossing is a mathematical
spectral statement, not a physical massless-fermion selection rule.
\end{theorem}
\begin{proof}
For the left-invariant vector field $E_i$, the integral curve through
$q$ is $q\exp(tT_i/a_i)$.  Differentiating the defining spinors gives
\[
  E_i\psi_+(u)=0,\qquad
  E_i\psi_-(v)=-\frac{T_i}{a_i}q^{-1}v
                  =\frac{i\tau_i}{2a_i}\psi_-(v).
\]
The signed spinor connection in the left-invariant trivialization
of Paper~I-A is
$D_{a,b,c}=\sum_i i\tau_iE_i+\kappa I_2$.
Since $\tau_i^2=I_2$, it follows \emph{term by term} that
$D\psi_+=\kappa\psi_+$ and
$D\psi_-=(\kappa-s/2)\psi_-$.
Both families have dimension two, are orthogonal for the normalized
Haar measure, and belong to the indicated Peter--Weyl degrees.
The normalized volume measure of every left-invariant metric is
normalized Haar measure, so these statements are valid for all
positive triples.

The $24$ Hurwitz elements satisfy $\sum_{q\in H}q=0$, and each is
unitary.  Consequently the finite normalized sample of each family
is isometric and the two sampled images are orthogonal; this is
also the Gram calculation in
Theorem~\ref{p1b:thm:v16-hurwitz-spin-band}.
The exact nearest-neighbor identities give
$A_H\psi_+=8\psi_+$ and $A_H\psi_-=4\psi_-$ on the finite samples.
Thus the displayed affine polynomial in $A_H$ acts as $\kappa$ on
$\mathscr K_+$ and as
$\kappa+\frac{s}{8}(4-8)=\kappa-s/2$ on
$\mathscr K_-$.  These are the same exact eigenvalues as the full
continuous Dirac operator on the two represented spaces, proving
the intertwining for every triple.  Substitution of
$a=b=c=R/2$ and $(a,b,c)=(r,r,t)$ gives the stated special cases.
\end{proof}
\begin{remark}[Non-round scope and global spectral ordering]
\label{p1b:rem:v17-aniso-spinor-scope}
For a generic anisotropic metric, these are persistent exact spinor
\emph{bands}, not necessarily the lowest eigenvalues, nor necessarily
of opposite signs.  The $24$-point quaternionic adjacency uses the
round reference positions, not the geodesic-neighbor operator of an
anisotropic metric.  It therefore makes no all-band or metric-dependent
finite Dirac claim.  Its value is an exact, input-matched sampling
intertwiner for two identified continuous eigenspaces and their
spectral-flow zero condition.  Physical state and interaction owners
remain downstream obligations.
\end{remark}


\section{Complete local octahedral response of an unmarked tetrahedral frame}
\label{p1b:sec:v18-complete-octahedral-response}
The round momentum parent remains $S^3\simeq\operatorname{Sp}(1)$.  The
following finite response stabilizer is not a replacement of that parent,
and the group stabilizer of an observable is not, by itself, an occupied
physical particle Hilbert space.

\begin{theorem}[First Born-even anisotropic response of the four Clifford axes]
\label{p1b:thm:v18-born-even-octahedral-response}
Let $H\cong2T$ be a fixed Hurwitz subgroup and choose the four
unoriented unit axes represented by
$u_1=(1,1,1)/\sqrt3$, $u_2=(1,-1,-1)/\sqrt3$,
$u_3=(-1,1,-1)/\sqrt3$, and $u_4=(-1,-1,1)/\sqrt3$ in
$V=\operatorname{Im}\mathbb H$.  For $v=(x,y,z)\in V$ put
\[
 R_a(v)=\frac14\bigl(\|v\|^2-(u_a\cdot v)^2\bigr),\quad
 \overline R(v)=\frac14\sum_{a=1}^4R_a(v),\quad
 \mathcal V_H(v)=\sum_{a=1}^4(R_a(v)-\overline R(v))^2.
\]
Then
\[
 \boxed{\overline R(v)=\tfrac16\|v\|^2,\qquad
 \mathcal V_H(v)=\tfrac19(x^2y^2+y^2z^2+z^2x^2),}\qquad
 \boxed{\operatorname{Stab}_{\operatorname{Sp}(1)}(\mathcal V_H)
       =N_{\operatorname{Sp}(1)}(H)=2O.}
\]
The response is even under $v\mapsto-v$ and is the first
nonconstant, nonradial, \emph{even harmonic response class} of the
unmarked tetrahedral frame: its harmonic part is proportional to
$x^4+y^4+z^4-3\|v\|^4/5$.
\end{theorem}
\begin{proof}
The tight-frame identity $\sum_a u_au_a^{\mathsf T}=4I/3$
gives $\overline R=\|v\|^2/6$.  Expanding the four squared deviations
gives the quartic displayed above.  Its zero locus on the unit sphere
consists of the six coordinate directions $\{\pm e_1,\pm e_2,\pm e_3\}$.
Any proper rotation preserving the quartic must permute these three
unoriented axes and therefore belongs to the octahedral rotation group.
Conversely every orientation-preserving signed coordinate permutation
preserves the quartic.  Its inverse image in $\operatorname{Sp}(1)$ is
$2O=N_{\operatorname{Sp}(1)}(H)$ by the earlier normalizer theorem.
Evenness is immediate.  The tetrahedral invariant harmonic spaces in
degrees one and two vanish; the degree-three invariant $xyz$ is odd,
and the degree-four nonradial even invariant is one-dimensional.
\end{proof}

\begin{proposition}[Complete octahedral coset and the attached sixfold normalizer]
\label{p1b:prop:v18-octahedral-sixfold-attachment}
With $H=2T$ and $X=2O\setminus H$, the $48$ elements split as
\[
 \boxed{2O=\underbrace{Q_8\sqcup(H\setminus Q_8)}_{8+16=24}
          \sqcup\underbrace{X_+\sqcup X_0\sqcup X_-}_{6+12+6=24},}
\]
where $X_{\pm}=\{q\in X:\operatorname{Re}q=\pm1/\sqrt2\}$
and $X_0=\{q\in X:\operatorname{Re}q=0\}$.  For
$u=(i+j+k)/\sqrt3$, put $c=(1+i+j+k)/2$ and
\[
 B_6(u)=\left\{\tfrac{\pm(i-j)}{\sqrt2},
                   \tfrac{\pm(j-k)}{\sqrt2},
                   \tfrac{\pm(k-i)}{\sqrt2}\right\}.
\]
Then
\[
 C_6(u)=\langle c\rangle\subset H,\quad
 N_{2O}(C_6(u))=C_6(u)\sqcup B_6(u)\cong2D_3,
 \quad B_6(u)\subset X_0,\quad
 \operatorname{span}_{\mathbb R}B_6(u)=u^\perp.
\]
Thus the attached $6$ is \emph{not an additional six group elements
outside the complete $48$}.  It is a sixfold normalizer interface
across the two existing $24$-element pieces.
\end{proposition}
\begin{proof}
The binary octahedral complement consists of the $24$ unit
quaternions $(\pm e_a\pm e_b)/\sqrt2$ with $0\le a<b\le3$
and $(e_0,e_1,e_2,e_3)=(1,i,j,k)$.  Separating the real
coordinate gives the stated $6+12+6$ counts.  The Hurwitz subgroup
has commutator subgroup $Q_8$ of order eight.  A member of $B_6(u)$
reverses $u$ by quaternionic conjugation, has square $-1=c^3$,
and conjugates $c$ to $c^{-1}$.  These six members form the nontrivial
coset of $C_6(u)$ in its order-twelve normalizer; they span the
imaginary plane perpendicular to $u$.  The direct group calculation
establishes every displayed identity.
\end{proof}

\begin{remark}[Scope of reciprocal and particle readings]
\label{p1b:rem:v18-particle-response-frontier}
The quartic stabilizer is an exact \emph{mathematical} response symmetry.
Born inversion of the relative coordinate makes odd scalar Taylor
terms invisible only for a source that is genuinely Born-even; it does
not forbid higher tetrahedral-even terms or physically actualize all
$48$ octahedral elements as particles.  The five $2T$ cores of a
supplied $2I$ packet have five ambient $2O$ normalizers, but
$2O\not\subset2I$; mobile $2I$ transport and local $2O$ reorientation
are different operations.  Physical occurrence remains an explicitly
typed downstream source condition, not a new foundational GP.
\end{remark}



\section{Finite exceptional incidence, cubature and operator constructions}
This section collects exact finite-geometry statements and conditional
operator comparisons.  A physical source, coupling, decay law or observer
is not determined by their finite algebra alone.


\subsection{Directed core-pair polarisation and the coherent nonet Gram}
\label{p1:sec:cc12-directed-coherent-nonet}

\begin{theorem}[Twenty directed axes and an exact rank-nine coherent Gram]
\label{p1:thm:cc12-directed-coherent-gram}
Take the ten exact projective $C_6$ axes $u_{ij}$ ($i<j$) in
Theorem~\ref{p1:thm:v285-ten-axis-quadrupole-exact-sequence}, in the
printed order $12,13,14,15,23,24,25,34,35,45$. Choose the signs
\[
 (\epsilon_{12},\ldots,\epsilon_{45})
 =(-,+,+,-,+,-,+,-,+,-),\quad
 n_{ij}=\epsilon_{ij}u_{ij},\quad n_{ji}=-n_{ij}.
\]
This choice is $A_5$-equivariant under the simultaneous ordinary
icosahedral action on the vector axes and the ordered core-pair labels;
there are precisely two such choices, exchanged by reversing all
$n_{ij}$. For every $i$ the outgoing four vectors form a regular
oriented tetrahedron:
\[
 \sum_{j\ne i}n_{ij}=0,\quad
 n_{ij}\cdot n_{ik}=-\tfrac13\ (j\ne k),\quad
 \sum_{j\ne i}n_{ij}n_{ij}^{\mathsf T}=\tfrac43I_3.
\]
Write $A(n)w=n\times w$. On $\mathbb C^3$ define
\[
 H_{ij}=\tfrac12\big(I_3-n_{ij}n_{ij}^{\mathsf T}+iA(n_{ij})\big).
\]
Every $H_{ij}$ is a positive Hermitian \emph{rank-one projection}.
The real synthesis $S:\mathbb R^{20}\to\operatorname{Herm}_3(\mathbb C)$,
$Se_{ij}=H_{ij}$, has Gram matrix
\[
 G_{ij,kl}=\operatorname{Tr}(H_{ij}H_{kl})
 =\frac{(1+n_{ij}\cdot n_{kl})^2}{4},\qquad
 \boxed{\operatorname{Spec}(G)=
 \{(20/3)^{(1)},(10/3)^{(3)},(2/3)^{(5)},0^{(11)}\}.}
\]
Its image is exactly $\mathbf1\oplus\mathbf3\oplus\mathbf5$
under $A_5$ and its kernel is $\mathbf3'\oplus2\mathbf4$.
\end{theorem}
\begin{proof}
The projective-axis stabilizer is $D_3$, while the stabilizer of
an ordered pair is its order-three rotation subgroup.  Choosing the
sign of one fixed axis and transporting by $A_5$ therefore fixes the
entire equivariant directed orientation, with only its global reversal
remaining. Substitution into the ten displayed exact axes proves the
four-vector tetrahedral identities. For $B=iA(n)$, $B^2=I-nn^T$,
$Bn=0$ and $B^3=B$. Hence $H=(B^2+B)/2$ is the spectral projector
onto the $+1$ eigenline of $B$, of rank one. Direct multiplication
using $\operatorname{Tr}(iA(n)iA(m))=2n\cdot m$ gives the Gram formula.

Let $U$ be the $3\times10$ matrix of directed unit columns for $i<j$,
let $\mathsf J$ be the $10\times10$ all-one matrix, and let $P_5$ be
the already-established Petersen quadrupole rank-five projection.
In the normalized even and odd combinations
$(e_{ij}\pm e_{ji})/\sqrt2$, the Gram splits exactly as
\[
 \boxed{G_+=\tfrac23(\mathsf J+P_5),\qquad G_-=U^TU.}
\]
Paper I-B gives $P_5^2=P_5$, $P_5\mathsf J=0$ and
$UU^T=\tfrac{10}{3}I_3$.  The even block has eigenvalues $20/3$
(multiplicity one), $2/3$ (multiplicity five) and zero (multiplicity
four); the odd block has $10/3$ (multiplicity three) and zero
(multiplicity seven). Its kernel therefore has $\mathbf4$ on the
even side and $\mathbf4\oplus\mathbf3'$ on the odd side, as claimed.
\end{proof}

\begin{corollary}[A single Gram polynomial produces the existing nonet]
\label{p1:cor:cc12-nonet-projector-polynomial}
With $G$ as above, the exact orthogonal rank-nine projector is
\[
 \boxed{P_9=I_{20}-
   (I_{20}-\tfrac3{20}G)
   (I_{20}-\tfrac3{10}G)
   (I_{20}-\tfrac32G).}
\]
It is independent of the common reversal of the twenty oriented axes.
The polar map
\[
 W=(SS^*)^{-1/2}S:\mathbb R^{20}\longrightarrow
 \operatorname{Herm}_3(\mathbb C)
\]
(where $(SS^*)^{-1/2}$ is the positive inverse on the full
nine-dimensional target) obeys $W^*W=P_9$, $WW^*=I_9$.
Thus the geometric projectors furnish an explicit equivariant isometry
on the nine-dimensional comparison quotient, fixing its relative
Hilbert--Schmidt normalization. This is not a physical meson
Hamiltonian or a symmetry-breaking assignment.
\end{corollary}
\begin{proof}
The displayed polynomial is one at each of the three positive Gram
eigenvalues and zero at the zero eigenvalue. The polar identities
follow from $SS^*>0$. Reversing every $n_{ij}$ conjugates the
Hermitian projector synthesis by matrix complex conjugation, an
orthogonal real transformation, leaving $S^*S$ and hence $P_9$
unchanged.
\end{proof}

\subsection{Exact exponential memory after the 600-chamber register observation}
\label{p1b:sec:cc13-register-memory}

\begin{theorem}[Closed 25-register Volterra law with no new coefficient]
\label{p1b:thm:cc13-register-memory}
Retain the $600\times600$ symmetric four-regular face adjacency $A$, the
$600\times25$ chamber-to-register incidence $U$, and the three orthogonal
register projectors $P_0,P_a,P_c$ of
Theorem~\ref{p1b:thm:v16-register-resolvent}.  Put
$S=U/\sqrt{24}$, $R=(J_5-I_5)\otimes(J_5-I_5)$,
$\mathsf L=A/4-I_{600}$, and define the observed, initially block-uniform
continuous-time evolution
\[
 M(t)=S^{\mathsf T}e^{t\mathsf L}S,\qquad x(t)=M(t)x_0,
 \qquad t\geq0.
\]
Then the complete $25$-component observation satisfies the exact
convolution equation
\begin{equation}\label{p1b:eq:cc13-volterra}
 \boxed{\frac{d}{dt}x(t)=\left(\frac R{16}-I_{25}\right)x(t)
 +\frac{75}{256}\int_0^t e^{-21(t-s)/16}P_cx(s)\,ds,\quad x(0)=x_0.}
\end{equation}
The initial hidden component is zero because the preparation is in
$\operatorname{im} S$.  The memory operator is self-adjoint and positive
semidefinite for $t\geq0$, has rank $16$, and has one time-decay rate
$21/16$.  Neither a Markov closure on the registers nor a physical time
parameter is asserted.
\end{theorem}
\begin{proof}
The source theorem supplies $AS=4S$ on $P_0$, $AS=-S$ on $P_a$, and,
for each unit $v\in P_c\mathbb R^{25}$, an orthonormal source--hidden pair
$\{s=Sv,h\}$ on which
\[
 A\big|_{\operatorname{span}\{s,h\}}=
 \begin{pmatrix}1/4&5\sqrt3/4\\5\sqrt3/4&-5/4\end{pmatrix}.
\]
The $600$-state Markov generator $\mathsf L=A/4-I$ therefore restricts to
\[
 \begin{pmatrix}-15/16&5\sqrt3/16\\5\sqrt3/16&-21/16\end{pmatrix}.
\]
Writing the pair components as $x_c,y_c$ with $y_c(0)=0$ gives
$\dot x_c=-15x_c/16+(5\sqrt3/16)y_c$ and
$\dot y_c=(5\sqrt3/16)x_c-21y_c/16$.  Variation of constants gives
$y_c(t)=(5\sqrt3/16)\int_0^t e^{-21(t-s)/16}x_c(s)ds$.
Inserting this expression yields the coefficient
$(5\sqrt3/16)^2=75/256$.  On $P_0$ and $P_a$ the equations are
$\dot x_0=0$ and $\dot x_a=-5x_a/4$; the direct register generator
$R/16-I$ has exactly the scalar values $0,-5/4,-15/16$ on
$P_0,P_a,P_c$.  Thus the three blocks combine into
\eqref{p1b:eq:cc13-volterra}.  Positivity and rank follow because
$P_c$ is an orthogonal rank-sixteen projector.
\end{proof}

\begin{corollary}[Exact two-rate relaxation and provable slowdown]
\label{p1b:cor:cc13-relaxation}
The observed semigroup-compression family (which is \emph{not} itself a
semigroup) is
\[
 \boxed{M(t)=P_0+e^{-5t/4}P_a+f_c(t)P_c,}
\]
where
\[
 f_c(t)=w_+e^{r_+t}+w_-e^{r_-t},\quad
 r_\pm=\frac{-9\pm\sqrt{21}}8,\qquad
 w_\pm=\frac12\pm\frac{3}{4\sqrt{21}}.
\]
For every $t>0$ one has
\[
 \boxed{e^{-15t/16}<f_c(t)\le e^{(-9+\sqrt{21})t/8},\qquad
 \|M(t)-P_0\|_{2\to2}=f_c(t).}
\]
Thus replacing the hidden chamber dynamics by the one-step register
Markov generator $R/16-I$ strictly underestimates the centred
$P_c$ return amplitude at every positive time.
\end{corollary}
\begin{proof}
Diagonalize the two-by-two block in the preceding proof; the source-vector
spectral weights are exactly the displayed $w_\pm$, both positive, and
their weighted mean rate equals $-15/16$.  Strict Jensen convexity of
$e^{rt}$ gives $f_c(t)>e^{-15t/16}$; the maximal-rate bound gives
$f_c(t)\le e^{r_+t}$.  Since $e^{-5t/4}<e^{-15t/16}<f_c(t)$, the norm
of the centred register return operator is $f_c(t)$.
\end{proof}

\begin{remark}[Epistemic and dynamical boundary]
This is an exact $600$-chamber \emph{comparison walk}, not the Hamiltonian
or kinetic generator of a physical system.  The positive-semidefinite
matrix-valued Volterra kernel need not be entrywise nonnegative.  A
physical memory or relaxation-time prediction requires an existing
same-action process owner and a proved relation between its clock and
this walk.  Neither a decay width nor an additional GP principle follows.
\end{remark}


\begin{theorem}[Failure of time-inhomogeneous Markov divisibility]
\label{p1b:thm:cc14-no-time-local-markov}
On the exact $25$-register coarse return $M(t)$ of
Theorem~\ref{p1b:thm:cc13-register-memory}, with the continuous-time
full-chamber generator $A/4-I_{600}$, let
$P=J_5/5$, $Q=I_5-P$ and
$P_a=P\otimes Q+Q\otimes P$, $P_c=Q\otimes Q$.
For every finite $t\ge0$, $M(t)$ is an invertible nonnegative
bistochastic matrix.  Its unique time-local generator is
\[
 \boxed{L_{\rm eff}(t)=\dot M(t)M(t)^{-1}
 =-\tfrac54P_a-h(t)P_c,\qquad
 h(t)=-\frac{d}{dt}\log f_c(t).}
\]
Let $r_\pm=(-9\pm\sqrt{21})/8$ and
$w_\pm=\tfrac12\pm3/(4\sqrt{21})$ so that
$f_c(t)=w_+e^{r_+t}+w_-e^{r_-t}$.  Then $h(0)=15/16$ and
$h'(t)=-\operatorname{Var}_{\omega(t)}(r_+,r_-)<0$
for every finite $t$, where
$\omega_\pm(t)=w_\pm e^{r_\pm t}/f_c(t)$.
In particular, for distinct registers $(i,j)$ and $(i,k)$
sharing the same first index,
\[
 \boxed{(L_{\rm eff}(t))_{(i,j),(i,k)}
 =\frac{4h(t)-15/4}{25}<0\quad(t>0).}
\]
Thus for every $t>0$ and all sufficiently small $\varepsilon>0$,
$M(t+\varepsilon)M(t)^{-1}$ has a negative off-diagonal entry.
The exact coarse family is not divisible by stochastic transition
matrices even with time-dependent Markov rates.
\end{theorem}
\begin{proof}
$M(t)=S^{\mathsf T}e^{t(A/4-I)}S$ averages a finite regular
Markov kernel over equal-sized chamber blocks, hence is
nonnegative and bistochastic.  The exact spectral form in the
preceding theorem is
$M(t)=P_0+e^{-5t/4}P_a+f_c(t)P_c$; its three
eigenvalues are positive for finite $t$, so it is invertible.
Differentiating the spectral expression gives $L_{\rm eff}$.
The derivative of a log-sum-exp is the negative variance of its
rates after the sign in $h$, and $w_\pm>0$, $r_+\ne r_-$,
so $h'(t)<0$ and $h(t)<h(0)=15/16$ for $t>0$.
For $j\ne k$, direct multiplication of the tensor-product
projectors gives $(P_a)_{(i,j),(i,k)}=3/25$ and
$(P_c)_{(i,j),(i,k)}=-4/25$.  Hence the stated entry is
$-\frac54(3/25)+4h(t)/25$ and is negative.
The unique intermediate propagator is $M(t+\varepsilon)M(t)^{-1}$;
its off-diagonal entry is $\varepsilon (L_{\rm eff})_{uv}
+O(\varepsilon^2)$, proving negative entries for sufficiently
small positive increments.
\end{proof}
\begin{remark}[Time-local negativity is not physical negative probability]
Each $M(t)$ is a stochastic matrix with nonnegative entries.
Only the would-be \emph{intermediate} register-only transition has
negative components.  The full $600$-chamber dynamics remains
Markovian.  This classical non-divisibility is not by itself a
statement about quantum CP divisibility, physical measurement time,
particle decay, or entropy production.
\end{remark}

\begin{corollary}[Monotone quadratic distinguishability with non-Markovian memory]
\label{p1b:cor:cc15-no-quadratic-backflow}
For any two initial probability vectors $p,q$ in the exact
25-register comparison of Theorem~\ref{p1b:thm:cc14-no-time-local-markov},
let $\Delta=p-q$ and $M(t)$ be the block-averaged coarse return. Then
\[
 \boxed{\|M(t)\Delta\|_2^2=
 e^{-5t/2}\|P_a\Delta\|_2^2+
 f_c(t)^2\|P_c\Delta\|_2^2.}
\]
For every $p\ne q$ this expression is strictly decreasing on
$[0,\infty)$, even though $M$ is \emph{not} classically
P-divisible for any positive intermediate time.
In particular, no witness based only on an increase of this quadratic
distance can detect the proved failure of Markov divisibility.
\end{corollary}
\begin{proof}
A difference of normalized probability vectors has zero total mass,
so $P_0\Delta=0$. The three mutually orthogonal projectors in the
source spectral decomposition give the displayed squared norm.
Since $r_+,r_-<0$ and $w_+,w_->0$, the function
$f_c(t)=w_+e^{r_+t}+w_-e^{r_-t}$ is strictly positive with
$f'_c(t)<0$.  Likewise $e^{-5t/2}$ strictly decreases. At least one
of $P_a\Delta,P_c\Delta$ is nonzero when $\Delta\ne0$.
Thus the derivative is strictly negative at every finite time.
Nondivisibility remains exactly the negative-intermediate-rate
result of the preceding theorem; monotonicity of one norm is not
a converse of stochastic divisibility.
\end{proof}
\begin{remark}[No assertion about trace norm, entropy, or physical decoherence]
This is an exact Euclidean-distance statement for the registered
classical probabilities. It does not prove monotonicity of Shannon
entropy, total variation, or quantum trace distance, and the time
parameter remains that of the auxiliary chamber walk.
\end{remark}


\begin{theorem}[Minimal 41-state linear realization of the 600-chamber register return]
\label{p1b:thm:cc16-minimal-41state-memory}
Let $M(t)=P_0+e^{-5t/4}P_a+f_c(t)P_c$ be the entire
$25\times25$ registered-return family in
Theorem~\ref{p1b:thm:cc13-register-memory}, with
$f_c(t)=w_+e^{r_+t}+w_-e^{r_-t}$ and the source-normalized
$P_c$ rank sixteen. Among all finite-dimensional real linear
\emph{time-invariant} realizations of the full matrix-valued family
$M(t)=C e^{tG}B$ for every $t\ge0$, the smallest state-space
dimension is exactly
\[
 \boxed{N_{\min}=1+8+16+16=41.}
\]
In particular, exactly sixteen hidden linear coordinates suffice
beyond the 25 observed register coordinates, and fewer than sixteen
cannot reproduce the complete family within that realization class.
An explicit 41-state realization takes $x\in\mathbb R^{25}$,
$y\in\operatorname{im}P_c\cong\mathbb R^{16}$,
$y(0)=0$, and $b=5\sqrt3/16$, with
\[
 \boxed{\dot x=(-\tfrac54P_a-\tfrac{15}{16}P_c)x+b y,\qquad
        \dot y=bP_cx-\tfrac{21}{16}y.}
\]
Here $y$ is identified isometrically with the rank-sixteen subspace
of $\mathbb R^{25}$. This is a minimal \emph{linear} realization;
it is not asserted to be a 41-state positive Markov chain or a
reconstruction of all unobserved chamber correlations.
\end{theorem}
\begin{proof}
The Laplace transform of the full return family has the simple-pole
partial fractions
\[
 \widehat M(s)=\frac{P_0}{s}+\frac{P_a}{s+5/4}
     +\frac{w_+P_c}{s-r_+}+\frac{w_-P_c}{s-r_-}.
\]
The four poles $0,-5/4,r_+,r_-$ are distinct and both $w_\pm$
are positive. For a rational matrix transfer function with distinct
simple poles, the McMillan degree is the sum of the ranks of its
residue matrices, which here are $1,8,16,16$. This proves the
41-dimensional lower bound for every time-invariant realization.
The displayed coupled system has 25 observed and 16 hidden
coordinates. On $P_0$ it is stationary, on $P_a$ it has rate $-5/4$,
and each of the 16 centered modes has exactly the symmetric $2\times2$
generator block of Theorem~\ref{p1b:thm:cc13-register-memory}.
Preparing $y(0)=0$ recovers the complete family $M(t)$; therefore
the 41-dimensional lower bound is attained.
\end{proof}
\begin{remark}[Compression does not imply a positive hidden-state model]
The minimality is for real finite-dimensional linear, time-invariant
input/output realizations of the complete 25-by-25 family. It does
not claim that arbitrary chamber initial distributions, multitime
measurements, or an entrywise-positive 41-state Markov representation
are reproduced. A physical memory kernel still requires a
source-preserving time and process identification.
\end{remark}

\subsection{A finite isometric frame for degree-five angular and Hopf-basic modes}
\label{p1b:sec:cc13-band-frame}

\begin{theorem}[Exact $36\to62$ isometric spherical sampling]
\label{p1b:thm:cc13-band-frame}
Let $\{(u_i,w_i)\}_{i=1}^{62}$ be the positive weighted
$A_5$-orbit rule of
Theorem~\ref{p1b:thm:new-62-point-angular-cubature}, with
\[
 w_i=\begin{cases}
 81/(462\cdot20)&u_i\in\mathcal F,\\
 256/(462\cdot30)&u_i\in\mathcal E,\\
 125/(462\cdot12)&u_i\in\mathcal V.
 \end{cases}
\]
On $L^2(S^2,d\Omega/(4\pi))$ define
$\mathscr Y_5=\bigoplus_{\ell=0}^5\mathscr H_\ell$, of dimension
$\sum_{\ell=0}^5(2\ell+1)=36$.  The finite sampling operator
\[
 (Sf)_i=\sqrt{w_i}\,f(u_i),\qquad S:\mathscr Y_5\to\mathbb C^{62},
\]
is an isometry, $S^*S=I_{36}$, so its $62\times62$ Gram $SS^*$
is an orthogonal rank-$36$ projector with kernel dimension $26$.
Its reproducing formula is
\begin{equation}\label{p1b:eq:cc13-frame-reconstruct}
 \boxed{f(n)=\sum_{i=1}^{62}w_if(u_i)
 \sum_{\ell=0}^5(2\ell+1)P_\ell(n\cdot u_i),\qquad f\in\mathscr Y_5.}
\end{equation}
Furthermore, for any $f\in\mathscr Y_6$, every coefficient
$\langle Y_{\ell m},f\rangle_{L^2(d\Omega)}$ with $\ell\le5$ is
computed exactly by the $62$-point rule:
\[
 \boxed{a_{\ell m}=4\pi\sum_i w_i f(u_i)
 \overline{Y_{\ell m}(u_i)},\quad \ell\le5,\qquad
 f\in\mathscr Y_6.}
\]
The latter statement is \emph{not} an isometry on $\mathscr Y_6$.
\end{theorem}
\begin{proof}
The product of a degree-$\ell$ and a degree-$k$ spherical harmonic
is the restriction of a polynomial on $\mathbb R^3$ of total degree
at most $\ell+k$.  For $\ell,k\le5$, the degree is at most ten,
and the source cubature is exact through eleven.  Therefore
$\sum_iw_i\overline{f(u_i)}g(u_i)
=\int\overline f g\,d\Omega/(4\pi)$, proving the isometry.
The normalized-sphere reproducing kernel on $\mathscr Y_5$ is
$K_5(n,u)=\sum_{\ell\le5}(2\ell+1)P_\ell(n\cdot u)$, by the
spherical-harmonic addition theorem.  Applying $S^*S=I$ proves the
reconstruction formula.  If $f\in\mathscr Y_6$, multiplication with
$\overline{Y_{\ell m}}$ for $\ell\le5$ has degree at most eleven,
proving the coefficient identity and exact cancellation of the
$\ell=6$ contribution.
The nonzero degree-twelve cubature defect
$\sum_iw_iP_{12}(n\cdot u_i)=7429/33000$ for an icosahedral vertex $n$
is given in the source cubature proof.  Degree-twelve harmonics occur
in products of two degree-six harmonics, so the norm-isometry theorem
cannot be extended to the complete $\mathscr Y_6$.
\end{proof}

\begin{corollary}[Hopf-basic isometry with the contact host unchanged]
\label{p1b:cor:cc13-hopf-basic-frame}
On the normalized Hopf contact measure
$d\mu_H=(4\pi^2)^{-1}\eta_0\wedge d\eta_0$, the pullback
$f\mapsto f\circ\pi_H$ is an isometry from normalized spherical
$L^2(S^2)$ into the Hopf-basic scalar functions on $S^3$.  For
$f,g\in\mathscr Y_5$,
\[
 \boxed{\int_{S^3}\overline{(f\circ\pi_H)}(g\circ\pi_H)
 \,d\mu_H=\sum_{i=1}^{62}w_i\overline{f(u_i)}g(u_i).}
\]
This is a finite exact evaluation of the chosen basic sector, not
an alteration of the total contact measure $4\pi^2$ or a sampling
formula for charged/non-basic Hopf modes.
\end{corollary}
\begin{proof}
Apply the Hopf-pushforward identity of
Corollary~\ref{p1b:cor:new-hopf-pullback-11-cubature} to the
polynomial $\overline f g$ of degree at most ten and divide by $4\pi^2$.
\end{proof}

\section{Alternating exceptional finite averages and continuous Haar reconstruction}
\label{p1b:sec:cc15-exceptional-symmetrizer}
\begin{theorem}[Bandwise Haar limit from two finite exceptional symmetrizers]
\label{p1b:thm:cc15-alternating-haar}
Fix the \emph{aligned} binary-octahedral and binary-icosahedral subgroups
$O=2O$, $I=2I\subset SU(2)$ containing the same $H=2T$ of
Theorem~\ref{p1:thm:2O-2I-generate-SU2}.  Let $\rho$ be a continuous
unitary representation on a finite-dimensional space $V$ and set
\[
 P_O=\frac1{48}\sum_{g\in O}\rho(g),\quad
 P_I=\frac1{120}\sum_{h\in I}\rho(h),\quad
 P_{SU2}=\int_{SU(2)}\rho(g)\,dg.
\]
Then $P_O,P_I$ are orthogonal projections,
$\operatorname{Ran}P_O\cap\operatorname{Ran}P_I
=\operatorname{Ran}P_{SU2}$, and
\[\boxed{(P_OP_I)^m\xrightarrow[m\to\infty]{\|\cdot\|}P_{SU2}.}\]
On the orthogonal complement of the $SU(2)$-fixed space the convergence
is geometric, with contraction determined by the nontrivial principal angles
between the two finite-group invariant spaces. The statement is for the
\emph{same} representation, and does not combine unrelated physical owners.
\end{theorem}
\begin{proof}
Finite-group averages of unitary representations are self-adjoint
idempotents, hence orthogonal projections.  A vector fixed by both
is fixed by the generated subgroup and, by continuity, by its dense closure
$SU(2)$. The converse is immediate. Orthogonal alternating projections
converge in norm onto their common range in finite dimension: after splitting
off the common range, no principal cosine is one, and each nonzero
principal-angle two-plane contracts geometrically.  This is an explicit
finite-dimensional result, with no infinite-completion claim uniform in
representation degree.
\end{proof}

\begin{theorem}[Exact first overlap and sharp low-band convergence rates]
\label{p1b:thm:cc15-exceptional-character-overlap}
On the irreducible $SU(2)$ representation $V_n=\operatorname{Sym}^n\mathbb C^2$
write $P_O^{(n)},P_I^{(n)}$ for the finite averages. The invariant
ranks are coefficients of the source's binary-polyhedral Molien series
\[
 M_O(t)=\frac{1+t^{18}}{(1-t^8)(1-t^{12})},\qquad
 M_I(t)=\frac{1+t^{30}}{(1-t^{12})(1-t^{20})}.
\]
In degrees $n=12,20$ both invariant spaces are lines. Their exact
squared principal cosines are
\[\boxed{c_{12}^2=\operatorname{Tr}(P_O^{(12)}P_I^{(12)})=\frac{11}{32},
\qquad c_{20}^2=\operatorname{Tr}(P_O^{(20)}P_I^{(20)})=\frac{95}{128}.}\]
At degree $24$, where the $2O$ invariant space has dimension two but
the $2I$ invariant space is a line,
\[\boxed{\operatorname{Tr}(P_O^{(24)}P_I^{(24)})=\frac{761}{2048}.}\]
For the full scalar harmonic band $\mathcal P_N(S^3)=\oplus_{n=0}^N\mathcal H_n(S^3)$,
with both groups acting by the same chosen right regular action, let
$P_0$ project onto constants. For every integer $m\ge1$ one has the
\emph{sharp} operator norms
\[
 \boxed{\left\|(P_OP_I)^m-P_0\right\|_{\mathcal P_{12}}
   =\left(\frac{11}{32}\right)^{m-1/2},\qquad
 \left\|(P_OP_I)^m-P_0\right\|_{\mathcal P_{20}}
   =\left(\frac{95}{128}\right)^{m-1/2}.}
\]
Thus the degree-twelve $2I$ harmonic alias is not fixed by the $2O$
projection: repeated alternating exceptional averages suppress it at
the first displayed exact rate.
\end{theorem}
\begin{proof}
The irreducible character is $\chi_n(q)=U_n(\operatorname{Re}q)$,
where $U_n$ is the second-kind Chebyshev polynomial. The finite sum
\[
 \operatorname{Tr}(P_O^{(n)}P_I^{(n)})=
 \frac1{48\cdot120}\sum_{g\in O}\sum_{h\in I}
 U_n\bigl(\operatorname{Re}(g^{-1}h)\bigr)
\]
may be evaluated using the source's exact quaternionic representatives.
The two sets share precisely the 24 Hurwitz elements; the $5,760$
real-part pair values fall into 23 distinct algebraic values
(in $\mathbb Q(\sqrt2,\sqrt5)$).  Substitution of $U_{12},U_{20},U_{24}$
gives exactly $11/32,95/128,761/2048$, respectively; the accompanying
reproducibility script evaluates these sums in exact radical arithmetic.
The displayed Hilbert series show that below degree 21 the common
nonconstant degrees at which \emph{both} finite projectors are nonzero
are exactly $12,20$, and they have rank one at both degrees. On two
rank-one projection lines with squared overlap $c^2$, the $m$th
product has norm $c^{2m-1}=(c^2)^{m-1/2}$.
Peter--Weyl splits $\mathcal H_n(S^3)$ into $(n+1)$ copies of this
right representation, so the same singular values occur on the scalar band.
Taking the maximum of the nonconstant blocks gives the sharp norms.
\end{proof}
\begin{remark}[No physical smoothing or infinitesimal spectral gap inferred]
The two finite operations are comparison averages, not an action-generated
relaxation law. Convergence is in every \emph{fixed finite} representation
band. Without a uniform-in-degree angle bound it is not asserted to be
operator-norm convergence on the whole $L^2(S^3)$. The harmonic alias
relation does not select an actual cosmological primordial state.
\end{remark}

\subsection{Sharp exceptional heat-smoothed scalar cubature bounds}
\label{p1b:sec:cc16-heat-cubature}
\begin{theorem}[Exact norm of the heat-smoothed exceptional averaging defect]
\label{p1b:thm:cc16-heat-cubature-error}
Let $\Gamma\in\{2T,2O,2I\}$ be one of the fixed binary-polyhedral
subgroups of $SU(2)=S^3(R)$, with its uniform group measure.
Let $\Delta_+\geq0$ be the round scalar Laplacian, and define, for $t>0$,
\[
 \mathcal E_{\Gamma,t}f:=\frac1{|\Gamma|}\sum_{g\in\Gamma}
       (e^{-t\Delta_+}f)(g)-\int_{S^3(R)}f\,d\mu_H.
\]
Write $m_\Gamma(n)=\dim\operatorname{Sym}^{n}(\mathbb C^2)^\Gamma$,
whose generating functions are exactly
\eqref{p1:eq:molien-2T}--\eqref{p1:eq:molien-2I}.
As a continuous functional on $L^2(S^3(R),d\mu_H)$,
\begin{equation}\label{p1b:eq:cc16-heat-cubature-norm}
 \boxed{\|\mathcal E_{\Gamma,t}\|_{L^2\to\mathbb C}^{\,2}
 =\sum_{n=1}^\infty (n+1)m_\Gamma(n)
      e^{-2tn(n+2)/R^2}.}
\end{equation}
The first nonzero terms occur at $n_\Gamma=6,8,12$ for
$2T,2O,2I$, respectively, each with multiplicity one. Therefore
\[
 \|\mathcal E_{2T,t}\|\sim\sqrt7\,e^{-48t/R^2},\quad
 \|\mathcal E_{2O,t}\|\sim3e^{-80t/R^2},\quad
 \|\mathcal E_{2I,t}\|\sim\sqrt{13}\,e^{-168t/R^2}
 \qquad(t/R^2\longrightarrow\infty).
\]
All equalities concern only Haar-normalized scalar functions on the
fixed round comparison shell. In particular no interacting QCD,
spin Dirac, metric-transport, or physical observer error bound follows
without a separate same-operator identification.
\end{theorem}
\begin{proof}
Peter--Weyl yields the Haar-orthonormal basis
$\sqrt{n+1}\,D^{(n)}_{ab}(g)$, $1\leq a,b\leq n+1$, of the
scalar degree-$n$ harmonic sector with eigenvalue
$n(n+2)/R^2$. Averaging its matrix coefficients over $\Gamma$
gives $\sqrt{n+1}\,(P_\Gamma^{(n)})_{ab}$, where
$P_\Gamma^{(n)}=|\Gamma|^{-1}\sum_{g\in\Gamma}D^{(n)}(g)$
is an orthogonal projector of rank $m_\Gamma(n)$.
Thus the squared norm of the group-averaging functional on the
$n$th harmonic sector is
$(n+1)\operatorname{tr}(P_\Gamma^{(n)})
=(n+1)m_\Gamma(n)$. Haar averaging annihilates every nonconstant
sector, so applying the heat multiplier and summing the mutually
orthogonal sectors proves \eqref{p1b:eq:cc16-heat-cubature-norm}.
The series converges for each $t>0$ because
$0\le m_\Gamma(n)\le n+1$.
Reading the first positive nonconstant coefficients from the three
already established Molien series proves the leading terms. The
higher terms are exponentially smaller as $t/R^2\to\infty$.
\end{proof}
\begin{corollary}[Certified scalar source integration]
\label{p1b:cor:cc16-heat-source-bound}
For every $f\in L^2(S^3(R))$, the error in heat-smoothed
$\Gamma$-point integration satisfies the sharp estimate
\[
 \left|\frac1{|\Gamma|}\sum_{g\in\Gamma}(e^{-t\Delta_+}f)(g)
 -\int f\,d\mu_H\right|
 \leq\left(\sum_{n\ge1}(n+1)m_\Gamma(n)e^{-2tn(n+2)/R^2}\right)^{1/2}
 \|f\|_2.
\]
The norm constant cannot be decreased for the specified operator
and time. This may be consumed by a downstream finite-window scalar
source or heat comparison if its operator, measure and radius match.
\end{corollary}


\subsection{A positive vertex--facet nineteen-design from the two exceptional siblings}
\label{p1b:sec:cc17-vertex-facet-positive-rule}

\begin{theorem}[The 600 chamber centres are one explicit double coset]
\label{p1b:thm:cc17-facet-double-coset}
Fix the same marked Hurwitz subgroup $H=2T\subset I=2I\subset SU(2)$
used in the 600-cell chamber construction, and fix
$u=(1+\mathbf i)/\sqrt2\in O=2O\setminus H$.
Then $uHu^{-1}=H$, the twin $I'=uIu^{-1}$ has
$I\cap I'=H$, and
\[
\boxed{\mathscr C=IuI,\qquad|\mathscr C|=\frac{120^2}{24}=600.}
\]
The elements of $\mathscr C$ are precisely the \emph{normalized centres}
of the $600$ tetrahedral facets of the 600-cell $I\subset S^3$.
Moreover $\mathscr C$ splits into five disjoint ordinary $I$ cosets,
and the existing 25 chamber blocks are bi-translated $24$-element
Hurwitz orbits $aHb$ (with the relevant $u$ inserted).
\end{theorem}
\begin{proof}
The twin-completion theorem and its explicit shared Hurwitz normalizer
give $I\cap uIu^{-1}=H$. The orbit--stabilizer formula for the
$I\times I$ left/right action therefore gives $|IuI|=120^2/24$.
To identify the geometric orbit, the four 600-cell vertices
\[
 (1/2,\varphi/2,0,\pm 1/(2\varphi)),\qquad
 (\varphi/2,1/2,\pm1/(2\varphi),0)
\]
are pairwise at scalar product $\varphi/2$, hence are a tetrahedral
facet; their normalized sum is $u$. Every left/right quaternionic
multiplication by elements of $I$ is an isometry preserving the
vertex set, so all points of $IuI$ are facet centres. The source's
complete facet census contains exactly 600, giving equality.
The five-coset and 25-block statements are the respective indices
$[I:H]=5$ and $[I\times I:H_{\mathrm{diag}}]=600$, grouped by
left/right $I/H$ labels. Since every block is a bi-translate of $H$,
it is a rotated Hurwitz spherical five-design.
\end{proof}

\begin{theorem}[Explicit \emph{positive} degree-nineteen vertex--facet cubature]
\label{p1b:thm:cc17-positive-vertex-facet-19}
Let $V=I$ (120 vertices), $F=IuI$ (600 tetrahedral facet centres),
with probability-normalized Haar measure $d\sigma$ on the round unit $S^3$.
For every restricted scalar polynomial $f\in\mathcal P_{19}(S^3)$,
\[
\boxed{\int_{S^3} f\,d\sigma
=\frac5{21}\frac1{120}\sum_{v\in V}f(v)
 +\frac{16}{21}\frac1{600}\sum_{c\in F}f(c).}
\]
The rule is positive, invariant under the declared marked $I$ left/right
operations, and supported on six $I$ cosets. Its weights per coset
are $25/105$ on $I$ and $16/105$ on each of the other five cosets;
it has exactly 720 distinct nodes. The strength is sharp: on the
nonconstant character $U_{20}(\operatorname{Re}q)$ its average is
\[\boxed{\frac{17}{28}\ne0.}\]
\end{theorem}
\begin{proof}
Write $\mu_I$ for uniform counting on $I$ and
$\mu_F=\mu_I*\delta_u*\mu_I$ for uniform counting on the 600-element
double coset. For any $SU(2)$ irrep
$V_n=\operatorname{Sym}^n\mathbb C^2$, their Fourier transforms are
$P_I^{(n)}$ and $P_I^{(n)}\rho_n(u)P_I^{(n)}$.
The source Molien series has no $I$-fixed vector for $1\le n\le19$
except one line at $n=12$. On that line the independently source-proved
Hurwitz-twin overlap is exactly $-5/16$, so the positive convex combination
$\frac5{21}P_I^{(12)}+\frac{16}{21}P_I^{(12)}\rho_{12}(u)P_I^{(12)}$
vanishes. Therefore all nonconstant Fourier coefficients through
19 vanish, proving the Haar identity by Peter--Weyl.
At degree 20 the $I$-fixed space is again a line and the exact
Chebyshev character sum
\[
\frac1{120}\sum_{h\in I}U_{20}(\operatorname{Re}(hu))=\frac{31}{64}
\]
gives the residual $5/21+(16/21)(31/64)=17/28$.
All node weights are strictly positive; the reduction of the
600-point orbit to five $I$ cosets uses $[I:H]=5$.
\end{proof}

\begin{corollary}[Exact heat-smoothed scalar quadrature error]
\label{p1b:cor:cc17-720-heat-error}
Let $\nu$ be the displayed positive 720-point probability and
$\Delta_+$ the round scalar positive Laplacian on $S^3(R)$.
For $t>0$ the exact squared $L^2\to\mathbb C$ error norm is
\[
 \boxed{\|\nu\circ e^{-t\Delta_+}-\sigma\circ e^{-t\Delta_+}\|^2
 =\sum_{n=20}^\infty(n+1)e^{-2tn(n+2)/R^2}
       \|\widehat\nu(n)\|_{HS}^2,}
\]
where $\widehat\nu(n)=\frac5{21}P_I^{(n)}+
 \frac{16}{21}P_I^{(n)}\rho_n(u)P_I^{(n)}$.
Its first nonzero term has exact amplitude
\[
\boxed{\|\nu\circ e^{-t\Delta_+}-\sigma\circ e^{-t\Delta_+}\|
 \sim\frac{17\sqrt{21}}{28}e^{-440t/R^2}
 \quad(t/R^2\to\infty).}
\]
This sharpens the uncorrected $2I$ scalar leading term
$\sqrt{13}e^{-168t/R^2}$ on the same round scalar comparison.
\end{corollary}
\begin{proof}
Peter--Weyl decomposes the scalar $n$-shell into $n+1$ copies of
$V_n$ and the heat operator multiplies it by
$e^{-tn(n+2)/R^2}$. The dual norm squared of averaging on the
shell is $(n+1)\|\widehat\nu(n)\|_{HS}^2$.
All nonconstant terms below 20 vanish by the theorem, and at
$n=20$ the Fourier matrix is rank one with eigenvalue $17/28$.
The remaining series is exponentially smaller as $t/R^2\to\infty$.
\end{proof}

\begin{corollary}[25 simultaneous low-band null measurements]
\label{p1b:cor:cc17-register-lowband-blindness}
For each existing chamber block $B_{ab}$ of 24 facet centres, any
$f\in\mathcal P_5(S^3)$ obeys
\[
\boxed{\frac1{24}\sum_{c\in B_{ab}} f(c)=\int f\,d\sigma
\quad\text{for every }(a,b)\in\{1,\ldots,5\}^2.}
\]
Thus the centered 24-dimensional quotient of the 25-register readout
annihilates every nonconstant harmonic of degrees 1 through 5.
In particular the existing chamber-memory register contrasts do
\emph{not} measure these low scalar harmonics.
\end{corollary}
\begin{proof}
Each $B_{ab}$ is an isometric bi-translate of the Hurwitz
24-cell $H$, which is a spherical five-design by the source.
The degree-five polynomial mean is therefore its Haar mean on every
block separately. Subtracting any two block means gives zero.
\end{proof}

\begin{remark}[Strongly limited physical scope]
The 720 nodes are an exact mathematical cubature, not a new finite
physical carrier, matter interaction, or extra geometric axiom.
The rule is \emph{marked} by an aligned $H\subset I,O$, and cannot be
inserted as an unmarked GP3 selector. The separately source-proved
positive degree-nineteen rotated-icosahedral rule remains valid and
is not superseded; the advance here is a rational, explicit,
vertex--existing-chamber-centre construction with 720 points.
\end{remark}


\section{Hopf-projected chamber cubature and exact angular response}
\label{p1b:sec:cc18-hopf-chamber-response}
All statements of this section are about the marked mathematical
$2I$--$2T$ vertex/facet construction and its existing finite comparison
operators.  In particular, they do not select a sky instrument,
physical time parameter, meson channel, or new GP axiom.

\begin{theorem}[Positive 150-direction degree-nine Hopf rule]
\label{p1b:thm:cc18-hopf-150-rule}
Let $I=2I$, $H=2T$, $u=(1+\mathbf i)/\sqrt2$ and $F=IuI$ be the
120 vertices and 600 tetrahedral facet centres of
Theorem~\ref{p1b:thm:cc17-facet-double-coset}.  Mark the Hopf map
$\pi(q)=q\mathbf i q^{-1}:S^3\to S^2$.
The measure of Theorem~\ref{p1b:thm:cc17-positive-vertex-facet-19}
pushes forward to a \emph{positive 150-direction} cubature, exact on
all real spherical polynomials through degree nine.  There are 30
shared vertex/facet directions with weight $41/3150$ each, and 120
facet-only directions with weight $8/1575$ each.  With probability-area
orthonormal harmonics $Y_{\ell m}$, the matrix
$S_{r,\ell m}=\sqrt{w_r}Y_{\ell m}(z_r)$ on
$\mathcal Y_4=\bigoplus_{\ell=0}^4\mathcal H_\ell(S^2)$ obeys
\[
\boxed{\dim\mathcal Y_4=25,\qquad S^*S=I_{25},\qquad
S_{\le4}^*S_{\ell=5}=0.}
\]
The isometry concerns \emph{complex amplitudes}, not quantum-state
tomography or an operational cosmological sky.
\end{theorem}
\begin{proof}
Each component of $\pi(q)$ is quadratic in the real quaternion
coordinates.  Pullback of a degree-nine spherical polynomial therefore
has degree at most eighteen on $S^3$, inside the source's degree-nineteen
exact cubature.  Haar probability measure pushes forward to normalized
spherical area.  Right multiplication by $1,\mathbf i,-1,-\mathbf i$
fixes $\pi$, and the exact $2I$ and $IuI$ orbit/coset enumeration gives
30 and 150 fibres respectively, of size four each.  Since
$u\mathbf i u^{-1}=\mathbf i$, $\pi(g)=\pi(gu)$ for each $g\in I$;
thus all 30 vertex directions also occur in $F$.  The two per-direction
weights are
$4/504+4(2/1575)=41/3150$ and $4(2/1575)=8/1575$.
They are positive and sum to $30(41/3150)+120(8/1575)=1$.
The product of two harmonics of degrees at most four is a polynomial
of degree at most eight; the product of a degree-five and degree-at-most-four
harmonic has degree at most nine.  Cubature exactness and orthogonality
give both matrix identities.  The explicit exact fibre and polynomial
checks are retained in the linked CC18 verification ledger.
\end{proof}

\begin{theorem}[Exact 25-block angular selection and singular values]
\label{p1b:thm:cc18-angular-selection}
Let $B_{ab}$, $a,b\in\{0,\ldots,4\}$, be the ordered 24-centre
blocks of the existing marked 600-chamber decomposition, and put
\[
R_\ell:\mathcal H_\ell(S^2)\to\mathbb C^{25},\qquad
(R_\ell f)_{ab}=\frac1{24}\sum_{q\in B_{ab}} f(\pi(q)).
\]
With probability-normalized input harmonics the exact ranks for
$\ell=0,1,\ldots,6$ are
\[
\boxed{\bigl(\operatorname{rank}R_0,\ldots,
\operatorname{rank}R_6\bigr)=(1,0,0,4,4,0,5).}
\]
The nonzero squared singular values of both $R_3$ and $R_4$ are
$525/64$ with multiplicity four; moreover $R_3^*R_4=0$.
In particular the total image of the $25$-dimensional angular band
$\mathcal Y_4$ in the $25$ chamber registers has rank exactly nine.
\end{theorem}
\begin{proof}
The marked left $A_5=2I/\{\pm1\}$ action on the five row cosets
carries $\mathbf1\oplus\mathbf4$, with right label multiplicity five.
The input harmonic decompositions for $\ell=1,\ldots,6$ are
$\mathbf3$, $\mathbf5$, $\mathbf3'\oplus\mathbf4$,
$\mathbf4\oplus\mathbf5$,
$\mathbf3\oplus\mathbf3'\oplus\mathbf5$, and
$\mathbf1\oplus\mathbf3\oplus\mathbf4\oplus\mathbf5$.
Equivariance therefore kills all pieces other than $\mathbf1$ and
$\mathbf4$.  The 24-point five-design additionally shows that
$R_1=R_2=0$ directly.  To certify the allowed components as nonzero,
let $P_\ell$ denote the Legendre polynomial and define the ordered
pair sums
\[
C_{\ell b}=\sum_{q,r\in B_{0b}}P_\ell\bigl(\pi(q)\cdot\pi(r)\bigr).
\]
Exact arithmetic in $\mathbb Q(\sqrt5)$ on the source-defined blocks gives
\[
\begin{array}{c|rr}
\ell&C_{\ell0}& C_{\ell b}\ (b\ne0)\\\hline
3&0&135\\
4&336&21\\
6&72&1773/64
\end{array}
\]
The harmonic addition theorem yields
\[
\|R_\ell\|_{\mathrm{HS}}^2
=\frac{5(2\ell+1)}{24^2}(C_{\ell0}+4C_{\ell1}),
\]
which is $525/16$ for each of $\ell=3,4$ and $21125/1024$
for $\ell=6$.  Schur's lemma makes the four nonzero squared singular
values equal, establishing $525/64$ for degrees three and four.
The right Hurwitz element $\mathbf j$ permutes columns and reverses the
Hopf axis, $\pi(q\mathbf j)=-\pi(q)$.  It consequently acts with
opposite parity on the degree-three and degree-four output subspaces,
which are orthogonal.  The degree-six invariant and $\mathbf4$
components are shown nonzero by the exact strengths in
Theorem~\ref{p1b:thm:cc18-degree-six-sector-split} below.
\end{proof}

\begin{corollary}[Angular sectors resolve the existing hidden chamber memory]
\label{p1b:cor:cc18-angular-memory}
For $Z=J_5/5$, $Y=I_5-Z$, and the existing orthogonal register
projectors $P_0=Z\otimes Z$, $P_a=Z\otimes Y+Y\otimes Z$,
$P_c=Y\otimes Y$, let $M(t)$ be the exact return
operator of Theorem~\ref{p1b:thm:cc13-register-memory}.  Then
\[
\boxed{M(t)R_3=f_c(t)R_3,\qquad M(t)R_4=f_c(t)R_4,}
\]
where the established hidden-memory kernel is
\[
f_c(t)=\left(\frac12+\frac{3}{4\sqrt{21}}\right)
 e^{(-9+\sqrt{21})t/8}
+\left(\frac12-\frac{3}{4\sqrt{21}}\right)
 e^{(-9-\sqrt{21})t/8}.
\]
The coarse Markov closure $e^{-15t/16}$ differs by
\[
\boxed{f_c(t)-e^{-15t/16}=\frac{75}{512}t^2+O(t^3).}
\]
\end{corollary}
\begin{proof}
The degree-three and degree-four outputs have vanishing row and
column marginals: the former by the absence of the corresponding
icosahedral invariant, and the latter also by the source's low-orbit
design.  Hence $R_3=P_cR_3$, $R_4=P_cR_4$.
Substitute into the source's identity
$M(t)=P_0+e^{-5t/4}P_a+f_c(t)P_c$.
The exact hidden centered generator has off-diagonal entries
$5\sqrt3/16$; its squared coupling $75/256$ is the difference in
second derivatives of the true and coarse return at zero.  Taylor's
theorem supplies the stated $75/512$ coefficient.
\end{proof}

\begin{theorem}[Degree-six angular invariant and complete register sector split]
\label{p1b:thm:cc18-degree-six-sector-split}
Use the ten marked projective axes of
Theorem~\ref{p1:thm:v285-ten-axis-quadrupole-exact-sequence} and put
$T_6(n)=\tfrac1{10}\sum_{a=1}^{10}P_6(n\cdot u_a)$.
Define the vertex and facet zonal averages
\[
F_6(n)=\frac1{120}\sum_{g\in I}P_6(n\cdot\pi(g)),\qquad
G_6(n)=\frac1{600}\sum_{q\in F}P_6(n\cdot\pi(q)).
\]
These satisfy the exact functional identities
\[
\boxed{F_6=\frac9{16}T_6,\qquad G_6=-\frac{45}{256}T_6.}
\]
On the existing 25 chamber blocks,
\[
\boxed{R(T_6)_{ab}=\begin{cases}
 11/144,&b=0,\\ -451/9216,&b\ne0,
\end{cases}}
\]
and equivalently
\[
R(T_6)=-\frac{55}{2304}\mathbf1_{25}
+\frac{77}{3072}\mathbf1_5\otimes(4,-1,-1,-1,-1).
\]
For $R_6$ the complete nonzero squared singular-value spectrum is
\[
\boxed{\operatorname{spec}_{>0}(R_6^*R_6)
=\left\{\frac{1934075}{262144}\;(1),\ 
\frac{3473925}{1048576}\;(4)\right\}.}
\]
The squared register-sector strengths on the input $\mathbf1$ and
$\mathbf4$ irreducibles respectively are
\[
\begin{array}{c|cc}
 &\mathbf1&\mathbf4\ \text{(per input direction)}\\\hline
P_0&89375/65536&0\\
P_a&1576575/262144&825825/262144\\
P_c&0&170625/1048576
\end{array}
\]
In particular the invariant $T_6$ has only stationary plus
$e^{-5t/4}$ register return, whereas a generic visible degree-six
input also has a nonzero $f_c(t)$ memory component.
\end{theorem}
\begin{proof}
The first icosahedral invariant harmonic on $S^2$ is unique in degree
six.  Hence $F_6$ and $G_6$ are multiples of $T_6$.  The exact marked
axis values are
$F_6(\mathbf i)=11/256$, $G_6(\mathbf i)=-55/4096$,
$T_6(\mathbf i)=11/144$, proving both identities.  The exact
block averages over the source's $25$ bi-Hurwitz cosets are
$11/144$ for the distinguished right column and $-451/9216$
elsewhere, which algebraically yields the $P_0\oplus P_a$ vector
above; these values are independently re-enumerated in the CC18
exact-arithmetic certificate.  The source gives
$\|T_6\|_{L^2(S^2)}^2=11/1053$.
Projecting $R(T_6)$ onto $P_0$ and $P_a$ and dividing its squared
lengths by this norm gives
$89375/65536$ and $1576575/262144$.
For the nontrivial $\mathbf4$ input, the addition theorem applied to
the five 120-point row-orbit means and the global facet mean gives
\[
\|(Y\otimes Z)R_6\|_{\mathrm{HS}}^2
=25\cdot13^2\left[\left(\frac9{16}\right)^2
-\left(\frac{45}{256}\right)^2\right]\frac{11}{1053}
=\frac{825825}{65536}.
\]
Schur's lemma divides this four-dimensional trace by four.
Theorem~\ref{p1b:thm:cc18-angular-selection} gives
$\|R_6\|_{\mathrm{HS}}^2=21125/1024$.  Subtract the $P_0$
trace, the invariant $P_a$ trace and the four-dimensional $P_a$
trace to obtain a four-dimensional $P_c$ trace
$4\cdot170625/1048576$.  The resulting positive strengths show
that the allowed $\mathbf1\oplus\mathbf4$ occurs with full rank
five.  Summing sector strengths gives the stated squared singular
values; $\dim\ker R_6=8$.  Since the source's $M(t)$ acts by
$1,e^{-5t/4},f_c(t)$ on the three projectors, the final return
statement is immediate.
\end{proof}

\begin{remark}[Scope and comparison boundaries]
\label{p1b:rem:cc18-physical-gates}
The degree-nine sky rule is a marked finite mathematical integration
identity.  The register dynamics uses its own dimensionless chamber
walk time and is not inferred to be the observed sky, a particle decay
generator, or a physical thermal clock.  An action-owned source,
transport/window, and detector mapping would be required to promote
any of these comparisons to an observable prediction.  The original
cubature strength, physical owner statements, and open audit gates
remain unchanged.
\end{remark}



\section{Strong core lumpability and two-sided finite Hopf incidence}
\label{p1b:sec:cc20-core-lumping-hopf}
\begin{theorem}[Two exact Markov core quotients of the chamber walk]
\label{p1b:thm:cc20-two-core-lumping}
Let $A$ be the existing four-regular adjacency on the $600$ tetrahedral
chambers, and let $U_L,U_R$ pull back five-core functions along the two
registered core labels.  Then
\[
 AU_L=U_L(J_5-I_5),\qquad AU_R=U_R(J_5-I_5).
\]
For $Q=A/4-I_{600}$ the two core observations are therefore strongly
lumpable for \emph{every} initial chamber law, each with generator
$Q_5=(J_5-5I_5)/4$ and transition matrix
\[
 P_5(t)=J_5/5+e^{-5t/4}(I_5-J_5/5).
\]
For the joint $25$-register return, however, the exact source operator is
$M(t)=P_0+e^{-5t/4}P_a+f_c(t)P_c$, where
$\operatorname{rank}(P_0,P_a,P_c)=(1,8,16)$.  Consequently, two strongly
Markovian core marginals coexist with a non-Markovian joint register return.
\end{theorem}
\begin{proof}
The source chamber-resolvent theorem gives eigenvalue $4$ on constant
chamber functions and eigenvalue $-1$ on each additive centred core
function $a_i+b_j$.  Any five-core function decomposes into these
parts, giving both intertwining identities.  Exponentiation proves the
strong lumpability and the stated five-state semigroup.  The established
joint-return decomposition has a nontrivial $P_c$ memory kernel and
fails the register-only Chapman--Kolmogorov identity; neither row nor
column marginal detects $P_c$.
\end{proof}
\begin{theorem}[Dual ten-point Hopf-fibre incidence]
\label{p1b:thm:cc20-dual-hopf-incidence}
Mark an order-ten cyclic subgroup of the existing binary-icosahedral
quaternions with the same left and right action conventions as the
source $25$ Hurwitz-core registers.  The $600$ facet centres split into
$60$ left and $60$ right ten-point orbits.  Each orbit meets five
registers of one fixed opposite core label in exactly two points each.
If $B_L,B_R\in\{0,1\}^{25\times60}$ record the five incident register
labels per fibre, then
\[
 B_LB_L^{\mathsf T}=12J_5\otimes I_5,\qquad
 B_RB_R^{\mathsf T}=12I_5\otimes J_5.
\]
For $C_L=B_L^{\mathsf T}/5$, $C_R=B_R^{\mathsf T}/5$ and
$Z=J_5/5$, $Y=I_5-Z$,
\[
 C_L^{\mathsf T}C_L+C_R^{\mathsf T}C_R
 =\frac{24}{5}(Z\otimes Z)
 +\frac{12}{5}(Z\otimes Y+Y\otimes Z).
\]
Its rank is nine and its kernel is precisely
$P_c=Y\otimes Y$, of dimension sixteen; the condition number on its
nonzero image is $\sqrt2$.  The statement concerns marked
phase-blind finite Hopf averages, not a constituted physical apparatus.
\end{theorem}
\begin{proof}
Every marked left ten-point fibre fixes the opposite right core label,
and conversely.  The exact $IuI$ double-coset incidence gives twelve
fibres per register and five registers per fibre, with two facets at
each incidence. Thus each row in a fixed-core block of $B$ has twelve
ones, each of its four partners has the same twelve ones, and rows
outside that block have disjoint support.  These counts give the
stated two Gram matrices.  Since $J_5=5Z$ and $I_5=Z+Y$, adding and
normalizing gives the projector expression and its spectrum.
\end{proof}
\begin{remark}[Source domain]
The chamber time in the first theorem is the dimensionless random-walk
parameter.  The marked finite Hopf circles in the second theorem are
not GP-selected physical preferred directions.  No equation of
physical particle rates or detector channels follows from these results.
\end{remark}


\subsection{Two-mode obstruction to a common diffusion clock}
\begin{proposition}[Finite-exceptional symmetrizer is not a common heat step]
\label{cc19:prop:two-mode-no-heat-clock}
Let $P_I,P_O$ denote right finite-subgroup averaging operators on the
same round scalar $SU(2)$ carrier, and set $T=P_IP_OP_I$.
On the one-dimensional right-$I$-fixed subspaces of scalar irreducible
harmonic degrees twelve and twenty, $T$ acts by $11/32$ and $95/128$,
respectively. There is no common $t>0$ such that its restriction to
both is equal to the scalar heat operator $e^{-t\Delta_{S^3(R)}}$.
\end{proposition}
\begin{proof}
The two exact principal-cosine squares are established by
Theorem~\ref{p1b:thm:cc15-exceptional-character-overlap}.  The round
scalar Laplacian has eigenvalues $168/R^2$ and $440/R^2$ on these
respective shells. Equality would require simultaneously
$t/R^2=\log(32/11)/168$ and $t/R^2=\log(128/95)/440$,
which are unequal. Matching one line does not match the other.
\end{proof}



\subsection{A positive hierarchy of exceptional-group orbit cubatures (mathematical comparison)}
\label{p1b:sec:cc110-exceptional-orbit-ladder}

Let $I=2I\subset SU(2)=S^3$, $\varphi=(1+\sqrt5)/2$, and choose
$g\in I$ with $\operatorname{Re}g=\varphi/2$. Define the six point orbits
\[
\begin{gathered}
V=I,\qquad F=I\frac{1+\mathbf i}{\sqrt2}I,\qquad
E=I\frac{1+g}{\|1+g\|}I,\\
Q_5=I\frac{2+g}{\|2+g\|}I,\qquad
Q_3=I\frac{7+\mathbf i+\mathbf j+\mathbf k}{2\sqrt{13}}I,
\qquad Q_2=I\frac{2\sqrt2+\mathbf i}{3}I.
\end{gathered}
\]
The six respective cardinalities are $120,600,720,1440,2400,3600$;
write $M_X(f)=|X|^{-1}\sum_{q\in X}f(q)$ for each orbit.
For even $n$, write
\[
C_n(q)=\frac1{120}\sum_{h\in I}U_n(\operatorname{Re}(h^{-1}q)),
\quad C_n(1)=1.
\]

\begin{theorem}[Positive 23-, 29-, 31- and 35-degree exceptional cubatures]
\label{p1b:thm:cc110-positive-orbit-cubatures}
For the successive orbit strings
$\mathcal X_3=(V,F,E)$,
$\mathcal X_4=(V,F,E,Q_5)$,
$\mathcal X_5=(V,F,E,Q_5,Q_3)$ and
$\mathcal X_6=(V,F,E,Q_5,Q_3,Q_2)$, define exact weights
$w^{(k)}$ as the unique solutions of
\[
\sum_{j=1}^k w_j^{(k)}=1,\qquad
\sum_{j=1}^k w_j^{(k)}C_n(x_j)=0
\quad(n\in D_k),
\]
where $x_j$ is the displayed representative of the $j$th orbit and
\[
D_3=\{12,20\},\quad D_4=\{12,20,24\},\quad
D_5=\{12,20,24,30\},\quad
D_6=\{12,20,24,30,32\}.
\]
Each moment matrix is nonsingular and all entries of every $w^{(k)}$
are strictly positive elements of $\mathbb Q(\sqrt5)$. Then
\[
\boxed{\int_{S^3} f\,d\sigma=\sum_{j=1}^k w_j^{(k)}M_{X_j}(f)
\quad\bigl(f\in\mathcal P_{d_k}(S^3)\bigr),\quad
(d_3,d_4,d_5,d_6)=(23,29,31,35).}
\]
The node counts are $(1440,2880,5280,8880)$. The first nonzero
harmonic averages occur in degrees $(24,30,32,36)$, so each strength is
sharp for the specified weights.
In particular, the degree-23 weights are exactly
\[
\boxed{w^{(3)}=\left(\frac5{36},\frac{128}{693},\frac{625}{924}\right)}.
\]
For the degree-29 weights, in the orbit order $(V,F,E,Q_5)$,
\[
\begin{split}
w_V^{(4)}&=-\frac{331}{4368}+\frac{4961\sqrt5}{109200},\\
w_F^{(4)}&=\frac{1024}{3003},\\
w_E^{(4)}&=\frac{14271875-3459375\sqrt5}{63599536},\\
w_{Q_5}^{(4)}&=\frac{3745405}{7338408}
+\frac{10687921\sqrt5}{1192491300}.
\end{split}
\]
The larger exact weights are specified without rounding by the displayed
finite moment systems and the accompanying $\mathbb Q(\sqrt5)$ certificate.
\end{theorem}

\begin{proof}
The source's binary-icosahedral Molien series is
$M_I(t)=(1-t^{60})/((1-t^{12})(1-t^{20})(1-t^{30}))$.
For $0<n\le35$, its only positive coefficients occur at
$n=12,20,24,30,32$, each with multiplicity one.
The round $n$th scalar harmonic shell is $V_n\otimes V_n^*$,
so its left- and right-$I$ invariant subspace has dimension $m_I(n)^2$.
Every double-orbit mean annihilates the complementary representation
components, and on each one-dimensional invariant line its value is
$C_n(x)$. The stated finite moment constraints therefore annihilate all
nonconstant scalar harmonics through $d_k$, giving exactness.

The facet cardinality is source-owned. For the edge-centre orbit, the
120 vertex neighbours at scalar product $\varphi/2$ give 720 edges;
the left/right $I$ action is transitive. Exact quaternionic intersection
checks give $|I\cap q_5Iq_5^{-1}|=10$ and
$|I\cap q_3Iq_3^{-1}|=6$, hence the 1440- and 2400-point orbits.
For $q_2=(2\sqrt2+\mathbf i)/3$, conjugation fixes the
$\mathbf i$ axis and rotates the $\mathbf j,\mathbf k$ components by
$\sin(2\theta)=4\sqrt2/9$. Since all coordinates of elements of $I$
lie in $\mathbb Q(\sqrt5)$ and $\sqrt2\notin\mathbb Q(\sqrt5)$,
only the subgroup $\{\pm1,\pm\mathbf i\}$ survives in the intersection,
of order four. Thus $|Iq_2I|=3600$.

For every representative with coordinates in $\mathbb Q(\sqrt5)$ up to
one common normalization, the even $U_n$ character sum reduces to exact
$\mathbb Q(\sqrt5)$ arithmetic through the recurrence
$U_{2j}(\sqrt y)=(4y-2)U_{2j-2}(\sqrt y)-U_{2j-4}(\sqrt y)$.
For $q_2$, cancellation of the odd mixed powers under conjugation by
$\mathbf j\in I$ reduces the moments to rational polynomials in
$\sin^2\theta=1/9$. Exact inversion of the displayed finite matrices
shows positive weights and gives the first surviving moments.
The independent finite certificates record each arithmetic step and
verify positivity with exact comparison in $\mathbb Q(\sqrt5)$.
\end{proof}

\begin{proposition}[The natural 600-cell skeleton cannot reach degree 24]
\label{p1b:prop:cc110-skeleton-24-obstruction}
Let $T$ be the 1200 normalized centres of the triangular faces of the
600-cell. On the union $V\cup F\cup E\cup T$ one has the exact
orbitwise identity
\[
\boxed{C_{24}=\frac{437}{840}+\frac{36179}{28560}C_{12}
-\frac{3211}{4080}C_{20}.}
\]
Therefore no signed or positive orbit-constant weights on these four
skeleton orbits can integrate every scalar harmonic through degree 24:
normalization and vanishing $C_{12},C_{20}$ would force the
$C_{24}$-average to equal $437/840\ne0$.
\end{proposition}

\begin{proof}
Direct exact $\mathbb Q(\sqrt5)$ character evaluation gives
\[
\begin{array}{c|ccc}
& C_{12}&C_{20}&C_{24}\\\hline
V&1&1&1\\
F&-5/16&31/64&-263/1024\\
E&-3/25&-211/625&1981/3125\\
T&-65/243&5207/19683&-1583/59049.
\end{array}
\]
Inserting each row yields the stated identity. Taking an orbit-weighted
mean proves the obstruction, regardless of the signs of the weights.
\end{proof}

\begin{remark}[Research scope and invariant-moment search]
These are marked, finite **mathematical** cubatures, not newly selected
physical momentum carriers or action coefficients. A general positive
orbit cubature on a band with $r$ independent nonconstant invariant
moment coordinates exists with at most $r+1$ double-orbit components:
Haar averaging places the origin in the convex hull of the compact set
of moment vectors, and Carath\'eodory's theorem selects at most $r+1$
points. For the degrees above, $r=2,3,4,5$. At larger degrees, fixed
spaces can have multiplicity greater than one; then full matrix-valued
Fourier constraints, not only scalar character moments, are required.
\end{remark}


\subsection{Positive double-icosahedral cubatures through degrees 39 and 41}
\label{p1b:sec:cc111-positive41}
\begin{theorem}[Explicit positive seven- and eight-orbit scalar cubatures]
\label{p1b:thm:cc111-positive41}
Let $I=2I\subset SU(2)$, let $g\in I$ have real part
$\varphi/2=(1+\sqrt5)/4$, and let $M_q$ be uniform averaging
on the distinct points of $IqI$. In addition to the six orbits
\[
q\in\left\{1,\frac{1+\mathbf i}{\sqrt2},
\frac{1+g}{\|1+g\|},\frac{2+g}{\|2+g\|},
\frac{7+\mathbf i+\mathbf j+\mathbf k}{2\sqrt{13}},
\frac{2\sqrt2+\mathbf i}{3}\right\},
\]
introduce $q'_5=(7/2+g)/\|7/2+g\|$ and
$q'_3=(6+\mathbf i+\mathbf j+\mathbf k)/\sqrt{39}$.
The new double orbits contain $1440$ and $2400$ points, respectively,
and all eight orbits are disjoint. For $n$ in the consecutive sets
$D_{39}=(12,20,24,30,32,36)$ and
$D_{41}=(12,20,24,30,32,36,40)$, put
\[
C_n(q)=\frac1{120}\sum_{h\in I}
U_n(\operatorname{Re}(h^{-1}q)),\qquad
A_{d}=\begin{pmatrix}1&\cdots&1\\
(C_n(q_X))_{n\in D_d,\,X\in\mathcal O_d}\end{pmatrix},
\quad w_d=A_d^{-1}e_0.
\]
Then each $A_d$ is invertible over $\mathbb Q(\sqrt5)$,
and every component of $w_d$ is strictly positive. Consequently
\[
\int_{S^3}f\,d\mu=\sum_{X\in\mathcal O_d}(w_d)_X M_X(f),
\qquad \deg f\le d,\qquad d\in\{39,41\}.
\]
The degree-39 rule has $10320$ nodes and the degree-41 rule has
$12720$ nodes. The first surviving averaged harmonics are $C_{40}$
and $C_{42}$, respectively, with exact nonzero values as recorded
in the attached \texttt{ExactCertificates-v1\_11.json}.
\end{theorem}
\begin{proof}
For the new representatives the exact intersections
$I\cap q'_5I(q'_5)^{-1}$ and $I\cap q'_3I(q'_3)^{-1}$ have sizes
$10$ and $6$. Orbit--stabilizer gives the displayed counts;
the exact character moment columns distinguish the eight orbits.
The existing degree-$\le41$ Molien series has one invariant line
exactly at each of $12,20,24,30,32,36,40$, and no other
nonconstant invariant lines. Double averaging annihilates every
noninvariant harmonic and turns the remaining ones into the
$C_n$ tests. The supplied exact $\mathbb Q(\sqrt5)$ certificate
replays Gaussian elimination, all zero moments, positive signs, and
nonzero first survivors. These equalities prove the two rules.
\end{proof}
\begin{remark}
This is a scalar round-$S^3$ comparison rule. No physical source,
spinorial identification or interacting spectral action is asserted.
\end{remark}


\subsection{Differential reconstruction of the icosahedral angular nodes}
\label{p1b:sec:cc111-klein-root}
\begin{theorem}[Klein generator differential closure and the 62 orbit nodes]
\label{p1b:thm:cc111-klein-root}
In standard Klein coordinates for a conjugate realization of the
fixed binary-icosahedral subgroup, let
\[
F_{12}(z,w)=zw(z^{10}+11z^5w^5-w^{10}),\quad
H_{20}=-\frac1{121}(F_{zz}F_{ww}-F_{zw}^2),
\quad T_{30}=-\frac1{20}(F_zH_w-F_wH_z).
\]
The three nonzero binary forms are invariant of degrees $12,20,30$,
and obey the exact degree-$60$ relation
\[ \boxed{T_{30}^2=H_{20}^3+1728F_{12}^5.} \]
Their projective root sets contain exactly $12,20,30$ distinct points,
respectively. Under the Hopf identification $\mathbb CP^1\cong S^2$
these are the three icosahedral special orbits of vertices, face centres
and edge-midpoint directions, up to a global rotation. Therefore the
source-owned 62-point degree-eleven cubature can be written
\[
\int_{S^2}f\,d\sigma=
\frac{125}{462}\langle f\rangle_{Z(F_{12})}+
\frac{81}{462}\langle f\rangle_{Z(H_{20})}+
\frac{256}{462}\langle f\rangle_{Z(T_{30})},
\qquad\deg f\le11.
\]
\end{theorem}
\begin{proof}
For $SL_2$ changes of coordinates the Hessian determinant of an
invariant binary form and the Jacobian determinant of two such forms
remain invariant, because $\det g=1$. Direct symbolic expansion gives
$H_{20}=z^{20}-228z^{15}w^5+494z^{10}w^{10}
+228z^5w^{15}+w^{20}$ and
$T_{30}=z^{30}+522z^{25}w^5-10005z^{20}w^{10}
-10005z^{10}w^{20}-522z^5w^{25}+w^{30}$,
and verifies the stated syzygy. Their dehomogenized polynomials are
square-free and pairwise coprime. An $A_5$-invariant root set with these
cardinalities must be the corresponding exceptional projective orbit.
The final equality is exactly the already source-established 62-point
formula rewritten using those same three node sets.
\end{proof}
\begin{corollary}[Optional chamber-to-Klein Fourier reconstruction]
\label{p1b:cor:cc111-chamber-klein}
Suppose the independently proposed chamber harmonic identity
$\sum_{a=1}^{16}|f_{6,a}(q)|^2=16+12C_{12}(q)$ holds with the
stated normalized $C_{12}$ and its rank-one $I$-invariant representation
vector $v_{12}$. Then
\[
13\int_{SU(2)}C_{12}(q)\rho_{12}(q)^\dagger\,d\mu(q)
=|v_{12}\rangle\langle v_{12}|.
\]
Hence the chamber intensity determines the projective Klein form,
its Hessian and Jacobian, and the three icosahedral node sets up to
basis rotation and an irrelevant overall phase.
\end{corollary}


\subsection{The first matrix-valued binary-icosahedral moment obstruction}
\label{p1b:sec:cc111-degree60-matrix}
\begin{theorem}[Failure of scalar-only moment tests at degree 60]
\label{p1b:thm:cc111-degree60-matrix}
For $I=2I$ acting on $V_n=\operatorname{Sym}^n\mathbb C^2$, put
$m_I(n)=\dim V_n^I$. The source Molien series gives fourteen
nonconstant degrees below $60$, each of multiplicity one, and
\[\boxed{m_I(60)=2.}\]
Thus $(\mathcal H_{60}(S^3))^{I_L\times I_R}$ has real dimension four.
Choose a real orthonormal basis $u_1,u_2$ of $V_{60}^I$ and set
\[
A_{60}(\nu)_{ab}=
\int_{SU(2)}\langle u_a,\rho_{60}(q)u_b\rangle\,\nu(dq).
\]
For a bi-$I$-invariant probability measure, exactness at degree $60$
requires $A_{60}(\nu)=0$ entrywise. Requiring only
$\operatorname{tr}A_{60}(\nu)=0$ misses three independent real
moments, even within positive bi-invariant probability measures.
The vectors associated with $F_{12}^5$ and $H_{20}^3$ span $V_{60}^I$.
\end{theorem}
\begin{proof}
The $t^{60}$ coefficient in
$(1-t^{60})/[(1-t^{12})(1-t^{20})(1-t^{30})]$
is $3-1=2$, and no earlier positive coefficient exceeds one.
Peter--Weyl gives the four invariant matrix coefficients
$\langle u_a,\rho_{60}(q)u_b\rangle$.
They are independent by Schur orthogonality. The trace is only their
sum on the diagonal. Pick a real bi-invariant degree-$60$ harmonic
$\Psi$ orthogonal to the trace. For sufficiently small
$\epsilon\ne0$, $(1+\epsilon\Psi)d\mu$ is a strictly positive
bi-invariant probability measure, exact on every lower harmonic and
on the trace test, but not on $\Psi$.
The invariant degree-$60$ relation is
$T_{30}^2=H_{20}^3+1728F_{12}^5$; the last two powers are independent.
\end{proof}
\begin{corollary}[Positive existence with at most nineteen double orbits]
\label{p1b:cor:cc111-degree61-existence}
There exists a positive bi-$I$-invariant scalar cubature exact for all
restricted polynomials through degree $61$, supported on at most
$19$ double orbits. This is an existence statement; it does not produce
explicit orbit coordinates or algebraic weights.
\end{corollary}
\begin{proof}
There are fourteen independent real nonconstant invariant moments
below degree $60$ and four at degree $60$. The continuous moment map
$SU(2)\to\mathbb R^{18}$ has zero Haar barycentre. Carath\'eodory's
theorem provides a convex combination of at most $19$ image points
with zero moments. Replacing each point by its $I$ double orbit
maintains those moments and gives the asserted positive finite rule.
Odd degree $61$ has no invariant moment because $-1\in I$.
\end{proof}


\subsection{Matrix-valued exceptional heat-smoothed quadrature error}
\label{p1b:sec:cc111-matrix-heat}
\begin{theorem}[Exact $L^2$ heat-error norm for a bi-invariant rule]
\label{p1b:thm:cc111-matrix-heat}
Let $\nu$ be a finite bi-$I$-invariant probability measure on the round
$S^3(R)$, and let $A_n(\nu)$ be its Fourier compression to
$V_n^I=\operatorname{Sym}^n(\mathbb C^2)^I$. For every $t>0$,
\[
\boxed{\|f\mapsto \nu(e^{-t\Delta}f)-\mu(f)\|_{L^2\to\mathbb C}^2
=\sum_{n\ge1}(n+1)e^{-2tn(n+2)/R^2}\|A_n(\nu)\|_{HS}^2.}
\]
In degrees of invariant multiplicity one, the block is the scalar
character moment. At degree $60$, the block has size $2\times2$, so
its full Hilbert--Schmidt norm must be used instead of its trace.
For the degree-39 and degree-41 research rules the respective leading
coefficients are exactly $\sqrt{41}|\nu_{39}(C_{40})|$ and
$\sqrt{43}|\nu_{41}(C_{42})|$, with heat exponents $1680/R^2$ and
$1848/R^2$.
\end{theorem}
\begin{proof}
The Haar-orthonormal Peter--Weyl coefficients are
$\sqrt{n+1}D^{(n)}_{ab}$. Bi-invariance restricts the Fourier
transform of $\nu$ to the invariant corner $V_n^I$.
By Schur orthogonality the squared norm of the degree-$n$ contribution
to the functional is exactly $(n+1)\|A_n(\nu)\|_{HS}^2$.
The scalar heat semigroup multiplies this sector by
$e^{-tn(n+2)/R^2}$. Orthogonal summation proves the theorem.
\end{proof}


\section{Matrix-aware positive bi-icosahedral cubature through degree sixty-one}
\label{p1b:sec:cc112-positive-degree61}

Let $I=2I\subset \mathrm{SU}(2)=S^3$ be the source-owned binary
icosahedral quaternion group.  Denote by $\mathcal P_d(S^3)$ the
restrictions of real or complex polynomials of degree at most $d$.
The source invariant Molien series is
\[
 M_I(t)=\frac{1-t^{60}}{(1-t^{12})(1-t^{20})(1-t^{30})}.
\]
For $V_n=\mathrm{Sym}^n(\mathbb C^2)$ and $W_n=V_n^I$, the left--right
$I$-invariant harmonic sector at degree $n$ has dimension
$(\dim W_n)^2$.  Below degree $60$, its nonconstant degrees are
\[
 12,20,24,30,32,36,40,42,44,48,50,52,54,56,
\]
all of multiplicity one. At degree $60$, $\dim W_{60}=2$ and \textbf{four} real Fourier matrix
coordinates are required; the trace alone is insufficient.

Fix the group element
\[
 g=\Bigl(\frac{1+\sqrt5}4,0,-\frac12,\frac{1-\sqrt5}4\Bigr)\in I
\]
and define the unit-quaternion representatives
\[
 q_2(a)=\frac{a+i}{\sqrt{a^2+1}},\quad
 q_3(a)=\frac{a+i+j+k}{\sqrt{a^2+3}},\quad
 q_5(a)=\frac{a+g}{\|a+g\|}.
\]
Take the following seventeen axis representatives:
\[
\begin{split}
 q_2:&\quad a=\tfrac34,\tfrac54,\tfrac32,\tfrac94,\tfrac72,
 \tfrac92,\tfrac{23}4,\tfrac{29}4;\\
 q_3:&\quad a=\tfrac{11}4,\tfrac72,\tfrac{21}4,\tfrac{33}4;\\
 q_5:&\quad a=\tfrac32,\tfrac94,\tfrac{21}4,\tfrac{29}4,\tfrac{19}2.
\end{split}
\]
Append the two generic representatives
\[
 q_A=\frac{4+6i+3j-6k}{\sqrt{97}},\qquad
 q_B=\frac{-6-i+j+5k}{\sqrt{63}}.
\]
Write $O_j=Iq_jI$ and $M_j(f)=|O_j|^{-1}\sum_{q\in O_j}f(q)$.

\begin{theorem}[Positive sixty-thousand-node scalar cubature;
 finite algebraic certificate]
\label{p1b:thm:cc112-degree61-positive}
There exist uniquely determined, strictly positive weights
$w_j\in\mathbb Q(\sqrt5)$, on the nineteen specified double orbits,
such that
\[
 \boxed{\int_{S^3}f\,d\mu=\sum_{j=1}^{19}w_j M_j(f)
 \qquad\forall f\in\mathcal P_{61}(S^3).}
\]
The $19$ orbits are pairwise distinct and contain exactly
$5(1440)+4(2400)+8(3600)+2(7200)=60000$ distinct points.
\end{theorem}

\begin{proof}[Finite exact algebraic certificate]
For each multiplicity-one degree $n$, use the exact group character
\[
 C_n(q)=\frac1{120}\sum_{h\in I}
 U_n\bigl(\operatorname{Re}(h^{-1}q)\bigr).
\]
Let $P_{60}=120^{-1}\sum_{h\in I}\rho_{60}(h)$, and set
\[
 p_0=P_{60}z^{60},\quad
 p_1=\frac12P_{60}\bigl(\rho_{60}(r)z^{60}
 +\rho_{60}(rj)z^{60}\bigr),\quad r=\frac{1+i+j}{\sqrt3}.
\]
These span the two-dimensional real invariant space; the exact
Gram determinant is
\[
 \det\bigl(\langle p_a,p_b\rangle\bigr)_{a,b=0,1}
 =\frac{33056706500737646485082176663440359954833984375}
 {131661431568318139081706962895581110024987740209152}>0.
\]
Build the $19\times19$ exact algebraic moment matrix $A$ with first
row $1$, fourteen scalar character rows in the stated degrees, and
four rows $\langle p_a,\rho_{60}(q_j)p_b\rangle$, $a,b\in\{0,1\}$.
All entries lie in $\mathbb Q(\sqrt5)$.  Define $w=A^{-1}e_0$.
The accompanying certificate supplies the full exact moment matrix
and a rational matrix $R$ with denominator $10^{12}$.
The rational enclosure of $\sqrt5$ verifies
\[
 \|I-RA\|_\infty<9.975\cdot10^{-12},\qquad
 \|w-Re_0\|_\infty<1.300\cdot10^{-12}.
\]
The smallest entry of $Re_0$ exceeds $0.004111620393$, so all exact
weights are positive.  Because every Fourier coordinate on the
bi-$I$-invariant harmonic space through degree $61$ vanishes and
normalization holds, the rule is exact on $\mathcal P_{61}(S^3)$.
The cardinalities were checked independently using exact
$|I|^2/|I\cap qIq^{-1}|$ stabilizer calculations.  Distinctness
follows from invertibility of $A$ (otherwise two columns coincide).
\end{proof}

\begin{corollary}[First nonzero mode and matrix handedness]
\label{p1b:cor:cc112-first-alias62}
All seventeen aligned orbits have symmetric degree-$60$ invariant
Fourier blocks. The two generic orbits have opposite nonzero
antisymmetric entries, and their weighted sum cancels exactly. The
first nonconstant invariant not cancelled by the rule is degree $62$,
with certified
\[
 \nu_{61}(C_{62})\approx-0.03773628967837215\neq0.
\]
Consequently the rule is not degree-$62$ exact. Its leading
heat-smoothed scalar averaging error for the round radius-$R$ sphere
is asymptotic to
\[
 \sqrt{63}\,|\nu_{61}(C_{62})|e^{-3968t/R^2}
 \approx 0.299522514e^{-3968t/R^2}.
\]
\end{corollary}

\noindent\emph{Scope:} These are exact statements in the finite scalar
comparison geometry. Neither a physical source map nor a spinorial,
gauge, QCD or whole-action identification is asserted.


\begingroup
\small
\printbibliography
\endgroup

\section{Five-coset geometry of mobile icosahedral source charts}
\label{p1b:sec:v3023-five-coset-incidence}
The five disjoint \emph{left cosets} below are distinct from the five
conjugate binary-tetrahedral subgroups and the twenty-five overlapping
left--right transporter $24$-cells. No physical QCD link, charge,
source location or gauge field is introduced by this decomposition.

\begin{theorem}[Exact five-coset partition and inter-coset 600-cell adjacency]
\label{p1b:thm:v3023-five-coset-adjacency}
Let $G=2I\subset SU(2)$ and let $H=2T\subset G$ be a marked Hurwitz
subgroup. Choose left-coset representatives $a_1,\ldots,a_5$, and put
$C_i=a_iH$. For the nearest-neighbour graph of the 600-cell, with
quaternionic vertex product $x\cdot y=\operatorname{Re}(x^{-1}y)$,
\[
\boxed{G=\bigsqcup_{i=1}^5C_i,\quad |C_i|=24,\quad
 A_{600}[C_i,C_i]=0.}
\]
Each pair $C_i,C_j$ with $i\ne j$ supports exactly $72$ undirected
nearest-neighbour edges. Every vertex has exactly three nearest neighbours
in \emph{each} of the four other cosets. Consequently
\[
\boxed{|V|=5\cdot24=120,\quad |E|=\binom52\,72=720,\quad
 \deg_{600}=4\cdot3=12.}
\]
These are exact incidence identities, not new internal matter copies.
\end{theorem}
\begin{proof}
If $x=a_ih_1$ and $y=a_ih_2$ are distinct vertices of the same coset,
$x\cdot y=\operatorname{Re}(h_1^{-1}h_2)\le1/2$ because the Hurwitz
units form a regular 24-cell. The 600-cell nearest-neighbour threshold
is $\varphi/2>1/2$, so no such pair is adjacent. The $G$ action on
$G/H$, after quotienting the central sign, is the two-transitive
five-point action of $A_5$ on $A_5/A_4$. It preserves 600-cell adjacency;
hence all ten unordered coset-pair edge counts are equal. Since the
600-cell has $720$ nearest-neighbour edges and none is internal to a
coset, each pair has $720/10=72$. Right multiplication by $H$ preserves
each left coset, acts transitively on its $24$ vertices and is an
isometry. Thus the $72$ links between a fixed pair of cosets distribute
as $72/24=3$ neighbours per vertex in either coset. Summing the four
other cosets gives degree $12$.
\end{proof}

\begin{remark}[Mobility and source-realisation boundary]
\label{p1b:rem:v3023-coset-physical-boundary}
The native Hurwitz $24$-cell has $96$ edges at dot product $1/2$;
none is a nearest edge of the 600-cell at dot product $\varphi/2$.
Accordingly five copied Hurwitz-QCD stencils do not themselves
provide the $720$ cross-coset nearest-edge holonomies. Nor are the
five disjoint cosets five conjugate $2T$ subgroups. A common-action
cross-coset transporter, or a licensed nonlocal/Born--Fourier
operator correspondence, is a separate physical input.
\end{remark}


\section{Exact return paths from the 600-cell to Hurwitz cosets}
\label{p1b:sec:v3027-600-hurwitz-return}
Let $G=2I$ be the $120$ normalized icosians, $H=2T$ one of its
five disjoint left cosets, and $A=A_{600}$ the exact nearest-neighbour
adjacency from Theorem~\ref{p1b:thm:v3023-five-coset-adjacency}.
Order the vertices of one coset first and write
\[
 A=\begin{pmatrix}0&B\\B^*&C\end{pmatrix},\qquad
 B\in\{0,1\}^{24\times96},\quad C\in\{0,1\}^{96\times96}.
\]
On the coset let $I,N,Z,D,J$ denote the mutually disjoint symmetric
shell matrices at quaternionic scalar products $1,1/2,0,-1/2,-1$,
respectively. Their row sums are $1,8,6,8,1$. The stars above are
ordinary matrix adjoints, not physical adjoints of a decay vertex.

\begin{theorem}[Exact first four complementary return shells]
\label{p1b:thm:v3027-600-hurwitz-returns}
For \emph{each} of the five marked Hurwitz cosets the following
integer-matrix identities hold:
\[
\boxed{\begin{aligned}
 BB^*&=12I+3N,\\
 BCB^*&=60I+27N+8Z,\\
 BC^2B^*&=420I+225N+92Z+18D,\\
 BC^3B^*&=3024I+1845N+952Z+342D+12J.
\end{aligned}}
\]
In particular each native $24$-cell nearest pair is joined by exactly
three shortest $600$-cell two-step paths, while all four distinct
Hurwitz separation shells enter the complementary return series.
The complementary graph $C$ is $9$-regular.
\end{theorem}
\begin{proof}
Use the finite $\mathbb Z[\varphi]$ vertex model already fixed in
this paper: $8$ coordinate axis vectors, $16$ half-coordinate sign
vectors and $96$ even-permutation icosians with entries proportional
to $(0,1,\varphi,\varphi^{-1})$. Two distinct normalized quaternions
are nearest if their scalar product is $\varphi/2$. Quaternionic left
multiplication carries each marked coset and its five Hurwitz shell
relations into an equivalent one. For a fixed coset, partition the
$24^2$ ordered endpoint pairs into the five scalar-product classes
$1,1/2,0,-1/2,-1$ and count the $24\to96\to\cdots\to96\to24$
walks by direct multiplication of their $0$--$1$ adjacency matrices.
The resulting class coefficients, for $BC^kB^*$ at $k=0,1,2,3$,
are respectively
$ (12,3,0,0,0)$,
$ (60,27,8,0,0)$,
$ (420,225,92,18,0)$, and
$ (3024,1845,952,342,12)$.
This calculation uses exact integer and $\mathbb Z[\varphi]$
comparisons, with no floating threshold. Right multiplication by
$H$ carries the representative checks to all rows, and left
multiplication transports the identity to all five cosets. Finally,
each of the $96$ outside vertices has three neighbours in each of
the other three cosets, so $C$ is $9$-regular.
\end{proof}

\begin{theorem}[Five rational complementary resolvent modes]
\label{p1b:thm:v3027-600-hurwitz-mode-resolvent}
The Hurwitz nearest adjacency has spectral decomposition
\[
 \operatorname{spec}(N)=\{8^{(1)},4^{(4)},0^{(9)},(-2)^{(8)},(-4)^{(2)}\}.
\]
Writing $E_\lambda$ for its orthogonal spectral projectors, the
complementary resolvent satisfies, wherever it is defined,
\[
\boxed{B(xI-C)^{-1}B^*=\sum_{\lambda}F_\lambda(x)E_\lambda}
\]
with the exact rational functions
\[
\begin{array}{c|c|l}
\lambda&\operatorname{rank}E_\lambda&F_\lambda(x)\\\hline
8&1&\dfrac{36}{x-9}\\[3pt]
4&4&\dfrac{24(x+3)}{x^2-4x-24}\\[3pt]
0&9&\dfrac{4x(3x+8)}{x^3-x^2-16x-16}\\[3pt]
-2&8&\dfrac{6}{x-1}\\[3pt]
-4&2&0
\end{array}.
\]
These are identities of the finite \emph{scalar graph comparison},
not physical propagators. In particular the uniform source $u$ and
uniform complement $v$ close on
$\bigl(\begin{smallmatrix}0&6\\6&9\end{smallmatrix}\bigr)$,
with spectrum $\{12,-3\}$.
\end{theorem}
\begin{proof}
The algebra spanned by $I,N,Z,D,J$ contains all the displayed return
moments and commutes with $N$. For each $\lambda$, calculate the
first moments of $BC^kB^*$ using the preceding integer matrices.
The complement Krylov vector $B^*E_\lambda$ satisfies the following
exact annihilating polynomials in $C$:
\[
\begin{array}{c|c}
\lambda&p_\lambda(C)B^*E_\lambda=0\\\hline
8&C-9I\\
4&C^2-4C-24I\\
0&C^3-C^2-16C-16I\\
-2&C-I
\end{array}.
\]
For $\lambda=-4$ the vector itself vanishes, as proved below.
The initial scalar moments are respectively
$(36,324,2916)$, $(24,168,1248)$, $(12,44,236)$ and $(6,6,6)$.
Summing the resulting unique moment recurrences yields the displayed
rational resolvents; the formulas are not inferred from a finite
truncation alone. The uniform subspace statement follows from
$Bv=6u$, $B^*u=6v$, $Cv=9v$.
\end{proof}

\begin{theorem}[Ten exact compactly supported $C_3$-character zero modes]
\label{p1b:thm:v3027-hurwitz-character-zero-modes}
For each marked coset, the two-dimensional Hurwitz
$\lambda=-4$ subspace is annihilated by $B^*$. Extending these
vectors by zero off their coset gives ten linearly independent,
compactly supported vectors in $\ker A_{600}$, all even under
quaternionic central inversion. These are not identified with
physical stable particles or odd spinor modes.
\end{theorem}
\begin{proof}
If $N\psi=-4\psi$, then the first return identity gives
$\|B^*\psi\|^2=\langle\psi,(12I+3N)\psi\rangle=0$.
As the within-coset adjacency is zero, $A(\psi,0)=(0,B^*\psi)=0$.
The coset supports are disjoint. For a second proof, use
$H/Q_8\simeq C_3$ and its two nontrivial characters
$\chi_\pm=\omega^{\pm k}$, $\omega^3=1$.
Each external vertex meets precisely one site in each of the three
$Q_8$-classes of any given Hurwitz coset, so its summed character
amplitude is $1+\omega+\omega^2=0$. The four nearest Hurwitz elements
in each nontrivial quotient class give
$N\chi_\pm=-4\chi_\pm$.
Finally $-1\in Q_8$, so the characters are even under the central sign.
\end{proof}

\begin{remark}[No scalar-hopping locality or decay promotion]
\label{p1b:rem:v3027-return-physical-boundary}
For a \emph{hypothetical} equal-edge Hamiltonian $gA_{600}$ with
intermediate scalar shift $\Delta$ the Schur self-energy is
$g^2B((z-\Delta)I-gC)^{-1}B^*$.
For $|z-\Delta|>9|g|$ it expands in the preceding return moments.
Already its next term contains the $Z$ shell, with $D$ and $J$
appearing at higher orders; it is not the native nearest-only Hurwitz
QCD stencil. Neither $gA_{600}$ nor $\Delta$ is asserted to be the
physical action-generated regulator, and graph zero modes do not
constitute measured particle-decay prohibitions.
\end{remark}


\section{Native exceptional graph Green functions and incidence descent}
\label{p1b:sec:gc-native-graph-green}
\begin{theorem}[Exact nearest-Hurwitz and $600$-cell graph resistance]
\label{p1b:thm:gc-native-resistance-classes}
Use the unit-edge, unweighted graph Laplacians $L_{24}=8I-N$ and
$L_{600}=12I-A_{600}$. Let $L^+$ denote their Moore--Penrose inverses.
The effective resistance between vertices $u,v$ is
$R(u,v)=(e_u-e_v)^*L^+(e_u-e_v)$.
For the $24$-cell, the scalar products $1/2,0,-1/2,-1$ have respectively
$96,72,96,12$ unordered pairs and the four exact resistance values
\[
 \boxed{\frac{115}{480},\quad \frac{132}{480},\quad
 \frac{139}{480},\quad\frac{144}{480}.}
\]
For the $600$-cell, the eight nontrivial inner-product shells ordered
by increasing angular distance have populations
$(12,20,12,30,12,20,12,1)$ and exact resistances
\[
 \boxed{\left(\frac{119}{720},\frac{14293}{75600},
 \frac{737}{3780},\frac{1903}{9450},\frac{37}{180},
 \frac{5231}{25200},\frac{3179}{15120},\frac{40}{189}\right).}
\]
These are comparison-graph electrical responses, not physical masses,
transition amplitudes or decay probabilities.
\end{theorem}
\begin{proof}
For any connected vertex-transitive graph of degree $d$ with adjacency
$A$ and $n$ vertices, solve
$(dI-A)G=I-J/n$, $G\mathbf1=0$, or equivalently
$G=(dI-A+J/n)^{-1}-J/n$. For the Hurwitz graph, the existing
five-eigenspace decomposition of $N$ fixes this inverse; restricting
its entries to the four quaternionic inner-product classes gives the
listed rational values. For the $600$-cell, use the exact nine-shell
equitable quotient $B_9$ on populations $(1,12,20,12,30,12,20,12,1)$.
The unique mean-zero solution of
$(12I-B_9)g=e_0-\mathbf1/120$ is
\[
 g=\Bigl(\frac{3701}{37800},\frac{2309}{151200},
 \frac{73}{21600},\frac2{4725},-\frac1{360},
 -\frac{23}{4725},-\frac{127}{21600},
 -\frac{1091}{151200},-\frac{299}{37800}\Bigr).
\]
Vertex transitivity gives $R_k=2(g_0-g_k)$, which is the displayed
list. Every computation is a finite rational matrix identity.
\end{proof}
\begin{corollary}[Canonical two-step core incidence and edge-curvature input]
\label{p1b:cor:gc-parent-core-incidence}
For each one of the five Hurwitz cosets with coordinate projection $P_c$,
\[
 \boxed{P_cA_{600}^{2}P_c=12P_c+3N_c,\qquad
 N_c=\frac13(P_cA_{600}^{2}P_c-12P_c).}
\]
Here $N_c$ is the native eight-regular $24$-cell nearest-neighbour
adjacency on the coset. Thus the oriented $24$ vertices, $96$ edges and
$96$ triangular clique faces, and their cellular boundary maps
$d_0:C^0\to C^1$, $d_1:C^1\to C^2$, are reconstructible from $A_{600}$
and the intrinsic Hurwitz-coset family. Their unit-incidence edge-curl
operator $d_1^*d_1$ has rank $73$, kernel dimension $23$, and first
positive eigenvalue $1$, in the convention of Paper~VII-A's native
Hurwitz cellular operator. The operator is mathematical and need not
be the physical Maxwell Hessian.
\end{corollary}
\begin{proof}
The first equality is the source-proved two-step return identity;
its diagonal path count is twelve, and each internal nearest pair has
exactly three common external neighbours. All other same-coset pairs
have none. The equality recovers the exact native graph, hence its
oriented clique complex and incidence operators. The edge-curvature
spectrum is the finite native clique-complex incidence calculation;
transport through this graph isomorphism preserves its spectrum.
\end{proof}

The following reduction is for the native unweighted $600$-cell adjacency;
it is not a same-action physical heat Hessian, electromagnetic coupling or
decay law.
\begin{theorem}[Exact Hurwitz boundary Dirichlet reduction of the $600$-cell]
\label{p1b:thm:graph-kron-600-hurwitz}
Use the existing native unweighted $600$-cell adjacency decomposition
$A_{600}=\bigl(\begin{smallmatrix}0&B\\B^*&C\end{smallmatrix}\bigr)$
relative to any of its five marked Hurwitz $24$-cosets. For the
unit-incidence graph Dirichlet form $\mathscr D(f)=\tfrac12f^*(12I-A_{600})f$,
fixing the $24$ boundary values and minimizing over all $96$ complement
values produces the uniquely determined boundary form
\[
 \boxed{L_{\rm eff}=12I_{24}-B(12I_{96}-C)^{-1}B^*.}
\]
Its spectrum is
\[
 \boxed{0^{(1)}\oplus7^{(4)}\oplus(450/43)^{(9)}
        \oplus(126/11)^{(8)}\oplus12^{(2)}.}
\]
If $I,N,Z,D,J_{\rm opp}$ denote the disjoint Hurwitz shell matrices
at scalar products $1,1/2,0,-1/2,-1$, then exactly
\[
 L_{\rm eff}=\frac{56245}{5676}I-\frac{115}{132}N
 -\frac{53}{172}Z-\frac{17}{132}D-\frac{343}{5676}J_{\rm opp}.
\]
Thus $\ker L_{\rm eff}=\mathbb R\mathbf1$,
$\lambda_{\min}^{+}(L_{\rm eff})=7$.  For the standalone Hurwitz
nearest graph Laplacian $L_{24}=8I-N$ one has the exact positive
semidefinite comparison $L_{\rm eff}\succeq L_{24}$ with difference spectrum
$0^{(3)}\oplus3^{(4)}\oplus(106/43)^{(9)}\oplus(16/11)^{(8)}$.
These assertions are finite unit-incidence graph statements, not physical
operator occurrence or a measured coupling.
\end{theorem}
\begin{proof}
The connected $600$-cell graph has degree $12$, and the complement Dirichlet
principal block $12I-C$ is strictly positive: its zero-extension has
strictly positive parent graph energy unless it vanishes.  The unique
Euler--Lagrange solution is $y=(12I-C)^{-1}B^*x$, giving the stated Schur
form.  The already proved
Theorem~\ref{p1b:thm:v3027-600-hurwitz-mode-resolvent} provides the five
exact functions $F_\lambda(z)$ on the spectral projectors $E_\lambda$
of $N$ at eigenvalues $8,4,0,-2,-4$.  Their values at $z=12$ are
$12,5,66/43,6/11,0$, respectively; subtracting from $12$ proves the
stated spectrum.  Lagrange interpolation of $E_\lambda$ as polynomials
in $N$, using the five Hurwitz distance-shell relations, gives the displayed
shell coefficients; direct multiplication confirms zero row sums.
The nonzero values imply a one-dimensional kernel and positive gap $7$.
On the same projectors $L_{24}=8I-N$ has eigenvalues $0,4,8,10,12$,
so subtracting gives $0,3,106/43,16/11,0$ with the stated multiplicities.
All five cosets are equivalent by the previously proved quaternionic
left translations. Since the complement graph $C$ is $9$-regular,
$\|C/12\|=3/4<1$; the full exterior return series converges in
operator norm without a supplementary spectral-domain condition. The
$k=0$ source identity $BB^*=12I+3N$ gives $\|B\|^2=36$, so the
error after retaining terms $k=0,\ldots,n$ is bounded by
$12(3/4)^{n+1}$. Equality holds for every $n\ge0$ on the constant
Hurwitz mode.  No stochastic walk or truncated return prescription
is used.
\end{proof}


The splitting below concerns finite adjacency and finite-walk operators,
not a selected physical action.
\begin{theorem}[Exact Hurwitz-local/collective splitting of the 600-cell zero band]
\label{p1b:thm:600-zero-split-collective15}
Let $A_{600}$ be the exact nearest-neighbour adjacency on the marked
$2I\supset2T$ 600-cell and let $\Pi_{25}$ be the exact projector
$\Pi_{25}(g,h)=U_4(\operatorname{Re}(g^{-1}h))/24$, where
$U_4(t)=16t^4-12t^2+1$ and $\operatorname{rank}\Pi_{25}=25$.
For each of the five marked left Hurwitz cosets $H_i$ let $E_{-4}^{(i)}$
be the internal Hurwitz nearest-neighbour $-4$-spectral projection,
extended by zero. Then
\[
P_{10}=\sum_{i=1}^5E_{-4}^{(i)},\qquad
P_{15}=\Pi_{25}-P_{10}
\]
are mutually orthogonal projections onto adjacency-zero modes, of
ranks $10$ and $15$, respectively. For every vertex $g$,
$(P_{15})_{gg}=1/8$. Under left $2I$ acting through $A_5$ the
complexified modules are $\operatorname{ran}P_{10}\cong2\mathbf5$
and $\operatorname{ran}P_{15}\cong3\mathbf5$.
Any adjacency-zero vector supported in a single Hurwitz coset is
contained in the local rank-two $-4$ subspace of that coset.
\end{theorem}
\begin{proof}
The existing exact return identity $BB^*=12I+3N$ implies
$B^*E_{-4}^{(i)}=0$, so each local rank-two range is contained in
$\ker A_{600}$. Their disjoint supports make the sum $P_{10}$ an
orthogonal projector of rank ten. Since $\Pi_{25}$ projects onto the
entire rank-25 kernel, $P_{15}=\Pi_{25}-P_{10}$ is its orthogonal
rank-15 complement. The exact Chebyshev diagonal is $5/24$,
while each local spectral projection has diagonal $2/24$ on its
coset, giving $1/8$ throughout. For each $g\in2I$, exact finite
left-translation character traces are
$\operatorname{tr}(L_gP_{10})=2U_4(\Re g)$ and
$\operatorname{tr}(L_gP_{15})=3U_4(\Re g)$; these identities follow
from the displayed source kernels and are independently verified
for all $120$ group elements in the attached integer certificate.
The character $U_4(\Re g)$ is the five-dimensional $A_5$
character. Finally, $\ker B^*$ within any coset has exact dimension
$2$, because $12I+3N$ vanishes only on its $N=-4$ eigenspace.
\end{proof}

\begin{theorem}[Canonical five nine-planes and ten pair-triplet projections]
\label{p1b:thm:fivecoset-nineplane-fusion}
Let $Q_i$ project onto the coordinates of coset $H_i$, and let
$P_{15}$ be the collective spectral projection above. Set
$S_i=3P_{15}Q_iP_{15}$. Then each $S_i$ is an orthogonal projector
of rank nine, $\sum_iS_i=3P_{15}$, and every pair of distinct
ranges has principal-angle squared cosines $1^{(3)},(1/4)^{(6)}$.
Writing $M_{ij}=S_iS_jS_i$, the formula
\[
T_{ij}=\frac13M_{ij}(4M_{ij}-I)
\]
projects orthogonally onto the three-dimensional intersection
$\operatorname{ran}S_i\cap\operatorname{ran}S_j$. The ten such
projectors satisfy
\[
\sum_{i<j}T_{ij}=2P_{15},\qquad
T_aT_bT_a=\frac19T_a\quad(a\ne b).
\]
\end{theorem}
\begin{proof}
Within one coset, the exact quaternionic kernel and its local
$-4$ component give
$24P_{15}|_{H_i}=3I-Z+3J_{\mathrm{opp}}$, where $Z$ is
the orthogonal Hurwitz separation shell. This matrix has eigenvalues
$8^{(9)},0^{(15)}$, proving the rank-nine projector identities.
For distinct cosets the unscaled integer cross-block
$H_{ij}=(24P_{15})_{H_i,H_j}(24P_{15})_{H_j,H_i}$ has annihilating
polynomial $t(t-16)(t-64)$, first trace $288$ and second trace
$13824$, giving eigenvalues $64^{(3)},16^{(6)},0^{(15)}$.
The $M_{ij}$ spectrum is consequently $1^{(3)},(1/4)^{(6)},0$
and the displayed polynomial is its $1$-spectral projector.
The source's exact integer quaternionic matrices further satisfy
$72T_{ij}=Z_{ij}\in M_{120}(\mathbb Z)$,
$\sum Z_{ij}=6(24P_{15})$ and
$Z_aZ_bZ_a=576Z_a$ for all 45 unordered distinct pairs; these
finite identities, exhaustively checked in the attached verifier,
prove the remaining assertions.
\end{proof}

\begin{theorem}[Petersen conference from collective three-plane holonomy]
\label{p1b:thm:petersen-cubic-conference}
For any three distinct unordered coset pairs $a,b,c$, the
canonical rank-three projections satisfy
\[
T_aT_bT_cT_a=\frac{\sigma_{abc}}{27}T_a,
\quad\sigma_{abc}=+1
\iff\{a,b,c\}\text{ is a three-edge path on four of the five labels},
\]
with sign $-1$ otherwise. There are 60 triplets of each sign.
These projectors are isometrically equivalent to the ten lines
of a real $(10,5)$ equiangular tight frame, tensored with $I_3$.
If $G$ is its scalar Gram matrix, then $G^2=2G$ and $\operatorname{rank}G=5$.
With $S=3(G-I)$, a diagonal sign switch $D$ gives
\[
DSD=2A_{\mathrm{Petersen}}+I_{10}-J_{10},\qquad S^2=9I_{10}.
\]
\end{theorem}
\begin{proof}
With the common integer numerators $Z_a=72T_a$, exact matrix
multiplication gives $Z_aZ_bZ_cZ_a=\sigma_{abc}13824Z_a$.
The 120 unordered triples split, by the multigraph degree pattern
of the three selected two-subsets, into 60 paths (positive),
20 stars, 10 triangles and 30 disconnected configurations
(all negative). These identities are checked over integers by the
attached finite certificate. Fix a reference rank-three range
$W_r$, let $V_r$ be its inclusion and define
$V_a=3T_a|_{W_r}$ for $a\ne r$. Pairwise isoclinicity makes
each $V_a$ an isometry. The cubic identity gives
$V_a^*V_b=g_{ab}I_3$ with $g_{aa}=1$, $|g_{ab}|=1/3$.
Since the ten ranges span the rank-fifteen collective space by
$\sum_aT_a=2P_{15}$, the scalar Gram has rank five and
satisfies $G^2=2G$. Its signs are fixed by the cubic products.
Ordering pair labels as $01,02,03,04,12,13,14,23,24,34$,
the exact switch
$D=\operatorname{diag}(1,-1,-1,-1,-1,-1,-1,1,1,1)$
gives the displayed Petersen conference matrix.
\end{proof}

\begin{theorem}[Full unit-edge first-return law on a marked Hurwitz coset]
\label{p1b:thm:600-first-return-law}
For the simple graph walk $P=A_{600}/12$, let $T$ be its first return
to a marked 24-vertex coset. For every starting coset vertex,
$\Pr(T=n)=\frac14(\frac34)^{n-2}$ for $n\ge2$, hence
$\mathbb ET=5$ and $\mathrm{Var}(T)=12$. The joint destination/time
kernel is
\[
K(z)=\frac z{12} B\bigl(\frac{12}{z}I-C\bigr)^{-1}B^T.
\]
Its destination kernel $K(1)$ has eigenvalues
$1^{(1)},(5/12)^{(4)},(11/86)^{(9)},(1/22)^{(8)},0^{(2)}$.
Every destination inside a given coset has positive first-return
probability, with the four distinct spatial shells first reached
at path lengths $2,3,4,5$ respectively.
\end{theorem}
\begin{proof}
There are no nearest edges within a coset, while every exterior
vertex has precisely three neighbours in it and nine elsewhere.
Thus each exterior step has first-return hazard $1/4$, giving the
geometric law. A return path with $k$ intermediate exterior edges
has weight $BC^kB^T/12^{k+2}$. Its generating series gives $K(z)$.
Inserting the existing exact rational complementary resolvent
$F_\lambda(x)$ at $x=12/z$ gives the eigenvalues at $z=1$.
Exact five-class shell interpolation gives strictly positive
coefficients $11867/68112$, $115/1584$, $53/2064$,
$17/1584$, $343/68112$ at inner products
$1,1/2,0,-1/2,-1$. The already proved first four return
matrices give the stated shortest path lengths.
\end{proof}

\begin{remark}[Scope firewall]
These are exact marked finite-graph and finite-representation
theorems. In particular adjacency-zero vectors have eigenvalue
$12$ under the vertex Laplacian $12I-A_{600}$; they are not
physical zero-energy states. The collective real rank-fifteen
module is not identified with the complex $\Spin(4)$ carrier
$\mathcal H_1$. Ten mathematical rank-three intersection
channels, their signed cubic identities, and graph-step return
times do not provide a physically occurring decay operator,
particle selection rule, branching fraction or lifetime.
\end{remark}

\section{Unmarked Sylow--quintet fusion and Petersen pair projectors}
The identities use the quaternionic $600$-cell adjacency and its specified
zero band; they carry no independent mass, decay or physical role assignment.
\label{p1b:sec:unmarked-sylow-fusion}

\begin{theorem}[Unmarked Sylow construction of the five local zero bands]
\label{p1b:thm:unmarked-five-sylow-projectors}
Let $G=2I$ and let $P_{25}$ be the rank-$25$ projector of
Corollary~\ref{p1:cor:v3026-twisted25-adjacency-kernel}.  Write
$\mathcal K_{25}\cong V_5\otimes V_5^*$ for its real $5$-quintet
matrix-coefficient representation under independent left and right
translations.  The $30$ elements of order four form the noncentral
parts of exactly five Sylow-$2$ quaternion subgroups $Q_{8,i}$, whose
normalizers are the five conjugate binary-tetrahedral subgroups.
Define
\[
  \pi_i=\frac18\sum_{h\in Q_{8,i}}\rho_5(h),\qquad
  P_i=I_5\otimes\pi_i,\qquad Q_i=P_{25}-P_i.
\]
Then $\pi_i$ has rank two, $P_i$ has rank ten, $Q_i$ has rank fifteen,
and the \emph{unordered families} $\{P_i\}$ and $\{Q_i\}$ are
canonical under the full $G$ right-action.  Each $P_i$ equals the
five-coset $C_3$ character-zero projector of
Theorem~\ref{p1b:thm:v3027-hurwitz-character-zero-modes} for the
corresponding $2T$ normalizer.  In particular
\[
   \sum_iP_i=2P_{25},\qquad \sum_iQ_i=3P_{25},\qquad
   P_iP_jP_i=\tfrac14P_i\quad(i\ne j).
\]
\end{theorem}
\begin{proof}
The Sylow-$2$ subgroups of $A_5$ are five Klein four-groups; their
preimages in $2I$ are five $Q_8$ subgroups.  The Sylow normalizer is
$A_4$, lifting to $2T$.  As $\chi_5(q)=U_4(\mathrm{Re}\,q)$,
\(\operatorname{tr}\pi_i=(2\cdot5+6\cdot1)/8=2\).
The $Q_8$-fixed quintet plane is the real sum of the two nontrivial
$C_3=2T/Q_8$ characters, identifying $P_i$ with the source's supported
character-zero projector under the five left $2T$ cosets.
Schur's lemma on the real irreducible quintet gives
$\sum\pi_i=2I_5$ by its trace.
The five-point action is two-transitive.  Its ordered pair stabilizer
is $C_3$, which acts on $\operatorname{ran}\pi_i$ as a nontrivial real
rotation.  Hence the self-adjoint compression $\pi_i\pi_j\pi_i$ is a
scalar multiple of $\pi_i$.  Its trace is
$4\operatorname{tr}(\pi_i\pi_j)=\operatorname{tr}\bigl(\pi_i(2I_5-\pi_i)\bigr)=2$,
so $\operatorname{tr}(\pi_i\pi_j)=1/2$.  Since $\operatorname{rank}\pi_i=2$,
the scalar is $1/4$.  Tensoring with $I_5$ proves the result.
\end{proof}

\begin{theorem}[Canonical rank-five pair channels and Petersen holonomy]
\label{p1b:thm:unmarked-ten-pair-five-planes}
For the five $Q_i$ above, the operator
\[
  T_{ij}:=\frac{4Q_iQ_jQ_i-Q_i}{3}\quad(1\le i<j\le5)
\]
is the rank-five orthogonal projector onto
$\operatorname{ran}Q_i\cap\operatorname{ran}Q_j$.  The ten projectors
are canonical as one $A_5$-covariant family and satisfy
\[
 \sum_{i<j}T_{ij}=2P_{25},\quad
 \sum_{j\ne i}T_{ij}=\frac43Q_i,\quad
 T_aT_bT_a=\frac19T_a\ (a\ne b).
\]
For three distinct pair labels $a,b,c$,
\[
 T_aT_bT_cT_a=\frac{\sigma_{abc}}{27}T_a,
 \qquad\sigma_{abc}=+1
 \iff\text{the three pair-edges form a path on four distinct labels}.
\]
In all other triple types the sign is negative.  There are $60$
positive and $60$ negative triples.  The scalar ten-line Gram has rank
five and Seidel matrix switching-equivalent to
$2A_{\mathrm{Petersen}}+I-J$.
\end{theorem}
\begin{proof}
The two-dimensional right quintet fixed planes $\operatorname{ran}\pi_i$
have squared principal cosines $1/4$.  Their three-dimensional
orthogonal complements intersect in a one-dimensional line and meet
with two other squared principal cosines $1/4$.  Tensoring the left
quintet gives the rank-five intersection and the polynomial spectral
projection formula.  The sum over all ten intersections commutes with
$A_5$ on the right factor, and has trace $10$ there, hence equals
$2I_5$; tensoring yields the first identity.  For a fixed i the four
intersection lines lie in its irreducible three-dimensional $A_4$
complement.  Their sum commutes with $A_4$ and has trace four, so it
is $(4/3)I_3$, proving the star identity.  This identity forces
squared overlap $1/9$ for the three other incident lines; the full
$10$-line tight identity then forces the same value for the remaining
three disjoint line labels, proving the pair formula.  The cubic signs and
switching class are determined by the four orbit types of unordered
triples of edges of $K_5$; the attached exact integer certificate
checks each of the 120 identities in the 120-vertex quaternion model.
\end{proof}

\begin{corollary}[Unmarked normalized positive-effect resolutions]
\label{p1b:cor:unmarked-pair-povms}
On $\mathcal K_{25}$ the families
$\{P_i/2\}$, $\{Q_i/3\}$ and $\{T_{ij}/2\}$ are POVMs.
On $\operatorname{ran}Q_i$ the four effects
$\{(3/4)T_{ij}:j\ne i\}$ form a POVM.  For any normalized state
supported on $P_i$ the five $P$-outcomes have probabilities
$1/2,1/8,1/8,1/8,1/8$; for any state supported on $T_a$ the ten
$T$-outcomes have probabilities $1/2$ and $1/18$ for each other label.
These statements do not assert an action-generated detector.
\end{corollary}
\begin{proof}
All effects are positive because each summand is a nonnegative
multiple of a projector.  Their normalizations are the preceding
fusion identities.  The probabilities follow by compressing the
effects onto the stated preparation projectors.
\end{proof}

\begin{remark}[Sharp symmetry and physical-identification boundary]
The $25$-band is irreducible under independent left and right $2I$
actions, being the external tensor product $V_5\boxtimes V_5^*$.
Consequently no nonzero proper equivariant projector, in particular
no invariant rank-fifteen projector, is selected by the unmarked
bi-action.  The family of five rank-fifteen projectors is covariant;
a preferred member needs further source data.  Nor is this fifteen-
dimensional *even* quintet bimodule thereby identified with the
Spin$(4)$ level-one $\mathcal H_1$, which has central-odd spinor
sectors.  No physical source, Hamiltonian, decay operator or clock is
claimed.
\end{remark}


\subsection{Two-shell pair transport and coherent distinction on the icosahedral quintet}
The shell comparison is an exact graph/operator calculation with no implied
physical evolution law, decay rate or source clock.
\label{p1b:sec:two-shell-coherent-transport}
Let $G=2I$, and let $\{Q_{8,i}\}_{i=1}^5$ be its five intrinsic Sylow-$2$
quaternion subgroups; their normalizers are the conjugate binary-tetrahedral
subgroups.  Let $\mathscr E=\binom{\{1,\ldots,5\}}2$, and let
$H\in\{0,1\}^{10\times5}$ be the incidence matrix
$H_{a,i}=\mathbf 1_{i\in a}$.  Introduce the two distinct twelve-element
conjugacy shells
\[
 S_+=\{g\in G:\operatorname{Re}g=(1+\sqrt5)/4\},\qquad
 S_-=\{g\in G:\operatorname{Re}g=(\sqrt5-1)/4\}.
\]
The first is the vertex-nearest $600$-cell class; the second is the
F40 shell of Theorem~\ref{p1:thm:v345-f40-edge-regular}.  Each acts on the
five $Q_8$ subgroups by twelve five-cycles.  The five-mark transfer
$3(J_5-I_5)$ is already established by the cited source theorem for
$S_-$; its independent $S_+$ check below is a cross-source bridge, not a
new claim to first discovery.

\begin{theorem}[Exact shell-independent ten-pair kernel]
\label{p1b:thm:two-shell-ten-pair-kernel}
For either $S\in\{S_+,S_-\}$, let
$M_S(a,b)=\#\{g\in S:g\cdot a=b\}$ on $a,b\in\mathscr E$.
Then
\[
 \boxed{M_S=2J_{10}-HH^{\mathsf T}},\qquad
 M_S(a,b)=\begin{cases}
 0,&a=b,\\1,&|a\cap b|=1,\\2,&a\cap b=\varnothing.
 \end{cases}
\]
In particular $K_{10}=M_S/12$ is independent of the selected shell,
\[
 \boxed{\operatorname{spec}(K_{10})=
 \{1^{(1)},(-1/4)^{(4)},0^{(5)}\}},\qquad
 K_{10}^n=P_0+(-1/4)^nP_4\quad(n\ge1),
\]
where $P_0=J_{10}/10$, $P_4=H(H^{\mathsf T}H)^{-1}H^{\mathsf T}-P_0$,
and $P_5=I_{10}-P_0-P_4$ are the canonical $(1,4,5)$ projectors.
The associated weighted unit-event network has exact effective resistances
$7/45$ for intersecting pair labels and $13/90$ for disjoint labels;
its one-way first hitting times are respectively $28/3$ and $26/3$
dimensionless steps.
\end{theorem}
\begin{proof}
Each shell projects to one complete twelve-element five-cycle conjugacy
class in $A_5$.  No five-cycle fixes a two-element subset.  For a fixed
$a=\{i,j\}$, transitivity of the point stabilizer implies that exactly
three shell elements carry $i$ to each distinct mark, and likewise for
$j$.  The six events $g(i)=j$ or $g(j)=i$ are disjoint because a five-cycle
cannot contain a transposition.  By pair-stabilizer transitivity the six
intersecting targets each occur once; the other six shell elements land
equally in the three disjoint targets, each twice.  This proves $M_S$.
Since $H^{\mathsf T}H=3I_5+J_5$, the ten-pair permutation representation
splits as $\mathbf1\oplus\mathbf4\oplus\mathbf5$, on which
$M_S$ acts as $12,-3,0$.  Hence the stated spectral projectors and
powers.  The weighted graph Laplacian $L=12I-M_S$ has pseudoinverse
$L^+=P_4/15+P_5/12$.  Direct compression of the incidence projectors
then gives effective resistance $7/45$ for one shared mark and $13/90$
for disjoint marks.  The graph is regular, symmetric, and pair-transitive;
the commute-time formula is $\mathbb E_aT_b+\mathbb E_bT_a=120R_{ab}$,
and symmetry makes the two hitting times equal.  The stated values follow.
All ninety ordered first hitting times are additionally checked by exact
rational absorbing-chain inversions in the attached finite certificate.
\end{proof}

\begin{theorem}[Basis-free two-shell quintet adjoint spectrum and triplet separation]
\label{p1b:thm:two-shell-adjoint-triplet}
Let $W=\ker H^{\mathsf T}$, carrying the real irreducible quintet of
$A_5$ with $\dim W=5$, and let $\rho_5$ be its orthogonal action.
For $X\in\operatorname{End}(W)$ define the mathematical maps
\[
 \Phi_\pm(X)=\frac1{12}\sum_{g\in S_\pm}
          \rho_5(g)X\rho_5(g)^{*}.
\]
These are unital, completely positive, trace preserving, self-adjoint
Hilbert--Schmidt channels.  In the $A_5$ isotypic resolution
\[
 \operatorname{End}(W)=\mathbf1\oplus\mathbf3\oplus\mathbf3'
                     \oplus2\mathbf4\oplus2\mathbf5,
\]
where $\mathbf3$ is the standard spatial icosahedral triplet, one has
\[
 \begin{array}{c|ccccc}
 &\mathbf1&\mathbf3&\mathbf3'&2\mathbf4&2\mathbf5\\ \hline
 \Phi_+&1&(1+\sqrt5)/6&(1-\sqrt5)/6&-1/4&0\\
 \Phi_-&1&(1-\sqrt5)/6&(1+\sqrt5)/6&-1/4&0
 \end{array}
\]
with total multiplicities $1,3,3,8,10$.  Put
$C=(\Phi_++\Phi_-)/2$ and $D=\Phi_+-\Phi_-$.  They satisfy
\[
 \boxed{CD=DC=D/6,\quad D^3=5D/9,\quad
 C(C-I)(C-I/6)(C+I/4)=0.}
\]
The two inequivalent rank-three sector projectors are fixed by
\[
 \boxed{\Pi_{3,\pm}=\frac12\left(\frac95D^2
                                   \pm\frac3{\sqrt5}D\right).}
\]
Thus the two shell transports distinguish the triplet sectors without
postulating a choice of quintet coordinates, three-dimensional bases,
relative weights, or a new splitting parameter.
\end{theorem}
\begin{proof}
The five-dimensional module $W$ is the orthogonal complement of
$\operatorname{im}H$ in the ten-pair permutation module.  Its character
is $\chi_W(g)=\#\mathrm{Fix}_{\mathscr E}(g)-\#\mathrm{Fix}_{\{1,\ldots,5\}}(g)$,
with values $(5,1,-1,0,0)$ on the five $A_5$ conjugacy classes,
identical to the already source-verified quintet character
$U_4(\operatorname{Re}g)$.  Its tensor square decomposes as displayed.
Each $S_\pm$ is a complete conjugacy class, so its normalized group
algebra class sum is central and acts on an irreducible module $V$
by $\chi_V(S_\pm)/\dim V$.  The relevant character values are
$1$, $\varphi$ and $1-\varphi$ (exchanged between the two classes),
$-1$ and $0$, with $\varphi=(1+\sqrt5)/2$.  The displayed spectra follow.
Each shell contains inverses and all summands are unitary conjugations;
this proves complete positivity, trace preservation, unitality, and
Hilbert--Schmidt self-adjointness.

The average $C$ acts with eigenvalues $1,1/6,-1/4,0$ on subspaces
of dimensions $1,6,8,10$, respectively.  The difference $D$ vanishes
outside the two triplets and acts on them with eigenvalues
$\pm\sqrt5/3$.  Thus the polynomial identities follow.  Defining
$Q=(9/5)D^2$ gives $Q^2=Q$, $QD=D$, $D^2=5Q/9$, $\operatorname{rank}Q=6$;
these identities prove that the two displayed $\Pi_{3,\pm}$ are
orthogonal idempotents of rank three.  An independent exact rational
$25\times25$ matrix replay constructs both class sums using the
quaternionic multiplication table and verifies all identities without
numerical eigensolvers.
\end{proof}

\begin{corollary}[Blindness of pair-projector marginals and exact comparison rates]
\label{p1b:cor:two-shell-pair-blindness}
Let $T_a$ denote the existing unmarked ten pair-channel projectors on
the rank-$25$ adjacency-zero band (or their corresponding rank-one
right-quintet factors).  For any state $\rho$, any $a\in\mathscr E$,
and any $n\ge0$, the single-time pair-effect statistics of the two
mathematical shell channels coincide:
\[
 \boxed{\operatorname{Tr}\bigl(T_a\Phi_+^n(\rho)\bigr)=
        \operatorname{Tr}\bigl(T_a\Phi_-^n(\rho)\bigr).}
\]
The coherent $\mathbf3$ and $\mathbf3'$ sectors can distinguish the two.
For the dimensionless Poissonized generator $\mathcal L_\pm=\Phi_\pm-I$,
the nonzero decay exponents on the right quintet for $\Phi_+$ are
\[
 \boxed{\frac{5-\sqrt5}{6},\quad\frac{5+\sqrt5}{6},
          \quad\frac54,\quad1},
 \qquad\frac{\gamma_{3'}}{\gamma_3}=
                 \frac{3+\sqrt5}{2}=\varphi^2.
\]
On the full $25$-dimensional zero band the fixed-point algebra of the
right-channel extension has dimension $25$ (the left factor is untouched).
The result is a dimensionless finite mathematical relaxation spectrum,
not a physical clock, lifetime, branching fraction, action-source
occurrence or instrument.
\end{corollary}
\begin{proof}
The ten pair effects span the ten-pair permutation representation
$\mathbf1\oplus\mathbf4\oplus\mathbf5$, and hence have zero projection
to either triplet.  By the previous theorem $D$ annihilates this span;
as the two channels commute, their powers agree on every such effect.
Hilbert--Schmidt adjointness transfers this to the trace expressions.
The Poissonized channel has exponents $1-\lambda$ for each eigenvalue
$\lambda$ of $\Phi_+$, giving the displayed values and ratio.  The
rank-five right quintet is irreducible, so the right-channel fixed
space is scalar.  The source's rank-$25$ coefficient representation
is left-quintet tensor right-quintet; conjugation on the right leaves
all left endomorphisms fixed, yielding dimension $5^2=25$.
\end{proof}

\begin{corollary}[No physical decay arrow from the bare shell kernel]
\label{p1b:cor:two-shell-bare-arrow-no-go}
The ten-channel kernels of the preceding theorem are symmetric,
irreducible and aperiodic, with strictly positive off-diagonal entries.
Their unique stationary distribution is uniform, all stationary
probability currents vanish, and no pair label is absorbing or forbidden.
The right-quintet channels $\Phi_\pm$ are unital and their unique
stationary right density matrix is $I_5/5$.  Consequently neither bare
shell transport can independently produce a preferred irreversible
mass-ordered particle decay direction; such a physical statement would
require a separately derived occurring source and physical-state map.
\end{corollary}
\begin{proof}
The exact entries of $K_{10}$ are $1/12$ or $1/6$ off diagonal, and
$K_{10}$ is symmetric.  Complete connectivity gives irreducibility;
triangles give period one; symmetry gives detailed balance for the
uniform law.  The adjoint channel spectrum has a single eigenvalue
$1$ on $\operatorname{End}(W)$, carried by the identity, proving the
stationary right-density claim.
\end{proof}

\begin{remark}[Sharp scope and non-promotion]
The $5$-core kernel and the Petersen $P_5$ provide the input for the exact
$10$-pair pushforward for both shell classes, its resistance/hitting law,
the shell-independent classical versus shell-distinct coherent resolution,
and the rational-polynomial projectors.  The uniqueness of equal weights
is scoped to conjugacy-invariant random-unitary averages supported on
one complete shell, and does not select stochastic versus unitary
physical evolution.  In particular no particle decay operator,
phenomenological rate, state/observable identification, or physical
source clock is proved here.  The existing GP3 and G0 proof-chain
conditions are neither modified nor silently discharged.
\end{remark}

\section{Native shell powers, coherent frame readouts, and observer compression}
\label{p1b:sec:single-shell-readout-proposal}

\begin{theorem}[Nearest quaternion powers generate the F40 shell]
\label{p1b:thm:native-shell-square}
In the fixed quaternionic $2I$ let
\[
 S_1=\{g:\operatorname{Re}g=(1+\sqrt5)/4\},\qquad
 S_2=\{g:\operatorname{Re}g=(\sqrt5-1)/4\}.
\]
The squaring map $g\mapsto g^2$ is a bijection $S_1\to S_2$.
For $g\in S_1$ one has $g^5=-1$, $g^{10}=1$; for $h\in S_2$
one has $h^5=1$. Hence centre-even representations see period five
under repeated application of one $S_1$ generator, whereas the
120-vertex left or right group walk detects the antipode at step five.
\end{theorem}
\begin{proof}
The quaternion identity $\operatorname{Re}(g^2)=2\operatorname{Re}(g)^2-1$
puts every square in $S_2$. Both conjugacy classes have size twelve, and
squaring maps $S_1$ onto a single conjugacy class; equivalently the
exact native 120-element multiplication certificate finds twelve distinct
squares and exact equality with $S_2$. Quaternion polar angles
$\pi/5$ and $2\pi/5$ give the order and central-power statements.
\end{proof}

\begin{theorem}[Single-shell polynomial generation of the coherent quintet algebra]
\label{p1b:thm:single-shell-adjoint-algebra}
On $\mathcal H=\operatorname{End}(V_5)$ let
\[
 A=\frac1{12}\sum_{g\in S_1}\operatorname{Ad}_{\rho_5(g)},\quad
 B=\frac1{12}\sum_{g\in S_2}\operatorname{Ad}_{\rho_5(g)}.
\]
The second operator lies in the unital algebra generated by the first:
\[
 \boxed{B=\frac{-432A^4+576A^3-81A^2-58A}{5}},
 \qquad A(A-I)(A+\tfrac14I)(A^2-\tfrac13A-\tfrac19I)=0.
\]
The rank-six paired-triplet projector is
\[
 P_6=-\frac{648}{5}A(A-I)(A-\tfrac12I)(A+\tfrac14I),
 \qquad P_6^2=P_6,\quad \operatorname{tr}P_6=6.
\]
Putting $D=(2A-\tfrac13I)P_6$, the two rank-three projectors are
$P_{3,\pm}=\tfrac12(P_6\pm 3D/\sqrt5)$.
\end{theorem}
\begin{proof}
The $A_5$-module $\mathcal H=\mathbf1\oplus\mathbf3\oplus\mathbf3'
\oplus2\mathbf4\oplus2\mathbf5$ is a direct sum of the five
class-sum eigenspaces of $A$ with eigenvalues
$1,(1+\sqrt5)/6,(1-\sqrt5)/6,-1/4,0$. Squaring the quaternion
generators exchanges the two five-cycle conjugacy classes; hence the
$B$ eigenvalues agree with those of $A$ outside the $3,3'$ blocks and
exchange the two golden eigenvalues within those blocks. The displayed
rational polynomial is the unique polynomial of degree at most four with
those five values. The projector polynomial takes value one on both
golden roots and zero on the other three eigenvalues. The two further
idempotents follow by the positive $\mathbb Q(\sqrt5)$ embedding.
All full $25\times25$ rational matrix identities have been replayed
independently. Note $B\ne A^2$: squaring a sampled generator before
averaging is different from multiplying two independently averaged
operators.
\end{proof}

\begin{theorem}[Unmarked 60-ray coefficient-frame POVM and exact character zeros]
\label{p1b:thm:coefficient-orbit-povm}
For the real, centre-even irreducible $A_5$ quintet $V_5$ put
$v_g=\operatorname{vec}(\rho_5(g))\in\mathcal H$, $g\in A_5$.
The diagonal-conjugation-fixed line in $\mathcal H$ is the identity line,
and its $A_5\times A_5$ orbit comprises exactly the sixty rays $[v_g]$.
Schur orthogonality implies
\[
 \sum_{g\in A_5}|v_g\rangle\langle v_g|=12I_{25},\qquad
 E_g=\frac1{12}|v_g\rangle\langle v_g|\succeq0,\quad
 \sum_g E_g=I_{25},\quad\operatorname{tr}E_g=5/12.
\]
For the normalized coefficient ray $\psi_h=v_h/\sqrt5$, the
probability distribution of this mathematical POVM is
\[
 \boxed{p(g\mid h)=\frac{|\chi_5(h^{-1}g)|^2}{60}},
\]
with relative identity probability $5/12$, relative order-two and
order-three probabilities $1/60$ each, and precisely twenty-four
zero-probability relative five-cycle outcomes. Its symmetric stochastic
kernel has exact spectrum
\[
 \boxed{\operatorname{spec}K=
 \{1^{(1)},(1/2)^{(16)},(2/5)^{(25)},(1/3)^{(18)}\}.}
\]
\end{theorem}
\begin{proof}
Absolute irreducibility of the quintet gives the unique identity-fixed
line by Schur's lemma. Left--right action generates all representation
matrices. The source-established full-group coefficient identity
$F^\dagger F=24I_{25}$ has the central duplicate $\pm1$;
dividing by two gives the displayed $60$-ray tight frame and fixes
its effect normalization. Frame-state Born probabilities equal the
square of the character as shown. The class character has values
$5,1,-1,0,0$ on $1,2A,3A,5A,5B$; their class sizes are
$1,15,20,12,12$. Finally
$\chi_5^2=\chi_1+\chi_3+\chi_{3'}+2\chi_4+2\chi_5$;
Schur orthogonality gives the scalar $m_\lambda/d_\lambda$
on each irreducible block in the regular representation, with
multiplicity $d_\lambda^2$. All $3600$ exact character-square
probabilities and the full rational frame identity were checked.
\end{proof}

\begin{corollary}[Canonical 15-dimensional pair-line blind complement]
\label{p1b:cor:pairline-blind15}
Let $Q_a$ be the ten unit-trace rank-one Petersen projected-line
projectors in $V_5$. They obey
\[
 \operatorname{tr}(Q_aQ_b)=\delta_{ab}+\frac{1-\delta_{ab}}9,
 \qquad\sum_{a=1}^{10}Q_a=2I_5.
\]
Their Hilbert--Schmidt span is ten-dimensional, with orthogonal projector
\[
 \mathscr P_{\rm vis}(X)=\frac98\sum_a\operatorname{tr}(Q_aX)Q_a
 -\frac14\operatorname{tr}(X)I_5.
\]
The complementary rank-fifteen projector
$\mathscr P_{\rm dark}=I_{\operatorname{End}(V_5)}-\mathscr P_{\rm vis}$
has diagonal $A_5$ type $\mathbf3\oplus\mathbf3'\oplus\mathbf4\oplus\mathbf5$.
It satisfies
$\mathscr P_{\rm vis}(A-B)=(A-B)\mathscr P_{\rm vis}=0$.
Thus the ten line effects cannot distinguish $A^n$ from $B^n$ via
single-time expectations for any integer $n\ge0$.
\end{corollary}
\begin{proof}
The ten-effect Gram is $G=8I_{10}/9+J_{10}/9$,
whose inverse is $9I_{10}/8-J_{10}/16$. This gives the projector
formula and its rank. Covariance permutes the ten line effects,
so their representation is $\mathbf1\oplus\mathbf4\oplus\mathbf5$.
Subtracting from the $25$-dimensional tensor-square character leaves
$\mathbf3\oplus\mathbf3'\oplus\mathbf4\oplus\mathbf5$.
Both class sums commute with the visible projector and differ only
on the triplets. Exact $25\times25$ rational matrix checks establish
idempotency, rank, commuting, and difference annihilation.
\end{proof}

\begin{proposition}[One-step shell symmetry alone cannot select Markov independence]
\label{p1b:prop:twostep-memory-no-go}
A uniform class generator drawn independently at each step and a
uniform class generator drawn once and reused have identical native
nearest-shell one-step marginals, equivariance, and induced five-mark
and ten-pair one-step kernels. Their respective two-step return
probabilities are $1/4$ versus $0$ on the five marks and $1/8$ versus
$0$ on the ten pair labels.
\end{proposition}
\begin{proof}
With $K_5=(J_5-I_5)/4$ and $K_{10}=(2J_{10}-HH^T)/12$,
independently sampled successive generators give $K_5^2$ and
$K_{10}^2$, with the stated diagonal entries. Reusing the first
generator twice gives the squared-generator F40 shell, whose induced
kernels equal $K_5$ and $K_{10}$ and have zero diagonal. The
native quaternion multiplication table verifies both paths exactly.
\end{proof}

\begin{remark}[No inferred physical occurrence, dynamics, or complete density tomography]
\label{p1b:rem:single-shell-observable-boundary}
All claims concern finite $2I$ representations and the specified
covariant mathematical maps. The $60$-effect POVM is normalized but
not informationally complete for arbitrary $25\times25$ quantum
density operators. Its exact five-cycle readout zeros are not forbidden
particle decays. The derived blind $15$ is representation-equivalent
to the source's Petersen $\mathfrak{sl}_4$ representation, not a
canonically identified Lie subalgebra or the physical $H_1$ momentum
carrier. One-step geometry does not select memoryless dynamics or a
physical time unit. The actual same-action source, observable and
physical state embeddings remain separate obligations.
\end{remark}

\section{Quadratic Petersen-pair readout on the intrinsic quintet}
The five-dimensional projected pair-lines below are not the
three-dimensional Hopf axes, and their mathematical measurement frame is not
a claim of an experimentally realized detector.
\label{p1b:sec:v17-coherent-pair-readout-proposal}
\begin{theorem}[Parameter-free informationally complete ninety-effect readout]
\label{p1b:thm:v17-petersen-pair-ic-povm}
Let $D$ be the $5\times10$ pair-incidence matrix and set
$P=I_{10}-D^\mathsf T(3I_5+J_5)^{-1}D$. On the five-dimensional
real quintet $W=\operatorname{im}P$, let
$Q_a=2Pe_ae_a^\mathsf TP$, $a=0,\ldots,9$. Then
$\sum_aQ_a=2I_W$, $Q_a^2=Q_a$, and
$\operatorname{tr}(Q_aQ_b)=1/9$ for $a\ne b$.
For $a<b$ put
\[
 C_{ab}=[Q_a,Q_b],\quad
 P_{ab}=\frac98(Q_a+Q_b-Q_aQ_b-Q_bQ_a),\quad
 J_{ab}=\frac9{2\sqrt2}iC_{ab},
 \quad R_{ab}^\pm=\frac12(P_{ab}\pm J_{ab}).
\]
Each $R_{ab}^\pm$ is rank-one orthogonal projection, and
\[
 \boxed{E_{ab}^\pm=\frac1{18}R_{ab}^\pm,\qquad
 \sum_{a<b,\pm}E_{ab}^\pm=I_W.}
\]
The ninety effects span $\operatorname{Herm}(W)$ over the reals,
with dimension $15+10=25$. The construction is covariance- and
sign-gauge invariant as a complete outcome set.
\end{theorem}
\begin{proof}
By the source Petersen projector formula the ten normalized projected
coordinate lines form a rank-five real equiangular tight frame, giving
the first three identities. On the plane of two distinct lines of squared
overlap $1/9$, direct rank-one multiplication gives
$C_{ab}^2=-(8/81)P_{ab}$, $P_{ab}^2=P_{ab}$, and
$\operatorname{tr}P_{ab}=2$. Hence $J_{ab}^2=P_{ab}$ and its two signs
split that plane into rank-one orthogonal complex eigenspaces.
Also
\[
 \sum_{a<b}P_{ab}=\frac98\Bigl(
 9\sum_aQ_a-\bigl[(\sum_aQ_a)^2-\sum_aQ_a^2\bigr]\Bigr)
 =\frac98(18I-2I)=18I.
\]
The original $Q_a$ have rank-ten real span, but the forty-five $P_{ab}$
have an exact nonzero $15\times15$ Gram minor $3^{14}/2^{32}$.
The forty-five $iC_{ab}$ have an exact nonzero $10\times10$ Gram minor
$2^{26}/3^{36}$. These occupy orthogonal symmetric and imaginary
antisymmetric Hermitian subspaces of dimensions 15 and 10; they
therefore span all Hermitian operators. The exact minor indices and
rational reconstruction are supplied in the v1.7 verifier.
\end{proof}

\begin{corollary}[Exact, well-conditioned tomography frame]
\label{p1b:cor:v17-petersen-ic-frame}
For $\mathcal F(X)=\sum_{a<b,\epsilon}E_{ab}^{\epsilon}
\operatorname{tr}(E_{ab}^{\epsilon}X)$ one has
\[
 \boxed{\operatorname{Spec}\mathcal F
 =(1/18)^{(1)}\oplus(43/5184)^{(9)}
 \oplus(1/576)^{(5)}\oplus(1/72)^{(10)}}.
\]
Thus the frame bounds are $1/576$ and $1/18$, its eigenvalue
condition number is 32, and every Hermitian $X$ obeys
$X=\mathcal F^{-1}\sum_{a<b,\epsilon}
\operatorname{tr}(E_{ab}^{\epsilon}X)E_{ab}^{\epsilon}$.
\end{corollary}
\begin{proof}
The real and imaginary contributions from the two signs are orthogonal.
For the 45 real plane projectors the exact rational frame Gram has
nonzero eigenvalues $1/18$ once, $43/5184$ nine times, and $1/576$
five times. The 45 real skew commutators yield the scalar frame
$1/72$ ten times. These follow from the exact $15\times15$ and
$10\times10$ compressed Gram matrices in the verifier. All eigenvalues
are positive, proving the inverse without assuming it.
\end{proof}

\begin{theorem}[Lie generation by the original ten line projectors]
\label{p1b:thm:v17-petersen-u5-generation}
On the complexification of $W$, the Lie algebra generated by the ten
skew-Hermitian operators $iQ_a$ is all $\mathfrak u(5)$, with 25 real
dimensions. Bracket depth at most three suffices. Its full associative
$\ast$-algebra is $M_5(\mathbb C)$ and has scalar commutant.
\end{theorem}
\begin{proof}
Choose five of the spanning real line vectors that are independent;
every off-diagonal inner product is nonzero. Their ten pair commutators
are nonzero multiples of a basis of $\Lambda^2W=\mathfrak{so}(5)$.
The $\mathfrak{so}(5)$ action on a nonzero traceless real rank-one
quadrupole generates $\operatorname{Sym}^2_0(W)$ under commutators,
and $\sum_a iQ_a=2iI_W$ supplies the center. This yields
$\mathfrak{so}(5)\oplus i\operatorname{Sym}^2(W)=\mathfrak u(5)$.
An exact explicit depth-three symmetric Gram minor is
$2^{40}/3^{42}\ne0$; ten first commutators independently span
$\mathfrak{so}(5)$. The full matrix algebra and scalar commutant follow.
\end{proof}

\begin{corollary}[Exact shell and memory witnesses in the derived readout]
\label{p1b:cor:v17-petersen-shell-discrimination}
Let $\Phi_{\rm near}$ and $\Phi_{\rm F40}$ be the source's
nearest and F40 class-averaged adjoint channels on $\operatorname{End}(W)$.
Enumerate the ten original line labels lexicographically as $0,\ldots,9$.
For preparation $R^+_{0,1}$ and effect $E^+_{0,5}$,
\[
 \operatorname{tr}(E^+_{0,5}\Phi_{\rm near}(R^+_{0,1}))=\frac3{256},\qquad
 \operatorname{tr}(E^+_{0,5}\Phi_{\rm F40}(R^+_{0,1}))=\frac{11}{2304}.
\]
The difference is exactly $1/144$. The independently resampled
nearest generator channel $\Phi_{\rm near}^2$ gives $13/1024$ in
the same test, whereas squaring the same generator before averaging
gives the F40 value $11/2304$; difference $73/9216$.
\end{corollary}
\begin{proof}
Conjugation by each of the twelve native quaternion generators acts
as its exact permutation on the ten unordered pairs. Apply the
permutation to the rational plane and commutator matrices and evaluate
$\operatorname{tr}(E\Phi(R))$. All $12$, $12$, and $144$ terms are exact
rational; the v1.7 source-data verifier enumerates each and checks the
printed fractions. The comparison distinguishes mathematically
specified channels, not the physical occurrence of one of them.
\end{proof}

\begin{remark}[No physical owner or false fifteen-dimensional Lie identification]
The ten *initial single* line effects still have a fifteen-dimensional
linear blind complement. The quadratic ninety-effect family repairs
that mathematical readout limitation; no actual detector or source
is established. That fifteen-dimensional linear complement contains
all $\mathfrak{so}(5)$ and a nonzero traceless symmetric component,
so its Lie envelope is $\mathfrak{sl}_5(\mathbb R)$, not a Lie copy of
the separate Petersen $\mathfrak{sl}_4(\mathbb R)$. This theorem has
no effect on the already proved $m_p=1$ or on any physical decay width.
\end{remark}


\clearpage
\appendix
\part*{Technical appendices}
\addcontentsline{toc}{part}{Technical appendices}
\section{Supplementary anchored exceptional roots in the Clifford carrier}
\label{p1:app:exroot-anchored-construction}

Use the real Clifford algebra already constructed in
Theorem~\ref{p1:thm:hurwitz-walsh-clifford}.  For its ordered orthonormal
negative generators \(e_1,e_2,e_3\), put
\[
 i=e_2e_3,\qquad j=e_3e_1,\qquad k=e_1e_2,\qquad
 \omega=e_1e_2e_3.
\]
The Clifford relations give \(i^2=j^2=k^2=ijk=-1\), central
\(\omega\) with \(\omega^2=1\), and the parity decomposition
\[
 A=\mathrm{Cl}(0,3)=\mathbb H\oplus\mathbb H\omega.
\]
Give the eight displayed Clifford monomials their coefficient inner product.
Then the scalar line is \(\mathbb R1\), its orthogonal complement is the
sum of the seven other monomial lines, and the quaternion list is
\eqref{p1:eq:2I-explicit}.  The following field anchoring is a specified
module construction on this entirely internal mathematical datum.

\begin{proposition}[Unique specified golden action]
\label{p1:prop:exroot-golden-action}
Use the orthonormal coordinates
\[
 (1,i,j,k,\omega,i\omega,j\omega,k\omega),\qquad
 \langle x+y\omega,x'+y'\omega\rangle
 =\operatorname{Re}(\bar x x'+\bar y y').
\]
There is a unique real-linear map \(\mathsf F\) commuting with every
even-quaternion left multiplication, satisfying
\(\mathsf F^2=\mathsf F+I\) and \(\mathsf F1=\omega\). It is self-adjoint,
with
\[
 \mathsf F=P_{\rm odd}+L_\omega,\qquad
 \mathsf F(x,y)=(y,x+y).
\]
For \(\varphi^2=\varphi+1\), the map
\[
 \rho_{\mathsf F}(a+b\varphi)=L_a+L_b\mathsf F,\qquad
 \iota(a+b\varphi)=a+b\omega
\]
is respectively a field-quaternion module representation and its
coefficient embedding. The latter transports the constant coefficient
of the field-valued inner product to the displayed real metric.
\end{proposition}
\begin{proof}
Equivariance and the unit condition give \(\mathsf F x=x\omega\)
for every even \(x\). The quadratic relation then forces
\(\mathsf F(x\omega)=x+x\omega\), proving uniqueness and the symmetric
two-by-two matrix. Commutation with \(L_a\) gives
\[
 (L_a+L_b\mathsf F)(L_c+L_d\mathsf F)
 =L_{ac+bd}+L_{ad+bc+bd}\mathsf F.
\]
This is multiplication in \(\mathbb H(\mathbb Q(\varphi))\).
For scalar coefficients, the constant coefficient of
\((a+b\varphi)(c+d\varphi)\) is \(ac+bd\); summing over the four
quaternion coordinates proves the metric statement.
\end{proof}

The coefficient embedding is not a Clifford algebra homomorphism:
\(\iota(\varphi)^2=1\), while \(\iota(\varphi^2)=1+\omega\).
Reversal of volume orientation gives
\(P_{\rm odd}-L_\omega\), transported by grade-sign conjugation.
These statements do not identify the golden lift with the native real
quaternion representation.

\begin{theorem}[Explicit unit-root construction and sections]
\label{p1:thm:exroot-explicit-e876}
Let \(U\) be the supplied 120-element \(2I\) list. Set
\[
 \Phi_8=\iota(U)\cup\iota(\varphi U),\quad
 \Phi_7=\Phi_8\cap1^\perp,\quad
 \Phi_6=\Phi_7\cap(i+j+k)^\perp.
\]
Then these are unit-root systems of types \(E_8,E_7,E_6\), of ranks
\(8,7,6\) and cardinalities \(240,126,72\).
The last section is memberwise for the retained local Hurwitz-axis
family; no member or weighting is selected by this theorem.
\end{theorem}
\begin{proof}
Index the displayed orthonormal coordinates by \(0,\ldots,7\).
Attach the binary point labels
\[
 (000,100,010,110,001,011,111,101).
\]
The fourteen affine-plane supports are
\[
\begin{array}{c|c}
0123&4567\\
0147&2356\\
0156&2347\\
0245&1367\\
0267&1345\\
0346&1257\\
0357&1246
\end{array}
\]
and the root set is exactly
\[
 \{\pm e_j:0\leq j<8\}\ \cup\
 \left\{\frac12\sum_{j\in B}\epsilon_j e_j:
 B\text{ one of these planes},\ \epsilon_j=\pm1\right\}.
\]
Indeed the coordinate and rational half-unit families in \(U\), together
with their \(\varphi\)-multiples, give the coordinate roots and the
planes \(0123,4567\). In the golden family
\(\tfrac12(0,\pm1,\pm\varphi,\pm\varphi^{-1})\), use
\(\varphi^{-1}=\varphi-1\). At even-permuted role positions
\((a,b,c,d)\), the two supports are
\(\{b,d,c+4,d+4\}\) and \(\{c,d,b+4,c+4\}\).
The even role permutation \((a,b,c,d)\mapsto(a,c,d,b)\) pairs the
opposite-sign and equal-sign cases at the repeated coordinate. Thus each
of the twelve remaining supports has all sixteen signs. This proves
\(|\Phi_8|=16+14\cdot16=240\).

For crystallographic closure, let \(C\) be the binary code of affine
functions on these eight points and put
\[
 \Lambda=\{z/2:z\in\mathbb Z^8,\ z\bmod2\in C\},\qquad
 B(v,w)=2\langle v,w\rangle.
\]
The code has dimension four and weights \(0,4,8\). Distinct plane
supports intersect in zero or two points, so \(C\) is doubly even and
self-dual. It follows that \(B\) is integral and even on \(\Lambda\).
The Euclidean covolume is \(16/2^8=1/16\); the determinant in \(B\)
is \(2^8(1/16)^2=1\). The shortest nonzero vectors have Euclidean
length one and are exactly the displayed coordinate and half-plane
roots. For each such root,
\(s_\alpha(v)=v-B(v,\alpha)\alpha\) preserves \(\Lambda\) and norm.
This proves reflection closure and crystallographic integrality.
The nonorthogonality graph is connected: every pair of coordinate
indices is in a plane, and every half-root meets its coordinate roots.

The scalar section retains fourteen coordinate roots and the seven
planes not containing coordinate zero, giving
\(126=14+7\cdot16\). It spans \(1^\perp\). The further axis section
retains eight odd coordinate roots, sixteen roots on \(4567\), and
eight sign choices on each of the six other planes: on their two
indices in \(\{1,2,3\}\) the signs must cancel. Hence its count is
\(8+16+6\cdot8=72\). Its span is the four-dimensional odd space plus
the two-dimensional even imaginary plane \(x_1+x_2+x_3=0\).
Reflections in roots of either section preserve the section.

The exact Cartan certificates have the following branching-arm lengths
and determinants:
\[
\begin{array}{c|c|c|c}
 &\text{rank}&\text{arm lengths}&\det C_{\rm Cartan}\\ \hline
 E_6&6&(1,2,2)&3\\
 E_7&7&(1,2,3)&2\\
 E_8&8&(1,2,4)&1
\end{array}
\]
They identify the connected simply laced types, rather than merely
matching counts. The supplied finite certificate lists the doubled
simple roots and the integral, same-sign expansion of every root;
its independent replayer reconstructs each vector and every reflection.
Equivalently, the lattice proof and the standard simply laced
classification give the same types.
\end{proof}

\begin{remark}[Construction scope and finite evidence]
\label{p1:rem:exroot-construction-scope}
The full root lists, simple-root expansions, native/golden character
comparison and replay scripts are retained in
\path{Paper00-ExceptionalRootReadouts-v1_4.zip}, in its frozen v1.1
subpacket. The lifted nearest-shell character differs from the native
one; an equality of frames does not identify the two group actions.
The inclusions \(2T\subset2O\) and \(2T\subset2I\) do not give
\(2O\subset2I\), since \(48\) does not divide \(120\).
The root sections do not change these subgroup relations. Paper~I-B,
Theorem~\ref{p1b:thm:exroot-response-norms}, gives the precise response
metrics of this construction.
\end{remark}




\section{Coordinate/Clifford obstructions and compression defects}
\label{p1:app:coordinate-clifford-obstructions}
These results delimit the positive coordinate-completion theorem; they do not
constitute or select a completion.  The main text retains each load-bearing
scope consequence and the level-one formulas it protects.

\begin{proposition}[Terminal-degree leakage for truncated Clifford multiplication]
\label{p1:prop:terminal-degree-leakage}
Let \(0\le m\le n\) and
\[
 C_a^{(m)}:=P_{\le m}(e_a\wedge+\iota_{e_a})P_{\le m}
\]
on \(\Lambda^{\le m}\RR^n\), where \(P_{\le m}\) is the orthogonal degree
projection.  On \(\Lambda^p\RR^n\),
\[
 \sum_{a=1}^n\bigl(C_a^{(m)}\bigr)^2
 =
 \begin{cases}
 nI,&p<m,\\[2mm]
 mI,&p=m.
 \end{cases}
\]
Moreover the full Clifford anticommutator
\[
 \{C_a^{(m)},C_b^{(m)}\}=2\delta_{ab}I
\]
holds on every degree \(p<m\).  Hence, when \(m<n\), no proper degree
truncation satisfies the full Clifford sphere relation on its terminal
degree.
\end{proposition}

\begin{proof}
If \(p<m\), then \(c_a\omega\) has degrees \(p-1\) and \(p+1\), both at most
\(m\), so the intermediate projection does nothing.  Thus
\(C_a^{(m)}C_b^{(m)}\omega=P_{\le m}c_ac_b\omega\), and the ordinary Clifford
relation gives both displayed claims below the terminal degree.  If \(p=m\),
the intermediate projection removes \(e_a\wedge\omega\) and retains only
\(\iota_{e_a}\omega\); consequently
\[
 \bigl(C_a^{(m)}\bigr)^2\omega
 =e_a\wedge\iota_{e_a}\omega.
\]
The degree operator identity
\(\sum_ae_a\wedge\iota_{e_a}=pI\) on \(\Lambda^p\) gives the terminal value
\(mI\).
\end{proof}

\begin{proposition}[Second-order invariant scalarity is not a level selector]
\label{p1:prop:quadratic-coordinate-scalarity-no-selector}
Let \(\mathcal F\) be a finite multiplicity-free \(\operatorname{Spin}(4)\)-module carrying
an equivariant coordinate quadruple \(\widehat X_a\) of vector type
\((\tfrac12,\tfrac12)\).  Then
\[
 Q:=\sum_{a=1}^4\widehat X_a^2
\]
is \(\operatorname{Spin}(4)\)-invariant and hence acts by a scalar on every irreducible
summand of \(\mathcal F\).  In particular, scalarity of
\(Q|_{\mathcal K}\) does not imply that the first exterior spinor sectors are
absent and cannot select the level-one ceiling.  If exterior channels are
retained and have nonzero reduced coefficients, they may change the scalar
coefficient through an outward-and-returning two-step path without breaking
scalarity.
\end{proposition}

\begin{proof}
Equivariance of the vector quadruple and orthogonality of the vector metric
make the contracted quadratic operator \(Q\) invariant.  Multiplicity-free
isotypic components are irreducible, so Schur's lemma makes the restriction of
\(Q\) scalar on each one.  At larger truncations a first application of
\(\widehat X_a\) may enter an exterior spinor sector and the second may return
to \(\mathcal K\); invariance constrains the resulting endomorphism to remain
scalar rather than forcing the intermediate channel to vanish.
\end{proof}

\begin{proposition}[Global Clifford divisibility obstruction]
\label{p1:prop:global-clifford-divisibility-obstruction}
Over \(\CC\),
\[
 \Cl_4(\CC)\cong M_4(\CC).
\]
Consequently every finite-dimensional unital exact
\(\Cl_4(\CC)\)-module has complex dimension divisible by four.  Since
\(\dim_{\CC}\mathcal H_1=15\), the full level-one carrier admits no global
exact unital \(\Cl_4(\CC)\)-module structure.  The protected spinor module
\(\mathcal K\) is a canonical four-dimensional exact Clifford block already
present in the chosen Clifford representative, but divisibility alone does
not assert that it is the only exact Clifford block available in every
completion.
\end{proposition}

\begin{proof}
The complex Clifford algebra in four Euclidean dimensions is the simple
matrix algebra \(M_4(\CC)\).  Its finite-dimensional unital modules are direct
sums of copies of the defining four-dimensional simple module, so their
dimensions are multiples of four.  The divisibility conclusion excludes
complex dimension fifteen.
\end{proof}

\begin{proof}[Proof of Proposition~\ref{p1:prop:finite-coordinate-compression-defect}]
Insert \(P+(1-P)=1\) between the two multipliers and use
\(\sum_ax_a^2=\Rmom^2\).  With
\(D_a=(1-P)M_{x_a}P\), the defect is \(\sum_aD_a^*D_a\).  If it vanished,
then \(\operatorname{Ran}P\) would be invariant under every \(M_{x_a}\).
The four commuting self-adjoint restrictions to this finite-dimensional space
would have a common eigenvector \(f\ne0\).  The equations
\((x_a-\lambda_a)f=0\) for all \(a\) would support \(f\) on a single point of
the sphere, a null set, contradicting \(f\in L^2\) and \(f\ne0\).
\end{proof}


\subsection{Quadratic products of the truncated coordinate representative}
\label{p1:app:truncated-quadratic-content}

\begin{proof}[Verification of Remark~\ref{p1:rem:symmetric-vs-clifford-quadratic}]
Use the representative of
Definition~\ref{p1:def:clifford-coordinate-representative}, with
\(\Rmom>0\), and write
\(C_a=P_{\le2}(e_a\wedge+\iota_{e_a})P_{\le2}\).
For the retained two-form \(\omega=e_2\wedge e_3\),
\[
 C_2\omega=e_3,\qquad C_1e_3=e_1\wedge e_3,\qquad C_1\omega=0.
\]
Consequently
\[
 \{\widehat X_1^{\rm Cl},\widehat X_2^{\rm Cl}\}\omega
 =\frac{\Rmom^2}{4}e_1\wedge e_3\ne0.
\]
The off-diagonal symmetric input is trace-free.  Equivariance and
irreducibility therefore make the product map injective on the
nine-dimensional \((1,1)\) summand of \(\operatorname{Sym}^2V\).
The scalar summand has nonzero image because
\(\sum_a(\widehat X_a^{\rm Cl})^2\ne0\).
For the alternating summand, the restriction to \(\Lambda^1_\CC\) is
\[
 [C_a,C_b]e_d=2(\delta_{bd}e_a-\delta_{ad}e_b).
\]
This is the faithful six-dimensional \(\mathfrak{so}(4)\) vector action,
so the alternating product map is injective as well.  The scalar,
\((1,1)\), and self-dual/anti-self-dual alternating summands are pairwise
inequivalent.  Their independent images account for all summands of the
\(16\)-dimensional domain \(V\otimes V\).  This proves the quadratic
operator-image decomposition and its dimension.  The trace-free symmetric
component vanishes on the lower exterior degrees and on the protected
Clifford block; its nonzero action is confined to the terminal degree.
\end{proof}

\subsection{Protected Cliffordization: detailed positive-stratum proof}
\label{p1:app:protected-cliffordization-proof}
The hypotheses are those of
Theorem~\ref{p1:thm:s3one-cliffordization}: the fixed equivariant Hermitian
completion \(\mathsf C_{\rm PI}(\rho_{\rm prot})\), its multiplicity-one
intertwiners, and the normalization with \(\rho_{\rm prot}>0\).

\begin{proof}[Proof of Theorem~\ref{p1:thm:s3one-cliffordization}]
\ProofPlan{We prove the protected Clifford relation in four steps: identify the retained Clebsch--Gordan channels, show that the finite coordinate action is off-diagonal, reduce the anticommutator equivariantly to its scalar trace part, and fix the scalar by the protected algebraic normalization.}

\begin{proofsteps}
\item \textbf{Retained Clebsch--Gordan channels.}
Write
\[
    S_+=(\tfrac12,0),
    \qquad
    S_-=(0,\tfrac12),
    \qquad
    V=V_{\rm coord}=(\tfrac12,\tfrac12).
\]
The Clebsch--Gordan decompositions are
\[
    V\otimes S_+
    =
    (\tfrac12,\tfrac12)\otimes(\tfrac12,0)
    =
    (0,\tfrac12)\oplus(1,\tfrac12),
\]
and
\[
    V\otimes S_-
    =
    (\tfrac12,\tfrac12)\otimes(0,\tfrac12)
    =
    (\tfrac12,0)\oplus(\tfrac12,1).
\]
Inside the protected block only the summands
\[
    (0,\tfrac12)=S_-,
    \qquad
    (\tfrac12,0)=S_+
\]
are retained.

\item \textbf{Off-diagonal finite coordinate action.}
Hence the finite coordinate action associated to
\(v\in V\) is necessarily off-diagonal:
\[
    X(v):S_+\to S_-,
    \qquad
    X(v):S_-\to S_+.
\]
Both retained Clebsch--Gordan channels occur with multiplicity one, so after
fixing the real self-adjoint structure their common scale is the only remaining
normalization freedom.

\item \textbf{Equivariant anticommutator reduction.}
Define the symmetric anticommutator bilinear form
\[
    B(v,w):=X(v)X(w)+X(w)X(v).
\]
Since each \(X(v)\) flips chirality, \(B(v,w)\) preserves chirality.  Moreover
\(B\) is symmetric in \(v,w\) and \(\operatorname{Spin}(4)\)-equivariant.

Now
\[
    \Sym^2 V
    =
    \Sym^2(\tfrac12,\tfrac12)
    =
    (0,0)\oplus(1,1).
\]
On the other hand,
\[
    \End(S_+)
    =
    S_+\otimes S_+^*
    =
    (0,0)\oplus(1,0),
\]
and
\[
    \End(S_-)
    =
    S_-\otimes S_-^*
    =
    (0,0)\oplus(0,1).
\]
There is no \(\operatorname{Spin}(4)\)-equivariant map from \((1,1)\) to either
\((1,0)\) or \((0,1)\).  Therefore the \((1,1)\)-part of
\(\Sym^2V\) is killed by \(B\), and only the scalar trace component can remain.
A priori the two chiral scalar components may have different coefficients, so
for some scalars \(\beta_+\) and \(\beta_-\),
\[
    B(v,w)=\langle v,w\rangle(\beta_+P_L+\beta_-P_R).
\]

\item \textbf{Protected algebraic normalization.}
Let \(e_1,\ldots,e_4\) be an orthonormal coordinate basis of \(V\), and write
\[
    X_a:=X(e_a)=\widehat X_{a,\mathcal K}^{\mathsf B}.
\]
The protected normalization gives
\[
    \sum_{a=1}^{4}X_a^2=\rho_{\rm prot}^2\mathbf1_{\mathcal K}.
\]
Since
\[
    B(e_a,e_a)=2X_a^2,
\]
we get
\[
    \sum_{a=1}^{4}B(e_a,e_a)
    =
    2\sum_{a=1}^{4}X_a^2
    =
    2\rho_{\rm prot}^2(P_L+P_R).
\]
But also
\[
    \sum_{a=1}^{4}B(e_a,e_a)
    =
    4\beta_+P_L+4\beta_-P_R.
\]
Hence
\[
    \beta_+=\beta_-=
    \frac{\rho_{\rm prot}^2}{2}.
\]
Therefore
\[
    \{X_a,X_b\}
    =
    \frac{\rho_{\rm prot}^2}{2}\delta_{ab}\mathbf1_{\mathcal K},
\]
as claimed.  Dividing by \(\rho_{\rm prot}/2\) gives
\[
    \Gamma_a:=\frac{2}{\rho_{\rm prot}}X_a,
    \qquad
    \{\Gamma_a,\Gamma_b\}=2\delta_{ab}\mathbf1_{\mathcal K}.
\]
\end{proofsteps}
\end{proof}

\subsection{Boundary quotients and the failure of ideal closure}
\label{p1:app:boundary-quotient-obstruction}
Fix the additional splitting and onset convention of
Definition~\ref{p1:def:oriented-boundary-face}.  The following verifies the
vector-space-only conclusion of
Proposition~\ref{p1:prop:boundary-quotient-not-algebra}.

\begin{proof}[Proof of Proposition~\ref{p1:prop:boundary-quotient-not-algebra}]
Choose nonzero \(\eta_\sigma\in N_\sigma\) and
\(0\ne e\in F_\sigma\).  The degree-three element
\(\eta_\sigma\odot e^2\) belongs to \(\mathcal L_3^\sigma\).  Multiplying it
by \(\eta_\sigma\in X\) gives
\[
    \eta_\sigma^2\odot e^2\in\Sym^4(X).
\]
Under the direct-sum decomposition \(X=N_\sigma\oplus F_\sigma\), this
element has normal degree two, whereas
\[
    \mathcal L_4^\sigma
    =N_\sigma\odot\Sym^3(F_\sigma)
\]
has normal degree one.  Hence the product does not lie in
\(\mathcal L_4^\sigma\), so \(\mathcal L^\sigma\) is not an ideal and no
quotient-algebra multiplication descends.
\end{proof}

\section{Supplementary Hurwitz finite-cell calculations}
\label{p1:app:hurwitz-supplementary}
The main finite-cell chain already states the design defect, Bose--Mesner rank
split, first invariant degree, and triality decomposition.  This appendix
records their explicit coordinate and projector calculations without adding a
new hypothesis or finite-cell datum.

\subsection{Natural left-action obstruction for the full carrier}
\label{p1:app:left-2t-carrier-obstruction}
Here \(2T\) acts through the \(SU(2)_L\) factor on \(\mathcal H_1\) and by
left translation on the scalar function space \(\CC^{2T}\), exactly as in
Proposition~\ref{p1:prop:h1-e15-left-2t-obstruction}.

\begin{proof}[Proof of Proposition~\ref{p1:prop:h1-e15-left-2t-obstruction}]
Under \(SU(2)_L\), the sectors of \(\mathcal H_1\) with trivial left factor
are
\[
    (0,0)\oplus(0,\tfrac12)\oplus(0,1),
\]
of total dimension \(1+2+3=6\).  The remaining left factors are the defining
two-dimensional representation and the three-dimensional tetrahedral
rotation representation; neither has a nonzero \(2T\)-fixed vector.
Thus \(\dim\mathcal H_1^{2T}=6\).

The fixed subspace of the left regular representation on
\(\CC^{2T}\) is the one-dimensional constant-function line.  The constant
vector is the valency-\(8\) eigenvector of \(A_+\), hence lies in
\(\mathcal E_{15}^{\rm BM}\), while the orthogonal zero-eigenspace
\(\mathcal E_9^{\rm BM}\) has no fixed vector.  Therefore
\(\dim(\mathcal E_{15}^{\rm BM})^{2T}=1\).

If \(F\) is equivariant, it maps the six-dimensional fixed subspace
\(\mathcal H_1^{2T}\) into the one-dimensional fixed subspace of
\(\CC^{2T}\).  Rank--nullity on that restriction gives a kernel of dimension
at least five, proving the stronger no-injection statement.
\end{proof}

\subsection{Cubature, projector, and triality calculations}

\begin{proposition}[Degree-six Hurwitz cubature defect]
\label{p1:prop:math-hurwitz-cubature-six-defect}
Let \(2T\subset S^3\subset\RR^4\) be the Hurwitz \(24\)-cell support in the
\(S^3_1\) cell.  Then \(2T\) is a spherical \(5\)-design, and in a Hurwitz
coordinate direction \(e_1\) its first scalar moment defect is
\[
  \delta_6:=
  \sum_{q\in 2T}\langle e_1,q\rangle^6
  -24\,\mathbb E_{S^3}(x_1^6)
  =\frac38.
\]
More generally, define
\[
 \Phi_6(x):=\sum_{q\in2T}\langle q,x\rangle^6
\]
and
\[
 H_6(x)
 :=\sum_{i=1}^4x_i^6
 -5\sum_{i\ne j}x_i^4x_j^2
 +30\sum_{i<j<k}x_i^2x_j^2x_k^2.
\]
Then \(H_6\) is harmonic, nonradial, and invariant under the full
\(24\)-cell symmetry, and
\[
 \Phi_6(x)-\frac{15}{8}|x|^6=\frac38H_6(x).
\]
Thus \(H_6\) is an explicit harmonic representative of the first
nonconstant full-cell invariant degree.
\end{proposition}

\begin{proof}
For the standard Hurwitz units
\[
2T=\{\pm1,\pm i,\pm j,\pm k\}
\cup
\left\{\frac{\pm1\pm i\pm j\pm k}{2}\right\},
\]
the real coordinate values are \(\pm1\), six zeros, and sixteen entries of
absolute value \(1/2\).  Hence
\[
  \sum_{q\in 2T}x_1(q)^6
  =2+16\left(\frac12\right)^6
  =\frac94.
\]
For the unit sphere \(S^3\subset\RR^4\),
\[
  \mathbb E_{S^3}(x_1^6)
  =\frac{1\cdot3\cdot5}{4\cdot6\cdot8}
  =\frac5{64},
\]
so \(24\mathbb E_{S^3}(x_1^6)=15/8\), and the defect is \(3/8\).
Exactness through degree five is the spherical-design property proved in the
main text.

Expanding the orbit sum over the eight quaternion units and sixteen
half-integer units gives
\[
 \Phi_6(x)=\frac{15}{8}|x|^6+\frac38H_6(x).
\]
Direct differentiation gives
\[
 \Delta\sum_i x_i^6=30\sum_i x_i^4,
\]
\[
 \Delta\sum_{i\ne j}x_i^4x_j^2
 =6\sum_i x_i^4+12\sum_{i\ne j}x_i^2x_j^2,
 \qquad
 \Delta\sum_{i<j<k}x_i^2x_j^2x_k^2
 =2\sum_{i\ne j}x_i^2x_j^2,
\]
so \(\Delta H_6=0\).  Invariance follows from the orbit-sum identity, and
\(H_6(e_1)=1\) while \(H_6(e_1+e_2)=-8\), proving nonradiality.
\end{proof}

\begin{corollary}[Explicit \(15+9\) Bose--Mesner projectors]
\label{p1:cor:math-twentyfourcell-explicit-projectors}
With \(A_\pi\) denoting the antipodal relation matrix, the orthogonal
projectors onto \(\mathcal E_9^{\rm BM}=\ker A_+\) and
\(\mathcal E_{15}^{\rm BM}=(\mathcal E_9^{\rm BM})^\perp\) are
\[
 P_9^{\rm BM}=\frac18(3I-A_0+3A_\pi),
 \qquad
 P_{15}^{\rm BM}=\frac18(5I+A_0-3A_\pi).
\]
They satisfy
\[
\begin{aligned}
(P_r^{\rm BM})^2&=P_r^{\rm BM},&
P_9^{\rm BM}P_{15}^{\rm BM}&=0,\\
\rank P_9^{\rm BM}&=9,&
\rank P_{15}^{\rm BM}&=15,\\
 A_+P_9^{\rm BM}&=0,&
 r&\in\{9,15\}.
\end{aligned}
\]
Equivalently,
\[
 P_9^{\rm BM}
 =\frac{(A_+-8I)(A_+-4I)(A_++2I)(A_++4I)}{256}.
\]
\end{corollary}

\begin{proof}
Identify \(A_\pi\) with the antipodal matrix in the multiplication table of
Theorem~\ref{p1:thm:s3one-numerical-data}\textup{(iv)}.  Substitution of
\[
A_0^2=6I+4A_0+6A_\pi,\qquad
A_0A_\pi=A_0,\qquad A_\pi^2=I
\]
shows directly that the displayed matrices are complementary idempotents.
The same table gives \(A_+A_0=3A_++3A_-\) and \(A_+A_\pi=A_-\), hence
\(A_+P_9^{\rm BM}=0\).  Their traces are \(9\) and \(15\), proving the
ranks.  The polynomial formula is the spectral projector onto the zero
eigenspace of \(A_+\), whose other eigenvalues are \(8,4,-2,-4\).
\end{proof}

\begin{proposition}[First nonconstant full-cell invariant degree]
\label{p1:prop:math-f4-sextic-obstruction}
For the unmarked Hurwitz \(24\)-cell with full reflection symmetry
\(W(F_4)\), the invariant polynomial ring has generator degrees
\[
 2,\;6,\;8,\;12.
\]
On \(S^3\), the degree-two invariant is constant and radial.  Hence the first
nonconstant scalar invariant compatible with the full unmarked \(24\)-cell
symmetry has degree \(6\).  After selecting a vertex, oriented diagonal, or
trimer orientation, the symmetry group is smaller; this proposition concerns
the full unmarked symmetry only.
\end{proposition}

\begin{proof}
The Chevalley theorem for finite reflection groups gives the displayed
fundamental degrees.  The degree-two generator is \(r^2\), which restricts to
\(1\) on the unit sphere.  Thus the first nonconstant spherical invariant is
sextic.  On \(S^3(\Rmom)\), unnormalized homogeneous degree-six coordinate
moments scale by \(\Rmom^6\).
\end{proof}

\begin{proposition}[Triality-frame octets and the \(K_{4,4,4}\) quotient]
\label{p1:prop:math-triality-octets-k444}
Use the explicit orthogonal map and \(\sqrt2\) rescaling of
Proposition~\ref{p1:prop:hurwitz-d4-coordinate-map}, and write
\[
 \Phi(D_4)=\{\pm e_i\pm e_j:1\le i<j\le4\}.
\]
After quotienting by antipodes, the twelve axes split into three tetrads
\[
\begin{aligned}
\mathcal T_A&=\{[e_1\pm e_2],[e_3\pm e_4]\},\\
\mathcal T_B&=\{[e_1\pm e_3],[e_2\pm e_4]\},\\
\mathcal T_C&=\{[e_1\pm e_4],[e_2\pm e_3]\}.
\end{aligned}
\]
Each tetrad consists of four mutually orthogonal axes.  Lifting back to
vertices gives three octets and hence \(24=8+8+8\).  The antipodal quotient
adjacency graph is \(K_{4,4,4}\), and \(D_4\) triality permutes the tetrads.
\end{proposition}

\begin{proof}
The three tetrads are the three perfect matchings of the coordinate labels
\(1,2,3,4\).  Within a tetrad the axes are pairwise orthogonal.  Between
distinct tetrads, representatives have normalized inner product \(\pm1/2\),
so quotient adjacency connects every axis in one tetrad to every axis in the
other two and no axes within a tetrad.  This is \(K_{4,4,4}\), and the outer
automorphism triality permutes the three matchings.
\end{proof}

\begin{theorem}[Exact chain complex of the Hurwitz \(24\)-cell boundary]
\label{p1:thm:v235-hurwitz-chain-complex}
Fix the short-root Hurwitz realization and an ordered triality-octet
partition \(V_0\sqcup V_1\sqcup V_2\) from
Proposition~\ref{p1:prop:math-triality-octets-k444}.  The boundary
of the \(24\)-cell has \(24\) vertices, \(96\) edges, \(96\)
triangular faces and \(24\) octahedral three-cells.  Its real cellular
boundary maps form
\[
 C_3\cong\mathbb R^{24}\xrightarrow{\partial_3}
 C_2\cong\mathbb R^{96}\xrightarrow{B=\partial_2}
 C_1\cong\mathbb R^{96}\xrightarrow{\partial_1}
 C_0\cong\mathbb R^{24},
 \qquad \partial_1B=B\partial_3=0.
\]
Orient each edge cyclically \(V_0\to V_1\to V_2\to V_0\) and
its incident triangular faces compatibly.  The nonzero entries
of \(B\) then all equal \(+1\).  The exact ranks are
\[
 \operatorname{rank}\partial_1=23,\qquad
 \operatorname{rank}B=73,\qquad
 \operatorname{rank}\partial_3=23,
 \quad\ker B=\operatorname{im}\partial_3.
\]
In particular, for any formal face vector \(f\in C_2\), the edge
sum \(D=Bf\) satisfies the \(23\) independent vertex-return
identities \(\partial_1D=0\), and precisely \(23\) independent
face-boundary directions are invisible to \(B\).
\end{theorem}
\begin{proof}
The cell counts follow by the Hurwitz coordinate incidence
\(\langle v,w\rangle=1/2\) on the unit short-root shell.
Every face joins one vertex of each triality octet, so the cyclic
edge orientation makes each oriented triangular boundary a sum of
three positive edges.  The boundary of a convex \(4\)-polytope is
homeomorphic to \(S^3\), hence \(H_1(S^3;\mathbb R)=
H_2(S^3;\mathbb R)=0\) and \(H_3\cong\mathbb R\).
Connectivity gives \(\operatorname{rank}\partial_1=24-1=23\);
exactness at \(C_1\) gives \(\operatorname{rank}B=96-23=73\),
and exactness at \(C_2\) gives \(\operatorname{rank}\partial_3=
96-73=23\).  The remaining one-dimensional kernel of
\(\partial_3\) is the oriented whole boundary cycle.
\end{proof}

\begin{corollary}[Dual three-octet cell incidence without physical owner selection]
\label{p1:cor:v235-dual-cell-triality}
The \(24\) octahedral three-cell normals, in the dual
\(D_4\) coordinates, are the \(24\) vectors
\(\pm e_i\pm e_j\) with \(1\le i<j\le4\).
The three perfect matchings of the four coordinate axes partition
them into three octets of eight cells.  Every Hurwitz edge lies in
exactly one cell of each octet.  This is a dual incidence
classification; it neither identifies the cell octets with color
representations nor supplies a canonical physical triality-owner map
from the signed Spin\((8)\) carrier of
Theorem~\ref{p1:thm:v567-signed-triality-channel-family}.
\end{corollary}
\begin{proof}
Each matching contains two disjoint coordinate pairs and four
sign choices per pair.  The unit-shell edge incidence gives three
adjacent octahedral normals per edge, belonging to the three
separate matchings.  This proves both the partition and incidence
assertions; no cross-carrier intertwiner is used.
\end{proof}

\section{Critical boundary near-hits and operator-domain distinctions}
\label{p1:sec:v167-critical-boundary-near-hits}

\begin{theorem}[Hardy constant mode and the zero Dirichlet-to-Neumann eigenvalue]
\label{p1:thm:v167-hardy-constant-dtn-zero}
Identify the complex ball \(B^2(\Rmom)\subset\CC^2\) with the real ball
\(B^4(\Rmom)\), whose boundary is \(S^3(\Rmom)\).  Let
\(\mathcal N_0\) be the scalar Dirichlet-to-Neumann operator for the Euclidean
Laplacian on this ball, with outward-normal convention
\[
 \mathcal N_0f:=\left.\partial_n\mathcal E_0f\right|_{S^3(\Rmom)},
 \qquad
 \Delta\mathcal E_0f=0,
 \qquad
 \left.\mathcal E_0f\right|_{S^3(\Rmom)}=f.
\]
Then
\[
 \boxed{\mathcal N_0\mathbf1=0.}
\]
Moreover \(\mathbf1\) belongs to the Hardy--Szeg\H{o} range and, as a real
CR-pluriharmonic function, belongs to the kernel of the standard CR Paneitz
operator.
\end{theorem}

\begin{proof}
The constant extension \(\mathcal E_0\mathbf1=\mathbf1\) is harmonic, is the
unique harmonic extension of its boundary value, and has zero normal
derivative.  The constant holomorphic polynomial on \(B^2\) has boundary
trace \(\mathbf1\), so Theorem~\ref{p1:thm:hardy-szego-projection} places it
in the Hardy--Szeg\H{o} range.  It is the real part of a CR-holomorphic
constant, and hence lies in the CR Paneitz kernel by the standard-sphere
kernel theorem of Section~\ref{p1:subsec:cr-paneitz}.
\end{proof}

\begin{corollary}[The scalar zero mode and the inverse contact operator have different domains]
\label{p1:cor:v167-critical-mode-owner-firewall}
The scalar constant lies in the boundary operator kernels specified above.
The inverse contact operator $\mathcal A_H=8\Delta_{b,0}^{-2}$ acts on the
orthogonal complement of that kernel. Its domain therefore omits the scalar
zero mode; it cannot represent that mode as one of its eigenvectors.
\end{corollary}
\begin{proof}
The inverse is defined after restriction to the complement of the kernel.
The preceding kernel identities give the stated domain distinction.
\end{proof}

\section{Axial boundary trace and determinant response}
\label{p1:sec:v557-boundary-process-package}

The preceding critical-boundary near-hit isolates a scalar zero mode but does
not yet package the corresponding oriented boundary geometry and its natural
scalar descendants.  The following results complete that mathematical layer.
Their domains are the round sphere, its great-sphere boundary, and the
displayed boundary tangent and determinant bundles.

\begin{definition}[Axial great-sphere boundary]
\label{p1:def:v557-axial-great-sphere-boundary}
Identify the unit-radius model of $S^3$ with the unit quaternions
$q=a+v\in\mathbb R\oplus\operatorname{Im}\mathbb H$.  Fix a unit
$u\in\operatorname{Im}\mathbb H$ and define
\[
 h_u(q):=\langle v,u\rangle,
 \qquad
 \Sigma_u:=h_u^{-1}(0)\subset S^3.
\]
Let
\[
 H_u^\pm:=\{q\in S^3:\ \pm h_u(q)\ge0\}
\]
be the two closed hemispheres bounded by $\Sigma_u$.
\end{definition}

\begin{theorem}[Regular axial boundary and canonical normal]
\label{p1:thm:v557-axial-boundary-normal}
For every unit $u\in\operatorname{Im}\mathbb H$, the set $\Sigma_u$ is a
round totally geodesic $S^2\subset S^3$.  Along $\Sigma_u$,
\[
 \nabla_{S^3}h_u=(0,u),
 \qquad
 \|\nabla_{S^3}h_u\|=1.
\]
Hence $\Sigma_u$ has a canonical oriented unit normal
\[
 \nu_u:=(0,u),
\]
with the sign fixed by the choice of $h_u>0$ side.  The orientation of $S^3$
and $\nu_u$ induce an orientation and area form $\mathrm dA_u$ on $\Sigma_u$.
\end{theorem}

\begin{proof}
In the ambient Euclidean space $\mathbb R\oplus\operatorname{Im}\mathbb H$,
$h_u$ is the linear functional with constant gradient $(0,u)$.  Its spherical
gradient is the tangential projection
\[
 \nabla_{S^3}h_u=(0,u)-h_u(q)q.
\]
On $h_u=0$ this equals $(0,u)$ and has unit norm.  Thus zero is a regular
value.  Since $\Sigma_u$ is the intersection of the unit sphere with the
linear hyperplane $\langle v,u\rangle=0$, it is a great $2$-sphere and hence
totally geodesic.  The orientation statement is standard for an oriented
hypersurface.
\end{proof}

\begin{theorem}[Axial hemisphere constant mode and zero DtN eigenvalue]
\label{p1:thm:v557-axial-hemisphere-dtn-zero}
Let $\mathcal N_u^\pm$ be the scalar Dirichlet-to-Neumann operator for the
Laplace--Beltrami operator on the round hemisphere $H_u^\pm$, with outward
normal convention.  Then
\[
 \boxed{\mathcal N_u^\pm\mathbf1=0.}
\]
Thus the scalar constant line on the generated axial boundary is a canonical
critical boundary mode.
\end{theorem}

\begin{proof}
The constant function on $H_u^\pm$ is harmonic and is the unique harmonic
extension of its constant boundary value.  Its normal derivative vanishes.
\end{proof}

\begin{definition}[Boundary tangent determinant scalarization]
\label{p1:def:v557-boundary-determinant-scalarization}
\leavevmode\par
Let $A\in\operatorname{End}(T\Sigma_u)$ be an orientation-preserving positive
endomorphism.  Its top-exterior action is
\[
 \Lambda^2A^*:\Lambda^2T^*\Sigma_u\longrightarrow\Lambda^2T^*\Sigma_u.
\]
Using the boundary Hodge isomorphism
\[
 *_\Sigma:\Lambda^2T^*\Sigma_u\stackrel{\sim}{\longrightarrow}
 \Lambda^0T^*\Sigma_u,
\]
define the scalar descendant
\[
 \mathfrak d_\Sigma(A)
 :=*_\Sigma\,\Lambda^2A^*\,*_\Sigma^{-1}.
\]
\end{definition}

\begin{theorem}[Canonical determinant-to-scalar descent]
\label{p1:thm:v557-determinant-scalar-descent}
For every positive orientation-preserving $A\in\operatorname{End}(T\Sigma_u)$,
\[
 \boxed{\mathfrak d_\Sigma(A)=(\det A)I_{\Lambda^0}.}
\]
In particular, for the scalar tangent contraction
\[
 A_\ell=e^{-\ell/2}I_{T\Sigma_u},\qquad \ell\ge0,
\]
one has
\[
 \boxed{\mathfrak d_\Sigma(A_\ell)=e^{-\ell}I_{\Lambda^0}.}
\]
Moreover
\[
 \mathfrak d_\Sigma(A_{\ell_1+\ell_2})
 =\mathfrak d_\Sigma(A_{\ell_1})\,
  \mathfrak d_\Sigma(A_{\ell_2}).
\]
Hence an additive tangent-amplitude depth has a canonical multiplicative
scalar boundary descendant.
\end{theorem}

\begin{proof}
On the oriented two-plane $T_q\Sigma_u$, the top exterior power of $A^*$ acts
on the one-dimensional area line by multiplication with $\det A$.  Conjugation
by the Hodge isomorphism preserves this scalar.  For
$A_\ell=e^{-\ell/2}I$ the determinant is $e^{-\ell}$, and multiplicativity
follows from either determinants or $A_{\ell_1+\ell_2}=A_{\ell_1}A_{\ell_2}$.
\end{proof}

\begin{corollary}[Generated scalar semigroup on the boundary constant line]
\label{p1:cor:v557-boundary-scalar-semigroup}
Restricting the preceding scalar descendant to the constant boundary line
$\mathbb R\mathbf1\subset C^\infty(\Sigma_u)$ gives
\[
 \rho_\ell:=e^{-\ell}I_{\mathbb R\mathbf1},
\]
with
\[
 \rho_0=I,
 \qquad
 \rho_{\ell_1+\ell_2}=\rho_{\ell_1}\rho_{\ell_2},
 \qquad
 0<\rho_\ell\le I.
\]
This scalar semigroup is canonically descended from the oriented tangent
amplitude; it is not introduced as an independent scalar profile.
\end{corollary}

\begin{remark}[Domain of the determinant semigroup]
\label{p1:rem:v557-boundary-process-owner-firewall}
The scalar semigroup is the determinant of the specified tangent
amplitude semigroup.  Its generator on the constant line is \(-I\),
whereas the hemisphere Dirichlet-to-Neumann operator annihilates that line.
These are distinct operators with distinct definitions, as in
Corollary~\ref{p1:cor:v167-critical-mode-owner-firewall}.
\end{remark}

\section{Higher Hopf--Hurwitz invariant lines}
\label{p1:app:higher-hopf-hurwitz-invariant-lines}

These alternatives retain absolute finite-group invariant tensors, whereas
the main quadratic theorem retains a moving joint tuple.  Minimal degrees
are therefore indexed by the input signature.

\begin{proposition}[Invariant pencil and polar transport]
\label{p1:prop:higher-invariant-pencil}
In defining spin coordinates \(x,y\), choose
\[
 f_6=xy(x^4-y^4),\quad f_8=x^8+14x^4y^4+y^8,\quad
 g_{12}=x^{12}-33x^8y^4-33x^4y^8+y^{12},\quad u_{12}=f_6^2.
\]
The invariant spaces \(I_n=(\operatorname{Sym}^nW_R^*)^{2T}\), for
\(0\leq n\leq12\), have dimensions
\[
 (1,0,0,0,0,0,1,0,1,0,0,0,2).
\]
The degree-twelve space has basis \(u_{12},g_{12}\), with
\[
 f_8^3=g_{12}^2+108f_6^4,\qquad
 \|u_{12}\|^2=8/231,\quad \|g_{12}\|^2=32/5,\quad
 \langle u_{12},g_{12}\rangle=0.
\]
Evaluation at a spin line \(L\) gives a nowhere-zero map
\(\epsilon_L:L^{\otimes12}\to I_{12}^*\), with image line \(A_{12}(L)\).
Its polar part
\[
 U_{12,L}=\epsilon_L(\epsilon_L^\dagger\epsilon_L)^{-1/2}:
 L^{\otimes12}\longrightarrow A_{12}(L)
\]
is a natural unitary isomorphism.  Degree twelve is the least positive
degree of a base-point-free system of absolute \(2T\)-invariants.
\end{proposition}
\begin{proof}
Expansion of the Molien series in \eqref{p1:eq:molien-2T} gives the
displayed dimensions.  Substitution of defining matrices for \(i\) and
\((1+i+j+k)/2\), which generate \(2T\), verifies invariance of the three
forms.  Multiplying the polynomials gives the relation.  With the induced
symmetric-tensor metric,
\(\|\sum a_kx^{n-k}y^k\|^2=\sum|a_k|^2/\binom nk\); this proves the
norms and orthogonality, and also
\(\|f_6\|^2=1/3,\ \|f_8\|^2=24/5\).
If \(u_{12}=0\), either \(x=0\), \(y=0\), or \(x^4=y^4\).
The first two cases leave \(g_{12}\ne0\) at a projective point; the last
gives \(g_{12}=-64x^{12}\ne0\).  Thus evaluation never vanishes.
The norm squared of evaluation is
\[
 \frac{|u_{12}(z)|^2/(8/231)+|g_{12}(z)|^2/(32/5)}{\|z\|^{24}}>0.
\]
It is bounded away from zero on the compact projective line, so the
inverse square root and polar map are smooth and natural.  Below degree
twelve the only positive invariant spaces are the one-dimensional spaces
in degrees six and eight.  Each nonconstant binary form has a projective
zero, proving minimality for this signature.
\end{proof}

\begin{proposition}[Compressed readouts and connection distinction]
\label{p1:prop:invariant-compression-connection}
For the invariant projector \(\Pi_n\) and restriction
\(R_n:\operatorname{Sym}^nW_R^*\to (L^*)^{\otimes n}\),
\[
 R_n\Pi_nR_n^\dagger=k_n(L)\operatorname{Id},\qquad k_6+k_8>0.
\]
These scalar readouts discard the phase of the invariant tensors.
In compatible fixed spin and real axial frames, the non-common-zero
property gives the nonzero family
\((f_6|_L)^2\otimes\lambda_h^{12}\) and
\((f_8|_L)^2\otimes\lambda_h^{16}\), respectively in
\((L\otimes\ell)^{*12}\) and \((L\otimes\ell)^{*16}\).
These are the displayed frame-qualified sections.
The normalizer element \((1+i)/\sqrt2\) acts on the degree-twelve basis
as \((u_{12},g_{12})\mapsto(u_{12},-g_{12})\).
Over \(\mathbb{CP}^1\), \(A_{12}\simeq L_{\mathrm{taut}}^{\otimes12}\)
has first Chern class \(-12\).  Its pullback to the full unit spinor
\(S^3\) has the evaluation frame of weight twelve.
The polar isomorphism does not identify two independently prescribed
projection connections.
\end{proposition}
\begin{proof}
The compressed endomorphism acts on a line and is positive; the
non-common-zero property of \(f_6,f_8\) proves strict positivity of their
sum.  The normalizer action is obtained by substituting
\((x,y)\mapsto(e^{i\pi/4}x,e^{-i\pi/4}y)\).
The Chern class follows from the everywhere injective evaluation map.
In the affine coordinate \(x\), \(y=1\), differentiation gives
\[
 u_{12}'g_{12}-u_{12}g_{12}'=-2f_6 f_8^2.
\]
Thus the projective evaluation map has degree twelve, with six simple
ramification points over \(f_6=0\), and eight points of ramification
order two over \(f_8=0\).  The induced projective curvature vanishes at
these points, whereas the standard input Hopf curvature is nonzero.
Consequently equality of the inherited projection connections is false
in general.  The connection
\(U_{12}\nabla_{L^{12}}U_{12}^{-1}\) is instead the explicitly transported
connection.  Compression to \(k_n\) cannot prove that the full
\(\Pi_n\) is fixed by a Hopf Cartan transformation.
\end{proof}

The geometric map \(U_{12}\) needs a further endpoint pairing
\(\kappa:A_{12}\otimes\ell^{12}\to\CC\) to give a pointing of
\((L\otimes\ell)^{*12}\).  Holding an additional image frame fixed
adds a second retention condition.  Neither this appendix nor the
compressed norm constructs that endpoint or fixes its phase.
For \(h\in2T\), the actual finite action obeys
\(\epsilon_{\rho(h)L}\rho(h)^{12}=\epsilon_L\) and
\(U_{12,\rho(h)L}\rho(h)^{12}=U_{12,L}\), as follows by evaluating
invariant polynomials and then taking the polar part.



\section{Axial and Hopf transport alternatives}
\label{p1:app:axial-hopf-transport-alternatives}

\begin{proposition}[Naturality of the finite axial transport]
\label{p1:prop:axial-invariant-transport}
For \(h\in2T\), naturality gives
\[
 \epsilon_{\rho(h)L}\rho(h)^{\otimes12}=\epsilon_L.
\]
The image lines are the same subspace of \(I_{12}^*\).
For \(\ell=\CC h\), define
\[
 \Psi_{h,L}((th)\otimes z)=t\rho(h)z.
\]  The induced map
\[
 \Theta_{h,L}=U_{12,\rho(h)L}\Psi_{h,L}^{\otimes12}
                 (1\otimes U_{12,L}^{-1}):
 \ell^{12}\otimes A_{12}(L)\longrightarrow A_{12}(L)
\]
satisfies
\(\Theta_{h,L}((th)^{12}\otimes a)=t^{12}a\).
Its trace is a covector in \(\ell^{*12}\); no Hopf dual remains.
\end{proposition}
\begin{proof}
Every element of \(I_{12}\) is \(2T\)-invariant, proving the evaluation
identity.  Unitarity gives the same identity for its polar part.
Substitution into the definition of \(\Theta\) cancels the two \(U\)'s.
Taking the trace on its one-dimensional endomorphism factor leaves
exactly the coefficient covector \(\lambda_h^{12}\).
\end{proof}

\begin{proposition}[Opposite Hopf pairing and moving-axis twists]
\label{p1:prop:opposite-hopf-moving-axis}
Put \(Jq=iq\) on real quaternions and
\(P_\pm=(1\mp i_{\CC}J)/2\) on their complexification.
For a real unit \(q\), \(z_\pm(q)=\sqrt2P_\pm q\) are unit vectors,
\[
 z_\pm(e^{\theta J}q)=e^{\pm i\theta}z_\pm(q),\qquad
 b(z_+,z_-)=1,\quad b(z_\pm,z_\pm)=0,
\]
where \(b\) is the complex-bilinear extension of the real metric.
Both projections of a real nonzero axis are nonzero and reconstruct
that axis.  An actual quaternionic Hopf rotation
\(e^{\theta J}h=\cos\theta\,h+\sin\theta\,Jh\) therefore leaves
\(\CC h\) at generic \(\theta\); it is not merely a coefficient
phase in that fixed complex line.
Hence \(H_-\simeq H_+^*\), and
\[
 \kappa_H=b^{\otimes12}(U_{12}^{-1}\otimes1):
 A_{12}\otimes H_-^{12}\longrightarrow\CC
\]
is a nonzero geometric pairing.
For the moving axis \(E_{\mathrm{ax},q}=\CC q\), consider the lift
\[
 T_k(\theta)(cq)=e^{ik\theta}c\,e^{\theta J}q.
\]
Projection to \(H_-\) intertwines its action precisely for \(k=1\);
projection to \(H_+\) intertwines precisely for \(k=-1\), when the
respective projected target is transported without an extra scalar
character.  With the convention \(c_1(E_{\mathrm{ax},k}/U(1))=k\),
\[
 c_1(A_{12}\otimes E_{\mathrm{ax},k}^{12})=12(k-1).
\]
\end{proposition}
\begin{proof}
The spectral projections of the orthogonal complex structure \(J\)
give the character laws and the three metric identities directly.
On the negative projected line the lifted source supplies the factor
\(e^{i(k-1)\theta}\), which must be one for every \(\theta\); on the
positive line it supplies \(e^{i(k+1)\theta}\).
The Chern-class statement follows from the associated characters and
\(c_1(A_{12})=-12\).
\end{proof}

\begin{remark}[What the smooth pairing retains]
On full \(S^3\), the frame rule
\(\kappa_P(a_{12}(q)\otimes q^{12})=1\) is smooth.  If the ambient
image line is held fixed, its character is
\(e^{12i(k-1)\theta}\).  The \(k=1\) action is not an endomorphism
preserving the untwisted real-axis self-reality; conjugation transports
it to the opposite Hopf role.  With the standard projector conventions,
the real-axis connection is zero, the opposite Hopf connections are
\(\pm\eta\), and the image connection has the form
\(B_A=12\eta+\delta\), with \(\delta\) horizontal.  A connection making
the displayed twelfth-power pairing parallel must satisfy
\(B_\ell=-B_A/12\).  It need not be the inherited real-axis connection.
These are comparisons of explicitly specified bundles, actions and
connections; an unrelated scalar-circle action does not cancel their
momentum Hopf character.
\end{remark}



\section{Joint-axis comparison readouts}
\label{p1:app:joint-axis-comparison-readouts}

\begin{proposition}[Finite symmetry versus continuous symmetry]
\label{p1:prop:finite-quadratic-covariant-space}
For quaternion coordinates \(q=(r,x,y,z)\), the \(2T\)-equivariant
quadratic maps to the conjugation vector representation have basis
\[
 P_1=r(x,y,z),\qquad P_2=(yz,zx,xy).
\]
Full \(SU(2)\) conjugation retains only \(P_1\).  At a half-integer
Hurwitz vertex with signs \(\epsilon_i\),
\(P_2=(\epsilon_1\epsilon_2\epsilon_3)P_1\).
\end{proposition}
\begin{proof}
The three axial half-turns force the first component to be a
combination of \(rx\) and \(yz\), and analogously for the other two.
The tetrahedral cyclic permutation forces the two coefficients to be
common to all three components.  These generators suffice for \(2T\)
conjugation.  A general spatial rotation does not preserve \(P_2\);
the scalar-times-vector summand \(r(x,y,z)\) is the only vector in the
quadratic \(SO(3)\) decomposition.  Substitution of \(r=1/2\) and
\((x,y,z)=\epsilon/2\) gives the last identity.
\end{proof}

\begin{theorem}[Angular moments and finite comparison heat]
\label{p1:thm:joint-axis-moments-heat}
For \(|n|=1\) and \(\epsilon\in\{\pm1\}^3\), let
\[
 f_\epsilon(n)=\frac{3-(\epsilon\cdot n)^2}{4}.
\]
The eight quadratic comparisons obey
\[
 \sum_\epsilon f_\epsilon=4,\qquad
 Q(n):=\sum_\epsilon f_\epsilon^2=3-\sum_i n_i^4.
\]
The minima \(Q=2\) are the six coordinate axes, the maxima \(Q=8/3\)
are the eight body diagonals, and the twelve edge directions are
saddles.  If \(D_\epsilon\) is the self-adjoint two-line comparison
block with \(D_\epsilon^2=f_\epsilon\operatorname{Id}_2\), then
\[
 \operatorname{Tr}D^2=8,\quad \operatorname{Tr}D^4=2Q(n),\quad
 \operatorname{Tr}e^{-\tau D^2}
     =2\sum_\epsilon e^{-\tau f_\epsilon}
     \geq16e^{-\tau/2}\quad(\tau>0).
\]
Equality in the heat inequality holds exactly on the six axes.
\end{theorem}
\begin{proof}
Sign summation gives
\(\sum(\epsilon\cdot n)^2=8\) and
\(\sum(\epsilon\cdot n)^4=24-16\sum n_i^4\); expanding \(f^2\)
gives the moments.  Lagrange multipliers imply that all nonzero
coordinates at a critical point have equal magnitude.  One, two and
three nonzero coordinates give respectively the axis, edge and
diagonal classes; restriction of the quadratic variation to the
tangent plane gives the stated types.  The trace formulas follow
blockwise.  Strict convexity of \(e^{-\tau x}\) gives Jensen's
inequality, with equality only when all \(f_\epsilon=1/2\).
Expanding their common value in the independent sign products forces
\(n_in_j=0\) for \(i\ne j\), hence the axis condition.
\end{proof}

The six axes form one octahedral-normalizer orbit.  Retaining the
marked vertex \((1+i+j+k)/2\) leaves its \(C_3\) stabilizer, which
splits them into the positive and negative triples. The norms and heat
minima are invariants of these finite geometric data.

\begin{proposition}[Fixed-axis section and joint-orbit section differ]
\label{p1:prop:fixed-axis-quadratic-zeros}
Fix \(h=(1+i+j+k)/2\) and vary only the spin line over
\(\mathbb{CP}^1\).  The quadratic
\(a s_0^2+2b s_0s_1+c s_1^2\) from the comparison matrix has
\(\det S=3/4\) and discriminant \(-3\), so it has two simple
projective zeros, at the Hopf axes parallel and antiparallel to
\(\operatorname{Im}h\).  Pulling this particular section to full
\(S^3\) does not remove those zeros.
\end{proposition}
\begin{proof}
Insert \(r=x=y=z=1/2\) into the quadratic matrix
\eqref{p1:eq:joint-quadratic-matrix}.  Its determinant is \(3/4\),
and \((2b)^2-4ac=-4\det S=-3\).  A root is an eigenspin line
of the imaginary quaternion, giving the two indicated axes.
The inverse image of a zero under the surjective Hopf map is nonempty.
In Theorem~\ref{p1:thm:joint-orbit-quadratic-descent} the complete
nonparallel tuple moves together; it is a different base and section.
\end{proof}



\section{The recentered exceptional two-plane quadratic}
\label{p1:app:v179-exceptional-two-plane}

This supplementary quadratic response is constructed from the coset maps.

\begin{proposition}[Quadratic pencil on the Hurwitz-fixed plane]
\label{p1:prop:v179-exceptional-quadratic-pencil}
Use the positive recentering of
Theorem~\ref{p1:thm:exceptional-core-visibility}.  If \(u_0\) is the
selected-cell vector and \(u_{\rm out}\) the positive unit sum of the
four outer cells, put
\[
 u_I=(u_0+2u_{\rm out})/\sqrt5,\quad
 v_I=(2u_0-u_{\rm out})/\sqrt5,\quad
 u_\pm=(u_0\pm u_{\rm out})/\sqrt2 .
\]
Then \(u_+=(3u_I+v_I)/\sqrt{10}\) and
\(u_-=(-u_I+3v_I)/\sqrt{10}\).
For \(P_I=|v_I\rangle\langle v_I|\) and
\(P_O=|u_-\rangle\langle u_-|\), one has
\(\operatorname{tr}(P_IP_O)=9/10\).  For real \(a,b\),
\[
 K(a,b)=aP_I+bP_O,\qquad
 \operatorname{tr}K=a+b,\qquad \det K=ab/10,
\]
\[
 \operatorname{spec}K=
 \left\{\tfrac12\left(a+b\pm
       \sqrt{(a-b)^2+\tfrac{18}{5}ab}\right)\right\}.
\]
\end{proposition}
\begin{proof}
The given basis change is orthogonal.  The component of \(u_-\) along
\(v_I\) is \(3/\sqrt{10}\), giving the projector overlap.
The two-dimensional determinant is \(ab(1-\operatorname{tr}P_IP_O)\);
the characteristic polynomial gives the eigenvalues.
\end{proof}

The coset-walk Laplacians of
Corollary~\ref{p1:cor:nearest-shell-inclusion-consequences} restrict to
\(D_I=\tfrac54P_I\) and to the transported \(D_O=2P_O\).
Their unit codiagonal sum has trace \(13/4\), determinant \(1/4\) and
spectrum \(\{(13-3\sqrt{17})/8,(13+3\sqrt{17})/8\}\).
This exact occurrence of \(\sqrt{17}\) is a dimensionless walk result.
The recentered \(2O\) action is transported from its two-coset
representation; it is not the restriction of the \(2I\) action.
Consequently this two-plane construction alone does not satisfy the
common continuous-representation hypothesis of
Theorem~\ref{p1:thm:2O-2I-generate-SU2}.


\section{Noncentral-element axes and extremal-pair geometry}
\label{p1:app:v191-occurrence-axis-bw}

\begin{proposition}[A noncentral $SU(2)$ element determines an inversion-even torus axis]
\label{p1:prop:v191-noncentral-maximal-torus-axis}
Let $q\in SU(2)$ with $q\neq\pm1$.  Then its centralizer is the unique
maximal torus containing it,
\[
 \boxed{Z_{SU(2)}(q)\cong U(1).}
\]
Writing $q=\cos\chi+\mathbf n\sin\chi$ with $0<\chi<\pi$ identifies the
corresponding unoriented axis $[\mathbf n]\in\mathbb{RP}^2$.  Inversion sends
$\mathbf n\mapsto-\mathbf n$ but leaves the torus fixed because
$Z(q^{-1})=Z(q)$.  Thus a noncentral element supplies an
inversion-even \emph{unoriented} axis.  The unmarked round sphere still has no
canonical preferred Hopf axis.
\end{proposition}

\begin{proposition}[Borel--Weil extremal-pair constructor on one fixed $SU(2)$ carrier]
\label{p1:prop:v191-bw-extremal-pair}
Fix \(\ell\in\tfrac12\ZZ_{\ge0}\) and one $SU(2)$ representation carrier
$V_\ell\cong H^0(\mathbb{CP}^1,\mathcal O(2\ell))^*$.  A point
$\mathbf n\in S^2\cong\mathbb{CP}^1$ determines the coherent extremal ray
$|\ell,+\ell\rangle_{\mathbf n}$, while its antipode determines
$|\ell,-\ell\rangle_{\mathbf n}$.  Hence an unoriented axis determines the
phase-independent extremal subspace
\[
 \boxed{\mathscr P_{\ell,[\mathbf n]}
 =\operatorname{span}\{|\ell,+\ell\rangle_{\mathbf n},
 |\ell,-\ell\rangle_{\mathbf n}\}.}
\]
Its dimension is two for \(\ell>0\) and one for \(\ell=0\), where the two
extremal rays coincide.  Equivalently, with $J_{\mathbf n}$ the rotation
generator in this same $SU(2)$ carrier, normalized with \(\hbar>0\),
\[
 P_{\ell,[\mathbf n]}
 =\mathbf1_{\{\hbar^2\ell^2\}}(J_{\mathbf n}^{\,2}).
\]
The construction is equivariant under rotation of the axis.
\end{proposition}

\begin{proof}
For any unit axis, \(J_{\mathbf n}\) is unitarily conjugate to \(J_3\), whose
eigenvalues are \(\hbar m\) for \(m=-\ell,-\ell+1,\ldots,\ell\).
For \(\ell>0\), the equation \(m^2=\ell^2\) selects two distinct
one-dimensional weight spaces.  At \(\ell=0\), it selects the single space
\(V_0=\CC\), so \(P_{0,[\mathbf n]}=I_{V_0}\).  Squaring the generator
removes the sign of the axis, and its spectral projector transforms by
unitary conjugation under rotations.
\end{proof}

\begin{remark}[Axis-dependent representation data]
The axis and its extremal projector are equivariant representation data.
The construction retains the family over all axes.
\end{remark}




\section{Primitive finite vector and bivector constructions}
\label{p1:sec:v558-primitive-owner-geometry}

{\emergencystretch=2em
The following results record owner types already implicit in the rooted
four-generator and exceptional-shell geometry, with their maps,
inner products and symmetry actions stated explicitly.\par}

\begin{theorem}[Rooted triality tetrad realizes the vector owner]
\label{p1:thm:v558-triality-tetrad-vector-owner}
Let \(\mathcal T=\{[a_1],\ldots,[a_4]\}\) be one of the tetrads of
Proposition~\ref{p1:prop:math-triality-octets-k444}.  After the already
constituted root/orientation data choose unit representatives \(a_i\).  If
\(V\) is the oriented Euclidean four-generator space, then
\[
 U_4:\mathbb R^{\mathcal T}\longrightarrow V,
 \qquad U_4\delta_i=a_i
\]
is an orthogonal isomorphism.  For every symmetry preserving the rooted
tetrad,
\[
 U_4\rho_{\mathcal T}(g)=\rho_V(g)U_4.
\]
Triality permutes the three tetrads, so they constitute one rank-four owner
type rather than three independent owner choices.
\end{theorem}
\begin{proof}
The four axes in each tetrad are mutually orthogonal by
Proposition~\ref{p1:prop:math-triality-octets-k444}; unit representatives are
therefore an orthonormal basis of \(V\).  The stabilizer action is the same
signed permutation on the source basis and on the corresponding vectors,
which proves equivariance.  Triality permutes the three perfect matchings.
\end{proof}

\begin{theorem}[Coordinate-pair carrier realizes the bivector owner]
\label{p1:thm:v558-coordinate-pair-bivector-owner}
Let
\[
 \mathcal P_6=\{01,02,03,12,13,23\}.
\]
For an oriented orthonormal basis \(e_0,\ldots,e_3\) of \(V\), define
\[
 U_6:\mathbb R^{\mathcal P_6}\longrightarrow\Lambda^2V,
 \qquad U_6\delta_{ij}=e_i\wedge e_j.
\]
Then \(U_6\) is an orthogonal isomorphism and is natural for every oriented
orthogonal symmetry preserving the rooted coordinate-pair presentation:
\[
 U_6\rho_{\mathcal P_6}(g)=(\Lambda^2\rho_V(g))U_6.
\]
Thus the primitive rank-six pair carrier is the bivector owner
\(\Lambda^2V\).
\end{theorem}
\begin{proof}
The six displayed bivectors are the standard orthonormal basis of
\(\Lambda^2V\).  Naturality is the defining functoriality of the exterior
square.
\end{proof}

\begin{proposition}[Fixed-point-free antipodal twelve-layer is rank-six composite]
\label{p1:prop:v558-antipodal-twelve-six-six}
Let \(X_0\) be a twelve-point middle incidence layer carrying the generated
fixed-point-free antipodal involution \(A\).  On \(\mathbb R[X_0]\),
\[
 A^2=I,\qquad \operatorname{tr}A=0,
\]
and hence
\[
 P_\pm=\frac{I\pm A}{2},\qquad
 \operatorname{rank}P_+=\operatorname{rank}P_-=6.
\]
Therefore
\[
 \boxed{\mathbb R[X_0]=\mathcal H_6^+\oplus\mathcal H_6^-}
\]
and the rank-twelve layer is composite relative to the generated antipodal
refinement. This proves the two ranks, not an identification of either
representation with the bivector owner \(\Lambda^2V\).
\end{proposition}
\begin{proof}
A fixed-point-free involution on twelve points is a product of six
transpositions, so its permutation trace is zero.  The rank formula follows
from the traces of \((I\pm A)/2\).
\end{proof}

\begin{proposition}[The two antipodal six-spaces have distinct rooted types]
\label{p1:prop:v214-antipodal-six-representation-distinction}
For the dual middle layer at \(u_+=(1,1,1,1)/2\),
\[
 X_0=\{\pm(e_i-e_j)/\sqrt2:1\le i<j\le4\},
\]
let \(A_4\) act by even coordinate permutations. It fixes \(u_+\),
preserves the layer, and commutes with antipodality. Its two eigenspaces
have complexified restrictions
\[
 (\mathcal H_6^+)_{\mathbb C}
 \cong\mathbf1\oplus\mathbf1'\oplus\mathbf1''\oplus\mathbf3,
 \qquad
 (\mathcal H_6^-)_{\mathbb C}\cong\mathbf3\oplus\mathbf3.
\]
In contrast,
\((\Lambda^2V)_{\mathbb C}|_{A_4}\cong\mathbf3\oplus\mathbf3\).
The odd space has the natural exterior-square identification
\[
 \delta_{(e_i-e_j)/\sqrt2}-\delta_{-(e_i-e_j)/\sqrt2}
 \longmapsto\sqrt2\,e_i\wedge e_j
 \quad(i<j),
\]
extended linearly on antisymmetric vertex coefficients; this is an
isometry for the counting metric. The even six-space is not
\(A_4\)-equivariantly isomorphic to it. Ranks alone therefore do not
exhaust owner representation types.
\end{proposition}
\begin{proof}
The even space is the permutation representation on the six unordered
coordinate pairs. On the identity, a double transposition and either
class of three-cycles its character is respectively \(6,2,0,0\).
Using the \(A_4\) characters gives
\(\mathbf1+\mathbf1'+\mathbf1''+\mathbf3\).
The normalized odd basis is
\((\delta_{(e_i-e_j)/\sqrt2}-\delta_{-(e_i-e_j)/\sqrt2})/\sqrt2\).
Swapping the ordered pair reverses its sign, exactly as on
\(e_i\wedge e_j\), proving the isometric intertwiner. The permutation
four-space restricts as \(\mathbf1\oplus\mathbf3\), whose exterior
square is \(\mathbf3\oplus\Lambda^2\mathbf3
=\mathbf3\oplus\mathbf3\), since \(A_4\) acts by oriented rotations
on its three-space. The even space has a fixed line and the odd space
has none, proving inequivalence.
\end{proof}

\section{Oriented Hurwitz--\texorpdfstring{\(D_4\)}{D4} frame determinant return}
\label{p1:sec:v559-d4-frame-determinant-return}

The following theorem is a mathematical descendant of the already constituted
Hurwitz--\(D_4\) cell.  It uses the complete orbit of oriented vertex--diagonal
local frames and the cyclic trimer carried by each such frame.

\begin{theorem}[Oriented-frame determinant return]
\label{p1:thm:v559-d4-oriented-frame-determinant-return}
Let \(\mathcal F_{D_4}\) be the complete orbit of oriented vertex--diagonal
local frames of the Hurwitz \(24\)-cell.  Then
\[
 |\mathcal F_{D_4}|=24\cdot8=192=|W(D_4)|.
\]
For every \(f\in\mathcal F_{D_4}\), let \(T_f\) be the oriented transverse
trimer and
\[
 G_f:=\CC^{T_f},\qquad D_f:=\Lambda^3G_f.
\]
The oriented-diagonal stabilizer \(C_3\) acts on \(G_f\) by a three-cycle and
therefore acts trivially on \(D_f\).  With the cyclic trimer basis
\((e_{f,0},e_{f,1},e_{f,2})\), set
\[
 \omega_f:=e_{f,0}\wedge e_{f,1}\wedge e_{f,2}\in D_f.
\]
This vector is independent of cyclic relabelling.  On
\[
 \mathscr D_{D_4}:=\bigoplus_{f\in\mathcal F_{D_4}}D_f
\]
use the inherited counting metric for which the \(\omega_f\) are orthonormal,
and define the primitive invariant column
\[
 J_{D_4}:\CC\longrightarrow\mathscr D_{D_4},\qquad
 J_{D_4}(1)=\sum_{f\in\mathcal F_{D_4}}\omega_f.
\]
Then
\[
 \boxed{J_{D_4}^{\dagger}J_{D_4}=192.}
\]
Equivalently, the closed determinant-line return of the complete oriented
frame orbit is the scalar \(192\).
\end{theorem}

\begin{proof}
The stabilizer chain in
Theorem~\ref{p1:thm:s3one-structural-invariants} gives
\(576\supset24\supset3\), hence the oriented frame orbit has cardinality
\(576/3=192\), equivalently \(24\cdot8\).  Proposition~\ref{p1:prop:diagonal-equivalence}
identifies the oriented-diagonal stabilizer with \(C_3\), acting by cyclic
permutation of the transverse trimer.  A three-cycle is even, so its action on
\(\Lambda^3G_f\) has determinant one and fixes \(\omega_f\).  Finally,
orthonormality of the summand determinant lines gives
\[
 J_{D_4}^{\dagger}J_{D_4}
 =\sum_{f\in\mathcal F_{D_4}}\langle\omega_f,\omega_f\rangle
 =|\mathcal F_{D_4}|=192.
\]
\end{proof}

\begin{corollary}[Unique positive central rank-three lift]
\label{p1:cor:v559-positive-central-rank3-lift}
Let \(G\) be any three-dimensional complex carrier and suppose a positive
central response \(R\in Z(\End G)\) is required to induce the determinant-line
scalar \(192\):
\[
 \Lambda^3R=192\,I_{\Lambda^3G}.
\]
Then
\[
 \boxed{R=xI_G,\qquad x^3=192,\qquad x=192^{1/3}.}
\]
This is the unique positive scalar operator with the specified determinant
on a three-dimensional complex Hilbert space.
\end{corollary}

\begin{proof}
Centrality in \(\End(G)\cong M_3(\CC)\) gives \(R=xI_G\).  Positivity gives
\(x>0\), while \(\Lambda^3R=(\det R)I=x^3I\).  Hence \(x^3=192\), with the
unique positive solution \(x=192^{1/3}\).
\end{proof}

\section{Rooted family co-realization of the \texorpdfstring{\(D_4\)}{D4} determinant response}
\label{p1:sec:v560-d4-rooted-family-corealization}

The determinant return is constructed from the same oriented
Hurwitz--\(D_4\) frame whose complementary trimer supplies the rooted
three-dimensional multiplicity carrier.  The following corollary removes a
purely mathematical duplication: once that rooted frame and its family trimer
are co-realized, the positive central lift of the determinant return does not
require a second choice of a rank-three scalar operator.

\begin{corollary}[Frame-independent rooted family lift]
\label{p1:cor:v560-frame-independent-family-lift}
Let \(f_0\in\mathcal F_{D_4}\) be an occurring rooted oriented local frame and
let
\[
 G_{f_0}=\CC^{T_{f_0}}\cong\CC[C_3]\cong\CC^3
\]
be its generated complementary trimer carrier.  Let
\[
 \rho_{D_4}:=J_{D_4}^{\dagger}J_{D_4}=192.
\]
Then the unique positive central operator on \(G_{f_0}\) whose determinant is
\(\rho_{D_4}\) is
\[
 \boxed{R_{D_4}^{\rm fam}:=\rho_{D_4}^{1/3}I_{G_{f_0}}
       =192^{1/3}I_{G_{f_0}}.}
\]
It is invariant under cyclic relabelling of the trimer and, under replacement
of the rooted frame by a symmetry-related frame, is transported to the
corresponding operator by the induced unitary identification.  Hence its
unitary-equivalence class is a generated rooted-family descendant of the
Hurwitz--\(D_4\) occurrence and contains no additional family-direction
selector.
\end{corollary}

\begin{proof}
The determinant scalar is generated by
Theorem~\ref{p1:thm:v559-d4-oriented-frame-determinant-return}.  The
complementary trimer of an oriented frame carries the regular \(C_3\)
representation, so a cyclic relabelling is unitary on \(G_{f_0}\).  By
Corollary~\ref{p1:cor:v559-positive-central-rank3-lift}, positivity and
centrality force the lift to be \(xI\) with \(x^3=192\).  Scalar operators
commute with every unitary relabelling, proving presentation independence.
Thus the lift is defined on every rooted trimer and is covariant under
the induced unitary changes of rooted presentation.
\end{proof}

\begin{remark}[Rank-three lift and marking]
\label{p1:rem:v560-cross-role-boundary}
The determinant scalar fixes a unique positive central rank-three lift.
A noncentral positive operator with the same determinant requires additional
matrix data; its determinant alone does not recover its eigenspaces.
\end{remark}


\section{Oriented-edge triality on the Hurwitz--\texorpdfstring{\(D_4\)}{D4} frame orbit}
\label{p1:sec:v561-d4-triality-incidence-parent}

The determinant return and its rooted co-realization use the complete orbit of oriented
Hurwitz--\(D_4\) local frames.  The same finite cell also carries the ordered
weight support of the Spin(8) triality trilinear.  The statements below are
incidence and invariant-tensor theorems on the specified finite supports.

Let
\[
 V_8:=\{\pm e_i:1\le i\le4\},
\]
and let \(S_8^+\) and \(S_8^-\) be the two parity classes of half-integral unit
vectors
\[
 \frac12(\pm e_1\pm e_2\pm e_3\pm e_4),
\]
with respectively even and odd numbers of minus signs.  These are the three
eight-weight sets customarily denoted \(8_v,8_s,8_c\); triality permutes them.

\begin{theorem}[Oriented-edge triality-support correspondence]
\label{p1:thm:v561-oriented-edge-triality-support}
Let \(a\to b\) be an oriented edge of the short-root/dual \(24\)-cell, so
\(a,b\in V_8\sqcup S_8^+\sqcup S_8^-\), \(|a|=|b|=1\), and
\(\langle a,b\rangle=\tfrac12\).  Then \(a,b\) lie in distinct triality
sectors and
\[
 c:=b-a
\]
is a unit weight in the third sector.  Hence
\[
 \boxed{a+(-b)+c=0.}
\]
For each ordered pair of distinct sectors there are exactly \(32\) such
oriented edges; the six ordered sector pairs therefore give
\[
 \boxed{6\cdot32=192}
\]
ordered triality incidences. For a fixed slot order \((V,S,C)\), there
are \(32\) weight-zero support triples of the Spin(8)-invariant
trilinear \(8_v\otimes8_s\otimes8_c\to\mathbb R\). The \(192\)
incidences retain all six slot-order tags; they are not \(192\)
distinct basis tensors in one fixed tensor product.
\end{theorem}

\begin{proof}
Adjacency of short roots is \(\langle a,b\rangle=\tfrac12\).  Thus
\[
 |b-a|^2=|a|^2+|b|^2-2\langle a,b\rangle=1.
\]
A direct parity check in the displayed weight sets shows that two adjacent
weights cannot lie in the same sector and that \(b-a\) lies in the remaining
sector.  For a fixed weight in one sector, exactly four of its eight neighbors
lie in either specified other sector.  Hence an ordered sector pair contains
\(8\cdot4=32\) oriented edges.  The zero-sum identity is immediate from the
definition of \(c\), and it is the standard weight-selection rule for the
triality trilinear.
\end{proof}

\begin{definition}[Primitive ordered triality-incidence tensor]
\label{p1:def:v561-d4-triality-incidence-tensor}
Let \(\mathcal S=\{(v,s,c)\in V_8\times S_8^+\times S_8^-:
v+s+c=0\}\), so \(|\mathcal S|=32\). Give the three complex weight
carriers their orthonormal counting bases. In the coefficient comparison
below use the dual coefficient carriers:
\[
 \mathcal H_{\rm id}
 =(\mathbb C^{V_8})^*\otimes(\mathbb C^{S_8^+})^*
                              \otimes(\mathbb C^{S_8^-})^*.
\]
The notation \(e_v,e_s,e_c\) in this definition denotes their
orthonormal dual coordinate bases. For \(\pi\in S_3\), let
\(\mathcal H_\pi\) be the tensor product with the three slots in order
\(\pi\), and let \(U_\pi:\mathcal H_{\rm id}\to\mathcal H_\pi\)
permute the slots. On the explicit tagged target put
\[
 \mathcal H_{\rm ord}:=\bigoplus_{\pi\in S_3}\mathcal H_\pi,\qquad
 \Theta_{D_4}^{\rm ord}
 :=\bigoplus_{\pi\in S_3}U_\pi
       \sum_{(v,s,c)\in\mathcal S}e_v\otimes e_s\otimes e_c.
\]
The ordered support is \(\mathcal S_{\rm ord}=S_3\times\mathcal S\).
An element with ordered weights \((a,d,c)\) corresponds to the directed
edge \(a\to-d\), since \(c=-a-d\); this is the incidence bijection.
This is an unsigned support-counting tensor, not the invariant
coefficient tensor. The signed map to the latter is specified below.
\end{definition}

\begin{corollary}[Triality-parent norm and frame-return count]
\label{p1:cor:v561-triality-parent-norm-192}
With the inherited counting metric,
\[
 \boxed{\|\Theta_{D_4}^{\rm ord}\|^2=192.}
\]
Thus the same Hurwitz--\(D_4\) finite owner carries both the oriented-frame
closed return
\[
 J_{D_4}^{\dagger}J_{D_4}=192
\]
and the complete ordered triality-incidence norm \(192\).  These are sibling
closed contractions of one finite incidence geometry, not an equality inferred
from dimension alone.
\end{corollary}

\begin{proof}
The support tensors are orthonormal simple tensors and every support element
occurs once in its tagged summand, so the norm square is
\(6\cdot32=192\). The frame-return
identity is Theorem~\ref{p1:thm:v559-d4-oriented-frame-determinant-return}; both
constructions descend from the same oriented Hurwitz--\(D_4\) cell and use its
complete \(192\)-element oriented incidence orbit.
\end{proof}

\begin{theorem}[Explicit metric-normalized Spin(8) triality and restriction]
\label{p1:thm:v214-normalized-spin8-triality}
Let \(U=\mathbb C^4\) have ordered orthonormal basis \(f_1,\ldots,f_4\).
On \(S^+=\Lambda^{\rm even}U\) and \(S^-=\Lambda^{\rm odd}U\),
write \(f_I\) for the orthonormal increasing wedge basis.
The complex vector representation has isotropic orthonormal weight
basis \(\mathsf p_i,\mathsf q_i\), with bilinear metric
\(g(\mathsf p_i,\mathsf q_j)=\delta_{ij}\), and real structure
\(\overline{\mathsf p_i}=\mathsf q_i\). Its real orthonormal basis is
\[
 a_i=(\mathsf p_i+\mathsf q_i)/\sqrt2,\qquad
 b_i=i(\mathsf p_i-\mathsf q_i)/\sqrt2.
\]
Put
\[
 c(\mathsf p_i)=\sqrt2\,\varepsilon_i,\qquad
 c(\mathsf q_i)=\sqrt2\,\iota_i,\qquad
 B(\alpha,\beta)=[\tau(\alpha)\wedge\beta]_{f_{1234}},
 \qquad T(v,s,t)=B(c(v)s,t),
\]
where \(\tau\) reverses wedge order.
This is a Spin(8)-invariant trilinear on the vector and two half-spin
modules, with the metric Clifford normalization
\(\{c(v),c(w)\}=2g(v,w)I\). For
\[
 n_i(I)=|\{j\in I:j<i\}|,\qquad
 \beta(K)=(-1)^{\sum_{j\in K}(j-1)},
\]
its complete coefficients are
\[
 T(\mathsf p_i,f_I,f_J)=
 \begin{cases}
 \sqrt2(-1)^{n_i(I)}\beta(I\cup\{i\}),
  &i\notin I,\ J=(I\cup\{i\})^c,\\
 0,&\text{otherwise},
 \end{cases}
\]
\[
 T(\mathsf q_i,f_I,f_J)=
 \begin{cases}
 \sqrt2(-1)^{n_i(I)}\beta(I\setminus\{i\}),
  &i\in I,\ J=(I\setminus\{i\})^c,\\
 0,&\text{otherwise}.
 \end{cases}
\]
Here \(|I|\) is even and \(|J|\) odd. There are \(32\) nonzero
fixed-slot coefficients, each of modulus \(\sqrt2\), and hence
\[
 \boxed{\|T\|_{\rm HS}^2=64.}
\]
For \(U=U_A\oplus U_B\), with
\(U_A=\operatorname{span}(f_1,f_2)\), exterior factorization gives
\[
 c_8(v_A)=c_A(v_A)\otimes I,\qquad
 T|_{V_A\otimes(S_{A+}\otimes S_{B+})
                 \otimes(S_{A-}\otimes S_{B+})}
 =T_A\otimes\varepsilon_B.
\]
The alternating spectator form \(\varepsilon_B\) has
\(\|\varepsilon_B\|_{\rm HS}^2=2\). Consequently the precisely
normalized spectator extraction
\[
 \operatorname{Ext}_B(F)
 :=\tfrac12\sum_{\alpha,\beta}
   \overline{(\varepsilon_B)_{\alpha\beta}}\,
                 F_{\alpha\beta}
 \quad\hbox{satisfies}\quad
 \boxed{\operatorname{Ext}_B(T)=T_A.}
\]
After the nondegenerate pairing \(B_A\) identifies the last slot
with its dual, this is exactly \(c_A(v)\), with coefficient one.
The companion block gives the opposite chirality.
In the protected Paper-I-A normalization the coordinates are therefore
\(X(v)=(\rho_{\mathsf C}/2)c_A(v)\), not an independently
rescaled triality map.
\end{theorem}
\begin{proof}
The creation and annihilation relations give the Clifford relation and
Hermiticity of \(c(a_i)=\varepsilon_i+\iota_i\) and
\(c(b_i)=i(\varepsilon_i-\iota_i)\).
Direct reversal of wedge factors gives
\[
 B(f_K,f_J)=\beta(K)\,\mathbf1_{\{J=K^c\}},\qquad
 B(c(v)\alpha,\beta)=B(\alpha,c(v)\beta).
\]
The second identity makes every Clifford commutator skew for \(B\);
the connected spin group generated by these commutators preserves \(B\).
Clifford covariance then proves invariance of \(T\).
The first identity and the insertion/removal sign
\((-1)^{n_i(I)}\) prove all displayed coefficients.
Assign \(f_I\) the weight
\(\frac12\sum_j(2\mathbf1_{\{j\in I\}}-1)e_j\).
The displayed nonzero coefficients have precisely zero total weight.
For each of the eight vector weights there are four admissible even
subsets \(I\), proving the count and norm. Equivalently, for each real
unit vector \(a\), \(c(a)\) is unitary between the two eight-dimensional
half-spin spaces and has squared Hilbert--Schmidt norm eight; summing
over eight real orthonormal vector directions gives \(64\).

Under the ordered factorization
\(\Lambda^*U=\Lambda^*U_A\otimes\Lambda^*U_B\), the first two creation
and annihilation operators act on the first factor, with no spectator
coefficient. The Chevalley pairing factors as \(B_A B_B\) on the
displayed block: the possible reordering sign is
\((-1)^{|\alpha_B|(|\alpha_A|+|\beta_A|)}=1\) when the pairing is
nonzero. In two generators each chiral Chevalley pairing is the
alternating form of norm square two. This proves both the exact
factorization and the factor \(1/2\) in extraction.
Thus the coefficient is computed, rather than inferred from
multiplicity one. Contraction with the unnormalized inverse
alternating form without \(1/2\) would give \(2T_A\);
contraction with a unit singlet would give \(\sqrt2T_A\).
These are different explicitly normalized extraction operations.
\end{proof}

\begin{corollary}[Signed averaging from the tagged incidence tensor]
\label{p1:cor:v214-tagged-signed-triality-projection}
Using the weight phases of the preceding theorem, define on
\(\mathcal S\)
\[
 \sigma(v,s,c)=T(v,s,c)/\sqrt2\in\{1,-1\}.
\]
Let \(D_\sigma\) multiply each tagged support basis tensor by this
sign and extend by the identity off that support. The specified signed
averaging map is
\[
 \mathcal A=\tfrac16\sum_{\pi\in S_3}U_\pi^*:
 \mathcal H_{\rm ord}\longrightarrow\mathcal H_{\rm id}.
\]
Then
\[
 \boxed{\sqrt2\,\mathcal A D_\sigma\Theta_{D_4}^{\rm ord}=T.}
\]
The unsigned norm \(192\), the fixed-slot signed support norm \(32\),
and the metric-normalized invariant tensor norm \(64\) refer to the
three displayed targets and maps. None is substituted for another.
\end{corollary}
\begin{proof}
Each inverse slot permutation returns the same signed fixed-slot tensor
\[
 \sum_{\mathcal S}\sigma(v,s,c)e_v\otimes e_s\otimes e_c.
\]
Averaging removes the six tags and multiplication by
\(\sqrt2\) supplies exactly the coefficients of \(T\).
\end{proof}

\begin{remark}[Support and invariant coefficients]
\label{p1:rem:v561-triality-physical-firewall}
The tensor \(\Theta_{D_4}^{\rm ord}\) records tagged unsigned support.
The signed averaging map in
Corollary~\ref{p1:cor:v214-tagged-signed-triality-projection} produces the
metric-normalized invariant tensor in its fixed tensor-product target.
Keeping these maps explicit distinguishes their norms and coefficients.
\end{remark}


\section{Why the coordinate class does not fix a square's ground state}
\label{p1:app:v233-square-owner}
Let \(\lambda=\pm3/(2R)\) be either lowest spinor eigenvalue of the
parent Levi--Civita operator. The self-adjoint operator
\(\mathscr D'=\mathscr D_{\rm LC}-\lambda I\) has the same principal
symbol and
\[
 [\mathscr D',M_f]=[\mathscr D_{\rm LC},M_f].
\]
Thus it supplies exactly the same compressed intrinsic coordinate
class. But \(J_R^*(\mathscr D')^2J_R\ge0\) is zero on the retained
spinor eigenspace with eigenvalue \(\lambda\), while its value on the
constant scalar is \(\lambda^2>0\). Its ground space is consequently
not the scalar line. The geometric ground-state corollary applies to
the specified \(K_{\rm LC}\), not to every operator with that symbol.
A ground-space claim must therefore specify the operator square, including
its zeroth-order part, in addition to its principal symbol.






\section{Natural determinant lift and rooted axial triality contraction}
\label{p1:sec:v562-d4-natural-lift-axial-triality}

The preceding releases derived the scalar determinant return
\(J_{D_4}^{\dagger}J_{D_4}=192\) and its positive central lift.  Centrality is
not an extra datum: it follows from equivariance of the determinant-only lift
under the already generated determinant-one trimer automorphisms.  The same
rooted triality parent also admits a finite axial contraction that fixes the
complementary trimer pointwise.  The two constructions retain their
displayed determinant and incidence normalizations.

\begin{theorem}[Natural rooted determinant lift]
\label{p1:thm:v562-natural-rooted-determinant-lift}
Let \(G\cong\CC^3\) be the rooted trimer carrier with generated cyclic shift
\(X_G\) and generated rooted phase \(R_G=\operatorname{diag}(1,-i,i)\).  Let
\(R\in\End(G)\) be positive and suppose its constructor consumes only the
rooted Hilbert carrier and the determinant scalar
\[
 \det R=J_{D_4}^{\dagger}J_{D_4}=192,
\]
with no marked rank-one family projector or noncentral family direction.
Same-owner equivariance under the generated determinant-one automorphisms
\(X_G\) and \(R_G\) forces
\[
 X_GRX_G^\dagger=R,
 \qquad
 R_GRR_G^\dagger=R.
\]
Hence
\[
 \boxed{R=xI_G,\qquad x^3=192,\qquad x=192^{1/3}.}
\]
Thus centrality is a theorem of the unmarked determinant-lift word, not an
additional lift premise.
\end{theorem}

\begin{proof}
Because \(R_G\) has three distinct eigenvalues, commutation with \(R_G\)
forces \(R\) to be diagonal in the rooted trimer basis,
\(R=\operatorname{diag}(r_0,r_1,r_2)\).  Commutation with the cyclic shift
\(X_G\) then gives \(r_0=r_1=r_2=:x\).  Positivity gives \(x>0\), and the
determinant constraint gives \(x^3=192\).  No-insertion forbids replacing
this unmarked lift by a noncentral lift without adding a separately generated
family-direction input.
\end{proof}

\begin{theorem}[Axial triality contraction on the complementary trimer]
\label{p1:thm:v562-axial-triality-trimer-identity}
Use the rooted triality weights
\[
 u_+=\tfrac12(1,1,1,1),\qquad
 u_-=\tfrac12(1,-1,-1,-1),
\]
\[
 m_0=\tfrac12(1,1,-1,-1),\quad
 m_1=\tfrac12(1,-1,1,-1),\quad
 m_2=\tfrac12(1,-1,-1,1),
\]
\[
 p_0=\tfrac12(1,-1,1,1),\quad
 p_1=\tfrac12(1,1,-1,1),\quad
 p_2=\tfrac12(1,1,1,-1).
\]
For the oriented vector weight \(v=-e_1\), the triality support rule
\(v+s+c=0\) gives
\[
 u_+\mapsto u_-,\qquad m_j\mapsto p_j\quad(j=0,1,2).
\]
Consequently, in the ordered quartet bases
\((u_+,m_0,m_1,m_2)\) and \((u_-,p_0,p_1,p_2)\), the finite triality
contraction is
\[
 \boxed{C_{-e_1}=I_4.}
\]
After removing the axial pair, its complementary-trimer restriction is
\[
 \boxed{C_{-e_1}|_{T_+\to T_-}=I_3.}
\]
\end{theorem}

\begin{proof}
For \(v=-e_1\), the support equation gives \(c=e_1-s\).  Direct substitution
yields \(e_1-u_+=u_-\) and \(e_1-m_j=p_j\) for \(j=0,1,2\).  Each source
weight therefore maps to the correspondingly ordered target weight with unit
incidence coefficient.
\end{proof}

On the root link the support rule \(s\mapsto e_1-s\)
is the root involution \(\kappa_v\) of
Theorem~\ref{p1:thm:rooted-trimer-vector-overlap}, with the
first coordinate called \(e_1\) in the triality notation.
Thus the aligned finite support map agrees with the geometric
pairing of the two trimers. Its incidence coefficients and the
ambient Gram coefficients remain the distinct maps displayed above.

\begin{remark}[Axial contraction and half-step transport]
\label{p1:rem:v562-half-step-firewall}
The identity contraction is distinct from the two rooted antiprism
half-step matchings \(X_G\) and \(X_G^\dagger\).  Thus the axial map
does not select the orientation-odd operator
\(C_{2,G}=-i(X_G-X_G^\dagger)\).  Selecting a half-step requires the
corresponding oriented matching datum.
\end{remark}


\section{Affine rooted reflection and transport inversion}
\label{p1:sec:v563-affine-reflection-inversion}

The axial triality theorem fixes the unmarked axial contraction but does
not determine either oriented antiprism half-step.  The following theorem gives
the exact action of every rooted affine reflection on those already generated
half-step transports.

\begin{theorem}[Affine rooted reflection acts by cyclic-transport inversion]
\label{p1:thm:v563-affine-reflection-inversion}
Let the rooted trimer be labelled by $k\in\mathbb Z_3$ and let
\[
 X_G e_k=e_{k+1}.
\]
For any $c\in\mathbb Z_3$ define the orientation-reversing affine reflection
\[
 s_c:k\longmapsto c-k\pmod 3,
 \qquad S_c e_k=e_{c-k}.
\]
Then
\[
 S_c^2=I,
 \qquad
 \boxed{S_cX_GS_c^{-1}=X_G^{-1}=X_G^\dagger}
\]
for every value of the affine offset $c$.
\end{theorem}

\begin{proof}
For every basis vector,
\[
 S_cX_GS_c^{-1}e_k
 =S_cX_Ge_{c-k}
 =S_ce_{c-k+1}
 =e_{k-1}
 =X_G^{-1}e_k.
\]
The formula is therefore independent of the choice of rooted affine origin.
\end{proof}

\begin{corollary}[Even/odd transport split and normalized cyclic generator]
\label{p1:cor:v563-signline-parity-generator}
With
\[
 C_{1,G}=X_G+X_G^\dagger,
 \qquad
 C_{2,G}=-i(X_G-X_G^\dagger),
\]
one has
\[
 S_cC_{1,G}S_c^{-1}=C_{1,G},
 \qquad
 \boxed{S_cC_{2,G}S_c^{-1}=-C_{2,G}}.
\]
Moreover $\operatorname{spec}(C_{2,G})=\{0,\sqrt3,-\sqrt3\}$ and, with the
cyclic orientation chosen consistently with $X_G$,
\[
 \boxed{
 X_G=\exp\!\left(i\frac{2\pi}{3\sqrt3}C_{2,G}\right).}
\]
Thus $C_{2,G}$ is the normalized Hermitian generator of the rooted $C_3$
transport, not an independently inserted family matrix.
\end{corollary}

\begin{proof}
The parity identities follow from the theorem.  Diagonalizing $X_G$ gives
$\operatorname{spec}(X_G)=\{1,\omega,\omega^2\}$ with
$\omega=e^{2\pi i/3}$, and hence
$\operatorname{spec}(C_{2,G})=\{0,\sqrt3,-\sqrt3\}$.  Functional calculus
then gives the displayed exponential identity.
\end{proof}

\begin{remark}[Root alignment and transport inversion]
\label{p1:rem:v563-halfstep-source-firewall}
The coordinate expression \(k\mapsto-k\) requires aligned origins,
but every affine reflection \(k\mapsto c-k\) induces the same
operator inversion \(X_G\leftrightarrow X_G^\dagger\).  The unmarked
axial contraction remains \(I_G\), while the two half-step transports
are interchanged by reflection.
\end{remark}



\section{Rooted common-parent return algebra}
\label{p1:sec:v566-rooted-return-algebra}
The trimer relations arise as closed paths through a common finite
Hermitian carrier.  The return algebra retains these relations before
applying a scalar, determinant or singular-value readout.

\begin{definition}[Rooted return through a common owner]
\label{p1:def:v566-rooted-return}
Let $A,B:G\to\mathcal H$ be generated arrows into the same finite Hermitian
owner and let $K_{\mathcal H}>0$ be its inherited metric operator on the
relevant support.  Define
\[
 \operatorname{Ret}_G(A,B):=A^\dagger K_{\mathcal H}^{-1}B
 \in\operatorname{End}(G).
\]
For generated arrows in the opposite direction, $A,B:\mathcal H\to G$, the
corresponding return is $AB^\dagger$.  The \emph{rooted return algebra}
$\mathfrak R_G$ is the unital $*$-algebra generated by all such returns from
one fixed rooted geometry, with all inner products and arrows retained.
\end{definition}

\begin{theorem}[Existing Hurwitz returns generate the full family algebra]
\label{p1:thm:v566-rooted-return-full-family-algebra}
On the rooted Hurwitz branch of
Theorem~\ref{p1:thm:v556-rooted-hurwitz-family-return-algebra},
\[
 \boxed{\mathfrak R_G=\operatorname{End}(G)=M_3(\mathbb C).}
\]
In particular, the shell--trimer incidence returns and the rooted phase are
common-parent returns in the sense of
Definition~\ref{p1:def:v566-rooted-return}, and their generated algebra already
contains all family matrix units.
\end{theorem}

\begin{proof}
The shell incidence maps occurring in
Theorem~\ref{p1:thm:v556-rooted-hurwitz-family-return-algebra} have the form
$\mathcal I_S:\mathbb C^S\to G$, so
$\mathcal I_S\mathcal I_S^\dagger$ is a rooted return.  That theorem proves
\[
 C^*(K_7,K_9,R_G)=M_3(\mathbb C).
\]
Hence $M_3(\mathbb C)\subseteq\mathfrak R_G$.  The reverse inclusion is
immediate because every return acts on $G$.
\end{proof}

\begin{corollary}[Relation data precede scalar terminal readout]
\label{p1:cor:v566-return-before-scalar-readout}
Whenever a later scalar, norm, determinant, support quotient or singular-value
readout is taken from a generated owner arrow, that readout is a descendant of
the owner-level return data.  It cannot be used to reconstruct discarded
noncentral returns unless the readout is proved injective on the relations
required by the claim.
\end{corollary}

\section{Signed rooted triality channels and family return}
\label{p1:sec:v567-signed-triality-channel-family}
The oriented-edge theorem fixes the complete weight-zero \emph{support}
of the rooted triality geometry.  A Clifford/triality contraction also
retains the signed Clifford coefficients.  The
support count and the invariant coefficient tensor are therefore distinct
typed descendants of the same rooted cell.

\begin{remark}[Unsigned incidence support versus signed invariant coefficients]
\label{p1:rem:v567-v561-support-scope}
The tensor $\Theta_{D_4}^{\rm ord}$ of
Definition~\ref{p1:def:v561-d4-triality-incidence-tensor} is retained as the
unsigned ordered \emph{support-counting} tensor, and its norm-square $192$
remains the exact support count of Corollary~\ref{p1:cor:v561-triality-parent-norm-192}.
It is not, by itself, the normalized Spin$(8)$-invariant coefficient tensor.
The explicit coefficient calculation in
Theorem~\ref{p1:thm:v214-normalized-spin8-triality} supplies the
occupation/Koszul signs and the real-vector metric conversion.
Corollary~\ref{p1:cor:v214-tagged-signed-triality-projection} specifies
the signed map from the six-tagged support count to that tensor.
No support count is changed by this correction.
\end{remark}

\begin{theorem}[Signed rooted triality channel isometry]
\label{p1:thm:v567-signed-triality-channel-family}
Use the rooted quartets
\[
 S_+=(u_+,m_0,m_1,m_2),\qquad
 C_+=(u_-,p_0,p_1,p_2),\qquad
 C_-=(n_*,n_0,n_1,n_2),
\]
where $n_*=-u_-$ and $n_i=-p_i$.  Let
$I_0:S_+\to C_+$ be the rooted axial triality isometry
$u_+\mapsto u_-$, $m_j\mapsto p_j$, and put
\[
 \eta=(\eta_0,\eta_1,\eta_2)=(1,-1,1),
 \qquad
 D_\eta:=\operatorname{diag}(1,-1,1).
\]
Use the Fock normalization of
Theorem~\ref{p1:thm:v214-normalized-spin8-triality}.
Here \(T(v)\) means the coefficient map
\[
 T(v):S^+\longrightarrow(S^-)^*,\qquad
             s\longmapsto T(v,s,\,\cdot\,).
\]
The target kets \(u_-,p_i,n_*,n_i\) are orthonormal last-slot
dual coordinate vectors indexed by the displayed support weights.
They are not identified with the raw Clifford output vectors;
comparison with \(c(v)\) retains the Chevalley pairing \(B\).
For \(j=i+2\), let \(N_j\) be the occupation projector on the rooted
source quartet; its occupied vectors are \(u_+\) and \(m_i\).
Take the real unit vector \(a_j=(\mathsf p_j+\mathsf q_j)/\sqrt2\) and retain its
occupation-lowering part \(T(a_j)N_j\).
Thus the following are occupation-restricted transverse channels,
not the whole transverse four-by-four Clifford matrix:
\[
 \widetilde B_i
 =|n_i\rangle\langle u_+|
  +\eta_i|n_*\rangle\langle m_i|,
 \qquad i=0,1,2.
\]
Let
\[
 G=\operatorname{span}_{\mathbb C}\{|p_0\rangle,|p_1\rangle,|p_2\rangle\},
 \qquad E=\mathbb C^{\{0,1,2\}},
\]
and define
\[
 \mathcal V:G\longrightarrow E\otimes C_-,\qquad
 \mathcal Vg:=\sum_{i=0}^2|i\rangle\otimes
 \widetilde B_i I_0^\dagger g.
\]
Then
\[
 \boxed{\mathcal V|p_j\rangle
 =\eta_j|j\rangle\otimes|n_*\rangle,\qquad
 \mathcal V^\dagger\mathcal V=I_G.}
\]
Moreover, for every $W\in\operatorname{End}(E)$,
\[
 \boxed{
 \mathcal V^\dagger(W\otimes I_{C_-})\mathcal V
 =D_\eta W D_\eta.}
\]
Thus
\[
 \boxed{
 \chi:\operatorname{End}(E)\xrightarrow{\ \cong\ }\operatorname{End}(G),
 \qquad \chi(W)=D_\eta W D_\eta}
\]
is a rooted unital $*$-isomorphism.
\end{theorem}

\begin{proof}
The increasing Fock basis fixes an explicit compatible rooted phase
convention:
\[
\begin{aligned}
 |u_+\rangle&=f_{1234},&
 (|m_0\rangle,|m_1\rangle,|m_2\rangle)
 &=(-f_{12},f_{13},-f_{14}),\\
 |u_-\rangle&=f_1^*,&
 (|p_0\rangle,|p_1\rangle,|p_2\rangle)
 &=(f_{134}^*,f_{124}^*,f_{123}^*),\\
 |n_*\rangle&=f_{234}^*,&
 (|n_0\rangle,|n_1\rangle,|n_2\rangle)
 &=(f_2^*,f_3^*,f_4^*).
\end{aligned}
\]
The symbols \(n_i=-p_i\) denote opposite weights; they do not impose
an equality of independently phased complex basis vectors.
All group matrices and arrows are transported in this same phase
convention. A change of rooted phases conjugates those matrices and the
displayed channel-to-family isomorphism together; it does not require
both channel and family cycle matrices to be unsigned permutations
simultaneously.
For \(a_1\), the creation part vanishes on the source quartet.
The coefficient formulas of the cited theorem give
\(T(a_1)|_{S_+\to C_+}=I_0\) with all four coefficients one.
For \(j=2,3,4\), contraction of \(f_{1234}\), followed by the
Chevalley complementary-index pairing, has coefficient \(+1\)
on \(n_{j-2}\): its insertion sign and pairing sign cancel.
On \((f_{12},f_{13},f_{14})\) the corresponding removal has
coefficient \(-1\) on \(n_*\). Multiplying by the three source
phases \((-1,1,-1)\) gives exactly \(\eta=(1,-1,1)\).
Since \(N_j\) kills the other two source vectors, these are all
coefficients of the stated channel. The metric factor is also fixed:
\(c(\mathsf q_j)=\sqrt2\,\iota_j\), while the \(\mathsf q_j\) coefficient in the
real unit \(a_j\) is \(1/\sqrt2\); their product is one.

Since $I_0^\dagger|p_j\rangle=|m_j\rangle$, the first term in
$\widetilde B_i$ vanishes on $m_j$ and the second gives
$\widetilde B_i|m_j\rangle=\eta_i\delta_{ij}|n_*\rangle$.
This proves the displayed action of $\mathcal V$, hence
$\mathcal V^\dagger\mathcal V=I_G$.  For matrix units $E_{ij}$ on $E$,
\[
 \mathcal V^\dagger(E_{ij}\otimes I)\mathcal V
 =\eta_i\eta_j|p_i\rangle\langle p_j|
 =D_\eta E_{ij}D_\eta.
\]
Linearity proves the formula for arbitrary $W$.  Since $D_\eta^2=I$, the map
is a unital $*$-isomorphism with inverse itself in the aligned rooted bases.
\end{proof}

\begin{corollary}[No independent channel-to-family embedding]
\label{p1:cor:v567-no-independent-channel-family-map}
Every generated channel response $W\in\operatorname{End}(E)$ has one rooted
family representative $\chi(W)$, and every family operator has the unique
channel representative $\chi^{-1}(A)=D_\eta A D_\eta$.
The signed channel consequently fixes the isomorphism of these two
matrix algebras without an additional choice of a unitary embedding.
\end{corollary}

\begin{remark}[Dependence on the rooted channel]
\label{p1:rem:v567-signed-triality-source-firewall}
The isomorphism uses the displayed rooted bases and signs.
A simultaneous unitary change of those data transports the isomorphism;
discarding the signs would replace it by a different channel map.
\end{remark}

\subsection{Full-quartet phase lift and its complete channel return}
\label{p1:sec:v242-signed-phase-lift}
Keep the rooted Fock conventions and occupation projectors of
Theorem~\ref{p1:thm:v567-signed-triality-channel-family}. The following
calculations retain the seed column as well as the three family columns.

\begin{lemma}[The full seed determines a real vector]
\label{p1:lem:v242-seed-real-isometry}
For a real vector
\(v=\sum_{j=1}^4(x_ja_j+y_jb_j)\), put \(z_j=x_j-iy_j\).
In the last-slot dual coordinates of the signed tensor,
\[
 \boxed{T(v)u_+=z_1u_-+\sum_{i=0}^2z_{i+2}n_i.}
\]
Consequently \(v\mapsto T(v)u_+\) is a real-linear isometry from
\(V_8\) onto the realification of the complex span of the four
displayed target vectors.
In particular its value determines all eight real coefficients of \(v\).
\end{lemma}
\begin{proof}
On \(u_+=f_{1234}\), every creation operator vanishes. The coefficient
formula in Theorem~\ref{p1:thm:v214-normalized-spin8-triality} gives
\(T(a_1)u_+=u_-\), \(T(a_{i+2})u_+=n_i\), and
\(T(b_j)u_+=-iT(a_j)u_+\). The displayed formula follows.
Its squared norm is \(\sum_j|z_j|^2=\sum_j(x_j^2+y_j^2)\).
The four displayed target vectors are orthonormal, proving surjectivity
onto their realification and injectivity on real vectors.
\end{proof}

\begin{corollary}[Unique real-vector lift of the existing rooted Hopf phase]
\label{p1:cor:v242-joint-phase-vector-lift}
Let \(r_i=m_+(q_i)/m_+(q_0)\) be the already defined rooted joint-quadratic
ratios of Theorem~\ref{p1b:thm:v135-hopf-sector-phase}. Thus
\(r=(1,-i,i)\) and \(R_+=\operatorname{diag}(r_i)\).
The unique signed channel representative of this phase is
\[
 W_+=\chi^{-1}(R_+)=D_\eta R_+D_\eta=R_+.
\]
The unique real vector satisfying
\(T(v_i)u_+=r_i n_i\) is
\[
 v_i=\operatorname{Re}(r_i)a_{i+2}
       -\operatorname{Im}(r_i)b_{i+2},\qquad
 \boxed{(v_0,v_1,v_2)=(a_2,b_3,-b_4).}
\]
For \(L_i=T(a_{i+2})N_{i+2}\) on the rooted quartet, one has
\[
 T(v_i)N_{i+2}=r_iL_i,
 \qquad
 \boxed{\mathcal V^\dagger(W_+\otimes I)\mathcal V=R_+.}
\]
The negative Hopf line gives the conjugate ratios and inverse real
rotations. No extra channel-to-family isometry is chosen.
\end{corollary}
\begin{proof}
The signed channel theorem fixes \(\chi\). Since \(D_\eta\) and
\(R_+\) are diagonal, they commute. The seed isometry then determines
each \(v_i\) uniquely. On the occupied source,
\(T(b_j)N_j=-iT(a_j)N_j\), because only annihilation remains.
This proves the full-quartet channel identity, including its seed
column. The signed return formula follows from the channel theorem.
Complex conjugation reverses the angles and gives the other Hopf line.
\end{proof}

\begin{lemma}[Distinct channel phases cannot be moved to one spinor leg alone]
\label{p1:lem:v242-one-leg-nonabsorption}
On these fixed rooted supports write
\[
 L_i=|n_i\rangle\langle u_+|
       +\eta_i|n_*\rangle\langle m_i|.
\]
If a common source endomorphism \(A\) obeys
\(L_iA=r_iL_i\) for all \(i\), then all \(r_i\) are equal.
The same conclusion follows if a common target endomorphism \(B\)
obeys \(BL_i=r_iL_i\) for all \(i\). Thus the nonconstant rooted
Hopf phase cannot be absorbed into either one of these spinor legs alone.
\end{lemma}
\begin{proof}
Taking the \(n_i\)-component of the first equation on \(u_+\) gives
\(\langle u_+,Au_+\rangle=r_i\) for every \(i\).
For the second equation, apply it to \(m_i\) and divide by
\(\eta_i\ne0\); this gives \(Bn_*=r_in_*\) for every \(i\).
Either equation forces the three scalars to coincide.
Keeping only the family columns would discard the shared seed and
would not prove the first assertion.
\end{proof}

\begin{proposition}[Complete transverse return before occupation lowering]
\label{p1:prop:v242-complete-transverse-return}
Let \(A_j=T(a_j)|_{S_+\to C_-}\) and
\(B_j=T(b_j)|_{S_+\to C_-}\), \(j=2,3,4\), retaining both
creation and annihilation contributions. Set
\[
 \mathcal W_0=\begin{pmatrix}A_2\\A_3\\A_4\end{pmatrix},\qquad
 \mathcal W_+=\begin{pmatrix}A_2\\B_3\\-B_4\end{pmatrix}.
\]
Then
\[
 \boxed{\mathcal W_0^\dagger\mathcal W_0=3I_4,\qquad
 \mathcal W_0^\dagger\mathcal W_+
  =I_4-2i(N_3-N_4)
  =\operatorname{diag}(1,1,1-2i,1+2i).}
\]
Its family restriction in rooted relative units is
\[
 \boxed{\rho=\operatorname{diag}(1,1-2i,1+2i),\qquad
        \rho^\dagger\rho=\operatorname{diag}(1,5,5).}
\]
\end{proposition}
\begin{proof}
The full signed coefficient maps satisfy
\(A_j^\dagger A_j=I_4\) and
\(A_j^\dagger B_j=i(I_4-2N_j)\). The latter follows directly from
\((\varepsilon_j+\iota_j)i(\varepsilon_j-\iota_j)
=i(I-2\varepsilon_j\iota_j)\); the unitary Chevalley pairing cancels
in this product. Summing the three channel returns gives
\(I_4+i(I_4-2N_3)-i(I_4-2N_4)\).
The rooted occupations are
\(N_3=\operatorname{diag}(1,0,1,0)\) and
\(N_4=\operatorname{diag}(1,0,0,1)\), proving the formulas.
\end{proof}

\begin{proposition}[Signed transverse target-image and rooted-seed rank]
\label{p1:prop:signed-target-image-rank}
Fix the normalized signed-Fock source/target bases and the exact
transverse channel stacks $\mathcal W_0,\mathcal W_+$ of the existing
complete signed-channel return.  For any $U=(u_{ij})\in M_3(\mathbb C)$, using one-based matrix indices,
form the mathematical return
\[
\mathcal K_{I_4}(U)
 =\left.\mathcal W_0^\dagger(U\otimes I_4)\mathcal W_+\right|_G.
\]
Its matrix is exactly
\[
\begin{pmatrix}
 u_{11}+iu_{22}-iu_{33} & iu_{12}+u_{21} & iu_{13}-u_{31}\\
 iu_{12}-u_{21} & u_{11}-iu_{22}-iu_{33} & i(u_{32}-u_{23})\\
 iu_{13}+u_{31} & i(u_{32}-u_{23}) & u_{11}+iu_{22}+iu_{33}
\end{pmatrix}.
\]
Thus $\operatorname{rank}_{\mathbb C}\mathcal K_{I_4}=8$,
$\operatorname{im}\mathcal K_{I_4}=\{M:M_{32}=M_{23}\}$,
and $\ker\mathcal K_{I_4}=\mathbb C(E_{23}+E_{32})$.
For the existing signed-Fock target seed $P_*=|n_*\rangle\langle n_*|$,
\[
\mathcal K_{P_*}(U)
 =\left.\mathcal W_0^\dagger(U\otimes P_*)\mathcal W_+\right|_G
 = D_\eta U D_\eta R_+,
\quad D_\eta=\operatorname{diag}(1,-1,1),\quad
 R_+=\operatorname{diag}(1,-i,i).
\]
This map has rank nine, but maps every unitary $U$ to a unitary
family operator. No source-action occurrence follows from either
mathematical target contraction.
\end{proposition}
\begin{proof}
Direct multiplication of the frozen signed-Fock channel matrices gives
the displayed first return. Its second and third rows impose the single
linear relation $M_{32}=M_{23}$, and solving its kernel gives exactly
$\mathbb C(E_{23}+E_{32})$, hence rank eight. The rooted-seed
compression gives left/right multiplication by fixed unitary diagonal
matrices, an invertible linear map on $M_3(\mathbb C)$; it maps
unitaries to unitaries. These are exact finite statements in their
fixed source and target frames.
\end{proof}


\begin{corollary}[A traced quarter-turn does not cancel the complete rooted source Gram]
\label{p1:cor:v244-quarter-turn-full-gram}
On the fixed rooted source and target frames of
Proposition~\ref{p1:prop:v242-complete-transverse-return}, let
\(D=N_3-N_4=\operatorname{diag}(0,0,1,-1)\) and, for
\(t\in\mathbb R\), \(\theta\in\mathbb R\), put
\(\mathcal J_\theta(t)=\mathcal W_0+t e^{i\theta}\mathcal W_+\).
Then the \emph{complete} four-column source Gram is
\[
 \boxed{\mathcal J_\theta(t)^\dagger\mathcal J_\theta(t)
 =3(1+t^2)I_4+2t\cos\theta\,I_4+4t\sin\theta\,D.}
\]
In particular, at the relative quarter-turn \(\theta=\pi/2\),
\[
 \boxed{\mathcal J_{\pi/2}(t)^\dagger\mathcal J_{\pi/2}(t)
 =3(1+t^2)I_4+4tD,\qquad
 \operatorname{Tr}(\mathcal J_{\pi/2}^\dagger\mathcal J_{\pi/2})
 =12(1+t^2).}
\]
For \(t\ne0\) no \emph{single} phase \(\theta\) removes the
linear cross term on all four columns.  Cancellation in its trace
is not cancellation of the operator-valued family source.
\end{corollary}
\begin{proof}
The three signed rows of \(\mathcal W_+\) have unit column Grams,
so \(\mathcal W_+^\dagger\mathcal W_+=3I_4\), as for
\(\mathcal W_0\).  Insert
\(\mathcal W_0^\dagger\mathcal W_+=I_4-2iD\) and its adjoint.
The seed columns (\(D=0\)) require \(\cos\theta=0\) for
cancellation, while the two signed columns require
\(\sin\theta=0\); these conditions cannot hold together.
\end{proof}
The statement is a finite rooted phase comparison, not a physical
curvature--particle vertex.  Simultaneous transport of both spinor
legs and the source remains a change of presentation.

\begin{corollary}[Exact imaginary rooted coherence with neutral self-returns]
\label{p1:cor:v3030-rooted-imaginary-coherence}
On the fixed rooted four-column source of
Proposition~\ref{p1:prop:v242-complete-transverse-return}, set
\(D=N_3-N_4=\operatorname{diag}(0,0,1,-1)\) and
\(C=\mathcal W_0^\dagger\mathcal W_+=I_4-2iD\).
Then
\[
 \operatorname{Tr}(D\mathcal W_0^\dagger\mathcal W_0)
 =\operatorname{Tr}(D\mathcal W_+^\dagger\mathcal W_+)=0,
 \qquad \boxed{\operatorname{ImTr}(DC)=-4.}
\]
For the \emph{complete} source Gram
\(G_\theta(t)=\mathcal J_\theta(t)^\dagger\mathcal J_\theta(t)\)
of Corollary~\ref{p1:cor:v244-quarter-turn-full-gram},
\[
 \boxed{\operatorname{Tr}(DG_\theta(t))=8t\sin\theta.}
\]
For \(t\ne0\), the marked root return therefore recovers
\(\sin\theta\) even though the two individual self-Grams are
signed-neutral.  Under simultaneous changes of the rooted source
frame, transporting \(D\) and \(G_\theta\) by the same conjugation,
the displayed marked contraction is invariant.  It is \emph{not}
by itself a metric/source coupling or an independently measured
physical gravitational observable.
\end{corollary}
\begin{proof}
The preceding proposition and corollary give
\(\mathcal W_0^\dagger\mathcal W_0=
\mathcal W_+^\dagger\mathcal W_+=3I_4\),
\(C=I_4-2iD\), \(\operatorname{Tr}D=0\), and
\(\operatorname{Tr}D^2=2\). Thus
\(\operatorname{Tr}(DC)=-2i\operatorname{Tr}D^2=-4i\).
Multiplying the complete Gram identity by \(D\) and taking its
trace gives \(4t\sin\theta\operatorname{Tr}D^2=8t\sin\theta\).
Cyclicity proves invariance under common conjugation.
\end{proof}

\begin{proposition}[Quarter-turn samples do not select continuous depth transport]
\label{p1:prop:v3030-phase-sample-nonuniqueness}
Let \(q\in[0,7]\) be a \emph{comparison depth}, and assume only
the prescribed integer samples \(\theta(j)=\pi j/2\pmod{2\pi}\)
for \(j=0,\ldots,7\). These samples do not fix the smooth
phase between levels or its first derivative at \(q=6\).
Indeed, the two smooth nondecreasing transports
\[
 \theta_0(q)=\frac{\pi q}{2},\qquad
 \theta_1(q)=\frac{\pi q}{2}-\frac14\sin(2\pi q)
\]
have identical samples but
\[
 \boxed{(\sin\theta_0)'(6)=-\frac\pi2,
 \qquad (\sin\theta_1)'(6)=0.}
\]
In particular, the fixed-rooted phase comparison alone does not
identify \(q\) with physical curvature, select the linear
\(\theta(q)=\pi q/2\) transport, or predict a K48 transfer rate.
\end{proposition}
\begin{proof}
For each integer \(j\), \(\sin(2\pi j)=0\), so the sampled
values agree.  Direct differentiation gives
\(\theta_0'=\pi/2\) and
\(\theta_1'=\frac\pi2[1-\cos(2\pi q)]\ge0\), with
\(\theta_1'(6)=0\).  Since \(\cos\theta_0(6)=-1\), the
chain rule gives the displayed distinct derivatives.
\end{proof}

\begin{remark}[Relative frames, normalization and physical scope]
\label{p1:rem:v242-phase-lift-scope}
The vector lift is unique for the specified phase output with the rooted
spinor frames held fixed. A common Spin(8) transport still preserves the
whole tensor. In particular the exterior torus
\(g=\exp[\frac{\pi i}{2}(N_3-N_4)]\) implements the displayed vector
rotation; its two spinor actions must also be retained in a common-frame
comparison, where they cancel the relative effect. The one-leg lemma
does not prohibit that paired Ward identity.

Dividing each stacked map by \(\sqrt3\) changes its cross return to
\(\rho/3\). The displayed \(\rho\) uses the unchanged root entry as
relative unit. Its value \(5\) is a Gram value on a rank-two family
support; it does not identify this support with a rank-five counting
owner. These are geometric phase-port and return results. They do not
assert that a physical mass-source action inserts this port or fixes a
particular positive-path composition or dimensionful coefficient.
\end{remark}




\setlength{\emergencystretch}{2em}
\section{Protected coordinate variance and covariance}
\label{p1:sec:v2201-geometry-first-consequences}
\paragraph{Geometry ancestry.}
Marked round Hopf geometry $\to$ protected Clifford coordinate class $\to$ state variance.
\paragraph{Status.}
These are consequences on the domains stated below.
Every restriction retains the specified carrier, inner product,
coordinate class and operator domain of its inputs.

\subsection{Sharp total variance of the protected coordinates}
\label{p1:sec:v2201-derived-d02}
\noindent\textit{Status: proved result on stated domain.}

On a four-dimensional protected Clifford block let $X_a=c\Gamma_a$,
$c>0$, $a=1,\ldots,4$, with
$\{\Gamma_a,\Gamma_b\}=2\delta_{ab}I$.  The coefficient $c$
is the coefficient of each normalized generator, rather than the total
coordinate radius $\rho_{\rm prot}$ of the Paper-I-A representative.
For every density operator $\sigma$,
\[
 \boxed{\sum_{a=1}^4\operatorname{Var}_\sigma(X_a)
 \ge3c^2.}
\]
\begin{proof}
Put $r_a=\operatorname{Tr}(\sigma\Gamma_a)$. Since
$(v\cdot\Gamma)^2=|v|^2I$, choosing $v=r/|r|$ gives $|r|\le1$.
The variance sum is $c^2(4-|r|^2)$. The density
$(I+v\cdot\Gamma)/4$ for unit $v$ attains equality in the irreducible
four-dimensional block.
\end{proof}
For the coordinates of Theorem~\ref{p1:thm:s3one-cliffordization},
$\widehat X_a=(\rho_{\rm prot}/2)\Gamma_a$, and hence the bound is
$3\rho_{\rm prot}^2/4$.  On the distinct intrinsic tangential class the
protected coefficient is $c=\sqrt3\Rmom/4$, so its protected bound is
$9\Rmom^2/16$.  These are two substitutions into the same theorem;
they do not identify the two coordinate classes.

\begin{proposition}[Exact protected coordinate covariance]
\label{p1:prop:protected-coordinate-covariance}
With $X_a=c\Gamma_a$ as above, set
$\mu_a=\operatorname{Tr}(\sigma X_a)$ and
\[
 V_{ab}:=\tfrac12\operatorname{Tr}
 \bigl(\sigma\{X_a-\mu_aI,X_b-\mu_bI\}\bigr).
\]
Then the full real symmetric covariance matrix and its spectrum are
\[
 \boxed{V=c^2I_4-\mu\mu^{\mathsf T}},\qquad
 \operatorname{spec}V=\{c^2,c^2,c^2,c^2-|\mu|^2\},\qquad |\mu|\le c.
\]
Every $\mu$ in this ball is realized by a density operator.
\end{proposition}
\begin{proof}
The anticommutator gives $V_{ab}=c^2\delta_{ab}-\mu_a\mu_b$.
The preceding norm argument gives $|\mu|\le c$; conversely
$\sigma=(I+(\mu/c)\cdot\Gamma)/4$ is positive, has trace one, and has
the prescribed means.  The rank-one matrix $\mu\mu^{\mathsf T}$ has
eigenvalues $|\mu|^2,0,0,0$, proving the spectrum.
\end{proof}
The bound depends only on the Clifford relations. It does not identify
$\rho_{\rm prot}$ with the sphere radius or with the intrinsic tangential
coordinate coefficients; those classes have different square fingerprints.
\par\noindent\textit{Inherited inputs:} \ref{p1:thm:s3one-cliffordization}, \ref{p1:thm:intrinsic-tangential-coordinate-class}.




\section{Euler forms, weighted variations and spin-frame lifts}
\label{p1:sec:v2203-euler-work-geometry}

This section extends the already constituted spin-frame Hopf descendant of
Corollary~\ref{p1:cor:spinframe-hopf-shape-descendant}. All objects here are
defined by their orientations, connection bundles, weights and pullback
maps; these data remain explicit hypotheses throughout.

\begin{proposition}[Euler form and the primitive spin lift]
\label{p1:prop:v2203-euler-spin-lift}
Let \(\Sigma\cong S^2\) carry a smooth oriented metric \(h\) and its
adapted tangent spin-frame bundle \(\pi:P^{\rm spin}_\Sigma\to\Sigma\).
Use the positive orientation and primitive fibre normalization of
Corollary~\ref{p1:cor:spinframe-hopf-shape-descendant}. If \(K_h\) is Gaussian
curvature, \(R_h=2K_h\), and \(\eta_h\) is the matching induced spin
connection with \(\int_{S^1}\eta_h=2\pi\), then
\begin{equation}
\label{p1:eq:v2203-euler-normalization}
 e_h=\frac{K_h}{2\pi}\,\mathrm dA_h
     =\frac{R_h}{4\pi}\,\mathrm dA_h,
 \qquad c_{1,h}=\frac{e_h}{2},
 \qquad \mathrm d\eta_h=\pi^*(\pi e_h),
 \qquad \int_\Sigma e_h=2.
\end{equation}
The dependent lift and its primitive normalization are
\[
 J_{E,h}=\frac{\eta_h}{2\pi}\wedge\pi^*e_h
 =\frac{\eta_h\wedge\mathrm d\eta_h}{2\pi^2}
 =2\Omega_{{\rm spin},h},\qquad
 \pi_!J_{E,h}=e_h,\qquad \int_PJ_{E,h}=2,
\]
where \(\Omega_{{\rm spin},h}=\eta_h\wedge\mathrm d\eta_h/(4\pi^2)\)
has integral one. At the matching round representative,
\(e_h=\mathrm d\Omega_2/(2\pi)\) and
\(\Omega_{{\rm spin},h}\) is the atlas's primitive Hopf representative.
\end{proposition}
\begin{proof}
The existing spin-frame result fixes \(2c_1=e(T\Sigma)\), because the
\(\operatorname{Spin}(2)\to SO(2)\) covering doubles the circle angle.
Chern--Weil evaluation for the matching Levi-Civita connections gives
\eqref{p1:eq:v2203-euler-normalization} for every smooth metric on this
fixed oriented bundle. Fibre integration gives the lift identities.
\end{proof}
Normalizing the Euler form or lift by its total charge therefore returns
\(e_h/2=c_{1,h}\) or \(J_{E,h}/2=\Omega_{{\rm spin},h}\). The integer two
does not change an independently declared primitive current
\(Q_D\Omega_H\) with \(Q_D=1\). For a separately deformed metric, equality
to the previously fixed round representative is not a pointwise identity
unless the connection is transported as specified above. The signed total
also does not guarantee \(K_h\ge0\) pointwise.

\begin{theorem}[Projector curvature needs an orientation insertion]
\label{p1:thm:v2203-projector-euler}
On the oriented unit sphere in \(\mathbb R^3\), put
\(E_u=u^\perp\), \(P(u)=I-u\otimes u\), and \(J_uv=u\times v\).
The projected connection \(\nabla^E=P\mathrm d\) has curvature
\begin{equation}
\label{p1:eq:v2203-projector-pfaffian}
 \mathcal F^E=P(\mathrm dP\wedge\mathrm dP)P=-\omega_2J_u,
 \qquad \omega_2=\tfrac12\epsilon_{ijk}u_i\mathrm du_j\wedge\mathrm du_k,
 \qquad e(E)=\frac{\operatorname{Tr}_{E}(J_u\mathcal F^E)}{4\pi}
             =\frac{\omega_2}{2\pi}.
\end{equation}
In particular \(\operatorname{Tr}_{E}\mathcal F^E=0\); the Euler
contraction is the oriented rank-two Pfaffian, not the ordinary real trace.
\end{theorem}
\begin{proof}
For tangent parameter vectors \(X,Y\), write \(a=\mathrm du(X)\) and
\(b=\mathrm du(Y)\). Differentiating \(P\) and multiplying gives
\(\mathcal F^E(X,Y)v=a\langle b,v\rangle-b\langle a,v\rangle
=-\omega_2(X,Y)J_uv\). Since \(J_u^2=-I_E\), the oriented contraction
has trace \(2\omega_2\); the uninserted skew curvature has trace zero.
Integration of the unit sphere's area gives Euler number two.
\end{proof}

\begin{corollary}[Degree, orientation and moving-projector transgression]
\label{p1:cor:v2203-euler-degree-transgression}
For an actually supplied smooth map \(u:B\to S^2\) from a closed oriented
surface, \(\int_B e(u^*E)=2\deg(u)\). A fixed axis has zero pulled-back
curvature, despite the plane's rank two. For a smooth variation
\(s=\delta u\perp u\), set
\(\beta_s(X)=u\cdot(s\times\mathrm du(X))\). Then
\[
 \delta e(u^*E)=\mathrm d\xi_s,\qquad \xi_s=\beta_s/(2\pi),
 \qquad
 \delta\int_Bw e=\int_B\delta w\,e-\int_B\mathrm dw\wedge\xi_s.
\]
If \(B\) has boundary, add \(\int_{\partial B}w\xi_s\).
\end{corollary}
\begin{proof}
The degree identity is the pullback integral of \(\omega_2\).
Cartan's formula for this closed area form gives
\(\delta(u^*\omega_2)=\mathrm d\beta_s\); integration by parts proves
the weighted identity.
\end{proof}
The projector is unchanged by \(u\mapsto-u\), whereas \(J_u\) changes
sign. The universal plane over the unoriented axis space \(\mathbb{RP}^2\)
has nonzero first Stiefel--Whitney class and no ordinary global oriented
Euler two-form. An oriented lift is required; a finite plane type or rank
does not supply its degree or continuous connection. This does not exclude
the differently typed Euler class with local orientation coefficients.

\begin{theorem}[Complete weighted Euler metric variation]
\label{p1:thm:v2203-weighted-euler-work}
On a fixed closed oriented \(\Sigma\cong S^2\), let
\(\mathcal E[w,h]=\int_\Sigma we_h\) for a smooth independent real weight
\(w\). With \(\nabla^2=\nabla_i\nabla^i\), whose spectrum is nonpositive,
and \(\Delta_+=-\nabla^2\), its complete first variation is
\begin{equation}
\label{p1:eq:v2203-weighted-euler-variation}
 \boxed{\delta\mathcal E
 =\int_\Sigma e_h\delta w
 +\int_\Sigma\mathcal C^E_{ij}[w]\delta h^{ij}\,\mathrm dA_h,
 \qquad
 \mathcal C^E_{ij}[w]
 =\frac{h_{ij}\nabla^2w-\nabla_i\nabla_jw}{4\pi}.}
\end{equation}
Equivalently,
\(\mathcal C^E_{ij}=-(h_{ij}\Delta_+w+\nabla_i\nabla_jw)/(4\pi)\).
For \(h_\phi=e^{2\phi}h_0\), the finite conformal identity is
\begin{equation}
\label{p1:eq:v2203-conformal-euler}
 e_{h_\phi}=e_{h_0}-\frac{\nabla_0^2\phi}{2\pi}\,\mathrm dA_0.
\end{equation}
\end{theorem}
\begin{proof}
The variation of \(\int wR_h\mathrm dA_h\), after integration by parts,
is the weighted Einstein tensor plus
\((h_{ij}\nabla^2w-\nabla_i\nabla_jw)\delta h^{ij}\).
The two-dimensional Einstein tensor vanishes identically. Divide by
\(4\pi\) and add the independent weight derivative. The conformal
curvature formula
\(K_{e^{2\phi}h_0}\mathrm dA_{e^{2\phi}h_0}
=(K_0-\nabla_0^2\phi)\mathrm dA_0\)
proves the second statement.
\end{proof}
For constant \(w\), the geometric work tensor is zero pointwise.
Its divergence is
\(\nabla^i\mathcal C^E_{ij}=-(K_h/4\pi)\nabla_jw\), precisely the
weight contribution required by simultaneous diffeomorphism variation.
The fixed-weight conformal functional is linear in \(\phi\); its pure
second conformal derivative vanishes, while its weight--shape derivative
need not. If \(h=X^*g\) and the embedding moves by \(V\), insert
\(\delta h=X^*(\delta g+\mathcal L_Vg)\) and the actual \(\delta w\)
in \eqref{p1:eq:v2203-weighted-euler-variation}. No induced-metric,
embedding, or weight derivative is removed by topological protection.

\begin{corollary}[Conformal weight--shape leaf]
\label{p1:cor:v2203-conformal-weight-leaf}
On the unit round sphere let \(h_\sigma=e^{-2\sigma}h_0\) and
\(w_\sigma=e^\sigma\). The geometric weighted port obeys
\[
 \mathcal E[w_\sigma,h_\sigma]
 =\frac1{2\pi}\int e^\sigma(1-\Delta_+\sigma)\mathrm d\Omega_2
 =\frac1{2\pi}\int e^\sigma(1-|\nabla\sigma|^2)\mathrm d\Omega_2.
\]
Its first derivative at zero is \(\int s\mathrm d\Omega_2/(2\pi)\), and
\[
 D^2\mathcal E|_0[s_1,s_2]
 =\frac1{2\pi}\int s_1(1-2\Delta_+)s_2\,\mathrm d\Omega_2.
\]
\end{corollary}
\begin{proof}
Use \eqref{p1:eq:v2203-conformal-euler} with \(\phi=-\sigma\), integrate
\(e^\sigma\Delta_+\sigma\) by parts, and differentiate twice.
\end{proof}
This is a weighted geometric Hessian. Its sign is not a stability claim
about a later complete action. The nonzero constant-mode derivative comes
from the varying weight, not from a variation of Euler number.

\begin{theorem}[Lift transgression and conserved geometric charge]
\label{p1:thm:v2203-lift-transgression}
Identify a smooth family of the fixed oriented spin-frame bundles. Let
\(\delta e_h=\mathrm d\xi\) and let \(\delta\eta_h\) be the basic
variation of the normalized principal connection. Then
\begin{equation}
\label{p1:eq:v2203-lift-transgression}
 \delta J_{E,h}=-\mathrm d\mathcal T,
 \qquad \mathcal T=\frac{\eta_h}{2\pi}\wedge\pi^*\xi.
\end{equation}
For a fixed region \(\mathcal D\subseteq P\),
\[
 \delta\int_{\mathcal D}wJ_E
 =\int_{\mathcal D}\delta wJ_E
 +\int_{\mathcal D}\mathrm dw\wedge\mathcal T
 -\int_{\partial\mathcal D}w\mathcal T.
\]
All integrated geometric jets of \(\int_PJ_E=2\) vanish on the closed
total space. For the matching induced spin connection one may take
\(\pi^*\xi=\delta\eta_h/\pi\), so
\(\mathcal T=\eta_h\wedge\delta\eta_h/(2\pi^2)\).
\end{theorem}
\begin{proof}
The term \(\delta\eta_h\wedge\pi^*e_h\) vanishes because it contains
three horizontal factors over a two-dimensional base. The same dimension
argument removes \(\mathrm d\eta_h\wedge\pi^*\xi\); differentiating
\(\eta_h\wedge\pi^*\xi\) proves
\eqref{p1:eq:v2203-lift-transgression}. Stokes' theorem proves the weighted
and charge statements. A moving identification adds a Lie derivative of
the closed three-form, hence another exact transgression.
\end{proof}

\begin{corollary}[Pointwise homothety invariance]
\label{p1:cor:v2203-euler-homothety}
If \(h_r=c(r)h_0\) with \(c(r)>0\) spatially constant and with the
natural oriented frame identification, then
\[
 e_{h_r}=e_{h_0},\qquad \eta_{h_r}=\eta_{h_0},\qquad
 J_{E,h_r}=J_{E,h_0}
\]
pointwise. Every derivative in a regular reparameterization of this
homothety family is zero. For every fixed weight, homogeneous metric
scale work is consequently zero, even if the weight is nonuniform.
\end{corollary}
\begin{proof}
Gaussian curvature and area scale reciprocally; a spatially constant
homothety preserves the Levi-Civita connection under the natural frame
identification. The lift then remains unchanged. This also follows for
the primitive shape metric from
\eqref{p1:eq:primitive-shape-normalization}.
\end{proof}

\begin{proposition}[The protected spectral two has a specified operator owner]
\label{p1:prop:v2203-euler-spectral-owner}
For the full de Rham complex on this closed sphere,
\[
 \operatorname{Str}_{\Lambda^*T^*\Sigma}e^{-t\Delta_{\rm dR}}
 =b_0-b_1+b_2=2\qquad(t>0).
\]
The ordinary scalar heat trace instead tends to \(b_0=1\) as
\(t\to\infty\). On a positive-curvature round spin sphere, the untwisted
spin Dirac operator has no zero modes.
\end{proposition}
\begin{proof}
On each positive eigenspace, \(\mathrm d+\mathrm d^\dagger\) is an
invertible parity-reversing map, so the even and odd heat contributions
cancel. Harmonic forms leave \(b_0=b_2=1,b_1=0\).
The scalar limit retains only constants. The spin claim follows by
integrating \(D^2=\nabla^*\nabla+R_h/4\) with \(R_h>0\).
\end{proof}

\begin{corollary}[Tangent-line charge gap is not a doubled scalar impedance]
\label{p1:cor:v2203-tangent-charge-gap}
On the atlas's round radius-\(\Rmom\) Hopf host, the fixed integer charge
\(q\) horizontal spectrum is
\[
 \lambda_{H;k,q}=\frac{k(k+2)-q^2}{\Rmom^2},
 \qquad k\ge|q|,\quad k\equiv q\pmod2,
 \qquad \operatorname{mult}(k,q)=k+1.
\]
The weight-two tangent-line field therefore has lowest horizontal
eigenvalue \(4/\Rmom^2\) and multiplicity three. The invariant scalar
Euler-density coefficient is a different \(q=0\) owner and has a constant
mode with eigenvalue zero. If a separately supplied product collar actually
acts on these sections by
\(-\partial_s^2+C_7+\Delta_H\), with \(C_7\ge0\), its lowest normal
spectral root in the tangent sector is
\[
 \sqrt{C_7+4/\Rmom^2},
\]
not a universal \(2\sqrt{C_7}\). Its finite-collar boundary impedance
also retains its actual length and endpoint conditions.
\end{corollary}
\begin{proof}
The fixed-charge decomposition of
Theorem~\ref{p1:thm:scalar-spectrum-by-hopf-charge} and
Proposition~\ref{p1:prop:hopf-charge-line-bundle}, with the vertical Casimir
removed as in Section~\ref{p1:subsec:sub-laplacian-spectrum}, gives the
displayed spectrum. Its lowest allowed level is \(k=|q|\), giving
\(2|q|/\Rmom^2\) and multiplicity \(|q|+1\). Set \(q=2\) or \(q=0\)
and take the positive normal root of the stated product operator.
\end{proof}
The collar operator acts on associated tangent-line sections, whereas
the scalar Euler-density coefficient belongs to the charge-zero sector.
The untwisted basic Dolbeault package has
index one by Theorem~\ref{p1:thm:hopf-basic-dolbeault-index}; it is not the
full de Rham supertrace two proved above.

\begin{remark}[Distinct geometric and spectral contractions]
\label{p1:rem:v2203-euler-physical-scope}
A de Rham supertrace, an oriented curvature contraction, a circle-weight-two
associated representation, and an ordinary Dirac heat trace have different
domains and contractions.  The tangent line has degree two and the adapted
spin circle has degree one.  Restriction to a finite rotation group does not
recover a continuous circle weight or connection.  The ordinary real
two-plane curvature trace is zero; its positive curvature-square trace
differs from its Euler Pfaffian.  Each comparison therefore retains its
specified bundle, connection and pullback map.
\end{remark}


\section{Further exact harmonic and operator comparisons}
\label{p1:sec:v215-geometric-provider-completion}

\subsection{Exact norms, duality, and higher cubature errors of the Hurwitz sextic}
\label{p1:sec:v215-sextic-providers}

\begin{proposition}[Sextic duality and exact Haar moments]
\label{p1:prop:v215-sextic-moments}
Let \(H_6\) be the homogeneous harmonic polynomial of
Proposition~\ref{p1:prop:math-hurwitz-cubature-six-defect}, and let
\(E\) denote probability Haar integration on \(S^3(1)\). The
orientation-preserving orthogonal transformation
\[
 O(x_1,x_2,x_3,x_4)=\frac1{\sqrt2}
 (x_1-x_2,x_1+x_2,x_3-x_4,x_3+x_4)
\]
satisfies the polynomial identity \(H_6\circ O=-H_6\). Consequently
\[
 E H_6^{2k+1}=0\quad(k\geq0),\qquad
 E H_6^2=\frac17,\qquad E H_6^4=\frac{61}{1001}.
\]
In particular, the probability-normalized harmonic \(Y_6=\sqrt7H_6\)
has
\[
 \Delta_+Y_6=48Y_6,\qquad
 E Y_6=E Y_6^3=0,\qquad
 E Y_6^2=1,\qquad E Y_6^4=\frac{427}{143}.
\]
Here \(\Delta_+\) is the positive round Laplacian. On \(S^3(a)\),
with \(V_a=2\pi^2a^3\), regard \(Y_6(x/a)\) as the corresponding
dimensionless profile and put \(\psi_6(x)=Y_6(x/a)/\sqrt{V_a}\). Then
\[
 \int\psi_6^2\,d\mathrm{vol}=1,\qquad
 \int|\nabla\psi_6|^2\,d\mathrm{vol}=\frac{48}{a^2},\qquad
 \int\psi_6^3\,d\mathrm{vol}=0,\qquad
 \int\psi_6^4\,d\mathrm{vol}=\frac{427}{143V_a}.
\]
\end{proposition}
\begin{proof}
Write \(y_i=x_i^2\). Direct substitution in the displayed sextic gives
\[
 H_6(x_1-x_2,x_1+x_2,x_3-x_4,x_3+x_4)=-8H_6(x).
\]
Homogeneity proves the stated identity, and the two planar rotation
blocks of \(O\) each have determinant one. Haar invariance now proves
every odd-moment identity. For an explicit rational certificate for the
even moments, on \(\sum_i y_i=1\) one has
\[
 H_6=1-8e_2(y)+48e_3(y),\qquad
 E\prod_{i=1}^4x_i^{2a_i}
 =\frac{\prod_i(2a_i)!}
 {4^{|a|}(|a|+1)!\prod_i a_i!}.
\]
The latter follows, for example, by normalizing four independent squared
standard real Gaussians; their normalized squares have Dirichlet
parameters \((1/2,1/2,1/2,1/2)\). It gives the following complete table
needed for the fourth power:
\[
\begin{array}{c|ccccc}
 E(e_2^j e_3^k)&k=0&k=1&k=2&k=3&k=4\\\hline
 j=0&1&1/48&3/4480&11/430080&29/26836992\\
 j=1&1/4&1/160&23/107520&1/118272&\\
 j=2&11/160&69/35840&131/1892352&&\\
 j=3&89/4480&13/21504&&&\\
 j=4&85/14336&&&&
\end{array}
\]
Inserting these entries into the finite identity
\[
 E H_6^n=
 \sum_{j+k\leq n}\frac{n!(-8)^j48^k}{j!k!(n-j-k)!}
                  E(e_2^j e_3^k)
 \qquad(n=2,4)
\]
gives \(1/7\) and \(61/1001\), respectively. A harmonic homogeneous
polynomial of degree six restricts to eigenvalue \(6(6+2)=48\).
Integration by parts and metric and volume scaling give the last four
identities.
\end{proof}

\begin{proposition}[Degree-six addition and Dirichlet kernels]
\label{p1:prop:v215-degree-six-kernels}
Let \(\{Y_A\}_{A=1}^{49}\) be any real probability-orthonormal basis of
the degree-six scalar harmonic space on \(S^3(a)\). Then, pointwise,
\[
 \sum_{A=1}^{49}Y_A(x)^2=49,\qquad
 \sum_{A=1}^{49}(\nabla_iY_A)(\nabla_jY_A)
                 =\frac{784}{a^2}g_{ij}.
\]
More generally, if real coefficients \(c_A\) have an uncentered second
moment \(E(c_Ac_B)=s^2\delta_{AB}\), the random function
\(F=\sum_Ac_AY_A\) has
\[
 E F(x)^2=49s^2,\qquad
 E(\nabla_iF\nabla_jF)=\frac{784s^2}{a^2}g_{ij}.
\]
These are second-moment identities; they impose no fourth-moment
condition on the coefficients.
\end{proposition}
\begin{proof}
The harmonic-space dimension is \((6+1)^2=49\). The reproducing kernel
is invariant under the full round isometry group, so its diagonal is
constant, equal to its probability integral \(49\). Likewise its
diagonal differentiated tensor is invariant under the point stabilizer
\(SO(3)\), hence is a constant multiple of \(g\). Its trace has
probability integral \(49\cdot48/a^2\). Dividing by the tangent
dimension three gives \(784/a^2\). Contracting the basis sums with the
specified second moment proves the remaining assertions.
\end{proof}

\begin{proposition}[The degree-eight harmonic and exact low-power aliases]
\label{p1:prop:v215-eight-alias}
Let \(E_T\) denote the uniform average on the Hurwitz support \(2T\),
and put
\[
 \Phi_8(u)=\sum_{q\in2T}(q\cdot u)^8,\qquad
 H_8(u)=\frac{32}{3}\Phi_8(u)-14|u|^8-7|u|^2H_6(u).
\]
Then \(H_8\) is a homogeneous harmonic polynomial of degree eight,
invariant under the full \(24\)-cell symmetry, and \(H_8(q)=1\) for
\(q\in2T\). For a unit vector \(u\), define
\(f_u(x)=4(u\cdot x)^2-1\). The following are exact identities:
\[
 \begin{split}
 E_T(u\cdot x)^6-E(u\cdot x)^6&=\frac{H_6(u)}{64},\\
 E_T(u\cdot x)^8-E(u\cdot x)^8
                         &=\frac{7H_6(u)+H_8(u)}{256},\\
 E_T f_u^3-E f_u^3&=H_6(u),\\
 E_T f_u^4-E f_u^4&=3H_6(u)+H_8(u).
 \end{split}
\]
In particular \(Ef_u^3=1\) and \(Ef_u^4=3\). Uniform averaging on
\(2I\), instead, equals \(E\) on both \(f_u^3\) and \(f_u^4\).
For the scalar exponential one has the local expansion
\[
 E_Te^{-t(u\cdot x)^2}-Ee^{-t(u\cdot x)^2}
                  =-\frac{t^3}{384}H_6(u)+O(t^4).
\]
The last formula concerns this specified multiplication function, with
fixed \(u\); it does not identify a differential heat kernel.
\end{proposition}
\begin{proof}
The sextic orbit identity already gives the first line. In four
dimensions,
\[
 \Delta\Phi_8=56\Phi_6,\qquad
 \Delta|u|^8=80|u|^6,\qquad
 \Delta(|u|^2H_6)=32H_6.
\]
Using \(\Phi_6=(15/8)|u|^6+(3/8)H_6\), these terms cancel in
\(\Delta H_8\). The orbit sum and its harmonic subtraction are
full-cell invariant. At a coordinate vertex \(e_1\),
\(\Phi_8(e_1)=2+16/256=33/16\) and \(H_6(e_1)=1\), whence
\(H_8(e_1)=1\); transitivity gives the other vertices. The Haar
monomial rule above gives \(E(u\cdot x)^8=7/128\), so the
definition of \(H_8\) gives the second line. Lower moments through
degree five are exact on \(2T\). Expanding
\((4(u\cdot x)^2-1)^3\) and its fourth power gives the two alias
formulas and the Haar values. Their degrees are six and eight, so the
eleven-design \(2I\) integrates them exactly. The exponential agrees
through order two; its order-three difference is
\(-t^3H_6/(3!64)\). On the compact sphere the exponential Taylor
remainder is bounded uniformly for \(t\) near zero.
\end{proof}

\begin{theorem}[Odd-profile spectral parity]
\label{p1:thm:v215-odd-profile-spectral-parity}
Let \(A(q)=A_0+qB\) be a self-adjoint operator family on a common
domain, and let a unitary \(U\) preserve that domain and satisfy
\(UA_0U^{-1}=A_0\), \(UBU^{-1}=-B\). Then
\[
 UA(q)U^{-1}=A(-q),\qquad
 U F(A(q)^2)U^{-1}=F(A(-q)^2)
\]
for every functional calculus defined on the two operators. In finite
dimension, or whenever the displayed trace is defined,
\[
 Z(q):=\operatorname{Tr}F(A(q)^2)=Z(-q).
\]
Every existing odd derivative of \(Z\) at zero vanishes. The same
holds for a trace-class localized expression
\(\operatorname{Tr}(W F(A(q)^2))\) if \(UWU^{-1}=W\). The same parity
holds for a finite compressed family \(PA(q)P\) if \(P\) commutes with
\(U\) and its range lies in the common operator domain. In particular,
pullback by \(O\) conjugates multiplication by \(H_6\) to its negative.
\end{theorem}
\begin{proof}
Conjugate the affine family. Unitary covariance of the spectral measure
gives the functional-calculus identity; unitary invariance and cyclicity
of each defined trace give the evenness assertions. Odd derivatives of
an even function vanish whenever they exist. Finally,
\(H_6\circ O^{-1}=-H_6\) follows from
Proposition~\ref{p1:prop:v215-sextic-moments}, and proves the stated
multiplication identity. No positivity of even derivatives follows from
this argument.
\end{proof}

\subsection{Resolvent, functional calculus, and return on the minimal closure}
\label{p1:sec:v215-closure-return}

\begin{theorem}[Exact return identities on the twenty-five-dimensional closure]
\label{p1:thm:v215-closure-return}
Use the retained space \(\mathcal L_R\) and its minimal first-order
closure \(\widetilde{\mathcal L}_R\) from
Theorem~\ref{p1:thm:v214-first-order-retained-closure}. Identify the
original model with its isometric image, let \(P\) project onto it,
and write
\(\mathcal Q_R=\widetilde{\mathcal L}_R\ominus\mathcal L_R\).
Relative to \(\mathcal L_R\oplus\mathcal Q_R\),
\[
 \widetilde D=\begin{pmatrix}D_F&B\\ B^*&0\end{pmatrix},\qquad
 D_FB=0,\qquad
 L:=BB^*=K_{\rm LC}-D_F^2
       =\frac3{R^2}P_V+\frac4{R^2}P_{\Lambda^2}.
\]
Here \(P_V\) has rank four and \(P_{\Lambda^2}\) rank six;
\(P_{\rm sc}\) and \(P_{\rm sp}\) denote the remaining scalar and
spinor projections. For \(z\ne0\) outside the spectrum of
\(\widetilde D\),
\[
 P(z-\widetilde D)^{-1}P
       =\left(z-D_F-\frac{L}{z}\right)^{-1}
          \quad\hbox{on }\mathcal L_R.
\]
In particular, its vector and bivector restrictions are
\(z/(z^2-3/R^2)\) and \(z/(z^2-4/R^2)\), respectively.
For every function defined on the finite spectra,
\[
 P F(\widetilde D^2)P=F(K_{\rm LC})
       \quad\hbox{on }\mathcal L_R,
\]
and
\[
 \operatorname{Tr}_{25}F(\widetilde D^2)
 -\operatorname{Tr}_{15}F(K_{\rm LC})
            =4F(3/R^2)+6F(4/R^2).
\]
For a real parameter \(t\), the exact compressed unitary is
\[
 P e^{-it\widetilde D}P
 =P_{\rm sc}+\cos(\sqrt3t/R)P_V
    +\cos(2t/R)P_{\Lambda^2}+e^{-itD_F}P_{\rm sp}.
\]
The squared norm transferred from \(\mathcal L_R\) to its partner
space, averaged in the normalized density \(P/15\), is
\[
 \operatorname{Tr}\left[(1-P)e^{-it\widetilde D}
           \frac{P}{15}e^{it\widetilde D}(1-P)\right]
 =\frac4{15}\sin^2(\sqrt3t/R)+\frac6{15}\sin^2(2t/R).
\]
\end{theorem}
\begin{proof}
The ten partner modes in the closure theorem are four scalar first
harmonics and six Killing one-forms. The first-order operator exchanges
each of these with its original exact one-form or two-form partner.
Explicitly, with coordinate restrictions \(x_a\),
\(\alpha_{ab}=x_a\,dx_b-x_b\,dx_a\), and
\(\beta_{ab}=dx_a\wedge dx_b\),
\[
 D x_a=i\,dx_a,\qquad D(dx_a)=-\frac{3i}{R^2}x_a,
 \qquad D\alpha_{ab}=2i\beta_{ab},\qquad
 D\beta_{ab}=-\frac{2i}{R^2}\alpha_{ab}.
\]
Indeed the coordinate Hessian is
\(\nabla dx_a=-(x_a/R^2)g\), which gives
\(\delta dx_a=3x_a/R^2\), \(\delta\alpha_{ab}=0\), and
\(\delta\beta_{ab}=2\alpha_{ab}/R^2\); also
\(d\alpha_{ab}=2\beta_{ab}\) and \(d\beta_{ab}=0\).
Their squared norms in probability volume are respectively
\(R^2/4,3/4,R^2/2,1/2\), giving the pair frequencies
\(\sqrt3/R\) and \(2/R\). The constant mode is already present,
so the scalar enlargement has dimension four rather than five.
It has no diagonal block on \(\mathcal Q_R\), and its compression
\(D_F\) vanishes on the ten original form modes coupled by \(B\).
This gives the block form and \(D_FB=0\). The already proved
compression leakage gives \(BB^*=L\). Hence
\[
 \widetilde D^2=
 \begin{pmatrix}K_{\rm LC}&0\\0&B^*B\end{pmatrix},\qquad
 \operatorname{spec}(B^*B)
       =(3/R^2)^{(4)}\oplus(4/R^2)^{(6)}.
\]
The Schur complement of the invertible block \(zI_{\mathcal Q_R}\)
in \(z-\widetilde D\) is \(z-D_F-L/z\), proving the resolvent
identity. The block diagonal square proves the functional-calculus
identity and trace difference. On a pair of square eigenvalue
\(\lambda>0\), choose orthonormal partner bases so that the
first-order block is
\(\sqrt\lambda\bigl(\begin{smallmatrix}0&1\\1&0\end{smallmatrix}\bigr)\).
Exponentiation gives cosine on the original component and a sine
of equal frequency on the partner component. The scalar and spinor
blocks are already invariant. This proves the compressed exponential;
summing the transferred squared norms over four vector and six
bivector basis vectors and dividing by fifteen proves the final line.
\end{proof}

\begin{corollary}[Preserved projected heat trace and the full-trace difference]
\label{p1:cor:v215-projected-heat-trace}
For \(t\geq0\),
\[
 \operatorname{Tr}_{25}\bigl(Pe^{-t\widetilde D^2}P\bigr)
 =\operatorname{Tr}_{15}e^{-tK_{\rm LC}},\qquad
 \operatorname{Tr}_{25}e^{-t\widetilde D^2}
 -\operatorname{Tr}_{15}e^{-tK_{\rm LC}}
 =4e^{-3t/R^2}+6e^{-4t/R^2}.
\]
Thus the closure represents every specified retained square readout
exactly with its retained projection. The unprojected trace is a
different mathematical readout with the displayed difference.
\end{corollary}
\begin{proof}
Apply Theorem~\ref{p1:thm:v215-closure-return} to \(F(s)=e^{-ts}\).
\end{proof}

\subsection{Higher invariant candidates on the complete joint orbit}
\label{p1:sec:v215-higher-joint-dependencies}

\begin{proposition}[Exact higher-power dependencies on the joint orbit]
\label{p1:prop:v215-higher-joint-dependencies}
Use the complete co-transported tuple and sign line
\(N_{\mathcal O}=L\otimes\ell\) of
Theorem~\ref{p1:thm:joint-hopf-axial-orbit}, and let
\(m=m_{\mathcal O}\in\Gamma(N_{\mathcal O}^{*\otimes2})\) be its
quadratic section. At the reference point use
\[
 L_0=\mathbb C e_+,\quad e_+=(1,0),\qquad
 h_0=(1+\mathbf i+\mathbf j+\mathbf k)/2,\qquad
 \mu=m_0(e_+\otimes h_0,e_+\otimes h_0)=\frac{-1+i}{2}.
\]
For a spin lift \(g\), co-transport the polynomials of
Proposition~\ref{p1:prop:higher-invariant-pencil} by
\(f_{n,g}(s)=f_n(g^{-1}s)\), and likewise for \(g_{12}\).
Put \(h_g=gh_0g^{-1}\) and let \(\lambda_{h_g}\in\ell_g^*\)
be the complex-linear covector with \(\lambda_{h_g}(h_g)=1\).
The following even-degree candidate sections descend to \(\mathcal O\):
\[
 \begin{split}
 M_{12}^{(u)}&=(f_{6,g}|_{L_g})^2\otimes\lambda_{h_g}^{\otimes12},\\
 M_{12}^{(g)}&=(g_{12,g}|_{L_g})\otimes\lambda_{h_g}^{\otimes12},\\
 M_{16}&=(f_{8,g}|_{L_g})^2\otimes\lambda_{h_g}^{\otimes16}.
 \end{split}
\]
On this same complete joint orbit they satisfy exactly
\[
 M_{12}^{(u)}=0,\qquad M_{12}^{(g)}=-8i\,m^6,
 \qquad M_{16}=16m^8.
\]
If \(g_{12}\) and \(f_8\) are replaced by their unit-norm
symmetric-tensor representatives, the latter two coefficients are
\(-i\sqrt{10}\) and \(10/3\), respectively.
\end{proposition}
\begin{proof}
The two lifts \(g\) and \(-g\) give the same transported polynomials
because their degrees are even, and give the same \(h_g\). The
sections therefore descend. In the temporary unit frame
\(e_g=(ge_+)\otimes h_g\), their coefficients are the reference
values \(f_6(e_+)^2=0\), \(g_{12}(e_+)=1\), and
\(f_8(e_+)^2=1\). The quadratic descent theorem gives
\(m_g(e_g,e_g)=\mu\), independent of \(g\). Direct multiplication
gives \(\mu^6=i/8\) and \(\mu^8=1/16\), proving all three
section identities in these frames, hence intrinsically. Equivalently,
each even dual power has a one-dimensional equivariant-section space,
so equality of the reference values determines equality on the orbit.
Finally \(\|g_{12}\|^2=32/5\) and \(\|f_8\|^2=24/5\).
The first normalization divides \(-8i\) by \(\sqrt{32/5}\);
the second divides \(16\) by \(24/5\), proving the printed constants.
\end{proof}

\begin{remark}[The base and transport are part of the identities]
The nowhere-zero \(m\) trivializes every even dual power of the joint
sign line, and these co-transported higher candidates are its displayed
powers rather than independent equivariant directions. The statement
does not hold for an absolute tensor held independently fixed while the
other inputs move, or for the distinct fixed-axis projective-line base.
In particular it does not trivialize the evaluation-image line
\(A_{12}\to\mathbb{CP}^1\), whose Chern class is \(-12\), without
the specified actual pullback. No equality of independently prescribed
projection connections is implied.
\end{remark}

\section{The coordinates generate their native operator envelope}
\label{p1:sec:v226-coordinates}

Work in an allowed sector-phase representative and put $R=R_{\rm mom}>0$,
$Y_a=X_a^D/R$. There is an orthogonal decomposition
\[
 \mathcal H=\mathbb{C} e_0\oplus V\oplus\Lambda^2V\oplus S,
 \qquad \dim(V,\Lambda^2V,S)=(4,6,4).
\]
With $W_av=e_a\wedge v$, the intrinsic matrices are
\begin{equation}\label{p1:eq:v226-matrix}
 Y_a=
 \begin{pmatrix}
 0&\frac12 e_a^*&0\\
 \frac12e_a&0&\frac{\sqrt2}{3}W_a^*\\
 0&\frac{\sqrt2}{3}W_a&0
 \end{pmatrix}
 \oplus \frac{\sqrt3}{4}\gamma_a,
 \qquad \{\gamma_a,\gamma_b\}=2\delta_{ab}I_4.
\end{equation}
This is the existing intrinsic tuple, not an additional representative choice
with new invariant coefficients. Write $P_0,P_V,P_B,P_S$ for the four
summand projectors. The existing fingerprint gives
\begin{equation}\label{p1:eq:v226-Q}
 Q:=\sum_aY_a^2=P_0+\frac{11}{12}P_V+\frac49P_B+\frac34P_S.
\end{equation}

\begin{theorem}[Coordinate-only native algebra]\label{p1:thm:v226-coordinate-only}
For the intrinsic branch,
\[
 \mathcal A_D:=\operatorname{alg}_\mathbb{C}^*(I,Y_1,Y_2,Y_3,Y_4)
 =\operatorname{End}(\mathbb{C} e_0\oplus V\oplus\Lambda^2V)\oplus\operatorname{End}(S).
\]
In particular $\mathcal A_{\rm kin}\subset\mathcal A_D$; adjoining the
kinematic algebra does not change this algebra.
\end{theorem}
\begin{proof}
The four values $1,11/12,4/9,3/4$ are distinct. Each of the four projectors
is therefore the Lagrange polynomial
\begin{equation}\label{p1:eq:v226-projector}
 P_\lambda=\prod_{\mu\ne\lambda}\frac{Q-\mu I}{\lambda-\mu}
\end{equation}
in $Q$, and belongs to $\mathcal A_D$.
The scalar/vector edge gives the actual rank-one operators
\[
 2P_VY_aP_0=|e_a\rangle\langle e_0|.
\]
Their adjoints and products generate all matrix units on $\mathbb{C} e_0\oplus V$.
For $a\ne b$ the next edge gives
\[
 \frac{3}{\sqrt2}P_BY_aP_V|e_b\rangle\langle e_0|
 =|e_a\wedge e_b\rangle\langle e_0|.
\]
These six independent bivectors, their adjoints and products supply every
matrix unit on the full $11$-dimensional geometric summand. On $S$, the
matrices $P_SY_aP_S=(\sqrt3/4)\gamma_a$ generate $M_4(\mathbb{C})$: the sixteen
Clifford monomials indexed by subsets of $\{1,2,3,4\}$ form a trace-orthogonal
basis of its endomorphisms. All $Y_a$ are block diagonal between geometry
and $S$, so no off-block matrix can be generated. Both inclusions follow.
\end{proof}

\begin{theorem}[Polynomial reconstruction of all round rotations]
\label{p1:thm:v226-rotations}
For $a<b$, set
\begin{equation}\label{p1:eq:v226-rotation}
 \boxed{T_{ab}=\frac{36}{17}P_V[Y_a,Y_b]P_V
               +\frac92P_B[Y_a,Y_b]P_B
               +\frac43P_S[Y_a,Y_b]P_S.}
\end{equation}
These are the existing skew-Hermitian infinitesimal $\operatorname{Spin}(4)$ generators:
zero on $e_0$, the standard rotations on $V$ and $\Lambda^2V$, and
$[\gamma_a,\gamma_b]/4$ on $S$. In particular
\[
 [T_{ab},Y_c]=\delta_{bc}Y_a-\delta_{ac}Y_b.
\]
They satisfy the $\mathfrak{so}(4)$ brackets. The compatible Lie-algebra
representation with this covariance is unique.
\end{theorem}
\begin{proof}
Let $R_{ab}e_c=\delta_{bc}e_a-\delta_{ac}e_b$ and let
$R_{ab}^{\wedge2}$ be its induced bivector action. The identity
$W_a^*W_b=\delta_{ab}I-|e_b\rangle\langle e_a|$ gives
\[
 P_V[Y_a,Y_b]P_V
 =\left(\frac14+\frac29\right)R_{ab}
 =\frac{17}{36}R_{ab},\qquad
 P_B[Y_a,Y_b]P_B=\frac29R_{ab}^{\wedge2}.
\]
The spinor block gives
$P_S[Y_a,Y_b]P_S=(3/16)[\gamma_a,\gamma_b]$.
Equation~\eqref{p1:eq:v226-rotation} cancels these coefficients exactly. The direct
representation has the stated covariance and Lie brackets.

If another skew-Hermitian representation $T'_{ab}$ has the same action on
all four $Y_c$, every $T'_{ab}-T_{ab}$ commutes with their algebra. By
Theorem~\ref{p1:thm:v226-coordinate-only}, this commutant is
$\mathbb{C} P_g\oplus\mathbb{C} P_S$, where $P_g=P_0+P_V+P_B$. Both representations
commute with these two projectors. Their difference, as a linear map from
$\mathfrak{so}(4)$ to its abelian scalar commutant, vanishes on all Lie
brackets. Since $\mathfrak{so}(4)$ is equal to its derived algebra, the
difference is zero.
\end{proof}

\begin{theorem}[Separate geometric duality and spinor chirality at degree six]
\label{p1:thm:v226-chirality}
Let $\chi=\gamma_1\gamma_2\gamma_3\gamma_4$ in the fixed orientation, and
let $\mathsf B=P_B*P_B$, $\mathsf S=P_S\chi P_S$, extended by zero.
For the fully antisymmetrized degree-four coordinate word,
\[
 \Omega_4=\sum_{\pi\in S_4}\operatorname{sgn}(\pi)
             Y_{\pi(1)}Y_{\pi(2)}Y_{\pi(3)}Y_{\pi(4)},
\]
one has the exact identity
\begin{equation}\label{p1:eq:v226-omega}
 \Omega_4=-\frac{34}{81}\mathsf B+\frac{27}{32}\mathsf S.
\end{equation}
Consequently both operators are coordinate polynomials of degree at most six:
\begin{equation}\label{p1:eq:v226-chiral-recovery}
 \boxed{\mathsf S=\frac{128}{33}\left(Q-\frac49I\right)\Omega_4,
 \qquad
 \mathsf B=\frac{1458}{187}\left(Q-\frac34I\right)\Omega_4.}
\end{equation}
Thus $(P_B\pm\mathsf B)/2$ and $(P_S\pm\mathsf S)/2$ are the separate
geometric self-duality and protected chirality projectors. The unital span
of coordinate words through degree six is the entire $137$-dimensional
native algebra. Degree five does not suffice.
\end{theorem}
\begin{proof}
Write $\{A,B\}=AB+BA$. The antisymmetric word equals
\[
 \{[Y_1,Y_2],[Y_3,Y_4]\}
 -\{[Y_1,Y_3],[Y_2,Y_4]\}
 +\{[Y_1,Y_4],[Y_2,Y_3]\}.
\]
On spinors the Clifford relations give
$24(\sqrt3/4)^4\chi=27\chi/32$. On the geometric block, set
$a=1/2$, $b=\sqrt2/3$ and
$K_{ij}=|e_i\wedge e_j\rangle\langle e_0|
        -|e_0\rangle\langle e_i\wedge e_j|$.
Its commutator has $B$-diagonal block $b^2R_{ij}^{\wedge2}$ and the
$0/B$ off-diagonal block $2abK_{ij}$. In the displayed three-term
contraction, the rotation products give $-4b^4*$ and the scalar-mediated
products give $-4a^2b^2*$. All other geometric blocks vanish. Hence
$-4b^2(a^2+b^2)=-34/81$, proving \eqref{p1:eq:v226-omega}.
Since $Q\mathsf B=(4/9)\mathsf B$ and
$Q\mathsf S=(3/4)\mathsf S$, solving these two scalar equations gives
\eqref{p1:eq:v226-chiral-recovery}.

For the finite-degree claim, the exact modular elimination in the check
script gives $136$ independent words through degree five and $137$ through
degree six over $\mathbb{F}_{97}$, under
$\sqrt2\mapsto14$, $\sqrt3\mapsto10$, $i\mapsto22$.
All defining relations hold, all relevant denominators are invertible,
and a nonzero minor in this reduction is nonzero over the characteristic-zero
coefficient field. The algebra has dimension at most $137$ by its two
reducing blocks. This proves the degree-six spanning claim.

There is also a direct obstruction at degree five. Define the nonzero
Hermitian pseudoscalar
\[
 D=729\mathsf B+544\mathsf S.
\]
Its trace against every coordinate word of degree at most five is zero.
Indeed simultaneous orthogonal covariance of the vector, exterior and
Clifford models makes $D$ odd under each orientation-reversing reflection.
For even tensor rank below four there is no such pseudotensor. At rank
four, coordinate reflections require each index to occur once, and
permutations force proportionality to the alternating tensor. Its coefficient
vanishes because
\[
 \operatorname{Tr}(D\Omega_4)
 =729\cdot6\left(-\frac{34}{81}\right)
  +544\cdot4\left(\frac{27}{32}\right)=0.
\]
Odd words have zero trace against $D$ by the central orthogonal inversion
$Y\mapsto-Y$. Thus their span has dimension at most $136$. The obstruction
is resolved at degree six, since
$\operatorname{Tr}(DQ\Omega_4)=561\ne0$.
\end{proof}

\begin{corollary}[Exact scope of the reconstruction]
\label{p1a:corollary:atlas-exact-scope-of-the-reconstruction-001f1a9}% stable atlas label; provenance v2_277

The already represented rotation, transvection and native central-sector
operators are coordinate-algebra descendants on this branch. The full
$M_{15}(\mathbb{C})$ remains unavailable without an actual nonzero geometric/spinor
bridge. The reconstructed operators do not select a Hamiltonian or its
clock. A polynomial scalar-sector projector is not a chosen physical
vacuum or a proof about the ground state of an independently specified square.
\end{corollary}
\section{A complete native measurement with the least outcome count}
\label{p1:sec:v226-tomography}

\begin{theorem}[Minimal native tomography, with exact inversion]
\label{p1:thm:v226-tomography}
On $\mathcal A_D=M_{11}(\mathbb{C})\oplus M_4(\mathbb{C})$ there is an informationally
complete mathematical POVM with exactly $137$ outcomes, and none with fewer
outcomes can separate its full state space. It can be constructed entirely
inside the coordinate-generated algebra of Theorem~\ref{p1:thm:v226-coordinate-only}.
\end{theorem}
\begin{proof}
Let $d=136$ and choose an HS-orthonormal real basis
$F_1,\ldots,F_d$ of the traceless Hermitian part of $\mathcal A_D$.
Let $v_1,\ldots,v_{d+1}\in\mathbb{R}^d$ be unit regular-simplex vectors, so
\[
 \sum_jv_j=0,\qquad
 \sum_jv_jv_j^T=\frac{d+1}{d}I_d.
\]
Define $B_j=\sum_kv_{jk}F_k$. Since
$\|B_j\|\le\|B_j\|_{\rm HS}=1$, the effects
\begin{equation}\label{p1:eq:v226-effects}
 \boxed{E_j=\frac1{137}(I+\eta B_j),\qquad\eta=\tfrac12,}
\end{equation}
are strictly positive and sum to $I$. All $B_j$ are native Hermitian
coordinate polynomials: Theorem~\ref{p1:thm:v226-coordinate-only} identifies their
entire algebra, and its proof constructs matrix units explicitly.

Every native density has the unique form
$\rho=I/15+\sum_kr_kF_k$. Its probabilities satisfy
\[
 p_j=\operatorname{Tr}(\rho E_j)=\frac1{137}(1+\eta v_j\cdot r).
\]
The simplex frame identity gives exact inversion:
\begin{equation}\label{p1:eq:v226-reconstruction}
 \boxed{\rho=\frac I{15}+\frac{136}{\eta}\sum_{j=1}^{137}p_jB_j.}
\end{equation}
The traceless effect-frame operator is the positive scalar
\begin{equation}\label{p1:eq:v226-frame}
 \mathcal F=\frac{\eta^2}{137\cdot136}I
            =\frac1{74528}I.
\end{equation}
A POVM with $N$ outcomes has at most $N-1$ independent probability
differences. Since the native state space has affine dimension $136$,
$N\ge137$ is necessary, and the construction attains this count.
\end{proof}

\begin{remark}[Measurement scope]
This is complete tomography of the native block-diagonal state algebra,
not of the $88$ off-block coherences of a hypothetical full $M_{15}$ state.
It is not constrained to have the sharp-radius position outcomes or the
same first moments as the endpoint position POVM. Those are different
measurement design constraints. A simplex frame and HS basis are
presentation data for the mathematical construction; the choice does not
become a preferred physical apparatus. Its physical performance still
requires the ordinary G0 occurrence. The linear inverse
\eqref{p1:eq:v226-reconstruction} is not asserted to be a positive or completely
positive classical-to-quantum operation.
\end{remark}

\section{Lie closure of the intrinsic coordinate controls}
\label{p1:sec:v226-control-lie}
Keep the intrinsic dimensionless tuple $Y_a=X_a^D/R$ and
$Q=\sum_aY_a^2$ of Theorem~\ref{p1:thm:v226-coordinate-only}.
Put $Q_0=Q-(31/45)I$, so $\operatorname{Tr}Q_0=0$, and retain
$\Omega_4$ from Theorem~\ref{p1:thm:v226-chirality}.

\begin{theorem}[Exact coordinate-control Lie algebras]
\label{p1:thm:v226-control-lie}
The real skew-Hermitian Lie algebras generated by the following lists
have the stated dimensions and types:
\[
\begin{array}{c|c|c}
 \text{Hermitian generators}&\text{Lie algebra}&\dim_{\mathbb R}\\\hline
 Y_1,\ldots,Y_4&\mathfrak{so}(11)\oplus\mathfrak{usp}(4)&65\\
 Y_1,\ldots,Y_4,Q_0&\mathfrak{su}(11)\oplus\mathfrak{usp}(4)\oplus\mathbb R&131\\
 Y_1,\ldots,Y_4,\Omega_4&\mathfrak{su}(11)\oplus\mathfrak{su}(4)&135\\
 Y_1,\ldots,Y_4,Q_0,\Omega_4&\mathfrak{s}(\mathfrak u(11)\oplus\mathfrak u(4))&136\\
 Y_1,\ldots,Y_4,Q_0,\Omega_4,I&\mathfrak u(11)\oplus\mathfrak u(4)&137
\end{array}
\]
Types are up to the displayed blockwise unitary change of basis. Here
$\mathfrak{usp}(4)$ acts on a complex four-space and has real dimension ten.
The one-dimensional direction in the second row is central.
\end{theorem}
\begin{proof}
On the geometric block the unitary
$U_g=\operatorname{diag}(1,iI_4,I_6)$ makes every $U_g(iY_a)U_g^*$
real skew symmetric. On the protected block choose
\[
 \gamma_1=\sigma_1\otimes I,\quad
 \gamma_2=\sigma_2\otimes I,\quad
 \gamma_3=\sigma_3\otimes\sigma_1,\quad
 \gamma_4=\sigma_3\otimes\sigma_2.
\]
Then $J=\gamma_1\gamma_3$ satisfies $J^T=-J$, $J^*J=I$, and
$\gamma_a^TJ+J\gamma_a=0$. Thus $iY_a$ preserve this symplectic form,
giving the first row's inclusion and upper dimension $55+10=65$.
Adding $iQ_0$ allows the full geometric unitary algebra but only the
protected symplectic algebra plus its scalar line; total trace zero bounds
the dimension by $121+10+1-1=131$.
Adding $i\Omega_4$ instead leaves each block traceless, giving upper bound
$120+15=135$. Adding both has only the total-trace constraint, giving 136;
adding $iI$ gives the full native upper bound 137.

All these upper bounds are attained by exact finite commutator elimination.
For reproducibility, use the coordinate matrices in
Eq.~\ref{p1:eq:v226-matrix}, specialize
$\sqrt2\mapsto14$, $\sqrt3\mapsto10$, $i\mapsto22$ in $\mathbb F_{97}$,
and repeatedly adjoin $i[A,B]$ for each initial Hermitian generator $A$
and each newly retained matrix $B$. Row elimination of their flattened
matrices stabilizes at $65,131,135,136,137$, respectively. The full
algorithm is supplied in \texttt{Paper00-KinematicClosureChecks-v1\_1.py}.
A nonzero minor in this good reduction is nonzero over the characteristic-zero
coefficient field. Combined with the above real dimension bounds, this
proves equality. The central part of $Q_0$ is
$(-4P_g+11P_S)/180$, so the relative scalar direction is nonzero and rational.
The direct-sum types follow from the attained inclusions; in the second
row the geometric simple factor, protected symplectic factor and this
relative scalar direction have the full stated dimensions.
\end{proof}

\begin{remark}[Associative words and Hamiltonian controls]
The complete associative algebra is generated by four coordinates, whereas
their Hamiltonian Lie closure has dimension 65. Products, commutators,
allowed Hamiltonian coefficients, pulse reversals and elapsed controls are
different operation claims. The quadratic and quartic additions above are
explicit existing coordinate words. Their physical availability as independently
controlled Hamiltonians is an additional property of the actual process.
No chirality, state, Hamiltonian coefficient or control duration is selected.
\end{remark}

\begin{corollary}[All native inner automorphisms need only the quartic extension]
\label{p1:cor:v226-native-inner-controls}
The algebra's connected inner-automorphism group is
$PU(11)\times PU(4)$, of dimension 135. Thus the coordinate Hamiltonians
with $\Omega_4$ already generate every connected unitary conjugation on the
native observable algebra, under the actual control hypotheses.
The centered square additionally supplies a relative block phase at the
Hilbert-space unitary level; that central phase acts trivially on
$M_{11}\oplus M_4$ by conjugation.
\end{corollary}
\begin{proof}
Inner conjugation by $U(11)\times U(4)$ has kernel its two scalar circle
factors. The $SU(11)\times SU(4)$ row of
Theorem~\ref{p1:thm:v226-control-lie} maps onto this quotient.
The additional central direction $(-4P_g+11P_S)/180$ commutes with the
whole native algebra. A separately licensed off-diagonal field or event
extension can make relative block phase meaningful; that is a different
operator domain.
\end{proof}


\section{Common-frame invariance of a quaternionic source overlap}
\label{p1:sec:v229-source-overlap}
\begin{corollary}[Parent and source must be transported together]
\label{p1:cor:v229-common-frame-overlap}
For unit quaternions $u,h$, the common left/right spin-frame change
$c(q)\mapsto g_Lc(q)g_R^{-1}$ preserves
$\operatorname{Tr}(c(u)^\dagger c(h))/2=\langle u,h\rangle$.
In the retained axial/opposite-trimer coordinates, $u=(1,1,1,1)/2$
and $h=-e_1$ therefore have invariant overlap $-1/2$.
\end{corollary}
\begin{proof}
Unitarity cancels $g_L$ in the product, and cyclicity cancels $g_R$
in its trace. The quaternionic Clifford Gram is the real Euclidean
inner product. Substitution gives the value. A frame that sends $h$
to the coordinate scalar also transports $u$; it does not change this
overlap or identify two previously different directions.
\end{proof}




\section{The five-core Petersen incidence and oriented-pair closure}
\label{p1:sec:v338-five-core-incidence}
The results in this section concern one supplied, mathematically embedded
$2I$ packet, with its five $2T$ cores.  Core labels are not
quark flavors, color labels, or a selection of physical states.

\begin{theorem}[Ten markings resolve the five-core contrast]
\label{p1:thm:v338-petersen-tight-frame}
Let $W=\mathbb R^5$, $P=J_5/5$, $Q=I_5-P$ and, for each
$S\subset\{1,\ldots,5\}$ of cardinality two, put
$v_S=\mathbf1_S-(2/5)\mathbf1$.  If $V$ is the $5\times10$ matrix
with columns $v_S$ and $A_{\mathrm P}$ is the Petersen adjacency
matrix (disjointness of two-subsets), then
\[
 VV^{\mathsf T}=3Q,\qquad
 V^{\mathsf T}V=I_{10}-A_{\mathrm P}+\tfrac15J_{10}=3P_{4}.
\]
Here $P_4$ is the rank-four Petersen contrast projector.  The
original incidence label space $W\otimes W$ admits, for each split
$S\sqcup S^c$ of sizes $2+3$, the orthogonal four-block grading
\[
 (\mathbb R^S\otimes PW)\oplus(\mathbb R^{S^c}\otimes PW)
 \oplus(\mathbb R^S\otimes QW)\oplus(\mathbb R^{S^c}\otimes QW),
 \qquad 25=2+3+8+12.
\]
The grading concerns incidence \emph{labels}, not physical particles.
\end{theorem}
\begin{proof}
For $S,T$, $\langle v_S,v_T\rangle=|S\cap T|-4/5$.
Its values are $6/5,1/5,-4/5$ for equal, one-overlap, and
disjoint pairs, respectively, giving $V^{\mathsf T}V$.
Each coordinate of $VV^{\mathsf T}$ is $3$ times the orthogonal
projection onto $\mathbf1^\perp$ by permutation symmetry and trace.
The four-block grading follows by the two mutually orthogonal
splits in the tensor factors, with dimensions $(2,3)(1,4)$.
\end{proof}

\begin{theorem}[Petersen double-coset partition of the binary icosahedron]
\label{p1:thm:v338-petersen-120-strata}
For $\pi:2I\to A_5$, fix a two-core marking $S$.  The three
relations $|\pi(g)S\cap S|=2,0,1$ divide the group elements into
exactly $12,36,72$ elements, respectively:
\[120=12+36+72=12(1+3+6).\]
For either single core $H_a$ with $a\in S$, the common refinement
with the existing $24+48+48$ partition of
Proposition~\ref{p1:prop:petersen-120-strata} is
\[
 \begin{array}{c|rrr|r}
 &\mathrm{same}&\mathrm{disjoint}&\mathrm{overlap}&\mathrm{sum}\\\hline
 H_a&6&0&18&24\\
 (\cup_iH_i)\setminus H_a&6&12&30&48\\
 (2I)\setminus\cup_iH_i&0&24&24&48\\\hline
 \mathrm{sum}&12&36&72&120
 \end{array}.
\]
\end{theorem}
\begin{proof}
The $A_5$ action is transitive on the ten two-subsets.  Each target
pair has exactly $120/10=12$ binary preimages, while the Petersen
relation supplies $1,3,6$ target-pair counts.  The refinement follows
by enumerating the identity, double-transposition, three-cycle and
five-cycle classes under the two fixed markings, or by direct
even-permutation enumeration with two binary lifts.  The finite
integer identity can be checked by the indicated finite enumeration.
\end{proof}

\begin{theorem}[Signed Petersen adjacency is the four-dimensional Hodge star]
\label{p1:thm:v338-signed-petersen-hodge}
Write $W=\mathbb Ru\oplus V_4$, $u=\mathbf1/\sqrt5$.
The ten oriented two-core labels form $U=\Lambda^2W$, and
\[
 U=(u\wedge V_4)\oplus\Lambda^2_+V_4\oplus\Lambda^2_-V_4,
 \qquad 10=4+3+3.
\]
In the oriented pair basis $e_i\wedge e_j$ ($i<j$), set
\[
 S_{ij,kl}=\begin{cases}
  \varepsilon_{mijkl},&\{i,j\}\cap\{k,l\}=\varnothing,
    \ \{m\}=\{1,\ldots,5\}\setminus\{i,j,k,l\},\\
  0,&\text{otherwise}.
 \end{cases}
\]
Then $|S|=A_{\mathrm P}$ entrywise, $S^3=5S$, and
\[
 \operatorname{Spec}(S)=\{(+\sqrt5)^{(3)},(-\sqrt5)^{(3)},0^{(4)}\}.
\]
On this oriented-pair carrier (distinct from the unsigned pair-label
carrier in the preceding theorem), the three mutually orthogonal
spectral projections are
\[
 P_{4}=I-S^2/5,\quad
 P_{\pm}=\tfrac12(S^2/5\pm S/\sqrt5).
\]
Every unit basis pair $x=e_i\wedge e_j$ has fixed weights
\[
 \|P_4x\|^2=\tfrac25,\qquad
 \|P_+x\|^2=\|P_-x\|^2=\tfrac3{10}.
\]
These are source-channel projections, not rest masses.
\end{theorem}
\begin{proof}
The contraction $\iota_u(e_i\wedge e_j)=(e_j-e_i)/\sqrt5$
has squared norm $2/5$.  Its complement is a simple bivector in
$\Lambda^2V_4$, whose self-dual and anti-self-dual projections have
equal squared norms, each $3/10$.  The Levi-Civita formula for $S$
is $\sqrt5$ times the Hodge star on $\Lambda^2V_4$ and vanishes
on $u\wedge V_4$.  Thus $S^3=5S$ and the displayed spectral
projections follow.  All formulas are invariant under the even
core permutations, up to the global orientation convention.
\end{proof}
\begin{remark}[Two distinct occurrences of the number twenty-five]
The $25=2+3+8+12$ label grading does not identify that label
space with the separate 25-dimensional matrix-coefficient subspace
of the $120$-element carrier in Paper~V-B.  These two representations
have different diagonal $A_5$ character decompositions.  In
particular, the equality of dimensions is not a physical identification.
\end{remark}


\section{The Petersen edge cochain as a traceless five-core contrast operator}
\label{p1:sec:v339-petersen-endomorphism-bridge}
All assertions in this section are finite mathematical results for a
\emph{supplied} embedded $2I\subset SU(2)$.  They assert neither its
physical occurrence nor a particle, gauge, or energy identification.
Write $G=2I/\{\pm1\}\cong A_5$, let $G$ permute the five tetrahedral
cores, and set $\Omega=\{1,2,3,4,5\}$,
$V=\mathbf1^\perp\subset\mathbb R^5$, $P=J_5/5$, $Q=I_5-P$.

\begin{proposition}[Exact incidence support and its fifteen-dimensional contrast]
\label{p1:prop:v339-incidence-rank15-contrast}
The twenty-five registers
$\mathcal R_{ij}=\{g\in2I:\pi(g)(j)=i\}$ each have $24$
elements. They are distinct from the ten subgroups
$\mathcal A_{\{i,j\}}=H_i\cap H_j\cong C_6$.
For the coefficient map $B:\mathbb R^{5\times5}\to\mathbb R^{120}$,
\[
 (BM)(g)=\sum_{j=1}^5M_{\pi(g)(j),j},\qquad
 B^*B=120(P\otimes P)+30(Q\otimes Q),
\]
one has $\operatorname{rank}B=17$, $\dim\ker B=8$ and the
\emph{diagonal-conjugation} splitting
\[
 \operatorname{im}B\cong\mathbf1\oplus\operatorname{End}(V)
 =\mathbf1\oplus\mathbb RI_V\oplus\mathfrak{sl}(V),
 \qquad 17=1+1+15.
\]
The traceless double-centering operator
\[
 T(M):=QMQ-\tfrac14\operatorname{tr}(QMQ)\,Q
\]
is an orthogonal rank-fifteen projection on the register label
space. Its image lies in $Q\otimes Q$, is perpendicular to $\ker B$,
and is therefore visible in the $120$-element group-function carrier.
\end{proposition}
\begin{proof}
The five-point action of $A_5$ is two-transitive. Hence two
coefficient columns labelled $(i,j)$ and $(k,l)$ overlap in
$24$ elements if both labels agree, in zero elements if exactly
one agrees, and in $6$ elements if both differ. Consequently
$B^*B=24I_{25}+6(J_5-I_5)\otimes(J_5-I_5)$, which equals the
displayed expression. Its positive eigenspaces are $P\otimes P$
and $Q\otimes Q$. The second is naturally $\operatorname{End}(V)$.
Removing its scalar endomorphism line gives exactly $T$.
\end{proof}

\begin{theorem}[Explicit oriented-edge isometry with $\mathfrak{sl}_4(\mathbb R)$]
\label{p1:thm:v339-petersen-sl4-isometry}
Orient arbitrarily the fifteen edges $e=(A\to B)$ of the
Petersen graph $KG(5,2)$, where $A,B\subset\Omega$ are disjoint
two-subsets. Let $W_P=C^1(KG(5,2);\mathbb R)$ be the signed
edge-cochain module and put, for $A=\{i,j\}$,
\[
 d_A=e_i-e_j,\quad v_A=e_i+e_j-\tfrac25\mathbf1,\quad
 S_e=d_Ad_A^{\mathsf T}-d_Bd_B^{\mathsf T},\quad
 K_e=v_Av_B^{\mathsf T}-v_Bv_A^{\mathsf T}.
\]
The maps $S:e\mapsto S_e$ and $K:e\mapsto K_e$ take values in
$\operatorname{Sym}^2_0(V)$ and $\mathfrak{so}(V)$ respectively.
For the mutually orthogonal rank-$(4,5,6)$ projectors
$P_{E,4},P_{E,5},P_{E,6}$ of $W_P$, defined by the graph Hodge
splitting into $\operatorname{im}\partial^*$ and $\ker\partial$,
\[
 S^*S=25P_{E,4}+4P_{E,5},\quad
 K^*K=4P_{E,6},\quad S^*K=0.
\]
Thus
\[
 \Phi:=S+K:W_P\xrightarrow{\sim}\mathfrak{sl}(V),\qquad
 U:=\Phi\bigl(\tfrac15P_{E,4}+\tfrac12P_{E,5}
                       +\tfrac12P_{E,6}\bigr)
\]
is an isometric $S_5$-equivariant isomorphism for the Euclidean
edge norm and Hilbert--Schmidt matrix norm. Its fifteen edge
images form an exact orthonormal basis of $\mathfrak{sl}(V)$.
Here $S_5$ refers to the abstract five-label symmetry; the
inner group action supplied by $2I$ is only its $A_5$ subgroup.
\end{theorem}
\begin{proof}
Let $D$ take an unoriented graph vertex $A$ to $d_Ad_A^\mathsf T$
and choose the signed incidence boundary
$\partial(A\to B)=\delta_A-\delta_B$, so $S=D\partial$.
For the Petersen adjacency $A_{\mathrm P}$,
\[
 D^*D=3I_{10}+J_{10}-A_{\mathrm P},\quad
 \partial\partial^*=3I_{10}-A_{\mathrm P},\quad
 \operatorname{Spec}(A_{\mathrm P})=\{3^{(1)},1^{(5)},(-2)^{(4)}\}.
\]
Their commuting spectral factors give the squared singular values
$0^{(6)},4^{(5)},25^{(4)}$ of $S$. At a graph vertex $A$, the
three disjoint neighboring $B$ satisfy $\sum_{B\cap A=\varnothing}v_B=-2v_A$,
so $K\partial^*=0$. The six-dimensional $\Lambda^2V$ is an
irreducible real $S_5$ representation; hence $KK^*$ is scalar
on it. Each $\|K_e\|_{\mathrm{HS}}^2=8/5$, so
$\operatorname{tr}(KK^*)=15(8/5)=24$ and $KK^*=4I_6$.
Symmetric and skew matrices are orthogonal. Therefore
$\Phi^*\Phi=25P_{E,4}+4(P_{E,5}+P_{E,6})$, which is positive
and yields $U^*U=I_{15}$. Reversing edge orientation flips
both $S_e$ and $K_e$, proving full signed equivariance.
\end{proof}

\begin{corollary}[Graph Hodge sectors form an exact matrix symmetric pair]
\label{p1:cor:v339-petersen-matrix-symmetric-pair}
The isometry identifies the $9+6$ edge Hodge decomposition with
\[
 U(\operatorname{im}\partial^*)=\operatorname{Sym}^2_0(V)
 \cong\mathbf4\oplus\mathbf5,\qquad
 U(\ker\partial)=\mathfrak{so}(V)\cong\mathbf3\oplus\mathbf3'.
\]
For the \emph{mathematical} matrix commutator,
\[
 [6,6]\subseteq6,\qquad[6,9]\subseteq9,\qquad[9,9]\subseteq6,
\]
which is the symmetric pair
$\mathfrak{sl}_4(\mathbb R)=\mathfrak{so}(4)
\oplus\operatorname{Sym}^2_0(\mathbb R^4)$.
The two rank-three sectors are the Hodge eigenspaces of
$\Lambda^2V$ after choosing its orientation; an odd five-label
permutation interchanges them.
\end{corollary}

\begin{theorem}[Rank-fifteen incidence projector and exact source-stratum trace]
\label{p1:thm:v339-petersen-projector-strata}
Let $\chi_4(\sigma)=\#\operatorname{Fix}_{\Omega}(\sigma)-1$
be the standard four-dimensional $A_5$ character. For the
orthonormal edge matrices $U_e=U(e)$ put
\[
 F_e(g)=\frac1{\sqrt{30}}\sum_{j=1}^5(U_e)_{\pi(g)(j),j},
 \qquad g\in2I.
\]
The $F_e$ are orthonormal in $\ell^2(2I)$; the rank-fifteen
orthogonal projector onto their span has kernel
\[
 \boxed{\Pi_{15}(g,h)=\frac{4\chi_4(\pi(g)^{-1}\pi(h))
                         -\chi_4(\pi(g))\chi_4(\pi(h))}{120}.}
\]
Its diagonal leverage is $0$ on the two central lifts, $2/15$
on the thirty double-transposition lifts, and $1/8$ on both
the forty three-cycle lifts and forty-eight five-cycle lifts.
The corresponding total diagonal traces are $0+4+5+6=15$.
\end{theorem}
\begin{proof}
The first proposition gives $B^*B=30I$ on the trace-free
$Q\otimes Q$ space, hence the $F_e$ are orthonormal. Standard
finite-group matrix-coefficient orthogonality gives the kernel
$4\chi_4(\pi(g)^{-1}\pi(h))/120$ of the full rank-sixteen
four-irrep coefficient space. Its unique diagonal-conjugation
invariant character line is spanned by $\chi_4(\pi(g))$,
whose normalized rank-one projector has kernel
$\chi_4(\pi(g))\chi_4(\pi(h))/120$. Subtracting this line
produces the displayed rank-fifteen projection. On the diagonal,
$\Pi_{15}(g,g)=(16-\chi_4(\pi(g))^2)/120$.
The lifted class sizes are $(2,30,40,48)$ and their
$\chi_4$ values $(4,0,1,-1)$, yielding the exact table.
\end{proof}
\begin{remark}[Three distinct fifteens and the physical occurrence firewall]
\label{p1:rem:v339-three-fifteens-firewall}
The real fifteen-dimensional $W_P\cong\mathfrak{sl}(V)$ is not
the complex fifteen-dimensional $\Spin(4)$ momentum carrier
$H_1$ and is not the fifteen-dimensional Bose--Mesner block of
the 24-cell. The rank-fifteen $\Pi_{15}$ is a \emph{finite-group
coefficient projector}, not a selected state, action Hessian,
rest-energy operator, Yukawa coupling, or experimental probability.
It is invariant under diagonal conjugation but not the full
independent left--right group action. The distinct rank-25
$\mathbf5\otimes\mathbf5$ coefficient increment remains in
Paper~V-B; the number $25$ does not identify it with the
original register label space. No universal GP3 momentum
multiplicity-one certificate follows from this construction.
\end{remark}


\section{Full rank-two incidence of the marked $600$-cell and the five-core gluing}
\label{p1:sec:v340-ranktwo-incidence}
Throughout this section the binary-icosahedral marking $2I\subset SU(2)$
is a \emph{conditional mathematical input}, not an assertion of physical
occurrence. Set $G=2I/\{\pm1\}\cong A_5$; its sixty elements are the
antipodal vertex pairs of the $600$-cell. All incidence matrices below
use projective (sixty-point) counting measure. In particular a statement
about the original $120$ vertices requires doubling the point counts.

\begin{theorem}[The 722-polygon pairwise-balanced design]
\label{p1:thm:v340-cyclic-coset-design}
For $k\in\{2,3,5\}$ let $\mathscr L_k$ be the set of all left cosets
$g\langle h\rangle\subset G$ with $\operatorname{ord}(h)=k$, and let
$B_k:\mathbb R^{\mathscr L_k}\to\mathbb R^G$ be the unsigned incidence
map. The great-circle polygons in the binary lift are, respectively,
squares, hexagons and decagons. Their projective incidence parameters are
\[
\begin{array}{c|ccc}
 k&2&3&5\\\hline
 \#\text{cyclic subgroups}&15&10&6\\
 \#\mathscr L_k&450&200&72\\
 \#\text{points per polygon}&2&3&5\\
 \#\text{polygons through a point}&15&10&6
\end{array}.
\]
Every pair of distinct projective vertices belongs to exactly one of
these $722=450+200+72$ polygons. If $M_k=B_kB_k^{\mathsf T}$,
\[
\boxed{M_2+M_3+M_5=30I_{60}+J_{60}},\qquad
450\binom22+200\binom32+72\binom52=\binom{60}{2}.
\]
\end{theorem}
\begin{proof}
The nonidentity elements of $A_5$ have orders $2,3,5$ and class counts
$15,20,24$. Their cyclic subgroup counts are $15,20/2=10,24/4=6$;
each order-$k$ subgroup has $60/k$ left cosets. For $g\ne h$ the
unique cyclic coset containing them is $g\langle g^{-1}h\rangle$:
any cyclic coset containing both must contain $g^{-1}h$ in the same
subgroup, whose order is $2,3$ or $5$. Consequently each off-diagonal
entry of $\sum M_k$ is one. The row valencies $15+10+6=31$ give
diagonal entries $31$, proving the matrix identity.
\end{proof}

\begin{theorem}[Exact order-character incidence projectors]
\label{p1:thm:v340-incidence-character-projectors}
Write $\chi_d$ for the irreducible real $A_5$ character of dimension
$d=4,5$. The characters on the respective classes of orders
$1,2,3,5$ are
\[
 (\chi_4(1),\chi_4(2),\chi_4(3),\chi_4(5))=(4,0,1,-1),
 \quad
 (\chi_5(1),\chi_5(2),\chi_5(3),\chi_5(5))=(5,1,-1,0).
\]
The point Gram differences are central character kernels:
\[
 (M_2-M_3)_{g,h}=\chi_5(g^{-1}h),\qquad
 (M_3-M_5)_{g,h}=\chi_4(g^{-1}h).
\]
Hence the exact mutually orthogonal idempotents are
\[
 \boxed{\mathscr P_5=\tfrac1{12}(M_2-M_3),\quad
 \mathscr P_4=\tfrac1{15}(M_3-M_5)},\quad
 \operatorname{rank}(\mathscr P_5,\mathscr P_4)=(25,16),
 \quad \mathscr P_5\mathscr P_4=0.
\]
Their images are the distinct $\mathbf5\boxtimes\mathbf5^*$ and
$\mathbf4\boxtimes\mathbf4^*$ regular-representation blocks. After
lifting from $G$ to the centre-even part of $\mathbb C[2I]$,
$\operatorname{im}\mathscr P_5$ is the $25$-dimensional nearest-shell
kernel of Theorem~\ref{p1:thm:exceptional-nearest-shell-spectra}.
No equality with the separate twenty-five \emph{register labels} is asserted.
\end{theorem}
\begin{proof}
For $g\ne h$, $(M_k)_{g,h}$ is one exactly if $g^{-1}h$ has
order $k$; on the diagonal it is the number of order-$k$ cyclic
subgroups, namely $15,10,6$. Their differences thus have the displayed
character values including at the identity. The central character
idempotent on $\mathbb R[A_5]$ has kernel
$d\chi_d(g^{-1}h)/60$. The coefficients $1/12=5/60$ and
$1/15=4/60$ are therefore forced. Orthogonality and the ranks
follow from the regular character decomposition.
\end{proof}

\begin{proposition}[Transversal 24-cells and the centred transporter]
\label{p1:prop:v340-transporter-hexagon-intersection}
For $1\le i,j\le5$, let
$\overline{\mathcal C}_{ij}=\{\sigma\in A_5:\sigma(j)=i\}$.
These twenty-five 12-point cells lift to the twenty-five inscribed
24-cells of $2I$. Distinct pairs $(i,j)\ne(k,l)$ have intersection
size $3$ in $G$ if $i\ne k,j\ne l$, and zero otherwise. Every one of
the two hundred projective great hexagons is such an intersection,
\emph{exactly once}. The intersection graph of the twenty-five cells
is the complement of the $5\times5$ rook graph, with parameters
$(25,16,9,12)$ and adjacency
\[
 A_{\rm cell}=(J_5-I_5)\otimes(J_5-I_5).
\]
If $B_{\rm tr}:\mathbb R^{25}\to\mathbb R^{60}$ is the cell-incidence
matrix and $Q=I_5-J_5/5$, then
\[
 \boxed{B_{\rm tr}(Q\otimes Q)B_{\rm tr}^{\mathsf T}=M_3-M_5
 =15\mathscr P_4},\qquad
 \operatorname{rank}B_{\rm tr}=17,\quad
 B_{\rm tr}^{\mathsf T}\mathscr P_5=0.
\]
\end{proposition}
\begin{proof}
The pointwise intersection count specifies images of two distinct
labels under an even permutation; three projective permutations
realize each allowed pair. Their common set is a coset of its
order-three stabilizer. A cyclic $C_3$ fixes precisely two of the
five labels and so belongs to exactly two transporter cells; there
are $25\cdot16/2=200$ intersecting pairs. The row of
$B_{\rm tr}$ at $\sigma$ is the vectorisation of its $5\times5$
permutation matrix $P_\sigma$. Double-centering gives
$QP_\sigma Q=P_\sigma-J_5/5$, hence its inner products are
\[
\operatorname{tr}[(P_\sigma-J/5)^{\mathsf T}
(P_\tau-J/5)]=\#\operatorname{Fix}(\sigma^{-1}\tau)-1
=\chi_4(\sigma^{-1}\tau).
\]
Apply the previous theorem. The constant and centred blocks give
rank $1+16=17$; the $\mathbf5$ block is orthogonal to both.
\end{proof}

\begin{theorem}[Five-core Reye gluing and the ten missing mixed-core hexagons]
\label{p1:thm:v340-reye-70-80}
Let $H_i\le G$ be the five $A_4$ point stabilizers, regarded as the
antipodal quotients of the binary tetrahedral cores. A single
24-cell has twelve Reye points, correctly classified as
\[
12=1_{\mathrm{identity}}+3_{\mathrm{order}\,2}
   +8_{\mathrm{order}\,3}.
\]
The union of the five core point sets has 36 points:
\[
 U=\bigcup_i H_i=\{1\}\sqcup G_{(2)}\sqcup G_{(3)},
 \qquad |U|=1+15+20=36.
\]
The union of the five \emph{core-contained} Reye hexagon sets
comprises precisely $70=10+60$ distinct planes, of order compositions
$(1,3,3)$ and $(2,3,3)$. Its $36\times70$ incidence matrix
$R_{70}$ has rank $35$, with a one-dimensional left nullspace
spanned by
\[
 z(g)=\begin{cases}-2,&g=1\text{ or }\operatorname{ord}(g)=2,\\
                         1,&\operatorname{ord}(g)=3.
       \end{cases}
 \qquad z^{\mathsf T}R_{70}=0.
\]
The \emph{full induced} set of 600-cell hexagons whose three
projective points lie in $U$ contains $80$, rather than $70$, planes:
\[
 \boxed{80=10_{(1,3,3)}+60_{(2,3,3)}+10_{(2,2,2)}}.
\]
The extra ten planes are the three-edge stars of the Petersen graph;
they are not contained in any one marked core. Their inclusion
raises the incidence rank to $\operatorname{rank}R_{80}=36$.
The complete 70-plane Reye union has $210$ incidences and point degrees
$10$ for the identity, $4$ for each order-two point and $7$ for
each order-three point. The old 25-axis count is \emph{not} the
projective Reye-point count.
\end{theorem}
\begin{proof}
Each $A_4$ contains $1$, three involutions and eight 3-cycles.
The five subgroups share the identity and involution class, whereas
each 3-cycle fixes two labels and belongs to two cores; their union
is exactly $U$. A hexagon is a $C_3$ coset. Enumerating the cosets
by the order types of their three elements gives the five patterns
\[
(1,3,3)^{(10)},\quad(2,3,3)^{(60)},\quad(2,2,2)^{(10)},
\quad(2,5,5)^{(60)},\quad(3,5,5)^{(60)}.
\]
The first two types and only those belong to at least one $H_i$;
all three first types are contained in $U$. On every core-contained
hexagon the displayed $z$-weights sum to zero. The exact integer
incidence calculation in the accompanying certificate establishes
rank $35$ for $R_{70}$ and full rank $36$ for $R_{80}$; equivalently,
the $70$-plane Gram spectrum consists of
\[
0^{(1)},\ (14\pm2\sqrt7)^{(1)},\
\bigl((13\pm\sqrt{97})/2\bigr)^{(4)},\
3^{(4)},\ 4^{(10)},\ 6^{(5)},\ 8^{(6)}.
\]
The ten $(2,2,2)$ cosets are translates of Klein-four cosets.
Under the canonical involution--Petersen-edge identification they
are precisely the three incident edges at each of the ten Petersen
vertices. The degree check is $10+15(4)+20(7)=210=70(3)$.
\end{proof}

\begin{corollary}[The Hopf-axis quotient and dodecahedral second neighbours]
\label{p1:cor:v340-reye-dodeca-second-neighbour}
Discard the identity point and identify each 3-cycle with its inverse
(the same imaginary Hopf axis). The 36 Reye points become
$25=15_{\rm Petersen\ edges}+10_{\rm Petersen\ vertices}$,
while the $60$ mixed $(2,3,3)$ planes become $30$ triples, one for
each nonedge of the Petersen graph. Their unsigned $25\times30$
incidence matrix has rank $24$, with left null mode $-2$ on the
fifteen Petersen edges and $+1$ on the ten vertices.
Furthermore the twenty order-three points form the vertices of the
dodecahedral graph $D$, with the sixty mixed Reye planes identified
with its graph-distance-two edges:
\[
 \boxed{D_{(2)}=D^2-3I_{20}.}
\]
For the signed Petersen overlap adjacency $S$ (with the same axis
orientation throughout) and the signed complementary adjacency $T$,
\[
 \boxed{T=S^2-3I_{10}},\qquad
 S^3=5S,\quad T^2+T=6I,\quad ST=TS=2S,
\]
\[
 \operatorname{Spec}(T)=\{2^{(6)},(-3)^{(4)}\}.
\]
In particular $T(S)=T(-S)$: this unsigned two-step Reye coupling
alone does not distinguish the two rank-three mirror eigenspaces of $S$.
\end{corollary}
\begin{proof}
The two 3-cycle representatives $r,r^{-1}$ have opposite imaginary
axes after choosing their order-six quaternion lifts with real part
$1/2$. Among the sixty mixed planes each inversion pair produces the
same projected triple, giving thirty. At the dodecahedron level,
for such lifts $q_r=\tfrac12+\tfrac{\sqrt3}{2}u_r$, the pair condition
that $rs$ is an involution is equivalent to
$\operatorname{Re}(q_rq_s)=0$, or $u_r\cdot u_s=1/3$.
These are exactly the six distance-two neighbours of each dodecahedral
vertex. The $\mathbb Z_2$ sheet decomposition of $D$ is $A\oplus S$,
where $A$ is the ordinary Petersen adjacency; therefore
$D^2-3I\cong(A^2-3I)\oplus(S^2-3I)$. The signed block is $T$.
The stated polynomial identities follow from $S^3=5S$ and
$\operatorname{Spec}(S)=\{(+\sqrt5)^{(3)},0^{(4)},(-\sqrt5)^{(3)}\}$.
The incidence rank and its sole row relation are verified by exact
integer calculation in the certificate.
\end{proof}

\begin{remark}[The transverse overlap does not imply a fivefold angle]
\label{p1:rem:v340-reye-angle-firewall}
The transverse-frame overlap quoted elsewhere as
$\cos^2\alpha=(9+3\sqrt5)/16$ yields
$\alpha\simeq7.761^\circ$ (relative double angle $\simeq15.522^\circ$),
\emph{not} a $72^\circ$ twist. It is also false that
$((3+\sqrt5)/4)^2=(9+3\sqrt5)/16$. No fivefold angle, physical
chirality splitting, particle multiplicity or mass relation is inferred
from these incidence statements. The proof scope is finite geometry.
\end{remark}


\section{Coset--Hodge Gram intertwiners and the parity-partner Petersen cover}
\label{p1:sec:v341-coset-hodge-desargues}
Throughout this section a marked $2I\subset SU(2)$ is supplied as a
\emph{conditional finite datum}.  It is not derived from a single
Hurwitz cell, and none of the following operators is a physical
Hamiltonian, QCD coupling, mass or measurement.  Let
$G=2I$, $\bar G=G/\{\pm1\}\cong A_5$, and let $[5]$ be its
natural five-letter set.  The transporter cells
$\mathcal C_{ij}=\{g\in G:\bar g(j)=i\}$ are the binary lifts
of Proposition~\ref{p1:prop:v340-transporter-hexagon-intersection}.
We use \emph{binary} counting measure on all 120 elements here, not
the 60-point projective convention used for the rank-two polygon
design of Section~\ref{p1:sec:v340-ranktwo-incidence}.

\begin{theorem}[Symmetric and antisymmetric coset Gram identities]
\label{p1:thm:v341-coset-gram-hodge}
Let $B:\mathbb R^{25}\to\ell^2(G)$ be the indicator matrix
$B_{g,ij}={\bf1}[\bar g(j)=i]$.  Write $R=(J_5-I_5)\otimes(J_5-I_5)$.
In the ten unordered label-pair positions $a=\{i,j\}$, $i<j$,
define the two matrices with columns
\[
 v_{ij}={\bf1}_{\mathcal C_{ij}}+{\bf1}_{\mathcal C_{ji}},
 \qquad
 w_{ij}={\bf1}_{\mathcal C_{ij}}-{\bf1}_{\mathcal C_{ji}},
 \qquad V=(v_{ij}),\quad W=(w_{ij}).
\]
Let $A$ be the disjoint-pair Petersen adjacency,
$J=J_{10}$ and $S_H=*_5(c\wedge\cdot)$ on
$\Lambda^2\mathbb R^5$, with $c=(1,1,1,1,1)$ and the increasing
oriented-pair basis.  Then the \emph{integer} identities are
\[
 \boxed{B^{\mathsf T}B=24I_{25}+6R},\quad
 \boxed{V^{\mathsf T}V=48I_{10}+12A+12J},
 \quad\boxed{W^{\mathsf T}W=12S_H^2},\quad
 \boxed{V^{\mathsf T}W=0}.
\]
Their Gram spectra and ranks are
\[
 \begin{array}{c|c|c}
 &\operatorname{Spec}(X^{\mathsf T}X)&\operatorname{rank}X\\\hline
 B&120^{(1)},30^{(16)},0^{(8)}&17\\
 V&204^{(1)},60^{(5)},24^{(4)}&10\\
 W&60^{(6)},0^{(4)}&6
 \end{array}.
\]
Under diagonal $A_5$ conjugation,
\[
 \operatorname{im}V\cong\mathbf1\oplus\mathbf4\oplus\mathbf5,
 \qquad \operatorname{im}W\cong\mathbf3\oplus\mathbf3',
 \qquad \operatorname{im}V\perp\operatorname{im}W.
\]
Thus the ten Petersen vertex channels and the two Hodge triplets
are explicit images of the \emph{same} 25 transporter indicators,
not a dimension-only analogy.
\end{theorem}
\begin{proof}
An even permutation realizing two prescribed images has three
projective choices, hence six binary lifts.  Consequently two
transporter cells meet in 24 points if identical, in no points
if exactly one input/output index is common, and in six points
if both indices differ.  This gives the first matrix identity.
Expanding the symmetric and skew pair columns in that Gram matrix
proves the second and fourth equations.  The same expansion for
skew columns yields $12(5I-\iota_c^{\mathsf T}\iota_c)$,
which equals $12S_H^2$ by
Theorem~\ref{p1:thm:kneser-five-six-hodge-duality}.
The eigenvalues of $R$ are $16,1,-4$ with multiplicities
$1,16,8$, respectively.  On the Petersen vertex eigenspaces
$\mathbf1,\mathbf5,\mathbf4$, $A$ has eigenvalues $3,1,-2$.
Since $S_H^3=5S_H$, its square has eigenvalues $5^{(6)},0^{(4)}$.
These give all ranks and spectra.  On the ten symmetric columns,
label permutation is the standard action on two-subsets; on the six
skew columns, the nonzero Hodge subspace is
$\Lambda^2(c^\perp)\cong\mathbf3\oplus\mathbf3'$.
The last Gram identity also follows from inversion $g\mapsto g^{-1}$,
which fixes each $v_{ij}$ and negates each $w_{ij}$.
\end{proof}

\begin{corollary}[Constructive rank-fifteen coefficient projector]
\label{p1:cor:v341-coset-pi15}
Put $Q=I_5-J_5/5$ and
$\chi_4(\bar g)=\#\operatorname{Fix}_{[5]}(\bar g)-1$.
Under binary counting measure, the projector of
Theorem~\ref{p1:thm:v339-petersen-projector-strata} admits the
explicit 25-cell factorization
\[
 \boxed{\Pi_{15}
 =\frac1{30}B(Q\otimes Q)B^{\mathsf T}
  -\frac1{120}\chi_4\chi_4^{\mathsf T}}.
\]
Its image decomposes under diagonal conjugation as
$\mathbf4\oplus\mathbf5\oplus\mathbf3\oplus\mathbf3'$.
More concretely it is
\[
 \left(\operatorname{im}V\cap\{1,\chi_4\}^{\perp}\right)
   \oplus\operatorname{im}W,
 \qquad\dim=9+6=15.
\]
For the two three-dimensional Hodge sectors let
\[
 P_\pm=\tfrac12\left(S_H^2/5\pm S_H/\sqrt5\right).
\]
Each $P_\pm$ has rank three, and
$(1/\sqrt{60})WP_\pm$ is an isometric embedding on
$\operatorname{im}P_\pm$ into $\ell^2(G)$.
\end{corollary}
\begin{proof}
The centered transporter rows are the matrix coefficients of the
standard four-dimensional representation; by character orthogonality
$B(Q\otimes Q)B^{\mathsf T}$ has entries
$\chi_4(\bar g^{-1}\bar h)$ and equals $30$ times the rank-sixteen
orthogonal coefficient projector.  The character vector $\chi_4$
is a unit-norm line after division by $\sqrt{120}$, so subtracting
its projector leaves the traceless rank-fifteen coefficient space.
The full transporter image is the orthogonal direct sum of the
constant line and this sixteen-dimensional coefficient block.
The preceding theorem supplies a rank-ten plus rank-six orthogonal
image; its invariant part is the single line
$\mathbb R(4-\chi_4)$ inside $\operatorname{im}V$, while its
nontrivial part is the asserted rank-nine plus rank-six space.
Finally $W^{\mathsf T}W=60(P_++P_-)$ proves the two isometries.
This is a concrete realization of the signed-edge module already
identified with $\mathfrak{sl}_4(\mathbb R)$ in
Theorem~\ref{p1:thm:v339-petersen-sl4-isometry}; the carrier is
still distinct from $H_1$ and the 24-cell Bose--Mesner block.
\end{proof}

\begin{theorem}[Desargues parity cover from the two six-axis planes]
\label{p1:thm:v341-desargues-rp2-parity}
Retain the maps $f_\pm:\binom{[5]}2\to\mathcal F_\pm$ of
Theorem~\ref{p1:thm:kneser-five-six-hodge-duality}.  In the
\emph{cross}-incidence table of the two ten-face triangulations of
$\mathbb{RP}^2$, one has exactly
\[
 \boxed{|f_+(a)\cap f_-(b)|=1\ \Longleftrightarrow\ a\cap b=\varnothing.}
\]
Thus the graph joining $f_+(a)$ to $f_-(b)$ precisely for
one common fivefold-axis label has bipartite adjacency
$\bigl(\begin{smallmatrix}0&A\\A&0\end{smallmatrix}\bigr)$:
it is the Desargues graph $GP(10,3)$, the canonical bipartite
double cover of the Petersen graph.

For any signed Petersen adjacency $S$ with $|S|=A$, set
\[
 L(S)=\frac12\begin{pmatrix}A+S&A-S\\A-S&A+S\end{pmatrix},
 \qquad L(S)\simeq A\oplus S.
\]
The \emph{different} signing $S_H$ gives the dodecahedral graph
$GP(10,2)$ with six positive and six negative base pentagons,
whereas $-A$ gives Desargues with all twelve base pentagons negative.
The Desargues spectrum is
\[
 \{3^{(1)},2^{(4)},1^{(5)},(-1)^{(5)},(-2)^{(4)},(-3)^{(1)}\};
\]
the dodecahedral spectrum is
$\{3^{(1)},1^{(5)},(-2)^{(4)},(\sqrt5)^{(3)},0^{(4)},(-\sqrt5)^{(3)}\}$.

Under the degree-six action of $S_5$ obtained by conjugation on the
six Sylow-five subgroups of $A_5$, an even permutation preserves
$\mathcal F_\pm$, while an odd permutation exchanges them:
\[
 \pi f_+(a)=\begin{cases}f_+(\pi a),&\operatorname{sgn}\pi=+1,\\
                             f_-(\pi a),&\operatorname{sgn}\pi=-1.
             \end{cases}
\]
On oriented pairs $T_\pi=\Lambda^2P_\pi$ additionally satisfies
$T_\pi S_HT_\pi^{\mathsf T}=\operatorname{sgn}(\pi)S_H$.
These are precise \emph{parity transformations between signings},
not ordinary automorphisms of a fixed signed dodecahedral connection.
\end{theorem}
\begin{proof}
The stabilizer of a two-subset of $[5]$ fixes a unique face in
each orbit $\mathcal F_\pm$.  After fixing $a=\{1,2\}$, the three
possible pair relationships $b=a$, $|a\cap b|=1$, and
$a\cap b=\varnothing$ yield cross-intersection cardinalities
$0,2,1$, respectively.  The $A_5$ action is transitive on the
corresponding ordered-pair orbits, establishing the displayed
equivalence.  The two-block lift and its spectral decomposition
follow by passing to sheet-even and sheet-odd vectors.  Enumeration
of the twelve Petersen pentagons yields the two voltage counts.
On Sylow-five subgroup labels an odd permutation interchanges
the two face orbits; it suffices to check a transposition and use
$S_5=A_5\sqcup(12)A_5$.  The Hodge parity equation follows from
naturality of the oriented exterior product under $P_\pi$.
A fully enumerated certificate is
\texttt{Proofs/Paper00-PetersenBridges-v1\_3.py}.
For historical double-cover and exceptional six-letter contexts see
D.~A.~Waller, \emph{Double covers of graphs} (1976),
\url{https://doi.org/10.1017/S0004972700025053}, and
Howard--Millson--Snowden--Vakil, \emph{A description of the outer
automorphism of $S_6$} (2008),
\url{https://doi.org/10.1016/j.jcta.2008.01.004}.
The cross-incidence equality is an explicit marked-subgroup
calculation, not a claim of an intrinsic $S_6$ action on $2I$.
\end{proof}

\begin{proposition}[Walk-kernel ten-space is not the Petersen vertex module]
\label{p1:prop:v341-walk-kernel-no-petersen-isomorphism}
For the left-convolution convention of
Theorem~\ref{p1:thm:exceptional-nearest-shell-spectra}, the
$25$-dimensional kernel of the central $2I$ nearest-shell walk is
the regular coefficient block $V_5\otimes V_5^*$.
On restricting one factor to the selected $2T$ core, the
$15$-dimensional joint kernel is the $\mathbf3$ summand times a
fivefold multiplicity.  Its ten-dimensional orthogonal complement
inside $\ker T_{2I}$ carries
\[
 \boxed{V_5\oplus V_5}
\]
under the \emph{commuting right-regular} $A_5$ action, not under a
left-regular action in this convolution convention.  This is
inequivalent to the Petersen vertex module
$\mathbf1\oplus\mathbf4\oplus\mathbf5$.
Moreover the 17-dimensional transporter indicator image lies in the
trivial and four-dimensional regular coefficient blocks and is
orthogonal to the five-dimensional block $\ker T_{2I}$.
Thus $25-15=10$ is not a spectral identification; the diagonal
conjugation action does not preserve the fixed-core joint kernel
without transporting that core as well.
\end{proposition}
\begin{proof}
The restriction $V_5|_{2T}=\chi_\omega\oplus
\chi_{\bar\omega}\oplus V_3$ and the lifted shell eigenvalues
$-\tfrac12,-\tfrac12,0$ appear in the proof of the joint-kernel
corollary following Theorem~\ref{p1:thm:exceptional-nearest-shell-spectra}.
The left-convolution operator acts on one matrix-coefficient factor;
its full commuting action is the right regular one on the other.
This leaves two copies of $V_5$ on the rank-ten complement.
Orthogonality of inequivalent regular isotypic blocks is
Peter--Weyl orthogonality, and the transporter span is exactly the
$\mathbf1\boxtimes\mathbf1\oplus\mathbf4\boxtimes\mathbf4$
coefficient subspace by
Proposition~\ref{p1:prop:v340-transporter-hexagon-intersection}.
\end{proof}

\begin{remark}[Reye gluing, 80 induced hexagons, and scope]
\label{p1:rem:v341-reye-compatibility}
The \emph{36 projective points and 70 core-contained Reye planes}
of Theorem~\ref{p1:thm:v340-reye-70-80} remain unchanged;
the full induced set has \emph{80} planes because ten mixed-core
$(2,2,2)$ hexagons are additional.  The Desargues cross-incidence
is a \emph{different} 20-face $\mathbb{RP}^2$ construction and
must not be substituted for either of these $\mathbb{RP}^3$ counts.
None of these finite carrier equivalences yields a $\mathbb Z_3$
color-center map, a QCD pole, or an observable without separately
licensed occurrence, action and readout data.
\end{remark}


\section{Three-cycle reconstruction of Petersen adjacency and the fifteen-operator class frame}
\label{p1:sec:v342-three-cycle-operator-reconstruction}
All statements here are conditional on the previously specified marked
binary-icosahedral packet $2I\subset SU(2)$ and its quotient
$G=2I/\{\pm1\}\cong A_5$.  They are exact mathematical consequences
of that packet, not new postulates of physical occurrence.
Let $\Omega=\{1,\ldots,5\}$,
$V=\mathbf1^\perp\subset\mathbb R^5$ and
$Q=I_5-\mathbf1\mathbf1^{\mathsf T}/5$.
Write $P_g$ for the permutation matrix of $g\in A_5$ and
$\chi_4(g)=\#\operatorname{Fix}_{\Omega}(g)-1$.  Define the
trace-free, centre-even group-coefficient matrix
\[
 Z_g=QP_gQ-\tfrac14\chi_4(g)Q\in\mathfrak{sl}(V).
\]
We use $\langle X,Y\rangle=\operatorname{tr}(X^{\mathsf T}Y)$.

\begin{theorem}[Petersen adjacency from inverse-paired three-cycle overlaps]
\label{p1:thm:v342-three-cycle-petersen-overlap}
For each two-subset $A\in\binom\Omega2$, let $T=\Omega\setminus A$
and choose one of the two mutually inverse three-cycles $g_A$
supported on $T$.  Set
\[
 S_A=\tfrac12(Z_{g_A}+Z_{g_A^{-1}}),\qquad
 K_A=\tfrac12(Z_{g_A}-Z_{g_A^{-1}}).
\]
The first matrix is independent of the orientation choice and the
second changes sign under reversal.  With
$P_T=\sum_{i\in T}e_ie_i^{\mathsf T}
 -\tfrac13\mathbf1_T\mathbf1_T^{\mathsf T}$,
\[
 S_A=\tfrac34Q-\tfrac32P_T,\qquad
 S_A^2=\tfrac9{16}Q,\qquad
 \sum_AS_A=0.
\]
On the ten fixed-pair labels the Gram matrix satisfies
\[
 \langle S_A,S_B\rangle=
 \begin{cases}
 9/4&A=B,\\
 -5/4&A\cap B=\varnothing,\\
 1/4&|A\cap B|=1,
 \end{cases}
 \quad
 G_S=2I_{10}-\tfrac32A_{\rm P}+\tfrac14J_{10},
\]
where $A_{\rm P}$ is the Petersen adjacency operator.  Hence
\[
 \operatorname{Spec}(G_S)=\{5^{(4)},(1/2)^{(5)},0^{(1)}\},
 \quad \operatorname{span}\{S_A\}=\operatorname{Sym}^2_0(V),
\]
and, without separately inserting Petersen adjacency,
\[
 A\sim_{KG(5,2)}B\quad\Longleftrightarrow\quad
 \langle S_A,S_B\rangle=-\tfrac54.
\]
Furthermore
$\|[S_A,S_B]\|_{\rm HS}^2=5/2$ on Petersen-adjacent
pairs, $1$ on distinct nonadjacent pairs, and $0$ on the diagonal.
\end{theorem}
\begin{proof}
On the moved three-label subspace, the mean of the two three-cycle
permutation matrices is $I-\tfrac32P_T$, and
$\chi_4(g_A)=1$.  Subtracting the quarter-trace gives the displayed
$S_A$.  Since $Q$ and $P_T$ commute, $P_T^2=P_T$ and
$\operatorname{rank}(P_T)=2$, the involution and trace identities
follow.  Evaluating $\operatorname{tr}(P_TP_U)$ for equal, disjoint-fixed-
pair, and one-label-overlap cases yields the Gram formula.  The
Petersen spectrum $\{3^{(1)},1^{(5)},(-2)^{(4)}\}$ gives the
stated Gram spectrum.  Its rank nine equals
$\dim\operatorname{Sym}^2_0(V)$.  Direct multiplication of the
rational $5\times5$ matrices gives the two commutator norms;
the exact verification is included in
\texttt{Proofs/Paper00-ThreeCycleReconstruction-v1\_2.py}.
\end{proof}

\begin{theorem}[Triangle--Petersen signed partial isometry]
\label{p1:thm:v342-triangle-petersen-cycle-isometry}
Each $K_A$ is one half the oriented triangle-flow matrix on the
three labels in $T=\Omega\setminus A$.  The ten skew operators form
an exact tight frame for $\mathfrak{so}(V)$:
\[
 \operatorname{Spec}(G_K)=\{(5/2)^{(6)},0^{(4)}\},\qquad
 X=\tfrac25\sum_A\langle X,K_A\rangle K_A
 \quad(X\in\mathfrak{so}(V)).
\]
The four relations are the simplicial three-boundaries in the
four-simplex on $\Omega$.
For an arbitrarily oriented Petersen edge $e=(B\to C)$ with
$B\cap C=\varnothing$, take the previous Paper~I-B skew edge
matrix $L_e=v_Bv_C^{\mathsf T}-v_Cv_B^{\mathsf T}$, where
$v_B=\mathbf1_B-\tfrac25\mathbf1$, and define
\[
 J_{A,e}=\langle K_A,L_e\rangle.
\]
Then $J\in\{-1,0,1\}^{10\times15}$, with
\[
 J_{A,e}\ne0\ \Longleftrightarrow
 |A\cap B|=|A\cap C|=1,
\]
exactly six nonzeros per row and four per column.  If
$P_{\triangle}$ projects onto the six-dimensional triangle-cycle
subspace of $\mathbb R^{10}$ and $P_{\rm P,cyc}$ onto the six-
dimensional Petersen edge-cycle space in $\mathbb R^{15}$, then
\[
 JJ^{\mathsf T}=10P_{\triangle},\qquad
 J^{\mathsf T}J=10P_{\rm P,cyc}.
\]
Thus $J/\sqrt{10}$ is a signed partial isometry of rank six,
independent of all orientation choices up to row/column signs.
\end{theorem}
\begin{proof}
The antisymmetric part of a three-cycle is the signed three-edge
circulation on its moved triple and has norm squared $3/2$.
Its pairwise Gram values are $0$ for disjoint fixed pairs and
$\pm1/2$ for one-label overlap.  The exact Gram spectrum yields
a $5/2$-tight six-dimensional frame.  The earlier edge matrices
satisfy $\sum_e|L_e\rangle\langle L_e|=4I_{\mathfrak{so}(V)}$
by Theorem~\ref{p1:thm:v339-petersen-sl4-isometry}.  Consequently,
writing $K$ and $L$ for their synthesis maps, $J=K^{\mathsf T}L$
and $KK^{\mathsf T}=(5/2)I$, $LL^{\mathsf T}=4I$ imply
$JJ^{\mathsf T}=10P_{\triangle}$ and
$J^{\mathsf T}J=10P_{\rm P,cyc}$.  Evaluation of one representative
of each five-label orbit gives the claimed integer support and
row/column degrees.  The four kernel relations are the rank-four
image of the simplicial boundary $\partial_3:C_3(\Delta^4)\to
C_2(\Delta^4)$.
\end{proof}

\begin{theorem}[Complete three-cycle frame and golden five-cycle polarization]
\label{p1:thm:v342-class-frame-three-cycle-tomography}
For the $60$ elements of $A_5$ let
$f_g=Z_g/\sqrt{15}\in\mathfrak{sl}(V)$ and define the positive
frame operator $\mathcal A_C=\sum_{g\in C}|f_g\rangle\langle f_g|$
for each conjugacy class $C$.  Under diagonal conjugation,
$\mathfrak{sl}(V)=\mathbf4\oplus\mathbf5\oplus\mathbf3_+\oplus
\mathbf3_-$, with orthogonal projectors $P_4,P_5,P_+,P_-$;
put $P_6=P_++P_-$.  The exact class-frame eigenvalues are
\[
\begin{array}{c|cccc|cc}
 C & \mathbf4 & \mathbf5 & \mathbf3_+ & \mathbf3_- &
 \operatorname{rank} & \operatorname{tr}\\ \hline
 1\ (1)&0&0&0&0&0&0\\
 2A\ (15)&1/3&8/15&0&0&9&4\\
 3A\ (20)&2/3&1/15&1/3&1/3&15&5\\
 5A\ (12)&0&1/5&a&b&11&3\\
 5B\ (12)&0&1/5&b&a&11&3
\end{array}
\quad a=\frac{5+2\sqrt5}{15},\quad b=\frac{5-2\sqrt5}{15}.
\]
Here $\sum_C\mathcal A_C=I_{15}$, and the $20$ three-cycle
operators alone reconstruct every $X\in\mathfrak{sl}(V)$:
\[
 \mathcal A_{3A}=\tfrac23P_4+\tfrac1{15}P_5+\tfrac13P_6,
 \qquad \mathcal A_{3A}^{-1}=\tfrac32P_4+15P_5+3P_6,
\]
\[
 \boxed{X=\sum_{g\in3A}\langle X,f_g\rangle
              \mathcal A_{3A}^{-1}f_g.}
\]
For $D_5=\mathcal A_{5A}-\mathcal A_{5B}$,
\[
 D_5P_4=D_5P_5=0,\qquad
 D_5^2=\tfrac{16}{45}P_6,\qquad
 \frac ab=9+4\sqrt5=\varphi^6,
 \quad\varphi=\frac{1+\sqrt5}{2}.
\]
Pullback to all $120$ elements of $2I$ gives the \emph{same}
frame operators when the duplicated lifted rows are normalized
as $Z_{\pi(g)}/\sqrt{30}$.  The two $3$-dimensional sector names
and the ordering of $5A,5B$ depend on the orientation convention.
\end{theorem}
\begin{proof}
The standard four-dimensional real representation has
$\sum_{g\in A_5}\rho(g)_{ij}\rho(g)_{kl}=15\delta_{ik}\delta_{jl}$.
Removing the trace line gives $\sum_g|f_g\rangle\langle f_g|=I_{15}$.
Each class-frame operator commutes with diagonal conjugation;
finite exact character calculations in the displayed four
irreducibles give the table.  In particular its three-cycle row
has strictly positive eigenvalues on all fifteen dimensions,
which gives the inverse and reconstruction identity.
The difference of the two five-cycle rows vanishes on the
symmetric $\mathbf4\oplus\mathbf5$ sectors and has eigenvalues
$\pm4\sqrt5/15$ on the two skew three-dimensional sectors.
Squaring proves the $D_5$ identity.  The class and orientation
normalizations, including all twenty three-cycles, are independently
checked in the companion exact-arithmetic certificate.
\end{proof}

\begin{theorem}[Minimal three-generator path criterion]
\label{p1:thm:v342-minimal-path-lie-generators}
Under ordinary matrix commutation, the family $\{S_A\}$
generates the full real Lie algebra
$\operatorname{Lie}\langle S_A:A\in\binom\Omega2\rangle
=\mathfrak{sl}(4,\mathbb R)$.  No two distinct members generate
more than a four-dimensional Lie algebra.  Among the $120$
three-element subsets of the ten labels, precisely the $60$
whose three pairs form a four-vertex path $P_4$ in the complete
label graph $K_5$ generate all fifteen dimensions; the
$10$ triangles, $20$ three-edge stars and $30$ disjoint
wedge--edge arrangements each generate a nine-dimensional
Lie algebra.  A path triple has cumulative Lie-word dimensions
\[
 3,\quad6,\quad11,\quad14,\quad15
\]
at depths $1,\ldots,5$ respectively.
\end{theorem}
\begin{proof}
The preceding theorems imply that the ten $S_A$ span
$\operatorname{Sym}^2_0(V)$ and their pairwise brackets span
$\mathfrak{so}(V)$, so the full Lie closure is
$\mathfrak{sl}(V)$.  The symmetric operators are $S_5$-equivariant
in the five-label marking; the possible unordered pairs form
two $S_5$-orbits, and the triples form the four graph-shape orbits
listed in the statement.  Exact rational commutator closure on
one representative of each orbit gives dimensions $4$ for either
pair orbit and $(15,9,9,9)$ for the four triple orbits.
The path representative yields depth increments $(3,3,5,3,1)$.
These computations use only the displayed rational projection
formula for $S_A$ and are executed in
\texttt{Proofs/Paper00-ThreeCycleReconstruction-v1\_2.py}.
\end{proof}

\begin{remark}[Observable ownership and the three-fifteens firewall]
\label{p1:rem:v342-class-frame-observable-firewall}
The three-cycle $f_g$ are a linear Hilbert--Schmidt reconstruction
frame, not a positive-operator-valued measurement on an already
constituted physical quantum state space.  Selecting one minimal
three-generator path is an extra marking: the unmarked ten-pair
family is canonical, but a preferred path is not.  The mathematical
matrix commutator does not define a new physical gauge algebra.
The golden ratio $a/b$ is an exact class-frame eigenvalue ratio,
not a measured mass ratio or physical parity violation.
The real $A_5$ operator carrier and its doubled centre-even
$2I$ coefficient lift remain distinct from the complex
$\Spin(4)$ momentum carrier $H_1$, the 24-cell association
module and the unrelated nearest-shell right-regular $\mathbf5
\oplus\mathbf5$ ten-space.  G0 occurrence, GP2 action/readout
ownership and the independent GP3 whole-candidate comparison
remain separate requirements.
\end{remark}


\begin{proposition}[Fivefold class difference and the normalized signed Petersen contrast]
\label{p1:prop:v343-fivefold-signed-contrast}
Let $W_P$ denote the real oriented $15$-edge Petersen representation,
with $W_P\downarrow_{A_5}=\mathbf3\oplus\mathbf3'\oplus\mathbf4\oplus\mathbf5$,
and let $\Phi:W_P\rightarrow\mathfrak{sl}(4,\mathbb R)$ be the
orthogonal intertwiner of Theorem~\ref{p1:thm:v339-petersen-sl4-isometry}.
Choose consistent names $5A,5B$ for the two twelve-element
five-cycle classes and put
\[
 K_W=\sum_{g\in5A}\rho_W(g)-\sum_{g\in5B}\rho_W(g),\qquad
 L_W=\frac{K_W}{4\sqrt5}.
\]
With the sector labels chosen consistently with $D_5=\mathcal A_{5A}-
\mathcal A_{5B}$ in
Theorem~\ref{p1:thm:v342-class-frame-three-cycle-tomography},
\[
 \boxed{L_W=\Pi_{\mathbf3}-\Pi_{\mathbf3'},\qquad
 \Phi L_W\Phi^{-1}=-\frac{15}{4\sqrt5}D_5.}
\]
The class-sum and contrast obey $K_W^3=80K_W$,
$L_W^3=L_W$, $L_W^2=\Pi_{\mathbf3}+\Pi_{\mathbf3'}$,
$\operatorname{rank}L_W^2=6$ and
$\operatorname{Spec}L_W=\{(+1)^{(3)},(-1)^{(3)},0^{(9)}\}$.
For an odd relabeling $t\in S_5\setminus A_5$ one has
$\rho_W(t)L_W\rho_W(t)^{-1}=-L_W$; switching the $5A/5B$
convention reverses $L_W$, not $L_W^2$.
\end{proposition}
\begin{proof}
The two three-dimensional character values on the two five-cycle
classes are $\varphi$ and $1-\varphi$ in opposite order; hence their
class-sum differences are $\pm 12(2\varphi-1)/3=\pm4\sqrt5$.
The differences vanish on $\mathbf4\oplus\mathbf5$.  Exact
class-character sums give $\operatorname{tr}((\Phi K_W\Phi^{-1})D_5)=-32$,
$\operatorname{tr}((\Phi K_W\Phi^{-1})^2)=480$, and
$\operatorname{tr}(D_5^2)=32/15$.  Thus the Hilbert--Schmidt
Cauchy--Schwarz bound is saturated with negative sign, proving
$\Phi K_W\Phi^{-1}=-15D_5$, and hence
$\Phi L_W\Phi^{-1}=-15D_5/(4\sqrt5)$.  The minimal polynomial
and outer-class interchange rule follow from the sector eigenvalues.
\end{proof}
\begin{remark}[The contrast is not a physical chirality operator]
\label{p1:rem:v343-fivefold-class-scope}
The odd $S_5$ action here is the outer class interchange of the
abstract $A_5$ packet.  It need not be a metric isometry of a supplied
binary-icosahedral embedding and is not, by itself, spacetime parity.
The $15$-dimensional real $W_P$ remains distinct from the complex
level-one momentum module $H_1$.
\end{remark}

\paragraph{Independent exact certificate for the eleven Petersen bridge claims.}
The following source-matched verification has no additional geometric
occurrence premise: all inputs are the marked \(A_5=2I/\{\pm1\}\)
permutation packet and the formulas already displayed in
Sections~\ref{p1:sec:v339-petersen-endomorphism-bridge}--
\ref{p1:sec:v342-three-cycle-operator-reconstruction}.
For the fifteen signed Petersen edges, exact rational contraction gives
\[
 S^*S=25P_{E,4}+4P_{E,5},\qquad K^*K=4P_{E,6},\qquad
 P_{E,4}+P_{E,5}+P_{E,6}=I_{15}.
\]
The normalized \(U\) obeys \(U^*U=I_{15}\) and signed equivariance
under all 120 relabelings in \(S_5\).  The coefficient kernel of
\(\Pi_{15}\) is verified for all 60-by-60 projective pairs,
with both binary lifts accounted for exactly.  Hence this is a
complete finite-domain check of the stated isometry, symmetric pair,
and coefficient projector, not a dimension-matching analogy.

Independently, enumerating all \(200\) cyclic \(C_3\) cosets in
\(A_5\) gives types
\[
 (1,3,3)^{10},\ (2,3,3)^{60},\ (2,2,2)^{10},\
 (2,5,5)^{60},\ (3,5,5)^{60}.
\]
The \(36\times70\) core-contained and \(36\times80\) induced Reye
incidences have exact ranks \(35\) and \(36\); the order-three
inverse-pair quotient has rank \(24\).  The sixty mixed
\((2,3,3)\) planes form precisely the pairs \((r,s)\) for which
\(rs\) is an involution, and the resulting 20-vertex graph is
explicitly isomorphic to the distance-two graph of \(GP(10,2)\).
For their positive-real order-six lifts
\[
 q_r=\frac12+\frac{\sqrt3}{2}u_r,
\]
the condition \(rs\) involutive is equivalent to
\(u_r\cdot u_s=1/3\).  The signed-Hodge matrix has
\[
 S_H^3=5S_H,\qquad (S_H^2-3I)^2+(S_H^2-3I)=6I,
\]
exactly.

The two Sylow-five face orbits and the Desargues cross-incidence
are verified for all 100 pairs of five-core labels and all 120
five-label relabelings.  The right-regular
\(V_5\oplus V_5\) complement is distinguished from the Petersen
vertex module by its three-cycle character, and the fifth-character
projector annihilates the 25 transporter columns.  The inverse-paired
three-cycle operators satisfy their ten-by-ten Gram and commutator
identities; their ten-by-fifteen signed incidence \(J\) has
\(JJ^{\mathsf T}=10P_{\triangle}\) and
\(J^{\mathsf T}J=10P_{{\rm P},\mathrm{cyc}}\).
The complete \(A_5\) conjugacy-class frame identities, the fivefold
contrast \(K_W=-15D_5\), and the Lie closures of all 120
three-generator subsets (reduced to four exact \(S_5\)-orbits) agree
with the stated formulas.  The character, orientation, and finite
incidence checks together verify the eleven claims originally grouped
under the Petersen geometric-intertwiner evidence issue.

Four standalone exact certificates, reconstructing the objects from
permutations and rational matrices rather than loading precomputed
matrices, are named
\path{Paper1B-PetersenSL4ExactProof-v1_1.py},
\path{Paper1B-ReyeExactProof-v1_1.py},
\path{Paper1B-DesarguesWalkExactProof-v1_1.py}, and
\path{Paper1B-ThreeCycleExactProof-v1_1.py}; their output digests
and precise proof scopes are in the companion DAG audit layer.
The result does not establish a unique \emph{physical} realization,
a marked \(2I\) occurrence from G0, or any additional GP3
selection.  It does not cover unrelated reproduced or conditional
claims elsewhere in Paper~I-B.



\section{The nontrivial Hurwitz character and the full Petersen quintet}
\label{p1:sec:v344-twisted-hurwitz-petersen}
Throughout this section $G=2I$, $H=2T\subset G$, $Q_8\triangleleft H$,
and the supplied binary-icosahedral embedding and core are \emph{mathematical
comparison data}.  In particular $[G:H]=5$ and $H/Q_8\simeq C_3$.
Fix one nontrivial unitary character
$\chi:H\to\{1,\omega,\omega^2\}$, where $\omega^3=1\ne\omega$.
Because $\chi(-1)=1$, it descends to one of the two nontrivial
one-dimensional characters $\mathbf1',\mathbf1''$ of $A_4=H/\{\pm1\}$.

\begin{theorem}[The twisted five-core representation is irreducible]
\label{p1:thm:v344-induced-hurwitz-quintet}
The induced unitary representation satisfies
\[
 \boxed{\operatorname{Ind}_{2T}^{2I}\chi\simeq\mathbf5},
\]
where $\mathbf5$ is the five-dimensional irreducible representation
of $A_5=2I/\{\pm1\}$.  The choice $\bar\chi$ yields an equivalent
quintet, not a second independent rank-$25$ coefficient sector.
\end{theorem}
\begin{proof}
The restrictions of the $A_5$ irreducibles to $A_4$ are
$\mathbf1\downarrow=\mathbf1$,
$\mathbf3\downarrow=\mathbf3$,
$\mathbf3'\downarrow=\mathbf3$,
$\mathbf4\downarrow=\mathbf1\oplus\mathbf3$, and
$\mathbf5\downarrow=\mathbf1'\oplus\mathbf1''\oplus\mathbf3$.
Frobenius reciprocity therefore places $\mathbf1'$ with multiplicity
one in the restricted quintet and in no other center-even irreducible.
The induced representation has dimension $[G:H]=5$, so it is exactly
$\mathbf5$.  Center-odd irreducibles are excluded because $\chi$ is
trivial on $-1$.  The same argument applies to $\mathbf1''$.
\end{proof}

\begin{theorem}[Full orthogonal twisted twenty-five-register carrier]
\label{p1:thm:v344-twisted-25-gram}
Choose left-coset representatives $g_1,\ldots,g_5$ for $G/H$.
For $g\in G$ put
\[
 F_{ij}(g)=
 \begin{cases}
  \chi(g_i^{-1}gg_j),&g_i^{-1}gg_j\in H,\\
  0,&\text{otherwise}.
 \end{cases}
 \qquad 1\leq i,j\leq5.
\]
The matrix $[F_{ij}(g)]$ is the monomial induced quintet, and
\[
 \boxed{\sum_{g\in G}F_{ij}(g)\overline{F_{kl}(g)}
          =24\delta_{ik}\delta_{jl}.}
\]
Writing $F:\mathbb C^{25}\to\mathbb C[G]$ for these coefficient
columns gives
\[
 \boxed{F^\dagger F=24I_{25},\qquad
        \Pi_{25}=\frac1{24}FF^\dagger,\qquad
        (\Pi_{25})_{g,h}=\frac1{24}\chi_{\mathbf5}(g^{-1}h).}
\]
The twisted coefficient space has rank $25$ and is orthogonal to
the untwisted five-coset coefficient space, which has rank $17$
and Gram eigenvalues $120^{(1)},30^{(16)},0^{(8)}$.
For the already-normalized unsigned incidence transfer
$T_u=P_1+\tfrac14P_4+\tfrac15P_5$, each unit twisted coefficient
$\xi_{ij}=F_{ij}/\sqrt{24}$ satisfies
$\|T_u\xi_{ij}\|^2=1/25$.
\end{theorem}
\begin{proof}
Theorem~\ref{p1:thm:v344-induced-hurwitz-quintet} gives an
irreducible unitary quintet.  Matrix-coefficient Schur orthogonality
on a group of order $120$ gives $120/5=24$ and the central
character-projector formula.  The trivial inducing character yields
$\mathbf1\oplus\mathbf4$, so its coefficient space has dimension
$1^2+4^2=17$ and is orthogonal to the $\mathbf5$ coefficient block.
The eigenvalues of its Gram follow by Schur orthogonality on the
$\mathbf1$ and $\mathbf4$ blocks.  The last assertion uses
$T_u\Pi_{25}=\Pi_{25}/5$; it is a dimensionless norm identity.
\end{proof}

\begin{proposition}[Explicit isometry to the Petersen pair quintet]
\label{p1:prop:v344-petersen-quintet-intertwiner}
On the ten unordered pairs $S\in\binom{\{1,\dots,5\}}2$, let
$D_{i,S}=\mathbf1_{i\in S}$.  The Petersen quintet projector is
\[
 \boxed{P_{\rm Pet}=I_{10}-D^\top(3I_5+J_5)^{-1}D,\qquad
 (P_{\rm Pet})_{S,T}=\delta_{ST}-\frac{|S\cap T|}{3}+\frac16.}
\]
Write $\rho_{10}$ for the pair permutation representation and
$Q_\chi=\tfrac1{12}\sum_{h\in A_4}\overline{\chi(h)}\rho_{10}(h)$.
For any pair basis vector $e_S$ with $Q_\chi e_S\ne0$, let
$v=Q_\chi e_S/\|Q_\chi e_S\|$.  Then
\[
 J e_i=\rho_{10}(g_i)v,\qquad
 \boxed{J^\dagger J=I_5,\quad JJ^\dagger=P_{\rm Pet},\quad
 \rho_{10}(g)J=J U_\chi(g)}
\]
(with $g_i$ now viewed modulo the center) supplies an explicit
unitary intertwiner of the twisted Hurwitz quintet with the original
Petersen five-dimensional submodule.
\end{proposition}
\begin{proof}
The incidence identity $DD^\top=3I_5+J_5$ proves the first
projector formula and its displayed entries.  The character
projector $Q_\chi$ has rank one in the Petersen quintet and
vanishes on $\mathbf1\oplus\mathbf4$.  Since $v$ transforms
under $A_4$ by $\chi$, the coset relation
$gg_j=g_i h$ proves equivariance.  Schur's lemma makes
$J^\dagger J$ scalar; each column is unit, so the scalar is one.
The image is the full five-dimensional $\operatorname{im}P_{\rm Pet}$.
\end{proof}

\begin{proposition}[Oriented Hurwitz edges are eight regular orbits]
\label{p1:prop:v344-hurwitz-edge-eight-orbits}
The $24$ Hurwitz vertices of the $24$-cell form the left regular
$H$-set.  Each vertex has eight nearest-neighbour edges, giving
$96$ unoriented and $192$ oriented edges.  Left multiplication by
$H$ acts freely on oriented edges, and therefore partitions them
into exactly eight orbits of size $24$.
\end{proposition}
\begin{proof}
The Hurwitz unit shell is the $D_4$ root polytope with degree eight,
so double-counting incidence gives $24\cdot8/2=96$ edges.
A group element fixing an oriented edge fixes its initial vertex;
left regularity forces that element to be the identity.
\end{proof}

\begin{remark}[Scope of the mathematical character bridge]
\label{p1:rem:v344-twisted-physical-scope}
This bridge repairs the rank-$17$ \emph{untwisted} obstruction by a
different, character-twisted representation; it does not turn the
$600$-cell comparison into five realized physical shells.  The
coefficient identity $1/25$ is not a hadron mass eigenvalue, a
flavor label, or an action-selected coupling.  Physical realization,
source normalization and mass-pole extraction belong downstream.
\end{remark}


\clearpage
\section{Quaternionic F40 edges as regular five-core transitions}
\label{p1:sec:v345-f40-native-arrows}
Let $G=2I$, let $H_1,\ldots,H_5$ denote its five conjugate
binary tetrahedral subgroups, and let $\mathcal F_{40}$ denote the
forty-vertex cubic cover of the ten projective threefold axes.
Write $\mathcal Z=\ker\partial_{40}\subset\mathbb R^{60}$ for its
real oriented-edge cycle space.  Every edge is oriented from the
quaternion endpoint of real part $+\tfrac12$ to the endpoint of
real part $-\tfrac12$.  Its two quotient elements are
$g_R(e)=q^{-1}q'$ and $g_L(e)=q'q^{-1}$.

\begin{theorem}[A regular sixty-edge fivefold transport packet]
\label{p1:thm:v345-f40-edge-regular}
The two quotients belong to the twelve-element shell
$\mathcal S=\{g\in G:\operatorname{Re}g=(\sqrt5-1)/4\}$.
Each directed edge covers exactly one Petersen edge and hence omits
one label $m_e\in\{1,\ldots,5\}$.  The map
\[
 e\longmapsto(g_R(e),m_e):E(\mathcal F_{40})\longrightarrow
 \mathcal S\times\{1,\ldots,5\}
\]
is a bijection, equivariant under conjugation of the shell and
permutation of the five labels by $A_5=G/\{\pm1\}$.  The action on
sixty edges is regular.  As real $A_5$-representations,
\[
 \mathbb R^{E(\mathcal F_{40})}\cong\mathbb R[A_5],\qquad
 \boxed{\mathcal Z\cong\mathbf3\oplus\mathbf3'\oplus3\mathbf5,
 \quad\dim\mathcal Z=21.}
\]
Each $g\in\mathcal S$ induces a five-cycle $p_g$ on the five
subgroups by $gH_mg^{-1}=H_{p_g(m)}$; the core-transition labels
$(m_e,p_{g_R(e)}(m_e))$ attain every ordered pair of distinct cores
exactly three times.  Thus
\[
 \boxed{\sum_{g\in\mathcal S}\rho_5(g)=3(J_5-I_5),\qquad
 \operatorname{Spec}\Bigl(\tfrac1{12}\sum_g\rho_5(g)\Bigr)
 =\{1^{(1)},(-1/4)^{(4)}\}.}
\]
\end{theorem}
\begin{proof}
The finite quaternionic multiplication and five-core incidence
certificate in Proposition~\ref{p1c:prop:v345-f40-certificate}
checks the bijection in exact $\mathbb Q(\sqrt5)$ arithmetic.
Independently, a shell element projects to a five-cycle, with
centralizer $C_5$ in $A_5$.  That centralizer acts freely and
transitively on the five core labels, so the $A_5$ action on
$(g,m)$ has trivial stabilizer.  The two vertex classes each have
twenty elements with stabilizer $C_3$.  Subtraction of the vertex
permutation characters from the regular sixty-edge character, with
the connected-graph constant restored, gives the displayed
$\mathbf3\oplus\mathbf3'\oplus3\mathbf5$ cycle character.
The core walk has zero diagonal and three entries at each ordered
off-diagonal position; its spectrum follows from $J_5$.
\end{proof}

\begin{theorem}[Three intrinsic readouts close the F40 cycle space]
\label{p1:thm:v345-f40-triple-frame}
On oriented edge vectors define integer-valued readout maps
\[
 T_R(e)=\delta_{g_R(e)}\in\mathbb R^{12},\qquad
 T_L(e)=\delta_{g_L(e)}\in\mathbb R^{12},\qquad
 A(e)=\delta_{(m_e,p_{g_R(e)}(m_e))}\in\mathbb R^{20}.
\]
On $\mathcal Z$, every single map has rank $11$, every pair has rank
$16$, and $\mathcal R=(T_R,T_L,A)^{\mathsf T}|_{\mathcal Z}$
has rank $21$.  In the Euclidean graph-chain metrics,
\[
 \boxed{\det(zI_{21}-\mathcal R^*\mathcal R)
 =(z-1)^5(z^2-12z+31)^3(z^2-8z+6)^5.}
\]
Consequently the lower frame bound is $4-\sqrt{10}>0$.
The maps commute with the native $A_5$ action; no odd abstract
five-core relabeling is needed for geometric completeness.
\end{theorem}
\begin{proof}
Restrict the three integer matrices to $\ker\partial_{40}$.
Exact rational rank and characteristic-polynomial calculations,
recorded in Proposition~\ref{p1c:prop:v345-f40-certificate},
prove the assertions.  In particular the polynomial has positive
roots $1$, $6\pm\sqrt5$ and $4\pm\sqrt{10}$.
\end{proof}

\begin{proposition}[Balanced arrow and left/right mirror separation]
\label{p1:prop:v345-f40-balanced-mirror}
For $z\in\mathcal Z$, write $Az$ as an off-diagonal $5\times5$
matrix $Y$.  It has zero row and column sums.  The resulting
rank-$11$ image splits orthogonally as
\[
 \mathcal B_{20}=\mathcal B_{\rm sym}^{(5)}\oplus
 \mathcal B_{\rm skew}^{(6)},\qquad
 \mathcal B_{\rm skew}\cong\mathbf3\oplus\mathbf3'.
\]
If $P_Z$ is the cycle projector and $P_{\rm sym},P_{\rm skew}$
are the corresponding output projectors, then
\[
 \boxed{AP_ZA^{\mathsf T}=2P_{\rm sym}+P_{\rm skew}.}
\]
Let $A_{\rm ico}$ denote the twelve-vertex icosahedral adjacency,
$P_{\rm ant}$ its antipodal permutation, and $J=J_{12}$.  The two
shell readouts obey
\[
 \begin{split}
 T_RP_ZT_L^{\mathsf T}&=-\tfrac94I+\tfrac12A_{\rm ico}
 +\tfrac14P_{\rm ant}-\tfrac1{24}J,\\
 (T_R+T_L)P_Z(T_R+T_L)^{\mathsf T}
 &=I+P_{\rm ant}-\tfrac16J,\\
 (T_R-T_L)P_Z(T_R-T_L)^{\mathsf T}
 &=10I-2A_{\rm ico}.
 \end{split}
\]
The symmetric left/right readout has rank five and annihilates
both mirror modes; the antisymmetric readout has rank eleven and
is nonzero on both mirror modes.
\end{proposition}
\begin{proof}
Cycle conservation at each vertex gives both row/column
balance conditions.  The twenty off-diagonal entries have nine
independent balance constraints; the symmetric space has dimension
five and the skew space dimension six.  Exact evaluation of the
three readout Grams yields the displayed identities.  The
icosahedral adjacency eigenvalues establish the two ranks.
\end{proof}

\begin{remark}[Geometric completeness versus physical action]
\label{p1:rem:v345-f40-source-firewall}
These are identities of the conditional occurring $2I$ packet.
The rank-$21$ frame is a comparison map, not a measured POVM,
QCD Hamiltonian, particle multiplicity or mass source.  The full
$F40$ graph symmetry, native $A_5$ symmetry, and gauge symmetry
of the later color-singlet action are different requirements.
\end{remark}

\clearpage
\section{F40--Petersen completion, mirror transport and harmonic coefficient calculus}
\label{p1:sec:v349-f40-harmonic-completion}
The results in this section are exact comparison mathematics on the fixed
binary-icosahedral packet; none constitutes a new occurrence postulate,
physical QCD source, or extension of the level-one momentum carrier.
Write $G=2I/\{\pm1\}=A_5$, let $V_5$ be the center-even twisted
Hurwitz quintet, and let $\mathcal Z_{40}=\ker\partial_{40}$ be the
real F40 cycle module from Theorem~\ref{p1:thm:v345-f40-edge-regular}.

\begin{theorem}[Sharp twisted-25/F40 multiplicity obstruction and minimal repair]
\label{p1:thm:v349-f40-minimal-quintet}
For the \emph{diagonal conjugation} action of $A_5$ on
$\mathcal T_{25}=\operatorname{End}(V_5)$ one has
\[
 \mathcal T_{25}\cong\mathbf1\oplus\mathbf3\oplus\mathbf3'
 \oplus2\mathbf4\oplus2\mathbf5,\qquad
 \mathcal Z_{40}\cong\mathbf3\oplus\mathbf3'\oplus3\mathbf5.
\]
Every equivariant linear map $\mathcal Z_{40}\to\mathcal T_{25}$
has rank at most $16$, and the bound is attained.  Among additional
real $A_5$ modules $W$ admitting an equivariant injection
$\mathcal Z_{40}\hookrightarrow\mathcal T_{25}\oplus W$,
the minimum is $\dim W=5$, necessarily of type $\mathbf5$.
The Petersen adjacency $A_P$ on the ten unordered core pairs supplies
that module intrinsically:
\[
 P_{P,5}=-\frac{(A_P-3I)(A_P+2I)}6,\qquad
 \mathcal Z_{40}\hookrightarrow\mathcal T_{25}\oplus\operatorname{im}P_{P,5}.
\]
The representation type is uniquely minimal, but the intertwiner is
not unique merely from the multiplicity data.
\end{theorem}
\begin{proof}
The quintet character on the $(1,2A,3A,5A,5B)$ classes is
$(5,1,-1,0,0)$, so the diagonal endomorphism character is
$(25,1,1,0,0)$.  Character orthogonality with the five irreducible
characters gives the first decomposition.  The F40 cycle character
$(21,1,-3,1,1)$ gives the second.  Complete reducibility and
Schur's lemma bound the common multiplicities by one $\mathbf3$,
one $\mathbf3'$ and two $\mathbf5$ copies, total dimension
$3+3+10=16$.  Adding exactly one $\mathbf5$ is both necessary
and sufficient.  The Petersen eigenvalues $3,1,-2$, of multiplicities
$1,5,4$, give the displayed spectral projector, whose image is
precisely the required existing quintet.
\end{proof}

\begin{proposition}[Lossless native twenty-one-mode decomposition]
\label{p1:prop:v349-f40-reduced-frame}
Let $T_R,T_L:\mathcal Z_{40}\to\mathbb R^{12}$ and
$A:\mathcal Z_{40}\to\mathbb R^{20}$ denote the intrinsic
right, left and directed-core readouts.  If $P_{S,5}=
\frac12(I+\mathrm{Ant})-J_{12}/12$ and
$(J_{\rm pair}Y)_{\{i,j\}}=(Y_{ij}+Y_{ji})/2$, then
\[
 \mathcal Mz=\left(P_{S,5}(T_R+T_L)z,\;
 P_{S,5}(T_R-T_L)z,\;J_{\rm pair}Az,\;P_{\rm skew}Az\right)
\]
has rank $5+5+5+6=21$ on $\mathcal Z_{40}$.  Its third output
belongs to $\operatorname{im}P_{P,5}$, and its fourth realizes
$\mathbf3\oplus\mathbf3'$.  The cycle module has no nonzero
$A_5$-invariant linear scalar functional, while the real space of
invariant symmetric quadratic forms has dimension $8$.
\end{proposition}
\begin{proof}
Double balance of $Az$ makes the symmetric unordered-pair
projection a row-sum-zero ten-pair function.  The Petersen
adjacency of such a function equals the function itself: the sum
on pairs disjoint from $\{i,j\}$ is its value on $\{i,j\}$ by
the vanishing total and core-row sums.  The five-dimensional
symmetric arrows and the six-dimensional skew arrows, together
with the two independent shell quintets, give an injective
readout by the exact F40 frame certificate.  The decomposition
$\mathbf3\oplus\mathbf3'\oplus3\mathbf5$ contains no trivial
summand; symmetric bilinear invariants have dimensions
$1+1+\dim\operatorname{Sym}^2\mathbb R^3=8$.
\end{proof}

\begin{theorem}[Exact mirror gluing and golden conditioning]
\label{p1:thm:v349-f40-mirror-gluing}
For a fivefold shell element $g$, let $p_g$ be its permutation of
five tetrahedral cores, and set $p_g^{d_g(i,j)}(i)=j$ with
$d_g(i,j)\in\{1,2,3,4\}$ for $i\ne j$.  The rational matrix
\[
 K_{(i,j),g}=\frac{5-2d_g(i,j)}{40}
\]
satisfies, on every $z\in\mathcal Z_{40}$,
\[
 \boxed{P_{\rm skew}Az=K(T_R-T_L)z},\qquad
 K^{\mathsf T}K=\frac{5I+A_{\rm ico}}{40}P_{\rm odd}.
\]
On the six-dimensional antipodal-odd and balanced-skew spaces,
$K$ is invertible, with inverse
$K^{-1}=2(5I-A_{\rm ico})K^{\mathsf T}$; its nonzero singular
values are $\sqrt{(5+\sqrt5)/40}$ and
$\sqrt{(5-\sqrt5)/40}$, each threefold.  Hence
$\kappa(K)=(1+\sqrt5)/2$.
\end{theorem}
\begin{proof}
The five-cycle transport has precisely the four cyclic
distances; the entry formula is its skew, balanced transport
kernel.  Contract the native edge-to-shell and edge-to-arrow
matrices with the orthogonal cycle projector from
Theorem~\ref{p1:thm:v345-f40-triple-frame}.  The exact
$20\times12$ rational product is $K(T_R-T_L)$, and its squared
Gram reduces to the displayed icosahedral adjacency expression.
The antipodal-odd adjacency has eigenvalues $\pm\sqrt5$,
each threefold.  Their two positive Gram eigenvalues prove
invertibility and the inverse formula.  The exact integer/rational
certificate is recorded in Paper~I-D and in the distributed
\texttt{Proofs/Paper00-F40MirrorGluing-v1\_1-verify.py}.
\end{proof}

\begin{theorem}[Twisted 25 as a complete degree-four scalar harmonic shell]
\label{p1:thm:v349-twisted25-spin2-lift}
The quintet $V_5$ is the restriction to $2I$ of the $SU(2)$
spin-two representation $\rho_2$.  For normalized Haar measure,
\[
 U(X)(q)=\sqrt5\operatorname{tr}(X\rho_2(q)),\qquad
 U:\operatorname{End}(\mathbb C^5)\longrightarrow
 \mathcal H_4(S^3(R))
\]
is an isometry after choosing an intertwining quintet basis, and
$\Delta_+U(X)=24U(X)/R^2$.  Its normalized restriction to the
120 vertices of $2I$ is the twisted Hurwitz 25-coefficient frame,
with $F^\dagger F=24I_{25}$.  Under the diagonal $A_5$ action,
$\mathcal H_4$ has type $\mathbf1\oplus\mathbf3\oplus\mathbf3'
\oplus2\mathbf4\oplus2\mathbf5$.
The pre-existing Petersen quintet also has a degree-two scalar
harmonic realization:
\[
 \mathcal Q_5=\left\{\sqrt{12}\,\mathbf v^{\mathsf T}Q\mathbf v:
 Q=Q^{\mathsf T},\ \operatorname{tr}Q=0\right\}
 \subset\mathcal H_2(S^3(R)),\quad
 \Delta_+|_{\mathcal Q_5}=8/R^2.
\]
Consequently there exists an $A_5$-equivariant isometric injection
$\mathcal Z_{40}\hookrightarrow\mathcal H_4\oplus\mathcal Q_5$,
with comparison-space complement of type $\mathbf1\oplus2\mathbf4$
and dimension nine.
\end{theorem}
\begin{proof}
The spin-two character restricts to $(5,1,-1,0,0)$, the
unique quintet character; Schur orthogonality proves the isometry,
and Peter--Weyl gives $\mathcal H_4$ and its eigenvalue.
Finite Schur orthogonality on $2I$ gives the stated coefficient
Gram after the unitary basis identification.  For a unit
quaternion $q=a+\mathbf v$, the ambient quadratic polynomial
$\mathbf v^{\mathsf T}Q\mathbf v$ is harmonic; the fourth spherical
moment of $\mathbf v$ makes the displayed $\sqrt{12}$
normalization isometric for the Frobenius metric.  Conjugation
acts through the $SO(3)$ spin-two module, which restricts to
$\mathbf5$.  The first theorem then provides the stated
injection and complement.  It is an existence statement, not a
canonical physical spectral map.
\end{proof}

\begin{corollary}[The inherited nearest-shell spectrum fixes the exact 600-cell gap]
\label{p1:cor:v3026-600-hs-spectral-gap}
For $A_{600}$ the ordinary $120$-vertex nearest-neighbour adjacency
$\operatorname{Re}(q^{-1}r)=\varphi/2$ and
$L_{600}=12I-A_{600}$, the existing
Theorem~\ref{p1:thm:exceptional-nearest-shell-spectra} gives
\[
\boxed{A_{600}=12T_{2I},\qquad
\dim\ker A_{600}=25,\qquad
\lambda_2(L_{600})=9-3\sqrt5>0.}
\]
Here $T_{2I}$ is the \emph{normalized} central nearest-shell walk of
that theorem; no new independently chosen graph generator or
normalized operator is introduced. The first nonconstant Laplacian
eigenspace has dimension four, and its quaternion-coordinate vectors
attain the gap. Its finite adjacency polynomial and exact integer
moments can also be replayed with
\path{Proofs/Paper00-HilbertSchmidtChecks-v1_3.py}.
\end{corollary}
\begin{proof}
The twelve nearest-shell elements form an inverse-closed conjugacy
class. Left and right central class convolution agree, so the
undirected adjacency operator is precisely $12T_{2I}$ under the
paper's operator conventions. Multiply every eigenvalue in the
\emph{already proved} normalized $2I$ spectral table by twelve.
Its largest eigenvalue is $12$, its next largest is
$6\varphi=3+3\sqrt5$ with multiplicity four, and its zero
eigenspace has dimension $25$. The four ambient coordinate functions
obey $\sum_{r\sim q}r=6\varphi q$, giving the equality witness.
\end{proof}

\begin{corollary}[The twisted quintet is exactly the adjacency-zero band]
\label{p1:cor:v3026-twisted25-adjacency-kernel}
For the \emph{same} character projector $\Pi_{25}$ of
Theorem~\ref{p1:thm:v344-twisted-25-gram} and the spin-two
identification of Theorem~\ref{p1:thm:v349-twisted25-spin2-lift},
\[
\boxed{\operatorname{Ran}\Pi_{25}=\ker A_{600},\qquad
 A_{600}\Pi_{25}=0,\qquad L_{600}\Pi_{25}=12\Pi_{25},}
\]
\[
\boxed{\Pi_{25}(q,r)={1\over24}U_4\!
   \left(\operatorname{Re}(q^{-1}r)\right),\quad
 U_4(x)=16x^4-12x^2+1.}
\]
This makes explicit the earlier $25$-dimensional central-walk kernel
of Proposition~\ref{p1:prop:v341-walk-kernel-no-petersen-isomorphism}
as the \emph{specified twisted projector}.  It is an adjacency kernel,
not a zero-Laplacian mode or a massless physical state.  In addition,
for the existing tetrahedron--vertex incidence $T$,
$\|T f\|^2=20\|f\|^2$ on $\operatorname{Ran}\Pi_{25}$.
The full $25$-dimensional degeneracy cannot select either of the two
previously rooted $C_3$ triality vectors without their rooted data.
\end{corollary}
\begin{proof}
The twelve nearest neighbours of the identity form one conjugacy
class in $2I$, so their adjacency is central group convolution.
The quintet is the restriction of $\operatorname{Sym}^4(\mathbb C^2)$,
with character $U_4(\operatorname{Re}q)$.  On the nearest shell,
$\operatorname{Re}q=\cos(\pi/5)$, so
$U_4(\cos(\pi/5))=\sin(5\pi/5)/\sin(\pi/5)=0$.
Schur's lemma therefore gives $A_{600}\Pi_{25}=0$.
The projector has rank $25$ by
Theorem~\ref{p1:thm:v344-twisted-25-gram}, and Corollary~\ref{p1:cor:v3026-600-hs-spectral-gap}
gives $\dim\ker A_{600}=25$; hence equality.
The displayed kernel is the character formula in that theorem and
the spin-two character formula.  Finally use
$T^{\mathsf T}T=20I+5A_{600}$ from
Theorem~\ref{p1:thm:v365-600-cell-incidence}.
\end{proof}

\begin{corollary}[Sharp global Hilbert--Schmidt graph rigidity]
\label{p1:cor:v3026-600-hs-dirichlet}
Let $X_q$ be same-sort operators on one \emph{fixed} finite Hilbert
fibre for all $q\in2I$, and put
$\overline X=120^{-1}\sum_q X_q$. Then
\[
\boxed{\sum_{\{q,r\}\in E_{600}}
\|X_q-X_r\|_{\rm HS}^2\ge
(9-3\sqrt5)\sum_{q\in2I}
\|X_q-\overline X\|_{\rm HS}^2.}
\]
The constant is sharp and the equality space (away from the uniform
mode) is the four-dimensional quaternionic first-harmonic eigenspace
tensored with any fixed fibre-operator direction. This does \emph{not}
assert dissipation under physical time evolution or identify a gauge
connection between nontrivially different fibres.
\end{corollary}
\begin{proof}
For each orthonormal Hilbert--Schmidt operator component, the left
side is the scalar graph quadratic form $\langle x,L_{600}x\rangle$.
Apply the spectral gap of the preceding theorem to the centered
component and sum. The four quaternionic coordinate eigenvectors
attain equality.
\end{proof}

\begin{theorem}[Complete cubic/quartic finite-group invariant defect]
\label{p1:thm:v349-twisted25-vertex-defect}
On $V_5^{\otimes k}$, $k=1,2,3,4$, the $SU(2)$ and $2I$
Haar-averaging projectors satisfy
\[
 \begin{array}{c|rrrr}k&1&2&3&4\\\hline
 \operatorname{rank}\Pi_{SU(2),k}&0&1&1&5\\
 \operatorname{rank}\Pi_{2I,k}&0&1&2&11\\
 \operatorname{rank}(\Pi_{2I,k}-\Pi_{SU(2),k})&0&0&1&6
 \end{array}
\]
and the projector difference at $k=3,4$ is precisely the
projection onto the unique $2I$-fixed vector in each spin-six
isotypic multiplicity component (one copy at $k=3$, six at
$k=4$).  For $u_{ab}=\sqrt5D^{(2)}_{ab}$ the continuum cubic
vertex is exactly
\[
 \int u_{ab}u_{cd}u_{ef}\,dg=5^{3/2}
 \begin{pmatrix}2&2&2\\a&c&e\end{pmatrix}
 \begin{pmatrix}2&2&2\\b&d&f\end{pmatrix}.
\]
The two-holomorphic/two-antiholomorphic quartic vertex is the
finite Clebsch--Gordan contraction
\[
 \int u_{ab}u_{cd}\overline{u_{ef}u_{gh}}\,dg=
 25\sum_{J=0}^4\frac1{2J+1}\sum_{M,N=-J}^J
 C^{JM}_{2a,2c}C^{JN}_{2b,2d}
 \overline{C^{JM}_{2e,2g}C^{JN}_{2f,2h}}.
\]
\end{theorem}
\begin{proof}
The $A_5$ class sizes are $(1,15,20,12,12)$, giving
$\dim(V_5^{\otimes k})^{2I}=(5^k+15+20(-1)^k)/60$.
Continuous spin coupling gives invariant ranks $0,1,1,5$.
In tensor spins at most eight, the first nontrivial $2I$ fixed
line appears at spin six and no further such line occurs;
the spin-six multiplicities in the triple and quadruple tensors
are one and six.  Since the continuous fixed space is contained
in the finite fixed space, the projector difference is exactly
this spin-six contribution.  Orthogonality and triple-product
identities for Wigner coefficients yield the displayed vertices;
finite averaging instead adds the established projector defect.
\end{proof}

\begin{remark}[Constitution and physical meaning]
\label{p1:rem:v349-mathematical-physical-boundary}
The exact $25+5$ construction is a \emph{comparison-space}
completion and not an alteration of the GP1--GP3 constitution,
Paper II's actual $m_p=1$ momentum shell, or the physical color
source.  In particular the golden condition number and the
Laplacian eigenvalues are not particle mass factors.  The
spin-two cubic/quartic tensors are available for a common-action
calculation but their physical coefficients and occurrence remain
open.
\end{remark}

\clearpage
\section{Exact fixed-shell Hurwitz cubature and low-band spectral intertwining}
\label{p1:sec:v350-hurwitz-exact-spectral}
The following mathematical results are properties of the fixed, round Hurwitz
comparison shell.  They refine the existing sampling isometry of
Corollary~\ref{p1:cor:v185-2t-degree-two-sampling-isometry}; they do not
change the physical $m_p=1$ carrier or imply a spinorial or physical
QCD field equation.

\begin{theorem}[Hurwitz spherical five-design with explicit moments]
\label{p1:thm:v350-hurwitz-five-design}
In Euclidean quaternionic coordinates the unit Hurwitz shell is
\[
 V_H=\{\pm e_i:1\le i\le4\}\,\cup\,
 \{(s_1,s_2,s_3,s_4)/2:s_i\in\{\pm1\}\}\subset S^3.
\]
For the round probability measure $d\sigma$ on $S^3$ and any polynomial
$p$ on $\mathbb R^4$ of total degree at most five,
\[
 \frac1{24}\sum_{q\in V_H}p(q)=\int_{S^3}p(q)\,d\sigma(q).
\]
In particular,
\[
 \frac1{24}\sum_q q_iq_j=\frac{\delta_{ij}}4,\qquad
 \frac1{24}\sum_q q_iq_jq_kq_l
 =\frac{\delta_{ij}\delta_{kl}+\delta_{ik}\delta_{jl}+
 \delta_{il}\delta_{jk}}{24}.
\]
\end{theorem}
\begin{proof}
The shell contains eight axis points and sixteen half-coordinate points.
Independent coordinate sign changes and permutation symmetry reduce the
even moments through degree four to
$\langle q_i^2\rangle=1/4$,
$\langle q_i^4\rangle=(2+16/16)/24=1/8$, and
$\langle q_i^2q_j^2\rangle=(16/16)/24=1/24$ for $i\ne j$.
These are the corresponding round $S^3$ moments. Central inversion kills
all odd monomials through degree five for both measures. This exhausts
monomials of degree at most five, hence all such polynomials.
\end{proof}

\begin{theorem}[Exact sampled scalar Laplacian on degrees zero, one and two]
\label{p1:thm:v350-hurwitz-scalar-laplacian}
Let $A_H$ be the $24\times24$ adjacency matrix defined by
$A_H(q,q')=1$ if and only if $\langle q,q'\rangle=1/2$. Let
$\mathcal B_2=\mathcal Y_0\oplus\mathcal Y_1\oplus\mathcal Y_2$
be the $14$-dimensional scalar harmonic band on $S^3(R)$,
and let $Sf(q)=f(q)/\sqrt{24}$. Set
\begin{equation}
\label{p1:eq:v350-hurwitz-exact-laplacian}
 L_H=\frac1{R^2}\left(\frac{A_H^2}{16}-\frac32A_H+8I\right).
\end{equation}
Then $S$ is an isometry and
\[
 \boxed{L_H Sf=S(-\Delta_{S^3(R)})f\quad(f\in\mathcal B_2).}
\]
The eigenvalues of $A_H$ on the sampled $\mathcal Y_0,\mathcal Y_1,
\mathcal Y_2$ are exactly $8,4,0$, respectively. This identity holds
on $S\mathcal B_2$ and is not asserted on the entire 24-dimensional
vertex space or on the independent physical momentum carrier $\mathcal H_1$.
\end{theorem}
\begin{proof}
Isometry follows from Theorem~\ref{p1:thm:v350-hurwitz-five-design},
since products of degree-two polynomials have degree at most four.
For each vertex $q$, the eight neighbors can be written
$q'_a=q/2+(\sqrt3/2)u_a$, where
$u_a\in T_qS^3$, $\sum_a u_a=0$, and
$\sum_a u_au_a^T=(8/3)(I-qq^T)$. Consequently,
\[
 \sum_{q'\sim q}q'=4q,\qquad
 \sum_{q'\sim q}q'q'^T=2I_4.
\]
Hence $A_H$ acts by $8$ on constants and by $4$ on linear
coordinate functions. On $q^TBq$ for any symmetric trace-free
$B$, the second identity gives $\sum_{q'\sim q}q'^TBq'=0$,
so its degree-two eigenvalue is zero. The eigenvalues of
$-\Delta_{S^3(R)}$ are $\ell(\ell+2)/R^2$; substitution of
$8,4,0$ into the polynomial in Eq.~\eqref{p1:eq:v350-hurwitz-exact-laplacian}
gives $0,3/R^2,8/R^2$, respectively. Injectivity of $S$ completes
the intertwining proof.
\end{proof}

\begin{remark}[No momentum-truncation flow or Dirac equivalence]
\label{p1:rem:v350-hurwitz-spectral-domain}
The exact identity is an additional finite spectral comparison on the
existing $24$ vertices. The scalar $14$-dimensional band is not the
$15$-dimensional physical $m_p=1$ carrier, and no graph-polynomial
formula for the curved spin Dirac operator follows from it without an
explicit spinorial sampling/intertwining certificate. The radial
separation of adjacent vertices remains $\ell_H=\pi R/3$ and cannot
be made small merely by reciprocal-position refinement.
\end{remark}

\section{Minimum closed Hurwitz face supports and dual-cut certificate}
\label{p1:sec:v356-hurwitz-minimal-faces}

\begin{theorem}[Minimum no-dangling triangular support of the Hurwitz 24-cell]
\label{p1:thm:v356-hurwitz-minimal-faces}
Let $\mathscr F$ be the $96$ triangular faces of the Hurwitz
$24$-cell boundary in Theorem~\ref{p1:thm:v235-hurwitz-chain-complex}.
Call a nonempty subset $S\subseteq\mathscr F$ \emph{no-dangling}
when each edge incident to some face of $S$ is incident to at least
two faces of $S$ (the three-face edge incidence is retained).
Then $|S|\geq 8$.  The no-dangling subsets with $|S|\leq 8$
are exactly the $24$ eight-triangle boundaries of the octahedral
three-cells.  In particular every minimum support is an oriented
integral $2$-cycle and a closed topological $2$-sphere.
\end{theorem}
\begin{proof}
The vertex set is the $24$ Hurwitz short roots; two vertices are
adjacent precisely when their normalized scalar product is $1/2$.
This generates exactly $96$ edges, the $96$ three-cliques are the
triangular faces, and each edge belongs to exactly three faces.
The following exhaustive recurrence is a finite certificate not using
an assumed geometric picture.  Index the faces from $0$ to $95$;
for each seed $f_0$ begin with $S=\{f_0\}$, maintaining the exact
edge-incidence count $n_e(S)$.  If every $n_e(S)$ is $0$ or at least
$2$, record $S$.  Otherwise select the least edge with $n_e(S)=1$.
Every possible completion must add one of the two remaining faces
incident to that edge.  Branch on each such face, memoize the face
bitmask, and discard branches of size exceeding $8$.  This
exhausts every no-dangling set of size at most $8$: any such set
contains a seed, and at every dangling edge its completion must
choose at least one of the two prescribed branches.

With the indicated exact Hurwitz adjacency the $96$ seed searches
visit $23{,}584$ distinct within-seed search states.  Their deduplicated
completed masks are precisely $24$ sets, each of cardinality $8$.
Independently, enumerating the $24$ octahedral cells from six-vertex
cliques with octahedral incidence yields $24$ distinct eight-face
boundaries, each in the completed-mask set.  Thus no other completed
mask occurs.  The orientations and cycle property are furnished by
$\partial_2\partial_3=0$ from the original chain theorem.
\end{proof}

\begin{proposition}[Dual-graph certificate for closed mod-two surfaces]
\label{p1:prop:v356-hurwitz-dual-cuts}
For the mod-two reduction of the boundary complex, every face
$2$-cycle is the boundary of a selection of octahedral cells.
Its face support is the edge cut of the dual $24$-vertex,
$8$-regular Hurwitz adjacency graph.  Every nonzero cut has at
least eight edges; the $24$ minimum cuts isolate a single
cell (or its complement).
\end{proposition}
\begin{proof}
The boundary is $S^3$, so $H_2(S^3;\mathbb F_2)=0$.
The dual vertices can be indexed by $\pm e_i\pm e_j$,
$i<j$, with adjacency when the Euclidean scalar product is $1$.
Fix the last dual vertex outside a cut, and enumerate the remaining
$2^{23}$ bitmasks.  With $c(0)=0$, removal of the least set bit $i$
gives the exact integer recurrence
\[
 c(S)=c(S\setminus\{i\})+8
       -2\#\{j\in S\setminus\{i\}:j\sim i\}.
\]
The resulting minimum nonzero value is eight and occurs $24$ times
among these complement-normalized masks (the $23$ other singleton
cuts and the complement of the fixed vertex).  This independently
certifies the mod-two subclass.  The preceding theorem's
no-dangling classification is stronger, because a selected edge
may meet all three faces without being an even mod-two cycle.
\end{proof}

\begin{remark}[Kinematic and physical scope]
\label{p1:rem:v356-no-face-action-selection}
\leavevmode\par
The classification concerns the native fixed momentum-shell
incidence.  It does not identify the shell with a primitive
reciprocal-position lattice or select any gauge face action, heat
parameter, physical vacuum, or measured Wilson observable.  An
action- and source-preserving presentation bridge remains a separate
physical certificate.
\end{remark}


\clearpage
\section{Hurwitz face--edge complement and marked 600-cell source grammar}
\label{p1:sec:v365-hurwitz-hadron-geometry}
The constructions below are mathematical consequences of the already
marked $2T\subset 2I$ shell, not assertions about an occurring strong
Hamiltonian, physical hadron pole, or decay probability.  In particular,
the 25 chamber \emph{blocks} below are disjoint subsets of the 600
four-dimensional tetrahedral cells, whereas the twenty-five transporter
\emph{registers} $C_{ij}\subset 2I$ overlap as sets of vertices.

\begin{theorem}[Exact vertex--chamber incidence of the marked 600-cell]
\label{p1:thm:v365-600-cell-incidence}
Let $T\in\{0,1\}^{600\times120}$ be the tetrahedron--vertex
incidence of the standard unit 600-cell, and let $A_{\rm vert}$ be its
nearest-vertex adjacency.  Then
\[
  \boxed{T^{\mathsf T}T=20I_{120}+5A_{\rm vert},\qquad
  \operatorname{rank}T=120.}
\]
The rank statement concerns this exact finite incidence, not the
$15$-dimensional intrinsic momentum module $\mathcal H_1$.
\end{theorem}
\begin{proof}
The tetrahedral vertex figure of the 600-cell has twenty facets,
so each vertex belongs to twenty tetrahedra.  Each nearest-neighbour
edge has pentagonal link, hence lies in five tetrahedra; nonadjacent
vertices cannot be vertices of the same regular tetrahedron.  These
incidence counts give the displayed Gram matrix.  The full-column-rank
claim has an independent integer certificate: reduce the explicit
$600\times120$ incidence to $\mathbb F_{101}$ and perform Gaussian
elimination; the pivot count is 120.  A full-rank minor modulo 101
implies rank 120 over $\mathbb Q$ and $\mathbb R$.  The coordinate
and pivot generator is recorded in the accompanying current DAG proof
file \texttt{Proofs/Paper00-HurwitzHadronFiniteCheck-v1\_1.py}.
\end{proof}

\begin{theorem}[Face--edge Laplacian complement and its 23-dimensional kernel]
\label{p1:thm:v365-hurwitz-face-laplacian}
Let $M\in\{0,1\}^{96\times96}$ be the \emph{unsigned} edge--face
incidence matrix of the Hurwitz 24-cell boundary.  Join two distinct
triangular faces whenever they share an edge, and denote the resulting
adjacency by $A_F$ and graph Laplacian by $L_F$.  Then
\[
 \boxed{M^{\mathsf T}M=3I_{96}+A_F,\quad
 L_F=6I_{96}-A_F,\quad M^{\mathsf T}M+L_F=9I_{96}},
\]
\[
 \boxed{\operatorname{rank}M=73,\qquad\dim\ker M=23.}
\]
The kernel is the image of the oriented boundaries of the 24
combinatorial octahedral three-cells, with their sole whole-boundary
linear relation, after the natural triality edge-orientation gauge.
\end{theorem}
\begin{proof}
Every triangular face has three edges and every edge belongs to three
faces (Theorem~\ref{p1:thm:v235-hurwitz-chain-complex}).  Two distinct
faces share at most one edge.  Thus the diagonal of $M^{\mathsf T}M$
is 3 and its off-diagonal is $A_F$; each face has six distinct
face-neighbours, establishing the Laplacian formula.  The 24 vertices
split into the three Hurwitz triality octets (coordinate axes, even
half-sign vertices and odd half-sign vertices).  Every triangle has
one vertex of each kind; orient its boundary cyclically in that order.
This fixes the signs of all three incident oriented edges to be
positive, so $M$ is conjugate by diagonal row/column signs to the
oriented boundary operator $\partial_2$.  The cellular complex is
the boundary $S^3$ of the regular 24-cell.  Since $H_1(S^3)=H_2(S^3)=0$,
its chain ranks are $\operatorname{rank}\partial_1=23$,
$\operatorname{rank}\partial_2=96-23=73$ and
$\operatorname{rank}\partial_3=96-73=23$; the single relation
among 24 octahedral boundaries is $\ker\partial_3\cong\mathbb R$.
This also proves the claimed kernel spanning description.
\end{proof}

\begin{proposition}[Cubic-invariant selection and the purely algebraic color bridge]
\label{p1:prop:v365-cubic-algebra-color}
Let $V=\operatorname{Im}\mathbb H\cong\mathbb R^3$ with the rotational
actions induced by $2T$, $2O$, and $2I$.  In $V^{\otimes3}$ the
respective invariant dimensions are $2,1,1$: the Levi-Civita
alternating tensor $\epsilon$ survives all three actions, and the
additional \emph{symmetric} tetrahedral cubic is removed by the
octahedral or icosahedral rotations.  Moreover
\[
 \mathfrak{so}(V)\oplus i\operatorname{Sym}_0^2(V)
   \;=\;\mathfrak{su}(3)
\]
under commutators in the standard complexified three-space.  The
right-hand side identifies an eight-dimensional Lie \emph{algebra
comparison}; it is not a new derivation of the physically realized
$SU(3)_c$ connection already owned by Paper V-A.
\end{proposition}
\begin{proof}
The alternating invariant line is $\Lambda^3V$.  An invariant
rank-three tensor symmetric in any pair splits into its symmetric
cubic and mixed symmetry pieces.  The rotational three-space's mixed
symmetry has no singlet, whereas the cubic harmonic has a unique
$A_4$-invariant line (represented in tetrahedral axes by $xyz$).
A quarter-turn from the octahedral group reverses that cubic;
the icosahedral three-space has no invariant harmonic of degree
three by its character average.  The matrix identity follows because
a skew-real $3\times3$ matrix and an imaginary real-symmetric
trace-free matrix are both anti-Hermitian and trace-free; their
commutators close and their dimensions add to $3+5=8$.
\end{proof}

\begin{remark}[Candidate chamber-block complement retained at its actual evidence level]
\label{p1:rem:v365-chamber-register-candidate}
For a \emph{specified} 25-block chamber partition $U$ and 25
transporter-register matrix $B$, the proposed marked identities are
\[
\begin{aligned}
 T^{\mathsf T}U&=J_{120\times25}-B,\\
 TB&=J_{600\times25}-(I_{600}+A_{\rm chamber})U,\\
 U^{\mathsf T}A_{\rm chamber}U&=6[(J_5-I_5)\otimes(J_5-I_5)].
\end{aligned}
\]
These require a separate source-matched exact construction of the
marked $U$ and its relation to the already defined $B$; the
$600$-cell incidence theorem by itself does \emph{not} prove them.
The existing $\operatorname{rank}B=17$ and its Petersen
coefficients retain their separate original proofs.  The local
sixfold $C_6$ vertex representation, its centre-even quadratic
comparison, and the $8+16$, $6+12+6$, and $8+8+8$ decompositions
must not be identified by matching dimensions alone.
\end{remark}


\section{The exact twenty-five-block chamber--Hurwitz bridge}
\label{p1:sec:v366-chamber-hurwitz-completion}
This is an extension of, not a replacement for, the marked finite
$2I$ transporter/Petersen construction and the exact incidence theorem
\ref{p1:thm:v365-600-cell-incidence}.  The $120$ group \emph{vertices},
$600$ tetrahedral \emph{chambers}, and $25$ overlapping group-element
\emph{registers} have distinct types.  Fix the already marked $G=2I$,
$H=2T$, and the ordered quotient $G/H$ of size five.  The labels of the
five cosets and the marked left/right representative frames are part of
the packet; no physical color, spin or decay-channel interpretation is
added by this section.

\begin{theorem}[Constructed chamber partition and transporter complement]
\label{p1:thm:v366-chamber-transporter-complement}
Let $\mathscr F$ be the $600$ tetrahedral facets of the $600$-cell,
$T\in\{0,1\}^{600\times120}$ its chamber--vertex incidence, and
$B_{g,(i,j)}=\mathbf1_{\{g\in a_iHb_j^{-1}\}}$ the $25$ transporter
indicators in the fixed coset frame.  There exists a \emph{constructed},
not merely dimension-counted, partition
\[
 \mathscr F=\bigsqcup_{i,j=1}^{5}\mathscr F_{ij},\qquad
 |\mathscr F_{ij}|=24,\qquad
 U_{f,(i,j)}=\mathbf1_{\{f\in\mathscr F_{ij}\}},
\]
with
\begin{equation}\label{p1:eq:v366-TU}
 \boxed{T^{\mathsf T}U=J_{120\times25}-B},\qquad
 B^{\mathsf T}B=24I_{25}+6R,\quad
 R=(J_5-I_5)\otimes(J_5-I_5).
\end{equation}
In each block the $24$ tetrahedra use each of the $96$ vertices
outside $a_iHb_j^{-1}$ exactly once and use none of the $24$ vertices
inside it.  The $\mathscr F_{ij}$ are disjoint \emph{chamber} sets;
$a_iHb_j^{-1}$ are overlapping \emph{vertex} sets.
\end{theorem}
\begin{proof}
A tetrahedron stabilizer in the $G_L\times G_R$ action
$f\mapsto afb^{-1}$ has order $24$, and its left and right
projections are conjugate binary-tetrahedral subgroups.  Conjugate both
projections to the fixed $H$.  With $f_0$ the aligned chamber, define
\[
 \mathscr F_{ij}=\{a_i h_1 f_0(h_2^{-1}b_j^{-1}):h_1,h_2\in H\}.
\]
Here the expression means the vertexwise two-sided quaternion action.
Its stabilizer has order $24$, so each orbit has $24^2/24=24$
chambers.  The $25$ such orbits are disjoint and exhaust $600$.
For the standard exact Hurwitz--icosahedral coordinate bank, the
vertexwise orbit enumeration shows that the $24$ chambers in
$\mathscr F_{ij}$ miss exactly the transporter $a_iHb_j^{-1}$ and
cover its $96$-element complement once.  This finite identity can
be checked without floating-point matching using the integer-pair
$\mathbb Q(\sqrt5)$ coordinate and quaternion-multiplication algorithm
in \texttt{Proofs/Paper00-HurwitzHadronChamberExact-v1\_1.py}; it
verifies every one of the $3000$ incidence entries.  The transporter
Gram then follows from the already established marked $G/H$ action
(or directly from the same exact bank).
\end{proof}

\begin{theorem}[One face step realizes the transporter and rook-complement rules]
\label{p1:thm:v366-face-register-identity}
Let $A\in\{0,1\}^{600\times600}$ join tetrahedral chambers
sharing a triangular facet, so $A\mathbf1=4\mathbf1$.  Then
\begin{equation}\label{p1:eq:v366-face-identity}
 \boxed{TB=J_{600\times25}-(I_{600}+A)U},\qquad
 \boxed{U^{\mathsf T}AU=B^{\mathsf T}B-24I_{25}=6R}.
\end{equation}
The induced $25$-vertex register transition graph is the complement
of the five-by-five rook graph; its spectrum is
$\{16^{(1)},1^{(16)},(-4)^{(8)}\}$.  There are exactly six
face-sharing chamber pairs from block $(i,j)$ to block $(k,l)$
when $i\ne k$ and $j\ne l$, and zero otherwise.  These are geometric
adjacencies, \emph{not} a physical Hamiltonian or decay probabilities.
\end{theorem}
\begin{proof}
For each chamber $f$, precisely five transporter sets contain none
of its four vertices: the register of $f$ and those of its four
face-neighbors.  Each of the other twenty registers contains exactly
one vertex of $f$.  This proves the first entrywise identity.  Multiply
it by $U^{\mathsf T}$, use $U^{\mathsf T}U=24I$ and
$U^{\mathsf T}J=24J$ and invoke the preceding complement to obtain
the second.  As $R=(J_5-I_5)^{\otimes2}$, its eigenvalues are
$16$ once, $1$ sixteen times and $-4$ eight times.  The finite
chamber-face incidence and register identities have also been
independently checked in integer arithmetic by the exact certificate.
\end{proof}

\begin{theorem}[Petersen and register face-walk sectors]
\label{p1:thm:v366-petersen-face-walk}
Let $\Pi_{15}$ denote the \emph{existing} rank-$15$ Petersen
coefficient projector on $\mathbb R[G]$ (not the intrinsic momentum
module $\mathcal H_1$), and put $Q=I_5-J_5/5$.  Then
\[
 \Pi_{15}=\frac1{30}(T^{\mathsf T}U)
   (Q\otimes Q)(T^{\mathsf T}U)^{\mathsf T}
   -\frac1{120}\chi_4\chi_4^{\mathsf T},
\quad (A^2+A-5I)T\Pi_{15}=0.
\]
The $A$-cyclic space from $T\Pi_{15}$ has dimension $30$ and,
with orthonormal partner vectors, is represented by
\[
 I_{15}\otimes\begin{pmatrix}1&\sqrt3\\\sqrt3&-2\end{pmatrix},
 \qquad \lambda_{\pm}=\frac{-1\pm\sqrt{21}}2.
\]
Its normalized source resolvent is $(z+2)/(z^2+z-5)$.
For the full register source $U$,
\[
 (A^2+A-5I)U(Q\otimes Q)=0,\qquad
 \dim\operatorname{span}\{U,AU,A^2U,\ldots\}=41=1+8+2(16).
\]
An optional \emph{trial} register potential $V(M)=\sum M_{ij}
P_{\mathscr F_{ij}}$ can generate all $300$ antipodally even chamber
modes from $U$ for an explicit integer weight choice.  The remaining
$300$ are antipodally odd.  None of these dimensions or resolvent
roots is a hadron mass, width or observed channel count.
\end{theorem}
\begin{proof}
Substitute \eqref{p1:eq:v366-TU} into the earlier Petersen projector
identity.  Integer multiplication of the $600\times120$ and
$600\times25$ matrices gives the two displayed annihilating
polynomials exactly.  The two successive images in each centered
component are independent, giving the $2\times2$ matrix; the
uniform and eight additive register contrasts account for the
remaining $1+8$ directions.  The cyclic ranks and the separate
finite-field full-rank witness for the \emph{trial} potential are
recorded in the cumulatively replayed certificate
\texttt{Proofs/Paper00-HadronChamberSplitCertificate-v1\_7.py}.
\end{proof}

\begin{theorem}[Canonical dual Hurwitz facet and rooted radial refinement]
\label{p1:thm:v366-dual-hurwitz-refinement}
For every marked $(i,j)$ block, the normalized barycenters $n_f$ of
its $24$ chambers form the opposite Hurwitz 24-cell
$a_i(2O\setminus2T)b_j^{-1}$, dual to the transporter
$C_{ij}=a_iHb_j^{-1}$.  The signed incidence
\[
 \ell_f(q)=\sqrt2\langle n_f,q\rangle\in\{-1,0,1\}
 \qquad(q\in C_{ij})
\]
has layers $6_++12_0+6_-$.  If a chosen root $v\in C_{ij}$ has
its \emph{rooted nearest eight} $N_8(v)$ and $t=\ell_f(v)$,
then, across these three layers,
\[
 \#(N_8(v)\cap C_{+,0,-})=(2+2t,4,2-2t),\qquad
 \#((C_{ij}\setminus N_8(v))\cap C_{+,0,-})=(4-2t,8,4+2t).
\]
The rooted near/far complement further has cardinalities $7+9$.
These are simultaneous refinements on the same 24-point torsor.
The nearest-eight subset is \emph{not} the normal subgroup $Q_8$,
and the three normal-coset octets $8+8+8$ are yet another marking.
\end{theorem}
\begin{proof}
Normalize the sum of the four 600-cell vertices in each chamber.
The resulting $24$ quaternionic normals are exactly the $2O\setminus
2T$ coset in the aligned reference block; multiplication by the marked
left/right representatives transports this equality to all blocks.
The inner products with the complementary Hurwitz vertices are
$\pm1/\sqrt2,0$, with multiplicities $6,6,12$.
The three root cases $t=1,0,-1$ are an exhaustive orbit enumeration,
verified by exact integer-pair quaternion coordinates in the
supplementary certificate.  The $8+7+9$ radial statement itself was
proved earlier and is only being intersected with the facet layers.
\end{proof}

\begin{proposition}[Local Hurwitz Fourier filters and their blind sectors]
\label{p1:prop:v366-hurwitz-fourier-filters}
On one matched dual pair let
$J(f,q)=1$, $L(f,q)=\ell_f(q)$ and $E(f,q)=2\ell_f(q)^2-1$.
Their nonzero singular-value lists are respectively
$\{24^{(1)}\}$, $\{(6\sqrt2)^{(4)}\}$ and
$\{8^{(9)}\}$, with pairwise orthogonal images.  Consequently the
local regular coefficient space splits as
\[
 \boxed{24=1+4+9+10=1+4+9+8+2}.
\]
Under the local $2T$ Fourier decomposition the $8$ consists of the
two twisted doublet coefficient modules; the remaining $2$ consists
of the two nontrivial $2T/Q_8\cong C_3$ scalar-character lines.
The signed $L/(6\sqrt2)$ is a rank-four partial isometry.  Its
normalization is an abstract finite-geometric equality, not an
identity with an independently defined QCD coupling.
\end{proposition}
\begin{proof}
In the aligned reference 24-cell, the orthogonality relations
follow by summing $1,\ell$, and $2\ell^2-1$ over its $24$ vertices
and $24$ dual facets.  The real regular representation of $2T$
has irreducible coefficient dimensions
$1^2+1^2+1^2+2^2+2^2+2^2+3^2=24$.
The respective scalar, fundamental-doublet and adjoint coefficient
images account for $1+4+9$; center parity and the $C_3$ quotient
separate the ten-dimensional orthogonal complement into $8+2$.
The finite projectors and singular values are replayed in the
\texttt{v1\_7} exact-certificate arithmetic.
\end{proof}

\begin{theorem}[Signed and cubic vertex-to-chamber Fourier channels]
\label{p1:thm:v366-signed-cubic-chamber}
Assemble the above matched signed layer maps into
$\mathcal L\in\mathbb R^{600\times120}$.  Then
\[
 \mathcal L^{\mathsf T}U=0,\qquad
 \operatorname{Spec}(\mathcal L^{\mathsf T}\mathcal L)
 =\{360^{(8)},120^{(36)},0^{(76)}\},\qquad
 \operatorname{rank}\mathcal L=44.
\]
For all 600 chamber normals and all 120 vertices define
\[
 K_3(f,q)=8\langle n_f,q\rangle^3-4\langle n_f,q\rangle.
\]
With $P_{4'}$ the coefficient projector onto the odd
four-dimensional $2I$ representation, one has
\[
 \boxed{K_3^{\mathsf T}K_3=4500P_{4'}},\qquad
 \operatorname{rank}K_3=16,\qquad
 (A^2-3A-I)K_3=0.
\]
For each $f\in\mathscr F_{ij}$, $K_3(f,q)=0$ precisely for
$q\in C_{ij}$; it is nonzero at the other 96 vertices.
Its $A$-cyclic space has dimension $32$, with eigenvalues
$(3\pm\sqrt{13})/2$.  A single matched incidence has values
$\ell\in\{-1,0,1\}$ and $\ell^3=\ell$; higher polynomials in
that \emph{one} incidence cannot recover the cubic blind sector.
\end{theorem}
\begin{proof}
Antipodal inversion interchanges the two signed six-layers,
forcing $\mathcal L^{\mathsf T}U=0$.  Explicit exact-field sums
of their $120$ columns give the stated Gram eigenvalues.
The degree-three Chebyshev polynomial $8x^3-4x$ is the character
of $\operatorname{Sym}^3(\mathbb C^2)$; character orthogonality
on $2I$ gives the $16$-dimensional coefficient projector.
The geometric normalization $4500$ and the stated rowwise zero
pattern follow from exact chamber-normal/vertex inner products.
The face-walk polynomial and cyclic rank are exact matrix identities
certified by the full integer/$\mathbb Q(\sqrt5)$ replay.  All
statements are finite geometry, not observable resonances.
\end{proof}

\begin{theorem}[Two-sided representation ledger of the chamber module]
\label{p1:thm:v366-600-chamber-representation-ledger}
Under the vertexwise left/right action,
\[
 \mathbb C^{\mathscr F}\cong
 \operatorname{Ind}_{\Gamma(2T)}^{2I_L\times2I_R}{\bf1},
\qquad
 \boxed{600=(126+148)_{\rm diagonal}+(174+152)_{\rm off\text{-}diagonal}}
\]
where the first entries in each parenthesis are center-even and the
second center-odd.  The entire diagonal sector has dimension
\[
 \boxed{274=1+18+32+75+8+108+32}.
\]
It is generated by scalar, quadratic, matched-signed and cubic
vertex-source maps under pure face-sharing propagation; the
remaining off-diagonal $326$ directions have distinct two-sided
representation types.  These dimensions classify the finite
\emph{representation} and do not count hadron species.
\end{theorem}
\begin{proof}
The two-sided chamber stabilizer is a graph of a $24$-element
$2T$ subgroup.  Frobenius reciprocity assigns multiplicity
$\dim\operatorname{Hom}_{2T}(\rho\!\mid_{2T},\sigma\!\mid_{2T})$
to each $2I_L\times2I_R$ pair.  The restriction character table,
verified against the finite conjugacy bank, yields the stated
$600$-dimensional ledger.  The cyclic-space dimensions for the
four indicated source types add to the $274$ diagonal modes.
\end{proof}

\begin{proposition}[Pentagonal complement and cross-core multiplication filter]
\label{p1:prop:v366-cross-core-pentagonal-kernel}
The union of the five conjugate binary-tetrahedral cores has 72
elements; the other 48 are the binary lifts of the 24 five-cycles
in $A_5$.  For distinct cores $H_i,H_j$, with
$H_i\cap H_j\cong C_6$, put
\[
 \mathcal K_{ij}(h,k)=\mathbf1_{\{hk\notin\cup_m H_m\}},
 \qquad(h,k)\in H_i\times H_j.
\]
This $24\times24$ zero-one matrix has 216 ones, rank three,
and singular values $12,6,6$; its four row/column classes have
size six.  Its local relative-core cardinality refinement is
$24=6_{\rm shared}+6_{\rm incompatible}+12_{\rm pentagonal}$.
The number $216/576=3/8$ is uniform \emph{incidence}, not a
branching fraction.  The kernel is even under each antipodal flip
and kills both local center-odd and nontrivial scalar $C_3$
character coefficient spaces.  Character twists preserve its rank
and do not create an odd-spinor channel.  In particular the local
abelianization $2T_{\rm ab}\cong C_3$ does not extend to a nontrivial
scalar character of perfect $2I$.
\end{proposition}
\begin{proof}
Pass to $2I/\{\pm1\}\cong A_5$.  Stabilizers of five letters
contain the identity, 3-cycles and double transpositions with
respective cover-membership counts $5,2,1$; no five-cycle belongs
to a point stabilizer.  The group conjugacy and multiplication
enumeration gives the four $6\times6$ row/column blocks, one zero
class and the active $3\times3$ complement-of-diagonal pattern.
The singular values of that block matrix are $12,6,6$.
Antipodal invariance and $Q_8$-coset summation yield the stated
annihilators.  The exact permutation-group and quaternionic checks
appear in the $\texttt{v1\_9}$ reference certificate.
\end{proof}

\begin{proposition}[Induced Hurwitz quintet and finite spinorial eligibility]
\label{p1:prop:v366-twisted-quintet-spinorial-branch}
Retain the \emph{previously established} induction
$\operatorname{Ind}_{2T}^{2I}\chi\cong{\bf5}$ for a nontrivial
$C_3$ character.  Tensoring with the fundamental odd doublet of
$2I$ gives
\[
 \boxed{{\bf5}\otimes{\bf2}={\bf4}'\oplus{\bf6}},\qquad
 P_{4'}=(2I-X)/5,\quad P_6=(X+3I)/5,\quad
 X=2\mathbf J\cdot\mathbf S.
\]
The projectors act on the \emph{abstract} spin-$2\otimes1/2$
carrier; their ranks are 4 and 6.  On unprojected spin factors one
also has
\[
 {\bf5}\otimes({\bf2}\otimes{\bf2}^*)
 =\mathbf3\oplus\mathbf3'\oplus\mathbf4\oplus2\mathbf5,
\quad
 {\bf5}\otimes{\bf2}^{\otimes3}
 =\mathbf2\oplus\mathbf2'\oplus3\mathbf4'\oplus4\mathbf6.
\]
Neither abstract finite-group spin nor these dimensions make a
physical baryon or meson identification without the actual
rotation, gauge, flavor and Pauli projections.
\end{proposition}
\begin{proof}
Restrict the $SU(2)$ Clebsch--Gordan identities
$2\otimes\tfrac12=\tfrac32\oplus\tfrac52$ and the corresponding
raw two-/three-spinor products to the marked $2I$.  Their exact
character decompositions give the stated finite representations.
$X$ has eigenvalues $-3,2$, giving the projector formula.
\end{proof}


\begin{proposition}[Local twelve-by-two Hurwitz face-transition permutation]
\label{p1:prop:v366-local-face-permutation}
Within a fixed $24$-chamber block $\mathscr F_{ij}$, each tetrahedron
has four face-neighbors, and their destination registers determine a
bijection between the four remaining row labels and the four
remaining column labels.  There are precisely twelve such
permutations, each occurring at two antipodal chambers:
\[24=12\cdot2.\]
After fixing an orientation chart, the twelve form one parity class
in $S_4$, isomorphic as a torsor to $A_4$.  This finite destination
permutation is neither physical spacetime parity nor a count of four
final-state particles.
\end{proposition}
\begin{proof}
Use the previous chamber partition and list the four triangles of
each of its 24 tetrahedra.  The face-sharing destination labels
are four distinct complementary row/column pairings.
Antipodal inversion leaves the pairing unchanged and exchanges
the two lifts.  The exact permutation list has twelve elements
with fixed parity in the selected chart, as checked combinatorially
in the source-matched $\texttt{v1\_7}$ verifier.
\end{proof}

\begin{proposition}[Sixteen direct inter-core contrasts and fifteen mirror-active contrasts]
\label{p1:prop:v366-register-mirror-contrast}
For a trial register potential
$V(M)=\sum_{i,j}M_{ij}P_{\mathscr F_{ij}}$ with
$M\in\mathbb R^{5\times5}$, let
$K(M)=\Phi^{\dagger}V(M)\Phi$ be its compression to the
\emph{pre-existing} normalized rank-$15$ Petersen coefficient
isometry $\Phi$.  Let $K_{\mathrm{inter}}(M)$ retain entries
between distinct local three-dimensional core summands and put
$D=\Pi_3-\Pi_{3'}$ for the existing Petersen mirror contrast.
With $Q=I_5-J_5/5$, the exact equivalences are
\[
 K_{\mathrm{inter}}(M)=0\quad\Longleftrightarrow\quad
 QMQ=0\quad\Longleftrightarrow\quad M_{ij}=a_i+b_j,
\]
\[
 [K(M),D]=0\quad\Longleftrightarrow\quad QMQ=\alpha Q.
\]
Thus the first image has dimension $16$ (additive kernel
$9$), and the second commutator image dimension $15$.
For the scalar centered register tensor,
$K(Q)=-I_{15}/15+P_4/5$, where $P_4$ is the old
rank-four Petersen Hodge projector.  The statement concerns
direct compression of \emph{trial} chamber weights; virtual
transitions through the extra fifteen chamber partner directions
are not excluded by an additive first-order kernel.
\end{proposition}
\begin{proof}
Express the Petersen edge isometry in the already fixed
five-core pair chart.  The twenty-five coordinate matrices
$\Phi^{\dagger}P_{\mathscr F_{ij}}\Phi$ have off-core
coefficient-map rank sixteen over $\mathbb Q$, and their
nullspace is the nine-dimensional row-plus-column subspace.
Commuting these matrices with $D$ gives exact coefficient-map
rank fifteen and leaves the one-dimensional scalar centered
line in addition to the nine-dimensional additive kernel.
These ranks and $K(Q)$ are rational-matrix identities checked
by exact elimination in the independent v1.3 segment of the
cumulative \texttt{v1\_7} certificate.
\end{proof}

\begin{remark}[Geometry-to-observable non-implication]
\label{p1:rem:v366-chamber-physical-firewall}
The finite chamber count $600$, register count $25$, vertex count
$120$, the two $15$-dimensional mathematical constructions, all
adjacency poles and projector norm ratios above are \emph{not}
measured hadronic particles, decay channels, masses or widths.
The action-owned physical source-intertwiner, its gauge/rotation
qualification and the full Feshbach leakage are separate
requirements of the APBI hadron-presentation definition in Paper~VII-A.
The chamber-register identities are established \emph{only for their
three displayed finite identities} by Theorems
\ref{p1:thm:v366-chamber-transporter-complement} and
\ref{p1:thm:v366-face-register-identity}; its physical caveat remains.
\end{remark}

\section{Exceptional-root triality resolution of the existing Hurwitz shell}
\label{p1:sec:v370-e8-hurwitz-resolved}
This section is a \emph{derived mathematical comparison} of objects already
specified in Theorem~\ref{p1:thm:exroot-explicit-e876} and in the marked
$2T\subset2I$ finite geometry.  It does not declare $E_8$ to be a physical
gauge group, identify its weight sectors with particles or generations, or
introduce a source, action, mass or decay operator.  All multiplicities below
refer to exact root sets or existing finite incidence operators.

Let $H=2T\subset U=2I$ be the specified Hurwitz core, and write
\[
 \Phi_8=\iota(U)\sqcup\iota(\varphi U),\qquad
 D_0=\iota(H),\quad D_1=\iota(\varphi H),\qquad
 \mathscr X=\Phi_8\setminus(D_0\cup D_1).
\]
Here the Clifford coefficient inner product, the golden action and the
orthogonal root reflections are precisely those fixed in the exceptional-root
appendix.  The labels $0,1$ distinguish root sheets and are not fields or
physical bands.

\begin{theorem}[Two orthogonal Hurwitz $D_4$ systems and three triality sectors]
\label{p1:thm:v370-e8-d4-triality}
The subsets $D_0,D_1$ are orthogonal 24-root systems of type $D_4$.
Under their commuting reflection groups, the 192-root complement has three
orbits $X_0,X_1,X_2$ of size 64.  Each is the Cartesian product of one
8-weight orbit in each $D_4$ factor, and meets each golden sheet in 32 roots.
Consequently the complex-root branching has the representation type
\[
 \mathfrak e_8\downarrow_{\mathfrak{so}(8)\oplus\mathfrak{so}(8)}
 \cong(\mathbf{28},\mathbf1)\oplus(\mathbf1,\mathbf{28})
 \oplus(\mathbf8_v,\mathbf8_v)
 \oplus(\mathbf8_s,\mathbf8_s)
 \oplus(\mathbf8_c,\mathbf8_c),
 \qquad248=28+28+64+64+64.
\]
There are 2240 unordered distinct root triples with sum zero; 32 lie
within each $D_i$, 192 lie in each of the six classes $(D_i,X_a,X_a)$,
and 1024 lie in $(X_0,X_1,X_2)$.  The last class factors as the two
32-term $D_4$ triality weight supports, $1024=32\cdot32$.
These are root-addition supports; no signed physical interaction constant is
assigned by their counts.
\end{theorem}
\begin{proof}
The supplied 120 unit quaternions are the eight coordinate units, sixteen
half-coordinate units and ninety-six even-permuted golden units.  Encode a
coordinate $(a+b\varphi)/2$ as the integral pair $(a,b)$, and use the ordered
eight Clifford coefficient coordinates of
Proposition~\ref{p1:prop:exroot-golden-action}.  The two Hurwitz core root
banks occupy orthogonal coefficient four-spaces, and each is the standard
24-vector $D_4$ root bank.  For roots represented in the doubled integral
coordinates, their reflections are
\[
 s_\alpha(r)=r-\frac{r\cdot\alpha}{2}\alpha
 \quad\text{when}\quad\alpha\cdot\alpha=4.
\]
Exact reflection closure of the remaining 192 roots gives three orbits of
64; projecting each orbit to the two coefficient four-spaces gives two
eight-element sets whose full Cartesian product is the original orbit.
Counting triples by the exact equation $r+s+t=0$ gives
$32+32+6(192)+1024=2240$.  In each mixed triple the two four-space
components must independently sum to zero, and each admits exactly 32
fixed-slot triality weight triples.  These finite counts and the generating
reflection formulas are replayed with integer arithmetic in
Proposition~\ref{p1c:prop:v370-e8-finite-certificate}; no approximate root
matching or physically selected basis is used.
\end{proof}

\begin{theorem}[Equitable 600-cell resolution and native chamber norms]
\label{p1:thm:v370-e8-shell-chamber}
Intersect the $X_a$ with the ordinary golden-lift sheet and forget that
sheet label, obtaining three subsets $C_a\subset U$.  Then
\[
 U=H\sqcup C_0\sqcup C_1\sqcup C_2,
 \qquad(|H|,|C_0|,|C_1|,|C_2|)=(24,32,32,32).
\]
The four nontrivial left $H$-cosets each meet every $C_a$ in eight points.
For the \emph{existing} 600-cell nearest-neighbor adjacency
$A_{600}$, the indicated partition is equitable, with quotient
\[
 Q=\begin{pmatrix}
 0&4&4&4\\3&3&3&3\\3&3&3&3\\3&3&3&3
 \end{pmatrix},\qquad
 \operatorname{spec}Q=\{12,-3,0,0\}.
\]
Let $\zeta^3=1\ne\zeta$ and define normalized vertex functions
\[
 \psi_{\zeta}(g)=\frac1{\sqrt{96}}
 \begin{cases}
 0,&g\in H,\\1,&g\in C_0,\\\zeta,&g\in C_1,\\
 \zeta^2,&g\in C_2.
 \end{cases}
\]
Then $\psi_\zeta$ and $\psi_{\bar\zeta}$ are center-even, orthonormal,
are nontrivial $H/Q_8\cong C_3$ character modes, and satisfy
\[
 A_{600}\psi_\zeta=0,\qquad
 \Pi_{25}\psi_\zeta=\psi_\zeta,\qquad
 \mathscr P_4\psi_\zeta=0,\qquad
 T_u\psi_\zeta=\frac15\psi_\zeta.
\]
Here $\Pi_{25}$ is the established 25-dimensional twisted-quintet
coefficient projector, $\mathscr P_4$ the separate 16-dimensional
central-character projector, and $T_u$ the established unsigned transfer;
none is the fifteen-dimensional Petersen projector.  For the separately
defined $600\times120$ chamber--vertex incidence $T$,
\[
 T^\dagger T\big|_{\operatorname{span}\{1_H,1_{C_0},1_{C_1},1_{C_2}\}}
 =20I+5A_{600}\big|_{\operatorname{span}\{1_H,1_{C_0},1_{C_1},1_{C_2}\}},
\]
with eigenvalues $80,5,20,20$.  In particular,
$\|T_u\psi_\zeta\|^2=1/25$ and $\|T\psi_\zeta\|^2=20$.
\end{theorem}
\begin{proof}
The orbit-bank construction in the preceding theorem fixes the
$24+32+32+32$ partition without choosing a new particle labeling.
The exact quaternion product gives the five left cosets and their
$24+4(8+8+8)$ distribution.  Counting the 12 nearest neighbors of
each of the 120 roots against all four classes gives \emph{each row} of
$Q$, not only the average degrees.  This yields the displayed quotient
and annihilates the two nontrivial Fourier combinations.
Left multiplication by $H$ permutes $(C_0,C_1,C_2)$ through a cyclic
order-three image, whose kernel is precisely the already specified $Q_8$.
Thus the Fourier functions carry the nontrivial inducing characters
of Theorem~\ref{p1:thm:v344-induced-hurwitz-quintet}.  The existing
character-projector kernels of
Theorems~\ref{p1:thm:v340-incidence-character-projectors} and
\ref{p1:thm:v344-twisted-25-gram}, independently checked entrywise on
all 60 antipodal classes, give the $\Pi_{25}$ and $\mathscr P_4$
identities.  The last two norm identities follow from the already fixed
$T_u\Pi_{25}=\Pi_{25}/5$ and
Theorem~\ref{p1:thm:v365-600-cell-incidence}.  All finite orbit,
coset, adjacency and projector checks are recorded in the cited
Paper~I-D certificate.
\end{proof}

\begin{proposition}[Existing $E_6$ section and exact 27-weight cubic support]
\label{p1:prop:v370-e6-cubic-support}
In the \emph{already chosen} section
$\Phi_6=\Phi_8\cap1^\perp\cap(i+j+k)^\perp$,
the perpendicular $E_8$ roots form an $A_2$ system of six roots.
The remaining 162 roots project to six distinct $A_2$ weights, with
27 roots at each weight.  Each of the two triples of $A_2$ weights
has the same 27 projected $E_6$ weights, up to overall sign for the
opposite triple.  The $D_4$ subsystem $D_1\subset\Phi_6$
splits each 27-weight set into three fixed weights and three octets:
\[
 27=1+1+1+8+8+8=3+8_v+8_s+8_c.
\]
Among the 27 distinct projected weights there are exactly 45 unordered
triples of distinct weights summing to zero.  Their class support is
\[
 \boxed{45=1_{3\text{ scalar}}+
 4_{\text{scalar},8_v,8_v}+
 4_{\text{scalar},8_s,8_s}+
 4_{\text{scalar},8_c,8_c}+
 32_{8_v,8_s,8_c}.}
\]
The last 32 have the root-support pattern of the previously established
$D_4$ weight-triality tensor; the 45 triples are the \emph{support},
not a coefficient-matched physical $E_6$ cubic or a new source.
\end{proposition}
\begin{proof}
The exact normal coordinate pair of a root in the specified section is
$(r_0,r_2+r_4+r_6)$ in the doubled integral coefficient bank.  The
six nonzero normal classes have 27 members each; the orthogonal roots
are six $A_2$ roots.  Orthogonal projection to the $E_6$ subspace
identifies the 27 weights across each triple, giving the
$3+8+8+8$ $D_1$-orbit sizes.  Enumerating unordered distinct
projected-weight triples with zero sum yields 45 and the displayed
$1+3(4)+32$ decomposition; each weight participates in exactly
five triples.  Both the exact normal-projection and Weyl-orbit tests
are replayed in Proposition~\ref{p1c:prop:v370-e8-finite-certificate}.
\end{proof}

\begin{remark}[Root-addition comparisons and the physical typing boundary]
\label{p1:rem:v370-e8-readout-boundary}
For a simply laced root basis, $[E_\alpha,E_\beta]$ can be nonzero
when $\alpha+\beta$ is a root.  This statement explains the exact
root-triple supports above, but is neither a Born--Hopf decay channel
nor an action-owned amplitude.  The $E_8$ three-sector representation
type is fixed, while coefficient-by-coefficient matching of its signed
Lie bracket with the existing rooted Clifford/Fock phases is \emph{not}
claimed.  The six mathematical $A_2$ weights are not identified with
Paper V's physical $SU(3)_c$, and the three octets are not fermion
generations.  The auxiliary golden-sheet label exchange has symmetric
three-sector compression $(J_3-I_3)/2$, equal to the \emph{even}
$C_3$ return $(X_G+X_G^\dagger)/2$ after an abstract cyclic-basis
identification; it supplies no odd/phase-sensitive source term.
Furthermore the two-point shadow of the 1840 distinct projected
root triples is already in $\mathbb Q[A_{600}]$, as is its
pair-multiplicity matrix.  This is a redundancy statement for
\emph{two-point} observables, not for signed cubic coefficients.
Any physical use must separately satisfy the established action-owned
source intertwiner, same-action normalization, center/rotation gates,
pole/residue tests and decay thresholds in Papers V and VII.  No
current physical prediction or gap status is promoted by this section.
\end{remark}


\clearpage


% BEGIN BERGER/600-CELL MATHEMATICAL EXPORT v1.17
\section{Marked icosahedral quadrupoles as potential shape-response data}
\label{bh:berger:v117:p1b:quadrupoles}
This is a comparison on an \emph{occurring marked} binary-icosahedral
packet, not a physical position-space deformation selector.
\begin{proposition}[Ten-axis cancellation and first cubic shear invariant]
\label{bh:berger:v117:p1b:cubic}
For the ten projective unit Hopf axes $u_e$ in the established
Petersen construction put $Q_e=u_eu_e^{\mathsf T}-I_3/3$. Then
\[
 \sum_{e=1}^{10}Q_e=0.
\]
For $B\in\operatorname{Sym}^2_0(\mathbb R^3)$ define
\[
 \mathcal I_{\rm ico}(B)=\frac1{10}\sum_{e=1}^{10}
                 (u_e^{\mathsf T}Bu_e)^3
                -\frac8{105}\operatorname{tr}(B^3).
\]
It is invariant under the existing $A_5$-marked axis action, but
not under arbitrary rotations. On $B=s(nn^{\mathsf T}-I_3/3)$,
\[
 \mathcal I_{\rm ico}(B)=\frac{16}{231}s^3T_6(n),
\]
where $T_6$ is the existing normalised collective sixth-order
angular comparison in the downstream spectral/cosmology sector.
\end{proposition}
\begin{proof}
The second moment of the ten projective axes is $10I_3/3$ by
irreducibility of the real icosahedral vector representation and its
trace.  The isotropic cubic average of a traceless matrix is
$(8/105)\operatorname{tr}(B^3)$; subtracting it isolates the marked
cubic invariant.  Expanding the Berger $Q_n$, the fourth-order
angular moments cancel and the remaining axis sixth moment is
exactly the established $T_6$ normalization.
\end{proof}
\begin{remark}[No automatic physical locking]
\label{bh:berger:v117:p1b:scope}
A complete equal-weight ten-axis packet has no homogeneous linear
quadrupole force; the twelve signed Petersen pentagons are mathematical
source comparisons, not twelve particles or selected physical vacua.
Full Born midpoint holonomy precludes a nonzero globally parallel
traceless spin-two marking without an actual transported source history.
\end{remark}
% END BERGER/600-CELL MATHEMATICAL EXPORT v1.17




% BEGIN FRAMEWORK v3.068 GEOMETRY BRIDGE SOURCE INTEGRATION
\section{Rooted register and Petersen source transport: exact finite comparison}
\noindent\emph{Scope.} The statements below retain their declared finite geometric, source-chart, and action-relative hypotheses. They do not assert physical particle decay or unconditional source occupation.

\begin{proposition}[Sharp spectrum of the rooted 25-register comparison Laplacian]
\label{p1b:prop:v368-rooted-register-laplacian}
Let $R=(J_5-I_5)\otimes(J_5-I_5)$ and $L=6(16I_{25}-R)$
be the previously constructed register Laplacian.  For any marked
register vector $e_r$, put $Q_r=I_{25}-e_re_r^{\mathsf T}$ and restrict
$L_r=Q_rLQ_r$ to $e_r^\perp$. Then
\[
\det(\lambda I_{24}-L_r)
=(\lambda-90)^{15}(\lambda-120)^7
 (\lambda^2-114\lambda+432).
\]
Thus the remaining simple eigenvalues are $57\pm3\sqrt{313}$.
The operator and its eigenvalues concern a specified rooted graph
comparison, not a physical charged mass or decay rate.
\end{proposition}
\begin{proof}
The Cartesian five-by-five register adjacency has eigenvalues
$16,1,-4$ of multiplicities $1,16,8$. Its weighted Laplacian has
$0,90,120$ with those multiplicities.  Removing the marked register
restricts the nonconstant $(\pm)$ isotypic summands to dimensions
$15$ and $7$. On the orthogonal two-dimensional permutation-fixed
sector the characteristic factor is $\lambda^2-114\lambda+432$,
as follows by applying $L$ to the normalized equal-sum vectors on
the $8$ same-row/column and $16$ disjoint-index registers.
Multiplying these three mutually orthogonal factors proves the claim.
\end{proof}
\begin{theorem}[Three-face indirect transfer and exact unitary block outcomes]
\label{p1b:thm:v368-chamber-three-step-unitary}
Let $A,U,S,R,P_0,P_a,P_c$ be the full $600$-chamber face operator,
$25$-register isometry $S=U/\sqrt{24}$, and orthogonal register
projectors of Theorem~\ref{p1b:thm:v16-register-resolvent}.  For
register $u$ put $s_u=Se_u$ and let $P_v^{\rm ch}$ project onto all
$24$ chambers of register $v$.  If $v$ shares exactly one of its two
five-coset indices with $u$, then
\[
P_v^{\rm ch}As_u=P_v^{\rm ch}A^2s_u=0,\qquad
P_v^{\rm ch}A^3s_u=3s_v.
\]
For any $v$ with neither index shared, $P_v^{\rm ch}As_u$ has six
entries $1/\sqrt{24}$ and eighteen zero entries, so its norm-square
is $1/4$.  In the explicitly declared mathematical comparison
$H_{\rm cmp}=gA$ with $\theta=gt/\hbar$, the 24-chamber outcome
probabilities satisfy
\[
\begin{aligned}
p_{u\to v}(\theta)&=\|P_v^{\rm ch}e^{-i\theta A}s_u\|^2\\
&=\left|\left(e^{-4i\theta}P_0+e^{i\theta}P_a
+a(\theta)P_c\right)_{vu}\right|^2
+\frac{|b(\theta)|^2}{25}R_{vu},
\end{aligned}
\]
where $\delta=\sqrt{21}\theta/2$ and
\[
a(\theta)=e^{i\theta/2}\left(\cos\delta-
\frac{3i}{2\sqrt{21}}\sin\delta\right),\qquad
b(\theta)=-\frac{5i}{2\sqrt7}e^{i\theta/2}\sin\delta.
\]
The outcome probabilities sum to one for all $\theta$.
For any individual disjoint-index target
$p_{u\to v}=\theta^2/4+O(\theta^4)$; for any individual
same-index target $p_{u\to v}=\theta^6/4+O(\theta^8)$.
No physical time scale, elementary-particle identification, width,
branching fraction, or proper-action gauge/Fock owner is asserted.
\end{theorem}
\begin{proof}
The original source decomposition gives $AS P_0=4SP_0$,
$AS P_a=-SP_a$ and $(A^2+A-5)SP_c=0$.  Thus
$A^2S=20SP_0+5SP_c-AS$ and
$A^3S=40SP_0+5SP_a-5SP_c+6AS$.
The register projectors have pair entries
$(P_0,P_a,P_c)=(1,3,-4)/25$ when one index is shared, and
$(1,-2,1)/25$ when none is shared.  Zero direct face contacts
between same-index blocks give the first two vanishing powers
and the third-power identity.  The six direct face incidences give
the direct norm-square.
Put $T=4(AS-SR/4)/(5\sqrt3)$.  The original resolvent
identities imply $S^*T=0$, $T^*T=P_c$, and on the centred
source/hidden pairs $A$ has matrix
$\begin{psmallmatrix}1/4&5\sqrt3/4\\5\sqrt3/4&-5/4\end{psmallmatrix}$.
Exponentiation yields
$e^{-i\theta A}S=S(e^{-4i\theta}P_0+e^{i\theta}P_a+aP_c)+bT$.
Within each 24-chamber target the $S$ and $T$ terms are
orthogonal because $S^*T=0$.  Direct face-incidence variance gives
$\|P_v^{\rm ch}Te_u\|^2=R_{vu}/25$, proving the probability
formula.  Its sum over $v$ is one by unitarity.  The first nonzero
amplitude order, or an exact Taylor expansion, gives the two
stated leading probabilities.
\end{proof}
\begin{theorem}[The normalized Petersen source is the 60-degree face-walk rotation]
\label{p1b:thm:v368-petersen-chamber-rotation}
Let $A,U,B,T,R$ be the exact 600-chamber, 25-register and
120-vertex incidence operators of Theorems
\ref{p1:thm:v366-chamber-transporter-complement}--
\ref{p1:thm:v366-petersen-face-walk}. Let
$S=U/\sqrt{24}$, $Q=I_5-J_5/5$,
$q=\operatorname{vec}(Q)/2$ and
$P_{15}=Q\otimes Q-qq^{\mathsf T}$. Let
$\mathcal U:\mathbb R^{15}_{\rm Petersen}\to P_{15}\mathbb R^{25}$
be the existing oriented edge isometry of
Theorem~\ref{p1:thm:v339-petersen-sl4-isometry}.
Then
\[
 C_E=\frac25(I_{600}+A)S\mathcal U
\]
is an isometry of rank fifteen, equal up to a global source-orientation
sign to the normalized $T$-lift of the existing Petersen vertex coefficients.
Moreover
\[
 C_E^*S\mathcal U=\tfrac12I_{15},\quad
 C_E^*AC_E=I_{15},\quad (A^2+A-5I)C_E=0.
\]
The two source-first face-walk blocks $M_c$ and $M_P$ are related by
\[
 M_c=\begin{pmatrix}1/4&5\sqrt3/4\\5\sqrt3/4&-5/4\end{pmatrix},\quad
 M_P=\begin{pmatrix}1&\sqrt3\\\sqrt3&-2\end{pmatrix},\quad
 M_P=O^*M_cO,\quad
 O=\begin{pmatrix}1/2&\sqrt3/2\\\sqrt3/2&-1/2\end{pmatrix}.
\]
Thus the source vectors make an exact $60^\circ$ angle in each
centered cyclic fibre, and their self and mixed resolvents are respectively
\[
\begin{aligned}
 (S\mathcal U)^*(z-A)^{-1}S\mathcal U&=\frac{z+5/4}{z^2+z-5}I_{15},\\
 C_E^*(z-A)^{-1}C_E&=\frac{z+2}{z^2+z-5}I_{15},\\
 C_E^*(z-A)^{-1}S\mathcal U&=\frac{z+5}{2(z^2+z-5)}I_{15}.
\end{aligned}
\]
All operators in this theorem are finite geometric comparisons, not
physical decay amplitudes, particle masses or widths.
\end{theorem}
\begin{proof}
For $x\in P_{15}\mathbb R^{25}$, $\mathbf1^*x=0$, so the
old exact identity $TB=J-(I+A)U$ gives
$TF_e/\sqrt5=-(I+A)U\mathcal Ue/\sqrt{150}
=-2(I+A)S\mathcal Ue/5$.
On the centered register factor the old cyclic matrix is $M_c$.
Since $(2/5)(I+M_c)e_1=(1/2,\sqrt3/2)^T$, the
matrix $O$ in the theorem is orthogonal and conjugates $M_c$ to $M_P$.
The isometry, overlap and $A$ expectation follow by direct
multiplication; its minimal polynomial is $\lambda^2+\lambda-5$.
Inverting $zI-M_c$ in the two source vectors gives the
three displayed resolvents.\end{proof}

% END FRAMEWORK v3.068 GEOMETRY BRIDGE SOURCE INTEGRATION

\end{document}
